% Compile with: xelatex
% This template is consumed by Pandoc; the generated standalone source is
% prevention-medicine-past-exams.tex in the same directory.
\documentclass[UTF8,zihao=-4,openany,oneside,fontset=none]{ctexbook}

\usepackage[a4paper,top=2.55cm,bottom=2.35cm,inner=2.55cm,outer=2.15cm,headheight=23pt,headsep=0.55cm,footskip=1.15cm]{geometry}
\usepackage{fontspec}
\usepackage{xeCJK}
\usepackage{amsmath,amssymb,mathtools}
\usepackage{booktabs,longtable,array,calc}
\usepackage{graphicx,float}
\usepackage{enumitem}
\usepackage[dvipsnames]{xcolor}
\usepackage{tikz}
\usepackage[most]{tcolorbox}
\usepackage{titlesec,titletoc}
\usepackage{fancyhdr}
\usepackage{etoolbox}
\usepackage{needspace}
\usepackage{ragged2e}
\usepackage{microtype}
\usepackage{url}
\usepackage{newunicodechar}
\usepackage{hyperref}
\usepackage{bookmark}

\defaultfontfeatures{Ligatures=TeX}
\setmainfont[
  Path=/System/Library/Fonts/Supplemental/,
  UprightFont=Times New Roman.ttf,
  BoldFont=Times New Roman Bold.ttf,
  ItalicFont=Times New Roman Italic.ttf,
  BoldItalicFont=Times New Roman Bold Italic.ttf
]{Times New Roman}
\setsansfont[
  Path=/System/Library/Fonts/Supplemental/,
  UprightFont=Arial.ttf,
  BoldFont=Arial Bold.ttf,
  ItalicFont=Arial Italic.ttf,
  BoldItalicFont=Arial Bold Italic.ttf
]{Arial}
\setmonofont[
  Path=/System/Library/Fonts/Supplemental/,
  UprightFont=Courier New.ttf,
  BoldFont=Courier New Bold.ttf,
  ItalicFont=Courier New Italic.ttf,
  BoldItalicFont=Courier New Bold Italic.ttf,
  Scale=MatchLowercase
]{Courier New}
\setCJKmainfont[
  Path=/System/Library/Fonts/Supplemental/,
  FontIndex=6,
  BoldFont=Songti.ttc,
  BoldFeatures={FontIndex=1}
]{Songti.ttc}
\setCJKsansfont[
  Path=/System/Library/Fonts/,
  FontIndex=0,
  BoldFont=Hiragino Sans GB.ttc,
  BoldFeatures={FontIndex=2}
]{Hiragino Sans GB.ttc}
\setCJKmonofont[Path=/System/Library/Fonts/,FontIndex=0]{Hiragino Sans GB.ttc}
\setCJKfallbackfamilyfont{\CJKrmdefault}[Path=/System/Library/Fonts/Supplemental/]{Arial Unicode.ttf}
\setCJKfallbackfamilyfont{\CJKsfdefault}[Path=/System/Library/Fonts/Supplemental/]{Arial Unicode.ttf}
\setCJKfallbackfamilyfont{\CJKttdefault}[Path=/System/Library/Fonts/Supplemental/]{Arial Unicode.ttf}
\xeCJKsetup{CJKspace=true,PunctStyle=kaiming,CJKmath=true}
\setCJKmathfont[Path=/System/Library/Fonts/Supplemental/,FontIndex=6]{Songti.ttc}
\xeCJKDeclareCharClass{CJK}{
  "2070 -> "209F,
  "2100 -> "218F,
  "2200 -> "22FF,
  "2460 -> "24FF,
  "2600 -> "26FF
}
\newunicodechar{⁻}{\ensuremath{{}^{-}}}

\definecolor{Ink}{HTML}{172033}
\definecolor{TextBody}{HTML}{293246}
\definecolor{TextMuted}{HTML}{687386}
\definecolor{Paper}{HTML}{FBFCFE}
\definecolor{PageTint}{HTML}{F2F6FA}
\definecolor{SubjectBlue}{HTML}{2D5FD5}
\definecolor{CurrentAccent}{HTML}{2D5FD5}
\definecolor{Warning}{HTML}{C85C26}
\definecolor{RuleGray}{HTML}{DCE2EA}

\color{TextBody}
\urlstyle{same}
\hypersetup{
  colorlinks=true,
  hypertexnames=false,
  linkcolor=Ink,
  urlcolor=SubjectBlue,
  citecolor=CurrentAccent,
  pdfborder={0 0 0},
  pdftitle={预防医学大四下往年题（截止22级）},
  pdfsubject={预防医学往年题知识文档汇编},
  pdfkeywords={预防医学, 往年题, 复习资料, 北京大学公共卫生学院}
}

\newcommand{\CurrentSubject}{总览}
\newcommand{\SetSubjectAccent}[2]{%
  \setcounter{chapter}{0}%
  \definecolor{CurrentAccent}{HTML}{#1}%
  \renewcommand{\CurrentSubject}{#2}%
  \hypersetup{urlcolor=CurrentAccent}%
}

\setlength{\parindent}{2em}
\setlength{\parskip}{0.22em}
\setlength{\emergencystretch}{2.5em}
\setlist{leftmargin=2.2em,itemsep=0.2em,topsep=0.45em,parsep=0.1em}
\setlist[itemize,1]{label=\textcolor{CurrentAccent}{\small\textbullet}}
\setlist[enumerate,1]{label=\textcolor{CurrentAccent}{\arabic*.}}
\raggedbottom
\clubpenalty=10000
\widowpenalty=10000
\displaywidowpenalty=10000

\providecommand{\tightlist}{%
  \setlength{\itemsep}{0pt}\setlength{\parskip}{0pt}}
\providecommand{\passthrough}[1]{#1}
\providecommand{\pandocbounded}[1]{#1}
\providecommand{\real}[1]{#1}

\pagestyle{fancy}
\fancyhf{}
\fancyhead[L]{\sffamily\footnotesize\bfseries\color{CurrentAccent}\CurrentSubject}
\fancyhead[R]{\sffamily\footnotesize\color{TextMuted}\nouppercase{\leftmark}}
\fancyfoot[C]{\sffamily\footnotesize\color{TextMuted}\thepage}
\renewcommand{\headrulewidth}{0.35pt}
\renewcommand{\headrule}{\hbox to\headwidth{\color{CurrentAccent!55}\leaders\hrule height \headrulewidth\hfill}}
\renewcommand{\footrulewidth}{0pt}
\fancypagestyle{plain}{%
  \fancyhf{}
  \fancyfoot[C]{\sffamily\footnotesize\color{TextMuted}\thepage}
  \renewcommand{\headrulewidth}{0pt}}

\renewcommand{\chaptermark}[1]{\markboth{#1}{}}
\pretocmd{\part}{\setcounter{chapter}{0}}{}{}

\titleformat{\part}[display]
  {\thispagestyle{empty}\centering\sffamily}
  {\vspace*{2.7cm}\color{CurrentAccent}\fontsize{15}{18}\selectfont\bfseries PART\enspace\Roman{part}}
  {1.1em}
  {\color{Ink}\fontsize{34}{43}\selectfont\bfseries}
  [\vspace{1.4em}\color{CurrentAccent}\rule{0.34\textwidth}{3pt}]
\titlespacing*{\part}{0pt}{0pt}{0pt}

\titleformat{\chapter}[display]
  {\normalfont\sffamily}
  {\color{CurrentAccent}\fontsize{12}{14}\selectfont\bfseries CHAPTER\enspace\thechapter}
  {0.55em}
  {\color{Ink}\fontsize{27}{34}\selectfont\bfseries}
  [\vspace{0.7em}\color{CurrentAccent}\titlerule]
\titleformat{name=\chapter,numberless}[display]
  {\normalfont\sffamily}
  {\color{CurrentAccent}\fontsize{11}{13}\selectfont\bfseries REFERENCE}
  {0.5em}
  {\color{Ink}\fontsize{25}{32}\selectfont\bfseries}
  [\vspace{0.65em}\color{CurrentAccent}\titlerule]
\titlespacing*{\chapter}{0pt}{-0.25cm}{1.15cm}

\titleformat{\section}[hang]
  {\normalfont\sffamily\Large\bfseries\color{Ink}}
  {\colorbox{CurrentAccent}{\makebox[2.4em][c]{\color{white}\small\thesection}}}
  {0.75em}
  {}
\titlespacing*{\section}{0pt}{1.4em}{0.65em}

\titleformat{\subsection}[hang]
  {\normalfont\sffamily\large\bfseries\color{Ink}}
  {\textcolor{CurrentAccent}{\thesubsection}}
  {0.65em}
  {}
\titlespacing*{\subsection}{0pt}{1.15em}{0.5em}

\titleformat{\paragraph}[runin]
  {\normalfont\sffamily\bfseries\color{CurrentAccent}}
  {}{0pt}{}[\quad]

\titlecontents{part}[0em]
  {\addvspace{1.0em}\sffamily\bfseries\color{CurrentAccent}}
  {第\thecontentslabel\enspace 部分\quad}
  {}
  {\hfill\contentspage}
\titlecontents{chapter}[1.2em]
  {\addvspace{0.22em}\sffamily\color{Ink}}
  {\contentslabel{5.8em}}
  {}
  {\titlerule*[0.55pc]{.}\contentspage}

\newlength{\MetaLabelWidth}
\setlength{\MetaLabelWidth}{6.4em}
\newcommand{\InfoLabel}[1]{%
  \par\Needspace{2\baselineskip}\vspace{0.18em}%
  \noindent\hangindent=\dimexpr\MetaLabelWidth+0.55em\relax\hangafter=1%
  \makebox[\MetaLabelWidth][l]{%
    \tcbox[on line,boxsep=0.5pt,left=4pt,right=4pt,top=1.5pt,bottom=1.5pt,
      arc=1.2mm,colback=CurrentAccent!9,colframe=CurrentAccent!22,boxrule=0.35pt]{%
      \sffamily\bfseries\footnotesize\color{CurrentAccent}#1}}%
  \ignorespaces}
\newcommand{\WarningLabel}[1]{%
  \par\Needspace{2\baselineskip}\vspace{0.18em}%
  \noindent\hangindent=\dimexpr\MetaLabelWidth+0.55em\relax\hangafter=1%
  \makebox[\MetaLabelWidth][l]{%
    \tcbox[on line,boxsep=0.5pt,left=4pt,right=4pt,top=1.5pt,bottom=1.5pt,
      arc=1.2mm,colback=Warning!9,colframe=Warning!25,boxrule=0.35pt]{%
      \sffamily\bfseries\footnotesize\color{Warning}#1}}%
  \ignorespaces}
\newcommand{\FrequencyPill}[1]{%
  \tcbox[on line,boxsep=0.5pt,left=5pt,right=5pt,top=1.5pt,bottom=1.5pt,
    arc=2mm,colback=CurrentAccent,colframe=CurrentAccent,boxrule=0pt]{%
    \sffamily\bfseries\footnotesize\color{white}考频\enspace#1}}
\newcommand{\ContentBand}[1]{%
  \ifstrequal{#1}{知识点原文摘取}{\setlength{\parindent}{0pt}}{\setlength{\parindent}{2em}}%
  \par\Needspace{4\baselineskip}\vspace{0.62em}%
  \noindent\begin{tikzpicture}[baseline=-0.55ex]
    \fill[CurrentAccent] (0,0) rectangle (0.085,0.52);
    \node[anchor=west,inner sep=0pt,text=CurrentAccent,font=\sffamily\bfseries] at (0.22,0.26) {#1};
  \end{tikzpicture}\par\vspace{0.2em}}
\newcommand{\QuestionLabel}[1]{{\sffamily\bfseries\color{CurrentAccent}#1}}
\newcommand{\PointSeparator}{%
  \setlength{\parindent}{2em}\par\vspace{0.38em}\noindent\textcolor{RuleGray}{\rule{\linewidth}{0.35pt}}\par\vspace{0.15em}}

\newtcolorbox{pastquestion}{
  enhanced jigsaw,
  breakable,
  colback=CurrentAccent!2,
  colframe=CurrentAccent!42,
  boxrule=0pt,
  borderline west={1.35pt}{0pt}{CurrentAccent},
  arc=0.8mm,
  left=2.1mm,right=2.1mm,top=1.35mm,bottom=1.35mm,
  before skip=3.8pt,after skip=3.8pt,
  fontupper=\small,
}

\newtcolorbox{correctionbox}{
  enhanced jigsaw,
  breakable,
  colback=Warning!5,
  colframe=Warning!55,
  boxrule=0.45pt,
  borderline west={2.2pt}{0pt}{Warning},
  arc=1mm,
  left=2.4mm,right=2.4mm,top=1.8mm,bottom=1.8mm,
  before skip=6pt,after skip=6pt,
}

\setlength{\LTpre}{0.65em}
\setlength{\LTpost}{0.65em}
\renewcommand{\arraystretch}{1.18}

\setcounter{tocdepth}{0}
\setcounter{secnumdepth}{2}

\begin{document}

\frontmatter
\begin{titlepage}
  \thispagestyle{empty}
  \begin{tikzpicture}[remember picture,overlay]
    \fill[PageTint] (current page.south west) rectangle (current page.north east);
    \fill[SubjectBlue] (current page.north west) rectangle ([xshift=1.35cm]current page.south west);
    \fill[SubjectBlue!8] ([xshift=-5.1cm,yshift=-5.0cm]current page.north east) circle (4.6cm);
    \draw[SubjectBlue!20,line width=1.2pt] ([xshift=-4.3cm,yshift=-4.2cm]current page.north east) circle (3.15cm);
    \fill[SubjectBlue!12] ([xshift=1.5cm,yshift=1.0cm]current page.south west) circle (3.1cm);
  \end{tikzpicture}

  \vspace*{1.7cm}
  \begin{flushleft}
    {\sffamily\bfseries\color{SubjectBlue}\small PREVENTIVE MEDICINE\enspace·\enspace REVIEW ARCHIVE}\par
    \vspace{1.05cm}
    {\sffamily\bfseries\color{Ink}\fontsize{33}{43}\selectfont
      预防医学大四下往年题\par}
    \vspace{0.15cm}
    {\sffamily\bfseries\color{SubjectBlue}\fontsize{22}{29}\selectfont
      （截止22级）\par}
    \vspace{0.75cm}
    {\color{SubjectBlue}\rule{3.2cm}{3.2pt}}\par
    \vspace{1.0cm}

    \begin{tcolorbox}[
      enhanced,
      colback=white,
      colframe=SubjectBlue!18,
      boxrule=0.5pt,
      arc=2mm,
      left=5mm,right=5mm,top=4mm,bottom=4mm,
      width=0.88\textwidth,
      shadow={1mm}{-1mm}{0mm}{black!7},
    ]
      \sffamily\small
      \textcolor{TextMuted}{网页地址}\quad
      \href{https://fxt-gw-pb.github.io/ph_study/}{\texttt{https://fxt-gw-pb.github.io/ph\_study/}}\par
      \vspace{0.3em}
      \textcolor{TextMuted}{GitHub 地址}\quad
      \href{https://github.com/fxt-gw-pb/ph_study}{\texttt{https://github.com/fxt-gw-pb/ph\_study}}
    \end{tcolorbox}

    \vspace{0.85cm}
    {\sffamily\large\bfseries\color{Ink}
      欢迎学弟学妹直接 fork 仓库并进行新的题目的补充}\par
  \end{flushleft}

  \vfill
  \begin{tikzpicture}
    \foreach \x/\metric/\captiontext in {
      0/7/门学科,
      4.0/502/个考点,
      8.0/1699/条题源记录
    }{
      \node[anchor=west,text=SubjectBlue,font=\sffamily\bfseries\fontsize{20}{24}\selectfont] at (\x,0.45) {\metric};
      \node[anchor=west,text=TextMuted,font=\sffamily\small] at (\x,0) {\captiontext};
    }
  \end{tikzpicture}
  \vspace{0.7cm}
\end{titlepage}

\chapter*{使用说明}
\addcontentsline{toc}{chapter}{使用说明}

本书完整汇编网页知识仓库中的七门课程文档，保留原始章节顺序、考频口径、知识点原文、匹配依据、题源记录、版本提示与勘误补充。为便于纸面复习，源文中的重复数字前缀由 LaTeX 自动编号取代；正文内容不作删节。

\begin{tcolorbox}[enhanced,breakable,colback=SubjectBlue!3,colframe=SubjectBlue!25,boxrule=0.45pt,arc=1.2mm,left=4mm,right=4mm,top=3mm,bottom=3mm]
  \begin{itemize}
    \item \textbf{考频}反映各源文采用的统计口径；不同课程的题源范围与计数方式并不完全相同，请先阅读每一部分的「说明」。
    \item \textbf{知识点原文摘取}用于系统复习；\textbf{对应往年题}用于定位真实题型与提问方式。
    \item 标注「未完全匹配到」「版本提示」或「勘误补充」的内容，应结合授课教师口径、最新版教材或规范再次核对。
  \end{itemize}
\end{tcolorbox}

\vspace{0.7em}
\noindent\textcolor{SubjectBlue}{\rule{1.2cm}{2pt}}\quad
\sffamily\bfseries 欢迎学弟学妹直接 fork 仓库并进行新的题目的补充。\normalfont

\clearpage
\tableofcontents
\clearpage

\mainmatter
\SetSubjectAccent{2D5FD5}{职业卫生学}

\part{职业卫生学}

\chapter*{说明}\label{users__fpb__26_spring__pkusph_study__tmp__pdfs__prevention-medicine-past-exams__01-ux804cux4e1aux536bux751fux5b66ux5f80ux5e74ux9898ux8003ux70b9ux6574ux7406.md__ux8bf4ux660e}
\addcontentsline{toc}{chapter}{说明}

\begin{itemize}
\tightlist
\item
  本文件根据《职业卫生学知识点总结》PDF（by 17 级杨贝妮）和往年题文件整理。
\item
  知识点正文均摘自《职业卫生学知识点总结》PDF。
\item
  往年题均摘自用户提供的往年题文件（04、09、10、11、14、21、24 级 + 04--14 级论述汇总）。
\item
  排序依据为考频：同一章节内按考频从高到低排列；同频按 PDF 中原始出现顺序排列。
\item
  一道考题若同时考查多个知识点，分别归入并各计 1 次。
\item
  仅纳入在往年题中匹配到的知识点；未被考查的小节/章节标注为''本章未在提供的往年题中匹配到明确考点''或省略。
\end{itemize}

\PointSeparator

\chapter{绪论}\label{users__fpb__26_spring__pkusph_study__tmp__pdfs__prevention-medicine-past-exams__01-ux804cux4e1aux536bux751fux5b66ux5f80ux5e74ux9898ux8003ux70b9ux6574ux7406.md__ux7b2cux4e00ux7ae0-ux7eeaux8bba}

\section{职业性有害因素的种类与来源}\label{users__fpb__26_spring__pkusph_study__tmp__pdfs__prevention-medicine-past-exams__01-ux804cux4e1aux536bux751fux5b66ux5f80ux5e74ux9898ux8003ux70b9ux6574ux7406.md__ux804cux4e1aux6027ux6709ux5bb3ux56e0ux7d20ux7684ux79cdux7c7bux4e0eux6765ux6e90}

\FrequencyPill{2 次}

\InfoLabel{对应小节} 一、基本概念 / 3、职业性有害因素 (occupational hazards or occupational harmful factors)
的基本概念和种类\\
\InfoLabel{匹配依据} 14 级简答直接问''职业卫生的主要活动及职业健康危险因素''；24
级简答直接问''职业有害因素的种类及来源''。

\ContentBand{知识点原文摘取}

·定义：生产工艺过程、劳动过程和工作环境(劳动条件)中存在和/或产生的，对职业人群健康、安全和作业能力可能造成不良影响的一切因素/条件统称为职业性有害因素\\
·产生条件：生产工艺过程、劳动过程和生产环境。\\
·生产工艺过程中产生的有害因素：化学因素(生产性毒物,生产性粉尘)；物理因素(异常气象条件,噪声、振动,辐射)；生物因素(病原微生物,如布氏杆菌、森林脑炎、艾滋病毒)\\
·劳动过程中的有害因素：组织制度、作息制度不合理；精神(心理)性职业紧张；劳动强度过大/生产定额不当；个别器官/系统过度紧张；不良生活方式；企业文化-监视器\\
·生产环境中的有害因素：自然环境中的因素(如高温、低温、紫外线)；厂房建筑/布局不合理、不符合职业卫生标准(有毒与无毒工艺的隔离)；由不合理生产过程/不当管理所致环境污染

\ContentBand{对应往年题}

\begin{pastquestion}

\QuestionLabel{【往年题1｜14级职卫考题.pdf｜简答】}\\
简述职业卫生的主要活动及职业健康危险因素

\end{pastquestion}

\begin{pastquestion}

\QuestionLabel{【往年题2｜24年完整试题回忆.docx｜三、简答1】}\\
简述职业有害因素的种类及来源

\end{pastquestion}

\PointSeparator

\section{工作有关疾病（职业相关疾病）的定义及特点}\label{users__fpb__26_spring__pkusph_study__tmp__pdfs__prevention-medicine-past-exams__01-ux804cux4e1aux536bux751fux5b66ux5f80ux5e74ux9898ux8003ux70b9ux6574ux7406.md__ux5de5ux4f5cux6709ux5173ux75beux75c5ux804cux4e1aux76f8ux5173ux75beux75c5ux7684ux5b9aux4e49ux53caux7279ux70b9}

\FrequencyPill{2 次}

\InfoLabel{对应小节} 一、基本概念 / 4、职业性有害因素所致健康损害（3）职业相关疾病\\
\InfoLabel{匹配依据} 04、09 简答均要求''工作相关疾病/职业相关疾病的定义及特点''。

\ContentBand{知识点原文摘取}

（3）职业相关疾病：职业性有害因素还能使机体的抵抗力下降，造成潜在的疾病显露或已患的疾病加重，从而表现为接触人群中某些常见病的发病率增高/病情加重。\\
·特点：病因往往是多因素的；潜在疾病暴露或病情加重；控制职业性有害因素，可减少其发生；不属于我国法定的职业病范围\\
·常见的职业性相关疾病：行为(精神)和身心疾病(如免疫功能紊乱、消化性溃疡、失眠、抑郁症状等)；非特异性呼吸系统疾患(鼻炎、咽炎等)；心脑血管疾病与代谢性疾病(肥胖)；腰背痛、肌肉骨骼损伤等疾患

\ContentBand{对应往年题}

\begin{pastquestion}

\QuestionLabel{【往年题1｜04级职卫考题回忆.doc｜简答】}\\
简答：职业流行病的作用、生殖损伤的表现、工作相关疾病定义及特点、职业肿瘤的特点

\end{pastquestion}

\begin{pastquestion}

\QuestionLabel{【往年题2｜09级职卫考题回忆.doc｜简答】}\\
简答：职业流行病的任务、生殖系统损害表现、中暑及表现、职业性肿瘤特点、 职业相关疾病及特点

\end{pastquestion}

\PointSeparator

\section{三级预防}\label{users__fpb__26_spring__pkusph_study__tmp__pdfs__prevention-medicine-past-exams__01-ux804cux4e1aux536bux751fux5b66ux5f80ux5e74ux9898ux8003ux70b9ux6574ux7406.md__ux4e09ux7ea7ux9884ux9632}

\FrequencyPill{1 次}

\InfoLabel{对应小节} 一、基本概念 /
2、职业病预防和控制------三级预防【掌握】；6、职业病防治的''三级预防''原则\\
\InfoLabel{匹配依据} 24 级单选第 1
题考''三级预防策略说法错误的是''。原答案（健康监护属一级预防）与另一种主流分类存在争议，详见下方勘误备注。

\ContentBand{知识点原文摘取}

2、职业病预防和控制------三级预防【掌握】\\
·职业卫生：识别、评价、预防、控制、确认\\
·一级预防(控制病因、保护敏感人群)→健康←健康监护\\
·二级预防(早发现、早诊断、早控制)→临床前期←早期诊断\\
·三级预防(改善健康状况、控制并发症)→疾病←康复治疗\\
·职业医学：健康监护、早期诊断、康复治疗

6、职业病防治的''三级预防''原则\\
·职业病防治方针：预防为主、防治结合\\
·第一级预防：病因源头控制\\
-职业病危害源头控制、劳动过程中的防护、危害治理与管理，对接触职业病危害、尚未发生健康危害的劳动者，通过综合控制措施预防损害的发生；\\
·第二级预防：三早(早发现、早干预、早治疗)\\
-职业健康检查，对从事接触职业病危害作业的劳动者，按照规定组织职业健康体检，早期发现职业病，对防护措施效果进行目标跟踪\\
·第三级预防：职业病诊断与治疗，对职业病人，按照规定，安排治疗、康复和定期检查，维护健康质量，控制并发症。

\begin{correctionbox}

【勘误】\\
·一级预防------环境监测、生物监测（目的：改善环境，促进健康）\\
·二级预防------健康监护（目的：早期发现病人）\\
·三级预防------康复治疗

\end{correctionbox}

\ContentBand{对应往年题}

\begin{pastquestion}

\QuestionLabel{【往年题1｜24年完整试题回忆.docx｜一、单选1】}\\
1、关于三级预防策略说法错误的是\\
A、从源头处控制属于一级预防\\
B、康复处理属于三级预防\\
C、诊断治疗属于二级预防\\
D、健康监护属于二级预防\\
E、\ldots\ldots{}\\
【勘误备注】本题答案受上述分类争议影响，暂不下结论；以课程教材/老师答案为准。

\end{pastquestion}

\PointSeparator

\section{职业卫生学概念}\label{users__fpb__26_spring__pkusph_study__tmp__pdfs__prevention-medicine-past-exams__01-ux804cux4e1aux536bux751fux5b66ux5f80ux5e74ux9898ux8003ux70b9ux6574ux7406.md__ux804cux4e1aux536bux751fux5b66ux6982ux5ff5}

\FrequencyPill{1 次}

\InfoLabel{对应小节} 一、基本概念 / 1、职业卫生学与职业医学的概念\\
\InfoLabel{匹配依据} 21 级名解 occupational health。

\ContentBand{知识点原文摘取}

·职业卫生学(occupational
health)：以职业人群为主要对象，主要研究劳动条件对职业人群健康的影响，主要任务是识别、评价、预测、控制和确认不良劳动条件，为保护劳动者健康、提高作业能力、改善劳动条件所应采取的措施提供科学依据。\\
-职业卫生学的内容包括职业从事者在生产工艺过程、劳动过程、生产环境接触的各种物理、化学、生物因素，作业组织安排、管理等的识别、评价、预测、控制，其主要内容主要属于一级预防。

（另见 GBZ/T 224-2009 定义，见第十四章）

\ContentBand{对应往年题}

\begin{pastquestion}

\QuestionLabel{【往年题1｜21级本科预防二班+胡润泽+职业卫生+主观题总结.docx｜名词解释】}\\
occupational health

\end{pastquestion}

\PointSeparator

\section{健康风险评估在职业卫生中的应用和意义}\label{users__fpb__26_spring__pkusph_study__tmp__pdfs__prevention-medicine-past-exams__01-ux804cux4e1aux536bux751fux5b66ux5f80ux5e74ux9898ux8003ux70b9ux6574ux7406.md__ux5065ux5eb7ux98ceux9669ux8bc4ux4f30ux5728ux804cux4e1aux536bux751fux4e2dux7684ux5e94ux7528ux548cux610fux4e49}

\FrequencyPill{1 次}

\InfoLabel{对应小节} 一、基本概念 / 10、风险评估\\
\InfoLabel{匹配依据} 21 级简答直接问''健康风险评估在职业卫生中的应用和意义''。

\ContentBand{知识点原文摘取}

10、风险评估：危害识别+暴露评估+剂量效应+风险表征→管理+控制\\
·Risk=hazard+exposure\\
11、职业卫生核心内容：预测、识别、评价、控制、确认职业有害因素，以保护和促进劳动者的健康

\ContentBand{对应往年题}

\begin{pastquestion}

\QuestionLabel{【往年题1｜21级本科预防二班+胡润泽+职业卫生+主观题总结.docx｜简答1】}\\
1、健康风险评估在职业卫生中的应用和意义

\end{pastquestion}

\PointSeparator

\section{职业卫生既是科学又是艺术（综合论述）}\label{users__fpb__26_spring__pkusph_study__tmp__pdfs__prevention-medicine-past-exams__01-ux804cux4e1aux536bux751fux5b66ux5f80ux5e74ux9898ux8003ux70b9ux6574ux7406.md__ux804cux4e1aux536bux751fux65e2ux662fux79d1ux5b66ux53c8ux662fux827aux672fux7efcux5408ux8bbaux8ff0}

\FrequencyPill{1 次}

\InfoLabel{对应小节} 综合（绪论 + 第十四章管理与实践相关概念）\\
\InfoLabel{匹配依据} 14
级论述题原题。匹配依据：题目''科学+艺术''对应于绪论中职业卫生学定义、核心内容与三级预防原则，以及第十四章对职业卫生的科学定义；属于综合性论述。疑似对应。

\ContentBand{知识点原文摘取}

11、职业卫生核心内容：预测、识别、评价、控制、确认职业有害因素，以保护和促进劳动者的健康

工(产)业卫生 Industrial
Health/hygiene：\ldots\ldots 是预期、识别、评价、交流和控制工作场所存在或产生的、能够造成劳动者和社区成员伤害、疾病、损伤的环境应激因素的艺术和科学这些因素分为生物性、化学性、物理性、人机工效学及心理性因素。（第十四章）

\ContentBand{对应往年题}

\begin{pastquestion}

\QuestionLabel{【往年题1｜14级职卫考题.pdf｜论述】}\\
举例论述为什么职业卫生既是一门科学，又是一门艺术

\end{pastquestion}

\PointSeparator

\chapter{职业生理与工效（何丽华）}\label{users__fpb__26_spring__pkusph_study__tmp__pdfs__prevention-medicine-past-exams__01-ux804cux4e1aux536bux751fux5b66ux5f80ux5e74ux9898ux8003ux70b9ux6574ux7406.md__ux7b2cux4e8cux7ae0-ux804cux4e1aux751fux7406ux4e0eux5de5ux6548ux4f55ux4e3dux534e}

\section{工效学（劳动/职业生理学）的概念与研究内容}\label{users__fpb__26_spring__pkusph_study__tmp__pdfs__prevention-medicine-past-exams__01-ux804cux4e1aux536bux751fux5b66ux5f80ux5e74ux9898ux8003ux70b9ux6574ux7406.md__ux5de5ux6548ux5b66ux52b3ux52a8ux804cux4e1aux751fux7406ux5b66ux7684ux6982ux5ff5ux4e0eux7814ux7a76ux5185ux5bb9}

\FrequencyPill{6 次}

\InfoLabel{对应小节} 一、概述 / 1、研究内容；2、概述 / 劳动(职业)生理学\\
\InfoLabel{匹配依据} ``工效学''或''职业工效学''作为名解或简答的概念/作用题，跨年度高频出现。

\ContentBand{知识点原文摘取}

【知识总结pdf没有，补充】职业功效学（occupational
ergonomics）：以人为中心，研究人、机器、设备和环境之间的相互关系，目标是健康、安全、舒适和效率。

·劳动(职业)生理学(work physiology)：研究在一定工作条件下机体器官和系统功能以及调节适应规律。\\
-工作条件(work condition)指工作任务、工作场所、对象、工作设备及生产环境等。\\
-工作条件对作业者的器官及系统产生一定的作用/效应，这种作用反之影响人的操作，二者关系是本学科研究及应用的核心

1、研究内容：\\
·体力劳动过程的生理变化与适应：体力劳动时机体的能量代谢；机体的调节与适应\\
·脑力劳动过程的生理变化与适应：脑力劳动生理特点；脑力劳动的职业卫生要求\\
·劳动负荷评价：作业类型；劳动负荷评价\\
·作业能力

\ContentBand{对应往年题}

\begin{pastquestion}

\QuestionLabel{【往年题1｜04级职卫考题回忆.doc｜名解】}\\
名解：工效学、TTS、职业肿瘤、氧债、热适应

\end{pastquestion}

\begin{pastquestion}

\QuestionLabel{【往年题2｜09级职卫考题回忆.doc｜名解】}\\
名解：全是英文，工效学、职业性肿瘤、氧债、生物监测、暂时性听阈位移

\end{pastquestion}

\begin{pastquestion}

\QuestionLabel{【往年题3｜10级职卫考题回忆.docx｜名解】}\\
全英，工效学，妇女职业卫生，热痉挛，粉尘分散度，短时间接触容许浓度，氧需

\end{pastquestion}

\begin{pastquestion}

\QuestionLabel{【往年题4｜14级职卫考题.pdf｜单选（回忆）】}\\
工效学的作用是什么。

\end{pastquestion}

\begin{pastquestion}

\QuestionLabel{【往年题5｜14级职卫考题.pdf｜简答】}\\
简述职业工效学的概念及其研究内容

\end{pastquestion}

\begin{pastquestion}

\QuestionLabel{【往年题6｜24年完整试题回忆.docx｜三、简答4】}\\
4、工效学的研究内容

\end{pastquestion}

\PointSeparator

\section{氧债（oxygen debt）/氧需（oxygen
demand）}\label{users__fpb__26_spring__pkusph_study__tmp__pdfs__prevention-medicine-past-exams__01-ux804cux4e1aux536bux751fux5b66ux5f80ux5e74ux9898ux8003ux70b9ux6574ux7406.md__ux6c27ux503aoxygen-debtux6c27ux9700oxygen-demand}

\FrequencyPill{3 次}

\InfoLabel{对应小节} 一、概述 / 4、职业活动过程中机体的生理变化与适应------作业时氧消耗的动态【考点】\\
\InfoLabel{匹配依据} 04、09 名解直接问''氧债''；10 名解问''氧需''。

\ContentBand{知识点原文摘取}

作业时氧消耗的动态【考点】\\
·氧需(oxygen
demand)：劳动一分钟所需要的氧量(l/min)。主要取决于循环系统的功能，其次是器官的功能。安静时成人为0.2-0.25
升/分；劳动时迅速上升\\
·氧上限(maximum oxygen
uptake)：血液在一分钟内能供应的最大氧量(l/min)。成年人一般不超过3L，经过体育锻炼的人可达4L。\\
·氧债(oxygen debt)：氧需和实际供氧之间的差值(l/min)。\\
氧需不超过氧上限：在作业刚开始的几分钟内，呼吸系统和循环系统还不能满足氧需，肌肉的活动所需能量是在缺氧条件下产生的。这一阶段氧需和实际供氧量出现了差额，此差值称氧债。几分钟后，呼吸和循环系统活动逐渐加强适应，氧需得到满足，即供氧稳定状态；这样的作业可维持很长时间。作业结束后，还有一个偿还氧债的过程\\
氧需超过氧上限：若劳动强度过大，氧需超过氧上限，机体则将继续处于供氧不足状态下作业，肌肉内的贮存物质(主要是糖原)迅速消耗，作业就无法持久；一般成年人的负氧债能力只有10
升左右，最大可达15 升。作业停止后一段时间内，机体继续消耗高于安静代谢的氧以偿还氧债，甚至1
小时以上才能完全补偿

\ContentBand{对应往年题}

\begin{pastquestion}

\QuestionLabel{【往年题1｜04级职卫考题回忆.doc｜名解】}\\
名解：工效学、TTS、职业肿瘤、氧债、热适应

\end{pastquestion}

\begin{pastquestion}

\QuestionLabel{【往年题2｜09级职卫考题回忆.doc｜名解】}\\
名解：全是英文，工效学、职业性肿瘤、氧债、生物监测、暂时性听阈位移

\end{pastquestion}

\begin{pastquestion}

\QuestionLabel{【往年题3｜10级职卫考题回忆.docx｜名解】}\\
全英，工效学，妇女职业卫生，热痉挛，粉尘分散度，短时间接触容许浓度，氧需

\end{pastquestion}

\PointSeparator

\section{提高作业能力的措施（白领工作中的职业危害与改善）}\label{users__fpb__26_spring__pkusph_study__tmp__pdfs__prevention-medicine-past-exams__01-ux804cux4e1aux536bux751fux5b66ux5f80ux5e74ux9898ux8003ux70b9ux6574ux7406.md__ux63d0ux9ad8ux4f5cux4e1aux80fdux529bux7684ux63aaux65bdux767dux9886ux5de5ux4f5cux4e2dux7684ux804cux4e1aux5371ux5bb3ux4e0eux6539ux5584}

\FrequencyPill{3 次}

\InfoLabel{对应小节}
一、概述（脑力劳动时的职业卫生要求；作业能力主要影响因素及改善措施【考点】；可能存在的职业危害★白领工作过程中可能存在的职业危害）\\
\InfoLabel{匹配依据} 04、09 论述题''白领''工作中可能存在的职业有害因素及其健康危害、提高作业能力的措施；11
论述''流水线可能存在的职业危害和工效学有害因素''。

\ContentBand{知识点原文摘取}

脑力劳动时的职业卫生要求【考点】\\
工作场所：噪声、照明色彩空间时的人机界面等。\\
噪声 :工作室应保持安静，环境噪声不超过 45dB ，超过 70dB 对整个脑力劳动的影响更大。\\
平稳、节奏舒缓的轻音乐或古典有利于脑力劳动。\\
大的办公室应用隔断开。\\
照明 室内光线应明亮。防止阳直射到桌面上，从左边 过来。\\
室温 20 ℃为宜。\\
墙壁颜色 明亮柔和，以淡绿、浅黄或灰色为好应避免用 刺眼或黑色的颜。

作业能力主要影响因素及改善措施【考点】\\
社会因素\\
和心理 个体因素\\
年龄、 性别身材健康及营养状况\\
环境因素\\
工作条件与性质\\
生产设备与工具\\
劳动强度与时间\\
劳动组织及制度\\
疲劳和休息\\
锻炼和练习（ training ）

可能存在的职业危害★白领工作过程中可能存在的职业危害

\ContentBand{对应往年题}

\begin{pastquestion}

\QuestionLabel{【往年题1｜04级职卫考题回忆.doc｜问答】}\\
问答：``白领''工作中可能存在的职业有害因素及其健康危害、提高作业能力的措施

\end{pastquestion}

\begin{pastquestion}

\QuestionLabel{【往年题2｜09级职卫考题回忆.doc｜问答】}\\
问答：白领工作过程中可能存在的职业危害\\
提高工作效率的方法

\end{pastquestion}

\begin{pastquestion}

\QuestionLabel{【往年题3｜11级职卫考题.docx｜应用】}\\
2 收集流水线可能存在的职业危害和工效学有害因素

\end{pastquestion}

\PointSeparator

\section{体力劳动时心血管系统血液再分配}\label{users__fpb__26_spring__pkusph_study__tmp__pdfs__prevention-medicine-past-exams__01-ux804cux4e1aux536bux751fux5b66ux5f80ux5e74ux9898ux8003ux70b9ux6574ux7406.md__ux4f53ux529bux52b3ux52a8ux65f6ux5fc3ux8840ux7ba1ux7cfbux7edfux8840ux6db2ux518dux5206ux914d}

\FrequencyPill{1 次}

\InfoLabel{对应小节} 一、概述 / 3、体力作业时机的调节和适应（神经-体液调节器官系统）\\
\InfoLabel{匹配依据} 24 级单选第 4 题''体力劳动时心血管系统的血液再分配情况''。疑似对应（PDF
仅提供框架性描述：神经体液调节各器官系统的协调、骨骼肌活动为主，并未直接列出''骨骼肌和心肌的血流量增大''，但属于本章生理调节范畴）。

\ContentBand{知识点原文摘取}

3、体力作业时机的调节和适应：在生产劳动过程中，为保证能量供应和各器官系统的协调，机体通过神经体液调节各器官系统的协调，机体通过神经体液调节各器官系统的生理功能，以适应生产劳动需要。\\
·工作种类、工作强度、作业姿势、轮班作业、个体差异→人在生产劳动过程←→神经-体液的调节和适应→完成作业、促进健康；作业能力下降、损害健康(负荷过高、劳动时间过长、恶劣的环境)\\
4、职业活动过程中机体的生理变化与适应\\
-体力工作时机的能量代谢：人类的工作是体力劳动与脑力劳动相结合进行的，不同类型的劳动有所偏重。骨骼肌约占人体中的40\%，以骨骼肌活动为主的体力劳动能量消耗较大。

\ContentBand{对应往年题}

\begin{pastquestion}

\QuestionLabel{【往年题1｜24年完整试题回忆.docx｜一、单选4】}\\
4、体力劳动时心血管系统的血液再分配情况\\
A、骨骼肌和内脏的血流量增大\\
B、骨骼肌和心肌的血流量增大\\
C、肝肾的血流量增大\\
D\&E、\ldots\ldots{}

\end{pastquestion}

\PointSeparator

\section{职业性肌肉骨骼疾患（WMSDs）}\label{users__fpb__26_spring__pkusph_study__tmp__pdfs__prevention-medicine-past-exams__01-ux804cux4e1aux536bux751fux5b66ux5f80ux5e74ux9898ux8003ux70b9ux6574ux7406.md__ux804cux4e1aux6027ux808cux8089ux9aa8ux9abcux75beux60a3wmsds}

\FrequencyPill{1 次}

\InfoLabel{对应小节} 第二章 职业生理与工效（手动补充）\\
\InfoLabel{匹配依据} 未匹配到，手动补充。

\ContentBand{知识点原文摘取}

【知识总结pdf没有，补充】职业性肌肉骨骼疾患（work related musculoskeletal
disorders，WMSDs）：由于暴露于工作场所的相关危险因素所导致的肌肉、神经、肌腱、关节、软骨和椎间盘等的损伤疾病，包括腰背痛、腕管综合征、其他肌肉骨骼系统或结缔组织疾病，多由于重复或过度用力或弯曲、攀爬、扭曲引起。

\ContentBand{对应往年题}

\begin{pastquestion}

\QuestionLabel{【往年题1｜21级本科预防二班+胡润泽+职业卫生+主观题总结.docx｜名词解释】}\\
职业性肌肉骨骼疾患（work related musculoskeletal disorders，WMSDs）

\end{pastquestion}

\PointSeparator

\chapter{不良气象条件}\label{users__fpb__26_spring__pkusph_study__tmp__pdfs__prevention-medicine-past-exams__01-ux804cux4e1aux536bux751fux5b66ux5f80ux5e74ux9898ux8003ux70b9ux6574ux7406.md__ux7b2cux4e09ux7ae0-ux4e0dux826fux6c14ux8c61ux6761ux4ef6}

\section{中暑（含分类、临床表现、热痉挛）}\label{users__fpb__26_spring__pkusph_study__tmp__pdfs__prevention-medicine-past-exams__01-ux804cux4e1aux536bux751fux5b66ux5f80ux5e74ux9898ux8003ux70b9ux6574ux7406.md__ux4e2dux6691ux542bux5206ux7c7bux4e34ux5e8aux8868ux73b0ux70edux75c9ux631b}

\FrequencyPill{4 次}

\InfoLabel{对应小节} 一、不良气象条件主要内容 /
5、高温作业（中暑）（5）高温作业------中暑：发病机制和临床表现；高温中暑的诊断分级\\
\InfoLabel{匹配依据} 09 简答''中暑及表现''、11 简答''中暑的分类和特点''、10 名解''热痉挛''、24 单选
6（热射病死亡表现）。

\ContentBand{知识点原文摘取}

（5）高温作业------中暑：高温环境下由于热平衡和/或水盐代谢紊乱等而引起的一种以中枢神经系统和/或心血管系统障碍为主要表现的急性热致疾病。\\
·病因：环境因素；劳动强度；劳动时间；个体因素\\
·致病因素：气温高；气湿大；气流小；热辐射强；劳动强度大；劳动时间过长\\
·诱发因素：体弱；肥胖；睡眠不足；未产生热适应\\
·发病机制和临床表现\\
-热射病(包括日射病sun
stroke)：人在热环境下，散热途径受阻，体温调节失调所致。其临床特点是，在高温环境中突然发病，体温可高达40
℃以上，开始时大量出汗，以后无汗，并伴有干热和意识障碍、嗜睡、昏迷等中枢神经系统症状。\\
-热痉挛(heat
cramp)：因高温过量出汗，体内钠、钾过量丢失所致。其临床特点是骨骼肌突然痉挛，并伴有收缩痛。痉挛以腓肠肌等四肢肌肉和腹肌为多见。痉挛发作多对称性，自行缓解，患者神志清醒，体温正常。\\
-热衰竭(heat
exhaustion)：脑部暂时供血不足。。发病机制不明确。多数认为在高温、高湿环境下，皮肤血流的增加不伴有内脏血管收缩或血容量的相应增加，导致脑部暂时供血减少而晕厥。发病一般迅速，先有头昏、头痛、心悸、出汗、恶心、呕吐、皮肤湿冷、面色苍白、血压下降，继而晕厥，体温不高或稍高；休息片刻即可清醒，一般不引起循环衰竭。\\
·高温中暑的诊断分级\\
-先兆中暑：在高温作业场所劳动一定时间后，出现大量出汗、口渴、头昏、耳鸣、胸闷、心悸、恶心、全身疲乏、四肢无力、注意力不集中等症状。体温正常或略有升高。\\
-轻症中暑：具备以下条件之一者，①头晕、胸闷、心悸、面色潮红、皮肤灼热；②有呼吸与循环衰竭的早期症状，大量出汗、面色苍白、血压下降、脉搏细弱而快；③肛温升高达38.5℃以上\\
-重症中暑：出现辐射热、痉挛或衰竭主要临床症状者

\ContentBand{对应往年题}

\begin{pastquestion}

\QuestionLabel{【往年题1｜09级职卫考题回忆.doc｜简答】}\\
简答：职业流行病的任务、生殖系统损害表现、中暑及表现、职业性肿瘤特点、 职业相关疾病及特点

\end{pastquestion}

\begin{pastquestion}

\QuestionLabel{【往年题2｜10级职卫考题回忆.docx｜名解】}\\
全英，工效学，妇女职业卫生，热痉挛，粉尘分散度，短时间接触容许浓度，氧需

\end{pastquestion}

\begin{pastquestion}

\QuestionLabel{【往年题3｜11级职卫考题.docx｜简答6】}\\
6 中暑的分类和特点

\end{pastquestion}

\begin{pastquestion}

\QuestionLabel{【往年题4｜24年完整试题回忆.docx｜一、单选6】}\\
6、下列哪种表现可能发生热作业人员死亡\\
A、体温可高达40℃以上，开始时大量出汗，以后无汗\\
B、骨骼肌突然痉挛，并伴有收缩痛\\
C\&D\&E、\ldots\ldots{}

\end{pastquestion}

\PointSeparator

\section{热适应（heat
acclimatization）}\label{users__fpb__26_spring__pkusph_study__tmp__pdfs__prevention-medicine-past-exams__01-ux804cux4e1aux536bux751fux5b66ux5f80ux5e74ux9898ux8003ux70b9ux6574ux7406.md__ux70edux9002ux5e94heat-acclimatization}

\FrequencyPill{2 次}

\InfoLabel{对应小节} 一、不良气象条件主要内容 / 5、高温作业（中暑）（4）热适应\\
\InfoLabel{匹配依据} 04 名解''热适应''、14 名解 Heat adaptation。

\ContentBand{知识点原文摘取}

（4）热适应(heat acclimatization)：人体在热环境一定时间后对热负荷产生的适应现象。\\
·热应激蛋白(heat shock protein, HSP)：机体细胞在受热及热适应诱导后合成的一组蛋白质，特别是分子量为27kD 和70kD
的HSP27 和HSP70，可以保护机体免受一定范围高温的致死性损伤。

\ContentBand{对应往年题}

\begin{pastquestion}

\QuestionLabel{【往年题1｜04级职卫考题回忆.doc｜名解】}\\
名解：工效学、TTS、职业肿瘤、氧债、热适应

\end{pastquestion}

\begin{pastquestion}

\QuestionLabel{【往年题2｜14级职卫考题.pdf｜名解】}\\
Heat adaptation

\end{pastquestion}

\PointSeparator

\section{物理因素特点}\label{users__fpb__26_spring__pkusph_study__tmp__pdfs__prevention-medicine-past-exams__01-ux804cux4e1aux536bux751fux5b66ux5f80ux5e74ux9898ux8003ux70b9ux6574ux7406.md__ux7269ux7406ux56e0ux7d20ux7279ux70b9}

\FrequencyPill{1 次}

\InfoLabel{对应小节} 一、不良气象条件主要内容 / 2、物理因素特点\\
\InfoLabel{匹配依据} 10 简答''物理因素的特点''。

\ContentBand{知识点原文摘取}

2、物理因素特点\\
·自然界存在，大部分是人体生理活动必须的\\
·有特定的参数\\
·有明确的来源，其存在与源的工作状态有关\\
·空间的强度不均匀，以源为中心，向四周传播，强度随距离增加呈指数关系衰减。\\
·存在状态不同，性质不同，对人体作用不同，噪声、微波等\\
·对人体的危害程度与物理参数不呈直线相关关系，在某一定范围对人无害，高/低于这一范围有不良影响\\
·物理因素还研究''适宜''范围，如至适温度，合理照明等\\
·绝大多数物理因素在脱离接触后，体内没有该物质的残留\\
·物理因素的预防，在各个环节都有可行、有效的方法。

\ContentBand{对应往年题}

\begin{pastquestion}

\QuestionLabel{【往年题1｜10级职卫考题回忆.docx｜简答】}\\
物理因素的特点、有机磷中毒的机制和临床表现、李涛所长之劳动者职业健康监护制度内容、手臂振动病的定义,临床表现和典型特征、常见职业伤害事故类型

\end{pastquestion}

\PointSeparator

\chapter{噪声与振动}\label{users__fpb__26_spring__pkusph_study__tmp__pdfs__prevention-medicine-past-exams__01-ux804cux4e1aux536bux751fux5b66ux5f80ux5e74ux9898ux8003ux70b9ux6574ux7406.md__ux7b2cux56dbux7ae0-ux566aux58f0ux4e0eux632fux52a8}

\section{暂时性听阈位移（TTS）/ 听觉适应 /
听觉疲劳}\label{users__fpb__26_spring__pkusph_study__tmp__pdfs__prevention-medicine-past-exams__01-ux804cux4e1aux536bux751fux5b66ux5f80ux5e74ux9898ux8003ux70b9ux6574ux7406.md__ux6682ux65f6ux6027ux542cux9608ux4f4dux79fbtts-ux542cux89c9ux9002ux5e94-ux542cux89c9ux75b2ux52b3}

\FrequencyPill{4 次}

\InfoLabel{对应小节} 三、噪声对人体的影响 / 1、听觉系统【掌握】（1）暂时性听阈位移(TTS)\\
\InfoLabel{匹配依据} 04、09 名解 TTS / 暂时性听阈位移；24 名解 temporary threshold shift；24 单选 7''听阈提高
10dB，1 分钟恢复''对应听觉适应。

\ContentBand{知识点原文摘取}

（1）暂时性听阈位移(temporary threshold shift,
TTS)：是指人或动物接触噪声后引起听阈变化，脱离噪声环境后经过一段时间听力可以恢复到原来水平。\\
·听觉适应(auditory
adaptation)：短时间暴露在强烈噪声环境中，引起听力明显下降，离开噪声环境后，听阈提高超过10\textasciitilde15dB，离开噪声环境1
分钟之内可以恢复。\\
·听觉疲劳（auditory
fatigue）：较长时间停留在强烈噪声环境中，引起听力明显下降，离开噪声环境后，听阈提高15\textasciitilde30dB，需要数小时甚至数十小时听力才能恢复。

\ContentBand{对应往年题}

\begin{pastquestion}

\QuestionLabel{【往年题1｜04级职卫考题回忆.doc｜名解】}\\
名解：工效学、TTS、职业肿瘤、氧债、热适应

\end{pastquestion}

\begin{pastquestion}

\QuestionLabel{【往年题2｜09级职卫考题回忆.doc｜名解】}\\
名解：全是英文，工效学、职业性肿瘤、氧债、生物监测、暂时性听阈位移

\end{pastquestion}

\begin{pastquestion}

\QuestionLabel{【往年题3｜24年完整试题回忆.docx｜二、名解2】}\\
2、temporary threshold shift, TTS

\end{pastquestion}

\begin{pastquestion}

\QuestionLabel{【往年题4｜24年完整试题回忆.docx｜一、单选7】}\\
7、一工人在噪声环境中工作，检测到听阈提高10dB，脱离噪声环境1分钟后恢复，属于\\
A、永久性听阈位移\\
B、听觉适应\\
C、听觉疲劳\\
D、听力损失\\
E、噪声性耳聋

\end{pastquestion}

\PointSeparator

\section{等响曲线（equal loudness
curve）}\label{users__fpb__26_spring__pkusph_study__tmp__pdfs__prevention-medicine-past-exams__01-ux804cux4e1aux536bux751fux5b66ux5f80ux5e74ux9898ux8003ux70b9ux6574ux7406.md__ux7b49ux54cdux66f2ux7ebfequal-loudness-curve}

\FrequencyPill{3 次}

\InfoLabel{对应小节} 二、噪声的物理特性及其评价 / 4、人对声音的主观感觉（1）等响曲线\\
\InfoLabel{匹配依据} 11、14、21 名解均出 equal loudness curve。

\ContentBand{知识点原文摘取}

（1）等响曲线：根据人耳对声音的感觉特性，使用声压级和频率，采用实验方法测出人耳对声音音响的主观感觉量，称为响度级(loudness
level),单位为㕫(phone)\\
·以1000Hz
纯音作为基准音，其他不同频率的纯音通过实验听起来与某一声压级的基准音响度相同时，即为等响，则该条件下的被测纯音响度级(㕫值)就等于基准音的声压级（dB
值）。\\
·利用与基准音比较的方法，得出听阈范围各种声频的响度级，将各个频率相同响度的数值用曲线连接，绘出各种响度的等响曲线图，称为等响曲线\\
·人耳对声音的主观感觉特点：对高频敏感，对低频不敏感。

\ContentBand{对应往年题}

\begin{pastquestion}

\QuestionLabel{【往年题1｜11级职卫考题.docx｜名解】}\\
妇女职业卫生、建设的''三同时''、AED、equal loudness curve、生物监测、occupational injury

\end{pastquestion}

\begin{pastquestion}

\QuestionLabel{【往年题2｜14级职卫考题.pdf｜名解】}\\
equal loudness curve

\end{pastquestion}

\begin{pastquestion}

\QuestionLabel{【往年题3｜21级本科预防二班+胡润泽+职业卫生+主观题总结.docx｜名词解释】}\\
equal loudness curve

\end{pastquestion}

\PointSeparator

\section{影响噪声对机体作用的因素}\label{users__fpb__26_spring__pkusph_study__tmp__pdfs__prevention-medicine-past-exams__01-ux804cux4e1aux536bux751fux5b66ux5f80ux5e74ux9898ux8003ux70b9ux6574ux7406.md__ux5f71ux54cdux566aux58f0ux5bf9ux673aux4f53ux4f5cux7528ux7684ux56e0ux7d20}

\FrequencyPill{2 次}

\InfoLabel{对应小节} 三、噪声对人体的影响 / 9、影响噪声对机体作用的因素【考点】\\
\InfoLabel{匹配依据} 14 简答''影响噪声的健康效应的因素''；21 简答''影响噪声对机体作用的因素''。

\ContentBand{知识点原文摘取}

9、影响噪声对机体作用的因素【考点】：噪声强度、频率；接触时间和接触方式；噪声的性质；其它有害因素共同存在；机体健康状况及个人敏感性；个体防护

\ContentBand{对应往年题}

\begin{pastquestion}

\QuestionLabel{【往年题1｜14级职卫考题.pdf｜简答】}\\
简述影响噪声的健康效应的因素有哪些

\end{pastquestion}

\begin{pastquestion}

\QuestionLabel{【往年题2｜21级本科预防二班+胡润泽+职业卫生+主观题总结.docx｜简答3】}\\
3、影响噪声对机体作用的因素

\end{pastquestion}

\PointSeparator

\section{生产性噪声的特征}\label{users__fpb__26_spring__pkusph_study__tmp__pdfs__prevention-medicine-past-exams__01-ux804cux4e1aux536bux751fux5b66ux5f80ux5e74ux9898ux8003ux70b9ux6574ux7406.md__ux751fux4ea7ux6027ux566aux58f0ux7684ux7279ux5f81}

\FrequencyPill{1 次}

\InfoLabel{对应小节} 一、噪声概念和分类 / 2、噪声 / （2）生产性噪声的特征【掌握】\\
\InfoLabel{匹配依据} 24 简答''生产性噪声的特征''。

\ContentBand{知识点原文摘取}

（2）生产性噪声的特征【掌握】：强度高(多超过80dB(A)，甚至高达110dB(A)以上)；高频音所占比例大(高频音危害大于中、低频音)；持续暴露时间长；往往与其他有害因素同时存在(振动、高温、毒物等)

\ContentBand{对应往年题}

\begin{pastquestion}

\QuestionLabel{【往年题1｜24年完整试题回忆.docx｜三、简答3】}\\
3、生产性噪声的特征

\end{pastquestion}

\PointSeparator

\section{手臂振动病}\label{users__fpb__26_spring__pkusph_study__tmp__pdfs__prevention-medicine-past-exams__01-ux804cux4e1aux536bux751fux5b66ux5f80ux5e74ux9898ux8003ux70b9ux6574ux7406.md__ux624bux81c2ux632fux52a8ux75c5}

\FrequencyPill{1 次}

\InfoLabel{对应小节} 五、振动 / 4、振动对机体的影响（2）局部振动·主要危害：手臂振动病\\
\InfoLabel{匹配依据} 10 简答''手臂振动病的定义,临床表现和典型特征''。

\ContentBand{知识点原文摘取}

·主要危害：手臂振动病Hand-arm vibration
disease---是长期从事手传振动作业而引起的以手部末梢和/或手臂神经功能障碍为主的疾病，并可引起手、臂骨关节-肌肉的损伤。\\
-典型表现：振动性白指，又称职业性雷诺现象\\
-GBZ 7---2014
职业性手臂振动病的诊断：根据一年以上连续从事手传振动作业的职业史，以手部末梢循环障碍、手臂神经功能障碍和/或骨关节肌肉损伤为主的临床表现，结合末梢循环功能、神经－肌电图检查结果，参考作业环境的职业卫生学资料，综合分析，排除其他病因所致类似疾病，方可诊断。

\ContentBand{对应往年题}

\begin{pastquestion}

\QuestionLabel{【往年题1｜10级职卫考题回忆.docx｜简答】}\\
物理因素的特点、有机磷中毒的机制和临床表现、李涛所长之劳动者职业健康监护制度内容、手臂振动病的定义,临床表现和典型特征、常见职业伤害事故类型

\end{pastquestion}

\PointSeparator

\section{振动评价的常用物理参量}\label{users__fpb__26_spring__pkusph_study__tmp__pdfs__prevention-medicine-past-exams__01-ux804cux4e1aux536bux751fux5b66ux5f80ux5e74ux9898ux8003ux70b9ux6574ux7406.md__ux632fux52a8ux8bc4ux4ef7ux7684ux5e38ux7528ux7269ux7406ux53c2ux91cf}

\FrequencyPill{1 次}

\InfoLabel{对应小节} 五、振动 / 2、振动的物理参量\\
\InfoLabel{匹配依据} 04-14级职业卫生论述往年题汇总.docx 中 14 级附录列出''振动评价：振动频谱、共振频谱、4
小时等能量频率计权加速度有效值''，对应名解题/复习题。疑似对应。

\ContentBand{知识点原文摘取}

·振动对人体健康的影响是振动位移、速度和加速度联合作用及其与机体相互作用的结果\\
·振动评价常用物理参量多采用：振动频谱；共振频谱；4 小时等能量频率计权加速度有效值

\ContentBand{对应往年题}

\begin{pastquestion}

\QuestionLabel{【往年题1｜04-14级职业卫生论述往年题汇总.docx｜14级附录】}\\
振动评价：振动频谱、共振频谱、4小时等能量频率计权加速度有效值

\end{pastquestion}

\PointSeparator

\chapter{电离辐射与非电离辐射}\label{users__fpb__26_spring__pkusph_study__tmp__pdfs__prevention-medicine-past-exams__01-ux804cux4e1aux536bux751fux5b66ux5f80ux5e74ux9898ux8003ux70b9ux6574ux7406.md__ux7b2cux4e94ux7ae0-ux7535ux79bbux8f90ux5c04ux4e0eux975eux7535ux79bbux8f90ux5c04}

\section{静态作业（static
work）}\label{users__fpb__26_spring__pkusph_study__tmp__pdfs__prevention-medicine-past-exams__01-ux804cux4e1aux536bux751fux5b66ux5f80ux5e74ux9898ux8003ux70b9ux6574ux7406.md__ux9759ux6001ux4f5cux4e1astatic-work}

\FrequencyPill{1 次}

\InfoLabel{对应小节} 第二/五章工效学作业类型\\
\InfoLabel{匹配依据} 24 名解 static
work。题目对应''作业类型''------属于第五章/职业生理与工效中作业类型分类内容。疑似对应（PDF
在第二章一节''研究内容''中提到''劳动负荷评价：作业类型''，static work 为标准教科书概念）。

\ContentBand{知识点原文摘取}

【知识总结pdf没有，补充】静态作业/静力作业(static
work)：主要依靠肌肉的等长性收缩来维持一定的体位，使躯体和四肢关节保持不动时所进行的作业。

1、研究内容：\\
·体力劳动过程的生理变化与适应：体力劳动时机体的能量代谢；机体的调节与适应\\
·脑力劳动过程的生理变化与适应：脑力劳动生理特点；脑力劳动的职业卫生要求\\
·劳动负荷评价：作业类型；劳动负荷评价\\
·作业能力

\ContentBand{对应往年题}

\begin{pastquestion}

\QuestionLabel{【往年题1｜24年完整试题回忆.docx｜二、名解4】}\\
4、static work

\end{pastquestion}

\PointSeparator

\section{手控制器类型}\label{users__fpb__26_spring__pkusph_study__tmp__pdfs__prevention-medicine-past-exams__01-ux804cux4e1aux536bux751fux5b66ux5f80ux5e74ux9898ux8003ux70b9ux6574ux7406.md__ux624bux63a7ux5236ux5668ux7c7bux578b}

\FrequencyPill{1 次}

\InfoLabel{对应小节} 第五章 工效学（控制器设计相关）\\
\InfoLabel{匹配依据} 24 单选 5''不属于手控制器的是（速度计）``；属于工效学控制器分类内容。疑似对应（PDF
中第五章主要为电磁辐射、未直接列出手控制器分类，应来源于课件 PPT，故仅做疑似对应记录）。

\ContentBand{知识点原文摘取}

（PDF 未直接列出''手控制器''分类内容，疑似对应于工效学课件相关内容；不补充外部知识）

\ContentBand{对应往年题}

\begin{pastquestion}

\QuestionLabel{【往年题1｜24年完整试题回忆.docx｜一、单选5】}\\
5、不属于手控制器的是\\
A、按压式控制器\\
B、旋转式控制器\\
C、移动式控制器\\
D、轮盘\\
E、速度计

\end{pastquestion}

\PointSeparator

\section{电磁辐射 /
比吸收率（SAR）}\label{users__fpb__26_spring__pkusph_study__tmp__pdfs__prevention-medicine-past-exams__01-ux804cux4e1aux536bux751fux5b66ux5f80ux5e74ux9898ux8003ux70b9ux6574ux7406.md__ux7535ux78c1ux8f90ux5c04-ux6bd4ux5438ux6536ux7387sar}

\FrequencyPill{1 次}

\InfoLabel{对应小节} 第五章 / 一、电磁辐射(Electromagnetic radiation)概述\\
\InfoLabel{匹配依据} 04-14 级论述汇总中作为定义类辅助考点出现（用于辨析名解或单选）。疑似对应。

\ContentBand{知识点原文摘取}

·电荷产生电场，电荷流动产生磁场，两者合成为电磁场；电磁场以波的形式（电场和磁场的振幅相互垂直）向外传递电磁能量，形成了电磁辐射

·确定性效应(deterministic
effects)：指受到辐射的组织，当损伤细胞足够多，表现出组织或器官功能呈现不同程度丧失的表现。这种效应具有明确的剂量-效应关系，即效应的严重程度与剂量成正比，有剂量阈值。\\
·随机性效应(stochastic
effects)：损伤的发生率与剂量成正比，损伤的严重程度与剂量无关，损伤效应无剂量阈值的电离辐射效应。

\ContentBand{对应往年题}

\begin{pastquestion}

\QuestionLabel{【往年题1｜04-14级职业卫生论述往年题汇总.docx｜14级附录】}\\
电磁辐射 Electromagnetic
radiation：电荷产生电场，电荷流动产生磁场，两者合称为电磁场；电磁场以波的形式（电场和磁场的振幅相互垂直）向外传递电磁能量，形成了电磁辐射。随机效应：发生率与剂量相关，严重程度与剂量无关，无剂量阈值。确定效应：严重程度具有明确的剂量-效应关系，和阈剂量。\\
电磁辐射剂量=比吸收率 specific absorption rate
：生物体每单位质量吸收的电磁辐射功率，单位为W/kg。对生物系统实际吸收的电磁辐射能量进行定量化

\end{pastquestion}

\PointSeparator

\chapter{粉尘与尘肺病}\label{users__fpb__26_spring__pkusph_study__tmp__pdfs__prevention-medicine-past-exams__01-ux804cux4e1aux536bux751fux5b66ux5f80ux5e74ux9898ux8003ux70b9ux6574ux7406.md__ux7b2cux516dux7ae0-ux7c89ux5c18ux4e0eux5c18ux80baux75c5}

\section{生产性粉尘的理化特性及其卫生学意义（含粉尘分散度、AED）}\label{users__fpb__26_spring__pkusph_study__tmp__pdfs__prevention-medicine-past-exams__01-ux804cux4e1aux536bux751fux5b66ux5f80ux5e74ux9898ux8003ux70b9ux6574ux7406.md__ux751fux4ea7ux6027ux7c89ux5c18ux7684ux7406ux5316ux7279ux6027ux53caux5176ux536bux751fux5b66ux610fux4e49ux542bux7c89ux5c18ux5206ux6563ux5ea6aed}

\FrequencyPill{6 次}

\InfoLabel{对应小节} 二、生产性粉尘的理化特性及其卫生学意义；2、粉尘的分散度（AED）\\
\InfoLabel{匹配依据} 10 名解''粉尘分散度''；11、21 名解''AED''；14 简答''评价粉尘危害性需要考虑的因素''；21
简答''生产性粉尘的理化性质''；24 单选 8''生产性粉尘理化特性及其卫生学意义说法错误''。

\ContentBand{知识点原文摘取}

二、生产性粉尘的理化特性及其卫生学意义：根据生产性粉尘来源、分类及其理化特性可初步判断其对人体的危害性质和程度。从卫生学角度出发，主要应考虑的粉尘理化特性包括------化学成分、浓度、接触时间，分散度，硬度，溶解性，荷电性，爆炸性，放射性。

1、粉尘的化学成分、浓度和接触时间：同一种粉尘，一般浓度越高，接触时间越长，危害越大\\
·例：二氧化硅粉尘→肺纤维化（矽肺）；金属粉尘→粉尘沉着症、过敏和中毒等

2、粉尘的分散度：指物质被粉碎的程度，以粉尘粒径大小(μm)的数量或质量组成百分比来表示。前者称为粒子分散度，后者称为质量分散度。\\
·意义：分散度越高，在空气中停留时间越长，被人体吸收的机会越多；分散度越高，比表面积越大，越易参与理化反应，对人体危害越大。\\
·不同种类的粉尘由于粉尘的密度和形状不同，同一粒径的粉尘在空气中的沉降速度不同，为了互相比较，引入空气动力学直径（aerodynamic
equivalent diameter，AED）：不考虑尘粒的形状，大小和密度，如果其在空气中的沉降速度与一种密度为1
的球形粒子的沉降速度一样时，则这种球形粒子的直径即为该种粉尘粒子的AED。\\
-同一空气动力学直径的尘粒，在空气中具有相同的沉降速度和悬浮时间，并趋向于沉降在人体呼吸道内的相同区域。一般认为，AED
小于15μm 的粒子可进入呼吸道，其中10～15μm 的粒子主要沉积在上呼吸道，因此把直径小于15μm
的尘粒称为可吸入性粉尘（inhalable dust）；5μm 以下的粒子可到达呼吸道深部和肺泡区，称之为呼吸性粉尘（respirable
dust）。

3、粉尘的硬度：\ldots\ldots{}\\
4、粉尘的溶解度：某些有毒粉尘，如含铅、砷等的粉尘可在上呼吸道溶解吸收，其溶解度越高对人体毒作用越强；相对无毒的粉尘如面粉，其溶解度越高作用越低；石英粉尘等很难溶解，在体内持续产生危害作用。\\
5、粉尘的荷电性\ldots\ldots{}\\
6、粉尘的爆炸性：可氧化的粉尘在适宜的浓度下一旦遇到明火、电火花和放电时，可发生爆炸。\\
7、粉尘的放射性\ldots\ldots{}

\ContentBand{对应往年题}

\begin{pastquestion}

\QuestionLabel{【往年题1｜10级职卫考题回忆.docx｜名解】}\\
全英，工效学，妇女职业卫生，热痉挛，粉尘分散度，短时间接触容许浓度，氧需

\end{pastquestion}

\begin{pastquestion}

\QuestionLabel{【往年题2｜11级职卫考题.docx｜名解】}\\
妇女职业卫生、建设的''三同时''、AED、equal loudness curve、生物监测、occupational injury

\end{pastquestion}

\begin{pastquestion}

\QuestionLabel{【往年题3｜14级职卫考题.pdf｜简答】}\\
简述评价粉尘危害性需要考虑的因素有哪些

\end{pastquestion}

\begin{pastquestion}

\QuestionLabel{【往年题4｜21级本科预防二班+胡润泽+职业卫生+主观题总结.docx｜名词解释】}\\
aerodynamic equivalent diameter，AED

\end{pastquestion}

\begin{pastquestion}

\QuestionLabel{【往年题5｜21级本科预防二班+胡润泽+职业卫生+主观题总结.docx｜简答2】}\\
2、从卫生学角度，简述生产性粉尘的理化性质

\end{pastquestion}

\begin{pastquestion}

\QuestionLabel{【往年题6｜24年完整试题回忆.docx｜一、单选8】}\\
8、关于生产性粉尘理化特性及其卫生学意义说法错误的是\\
A、同一空气动力学直径的尘粒，在空气中具有相同的沉降速度和悬浮时间，并趋向于沉降在人体呼吸道内的相同区域\\
B、溶解度越高毒作用越强\\
C、可氧化的粉尘如煤、面粉、糖、亚麻、硫磺、铝等，在适宜的浓度下一旦遇到明火、电火花和放电时，可发生爆炸。\\
D、同一种粉尘，浓度越高，危害越大\\
E、同一种粉尘，接触时间越长，危害越大

\end{pastquestion}

\PointSeparator

\section{矽肺（silicosis）与尘肺病分类}\label{users__fpb__26_spring__pkusph_study__tmp__pdfs__prevention-medicine-past-exams__01-ux804cux4e1aux536bux751fux5b66ux5f80ux5e74ux9898ux8003ux70b9ux6574ux7406.md__ux77fdux80basilicosisux4e0eux5c18ux80baux75c5ux5206ux7c7b}

\FrequencyPill{3 次}

\InfoLabel{对应小节} 三、尘肺病 / 1、尘肺病的分类；5、几种主要的尘肺病（1）游离二氧化硅与矽肺\\
\InfoLabel{匹配依据} 24 名解 silicosis；24 单选 9''尘肺病按病因分为五类，下列错误的是''；14
单选''哪个不是职业病规定的尘肺病''。

\ContentBand{知识点原文摘取}

三、尘肺病（pneumoconiosis）：是由于生产过程中长期吸入粉尘而发生的以肺组织纤维化为主的疾病。\\
1、尘肺病的分类：按病因分为五类\\
·矽肺：吸入含有游离二氧化硅粉尘引起，如石英尘\\
·硅酸盐肺：吸入结合型二氧化硅粉尘，如石棉、滑石、云母等粉尘\\
·炭尘肺：吸入煤炭、石墨、碳黑、活性炭等粉尘\\
·金属尘肺：吸入某些金属粉尘，如铁、铝尘\\
·混合性尘肺：吸入含游离二氧化硅粉尘和其他粉尘等引起，如煤尘\\
·我国的法定尘肺病包含十三种：矽肺、煤工尘肺、石墨尘肺、碳黑尘肺、石棉肺、滑石尘肺、水泥尘肺、云母尘肺、陶工尘肺、铝尘肺、电焊工尘肺、铸工尘肺、根据《尘肺病诊断标准》和《尘肺病理诊断标准》可以诊断的其他尘肺病

（1）游离二氧化硅与矽肺\\
·矽肺(silicosis)：在生产过程中长期吸入游离SiO2 粉尘而引起的以肺部弥漫性纤维化为主的全身性疾病\\
-我国矽肺病例占尘肺总病例的40\%左右，位居第二，是尘肺中危害最严重的一种。

\ContentBand{对应往年题}

\begin{pastquestion}

\QuestionLabel{【往年题1｜14级职卫考题.pdf｜单选（回忆）】}\\
哪个不是职业病规定的尘肺病

\end{pastquestion}

\begin{pastquestion}

\QuestionLabel{【往年题2｜24年完整试题回忆.docx｜二、名解3】}\\
3、silicosis

\end{pastquestion}

\begin{pastquestion}

\QuestionLabel{【往年题3｜24年完整试题回忆.docx｜一、单选9】}\\
9、尘肺病按病因分为五类，下列错误的是：\\
A、矽肺\\
B、煤工尘肺\\
C、硅酸盐肺\\
D、炭尘肺\\
E、金属尘肺

\end{pastquestion}

\PointSeparator

\section{我国综合防尘和降尘措施（``八字方针''）}\label{users__fpb__26_spring__pkusph_study__tmp__pdfs__prevention-medicine-past-exams__01-ux804cux4e1aux536bux751fux5b66ux5f80ux5e74ux9898ux8003ux70b9ux6574ux7406.md__ux6211ux56fdux7efcux5408ux9632ux5c18ux548cux964dux5c18ux63aaux65bdux516bux5b57ux65b9ux9488}

\FrequencyPill{1 次}

\InfoLabel{对应小节} 五、粉尘危害的控制管理 / 5、我国的综合防尘和降尘措施\\
\InfoLabel{匹配依据} 24 简答''论述我国综合防尘治尘的'八字方针'的具体内涵''。

\ContentBand{知识点原文摘取}

5、我国的综合防尘和降尘措施：可概括为''革水密风护管教查''八字方针，即工艺改革和技术革新、湿式作业、将发尘源密闭、加强通风及抽风措施、个人防护、经常性地维修和管理工作、加强宣传教育、定期检查环境空气中粉尘浓度和接触者的体格检查

\ContentBand{对应往年题}

\begin{pastquestion}

\QuestionLabel{【往年题1｜24年完整试题回忆.docx｜三、简答5】}\\
5、论述我国综合防尘治尘的''八字方针''的具体内涵

\end{pastquestion}

\PointSeparator

\section{水泥工厂职业有害因素的健康教育与改进措施}\label{users__fpb__26_spring__pkusph_study__tmp__pdfs__prevention-medicine-past-exams__01-ux804cux4e1aux536bux751fux5b66ux5f80ux5e74ux9898ux8003ux70b9ux6574ux7406.md__ux6c34ux6ce5ux5de5ux5382ux804cux4e1aux6709ux5bb3ux56e0ux7d20ux7684ux5065ux5eb7ux6559ux80b2ux4e0eux6539ux8fdbux63aaux65bd}

\FrequencyPill{1 次}

\InfoLabel{对应小节}
综合（第三章不良气象、第四章噪声振动、第六章粉尘------水泥行业涉及粉尘、高温、噪声、振动）\\
\InfoLabel{匹配依据} 10 论述题原题（题干并附标准答案要点：粉尘、高温、噪声、振动）。

\ContentBand{知识点原文摘取}

（详见各对应章节关于粉尘控制------五、粉尘危害的控制管理；高温防暑降温------一、不良气象条件主要内容 /
5、(防暑降温措施)；噪声防范------四、噪声防范措施；振动职业危害预防------五、振动 /
6、振动职业危害的预防措施）

\ContentBand{对应往年题}

\begin{pastquestion}

\QuestionLabel{【往年题1｜10级职卫考题回忆.docx｜大题】}\\
对讨论课中的水泥工厂的职业有害因素进行健康教育和改进措施。

\end{pastquestion}

\begin{pastquestion}

\QuestionLabel{【往年题2｜04-14级职业卫生论述往年题汇总.docx｜10级】}\\
对讨论课中的水泥工厂的职业有害因素进行健康教育和改进措施。\\
粉尘、高温、噪声、振动

\end{pastquestion}

\PointSeparator

\chapter{高分子化合物}\label{users__fpb__26_spring__pkusph_study__tmp__pdfs__prevention-medicine-past-exams__01-ux804cux4e1aux536bux751fux5b66ux5f80ux5e74ux9898ux8003ux70b9ux6574ux7406.md__ux7b2cux4e03ux7ae0-ux9ad8ux5206ux5b50ux5316ux5408ux7269}

\section{高分子化合物的概念、特点与单体毒性}\label{users__fpb__26_spring__pkusph_study__tmp__pdfs__prevention-medicine-past-exams__01-ux804cux4e1aux536bux751fux5b66ux5f80ux5e74ux9898ux8003ux70b9ux6574ux7406.md__ux9ad8ux5206ux5b50ux5316ux5408ux7269ux7684ux6982ux5ff5ux7279ux70b9ux4e0eux5355ux4f53ux6bd2ux6027}

\FrequencyPill{3 次}

\InfoLabel{对应小节} 一、高分子化合物概述；三、3-4 高分子化合物对健康的影响\\
\InfoLabel{匹配依据} 11 简答''高分子材料生产过程中可能的健康损伤''；24 单选
10''关于高分子化合物的说法错误的是''；14 单选''氯丙烯引起的肿瘤是哪种''（涉及单体毒性，疑似对应于氯乙烯）。

\ContentBand{知识点原文摘取}

1、概念：高分子化合物是由一种或几种单体，经聚合或缩聚而成、分子量达数千至数百万的化合物。\\
·聚合：许多单体连接起来形成高分子化合物的过程。在此过程中，不析出任何副产品。\\
·缩聚：单体间先缩合析出一个分子水、氨、氯化氢或醇，再聚合成高分子化合物。\\
2、特点：高强度、耐腐蚀、绝缘性能好、成品无毒或毒性很小。\\
3、种类：塑料(plastic)；合成纤维(synthetic fiber)；合成橡胶(synthetic
rubber)；涂料(caotings)；粘胶剂(adhesive)

1、概述\\
·高分子化合物本身无毒或毒性很少\\
·在高分子化合物生产过程的每个阶段，作业者均可接触到不同类型的毒物

4、单体毒性（所有单体对人体均有毒性）\\
·石油裂解气：乙烯、丙烯、丁二烯；苯系化合物\\
·使用单体：乙烯、丙烯、氯乙烯、苯乙稀、丙烯腈\\
·在体内的代谢方式大致相同，被氧化成为环氧化物，与核酸结合，致癌和致突变。

6、热解产物毒性：高分子化合物与空气中氧接触，并受热、紫外线照射及机械作用，可被氧化。除了释放单体外(未聚合)，同时产生裂解产物CO、HCI、光气等。

\ContentBand{对应往年题}

\begin{pastquestion}

\QuestionLabel{【往年题1｜11级职卫考题.docx｜简答3】}\\
3 高分子材料生产过程中可能的健康损伤

\end{pastquestion}

\begin{pastquestion}

\QuestionLabel{【往年题2｜14级职卫考题.pdf｜单选（回忆）】}\\
氯丙烯引起的肿瘤是哪种

\end{pastquestion}

\begin{pastquestion}

\QuestionLabel{【往年题3｜24年完整试题回忆.docx｜一、单选10】}\\
10、关于高分子化合物的说法错误的是\\
A、高分子化合物本身无毒或毒性很少\\
B、在高分子化合物生产过程的每个阶段，作业者均可接触到不同类型的毒物\\
C、生产高分子化合物时，助剂对接触者影响较大\\
D、单体在体内的代谢方式大致相同，被氧化成为环氧化物\\
E、高分子化合物与空气中氧接触，并受热、紫外线照射及机械作用，可被氧化，产生裂解产物

\end{pastquestion}

\PointSeparator

\section{氯乙烯（VC）---肢端溶骨症 /
肝血管肉瘤}\label{users__fpb__26_spring__pkusph_study__tmp__pdfs__prevention-medicine-past-exams__01-ux804cux4e1aux536bux751fux5b66ux5f80ux5e74ux9898ux8003ux70b9ux6574ux7406.md__ux6c2fux4e59ux70efvcux80a2ux7aefux6eb6ux9aa8ux75c7-ux809dux8840ux7ba1ux8089ux7624}

\FrequencyPill{1 次}

\InfoLabel{对应小节} 四、常见高分子化合物 / 1、氯乙烯（VC）（4）毒性作用（慢性毒性）\\
\InfoLabel{匹配依据} 14 单选''氯丙烯引起的肿瘤是哪种''（疑似对应氯乙烯所致肝血管肉瘤）；并见 04-14
级论述汇总附录中''肢端溶骨症''等定义为名解辅助考点。

\ContentBand{知识点原文摘取}

-慢性毒性：氯乙烯长期接触，对人体的健康可产生多系统不同程度的影响，我们通常将这些症状称为氯乙烯病或氯乙烯综合症。神经系统；消化系统；肝血管肉瘤(hepatic
angiosarcoma)；肢端溶骨症(acroosteolysis, AOL)；血液系统；皮肤

\ContentBand{对应往年题}

\begin{pastquestion}

\QuestionLabel{【往年题1｜04-14级职业卫生论述往年题汇总.docx｜14级附录】}\\
肢端溶骨症
acroosteolysis，AOL：氯乙烯引起机体全身性改变在指端局部的一种表现。特点为末节指骨骨质溶解性损害，早期为手指麻木疼痛肿胀等雷诺综合征。

\end{pastquestion}

\PointSeparator

\section{氟聚合物烟尘热}\label{users__fpb__26_spring__pkusph_study__tmp__pdfs__prevention-medicine-past-exams__01-ux804cux4e1aux536bux751fux5b66ux5f80ux5e74ux9898ux8003ux70b9ux6574ux7406.md__ux6c1fux805aux5408ux7269ux70dfux5c18ux70ed}

\FrequencyPill{1 次}

\InfoLabel{对应小节} 四、3、含氟塑料（5）临床表现 / 氟聚合物烟尘热\\
\InfoLabel{匹配依据} 04-14 级论述汇总附录列出''氟聚合物烟尘热 fluropolymer fume fever''。疑似对应。

\ContentBand{知识点原文摘取}

·氟聚合物烟尘热：常于聚四氟乙烯、聚六氟丙烯热加工成型吸入热解物所致。\\
-临床表现：①畏寒发热、乏力、头昏、头痛、恶心、呕吐、眼及咽喉干燥、胸闷等；②体温多37.5-39.5℃，一般12-48h
消退；③眼及咽部充血或扁桃体肿大，白细胞总数及中性白细胞增高；1-2d 自愈。

\ContentBand{对应往年题}

\begin{pastquestion}

\QuestionLabel{【往年题1｜04-14级职业卫生论述往年题汇总.docx｜14级附录】}\\
氟聚合物烟尘热fluropolymer fume
fever：聚四氟乙烯、聚全氟丙烯热加工成型过程中，作业工人吸入热解物微粒所致，引起以发热为主要表现的疾病

\end{pastquestion}

\PointSeparator

\chapter{重金属中毒与职业心理学}\label{users__fpb__26_spring__pkusph_study__tmp__pdfs__prevention-medicine-past-exams__01-ux804cux4e1aux536bux751fux5b66ux5f80ux5e74ux9898ux8003ux70b9ux6574ux7406.md__ux7b2cux516bux7ae0-ux91cdux91d1ux5c5eux4e2dux6bd2ux4e0eux804cux4e1aux5fc3ux7406ux5b66}

\section{慢性铅中毒：临床表现、检测指标、防治、女工铅中毒}\label{users__fpb__26_spring__pkusph_study__tmp__pdfs__prevention-medicine-past-exams__01-ux804cux4e1aux536bux751fux5b66ux5f80ux5e74ux9898ux8003ux70b9ux6574ux7406.md__ux6162ux6027ux94c5ux4e2dux6bd2ux4e34ux5e8aux8868ux73b0ux68c0ux6d4bux6307ux6807ux9632ux6cbbux5973ux5de5ux94c5ux4e2dux6bd2}

\FrequencyPill{5 次}

\InfoLabel{对应小节} 第八章一、（三）常见重金属职业中毒表现 / 1、铅中毒 +（连同第十一章 6、铅）\\
\InfoLabel{匹配依据} 10 论述''女性职工铅中毒的临床表现''；14 论述''焦炉场职业卫生调查方案（铅等重金属）``；24
单选 14''哪一项不作为铅的检测指标''；24 单选 16''经乳汁排出引起婴儿中毒（铅汞镉）``；24 简答 6''防治铅中毒''。

\ContentBand{知识点原文摘取}

1、铅中毒\\
（3）毒理\\
·呼吸道是主要途径，其次是消化道，但四乙基铅可经皮肤和粘膜吸收\\
·血循环中铅早期主要分布于肝、肾、脑、皮肤和骨骼肌中，数周后，由软组织转移于骨\\
·人体内90\%-95\%的铅储存于骨内（{[}Pb3(PO4)2{]}，稳定）\\
·主要通过肾脏排出：体内铅排出缓慢，半减期5-10 年\\
·铅作业工人接触指标：血铅、尿铅；血铅水平用于职业性铅中毒诊断依据\\
·铅作用于全身各器官和系统，主要累及造血系统、神经系统、消化系统、血管及肾脏\\
·血铅可通过胎盘进入胎儿体内影响子代；乳汁内铅可影响婴儿

（5）慢性铅中毒临床表现\\
·神经系统：类神经症：头昏、头痛、乏力、失眠、多梦、记忆力减退等\\
周围神经病：感觉型（肢端麻木）、运动型（铅麻痹、垂腕、垂足）、混合型\\
严重者出现中毒性脑病\\
·消化系统：食欲不振、恶心、腹胀、便秘/腹泻、严重者腹绞痛，在脐周部位，一般镇痛药无效。\\
·血液及造血系统：贫血,点彩和网织红细胞增多。\\
·其他：齿龈蓝黑色''铅线''；肾脏损害；女工不孕、畸胎等\\
·铅为可疑人类致癌物

（8）治疗\\
·治疗原则：中毒患者应根据具体情况，使用金属络合剂驱铅治疗，如依地酸二钠钙、二巯丁二酸钠等注射，或二巯丁二酸口服，辅以对症治疗。

6、铅\\
·血铅含量只反映机体近期（1 个月左右）铅暴露的情况。\\
·骨铅水平才能反映较长时间的慢性铅暴露情况，只有在妊娠期和绝经期，骨铅才大量释放。\\
（1）铅作业对女性生殖功能的影响\\
·对月经的影响：周期延长，经量减少，经期缩短\\
·对妊娠经过和妊娠结局的影响：Pb
作业女工妊娠并发症---先兆流产、妊娠高血压综合征。\ldots\ldots 早产率、低出生体重儿率增高，Pb
作业女工子代智商低于对照组儿童。\\
·经胎盘转运：\ldots\ldots{}\\
·母乳排Pb：Pb 可经乳汁排出\\
·母源性小儿Pb 中毒：Pb 经母体进入乳儿体内，可引起乳儿Pb 中毒。

\ContentBand{对应往年题}

\begin{pastquestion}

\QuestionLabel{【往年题1｜10级职卫考题回忆.docx｜大题】}\\
女性职工铅中毒的临床表现。

\end{pastquestion}

\begin{pastquestion}

\QuestionLabel{【往年题2｜14级职卫考题.pdf｜论述】}\\
假设要对某焦炉场进行职业卫生调查（已知该工厂主要的职业有害因素是铅等重金属），请设计职业卫生调查方案，包括：调查的内容、监测的内容、监测的指标、如何进行危害评价、提出健康促进的建议

\end{pastquestion}

\begin{pastquestion}

\QuestionLabel{【往年题3｜24年完整试题回忆.docx｜一、单选14】}\\
14、下列哪一项不作为铅的检测指标\\
A、血铅\\
B、尿铅\\
C、骨铅\\
D、发铅\\
E、尿δ-ALA

\end{pastquestion}

\begin{pastquestion}

\QuestionLabel{【往年题4｜24年完整试题回忆.docx｜一、单选16】}\\
16、能经乳汁排出并可能引起婴儿中毒的生产性毒物有\\
A、铅、汞、镉\\
B、石棉、二氧化硅\\
C、乙烯、丙烯\\
D、氰化钠、氰化钾\\
E、氯气、溴化氢

\end{pastquestion}

\begin{pastquestion}

\QuestionLabel{【往年题5｜24年完整试题回忆.docx｜三、简答6】}\\
6、简述如何防治铅中毒

\end{pastquestion}

\PointSeparator

\section{职业紧张（work stress）/ 职业倦怠 /
精疲力竭}\label{users__fpb__26_spring__pkusph_study__tmp__pdfs__prevention-medicine-past-exams__01-ux804cux4e1aux536bux751fux5b66ux5f80ux5e74ux9898ux8003ux70b9ux6574ux7406.md__ux804cux4e1aux7d27ux5f20work-stress-ux804cux4e1aux5026ux6020-ux7cbeux75b2ux529bux7aed}

\FrequencyPill{3 次}

\InfoLabel{对应小节} 二、职业心理与职业紧张（一）概念；（三）职业紧张 /
4、职业紧张反应的表现；（4）精疲力竭症\\
\InfoLabel{匹配依据} 21 论述''基于医护人员高强度工作，谈谈怎么减少职业倦怠和职业紧张''；24 名解 work
stress；并间接出现于 14 级附录定义''职业紧张''。

\ContentBand{知识点原文摘取}

·职业紧张(occupational stress)/工作紧张(work
stress)：是指在某种职业条件下，客观需求与个人适应能力之间的失衡所带来的生理和心理压力

·根据发生的时间特点不同，分为：急性紧张反应、创伤后紧张反应、慢性紧张反应\\
·职业紧张引起的特征性问题：精疲力竭症、过劳死、伤亡事故

（4）精疲力竭症：精疲力竭的发生是职业紧张的直接后果，是个体不能应对职业紧张的最重要的表现之一\\
·Maslach 提出精疲力竭症三维模式，主要内容是：①情绪耗竭(emotional
exhaustion)：\ldots\ldots；②人格解体(depersonalization)：\ldots\ldots；③职业效能下降：\ldots\ldots{}

5、职业紧张的控制和干预：预防职业紧张首先应探寻和确定紧张源，可从个人和组织两个方面采取干预措施。对个体应增强应对能力，对组织则努力消除紧张源。但无论从哪方面干预，都需要采取综合性措施。\\
·增强个体应对能力：社会支持\\
·增强个体应对反应：预估工作中可能出现的困难并做一定准备\\
·组织措施：工作方式和劳动组织符合卫生学要求\\
·培训和教育：有针对性进行职业指导和培训，适应职业环境\\
·法律保障：\ldots\ldots{}\\
·健康促进：开展健康教育和健康促进活动，增强个体应对职业紧张的能力

\ContentBand{对应往年题}

\begin{pastquestion}

\QuestionLabel{【往年题1｜21级本科预防二班+胡润泽+职业卫生+主观题总结.docx｜论述2】}\\
2、基于医护人员高强度工作，谈谈怎么减少职业倦怠和职业紧张

\end{pastquestion}

\begin{pastquestion}

\QuestionLabel{【往年题2｜24年完整试题回忆.docx｜二、名解6】}\\
6、work stress

\end{pastquestion}

\begin{pastquestion}

\QuestionLabel{【往年题3｜04-14级职业卫生论述往年题汇总.docx｜14级附录】}\\
职业紧张：在某种职业条件下，客观需求与个人适应能力之间的失衡所带来的生理和心理压力。

\end{pastquestion}

\PointSeparator

\section{心身疾病（psychosomatic
diseases）}\label{users__fpb__26_spring__pkusph_study__tmp__pdfs__prevention-medicine-past-exams__01-ux804cux4e1aux536bux751fux5b66ux5f80ux5e74ux9898ux8003ux70b9ux6574ux7406.md__ux5fc3ux8eabux75beux75c5psychosomatic-diseases}

\FrequencyPill{3 次}

\InfoLabel{对应小节} （四）心身疾病 / 1、疾病分类；2、常见的心身疾病\\
\InfoLabel{匹配依据} 14 单选''不属于心身疾病的是''；21 名解 psychosomatic diseases；并见 04-14
级论述汇总附录中''心身疾病''定义。

\ContentBand{知识点原文摘取}

（四）心身疾病(psychosomatic
diseases)/心理生理障碍：指一组与心理和社会因素密切相关，但以躯体症状表现为主的疾病。心身疾病的范围甚为广泛，可以累及人体的各个器官和系统。目前包括了由情绪因素所引起的，以躯体症状为主要表现，受自主神经所支配的系统或器官的多种疾病。

2、常见的心身疾病：支气管哮喘；消化性溃疡；原发性高血压；癌症；甲状腺功能亢进

\ContentBand{对应往年题}

\begin{pastquestion}

\QuestionLabel{【往年题1｜14级职卫考题.pdf｜单选（回忆）】}\\
不属于心身疾病的是

\end{pastquestion}

\begin{pastquestion}

\QuestionLabel{【往年题2｜21级本科预防二班+胡润泽+职业卫生+主观题总结.docx｜名词解释】}\\
psychosomatic diseases

\end{pastquestion}

\begin{pastquestion}

\QuestionLabel{【往年题3｜04-14级职业卫生论述往年题汇总.docx｜14级附录】}\\
心身疾病：又称心理生理障碍，是指一组与心理和社会因素密切相关，由情绪因素引起但以躯体症状表现为主的疾病

\end{pastquestion}

\PointSeparator

\section{职业性金属（与类金属）中毒的特点}\label{users__fpb__26_spring__pkusph_study__tmp__pdfs__prevention-medicine-past-exams__01-ux804cux4e1aux536bux751fux5b66ux5f80ux5e74ux9898ux8003ux70b9ux6574ux7406.md__ux804cux4e1aux6027ux91d1ux5c5eux4e0eux7c7bux91d1ux5c5eux4e2dux6bd2ux7684ux7279ux70b9}

\FrequencyPill{1 次}

\InfoLabel{对应小节} （五）职业性金属中毒的主要特点\\
\InfoLabel{匹配依据} 21 简答''职业性金属中毒的特点''。

\ContentBand{知识点原文摘取}

（五）职业性金属中毒的主要特点\\
·中毒对象：有职业接触史的职业人群（涉及行业广泛）\\
·生物转运：呼吸道吸入为主，可随血液转运分布于全身各组织器官；不易被组织代谢降解，易在体内蓄积（达到临界作用水平）导致慢性毒作用，但不同金属有不同蓄积器官和靶器官；主要经肾脏排泄，但排泄速率各有不同，除汞外生物半衰期都较长\\
·中毒机制：金属与体内巯基、配基结合，使蛋白质结构改变从而丧失功能\\
·健康影响：毒作用表现形式多样、累及范围广，可致多器官损伤；毒性与接触剂量、化合物结构、金属氧化状态有关；主要是在体内蓄积导致慢性中毒，不同金属因其靶器官和毒性不同而出现不同的临床表现

\ContentBand{对应往年题}

\begin{pastquestion}

\QuestionLabel{【往年题1｜21级本科预防二班+胡润泽+职业卫生+主观题总结.docx｜简答4】}\\
4、职业性金属中毒的特点

\end{pastquestion}

\PointSeparator

\section{慢性汞中毒的典型症状}\label{users__fpb__26_spring__pkusph_study__tmp__pdfs__prevention-medicine-past-exams__01-ux804cux4e1aux536bux751fux5b66ux5f80ux5e74ux9898ux8003ux70b9ux6574ux7406.md__ux6162ux6027ux6c5eux4e2dux6bd2ux7684ux5178ux578bux75c7ux72b6}

\FrequencyPill{1 次}

\InfoLabel{对应小节} （三）2、慢性汞中毒\\
\InfoLabel{匹配依据} 14 单选''慢性汞中毒的三大典型症状''。

\ContentBand{知识点原文摘取}

2、慢性汞中毒：液态汞、挥发性\\
（2）临床表现\\
·易兴奋征：精神异常是慢性汞中毒的重要特征之一\\
·震颤：手指、眼睑、唇、舌及四肢等意向性震颤\\
·口腔-牙龈炎：流涎多、牙龈肿胀、易出血、牙齿松动甚至脱落、有的齿龈''汞线''.\\
·少数患者肾损伤（汞在肾脏蓄积）

\ContentBand{对应往年题}

\begin{pastquestion}

\QuestionLabel{【往年题1｜14级职卫考题.pdf｜单选（回忆）】}\\
慢性汞中毒的三大典型症状

\end{pastquestion}

\PointSeparator

\section{夜班作业 /
职业心理学（综合）}\label{users__fpb__26_spring__pkusph_study__tmp__pdfs__prevention-medicine-past-exams__01-ux804cux4e1aux536bux751fux5b66ux5f80ux5e74ux9898ux8003ux70b9ux6574ux7406.md__ux591cux73edux4f5cux4e1a-ux804cux4e1aux5fc3ux7406ux5b66ux7efcux5408}

\FrequencyPill{1 次}

\InfoLabel{对应小节} 二、职业心理与职业紧张 /（一）概念；（二）1、与职业有关的心理因素 / 夜班作业\\
\InfoLabel{匹配依据} 04-14 级论述汇总附录列出''职业心理学''\,``夜班作业''作为辅助名解考点。疑似对应。

\ContentBand{知识点原文摘取}

·职业心理学(Occupational
psychology)：是从人与职业的心理与社会环境关系角度研究人在职业过程中心理活动的特点和规律的学科，是应用心理学的一个分支学科。

-夜班作业(night
work)：指在一天中通常在睡眠的这段时间里进行的职业活动，是轮班劳动中对劳动者身心影响最大的作业。

\ContentBand{对应往年题}

\begin{pastquestion}

\QuestionLabel{【往年题1｜04-14级职业卫生论述往年题汇总.docx｜14级附录】}\\
职业心理学：从人与职业的心理与社会环境关系角度，研究人在职业过程中心理活动的特点和规律的学科。夜班作业 night
work：通常在睡眠的这段时间里进行的职业活动，是轮班劳动中对劳动者身心影响最大的作业。

\end{pastquestion}

\PointSeparator

\chapter{有机溶剂}\label{users__fpb__26_spring__pkusph_study__tmp__pdfs__prevention-medicine-past-exams__01-ux804cux4e1aux536bux751fux5b66ux5f80ux5e74ux9898ux8003ux70b9ux6574ux7406.md__ux7b2cux4e5dux7ae0-ux6709ux673aux6eb6ux5242}

\section{苯：急性中毒靶器官与对血液的影响}\label{users__fpb__26_spring__pkusph_study__tmp__pdfs__prevention-medicine-past-exams__01-ux804cux4e1aux536bux751fux5b66ux5f80ux5e74ux9898ux8003ux70b9ux6574ux7406.md__ux82efux6025ux6027ux4e2dux6bd2ux9776ux5668ux5b98ux4e0eux5bf9ux8840ux6db2ux7684ux5f71ux54cd}

\FrequencyPill{2 次}

\InfoLabel{对应小节} 四、常见有机溶剂职业中毒表现 / 1、苯\\
\InfoLabel{匹配依据} 14 单选''苯急性中毒的靶器官''；11 简答''苯和氨基苯对血液的影响''（其中苯部分）。

\ContentBand{知识点原文摘取}

1、苯\\
（1）急性中毒（短时高剂量）\\
·中枢神经系统麻醉作用\\
·实验室检查：尿酚↑ 、血苯↑\\
（2）慢性中毒（长期低剂量）：主要损害造血系统------白细胞系统↓→血小板↓→全血系统↓\\
·白细胞系统↓：中性粒细胞↓，淋巴细胞相对值↑；白细胞有较多的毒性颗粒、空泡、破碎细胞等\\
·血小板↓：出血倾向↑，女性月经过多；血小板形态异常\\
·全血系统↓：再生障碍性贫血，骨髓增生异常综合征(MDS)，白血病(急性粒细胞白血病)\\
·IARC 确认苯为人类致癌物

\ContentBand{对应往年题}

\begin{pastquestion}

\QuestionLabel{【往年题1｜11级职卫考题.docx｜简答2】}\\
2 苯和铵苯对血液的影响

\end{pastquestion}

\begin{pastquestion}

\QuestionLabel{【往年题2｜14级职卫考题.pdf｜单选（回忆）】}\\
苯急性中毒的靶器官

\end{pastquestion}

\PointSeparator

\section{苯的氨基化合物（苯胺）---高铁血红蛋白血症 /
化学性发绀}\label{users__fpb__26_spring__pkusph_study__tmp__pdfs__prevention-medicine-past-exams__01-ux804cux4e1aux536bux751fux5b66ux5f80ux5e74ux9898ux8003ux70b9ux6574ux7406.md__ux82efux7684ux6c28ux57faux5316ux5408ux7269ux82efux80faux9ad8ux94c1ux8840ux7ea2ux86cbux767dux8840ux75c7-ux5316ux5b66ux6027ux53d1ux7ec0}

\FrequencyPill{2 次}

\InfoLabel{对应小节} 六、苯的氨基和硝基化合物 / 1、毒作用特点\\
\InfoLabel{匹配依据} 11 简答''苯和铵苯对血液的影响''（氨基苯部分）；24 单选
12''工人出现头晕、头痛、口唇发绀，可能是因为吸入了\ldots 苯胺''。

\ContentBand{知识点原文摘取}

1、毒作用特点\\
·血液损害：（苯胺）\\
-形成高铁血红蛋白，缺氧，化学性发绀：少数为直接作用，大多为间接作用（代谢中间物）\\
-溶血作用：间接毒性作用。破坏了红细胞膜的稳定性；红细胞内形成赫恩小体\\
·肝脏损害：1）直接损害：中毒性肝病；2）间接损害：溶血作用导致的继发性肝损伤\\
·晶体损害：中毒性白内障（三硝基甲苯）\\
·泌尿系统损害：直接作用、溶血导致的间接损伤\\
·神经系统损害：神经细胞脂肪变性、神经炎\\
·皮肤损害和致敏作用：接触性皮炎、过敏性皮炎\\
·致癌作用：职业性膀胱癌

\ContentBand{对应往年题}

\begin{pastquestion}

\QuestionLabel{【往年题1｜11级职卫考题.docx｜简答2】}\\
2 苯和铵苯对血液的影响

\end{pastquestion}

\begin{pastquestion}

\QuestionLabel{【往年题2｜24年完整试题回忆.docx｜一、单选12】}\\
12、工人出现头晕、头痛、口唇发绀，可能是因为吸入了\\
A、苯\\
B、二硫化碳\\
C、二氯乙烷\\
D、正己烷\\
E、苯胺

\end{pastquestion}

\PointSeparator

\section{有机溶剂的毒作用特点（脂溶性 /
神经毒性）}\label{users__fpb__26_spring__pkusph_study__tmp__pdfs__prevention-medicine-past-exams__01-ux804cux4e1aux536bux751fux5b66ux5f80ux5e74ux9898ux8003ux70b9ux6574ux7406.md__ux6709ux673aux6eb6ux5242ux7684ux6bd2ux4f5cux7528ux7279ux70b9ux8102ux6eb6ux6027-ux795eux7ecfux6bd2ux6027}

\FrequencyPill{1 次}

\InfoLabel{对应小节} 二、有机溶剂的理化特性与毒作用特点 / 2、毒作用特点（1）与理化特性密切相关\\
\InfoLabel{匹配依据} 24 单选 11''有机溶剂的神经毒性与哪一个理化特性密切相关（脂溶性）``。

\ContentBand{知识点原文摘取}

·挥发性、可溶性、易燃性：①易挥发，接触途径以吸入为主；②脂溶性是有机溶剂的重要特性，进入体内易与富含脂质的神经组织亲和而具有麻醉作用；又具一定的水溶性，易经皮肤进入体内③大多具有可燃性，如汽油、乙醇等，可用作燃料；有些属非可燃物而用作灭火剂，如卤代烃类化合物

\ContentBand{对应往年题}

\begin{pastquestion}

\QuestionLabel{【往年题1｜24年完整试题回忆.docx｜一、单选11】}\\
11、有机溶剂的神经毒性与哪一个理化特性密切相关\\
A、挥发性\\
B、脂溶性\\
C、可燃性\\
D、化学结构\\
E、\ldots\ldots{}

\end{pastquestion}

\PointSeparator

\section{甲苯/二甲苯（``三强一弱''）与赫恩小体}\label{users__fpb__26_spring__pkusph_study__tmp__pdfs__prevention-medicine-past-exams__01-ux804cux4e1aux536bux751fux5b66ux5f80ux5e74ux9898ux8003ux70b9ux6574ux7406.md__ux7532ux82efux4e8cux7532ux82efux4e09ux5f3aux4e00ux5f31ux4e0eux8d6bux6069ux5c0fux4f53}

\FrequencyPill{1 次}

\InfoLabel{对应小节} 四、2、甲苯、二甲苯；六、1、毒作用特点（赫恩小体）\\
\InfoLabel{匹配依据} 04-14
级论述汇总附录列出''甲苯相比于苯，呈现'三强一弱'\,``以及''赫恩小体''作为名解辅助考点。疑似对应。

\ContentBand{知识点原文摘取}

2、甲苯、二甲苯\\
·临床表现：相比于苯，呈现''三强一弱''的特点\\
-急性毒作用强，中枢神经系统麻醉作用强，皮肤黏膜刺激作用强，对血液系统影响弱\\
-合并心、肾、肝、肺等器官的损害

-溶血作用：间接毒性作用。破坏了红细胞膜的稳定性；红细胞内形成赫恩小体

\ContentBand{对应往年题}

\begin{pastquestion}

\QuestionLabel{【往年题1｜04-14级职业卫生论述往年题汇总.docx｜14级附录】}\\
甲苯相比于苯，呈现''三强一弱''的特点。急性毒作用强，中枢神经系统麻醉作用强，皮肤黏膜刺激作用强，对血液系统影响弱，合并心、肾、肝、肺等器官的损害。\\
赫恩小体

\end{pastquestion}

\PointSeparator

\chapter{农药中毒}\label{users__fpb__26_spring__pkusph_study__tmp__pdfs__prevention-medicine-past-exams__01-ux804cux4e1aux536bux751fux5b66ux5f80ux5e74ux9898ux8003ux70b9ux6574ux7406.md__ux7b2cux5341ux7ae0-ux519cux836fux4e2dux6bd2}

\section{有机磷农药与氨基甲酸酯类农药中毒（机制、临床表现及二者异同）}\label{users__fpb__26_spring__pkusph_study__tmp__pdfs__prevention-medicine-past-exams__01-ux804cux4e1aux536bux751fux5b66ux5f80ux5e74ux9898ux8003ux70b9ux6574ux7406.md__ux6709ux673aux78f7ux519cux836fux4e0eux6c28ux57faux7532ux9178ux916fux7c7bux519cux836fux4e2dux6bd2ux673aux5236ux4e34ux5e8aux8868ux73b0ux53caux4e8cux8005ux5f02ux540c}

\FrequencyPill{2 次}

\InfoLabel{对应小节} 二、有机磷酸酯类农药 / 6、靶器官毒性表现及中毒机理；7、临床表现；三、氨基甲酸酯类农药 /
6、有机磷农药与氨基甲酸酯类农药的区别\\
\InfoLabel{匹配依据} 10 简答''有机磷中毒的机制和临床表现''；14
简答''氨基甲酸酯类农药与有机磷类农药的毒性作用异同点''。

\ContentBand{知识点原文摘取}

6、靶器官毒性表现及中毒机理：有机磷化合物是神经毒物，乙酰胆碱酯酶是它的靶点。进入体内的有机磷化合物与乙酰胆碱酯酶有高度的亲和性，结合成磷酰化胆碱酯酶，乙酰胆碱酯酶失活\\
·毒性机制：抑制胆碱酯酶活性\\
·抑制胆碱酯酶活性，导致结果\\
-乙酰胆碱蓄积→毒蕈碱样症状、烟碱样症状和中枢神经系统症状表现。\\
-酶的老化：磷酰化胆碱酯酶不能重活化15-30 天

7、临床表现\\
（1）急性中毒：通常2～6 小时。\\
·毒蕈碱样症状：恶心，呕吐，多汗，流涎， 瞳孔缩小，严重者呼吸困难，肺水肿。\\
·烟碱样症状：肌肉震颤，痉挛（胸、面颊、 上肢）；呼吸肌麻痹死亡。\\
·中枢神经系统症状：头痛，头晕，多梦，失 眠，昏迷；一般无后遗症，迟发性周围神经
病；心肌毒性作用，发生''电击样''死亡；``癔 病样''发作。

6、有机磷农药与氨基甲酸酯类农药的区别\\
有机磷农药 / 氨基甲酸酯农药\\
磷酰基与CHE 结合，氧化速度↑，毒性越大↑ / 以分子与CHE 结合\\
可用复能剂CHE 复能 / 不用复能剂，CHE 抑制\\
CHE''酶老化''易引起严重中毒（磷酰基结合）/ CHE''易复能''，不容易引起严重中毒（分子结合）

7、急性有机磷中毒与氯基甲酸酯类农药中毒的鉴别要点\\
接触史：有机磷 / 氨基甲酸酯类农药\\
体表或呕吐物气味：蒜臭味 / 无蒜臭味\\
病程：相对较长 / 相对较短\\
阿托品用量：大 / 较小\\
复能剂治疗效果：显著 / 不明显\\
中毒胆碱酯酶活性：恢复缓慢 / 恢复较快

\ContentBand{对应往年题}

\begin{pastquestion}

\QuestionLabel{【往年题1｜10级职卫考题回忆.docx｜简答】}\\
物理因素的特点、有机磷中毒的机制和临床表现、李涛所长之劳动者职业健康监护制度内容、手臂振动病的定义,临床表现和典型特征、常见职业伤害事故类型

\end{pastquestion}

\begin{pastquestion}

\QuestionLabel{【往年题2｜14级职卫考题.pdf｜简答】}\\
简述氨基甲酸酯类农药与有机磷类农药的毒性作用异同点

\end{pastquestion}

\PointSeparator

\chapter{妇女职业卫生}\label{users__fpb__26_spring__pkusph_study__tmp__pdfs__prevention-medicine-past-exams__01-ux804cux4e1aux536bux751fux5b66ux5f80ux5e74ux9898ux8003ux70b9ux6574ux7406.md__ux7b2cux5341ux4e00ux7ae0-ux5987ux5973ux804cux4e1aux536bux751f}

\section{妇女职业卫生（定义/任务）}\label{users__fpb__26_spring__pkusph_study__tmp__pdfs__prevention-medicine-past-exams__01-ux804cux4e1aux536bux751fux5b66ux5f80ux5e74ux9898ux8003ux70b9ux6574ux7406.md__ux5987ux5973ux804cux4e1aux536bux751fux5b9aux4e49ux4efbux52a1}

\FrequencyPill{2 次}

\InfoLabel{对应小节} 第十一章总论（妇女职业卫生 Occupational Health for Women）\\
\InfoLabel{匹配依据} 10、11 名解''妇女职业卫生''。

\ContentBand{知识点原文摘取}

·妇女职业卫生：是研究劳动条件（或职业条件）对妇女健康，特别是对生殖健康的影响以及保护妇女劳动者健康的预防保健的科学，是职业卫生的重要组成部分，是对职业妇女进行劳动保护的理论基础。其基本任务在于防止职业有害因素对女职工健康尤其是生殖健康的不良影响，使妇女能够健康、持久、高效地从事生产劳动以及各种职业活动，且不影响子代健康。

\ContentBand{对应往年题}

\begin{pastquestion}

\QuestionLabel{【往年题1｜10级职卫考题回忆.docx｜名解】}\\
全英，工效学，妇女职业卫生，热痉挛，粉尘分散度，短时间接触容许浓度，氧需

\end{pastquestion}

\begin{pastquestion}

\QuestionLabel{【往年题2｜11级职卫考题.docx｜名解】}\\
妇女职业卫生、建设的''三同时''、AED、equal loudness curve、生物监测、occupational injury

\end{pastquestion}

\PointSeparator

\section{生殖损伤 / 生殖毒性 /
生殖系统损害}\label{users__fpb__26_spring__pkusph_study__tmp__pdfs__prevention-medicine-past-exams__01-ux804cux4e1aux536bux751fux5b66ux5f80ux5e74ux9898ux8003ux70b9ux6574ux7406.md__ux751fux6b96ux635fux4f24-ux751fux6b96ux6bd2ux6027-ux751fux6b96ux7cfbux7edfux635fux5bb3}

\FrequencyPill{3 次}

\InfoLabel{对应小节} 三、有害因素对女性生殖健康的影响 /
1、（1）生殖毒性；2、职业性有害因素对女性生殖机能的影响；3、标志生殖损伤的生殖结局\\
\InfoLabel{匹配依据} 04 简答''生殖损伤的表现''；09 简答''生殖系统损害表现''；并 04-14
级论述汇总附录列''生殖毒性''作为名解辅助考点。

\ContentBand{知识点原文摘取}

·生殖损伤(reproductive impairment)：生殖机能的全部过程中任何一个环节受到影响而造成生殖机能障碍或不良生殖结局。

1、职业有害因素的生殖毒性和发育毒性\\
（1）生殖毒性：主要指环境因素对生殖系统的不利影响，可表现为生殖器官及内分泌系统的变化，以及对性周期和性行为，生育力，妊娠结局、分娩、哺乳变化及生殖早衰。

2、职业性有害因素对女性生殖机能的影响\\
·受孕前：月经异常、性欲改变、不孕(受孕力)、生殖细胞突变\\
·妊娠时：母体方面---对毒物的敏感性增强、妊娠并发症增加\\
胎儿方面---胚胎或胎儿死亡、畸胎、生长迟缓、出生后功能发育不全、致癌\\
·分娩后：由于乳汁排毒或其他原因的暴露影响发育

3、标志生殖损伤的生殖结局：性欲改变、阳痿；低出生体重；不育；新生儿死亡；生育力下降；出生后发育障碍；自然流产；遗传病；死胎；婴儿期死亡率增高；畸胎；儿童期死亡率增高；性比改变；儿童期恶性肿瘤

\ContentBand{对应往年题}

\begin{pastquestion}

\QuestionLabel{【往年题1｜04级职卫考题回忆.doc｜简答】}\\
简答：职业流行病的作用、生殖损伤的表现、工作相关疾病定义及特点、职业肿瘤的特点

\end{pastquestion}

\begin{pastquestion}

\QuestionLabel{【往年题2｜09级职卫考题回忆.doc｜简答】}\\
简答：职业流行病的任务、生殖系统损害表现、中暑及表现、职业性肿瘤特点、 职业相关疾病及特点

\end{pastquestion}

\begin{pastquestion}

\QuestionLabel{【往年题3｜04-14级职业卫生论述往年题汇总.docx｜14级附录】}\\
生殖毒性：主要指环境因素对生殖系统的不利影响，可表现为生殖器官和内分泌系统，性周期和性行为，生育能力，妊娠结局、分娩、哺乳的变化及生殖早衰。

\end{pastquestion}

\PointSeparator

\chapter{职业肿瘤}\label{users__fpb__26_spring__pkusph_study__tmp__pdfs__prevention-medicine-past-exams__01-ux804cux4e1aux536bux751fux5b66ux5f80ux5e74ux9898ux8003ux70b9ux6574ux7406.md__ux7b2cux5341ux4e8cux7ae0-ux804cux4e1aux80bfux7624}

\section{职业性肿瘤（定义、特点）}\label{users__fpb__26_spring__pkusph_study__tmp__pdfs__prevention-medicine-past-exams__01-ux804cux4e1aux536bux751fux5b66ux5f80ux5e74ux9898ux8003ux70b9ux6574ux7406.md__ux804cux4e1aux6027ux80bfux7624ux5b9aux4e49ux7279ux70b9}

\FrequencyPill{5 次}

\InfoLabel{对应小节} 一、概述 / 2、定义；三、职业肿瘤的发病特点\\
\InfoLabel{匹配依据} 04、09 名解 + 简答''职业(性)肿瘤特点''；14 名解 Occupational tumor + 14
单选''职业肿瘤最常见的是哪个''。

\ContentBand{知识点原文摘取}

·职业性致癌因素(occupational
carcinogen)：与职业较长时间接触有关的，在一定条件下能使正常细胞转化为肿瘤细胞，且经过较长潜隐期能发展为可以检出肿瘤的致病因素。\\
-职业性致癌因素：机械刺激、物理因素、化学因素、生物因素（EB、HBV）\\
·职业肿瘤(Occupational
tumor)：在工作环境中长期接触职业致癌因素，经过较长的潜隐期而患某种特定肿瘤，职业癌(Occupational cancer)。

三、职业肿瘤的发病特点\\
1、病因明确：职业性致癌因素\\
2、好发部位较固定（接触部位）。\\
·皮肤\\
·呼吸道：远隔部位(联苯胺---膀胱癌)；一种接触致多种肿瘤(放射线白血病、肺癌、皮肤癌等)\\
3、潜隐期：发病年龄\\
·潜隐期较长，10 年以上。\\
·职业肿瘤的发病年龄较非职业性同类肿瘤病人年龄小。\\
4、病理类型：恶性度较高\\
·特定的病理类型。如氯甲醚肺癌大部分为未分化小细胞癌，\\
·苯：白血病的类型不同，无规律\\
5、发病条件：接触物质的理化性质、接触途径、接触剂量和接触时间、个人生活习惯、易感性

\ContentBand{对应往年题}

\begin{pastquestion}

\QuestionLabel{【往年题1｜04级职卫考题回忆.doc｜名解 + 简答】}\\
名解：工效学、TTS、职业肿瘤、氧债、热适应\\
简答：\ldots\ldots、职业肿瘤的特点

\end{pastquestion}

\begin{pastquestion}

\QuestionLabel{【往年题2｜09级职卫考题回忆.doc｜名解 + 简答】}\\
名解：全是英文，工效学、职业性肿瘤、氧债、生物监测、暂时性听阈位移\\
简答：\ldots\ldots、职业性肿瘤特点、 职业相关疾病及特点

\end{pastquestion}

\begin{pastquestion}

\QuestionLabel{【往年题3｜14级职卫考题.pdf｜名解】}\\
Occupational tumor

\end{pastquestion}

\begin{pastquestion}

\QuestionLabel{【往年题4｜14级职卫考题.pdf｜单选（回忆）】}\\
职业肿瘤最常见的是哪个

\end{pastquestion}

\begin{pastquestion}

\QuestionLabel{【往年题5｜04-14级职业卫生论述往年题汇总.docx｜14级附录】}\\
职业性致癌因素：与职业较长时间接触有关的，在一定条件下能使正常细胞转化为肿瘤细胞，且经过较长潜隐期能发展为可以检出肿瘤的致病因素。

\end{pastquestion}

\PointSeparator

\section{以石棉为例：职业性肿瘤（致癌因素）的识别与诊断}\label{users__fpb__26_spring__pkusph_study__tmp__pdfs__prevention-medicine-past-exams__01-ux804cux4e1aux536bux751fux5b66ux5f80ux5e74ux9898ux8003ux70b9ux6574ux7406.md__ux4ee5ux77f3ux68c9ux4e3aux4f8bux804cux4e1aux6027ux80bfux7624ux81f4ux764cux56e0ux7d20ux7684ux8bc6ux522bux4e0eux8bcaux65ad}

\FrequencyPill{4 次}

\InfoLabel{对应小节} 二、职业致癌因素的识别和确认；四、常见的职业肿瘤 / 1、职业性呼吸道肿瘤·石棉\\
\InfoLabel{匹配依据} 11、21 简答''以石棉为例说明职业性癌症的诊断（原则）``；14
简答''以石棉为例，简述职业致癌因素的识别和确认过程''；24 单选 13''目前比例最高的职业性肿瘤（石棉致肺癌）``。

\ContentBand{知识点原文摘取}

二、职业致癌因素的识别和确认\\
1、临床观察：发现职业性肿瘤------临床病例分析与观察是识别职业致癌因素的第一线索。\\
2、实验研究------动物诱癌实验(首选)；体外试验\\
3、人群流行病学调查\\
4、职业致癌物（职业性肿瘤）的判断方法（IARC 分组）

4、职业肿瘤诊断标准(GBZ94-2017)：肿瘤临床诊断明确{[}原发肿瘤,肿瘤发生部位与职业性致癌因素靶器官一致,组织病理检查(临床影像、内镜检查){]}；职业性致癌物接触史明确(接触时间\&状况)；潜隐期符合要求

·石棉：1934 年报道石棉可致肺癌，1955
年确认。大量的流行病学调查结果显示，肺癌是石棉工人死亡死因的20\%。潜伏期一般为20
年，石棉致肺癌作用的强弱及肺癌的种类与石棉的种类及纤维形态有关。石棉还可导致接触者的胸腹膜间皮瘤。

石棉肺：是硅酸盐尘肺中最常见、危害最严重的一种。（第六章）\\
-以青石棉致纤维化和致癌作用最强。

\ContentBand{对应往年题}

\begin{pastquestion}

\QuestionLabel{【往年题1｜11级职卫考题.docx｜简答5】}\\
5 以石棉为例说明职业性癌症的诊断

\end{pastquestion}

\begin{pastquestion}

\QuestionLabel{【往年题2｜14级职卫考题.pdf｜简答】}\\
以石棉为例，简述职业致癌因素的识别和确认过程

\end{pastquestion}

\begin{pastquestion}

\QuestionLabel{【往年题3｜21级本科预防二班+胡润泽+职业卫生+主观题总结.docx｜简答5】}\\
5、以石棉为例，简述职业性肿瘤的诊断原则，讨论

\end{pastquestion}

\begin{pastquestion}

\QuestionLabel{【往年题4｜24年完整试题回忆.docx｜一、单选13】}\\
13、目前比例最高的职业性肿瘤\\
A、苯所致白血病\\
B、氯乙烯致肝血管肉瘤\\
C、联苯胺致膀胱癌\\
D、石棉致肺癌\\
E、\ldots\ldots{}

\end{pastquestion}

\PointSeparator

\section{职业性肿瘤的诊断原则与预防原则}\label{users__fpb__26_spring__pkusph_study__tmp__pdfs__prevention-medicine-past-exams__01-ux804cux4e1aux536bux751fux5b66ux5f80ux5e74ux9898ux8003ux70b9ux6574ux7406.md__ux804cux4e1aux6027ux80bfux7624ux7684ux8bcaux65adux539fux5219ux4e0eux9884ux9632ux539fux5219}

\FrequencyPill{2 次}

\InfoLabel{对应小节} 一、4、职业肿瘤诊断标准；五、职业性肿瘤的预防原则\\
\InfoLabel{匹配依据} 24 简答''职业性肿瘤的诊断原则''；并见 04-14 级论述汇总附录列''职业肿瘤预防原则''。

\ContentBand{知识点原文摘取}

4、职业肿瘤诊断标准(GBZ94-2017)：肿瘤临床诊断明确{[}原发肿瘤,肿瘤发生部位与职业性致癌因素靶器官一致,组织病理检查(临床影像、内镜检查){]}；职业性致癌物接触史明确(接触时间\&状况)；潜隐期符合要求

五、职业性肿瘤的预防原则\\
1、严格执行职业卫生标准、法律：职业病防治法\\
2、建立并完善职业肿瘤监测体系（早发现、诊断、治疗）：健康档案、报告体系\\
3、加强对职业性致癌因素的控制和管理\\
（1）用非致癌物代替致癌物。\\
（2）改革操作工艺、减少直接接触机会。\\
（3）改善劳动条件，减少有害气体、粉尘在劳动场所的跑、冒、滴、漏，对于有致癌因素存在的生产环节，应采用密闭生产或有通风、回收装置，避免污染环境\\
（4）对致癌物采取严格管理措施，建立登记制度。\\
4、健康教育，提高自我防护能力\\
5、生活方式，合理膳食\\
6、建立预测制度\\
7、定期体检：定期体检，早期发现癌前病变

\ContentBand{对应往年题}

\begin{pastquestion}

\QuestionLabel{【往年题1｜24年完整试题回忆.docx｜三、简答2】}\\
2、职业性肿瘤的诊断原则

\end{pastquestion}

\begin{pastquestion}

\QuestionLabel{【往年题2｜04-14级职业卫生论述往年题汇总.docx｜14级附录】}\\
职业肿瘤预防原则：严格执行职业卫生标准、法律，建立并完善职业肿瘤监测体系，加强对职业性致癌因素的控制和管理，健康教育，提高自我防护能力，改善生活方式，合理膳食，建立预测制度，风险评估，定期体检

\end{pastquestion}

\PointSeparator

\section{我国应该重点控制的职业病及原因}\label{users__fpb__26_spring__pkusph_study__tmp__pdfs__prevention-medicine-past-exams__01-ux804cux4e1aux536bux751fux5b66ux5f80ux5e74ux9898ux8003ux70b9ux6574ux7406.md__ux6211ux56fdux5e94ux8be5ux91cdux70b9ux63a7ux5236ux7684ux804cux4e1aux75c5ux53caux539fux56e0}

\FrequencyPill{1 次}

\InfoLabel{对应小节} 综合（绪论 + 第六章尘肺病现状 + 本章肿瘤；亦含化学中毒、噪声等）\\
\InfoLabel{匹配依据} 11
应用''说明两个我国应该重点控制的职业病，原因''。疑似对应；与多个章节相关，但核心匹配第六章尘肺病现状（首位）+
本章职业肿瘤+ 化学中毒。

\ContentBand{知识点原文摘取}

2、尘肺病的危害现状：目前尘肺病仍是我国最严重的职业病，影响面最广、危害最严重，且已进入高发期；以煤工尘肺和矽肺为主；绝大多数的尘肺病例分布在煤炭行业；超过半数的病例分布在中、小型企业；尘肺病发病工龄有缩短的趋势；农民工成为受尘肺病危害的高危人群。

我国《职业病分类和目录》中职业性化学中毒60 项，其中有机溶剂中毒近30
项；有机溶剂种类多，在工业生产各个行业应用广泛，接触机会多

\ContentBand{对应往年题}

\begin{pastquestion}

\QuestionLabel{【往年题1｜11级职卫考题.docx｜应用1】}\\
1说明两个我国应该重点控制的职业病，原因

\end{pastquestion}

\PointSeparator

\chapter{职业流行病学}\label{users__fpb__26_spring__pkusph_study__tmp__pdfs__prevention-medicine-past-exams__01-ux804cux4e1aux536bux751fux5b66ux5f80ux5e74ux9898ux8003ux70b9ux6574ux7406.md__ux7b2cux5341ux4e09ux7ae0-ux804cux4e1aux6d41ux884cux75c5ux5b66}

\section{职业流行病的作用 /
任务（含利用职业人群开展流行病学研究的优势）}\label{users__fpb__26_spring__pkusph_study__tmp__pdfs__prevention-medicine-past-exams__01-ux804cux4e1aux536bux751fux5b66ux5f80ux5e74ux9898ux8003ux70b9ux6574ux7406.md__ux804cux4e1aux6d41ux884cux75c5ux7684ux4f5cux7528-ux4efbux52a1ux542bux5229ux7528ux804cux4e1aux4ebaux7fa4ux5f00ux5c55ux6d41ux884cux75c5ux5b66ux7814ux7a76ux7684ux4f18ux52bf}

\FrequencyPill{3 次}

\InfoLabel{对应小节} 一、职业流行病学的应用 / 5、利用职业人群开展流行病学研究优势\\
\InfoLabel{匹配依据} 04 简答''职业流行病的作用''；09 简答''职业流行病的任务''；24 单选
2''利用职业接触人群开展基础研究的优势不包括''。

\ContentBand{知识点原文摘取}

5、利用职业人群开展流行病学研究优势\\
·人群研究，暴露相对明确，可测量；\\
·健康监护，可以定期获得健康效应的资料；\\
·可以开展队列研究及干预研究；\\
·可以利用基础医学及生命科学最先进技术，开展早期健康效应研究：遗传、表观遗传及环境暴露、组学技术等等；\\
·人群研究，容易转化，产生政策，直接保护劳动者

6、具体应用举例\\
·暴露特征与致病机制：识别职业危害因素及其作用条件,进行接触人群健康风险评价、致病机制研究等；\\
·寻求职业危害因素接触水平-反应关系：为制订、修订职业卫生标准和职业病诊断标准提供依据，如有害因素的最高容许浓度、8
小时平均加权浓度等；\\
·职业环境监测及健康监护、防治策略研究：评价、制定及改进卫生干预措施\ldots\ldots；\\
·事故调查：突发安全事故的调查

\ContentBand{对应往年题}

\begin{pastquestion}

\QuestionLabel{【往年题1｜04级职卫考题回忆.doc｜简答】}\\
简答：职业流行病的作用、生殖损伤的表现、工作相关疾病定义及特点、职业肿瘤的特点

\end{pastquestion}

\begin{pastquestion}

\QuestionLabel{【往年题2｜09级职卫考题回忆.doc｜简答】}\\
简答：职业流行病的任务、生殖系统损害表现、中暑及表现、职业性肿瘤特点、 职业相关疾病及特点

\end{pastquestion}

\begin{pastquestion}

\QuestionLabel{【往年题3｜24年完整试题回忆.docx｜一、单选2】}\\
2、利用职业接触人群开展基础研究的优势不包括\\
A、人群研究，暴露相对明确\\
B、可以开展队列研究及干预研究\\
C、可以开展早期健康效应研究\\
D、容易转化，产生政策\\
E、选择人群偏倚相对较小

\end{pastquestion}

\PointSeparator

\section{健康工人效应（healthy worker
effect）}\label{users__fpb__26_spring__pkusph_study__tmp__pdfs__prevention-medicine-past-exams__01-ux804cux4e1aux536bux751fux5b66ux5f80ux5e74ux9898ux8003ux70b9ux6574ux7406.md__ux5065ux5eb7ux5de5ux4ebaux6548ux5e94healthy-worker-effect}

\FrequencyPill{1 次}

\InfoLabel{对应小节} 二、职业流行病学研究中的建议 /（9）职业暴露人群的偏性及来源·健康工人效应\\
\InfoLabel{匹配依据} 24 名解''健康工人效应''。

\ContentBand{知识点原文摘取}

·健康工人效应\\
-某岗位存在职业危害因素，受到影响的劳动者（易感者）调离了岗位，留在岗位的则是身体健康的，调查的结果，可能会得出随工作时间延长，工人越健康的结论；\\
-横断面调查中更容易遇到健康工人效应；\\
-尤其是致敏物或刺激性物质等以及上岗时为易感人群的，如三氯乙烯接触性皮炎、TDI 导致的哮喘。

\ContentBand{对应往年题}

\begin{pastquestion}

\QuestionLabel{【往年题1｜24年完整试题回忆.docx｜二、名解1】}\\
1、健康工人效应

\end{pastquestion}

\PointSeparator

\section{前瞻性队列研究观察期（暴露时间/潜伏期）的确定}\label{users__fpb__26_spring__pkusph_study__tmp__pdfs__prevention-medicine-past-exams__01-ux804cux4e1aux536bux751fux5b66ux5f80ux5e74ux9898ux8003ux70b9ux6574ux7406.md__ux524dux77bbux6027ux961fux5217ux7814ux7a76ux89c2ux5bdfux671fux66b4ux9732ux65f6ux95f4ux6f5cux4f0fux671fux7684ux786eux5b9a}

\FrequencyPill{1 次}

\InfoLabel{对应小节} 三、队列研究的部分问题 / 2、暴露时间的确定：职业病的潜伏期\\
\InfoLabel{匹配依据} 11 简答''前瞻性队列研究观察期长短确定原则''。

\ContentBand{知识点原文摘取}

2、暴露时间的确定：职业病的潜伏期\\
·例如：在研究癌症和职业暴露的关系时，从开始暴露到临床可检出癌症，需要一定潜伏期，短于该潜伏期而发生的癌症，一般认为与暴露无关。随访可以不从暴露的起点开始，可以在一个最短潜伏期之后开始\\
·但在实际工作中，如何来确定潜伏期是个难题。各自经验的不同，可能采用不同的潜伏值，可能会影响不同研究之间的可比性。\\
-在研究电焊工和肺癌的关系时，可以要求进入队列的工人至少要有5 年以上电焊工作业史，理由是基于如果电焊作业史不足5
年的工人患了肺癌，多认为与电焊作业无关的假设。因此，在队列研究开始时，首先需明确队列的定义；\\
-给暴露下定义时，还需要考虑暴露的累积影响，在相同累积暴露量的情况下，暴露时间的长短不同或开始暴露距观察终点的时间不同，都会影响暴露的作用。

\ContentBand{对应往年题}

\begin{pastquestion}

\QuestionLabel{【往年题1｜11级职卫考题.docx｜简答1】}\\
1前瞻性队列研究观察期长短确定原则

\end{pastquestion}

\PointSeparator

\chapter{职业卫生管理与实践}\label{users__fpb__26_spring__pkusph_study__tmp__pdfs__prevention-medicine-past-exams__01-ux804cux4e1aux536bux751fux5b66ux5f80ux5e74ux9898ux8003ux70b9ux6574ux7406.md__ux7b2cux5341ux56dbux7ae0-ux804cux4e1aux536bux751fux7ba1ux7406ux4e0eux5b9eux8df5}

\section{职业接触限值（OEL：PC-TWA / PC-STEL / MAC /
BEL）}\label{users__fpb__26_spring__pkusph_study__tmp__pdfs__prevention-medicine-past-exams__01-ux804cux4e1aux536bux751fux5b66ux5f80ux5e74ux9898ux8003ux70b9ux6574ux7406.md__ux804cux4e1aux63a5ux89e6ux9650ux503coelpc-twa-pc-stel-mac-bel}

\FrequencyPill{4 次}

\InfoLabel{对应小节} 第九章五、（6）；并参第十四章相关定义；BEL 见第十五章生物监测\\
\InfoLabel{匹配依据} 10 名解''短时间接触容许浓度''；14 名解''Occupation exposure limit''；21 名解''biological
exposure limit, BEL''；并见 04-14 论述汇总附录''OEL/PC-TWA/PC-STEL/MAC''。

\ContentBand{知识点原文摘取}

·职业接触限值occupational exposure limits,
OELs：职业性有害因素的接触限制量值。指劳动者在职业活动过程中长期反复接触，对绝大多数接触者的健康不引起有害作用的容许接触水平。\\
·化学有害因素的职业接触限制，包括：\\
-时间加权平均容许浓度（PC-TWA）：以时间为权数规定的8h 工作日、40h 工作周的平均容许接触浓度。\\
-短时间接触容许浓度（PC-STEL）：在遵守PC-TWA 前提下容许短时间（15min）接触的浓度。\\
-最高容许浓度（MAC）：工作地点、在一个工作日内、任何时间有毒化学物质均不应超过的浓度。\\
·物理因素的职业接触限制，包括：时间加权平均容许限值、最高容许限制

\ContentBand{对应往年题}

\begin{pastquestion}

\QuestionLabel{【往年题1｜10级职卫考题回忆.docx｜名解】}\\
全英，工效学，妇女职业卫生，热痉挛，粉尘分散度，短时间接触容许浓度，氧需

\end{pastquestion}

\begin{pastquestion}

\QuestionLabel{【往年题2｜14级职卫考题.pdf｜名解】}\\
Occupation exposure limit

\end{pastquestion}

\begin{pastquestion}

\QuestionLabel{【往年题3｜21级本科预防二班+胡润泽+职业卫生+主观题总结.docx｜名词解释】}\\
biological exposure limit, BEL

\end{pastquestion}

\begin{pastquestion}

\QuestionLabel{【往年题4｜04-14级职业卫生论述往年题汇总.docx｜14级附录】}\\
职业接触限值Occupation exposure
limit，OEL：\ldots\ldots 化学因素的职业接触限值可分为时间加权平均容许浓度、最高容许浓度和短时间接触容许浓度三类\\
时间加权平均容许浓度 PC-TWA\ldots\ldots{}\\
短时间接触容许浓度 PC-STEL\ldots\ldots{}\\
最高容许浓度（MAC）：工作地点，一个工作日内任何时间有毒化学物质均不应超过的浓度

\end{pastquestion}

\PointSeparator

\section{健康监护的内容与制度}\label{users__fpb__26_spring__pkusph_study__tmp__pdfs__prevention-medicine-past-exams__01-ux804cux4e1aux536bux751fux5b66ux5f80ux5e74ux9898ux8003ux70b9ux6574ux7406.md__ux5065ux5eb7ux76d1ux62a4ux7684ux5185ux5bb9ux4e0eux5236ux5ea6}

\FrequencyPill{2 次}

\InfoLabel{对应小节} 综合（第十四章管理实践 + 第十五章 健康监护）\\
\InfoLabel{匹配依据} 10 简答''劳动者职业健康监护制度内容''；11 简答''健康监护的内容和指标''。

\ContentBand{知识点原文摘取}

4、职业卫生工作常开展的三种监测工作\\
·环境监测、生物监测、健康监护→环境化学物→机体→排出；（转归全过程）\\
·环境监测及生物监测均属于暴露评价，是开展健康效应研究的基础；\\
·健康监护：识别早期健康效应，是采取干预措施及改进管理的依据；

5、健康监护(GBZ235 -2011《放射工作人员职业健康监护技术规范》)：建立健康档案、定期体检及监测\\
·个人剂量监测：包括外照射剂量、皮肤污染、体内污染监测\\
·工作场所监测：X、γ 射线、工作场所表面污染。空气中气载放射性核素浓度监测等。

4、卫生保健措施------接尘工人健康检查\\
·就业前检查\\
·定期检查：应根据接尘情况确定检查间隔；检查项目应包括职业史、自觉症状和拍摄后前位胸大片\\
·脱尘检查：根据接触粉尘性质确定检查间隔

\ContentBand{对应往年题}

\begin{pastquestion}

\QuestionLabel{【往年题1｜10级职卫考题回忆.docx｜简答】}\\
物理因素的特点、有机磷中毒的机制和临床表现、李涛所长之劳动者职业健康监护制度内容、手臂振动病的定义,临床表现和典型特征、常见职业伤害事故类型

\end{pastquestion}

\begin{pastquestion}

\QuestionLabel{【往年题2｜11级职卫考题.docx｜简答4】}\\
4健康监护的内容和指标

\end{pastquestion}

\PointSeparator

\section{职业卫生标准 /
不属于职业卫生法的是}\label{users__fpb__26_spring__pkusph_study__tmp__pdfs__prevention-medicine-past-exams__01-ux804cux4e1aux536bux751fux5b66ux5f80ux5e74ux9898ux8003ux70b9ux6574ux7406.md__ux804cux4e1aux536bux751fux6807ux51c6-ux4e0dux5c5eux4e8eux804cux4e1aux536bux751fux6cd5ux7684ux662f}

\FrequencyPill{1 次}

\InfoLabel{对应小节} 第十四章 / 二、职业卫生法律体系；并见 04-14 级附录''职业卫生标准''定义\\
\InfoLabel{匹配依据} 14 单选''不属于职业卫生法的是''；疑似对应。

\ContentBand{知识点原文摘取}

·职业卫生标准（occupational health
standards）：为实施职业病防治法律法规和有关政策，保护劳动者健康，预防、控制和消除职业病危害，防治职业病，由法律授权部门根据职业病防治法的规定制定的在全国范围内统一实施的技术要求。

\ContentBand{对应往年题}

\begin{pastquestion}

\QuestionLabel{【往年题1｜14级职卫考题.pdf｜单选（回忆）】}\\
不属于职业卫生法的是

\end{pastquestion}

\begin{pastquestion}

\QuestionLabel{【往年题2｜04-14级职业卫生论述往年题汇总.docx｜14级附录】}\\
职业卫生标准（occupational health standards）：为实施职业病防治法律法规和有关政策\ldots\ldots{}

\end{pastquestion}

\PointSeparator

\section{建设项目''三同时''}\label{users__fpb__26_spring__pkusph_study__tmp__pdfs__prevention-medicine-past-exams__01-ux804cux4e1aux536bux751fux5b66ux5f80ux5e74ux9898ux8003ux70b9ux6574ux7406.md__ux5efaux8bbeux9879ux76eeux4e09ux540cux65f6}

\FrequencyPill{1 次}

\InfoLabel{对应小节} 第十四章 / 四、职业卫生管理基本制度（建设项目职业卫生管理）\\
\InfoLabel{匹配依据} 11 名解''建设的'三同时'\,``。疑似对应（PDF
在第十四章仅列出条目''四、职业卫生管理基本制度''，未展开具体内容；该考点为常见职业病防治法基础概念，引用原文不补充）。

\ContentBand{知识点原文摘取}

四、职业卫生管理基本制度\\
五、职业接触控制\\
（PDF 此节仅列条目，未展开''三同时''细则）

\ContentBand{对应往年题}

\begin{pastquestion}

\QuestionLabel{【往年题1｜11级职卫考题.docx｜名解】}\\
妇女职业卫生、建设的''三同时''、AED、equal loudness curve、生物监测、occupational injury

\end{pastquestion}

\PointSeparator

\section{职业健康管理 / 健康教育与健康促进 /
半导体公司职业卫生工作要点}\label{users__fpb__26_spring__pkusph_study__tmp__pdfs__prevention-medicine-past-exams__01-ux804cux4e1aux536bux751fux5b66ux5f80ux5e74ux9898ux8003ux70b9ux6574ux7406.md__ux804cux4e1aux5065ux5eb7ux7ba1ux7406-ux5065ux5eb7ux6559ux80b2ux4e0eux5065ux5eb7ux4fc3ux8fdb-ux534aux5bfcux4f53ux516cux53f8ux804cux4e1aux536bux751fux5de5ux4f5cux8981ux70b9}

\FrequencyPill{3 次}

\InfoLabel{对应小节} 综合（第十四章 + 健康监护 + 各章预防原则）\\
\InfoLabel{匹配依据} 24 单选 3''关于职业健康管理说法错误的是（职业健康体检不能用一般体检代替）``；24 论述
1''如何更好地开展职业健康教育和健康促进工作''；24 论述 2''试述天科合达半导体公司的职业卫生工作要点''。

\ContentBand{知识点原文摘取}

8、工作场所职业卫生管理规定：卫生健康部门监督→用人单位主体责任→技术服务机构服务\\
9、可持续的职业健康战略：有害因素、暴露评价；一般健康体检、职业健康监护、健康促进；危害控制\\
10、风险评估：危害识别+暴露评估+剂量效应+风险表征→管理+控制

·健康促进：开展健康教育和健康促进活动，增强个体应对职业紧张的能力（第八章）

4、安全卫生管理（第八章------金属毒物的职业健康管理）\\
·遵守职业安全卫生法规\\
·实施安全生产监督管理制度：定期检修设备；杜绝违章操作\\
·加强职业卫生知识宣传教育\\
5、职业卫生服务\\
·定期监测作业场所空气中金属毒物浓度：应低于国家职业卫生标准规定的职业接触限值\\
·上岗前体检、定期体检\\
·有职业禁忌症者，严禁上岗

\ContentBand{对应往年题}

\begin{pastquestion}

\QuestionLabel{【往年题1｜24年完整试题回忆.docx｜一、单选3】}\\
3、关于职业健康管理说法错误的是\\
A、职业健康体检可以用一般健康体检代替\\
B-E、\ldots\ldots{}

\end{pastquestion}

\begin{pastquestion}

\QuestionLabel{【往年题2｜24年完整试题回忆.docx｜四、论述1】}\\
1、如何更好地开展职业健康教育和健康促进工作

\end{pastquestion}

\begin{pastquestion}

\QuestionLabel{【往年题3｜24年完整试题回忆.docx｜四、论述2】}\\
2、从职业有害因素和职业健康管理的角度，试述天科合达半导体公司的职业卫生工作要点

\end{pastquestion}

\PointSeparator

\chapter{职业卫生监测}\label{users__fpb__26_spring__pkusph_study__tmp__pdfs__prevention-medicine-past-exams__01-ux804cux4e1aux536bux751fux5b66ux5f80ux5e74ux9898ux8003ux70b9ux6574ux7406.md__ux7b2cux5341ux4e94ux7ae0-ux804cux4e1aux536bux751fux76d1ux6d4b}

\section{生物监测（biological
monitoring）}\label{users__fpb__26_spring__pkusph_study__tmp__pdfs__prevention-medicine-past-exams__01-ux804cux4e1aux536bux751fux5b66ux5f80ux5e74ux9898ux8003ux70b9ux6574ux7406.md__ux751fux7269ux76d1ux6d4bbiological-monitoring}

\FrequencyPill{4 次}

\InfoLabel{对应小节} 三、生物监测 / 1、定义；2、特点\\
\InfoLabel{匹配依据} 09、11、14 名解''生物监测''；24 名解 biological monitoring。

\ContentBand{知识点原文摘取}

三、生物监测(Biological Monitoring)\\
1、定义：定期、有计划地监测人体材料中化学物质及其代谢产物的含量或由它们所致的无害生物效应水平，与参比值比较，以评价人体接触化学物的程度及可能的健康效应。\\
2、特点\\
（1）优点\\
·反映不同途径（消化道、呼吸道、皮肤）和来源（食物、空气、水，职业与非职业的）的总的接触量，而环境监测只能反映环境中通过呼吸道进入机体的量。\\
·可以直接检测引起健康损害作用的内接触剂量或内负荷。\\
·综合了个体接触毒物的差异因素和毒物的典型动力学过程极其变异性。\\
·通过易感性指标的监测，早发现、确定易感人群。\\
·花费较少，可较早地检出对健康可能的损害，为及时采取预防措施提供依据。\\
（2）缺点\\
·不适用于所有毒物\\
·对象是人，存在监测对象是否合作的问题------知情同意\\
·检测方法有限\\
·毒物接触与内剂量的关系了解不多，监测结果不易解释

\ContentBand{对应往年题}

\begin{pastquestion}

\QuestionLabel{【往年题1｜09级职卫考题回忆.doc｜名解】}\\
名解：全是英文，工效学、职业性肿瘤、氧债、生物监测、暂时性听阈位移

\end{pastquestion}

\begin{pastquestion}

\QuestionLabel{【往年题2｜11级职卫考题.docx｜名解】}\\
妇女职业卫生、建设的''三同时''、AED、equal loudness curve、生物监测、occupational injury

\end{pastquestion}

\begin{pastquestion}

\QuestionLabel{【往年题3｜14级职卫考题.pdf｜名解】}\\
Biological monitoring

\end{pastquestion}

\begin{pastquestion}

\QuestionLabel{【往年题4｜24年完整试题回忆.docx｜二、名解5】}\\
5、biological monitoring

\end{pastquestion}

\PointSeparator

\chapter{职业安全与职业伤害}\label{users__fpb__26_spring__pkusph_study__tmp__pdfs__prevention-medicine-past-exams__01-ux804cux4e1aux536bux751fux5b66ux5f80ux5e74ux9898ux8003ux70b9ux6574ux7406.md__ux7b2cux5341ux516dux7ae0-ux804cux4e1aux5b89ux5168ux4e0eux804cux4e1aux4f24ux5bb3}

\section{职业伤害（occupational
injury）---定义、事故类型与预防策略}\label{users__fpb__26_spring__pkusph_study__tmp__pdfs__prevention-medicine-past-exams__01-ux804cux4e1aux536bux751fux5b66ux5f80ux5e74ux9898ux8003ux70b9ux6574ux7406.md__ux804cux4e1aux4f24ux5bb3occupational-injuryux5b9aux4e49ux4e8bux6545ux7c7bux578bux4e0eux9884ux9632ux7b56ux7565}

\FrequencyPill{4 次}

\InfoLabel{对应小节} 一、伤害概述 / 4、职业伤害；二、职业伤害事故类型与原因 /
5、常见职业伤害事故类型；八、职业安全事故预防策略\\
\InfoLabel{匹配依据} 11、14 名解 occupational injury；10 简答''常见职业伤害事故类型''；24 单选
15''工伤预防策略不包括''。

\ContentBand{知识点原文摘取}

4、职业伤害(occupational
injuries)/工作伤害/工伤：指在生产劳动过程中，由于外部因素直接作用而引起机体组织的突发性意外损伤，如因职业性事故（occupational
accidents）导致的伤亡及急性化学物中毒。由这种事故造成人身伤害或死亡称为职业性伤亡事故。\\
5、职业安全(occupational
safety)/劳动安全：是研究预防和控制职业伤害事故的一门专业，是指在生产过程中，为避免人身或设备事故，创建安全、健康的生产和操作环境而采取的各项措施及相应活动。

5、常见职业伤害事故类型及其主要原因\\
·物体打击：高空作业时工具等高处坠落；起重吊装、拆装时物件掉落；\ldots\ldots{}\\
·机械伤害：检修、检查机械时忽视安全操作规程，如带电作业，未挂''不准开闸''警示牌等；\ldots\ldots{}\\
·高处坠落：\ldots\ldots{}\\
·车辆伤害：\ldots\ldots{}\\
·电击伤害：\ldots\ldots{}\\
·操作事故所致伤害：\ldots\ldots{}

8、职业安全事故预防策略\\
（1）我国工伤事故的预防控制方针：安全第一、预防为主、综合治理\\
（2）安全生产管理体制：企业全面负责；行业管理；国家监察；群众监督；劳动者遵章守纪\\
（3）三级预防\\
·一级预防：应用全人群策略、高危人群策略和健康促进策略，强化主动保护和被动防护措施，防止和减少伤害的发生\\
·二级预防：旨在降低伤害的死亡率和致残率，第一时间的紧急救护\ldots\ldots{}\\
·三级预防：主要是使伤者恢复正常功能和社区康复

\ContentBand{对应往年题}

\begin{pastquestion}

\QuestionLabel{【往年题1｜10级职卫考题回忆.docx｜简答】}\\
物理因素的特点、有机磷中毒的机制和临床表现、李涛所长之劳动者职业健康监护制度内容、手臂振动病的定义,临床表现和典型特征、常见职业伤害事故类型

\end{pastquestion}

\begin{pastquestion}

\QuestionLabel{【往年题2｜11级职卫考题.docx｜名解】}\\
妇女职业卫生、建设的''三同时''、AED、equal loudness curve、occupational injury

\end{pastquestion}

\begin{pastquestion}

\QuestionLabel{【往年题3｜14级职卫考题.pdf｜名解】}\\
Occupational injury

\end{pastquestion}

\begin{pastquestion}

\QuestionLabel{【往年题4｜24年完整试题回忆.docx｜一、单选15】}\\
15、工伤预防策略不包括\\
A、危险作业频繁更换操作者\\
B-E、\ldots\ldots{}

\end{pastquestion}

\PointSeparator

\chapter*{章节备注}\label{users__fpb__26_spring__pkusph_study__tmp__pdfs__prevention-medicine-past-exams__01-ux804cux4e1aux536bux751fux5b66ux5f80ux5e74ux9898ux8003ux70b9ux6574ux7406.md__ux7ae0ux8282ux5907ux6ce8}
\addcontentsline{toc}{chapter}{章节备注}

\begin{itemize}
\tightlist
\item
  \textbf{实习（何丽华）}：本章未在提供的往年题中匹配到明确考点。
\item
  \textbf{第七章中其他单体（丙烯腈、TDI、含氟塑料具体诊断分级等）}：除''高分子化合物概念/单体毒性/氯乙烯''外，未在往年题中明确出现作为独立考点。
\item
  \textbf{第十章中拟除虫菊酯、有机磷与拟除虫菊酯混配的具体细节}：未在往年题中独立出现。
\item
  \textbf{第十一章中 CS₂
  对生殖的影响、女性五期保护具体规定}：未在往年题中独立出现（但''生殖损伤/生殖系统损害''已纳入）。
\item
  \textbf{第十三章中横断面研究、铬酸盐案例等}：未在往年题中独立出现。
\item
  \textbf{第十四章中职业卫生法律体系详条}：除''职业卫生标准''\,``三同时''``健康监护制度''``OEL''外，未独立出现具体法条考点。
\item
  \textbf{第十五章环境监测、皮肤污染监测}：本次题库中除''生物监测''外，未独立成考点。
\item
  \textbf{第十六章 Haddon 模型、五E综合干预}：未在往年题中独立出现作为名解或简答。
\end{itemize}

\SetSubjectAccent{0F8A78}{毒理学}

\part{毒理学}

\chapter*{说明}\label{users__fpb__26_spring__pkusph_study__tmp__pdfs__prevention-medicine-past-exams__02-ux6bd2ux7406ux5b66ux5f80ux5e74ux9898ux8003ux70b9ux6574ux7406.md__ux8bf4ux660e}
\addcontentsline{toc}{chapter}{说明}

\begin{itemize}
\tightlist
\item
  本文件根据 14 节毒理学课件 Markdown 文件（《1 绪论》《2 毒理学实验基础》《3 Biotransportation and
  Biotransformation\_WXIAO》《4 生物转化》《5 毒作用机制、急性毒性和局部毒性试验》《6
  亚慢性和慢性毒性试验》《7 遗传毒理学（一）》《8 遗传毒理学（二）》《9 化学致癌作用》《10 免疫毒理学》《11
  化学致畸作用、生殖毒性、发育毒性2021》《11 化学致畸作用和发育毒性2022》《12 神经和神经行为毒理学》《13
  安全性评价》《14 风险评估、管理毒理学》）和用户提供的往年题文件整理。
\item
  知识点正文均摘自课件 Markdown
  原文，未自行改写、扩写或总结；原文较长时，在不改变原意和原顺序的前提下摘取与考题最相关的部分。
\item
  往年题均摘自用户提供的往年题文件（07级毒理口试题回忆.doc、09级口试题目.docx、14级毒理考题汇总.xls、15级毒理学考试题目汇总.xls、17级毒理试题汇总.xlsx、18级毒理回忆by连昕瑶（搬运自经验分享）.docx、19级毒理期末口试回忆.xlsx、20级完整试题回忆.docx、21级毒理期末口试回忆.xlsx、毒理学-来自课本的题库补充byYYS.docx）。
\item
  排序依据为：先按课件章节顺序，再在同一章节内按考频从高到低排列；同频时按课件原始出现顺序排列。
\item
  一道考题若同时考查多个知识点，分别归入多个考点，并各计 1 次。
\item
  往年题口试回忆文件多为每位同学（按序号）的 2
  道口试题；同一考点在不同年份/序号反复出现的，记为多次考频，``对应往年题''按来源文件归并并保留题目原文与题型；个别年份附带的追问、备注一并保留以便复习。
\item
  如果某道题未在课件中完全匹配到对应知识点，则标记为''未完全匹配到''，并根据题意暂归入最合适章节。
\item
  已逐题核对所有往年题文件，确保最终文件中没有任何题目类型遗漏。
\end{itemize}

\PointSeparator

\chapter{绪论}\label{users__fpb__26_spring__pkusph_study__tmp__pdfs__prevention-medicine-past-exams__02-ux6bd2ux7406ux5b66ux5f80ux5e74ux9898ux8003ux70b9ux6574ux7406.md__ux7b2cux4e00ux7ae0-1-ux7eeaux8bba}

\section{剂量-效应关系、剂量-反应关系及 S
形曲线}\label{users__fpb__26_spring__pkusph_study__tmp__pdfs__prevention-medicine-past-exams__02-ux6bd2ux7406ux5b66ux5f80ux5e74ux9898ux8003ux70b9ux6574ux7406.md__ux5242ux91cf-ux6548ux5e94ux5173ux7cfbux5242ux91cf-ux53cdux5e94ux5173ux7cfbux53ca-s-ux5f62ux66f2ux7ebf}

\InfoLabel{考频} 12 次

\InfoLabel{对应小节} 三、Basic Concepts of Toxicology /《1
绪论》第25页''效应和反应''、第27页''剂量--效应（反应）关系的类型''

\InfoLabel{匹配依据} 多年口试反复直接考''dose-response relationship 的概念和毒理学意义、S
型曲线的意义''\,``剂量反应关系、剂量效应关系、S 型曲线的解释和意义''，语义完全对应本小节。

\ContentBand{知识点原文摘取}

效应（effect）指接触的外源化学物所引起的生物学改变。此种变化的程度用计量单位来表示。毒效应是化学物质对机体所致的不可逆或有害的生物学改变。反应（response）指接触的外源化学物的群体中出现某种效应的个体在群体中所占比例，一般以百分率或比值表示。

剂量--效应（反应）关系的类型：1. 直线型；2. 抛物线型（先陡峭后平缓的对数曲线）；3.
S形曲线，在外源化学物的剂量--反应关系中较为常见，呈两端平缓中间陡峭的''S''形，即在低剂量范围内随着剂量增加反应或效应强度增高较为缓慢，然后剂量较高时反应或效应强度也随之急速增加，但当剂量继续增加时反应或效应强度增高又趋向缓慢。可分为对称S形曲线和非对称S形曲线两种形式。

\begin{correctionbox}

【勘误补充】\\
剂量-效应关系（dose-effect
relationship）表示化学物质的剂量与个体或群体中发生的量效应强度之间的关系。其中，``效应（effect）''是指外源化学物引起的生物学改变，其变化程度可以用计量单位表示，因此剂量-效应关系强调的是效应强弱的变化。剂量-反应关系（dose-response
relationship）表示化学物质的剂量与某一群体中某种反应发生率之间的关系。其中，``反应（response）''是指接触外源化学物后，在一个群体中出现某种效应的个体所占的比例，通常以百分率或比值表示，因此剂量-反应关系强调的是发生反应的人群比例。二者的区别在于：剂量-效应关系关注效应强度的变化，而剂量-反应关系关注反应发生率的变化。

\end{correctionbox}

剂量-反应关系的毒理学意义包括：①说明毒作用与剂量之间存在规律性关系，一般剂量越大，毒作用越强，反应率越高；②是判断外源化学物与毒性效应之间因果关系的重要依据；③可用于确定
LD50、ED50、NOAEL、LOAEL、threshold dose
等重要毒性参数；④是安全性评价和风险评估的重要基础，可用于制定安全限值和接触标准；⑤可用于比较不同化学物毒性的强弱。

外源化学物的剂量-反应关系常表现为 S
型曲线，其特点为``两端平缓，中间陡峭''：低剂量区反应增加缓慢，中剂量区反应迅速增加，高剂量区又逐渐趋于平缓。S
型曲线的毒理学意义包括：①反映群体个体敏感性的差异；②有助于确定 LD50、ED50
和阈剂量等毒性参数；③可用于比较毒物毒性的强弱，曲线越靠左说明毒性越强，曲线越陡说明个体差异越小；④为安全性评价和风险评估提供依据。

剂量-反应关系之所以常表现为 S
型曲线，主要与阈值、个体差异和耐受性有关。低剂量时机体存在代偿、修复和解毒能力，只有少数敏感个体发生反应，因此曲线开始较平缓；随着剂量增加，更多不同敏感性的个体陆续发生反应，因此中段曲线变得陡峭；而在高剂量时，大多数敏感个体已经发生反应，剩余个体多为耐受性较强或不敏感个体，因此继续增加剂量时新增反应个体有限，曲线再次趋于平缓。

\InfoLabel{知识来源} 《1 绪论》第25页；《1 绪论》第27页

\ContentBand{对应往年题}

\begin{pastquestion}

\QuestionLabel{【往年题1｜07级毒理口试题回忆.doc｜口试/简答】}\\
does-response relationship
的概念和毒理学意义。S型曲线的意义（1号题、43号题）；剂量反应得概念和意义，S曲线的意义（22题）；剂量---反应关系定义及毒理学意义、s型剂量反应关系的意义，追问：为什么常见的剂量反应为s型曲线（阈值，个体差异，耐受性）（46号题）。PS：剂量反应关系和剂量效应关系要区分一下，郝老师追问S型曲线为什么变平，应该主要是个体差异的原因。

\end{pastquestion}

\begin{pastquestion}

\QuestionLabel{【往年题2｜09级口试题目.docx｜口试】}\\
dose-response relationship的概念和毒理学意义；S型dose-response relationship的意义（22）

\end{pastquestion}

\begin{pastquestion}

\QuestionLabel{【往年题3｜21级毒理期末口试回忆.xlsx｜口试】}\\
剂量-反应关系 剂量-效应关系
S型曲线（序号8）；剂量效应，计量反应关系概念及毒理学意义，s形剂量效应关系的意义（序号54）

\end{pastquestion}

\begin{pastquestion}

\QuestionLabel{【往年题4｜17级毒理试题汇总.xlsx｜口试】}\\
剂量反应关系、剂量效应关系、S型曲线的解释和意义（题号9）；dose-response和dose-effect概念和毒理学意义，s形剂量反应曲线的意义（题号88）

\end{pastquestion}

\PointSeparator

\section{毒理学的主要研究方法及其优缺点}\label{users__fpb__26_spring__pkusph_study__tmp__pdfs__prevention-medicine-past-exams__02-ux6bd2ux7406ux5b66ux5f80ux5e74ux9898ux8003ux70b9ux6574ux7406.md__ux6bd2ux7406ux5b66ux7684ux4e3bux8981ux7814ux7a76ux65b9ux6cd5ux53caux5176ux4f18ux7f3aux70b9}

\InfoLabel{考频} 12 次

\InfoLabel{对应小节} 二、Methods in Toxicology Study /《1 绪论》第12--14页、第45页

\InfoLabel{匹配依据}
``毒理学主要研究方法的优缺点''``四种毒理学方法的优缺点''``各种毒理学研究方法的优缺点''反复出现，对应本小节四种研究方法的优缺点比较表。

\ContentBand{知识点原文摘取}

1 In vivo study in experimental
animals：Advantages------化学物在实验动物产生的作用大多数情况下可以外推于人；易于控制暴露条件；能评价宿主特征的作用（如性别、年龄、遗传特征等）。Disadvantages------实验动物和人对外源化学物的反应敏感性不同，有时甚至存在着质的差别；暴露的浓度和时间的模式不同于人群的暴露；动物保护和动物福利的压力。

2 In vitro
study：Advantages------影响因素少，易于控制；可进行深入的研究（如机制、代谢）；人力物力花费较少。Disadvantages------不能全面反映毒作用；难以观察慢性毒作用。

3 Human clinical
study：Advantages------控制暴露条件；直接获得人体毒性资料。Disadvantages------伦理限制；限于暂时、微小、可逆的效应；一般不适于研究敏感的人群。

4 Epidemiologic study in
population：Advantages------真实的暴露条件；涵盖全部的人敏感性。Disadvantages------有混杂暴露；多为回顾性。

\InfoLabel{知识来源} 《1 绪论》第12页；《1 绪论》第13页；《1 绪论》第14页；《1 绪论》第45页

\ContentBand{对应往年题}

\begin{pastquestion}

\QuestionLabel{【往年题1｜07级毒理口试题回忆.doc｜口试】}\\
毒理学4种研究方法优缺点比较（2号题）；毒理学试验方法及优缺点（26号题）；毒理学主要研究方法的优缺点（39号题）。姚碧云对将流行病学方法放在第一个说很不满意，问毒理学最主要的方法是啥，回答了体内跟体外后就放走了。

\end{pastquestion}

\begin{pastquestion}

\QuestionLabel{【往年题2｜09级口试题目.docx｜口试】}\\
毒理的研究方法（追问流行病研究方法为什么没有人群保护）（2）；毒理学研究的方法和优缺点（26）；各种毒理学研究方法的优缺点（39）

\end{pastquestion}

\begin{pastquestion}

\QuestionLabel{【往年题3｜21级毒理期末口试回忆.xlsx｜口试】}\\
毒理学的主要研究方法及其优缺点（序号50）

\end{pastquestion}

\begin{pastquestion}

\QuestionLabel{【往年题4｜19级毒理期末口试回忆.xlsx｜口试】}\\
毒理学研究方法和优缺点（序号82）

\end{pastquestion}

\begin{pastquestion}

\QuestionLabel{【往年题5｜17级毒理试题汇总.xlsx｜口试】}\\
毒理学实验的方法和优缺点（题号10）；毒理学的四个主要方法和优缺点（题号58）

\end{pastquestion}

\begin{pastquestion}

\QuestionLabel{【往年题6｜15级毒理学考试题目汇总.xls｜口试】}\\
毒理学研究方法及优缺点（题号68）；毒理学研究的方法优缺点（题号30）；毒理学各种实验研究的优缺点（题号68）

\end{pastquestion}

\begin{pastquestion}

\QuestionLabel{【往年题7｜14级毒理考题汇总.xls｜口试】}\\
毒理学各种实验研究的优缺点（题号1）；毒理学方法和优缺点（题号19）；毒理学主要研究方法和优缺点（题号68）

\end{pastquestion}

\PointSeparator

\section{毒理学研究范畴（描述、机制、管理毒理学）及各领域间的关系}\label{users__fpb__26_spring__pkusph_study__tmp__pdfs__prevention-medicine-past-exams__02-ux6bd2ux7406ux5b66ux5f80ux5e74ux9898ux8003ux70b9ux6574ux7406.md__ux6bd2ux7406ux5b66ux7814ux7a76ux8303ux7574ux63cfux8ff0ux673aux5236ux7ba1ux7406ux6bd2ux7406ux5b66ux53caux5404ux9886ux57dfux95f4ux7684ux5173ux7cfb}

\InfoLabel{考频} 11 次

\InfoLabel{对应小节} 一、Categories of Toxicology /《1 绪论》第8--10页、第45页

\InfoLabel{匹配依据}
``毒理学研究范畴''``毒理学的领域及相互联系''``毒理学研究范畴和关系''反复出现，对应描述、机制、管理三大范畴及基础毒理学与各分支关系。

\ContentBand{知识点原文摘取}

Descriptive
描述毒理学：通过各种毒性试验研究环境因素的毒作用特征。以期为安全性评价和管理法规与措施的制订提供基础资料。Mechanistic
机制毒理学：研究环境因素对生物体毒作用的细胞、分子以及生化机制。Regulatory
管理毒理学：依据描述和机制毒理学对毒作用规律的研究成果，协助政府部门制订相关法规条例和管理措施并付诸实施，以确保对化学品、药品、食品等进行有效的管理，保障接触人群的健康和环境安全。

（基础毒理学与各分支以及相关学科的关系：环境和职业毒理学、食品毒理学、药物毒理学；基础毒理涉及遗传学、系统毒理、细胞与分子毒理、生化、分子生物、生理、病理、免疫。）

\InfoLabel{知识来源} 《1 绪论》第8页；《1 绪论》第9页；《1 绪论》第10页；《1 绪论》第45页

\ContentBand{对应往年题}

\begin{pastquestion}

\QuestionLabel{【往年题1｜21级毒理期末口试回忆.xlsx｜口试】}\\
毒理学研究范畴和关系（序号30）；毒理学研究领域及其联系（序号49）；毒理学的研究领域以及各个领域之间的关系（序号80）

\end{pastquestion}

\begin{pastquestion}

\QuestionLabel{【往年题2｜19级毒理期末口试回忆.xlsx｜口试】}\\
毒理学的研究领域，及相互之间的关系（序号49）；毒理学范畴及其联系（序号80）；毒理学的领域及相互联系（序号30）

\end{pastquestion}

\begin{pastquestion}

\QuestionLabel{【往年题3｜17级毒理试题汇总.xlsx｜口试】}\\
毒理学的几个范畴（题号57）；简述毒理学研究范畴（题号95）；毒理学研究的内容（题号9）

\end{pastquestion}

\begin{pastquestion}

\QuestionLabel{【往年题4｜15级毒理学考试题目汇总.xls｜口试】}\\
研究范畴（题号31）

\end{pastquestion}

\begin{pastquestion}

\QuestionLabel{【往年题5｜14级毒理考题汇总.xls｜口试】}\\
简述毒理学研究范畴（题号31）

\end{pastquestion}

\PointSeparator

\section{生物学标志（biomarker）的概念、分类及毒理学意义}\label{users__fpb__26_spring__pkusph_study__tmp__pdfs__prevention-medicine-past-exams__02-ux6bd2ux7406ux5b66ux5f80ux5e74ux9898ux8003ux70b9ux6574ux7406.md__ux751fux7269ux5b66ux6807ux5fd7biomarkerux7684ux6982ux5ff5ux5206ux7c7bux53caux6bd2ux7406ux5b66ux610fux4e49}

\InfoLabel{考频} 10 次

\InfoLabel{对应小节} 三、Basic Concepts of Toxicology /《1 绪论》第29--30页

\InfoLabel{匹配依据} ``biomarker
的概念、分类、毒理学意义''反复出现，完全对应生物学标志定义及接触/效应/易感性三类标志。

\ContentBand{知识点原文摘取}

生物学标志（biomarker）：指可反映外源化学物通过生物学屏障进入组织或体液（其原型及其代谢产物）以及引起的生物学后果（包括可测量的细胞、生化组分、结构及功能的改变和行为改变）的指标。

接触生物学标志（biomarker of
exposure）：在组织、体液或排泄物中测定的外源化学物或其代谢物，以及生物体内测定的外源化学物或其代谢物与靶分子或靶细胞之间相互作用的产物。可提供有关化学物质暴露的信息，可用于预测个体所接受的化学物质剂量，并能与导致的损害相关联。效应生物学标志（biomarker
of effect）：在生物体内可测量的生化、生理、行为、病理组织学或其他方面的变化。易感性生物学标志（biomarker of
susceptibility）：可反映机体对化学物质毒作用敏感程度的生理或生化状态改变的指标，对于易感人群的鉴定、区分和保护具有重要的意义。

\InfoLabel{知识来源} 《1 绪论》第29页；《1 绪论》第30页

\ContentBand{对应往年题}

\begin{pastquestion}

\QuestionLabel{【往年题1｜07级毒理口试题回忆.doc｜口试】}\\
biomarker的概念、分类、毒理学意义（30号、31号、52号题）

\end{pastquestion}

\begin{pastquestion}

\QuestionLabel{【往年题2｜09级口试题目.docx｜口试】}\\
biomarker的概念，分类及意义（12、31、55、52）

\end{pastquestion}

\begin{pastquestion}

\QuestionLabel{【往年题3｜19级毒理期末口试回忆.xlsx｜口试】}\\
biomarker的定义分类及意义（序号1）

\end{pastquestion}

\begin{pastquestion}

\QuestionLabel{【往年题4｜17级毒理试题汇总.xlsx｜口试】}\\
biomarker的概念、分类（具体内容都要说）、意义（题号97）；生物标志物定义分类毒理学意义（题号1）

\end{pastquestion}

\begin{pastquestion}

\QuestionLabel{【往年题5｜15级毒理学考试题目汇总.xls｜口试】}\\
biomarkers的定义、分类、意义（题号69）；biomarker概念 分类 意义（题号3）

\end{pastquestion}

\begin{pastquestion}

\QuestionLabel{【往年题6｜14级毒理考题汇总.xls｜口试】}\\
bromaker定义，分类，意义（题号80）；生物标志（题号69）

\end{pastquestion}

\PointSeparator

\section{毒作用效应谱、有阈毒性和无阈毒性作用}\label{users__fpb__26_spring__pkusph_study__tmp__pdfs__prevention-medicine-past-exams__02-ux6bd2ux7406ux5b66ux5f80ux5e74ux9898ux8003ux70b9ux6574ux7406.md__ux6bd2ux4f5cux7528ux6548ux5e94ux8c31ux6709ux9608ux6bd2ux6027ux548cux65e0ux9608ux6bd2ux6027ux4f5cux7528}

\InfoLabel{考频} 7 次

\InfoLabel{对应小节} 三、Basic Concepts of Toxicology /《1
绪论》第25页''毒作用效应谱''、第26页''有阈和无阈毒性作用''

\InfoLabel{匹配依据} ``毒效应谱、有阈毒性和无阈毒性的概念''``毒作用效应谱和有阈和无阈毒性''直接对应本小节。

\ContentBand{知识点原文摘取}

毒作用效应谱（spectrum of toxic
effects）：机体接触外源化学物后，依外源化学物的性质和剂量，可引起从微小的生理生化指标异常改变到明显的临床中毒表现，直至死亡的多种毒作用表现，称为毒效应谱。

有阈和无阈毒性作用：一般认为，外源化学物的一般毒性（器官毒性）和致畸作用是有阈值的，即达到一定的剂量水平才对机体产生毒性作用；而遗传毒性致癌作用和致突变作用则无阈值，即只要接触就可能产生有害作用。

\InfoLabel{知识来源} 《1 绪论》第25页；《1 绪论》第26页

\ContentBand{对应往年题}

\begin{pastquestion}

\QuestionLabel{【往年题1｜09级口试题目.docx｜口试】}\\
毒作用分类和效应谱（19）；毒效应谱和毒性作用的分类（9）；有阈和无阈化学物的剂量反应关系（21）

\end{pastquestion}

\begin{pastquestion}

\QuestionLabel{【往年题2｜21级毒理期末口试回忆.xlsx｜口试】}\\
毒效应谱、有阈毒性和无阈毒性的概念（序号23）

\end{pastquestion}

\begin{pastquestion}

\QuestionLabel{【往年题3｜17级毒理试题汇总.xlsx｜口试】}\\
毒效应谱和有阈和无阈毒性（题号98）；毒效应谱和有阈毒性和无阈毒性（题号73）

\end{pastquestion}

\begin{pastquestion}

\QuestionLabel{【往年题4｜20级完整试题回忆.docx｜单项选择】}\\
一般认为，化学物的哪种毒性作用是无阈值的（第10题）

\end{pastquestion}

\PointSeparator

\section{安全限值与安全系数（含
ADI、RfD、MAC；动物外推到人的方法）}\label{users__fpb__26_spring__pkusph_study__tmp__pdfs__prevention-medicine-past-exams__02-ux6bd2ux7406ux5b66ux5f80ux5e74ux9898ux8003ux70b9ux6574ux7406.md__ux5b89ux5168ux9650ux503cux4e0eux5b89ux5168ux7cfbux6570ux542b-adirfdmacux52a8ux7269ux5916ux63a8ux5230ux4ebaux7684ux65b9ux6cd5}

\InfoLabel{考频} 9 次

\InfoLabel{对应小节} 三、Basic Concepts of Toxicology /《1
绪论》第19页''安全限值''\,``安全系数''、第20页''ADI/RfD/MAC''

\InfoLabel{匹配依据} ``安全限值和安全系数概念、意义''``safety limit value和safety
factor''反复出现，对应本小节定义及安全系数 100 的来源。

\ContentBand{知识点原文摘取}

安全限值：对各种有害因素规定的接触限量要求，在低于此接触量时，根据现有的知识，不会观察到任何直接和/或间接的有害作用。通常通过最大未观察到有害作用剂量（NOAEL）除以安全系数计算得出。

安全系数（safety
factor）：是考虑到动物试验结果外推到人群时的不确定性及人群毒性资料本身所包含的不确定因素，确定的安全性界限范围。动物试验外推到人通常有三种基本的方法：利用不确定系数（安全系数）；利用药物动力学外推；利用数学模型。安全系数一般采用100，为物种间差异（10）和个体间差异（10）的积。

每日容许摄入量（ADI）：以体重表达的每日容许摄入量，以此量终生摄入无可测量的健康危险性（标准人为60kg）。参考剂量（RfD）/参考浓度（RfC）：人群（包括敏感亚群）终身暴露于该水平时，预期在一生中发生非致癌（或非致突变）性有害效应的危险度很低。最高容许浓度（MAC）：指某一外源化学物可以在环境中存在而不致对人体造成任何损害作用的浓度。

\InfoLabel{知识来源} 《1 绪论》第19页；《1 绪论》第20页

\ContentBand{对应往年题}

\begin{pastquestion}

\QuestionLabel{【往年题1｜21级毒理期末口试回忆.xlsx｜口试】}\\
安全限值和安全系数概念、意义（序号37）；安全限值和安全系数的概念和意义（序号81）

\end{pastquestion}

\begin{pastquestion}

\QuestionLabel{【往年题2｜19级毒理期末口试回忆.xlsx｜口试】}\\
安全系数，安全限值（序号97，备注：还要记得ADI（营养）和最高容许浓度（环境））；安全限值和安全系数（英文）概念和毒理学意义（序号37）；safety
limit value和safety factor以及它们的意义（序号82）

\end{pastquestion}

\begin{pastquestion}

\QuestionLabel{【往年题3｜17级毒理试题汇总.xlsx｜口试】}\\
安全限值和安全系数的概念和意义（英文）（题号37，魏老师追问安全系数是固定的吗，没有NOAEL的情况下如何制定安全限值）

\end{pastquestion}

\begin{pastquestion}

\QuestionLabel{【往年题4｜09级口试题目.docx｜口试】}\\
1RfD什么情况下用NOAEL什么情况下用LOAEL（50号备注追问）

\end{pastquestion}

\PointSeparator

\section{常用毒性参数：致死剂量（LD100/MLD/MTD/LD50）、LOAEL、NOAEL、阈剂量的概念及意义}\label{users__fpb__26_spring__pkusph_study__tmp__pdfs__prevention-medicine-past-exams__02-ux6bd2ux7406ux5b66ux5f80ux5e74ux9898ux8003ux70b9ux6574ux7406.md__ux5e38ux7528ux6bd2ux6027ux53c2ux6570ux81f4ux6b7bux5242ux91cfld100mldmtdld50loaelnoaelux9608ux5242ux91cfux7684ux6982ux5ff5ux53caux610fux4e49}

\InfoLabel{考频} 10 次

\InfoLabel{对应小节} 三、Basic Concepts of Toxicology /《1
绪论》第22页''致死剂量''、第23页''LOAEL/NOAEL/阈剂量''

\InfoLabel{匹配依据} ``阈值（threshold）的概念和毒理学意义''``NOAEL和LOAEL的概念和意义''``threshold
dose、NOAEL、LOAEL定义及毒理学意义''反复出现，对应本小节各毒性参数定义；致死参数定义亦出自本页。

\ContentBand{知识点原文摘取}

绝对致死剂量（LD100）：能引起一组观察生物体全部死亡的最低剂量。最小致死剂量（MLD）：在一个观察群体中，仅引起个别生物体发生死亡的最低剂量。最大耐受剂量（MTD，LD0）：在一个观察群体中，不引起生物体死亡的最高剂量。半数致死剂量（LD50）：在一个观察群体中能引起50\%生物体死亡的剂量。它是依据实验数据，经过计算得到的一个统计值。

观察到有害作用的最低剂量（LOAEL）：通过实验和观察，引起机体某种有害改变的最低剂量。最大未观察到有害作用剂量（NOAEL）：通过实验和观察，以现有的技术手段和检测指标，未观察到与受试物有关的毒性作用的最大剂量。阈剂量（threshold
dose）：化学物质引起受试对象中的少数个体出现某种最轻微的异常改变的最低剂量。

\InfoLabel{知识来源} 《1 绪论》第22页；《1 绪论》第23页

\ContentBand{对应往年题}

\begin{pastquestion}

\QuestionLabel{【往年题1｜07级毒理口试题回忆.doc｜口试】}\\
阈值（英文）的概念和毒理学意义（3号、25号题）；LOAEL
NOAEL的概念和意义（4号题）；阈值的定义和毒理学意义（25号题）

\end{pastquestion}

\begin{pastquestion}

\QuestionLabel{【往年题2｜09级口试题目.docx｜口试】}\\
threshold定义和毒理学意义（3）；noael loael的概念
毒理学意义（13）；阈值（英文）概念及意义（24）；阈值的概念和毒理学意义（25）

\end{pastquestion}

\begin{pastquestion}

\QuestionLabel{【往年题3｜21级毒理期末口试回忆.xlsx｜口试】}\\
LD50的毒理学意义与局限性（序号21）；threshold
dose、NOAEL、LOAEL定义及毒理学意义（序号55，备注：阈剂量和LOAEL的区别可以自己总结一下）

\end{pastquestion}

\begin{pastquestion}

\QuestionLabel{【往年题4｜17级毒理试题汇总.xlsx｜口试】}\\
阈值，loael，noael概念及其毒理学意义（题号5、87）；ld50概念意义局限（题号23）；threshold概念及其毒理学意义（题号3）；阈值的概念和毒理学意义（题号89）

\end{pastquestion}

\begin{pastquestion}

\QuestionLabel{【往年题5｜18级毒理回忆by连昕瑶（搬运自经验分享）.docx｜简答答案】}\\
急性毒性实验参数（含上限参数 LD100/LD50/LD01/LD0、下限参数 NOAEL/LOAEL，及 阈剂量 threshold dose 定义）

\end{pastquestion}

\PointSeparator

\section{毒理学发展趋势与现代毒理学新领域（组学、系统毒理学、毒性测试策略转变、计算/预测毒理学、转化毒理学、3Rs）}\label{users__fpb__26_spring__pkusph_study__tmp__pdfs__prevention-medicine-past-exams__02-ux6bd2ux7406ux5b66ux5f80ux5e74ux9898ux8003ux70b9ux6574ux7406.md__ux6bd2ux7406ux5b66ux53d1ux5c55ux8d8bux52bfux4e0eux73b0ux4ee3ux6bd2ux7406ux5b66ux65b0ux9886ux57dfux7ec4ux5b66ux7cfbux7edfux6bd2ux7406ux5b66ux6bd2ux6027ux6d4bux8bd5ux7b56ux7565ux8f6cux53d8ux8ba1ux7b97ux9884ux6d4bux6bd2ux7406ux5b66ux8f6cux5316ux6bd2ux7406ux5b663rs}

\InfoLabel{考频} 8 次

\InfoLabel{对应小节} 四、Progress in Toxicology /《1 绪论》第31--43页、第46页

\InfoLabel{匹配依据}
``毒理学发展的趋势和新领域''``现代毒理学的新领域及发展趋势''``替代毒理学的概念和主要方法''等，对应本小节六个发展趋势。

\ContentBand{知识点原文摘取}

毒理学发展趋势：毒理学研究规范化和标准化（GLP成为法规毒理学研究必备条件）；组学技术的应用和系统毒理学的发展；毒性测试策略的转变（Toxicity
Testing in the 21st
Century，将以整体动物试验为基础的传统毒性测试转变到主要以体外毒性试验为主）；3Rs在毒理学研究中的应用；计算/预测毒理学；转化毒理学。

系统毒理学：通过生物信息学和计算毒理学技术整合机体暴露后在不同剂量、不同时间的基因表达谱、蛋白质表达谱和毒物代谢谱的改变等分子、细胞、组织等不同研究层次的高通量信息，结合传统毒理学的研究资料，系统研究外源性化学物和环境应激等与机体的交互作用。转化毒理学：研究如何将毒理学基础研究成果迅速有效地转化为可在公共卫生与预防医学实际工作中应用的理论、技术和方法。

\InfoLabel{知识来源} 《1 绪论》第31页；《1 绪论》第33页；《1 绪论》第36页；《1 绪论》第37页；《1
绪论》第43页；《1 绪论》第46页

\ContentBand{对应往年题}

\begin{pastquestion}

\QuestionLabel{【往年题1｜21级毒理期末口试回忆.xlsx｜口试】}\\
毒理学发展的趋势和新领域（序号8）；替代毒理学的概念和主要方法（序号43）；现代毒理学的新领域及发展趋势（序号100）

\end{pastquestion}

\begin{pastquestion}

\QuestionLabel{【往年题2｜17级毒理试题汇总.xlsx｜口试】}\\
现代毒理学发展的新领域与发展趋势（题号38）；简述毒理学的新领域和未来的发展趋势（题号86）

\end{pastquestion}

\begin{pastquestion}

\QuestionLabel{【往年题3｜19级毒理期末口试回忆.xlsx｜口试】}\\
毒理学研究外源性物质与人体的交互作用，研究内容具体是什么（序号82，相关）

\end{pastquestion}

\PointSeparator

\section{毒理学替代法（Toxicology alternatives）与 3Rs
原则}\label{users__fpb__26_spring__pkusph_study__tmp__pdfs__prevention-medicine-past-exams__02-ux6bd2ux7406ux5b66ux5f80ux5e74ux9898ux8003ux70b9ux6574ux7406.md__ux6bd2ux7406ux5b66ux66ffux4ee3ux6cd5toxicology-alternativesux4e0e-3rs-ux539fux5219}

\InfoLabel{考频} 9 次

\InfoLabel{对应小节} 四、Progress in Toxicology /《1 绪论》第39--42页

\InfoLabel{匹配依据}
``毒理学替代实验的意义及方法''``替代法的概念和具体方法''``毒理学中替代法的概念及其应用（3R原则里的替代）''反复出现，对应本小节。

\ContentBand{知识点原文摘取}

3Rs Rules：在生物医学实验中减少（reduction）、替代（replacement）和优化（refinement）使用实验动物。

Toxicology
alternatives：在毒理学研究中用组织学、胚胎学、细胞学或计算机等方法取代整体动物实验，以低等动物取代高等级动物等。替代法的基本方法：体外技术及人类模型的使用（培养的细胞、组织、器官）；低等物种的利用（如植物、细菌、真菌、昆虫、软体动物、早期发育阶段脊椎动物）；物理化学方法与计算机的使用（定量结构-活性关系
QSAR------毒性预测评价）。ECVAM验证并被OECD列入测试指南的急性毒性替代方法：固定剂量法、急性毒性分级法、上/下移动法。

\InfoLabel{知识来源} 《1 绪论》第39页；《1 绪论》第40页；《1 绪论》第41页；《1 绪论》第42页

\ContentBand{对应往年题}

\begin{pastquestion}

\QuestionLabel{【往年题1｜19级毒理期末口试回忆.xlsx｜口试】}\\
毒理替代试验的概念和方法（序号82）；替代法的概念和具体方法（序号84，备注：操作不用紧张\ldots\ldots 局部毒性试验的常用替代方法）；毒理学中替代法的概念及其应用（就是3R原则里面的那个替代）（序号28）

\end{pastquestion}

\begin{pastquestion}

\QuestionLabel{【往年题2｜15级毒理学考试题目汇总.xls｜口试】}\\
毒理学替代实验的意义、方法（题号33）；替代法意义及主要方法（无题号）；毒理学替代实验的意义及方法（题号4）

\end{pastquestion}

\begin{pastquestion}

\QuestionLabel{【往年题3｜14级毒理考题汇总.xls｜口试】}\\
替代法的意义和主要方法（题号33，备注：替代法记得不要只答急毒和局毒的，其他也要提）；毒理学替代实验的意义及方法（题号4）

\end{pastquestion}

\PointSeparator

\section{毒理学定义、主要研究内容（外源化学物与生物机体的交互作用：毒动学与毒效学）}\label{users__fpb__26_spring__pkusph_study__tmp__pdfs__prevention-medicine-past-exams__02-ux6bd2ux7406ux5b66ux5f80ux5e74ux9898ux8003ux70b9ux6574ux7406.md__ux6bd2ux7406ux5b66ux5b9aux4e49ux4e3bux8981ux7814ux7a76ux5185ux5bb9ux5916ux6e90ux5316ux5b66ux7269ux4e0eux751fux7269ux673aux4f53ux7684ux4ea4ux4e92ux4f5cux7528ux6bd2ux52a8ux5b66ux4e0eux6bd2ux6548ux5b66}

\InfoLabel{考频} 6 次

\InfoLabel{对应小节} 一、Definition of Toxicology /《1 绪论》第4--5页（并见《9 化学致癌作用》第4页）

\InfoLabel{匹配依据}
``毒理学研究内容包括化学物质和生物机体的交互作用，其研究内容包括（指毒效学和毒动学）''``毒理学的主要研究内容（书上的内容）''，对应毒理学定义与研究内容。

\ContentBand{知识点原文摘取}

毒理学：从生物学角度研究外源因素和内源因素对生物体和生态系统的损害作用与机制，并进行安全性评价、危险度评定/风险评估、危险性管理与交流的科学。Toxicology
is a scientific discipline involving in the study of adverse effects of environmental agents on living
organisms and ecosystem and the mechanisms of its action. It is vitally important for the food safety, drug
safety and environment safety.

【翻译】毒理学是一门研究环境因素对生物体和生态系统的不良影响及其作用机制的科学学科。它对于食品安全、药品安全和环境安全都非常重要。

（《9
化学致癌作用》第4页：化学物质与生物机体的有害交互作用------外源化学物经暴露、吸收、分布、代谢、排泄形成终毒物，作用于靶器官；毒动学、毒效学；化学结构、毒性、交互作用、毒效应。）

\InfoLabel{知识来源} 《1 绪论》第4页；《1 绪论》第5页；《9 化学致癌作用》第4页

\ContentBand{对应往年题}

\begin{pastquestion}

\QuestionLabel{【往年题1｜17级毒理试题汇总.xlsx｜口试】}\\
毒理学研究内容包括化学物质和生物机体的交互作用，其研究内容包括（指的是毒效学和毒动学）（题号61）；毒理学主要研究化学物质与生物机体的交互作用，其主要研究内容包括？（题号36，备注：第一道题是书里原话，不要答分类）

\end{pastquestion}

\begin{pastquestion}

\QuestionLabel{【往年题2｜19级毒理期末口试回忆.xlsx｜口试】}\\
毒理学研究外源性物质与人体的交互作用，研究内容具体是什么？（序号82）；毒理学的主要研究内容（书上的内容）（序号84）

\end{pastquestion}

\begin{pastquestion}

\QuestionLabel{【往年题3｜21级毒理期末口试回忆.xlsx｜口试】}\\
简述机体对外源化学物的作用及其毒理学意义（序号26，相关，亦见第三章）

\end{pastquestion}

\PointSeparator

\section{治疗指数（Therapeutic
Index）}\label{users__fpb__26_spring__pkusph_study__tmp__pdfs__prevention-medicine-past-exams__02-ux6bd2ux7406ux5b66ux5f80ux5e74ux9898ux8003ux70b9ux6574ux7406.md__ux6cbbux7597ux6307ux6570therapeutic-index}

\InfoLabel{考频} 1 次

\InfoLabel{对应小节} 三、Basic Concepts of Toxicology /《1 绪论》第21页''治疗指数''

\InfoLabel{匹配依据} 07级口试关于 LD50
意义中含''计算药物的治疗指数（LD50/ED50）``，单列以备复习；本知识点在课件中明确。

\ContentBand{知识点原文摘取}

治疗指数（The Therapeutic
Index，TI）：T.I.=LD50/ED50。对于药物，常用治疗指数来计算其安全性，TI越大安全性越高。但治疗指数不够完善，没考虑药物最大有效量时的毒性和剂量---反应曲线斜率。故用安全界限来评价药物的安全性比治疗指数更佳。

\InfoLabel{知识来源} 《1 绪论》第21页

\ContentBand{对应往年题}

\begin{pastquestion}

\QuestionLabel{【往年题1｜07级毒理口试题回忆.doc｜口试】}\\
LD50的概念和优缺点（含''计算药物的治疗指数 LD50/ED50，了解药效学剂量和毒性剂量的距离''）（1号题、3号题等）

\end{pastquestion}

\PointSeparator

\section{靶剂量、靶器官的概念及影响靶器官形成的因素（久远题库补充）}\label{users__fpb__26_spring__pkusph_study__tmp__pdfs__prevention-medicine-past-exams__02-ux6bd2ux7406ux5b66ux5f80ux5e74ux9898ux8003ux70b9ux6574ux7406.md__ux9776ux5242ux91cfux9776ux5668ux5b98ux7684ux6982ux5ff5ux53caux5f71ux54cdux9776ux5668ux5b98ux5f62ux6210ux7684ux56e0ux7d20ux4e45ux8fdcux9898ux5e93ux8865ux5145}

\InfoLabel{考频} 2 次

\InfoLabel{对应小节} 三、Basic Concepts of Toxicology /《1 绪论》第24页''剂量''；化学物毒性分类 /《9
化学致癌作用》第4--5页

\InfoLabel{匹配依据}
相对久远题库反复列出''靶器官的概念及决定某器官成为靶器官的影响因素''，现有整理中仅在毒性评价目的和致癌试验目的中提及''确定靶器官''，未单列概念题。

\ContentBand{知识点原文摘取}

靶剂量（target
dose）：指到达靶器官（如组织、细胞）的外源化学物和/或其代谢产物的剂量。化学物经暴露、吸收、分布、代谢、排泄形成终毒物，作用于靶器官，在分子、细胞、器官水平产生毒效应。化学物毒性包括基础毒性、系统毒性、局部毒性和特殊毒性，其中系统毒性可表现为急性毒性、亚慢性毒性和慢性毒性，特殊毒性包括致癌、致畸、致突变、生殖毒性和靶器官毒性。

久远题库补充答案：靶器官是外源化学物可以直接发挥毒作用的器官。决定某器官成为靶器官的因素包括：器官在体内的解剖位置和功能，特别是吸收和排泄器官；器官血液供应；是否具有特殊摄入系统；该器官代谢毒物的能力以及活化/解毒系统平衡；是否存在特殊酶或生化途径；毒物是否与特殊生物大分子结合；组织修复能力；以及对特异性损伤的易感性。

\InfoLabel{知识来源} 《1 绪论》第24页；《9 化学致癌作用》第4页；《9
化学致癌作用》第5页；相对久远的总结《毒理题库-10年更新.doc》第2题

\ContentBand{对应往年题}

\begin{pastquestion}

\QuestionLabel{【往年题1｜相对久远的总结/毒理题库-10年更新.doc｜久远题库补充】}\\
靶器官的概念及决定某器官成为靶器官的影响因素。

\end{pastquestion}

\begin{pastquestion}

\QuestionLabel{【往年题2｜相对久远的总结/毒理学题库-整理版（精华）-久远.doc｜久远题库补充】}\\
靶器官的概念及影响因素。

\end{pastquestion}

\PointSeparator

\chapter{毒理学实验基础}\label{users__fpb__26_spring__pkusph_study__tmp__pdfs__prevention-medicine-past-exams__02-ux6bd2ux7406ux5b66ux5f80ux5e74ux9898ux8003ux70b9ux6574ux7406.md__ux7b2cux4e8cux7ae0-2-ux6bd2ux7406ux5b66ux5b9eux9a8cux57faux7840}

\section{实验动物的选择原则（物种、品系、微生物控制、个体选择）}\label{users__fpb__26_spring__pkusph_study__tmp__pdfs__prevention-medicine-past-exams__02-ux6bd2ux7406ux5b66ux5f80ux5e74ux9898ux8003ux70b9ux6574ux7406.md__ux5b9eux9a8cux52a8ux7269ux7684ux9009ux62e9ux539fux5219ux7269ux79cdux54c1ux7cfbux5faeux751fux7269ux63a7ux5236ux4e2aux4f53ux9009ux62e9}

\InfoLabel{考频} 8 次

\InfoLabel{对应小节} 六、实验动物的选择和处理 /《2
毒理学实验基础》第29页''实验动物物种的选择基本原则''、第35页''个体选择''

\InfoLabel{匹配依据}
``毒理学实验动物选择原则''``实验动物的选择''``动物实验，如何选动物（实验那一章）''反复出现，对应本小节。

\ContentBand{知识点原文摘取}

（一）实验动物物种的选择基本原则：①对受试物在代谢、生物化学和毒理学特征与人最接近的物种；②敏感的物种；③自然寿命不太长的物种；④易于饲养和实验操作的物种；经济并易于获得的物种。

（二）实验动物品系的选择：近交系、杂交群动物、封闭群。（三）对实验动物微生物控制的选择：Ⅰ级普通动物、Ⅱ级清洁动物、Ⅲ级无特定病原体动物（SPF）、Ⅳ级无菌动物。（四）个体选择：①性别；②年龄和体重；③生理状态；④健康状况。

\InfoLabel{知识来源} 《2 毒理学实验基础》第29页；《2 毒理学实验基础》第33页；《2 毒理学实验基础》第34页；《2
毒理学实验基础》第35页

\ContentBand{对应往年题}

\begin{pastquestion}

\QuestionLabel{【往年题1｜21级毒理期末口试回忆.xlsx｜口试】}\\
毒理学实验动物选择原则（序号19，备注：wxt老师提问很多其他内容 迟发神经毒性用什么实验动物 小鼠基因型的分组）

\end{pastquestion}

\begin{pastquestion}

\QuestionLabel{【往年题2｜07级毒理口试题回忆.doc｜口试】}\\
毒理学实验动物的选择原则（35号题）

\end{pastquestion}

\begin{pastquestion}

\QuestionLabel{【往年题3｜09级口试题目.docx｜口试】}\\
毒理学实验动物选择原则（46，备注：4考场姚碧云老师监考，会偶尔追问，如急性和慢性选择什么动物）

\end{pastquestion}

\begin{pastquestion}

\QuestionLabel{【往年题4｜17级毒理试题汇总.xlsx｜口试】}\\
实验动物选择原则（题号22，备注：背的好不如抽的好）

\end{pastquestion}

\begin{pastquestion}

\QuestionLabel{【往年题5｜19级毒理期末口试回忆.xlsx｜口试】}\\
实验动物的选择原则（序号58）

\end{pastquestion}

\begin{pastquestion}

\QuestionLabel{【往年题6｜14级毒理考题汇总.xls｜口试】}\\
动物实验，如何选动物（实验那一章）（题号12）；实验动物的选择（题号75）

\end{pastquestion}

\begin{pastquestion}

\QuestionLabel{【往年题7｜15级毒理学考试题目汇总.xls｜口试】}\\
毒理实验的动物选择（题号12）

\end{pastquestion}

\PointSeparator

\section{化学物的联合作用（概念、类型及毒理学意义）}\label{users__fpb__26_spring__pkusph_study__tmp__pdfs__prevention-medicine-past-exams__02-ux6bd2ux7406ux5b66ux5f80ux5e74ux9898ux8003ux70b9ux6574ux7406.md__ux5316ux5b66ux7269ux7684ux8054ux5408ux4f5cux7528ux6982ux5ff5ux7c7bux578bux53caux6bd2ux7406ux5b66ux610fux4e49}

\InfoLabel{考频} 6 次

\InfoLabel{对应小节} 八、毒作用的影响因素 /《2 毒理学实验基础》第64页''化学物的联合作用''、第65页''类型''

\InfoLabel{匹配依据} ``化学物联合作用的概念、分类及毒理学意义''``联合作用的概念、生物学意义''直接对应本小节。

\ContentBand{知识点原文摘取}

化学物的联合作用：一种化学物对机体的毒性作用，可以由于同时或先后接触另一种化学物而使其所表现的联合毒性比任一单一化学物的毒性增强或减弱，毒理学将两种或两种以上的化学物对机体的交互作用称为联合作用（joint
action）。

类型：相加作用（additional joint action）、协同作用（synergistic effect）、加强作用（potentiation joint
action）、拮抗作用（antagonistic joint action）、独立作用（independent effect）。

【翻译】联合作用的类型包括：相加作用、协同作用、加强作用、拮抗作用和独立作用。

\InfoLabel{知识来源} 《2 毒理学实验基础》第64页；《2 毒理学实验基础》第65页

\ContentBand{对应往年题}

\begin{pastquestion}

\QuestionLabel{【往年题1｜21级毒理期末口试回忆.xlsx｜口试】}\\
联合作用的概念、分类及毒理学意义（序号24，备注：郝老师追问了相加和协同的概念，以及可不可以举几个协同作用的例子）

\end{pastquestion}

\begin{pastquestion}

\QuestionLabel{【往年题2｜09级口试题目.docx｜口试】}\\
化学物联合作用的概念、分类及毒理学意义（41）

\end{pastquestion}

\begin{pastquestion}

\QuestionLabel{【往年题3｜19级毒理期末口试回忆.xlsx｜口试】}\\
简述联合作用的概念、生物学意义（序号24，备注：老师会很耐心地引导性地问你）

\end{pastquestion}

\begin{pastquestion}

\QuestionLabel{【往年题4｜17级毒理试题汇总.xlsx｜口试】}\\
化学物联合作用的概念 分类 毒理学意义（题号25，追问了协同作用可能的原因，还有什么替代实验）

\end{pastquestion}

\begin{pastquestion}

\QuestionLabel{【往年题5｜20级完整试题回忆.docx｜单项选择】}\\
联合作用包括以下4种，除了 A相加作用 B独立作用 C协同作用 D拮抗作用 E诱导作用（第4题）

\end{pastquestion}

\PointSeparator

\section{毒作用的影响因素（化学结构、理化性质、机体、暴露等）}\label{users__fpb__26_spring__pkusph_study__tmp__pdfs__prevention-medicine-past-exams__02-ux6bd2ux7406ux5b66ux5f80ux5e74ux9898ux8003ux70b9ux6574ux7406.md__ux6bd2ux4f5cux7528ux7684ux5f71ux54cdux56e0ux7d20ux5316ux5b66ux7ed3ux6784ux7406ux5316ux6027ux8d28ux673aux4f53ux66b4ux9732ux7b49}

\InfoLabel{考频} 7 次

\InfoLabel{对应小节} 八、毒作用的影响因素 /《2
毒理学实验基础》第51--67页（第53页化学结构、第56页理化性质、第60--67页综合）

\InfoLabel{匹配依据}
``影响化学物毒效应的因素''``影响毒作用效应的因素''``毒作用的影响因素''``影响毒性作用的因素''反复出现，对应本小节。

\ContentBand{知识点原文摘取}

化学结构对其毒性的影响：同系物的碳原子数（碳原子数愈多毒性愈大，甲醇与甲醛除外，超过7～9个时毒性反而迅速下降）；取代基团；分子中不饱和键；构型异构（顺反异构和手性异构）。

化学物的理化性质：溶解度（水溶性愈大毒性愈大；脂溶性物质易在脂肪中蓄积，易侵犯神经系统）；挥发度；分散度；纯度和杂质。其余影响因素综合包括：暴露（途径、时间/频率、溶剂、稀释度、交叉暴露、联合暴露）；机体个体因素（遗传因素、性别、年龄、生理状态、营养/习惯、疾病）；物种、品系；环境条件（气温、气湿等）。

\InfoLabel{知识来源} 《2 毒理学实验基础》第53页；《2 毒理学实验基础》第56页；《2 毒理学实验基础》第60页；《2
毒理学实验基础》第62页；《2 毒理学实验基础》第63页；《2 毒理学实验基础》第67页

\ContentBand{对应往年题}

\begin{pastquestion}

\QuestionLabel{【往年题1｜21级毒理期末口试回忆.xlsx｜口试】}\\
影响化学物毒性的因素（序号21）；化学物毒性的影响因素（序号56）

\end{pastquestion}

\begin{pastquestion}

\QuestionLabel{【往年题2｜15级毒理学考试题目汇总.xls｜口试】}\\
化学物毒作用影响因素（题号11）；影响外源化学物毒效应的因素（题号No.11）；毒作用的影响因素（题号40）；影响毒性作用的因素（题号20）

\end{pastquestion}

\begin{pastquestion}

\QuestionLabel{【往年题3｜17级毒理试题汇总.xlsx｜口试】}\\
影响化学物毒效应的因素有哪些（题号68）；毒效应影响因素（题号11）

\end{pastquestion}

\begin{pastquestion}

\QuestionLabel{【往年题4｜14级毒理考题汇总.xls｜口试】}\\
影响毒效应的因素（题号11、74）；影响毒作用效应的因素（题号40）

\end{pastquestion}

\PointSeparator

\section{毒理学试验的对照组类型及意义}\label{users__fpb__26_spring__pkusph_study__tmp__pdfs__prevention-medicine-past-exams__02-ux6bd2ux7406ux5b66ux5f80ux5e74ux9898ux8003ux70b9ux6574ux7406.md__ux6bd2ux7406ux5b66ux8bd5ux9a8cux7684ux5bf9ux7167ux7ec4ux7c7bux578bux53caux610fux4e49}

\InfoLabel{考频} 6 次

\InfoLabel{对应小节} 四、毒理学试验设计 /《2 毒理学实验基础》第20页''对照组问题''

\InfoLabel{匹配依据}
``对照的种类和意义''``毒理学试验的四种对照及意义''``对照组的类型和意义''直接对应本小节；18级回忆文件给出了完整答案。

\ContentBand{知识点原文摘取}

①未处理对照组（空白对照组）：确定本底值，进行质量控制（常用于遗传毒理学试验）；②阴性（溶剂）对照：是染毒组比较的基础；③阳性对照：检测试验体系的有效性；④历史性对照：检测试验体系的稳定性。毒理学实践的基本手段之一就是比较分析。

\InfoLabel{知识来源} 《2 毒理学实验基础》第20页

\ContentBand{对应往年题}

\begin{pastquestion}

\QuestionLabel{【往年题1｜07级毒理口试题回忆.doc｜口试】}\\
毒性试验中四种对照方法及各自意义（17号题）；毒理学常用的对照及意义（53号题）

\end{pastquestion}

\begin{pastquestion}

\QuestionLabel{【往年题2｜09级口试题目.docx｜口试】}\\
毒理学试验的四种对照及意义（48）；毒理学常用的对照及意义（53）

\end{pastquestion}

\begin{pastquestion}

\QuestionLabel{【往年题3｜21级毒理期末口试回忆.xlsx｜口试】}\\
对照组的类型和意义（序号17）

\end{pastquestion}

\begin{pastquestion}

\QuestionLabel{【往年题4｜19级毒理期末口试回忆.xlsx｜口试】}\\
对照的种类和意义（序号17）

\end{pastquestion}

\begin{pastquestion}

\QuestionLabel{【往年题5｜18级毒理回忆by连昕瑶（搬运自经验分享）.docx｜简答答案】}\\
对照组问题：①未处理对照组（空白对照组）；②阴性（溶剂）对照；③阳性对照；④历史性对照（含各对照组的定义、用途与质量控制意义）

\end{pastquestion}

\begin{pastquestion}

\QuestionLabel{【往年题6｜20级完整试题回忆.docx｜单项选择】}\\
3Rs原则是指（第3题，相关，对照/实验设计基础）

\end{pastquestion}

\PointSeparator

\section{毒理学实验的三个基本原则及局限性}\label{users__fpb__26_spring__pkusph_study__tmp__pdfs__prevention-medicine-past-exams__02-ux6bd2ux7406ux5b66ux5f80ux5e74ux9898ux8003ux70b9ux6574ux7406.md__ux6bd2ux7406ux5b66ux5b9eux9a8cux7684ux4e09ux4e2aux57faux672cux539fux5219ux53caux5c40ux9650ux6027}

\InfoLabel{考频} 6 次

\InfoLabel{对应小节} 二、毒理学实验的原则和局限性 /《2
毒理学实验基础》第6页''基本原则（三个原则）``、第8页''利用动物实验的局限性''

\InfoLabel{匹配依据}
``毒理学实验的三个原则及其局限性''``毒理学实验的基本原则及局限性''``毒理学实验设计的三个基本原则以及局限性''直接对应本小节。

\ContentBand{知识点原文摘取}

1.基本原则（三个原则）：第一原则------化学物在实验动物产生的作用，可以外推于人；第二原则------实验动物必须暴露于高剂量受试物；第三原则------选择成年的健康（雄性和雌性未孕）实验动物和人可能的暴露途径。

3.利用动物实验的局限性（不确定性）：①实验动物和人对化学物的反应敏感性不同，有时甚至存在着质的差别；②试验中选用较大的染毒剂量，有些化学物在高剂量和低剂量的毒性作用规律并不一定一致；③毒理学试验所用动物数量有限，那些发生率很低的毒性反应在少量动物中难以发现；④实验动物一般都是实验室培育的品系，一般选用成年健康动物，反应较单一。

\InfoLabel{知识来源} 《2 毒理学实验基础》第6页；《2 毒理学实验基础》第8页

\ContentBand{对应往年题}

\begin{pastquestion}

\QuestionLabel{【往年题1｜21级毒理期末口试回忆.xlsx｜口试】}\\
毒理学实验三大基本原则及局限性（序号20）；毒理学实验的基本原则及局限性（序号91）

\end{pastquestion}

\begin{pastquestion}

\QuestionLabel{【往年题2｜17级毒理试题汇总.xlsx｜口试】}\\
毒理学试验的三个原则和局限性（题号69）；毒理学实验的三个原则及其局限性（题号21）；毒理学实验设计的三个基本原则以及局限性（题号95）

\end{pastquestion}

\begin{pastquestion}

\QuestionLabel{【往年题3｜19级毒理期末口试回忆.xlsx｜口试】}\\
毒理学实验的三个原则及其局限性（序号20）；毒理学实验的基本原则和局限性（无序号）

\end{pastquestion}

\PointSeparator

\section{动物实验结果外推到人的不确定性及如何减少}\label{users__fpb__26_spring__pkusph_study__tmp__pdfs__prevention-medicine-past-exams__02-ux6bd2ux7406ux5b66ux5f80ux5e74ux9898ux8003ux70b9ux6574ux7406.md__ux52a8ux7269ux5b9eux9a8cux7ed3ux679cux5916ux63a8ux5230ux4ebaux7684ux4e0dux786eux5b9aux6027ux53caux5982ux4f55ux51cfux5c11}

\InfoLabel{考频} 6 次

\InfoLabel{对应小节} 二、毒理学实验的原则和局限性 /《2 毒理学实验基础》第8页（并见《1 绪论》第47页 Review
Question 1）

\InfoLabel{匹配依据}
``将动物实验结果外推到人有哪些不确定性，如何消除这些不确定性''``动物毒理试验外推的不确定性''直接对应本小节及绪论复习题。

\ContentBand{知识点原文摘取}

不确定性：用实验动物的毒理学试验资料外推到人群接触的安全性时，会有很大的不确定性。不确定性的产生：①实验动物和人对化学物的反应敏感性不同，有时甚至存在着质的差别；②试验中选用较大的染毒剂量，高剂量和低剂量的毒性作用规律并不一定一致；③试验所用动物数量有限，发生率很低的毒性反应在少量动物中难以发现；④实验动物一般为实验室培育品系，反应较单一。（《1
绪论》第47页 Review Question：用动物毒性测试的结果评定人体接触化学物安全性的不确定性有哪些？如何尽量避免？）

\InfoLabel{知识来源} 《2 毒理学实验基础》第8页；《1 绪论》第47页

\ContentBand{对应往年题}

\begin{pastquestion}

\QuestionLabel{【往年题1｜19级毒理期末口试回忆.xlsx｜口试】}\\
将动物实验结果外推到人有哪些不确定性，如何消除这些不确定性（序号41）；动物实验结果外推到人的不确定性，如何减少这种不确定性（无序号）

\end{pastquestion}

\begin{pastquestion}

\QuestionLabel{【往年题2｜07级毒理口试题回忆.doc｜口试】}\\
动物毒理试验外推的不确定性（36号题）

\end{pastquestion}

\begin{pastquestion}

\QuestionLabel{【往年题3｜17级毒理试题汇总.xlsx｜口试】}\\
动物实验外推到人的不确定性，如何规避这种不确定性（题号46）；毒理学动物实验结果外推到人的不确定性有哪些？如何减少这种不确定性？（题号55）

\end{pastquestion}

\PointSeparator

\section{毒性评价试验的基本目的}\label{users__fpb__26_spring__pkusph_study__tmp__pdfs__prevention-medicine-past-exams__02-ux6bd2ux7406ux5b66ux5f80ux5e74ux9898ux8003ux70b9ux6574ux7406.md__ux6bd2ux6027ux8bc4ux4ef7ux8bd5ux9a8cux7684ux57faux672cux76eeux7684}

\InfoLabel{考频} 5 次

\InfoLabel{对应小节} 三、毒性评价试验的基本目的 /《2 毒理学实验基础》第9页

\InfoLabel{匹配依据} ``毒性评价试验的目的''``毒理学评价试验的基本目的''``毒性评价实验的目的''直接对应本小节。

\ContentBand{知识点原文摘取}

三、毒性评价试验的基本目的：1.化学物毒作用的表现和性质；2.剂量-反应(效应)研究；3.确定毒作用的靶器官；4.确定损害的可逆性；5.其它（（1）毒作用的敏感检测指标和生物学标志；（2）毒作用机制研究；（3）受试物的毒物动力学和代谢研究；（4）中毒的解救措施）。

\InfoLabel{知识来源} 《2 毒理学实验基础》第9页

\ContentBand{对应往年题}

\begin{pastquestion}

\QuestionLabel{【往年题1｜07级毒理口试题回忆.doc｜口试】}\\
毒性评价的目的（43号题）；毒理学评价试验的基本目的（29号题）

\end{pastquestion}

\begin{pastquestion}

\QuestionLabel{【往年题2｜09级口试题目.docx｜口试】}\\
毒理学评价实验的目的是什么（29）

\end{pastquestion}

\begin{pastquestion}

\QuestionLabel{【往年题3｜17级毒理试题汇总.xlsx｜口试】}\\
毒性评价实验的目的（题号65）

\end{pastquestion}

\begin{pastquestion}

\QuestionLabel{【往年题4｜15级毒理学考试题目汇总.xls｜口试】}\\
毒性试验结果的准确解释/毒理学实验结果的解释（题号42、77，相关，见结果分析）

\end{pastquestion}

\PointSeparator

\section{受试物溶剂（助溶剂）选择原则}\label{users__fpb__26_spring__pkusph_study__tmp__pdfs__prevention-medicine-past-exams__02-ux6bd2ux7406ux5b66ux5f80ux5e74ux9898ux8003ux70b9ux6574ux7406.md__ux53d7ux8bd5ux7269ux6eb6ux5242ux52a9ux6eb6ux5242ux9009ux62e9ux539fux5219}

\InfoLabel{考频} 6 次

\InfoLabel{对应小节} 五、受试物的准备 /《2 毒理学实验基础》第25页''受试物的准备''

\InfoLabel{匹配依据}
``试验溶剂选择的原则''``毒理学试验溶剂的选择原则''``毒理实验设计中选择溶剂的原则''直接对应本小节。

\ContentBand{知识点原文摘取}

受试物的准备：剂型（水溶液、油溶液或混悬液）；对溶剂和助溶剂的要求------无毒，与受试物不起反应，稳定；制剂要点；染毒容积以受试的实验动物物种或制剂来确定。（受试样品准备：①纯度及杂质成分；②化学结构和理化性质；③相似化合物的毒性参考资料；④同种同一批号；⑤成分和配方必须固定；⑥稳定性受试物应一次备齐全部实验的用量。）

\begin{correctionbox}

【勘误补充】\\
对溶剂和助溶剂的要求：\\
1. 无毒或实际无毒\\
2. 与受试物不发生反应，受试物在溶液中稳定\\
3. 对受试物的毒动学和毒效学无显著影响\\
4. 无刺激性，无气味

\end{correctionbox}

\InfoLabel{知识来源} 《2 毒理学实验基础》第24页；《2 毒理学实验基础》第25页

\ContentBand{对应往年题}

\begin{pastquestion}

\QuestionLabel{【往年题1｜07级毒理口试题回忆.doc｜口试】}\\
受试物溶剂选择的原则（无毒、稳定，这两点很重要）（26号题）

\end{pastquestion}

\begin{pastquestion}

\QuestionLabel{【往年题2｜09级口试题目.docx｜口试】}\\
溶剂的选用原则（9）；试验溶剂选择的原则（26）

\end{pastquestion}

\begin{pastquestion}

\QuestionLabel{【往年题3｜21级毒理期末口试回忆.xlsx｜口试】}\\
毒理学试验溶剂的选择原则（序号18，备注：溶剂选择老师追问了一句常用溶剂------水溶性用水，脂溶性用玉米油或食用油）；溶剂的选择原则（序号64）

\end{pastquestion}

\begin{pastquestion}

\QuestionLabel{【往年题4｜19级毒理期末口试回忆.xlsx｜口试】}\\
毒理学实验溶剂选择原则（序号18，备注：追问了常用的溶剂有哪些）

\end{pastquestion}

\begin{pastquestion}

\QuestionLabel{【往年题5｜17级毒理试题汇总.xlsx｜口试】}\\
毒理学实验溶剂选择原则（题号78）；毒理实验设计中选择溶剂的原则（题号19）

\end{pastquestion}

\PointSeparator

\section{脏器系数（脏器相对重量）的概念、意义及影响因素}\label{users__fpb__26_spring__pkusph_study__tmp__pdfs__prevention-medicine-past-exams__02-ux6bd2ux7406ux5b66ux5f80ux5e74ux9898ux8003ux70b9ux6574ux7406.md__ux810fux5668ux7cfbux6570ux810fux5668ux76f8ux5bf9ux91cdux91cfux7684ux6982ux5ff5ux610fux4e49ux53caux5f71ux54cdux56e0ux7d20}

\InfoLabel{考频} 7 次

\InfoLabel{对应小节} 七、实验动物处死及生物标本采集 /《2 毒理学实验基础》第47页（并见《6
亚慢性和慢性毒性试验》第23--25页''脏器相对重量（脏器系数）``）

\InfoLabel{匹配依据}
``脏器系数的概念意义和影响因素''``脏器系数意义及影响因素''反复出现，对应本小节大体解剖称重计算脏器系数及亚慢性章节脏器系数定义。

\ContentBand{知识点原文摘取}

大体解剖在实验动物处死后半小时内进行，解剖方法采用胸腔、腹腔脏器联合取出法。应观察有关脏器的外形和表面情况、颜色、边界和大小、质地、切面。对指定的脏器称重，并计算脏器系数。

（《6
亚慢性和慢性毒性试验》第23页）脏器相对重量（脏器系数）概念：脏器重与体重比值（脏体比），脏器重与脑重比值（脏脑比）。

\InfoLabel{知识来源} 《2 毒理学实验基础》第47页；《6 亚慢性和慢性毒性试验》第23页；《6
亚慢性和慢性毒性试验》第25页

\ContentBand{对应往年题}

\begin{pastquestion}

\QuestionLabel{【往年题1｜07级毒理口试题回忆.doc｜口试】}\\
脏器系数的概念和意义（4号题）；脏器系数的概念意义以及什么因素会影响脏器系数（4号题）；脏器系数的意义及影响因素（57号题）

\end{pastquestion}

\begin{pastquestion}

\QuestionLabel{【往年题2｜09级口试题目.docx｜口试】}\\
简述脏器系数的意义及其影响因素（56）；脏器系数的毒理学意义和影响因素（57）

\end{pastquestion}

\begin{pastquestion}

\QuestionLabel{【往年题3｜21级毒理期末口试回忆.xlsx｜口试】}\\
脏器系数意义及影响因素（序号20、66，备注：脏器系数追问了什么情况会让脏器质量增加------水肿、肿瘤、充血）

\end{pastquestion}

\begin{pastquestion}

\QuestionLabel{【往年题4｜19级毒理期末口试回忆.xlsx｜口试】}\\
脏器系数变化的意义及其变化的影响因素（序号20）；脏器系数的概念意义和影响因素（序号66）

\end{pastquestion}

\begin{pastquestion}

\QuestionLabel{【往年题5｜17级毒理试题汇总.xlsx｜口试】}\\
脏器系数的意义和影响脏器系数的因素（题号76）

\end{pastquestion}

\PointSeparator

\section{毒理学实验结果的处理与分析（统计学意义与生物学意义）}\label{users__fpb__26_spring__pkusph_study__tmp__pdfs__prevention-medicine-past-exams__02-ux6bd2ux7406ux5b66ux5f80ux5e74ux9898ux8003ux70b9ux6574ux7406.md__ux6bd2ux7406ux5b66ux5b9eux9a8cux7ed3ux679cux7684ux5904ux7406ux4e0eux5206ux6790ux7edfux8ba1ux5b66ux610fux4e49ux4e0eux751fux7269ux5b66ux610fux4e49}

\InfoLabel{考频} 4 次

\InfoLabel{对应小节} 九、毒理学实验结果的处理和分析 /《2 毒理学实验基础》第68--74页（并见《6
亚慢性和慢性毒性试验》第33--36页）

\InfoLabel{匹配依据}
``毒性试验结果解释''``对毒理学实验结果的解释''``毒理学实验结果分析''对应本小节统计学意义与生物学意义、结果判定原则。

\ContentBand{知识点原文摘取}

认真贯彻实验设计的对照、随机和重复的原则。统计学意义和生物学意义：（1）纵向比较，此参数的改变有无剂量-反应关系；（2）横向比较，此参数的改变是否伴有其它相关参数的改变；（3）与历史性对照比较。

（《6》第34页）差别是否是处理相关效应：1.明显的剂量反应关系；2.个别动物结果推测是离群值；3.测定方法有固定的不精确性；4.在正常的生物学变异范围内；5.缺乏生物学合理性。

\InfoLabel{知识来源} 《2 毒理学实验基础》第70页；《2 毒理学实验基础》第73页；《2 毒理学实验基础》第74页；《6
亚慢性和慢性毒性试验》第34页

\ContentBand{对应往年题}

\begin{pastquestion}

\QuestionLabel{【往年题1｜15级毒理学考试题目汇总.xls｜口试】}\\
毒性试验结果的准确解释（题号42）；对毒理学实验结果的解释（题号77）

\end{pastquestion}

\begin{pastquestion}

\QuestionLabel{【往年题2｜14级毒理考题汇总.xls｜口试】}\\
毒性试验结果解释（题号14）；毒理学实验结果分析（题号42）

\end{pastquestion}

\begin{pastquestion}

\QuestionLabel{【往年题3｜13级（见07级口试题回忆.doc 引用）｜口试】}\\
毒理学试验结果的统计学/生物学意义判定（结合纵向/横向/历史性对照比较）

\end{pastquestion}

\PointSeparator

\section{受试样品与染毒制剂准备要点（久远题库补充）}\label{users__fpb__26_spring__pkusph_study__tmp__pdfs__prevention-medicine-past-exams__02-ux6bd2ux7406ux5b66ux5f80ux5e74ux9898ux8003ux70b9ux6574ux7406.md__ux53d7ux8bd5ux6837ux54c1ux4e0eux67d3ux6bd2ux5236ux5242ux51c6ux5907ux8981ux70b9ux4e45ux8fdcux9898ux5e93ux8865ux5145}

\InfoLabel{考频} 3 次

\InfoLabel{对应小节} 五、受试物的准备 /《2 毒理学实验基础》第24--26页

\InfoLabel{匹配依据}
久远题库中''准备染毒制剂时的要点''与现有''受试物溶剂选择原则''相邻但范围更宽，涉及样品纯度、稳定性、配方、pH、等渗和制剂一致性，宜单列补充。

\ContentBand{知识点原文摘取}

受试样品的准备：纯度及杂质成分；化学结构和理化性质（挥发性和溶解性）；相似化合物的毒性参考资料；同种同一批号；成分和配方必须固定；稳定性受试物应一次备齐全部实验的用量。受试物的准备：剂型可为水溶液、油溶液或混悬液；溶剂和助溶剂要求无毒、与受试物不起反应、稳定；染毒容积以实验动物物种或制剂确定。

制剂要点：制备时加热受试物不应接近或达到改变其化学性质或物理性质的温度；评价固体受试物皮肤毒性时，应保持其形状和颗粒大小；多成分受试物应按配方配制，使染毒制剂准确反映原混合物；制剂应保持化学稳定性和一致性；应减少总试验容积，溶剂或赋形剂不应过多；制剂应易于准确染毒；pH
应为 5--9；不应用酸碱解离受试物；非胃肠道染毒时受试物应尽可能接近等渗。

\InfoLabel{知识来源} 《2 毒理学实验基础》第24页；《2 毒理学实验基础》第25页；《2 毒理学实验基础》第26页

\ContentBand{对应往年题}

\begin{pastquestion}

\QuestionLabel{【往年题1｜相对久远的总结/毒理题库-10年更新.doc｜久远题库补充】}\\
准备染毒制剂时的要点。

\end{pastquestion}

\begin{pastquestion}

\QuestionLabel{【往年题2｜相对久远的总结/毒理题库-15年更新.doc｜久远题库补充】}\\
准备染毒制剂时的要点；受试物溶剂的选择原则。

\end{pastquestion}

\begin{pastquestion}

\QuestionLabel{【往年题3｜相对久远的总结/毒理题库-12年更新.doc｜久远题库补充】}\\
受试物溶剂的选择原则。

\end{pastquestion}

\PointSeparator

\chapter{/ 《3 Biotransportation and
Biotransformation\_WXIAO》（生物转运）}\label{users__fpb__26_spring__pkusph_study__tmp__pdfs__prevention-medicine-past-exams__02-ux6bd2ux7406ux5b66ux5f80ux5e74ux9898ux8003ux70b9ux6574ux7406.md__ux7b2cux4e09ux7ae0-3-biotransportation-and-biotransformation_wxiaoux751fux7269ux8f6cux8fd0}

\section{脂/水分配系数（lipid-water partition
coefficient）的概念及其在吸收、分布、转运、排泄中的意义}\label{users__fpb__26_spring__pkusph_study__tmp__pdfs__prevention-medicine-past-exams__02-ux6bd2ux7406ux5b66ux5f80ux5e74ux9898ux8003ux70b9ux6574ux7406.md__ux8102ux6c34ux5206ux914dux7cfbux6570lipid-water-partition-coefficientux7684ux6982ux5ff5ux53caux5176ux5728ux5438ux6536ux5206ux5e03ux8f6cux8fd0ux6392ux6cc4ux4e2dux7684ux610fux4e49}

\InfoLabel{考频} 12 次

\InfoLabel{对应小节} Disposition of xenobiotics: Biotransportation /《3 Biotransportation and
Biotransformation\_WXIAO》第7页''Lipid-water partition coefficient''

\InfoLabel{匹配依据}
``脂水分配系数（英文）的概念和毒理学意义''``脂水分配系数在吸收、转运、排泄过程中的意义''反复出现，与本小节完全对应；18级回忆文件给出完整答案。

\ContentBand{知识点原文摘取}

Lipid-water partition coefficient（LWPC）：{[}lipid phase{]}/{[}water phase{]} at equilibrium（non-ionized
form）。LWPC ↑ → Lipid solubility ↑ → Across biomembranes ↑ → Accumulation in lipids ↑；LWPC ↓ → Water
solubility ↑ → Across biomembranes ↓ → Urinary secretion ↑。

【翻译】脂/水分配系数是指非离子形式化学物在平衡状态下脂相浓度与水相浓度的比值。脂/水分配系数升高时，脂溶性增强，越容易通过生物膜，也越容易在脂质中蓄积；脂/水分配系数降低时，水溶性增强，通过生物膜能力下降，尿排泄倾向增加。

（18级回忆答案）脂/水分配系数是表示脂溶性（亲脂性）的一个重要参数，具体指当一种物质在脂相和水相中分配达到均衡时，其在脂相和水相中溶解度的比值。意义：脂/水分配系数高的化学物易于吸收，易于分布和蓄积，易经肝肠循环重吸收，不易排泄，而且多需要经历生物转化成水溶性代谢产物才能排泄。易在脂肪中蓄积，易侵犯神经系统。

\InfoLabel{知识来源} 《3 Biotransportation and Biotransformation\_WXIAO》第7页

\ContentBand{对应往年题}

\begin{pastquestion}

\QuestionLabel{【往年题1｜07级毒理口试题回忆.doc｜口试】}\\
脂水分配系数的概念，意义（2号题）；脂水分配系数（英文）的概念及意义（15号、72号题，并涉及ADME各过程意义）；脂水分配系数对ade的影响（28号题）

\end{pastquestion}

\begin{pastquestion}

\QuestionLabel{【往年题2｜09级口试题目.docx｜口试】}\\
脂水分配系数（英文）的概念和毒理学意义（7、23）；脂水分配系数对化学物在体内的转运、吸收和分布的影响（28，备注：脂溶性大的易经肝肠循环重吸收）

\end{pastquestion}

\begin{pastquestion}

\QuestionLabel{【往年题3｜21级毒理期末口试回忆.xlsx｜口试】}\\
lipid/water partition
coefficient的概念和毒理学意义（序号10）；脂水分配系数在吸收、转运、排泄过程中的意义（序号23、66）；脂水分配系数的概念和毒理学意义（序号56）

\end{pastquestion}

\begin{pastquestion}

\QuestionLabel{【往年题4｜17级毒理试题汇总.xlsx｜口试】}\\
脂水分配系数（英文）的概念和毒理学意义（题号86）；脂水分配系数的概念和在吸收分布排泄的意义（题号66）

\end{pastquestion}

\begin{pastquestion}

\QuestionLabel{【往年题5｜18级毒理回忆by连昕瑶（搬运自经验分享）.docx｜简答答案】}\\
脂/水分配系数（lipid/water partition
coefficient）：概念与意义（高脂水分配系数易吸收、分布蓄积、经肝肠循环重吸收、不易排泄、易在脂肪蓄积、易侵犯神经系统）

\end{pastquestion}

\PointSeparator

\section{生物（膜）屏障：血脑屏障、血-脑脊液屏障、血睾屏障、胎盘屏障、血眼屏障及毒理学意义}\label{users__fpb__26_spring__pkusph_study__tmp__pdfs__prevention-medicine-past-exams__02-ux6bd2ux7406ux5b66ux5f80ux5e74ux9898ux8003ux70b9ux6574ux7406.md__ux751fux7269ux819cux5c4fux969cux8840ux8111ux5c4fux969cux8840-ux8111ux810aux6db2ux5c4fux969cux8840ux777eux5c4fux969cux80ceux76d8ux5c4fux969cux8840ux773cux5c4fux969cux53caux6bd2ux7406ux5b66ux610fux4e49}

\InfoLabel{考频} 9 次

\InfoLabel{对应小节} Barriers: protective against xenobiotics toxicity /《3 Biotransportation and
Biotransformation\_WXIAO》第23--24页

\InfoLabel{匹配依据}
``简述生物屏障及毒理学意义''``重要的生物膜屏障举例及其毒理学意义''``血脑屏障和胎盘屏障的概念和毒理学意义''反复出现，对应本小节五种屏障。

\ContentBand{知识点原文摘取}

Barriers: protective against xenobiotics toxicity------Blood-brain barrier；Blood-cerebrospinal fluid
barrier；Blood-testis barrier；Blood-placental barrier；Blood-eye barrier。

【翻译】生物屏障可保护机体免受外源化学物毒性影响，包括血脑屏障、血-脑脊液屏障、血睾屏障、血胎盘屏障和血眼屏障。

BBB and BCSFB limit distribution to the brain；Active transporters determine the amount of xenobiotics in the
brain；BBB is not fully developed at birth → more susceptible in newborns。Blood-placental barrier------The
factors determine the permeability of the placenta: thickness, surface area, carrier systems, and
lipid-protein concentration of the membranes；Active transport / Passive diffusion / Permeable size
\textless{} 1 kDa / Biotransformation。

【翻译】血脑屏障和血-脑脊液屏障限制外源化学物向脑内分布；主动转运体决定脑内外源化学物的数量；出生时血脑屏障尚未完全发育，因此新生儿更易受影响。血胎盘屏障的通透性受胎盘厚度、表面积、载体系统以及膜中脂质和蛋白浓度等因素影响；相关过程包括主动转运、被动扩散、小于
1 kDa 的分子通过以及生物转化。

\begin{correctionbox}

【勘误补充】\\
生物屏障可限制外源化学物进入某些重要器官和组织，具有保护机体、减少靶器官毒性的重要意义。\\
血脑屏障和血-脑脊液屏障可限制毒物进入中枢神经系统，但新生儿血脑屏障发育不完善，对神经毒物更敏感。\\
血睾屏障可限制毒物进入生精上皮，保护生殖细胞，但一旦毒物通过，可能造成生殖毒性。\\
血胎盘屏障可部分阻止毒物进入胎儿体内，但并非绝对屏障，小分子、脂溶性或可经转运体转运的毒物仍可通过，导致胚胎毒性、发育毒性或致畸作用。\\
血眼屏障可限制毒物进入眼内组织，保护视觉器官，但也会影响眼部毒性物质的分布和清除。

\end{correctionbox}

\InfoLabel{知识来源} 《3 Biotransportation and Biotransformation\_WXIAO》第23页；《3 Biotransportation and
Biotransformation\_WXIAO》第24页

\ContentBand{对应往年题}

\begin{pastquestion}

\QuestionLabel{【往年题1｜07级毒理口试题回忆.doc｜口试】}\\
血脑屏障和胎盘屏障的概念和毒理学意义（5号题）

\end{pastquestion}

\begin{pastquestion}

\QuestionLabel{【往年题2｜09级口试题目.docx｜口试】}\\
胎盘屏障和血脑屏障的毒理学意义（5）

\end{pastquestion}

\begin{pastquestion}

\QuestionLabel{【往年题3｜21级毒理期末口试回忆.xlsx｜口试】}\\
简述生物屏障及毒理学意义（序号14）；重要的生物膜屏障举例及其毒理学意义（序号60，备注：郝老师追问了具体生物膜屏障，比如血脑屏障的生物学构成）；生物膜屏障和毒理学意义（序号92）

\end{pastquestion}

\begin{pastquestion}

\QuestionLabel{【往年题4｜19级毒理期末口试回忆.xlsx｜口试】}\\
重要的生物学屏障（序号60）

\end{pastquestion}

\begin{pastquestion}

\QuestionLabel{【往年题5｜17级毒理试题汇总.xlsx｜口试】}\\
生物学屏障（题号82）；机体内主要的生物膜屏障及毒理学意义（题号15）

\end{pastquestion}

\begin{pastquestion}

\QuestionLabel{【往年题6｜15级毒理学考试题目汇总.xls｜口试】}\\
生物屏障类型，毒理学意义（题号69）；生物屏障（题号35）；人体的生物屏障（题号69）

\end{pastquestion}

\begin{pastquestion}

\QuestionLabel{【往年题7｜14级毒理考题汇总.xls｜口试】}\\
生物屏障有哪些及其毒理学意义（题号35）；生理屏障在毒理学上的意义（题号6）；生物屏障有哪些？毒理学意义（题号81）；生理屏障在毒理学上的意义（题号6）

\end{pastquestion}

\PointSeparator

\section{影响经胃肠道（消化道）吸收的因素}\label{users__fpb__26_spring__pkusph_study__tmp__pdfs__prevention-medicine-past-exams__02-ux6bd2ux7406ux5b66ux5f80ux5e74ux9898ux8003ux70b9ux6574ux7406.md__ux5f71ux54cdux7ecfux80c3ux80a0ux9053ux6d88ux5316ux9053ux5438ux6536ux7684ux56e0ux7d20}

\InfoLabel{考频} 9 次

\InfoLabel{对应小节} Disposition of xenobiotics: Absorption /《3 Biotransportation and
Biotransformation\_WXIAO》第11页''影响胃肠道对毒物吸收的因素''

\InfoLabel{匹配依据}
``影响消化道吸收的因素''``概述影响胃肠道吸收的因素''反复出现，与本小节8项因素完全对应（往年答案另需补充环境因素气温气湿，见《2》第63页）。

\ContentBand{知识点原文摘取}

影响胃肠道对毒物吸收的因素：1.毒物的理化性质（脂溶性等）；2.胃肠道存留食物的多少；3.毒物在胃肠道的停留时间；4.肠道的吸收面积和吸收能力；5.胃肠道局部的pH；6.胃肠道的分泌能力；7.肠道菌群；8.化学物间的相互作用（EDTA）。胃肠道经吸收的主要方式：Simple
diffusion（the most common way）；Carrier transport；Cytosis（nanoparticles and particular matter）。

【翻译】胃肠道吸收的主要方式包括：简单扩散，这是最常见方式；载体转运；以及胞吞/胞吐，主要见于纳米颗粒和颗粒性物质。

\InfoLabel{知识来源} 《3 Biotransportation and Biotransformation\_WXIAO》第11页；《3 Biotransportation and
Biotransformation\_WXIAO》第12页

\ContentBand{对应往年题}

\begin{pastquestion}

\QuestionLabel{【往年题1｜07级毒理口试题回忆.doc｜口试】}\\
消化道吸收化学物影响因素（14号题）；影响消化道吸收的因素（33号题）

\end{pastquestion}

\begin{pastquestion}

\QuestionLabel{【往年题2｜09级口试题目.docx｜口试】}\\
影响胃肠道吸收的因素（14）；概述影响胃肠道吸收的因素（33）

\end{pastquestion}

\begin{pastquestion}

\QuestionLabel{【往年题3｜15级毒理学考试题目汇总.xls｜口试】}\\
消化道，呼吸道吸收影响因素（题号5）；影响消化道和呼吸道吸收的因素（题号34）；化学物消化道吸收与呼吸道吸收的影响因素（题号68）；呼吸道消化道吸收影响因素（题号34）

\end{pastquestion}

\begin{pastquestion}

\QuestionLabel{【往年题4｜14级毒理考题汇总.xls｜口试】}\\
经消化道和呼吸道吸收的影响因素（题号5）；消化道吸收和呼吸道吸收的影响因素（题号68）

\end{pastquestion}

\PointSeparator

\section{影响经呼吸道吸收的因素}\label{users__fpb__26_spring__pkusph_study__tmp__pdfs__prevention-medicine-past-exams__02-ux6bd2ux7406ux5b66ux5f80ux5e74ux9898ux8003ux70b9ux6574ux7406.md__ux5f71ux54cdux7ecfux547cux5438ux9053ux5438ux6536ux7684ux56e0ux7d20}

\InfoLabel{考频} 7 次

\InfoLabel{对应小节} 经呼吸道吸收（Inhalation）/《3 Biotransportation and Biotransformation\_WXIAO》第14--16页

\InfoLabel{匹配依据}
``呼吸道吸收的影响因素''``影响呼吸道对外源化学物吸收的因素''反复出现，对应本小节气体/蒸气/气溶胶颗粒物吸收影响因素。

\ContentBand{知识点原文摘取}

经呼吸道吸收（Inhalation）：Gases and vapors------Simple diffusion；Determining factors：1. Water
solubility（高水溶性如SO2、Cl2在上呼吸道；低水溶性如NO2在支气管、细支气管和肺泡）；2.
Reactivity（高反应性如Cl2在上呼吸道）；3. Vapor pressure；4. Blood-gas partition
coefficient（BGPC，血中与肺泡中ΔC）。Aerosols and Particles（PM、SiO2、Asbestos）：Mainly via
cytosis；Size决定靶位；Water solubility；Composition；Surface structure。

【翻译】经呼吸道吸收中，气体和蒸气主要通过简单扩散吸收，决定因素包括水溶性、反应性、蒸气压和血气分配系数。气溶胶和颗粒物主要通过胞吞/胞吐吸收，其靶部位受粒径、水溶性、组成和表面结构影响。

\InfoLabel{知识来源} 《3 Biotransportation and Biotransformation\_WXIAO》第14页；《3 Biotransportation and
Biotransformation\_WXIAO》第15页；《3 Biotransportation and Biotransformation\_WXIAO》第16页

\ContentBand{对应往年题}

\begin{pastquestion}

\QuestionLabel{【往年题1｜07级毒理口试题回忆.doc｜口试】}\\
呼吸道吸收的影响因素（39号题，备注：张宝旭老师问呼吸道那题还有没有其他因素，补充了环境方面的因素）

\end{pastquestion}

\begin{pastquestion}

\QuestionLabel{【往年题2｜09级口试题目.docx｜口试】}\\
外源化合物经呼吸道吸收的影响因素（34）；影响呼吸道对外源化学物吸收的因素（39）

\end{pastquestion}

\begin{pastquestion}

\QuestionLabel{【往年题3｜15级毒理学考试题目汇总.xls｜口试】}\\
呼吸道消化道吸收影响因素（题号34，备注：郝老师额外问了可吸收颗粒物的粒径和吸收部位）

\end{pastquestion}

\begin{pastquestion}

\QuestionLabel{【往年题4｜14级毒理考题汇总.xls｜口试】}\\
消化道吸收和呼吸道吸收的影响因素（题号68）

\end{pastquestion}

\PointSeparator

\section{影响经皮肤吸收的因素}\label{users__fpb__26_spring__pkusph_study__tmp__pdfs__prevention-medicine-past-exams__02-ux6bd2ux7406ux5b66ux5f80ux5e74ux9898ux8003ux70b9ux6574ux7406.md__ux5f71ux54cdux7ecfux76aeux80a4ux5438ux6536ux7684ux56e0ux7d20}

\InfoLabel{考频} 5 次

\InfoLabel{对应小节} 经皮肤吸收 /《3 Biotransportation and
Biotransformation\_WXIAO》第17--18页''影响皮肤吸收的因素''

\InfoLabel{匹配依据}
``外源化学物经皮吸收的影响因素''``影响皮肤吸收的因素''直接对应本小节（往年答案需补充环境因素气温气湿、染毒途径）。

\ContentBand{知识点原文摘取}

影响皮肤吸收的因素：1.皮肤------角质层完整性、水份含量、厚度、面积和酸碱度；2.外源化学物------脂溶性、溶剂、分子量（\textless400
Da）和浓度；3.皮肤的解剖和生理特性------温度、血流量、年龄、种族和种属。大小鼠及兔的皮肤较人的皮肤更易通透；豚鼠、猪和猴子的皮肤通透性则与人相似。外源化学物通过皮肤吸收的途径：角质层吸收（大部分亲脂性毒物，简单扩散）；毛囊吸收；汗腺管吸收；破损皮肤部位吸收；真皮吸收（简单扩散）。

\InfoLabel{知识来源} 《3 Biotransportation and Biotransformation\_WXIAO》第17页；《3 Biotransportation and
Biotransformation\_WXIAO》第18页

\ContentBand{对应往年题}

\begin{pastquestion}

\QuestionLabel{【往年题1｜07级毒理口试题回忆.doc｜口试】}\\
外源化学物经皮吸收的影响因素（还要答环境因素，气温气湿等）（未知题号第1题）

\end{pastquestion}

\begin{pastquestion}

\QuestionLabel{【往年题2｜09级口试题目.docx｜口试】}\\
外源化学物经皮肤吸收的影响因素（题号未知）；影响皮肤吸收的因素（题号未知，备注：魏老师说第4题答案不全，还要补充染毒途径）

\end{pastquestion}

\PointSeparator

\section{贮存库（storage
depot）的概念、毒理学意义及机体常见贮存库}\label{users__fpb__26_spring__pkusph_study__tmp__pdfs__prevention-medicine-past-exams__02-ux6bd2ux7406ux5b66ux5f80ux5e74ux9898ux8003ux70b9ux6574ux7406.md__ux8d2eux5b58ux5e93storage-depotux7684ux6982ux5ff5ux6bd2ux7406ux5b66ux610fux4e49ux53caux673aux4f53ux5e38ux89c1ux8d2eux5b58ux5e93}

\InfoLabel{考频} 5 次

\InfoLabel{对应小节} Target organ and storage/accumulation depot /《3 Biotransportation and
Biotransformation\_WXIAO》第21--23页

\InfoLabel{匹配依据} ``storage depot的概念和毒理学意义，机体常见的storage
depot有哪些''``贮存库概念意义''直接对应本小节。

\ContentBand{知识点原文摘取}

Storage depot：the site where a toxicant is accumulated but is NOT the major site of toxicity（e.g., Bone for
lead and cadmium；Fat for insecticide DDT）。Toxicological significance: Friend vs.~Foe------（1）buffer depot
of acute toxicity；（2）releasing depot of chronic toxicity。Plasma proteins as storage depot（α1-acid
glycoprotein、α/β-lipoproteins、β-globulin；Bound toxicants less toxic，Free toxicants more toxic）；Liver and
Kidney as storage depots；Fat as storage
depot（高亲脂性毒物，肥胖个体蓄积更多但毒性较轻，饥饿等扰动时可突然释放入血）；Bone as storage
depot（F-、Pb，约90\%铅存于骨）。

【翻译】贮存库是毒物蓄积但并非主要毒作用发生的部位，例如铅和镉可蓄积于骨，杀虫剂 DDT
可蓄积于脂肪。其毒理学意义具有双重性：一方面可缓冲急性毒性，另一方面也可作为慢性毒性的释放库。血浆蛋白也可作为贮存库，结合状态的毒物毒性较低，游离状态的毒物毒性较高；肝、肾、脂肪和骨也可作为不同毒物的贮存库。

\InfoLabel{知识来源} 《3 Biotransportation and Biotransformation\_WXIAO》第21页；《3 Biotransportation and
Biotransformation\_WXIAO》第22页；《3 Biotransportation and Biotransformation\_WXIAO》第23页

\ContentBand{对应往年题}

\begin{pastquestion}

\QuestionLabel{【往年题1｜07级毒理口试题回忆.doc｜口试】}\\
storage depot（37号题）

\end{pastquestion}

\begin{pastquestion}

\QuestionLabel{【往年题2｜09级口试题目.docx｜口试】}\\
贮存库（21）

\end{pastquestion}

\begin{pastquestion}

\QuestionLabel{【往年题3｜21级毒理期末口试回忆.xlsx｜口试】}\\
storage depot的概念和毒理学意义，机体常见的storage
depot有哪些（序号59）；郝老师追问了贮存库的危害（序号44备注）

\end{pastquestion}

\begin{pastquestion}

\QuestionLabel{【往年题4｜19级毒理期末口试回忆.xlsx｜口试】}\\
storage depot的概念、毒理学意义，以及主要的storage depot（序号59）

\end{pastquestion}

\begin{pastquestion}

\QuestionLabel{【往年题5｜17级毒理试题汇总.xlsx｜口试】}\\
贮存库概念意义（题号83）

\end{pastquestion}

\PointSeparator

\section{肠肝循环（enterohepatic cycling /
recirculation）的概念及毒理学意义}\label{users__fpb__26_spring__pkusph_study__tmp__pdfs__prevention-medicine-past-exams__02-ux6bd2ux7406ux5b66ux5f80ux5e74ux9898ux8003ux70b9ux6574ux7406.md__ux80a0ux809dux5faaux73afenterohepatic-cycling-recirculationux7684ux6982ux5ff5ux53caux6bd2ux7406ux5b66ux610fux4e49}

\InfoLabel{考频} 5 次

\InfoLabel{对应小节} Bile acid excretion / Enterohepatic cycling /《3 Biotransportation and
Biotransformation\_WXIAO》第31页（并见经粪便排泄第30页）

\InfoLabel{匹配依据} ``肠肝循环（英文）概念及意义''``entrohepatic
recirculation的概念及毒理学意义''直接对应本小节。

\ContentBand{知识点原文摘取}

Enterohepatic cycling：Bile acid de novo synthesis 0.2-0.6 g/day = bile acid secretion by feces = 5\% of total
bile acid pool；The remaining 95\% of bile acids in the pool（\textasciitilde3 g）are recycled 4 to 12
times/day。经粪便排泄：未吸收的化学物；胆汁排泄（大分子量化学物）；肠内排泄；肠道菌群（生物转化）。Bile acid
excretion：High molecular weight conjugates（e.g.~GSH-conjugated）；Rapidly excreted into bile acid（higher
concentration in bile vs.~plasma）。

【翻译】肠肝循环中，胆汁酸每日新合成量约 0.2-0.6 g，与经粪便排出的胆汁酸量相当，约占胆汁酸总池的 5\%；其余约
95\% 的胆汁酸（约 3 g）每天可循环 4-12
次。胆汁酸排泄主要涉及高分子量结合物，例如谷胱甘肽结合物，这些物质可快速排入胆汁，且胆汁中浓度高于血浆。

\begin{correctionbox}

【勘误补充】\\
肠肝循环 enterohepatic
circulation：在下段肠道中，在黏膜和肠道菌群水解酶的作用下，结合物会被分解，外源化学物再次游离，被肠道吸收，经肝门静脉重新入肝，使得外源化学物在血液中的持续时间延长，代谢变化增多

\end{correctionbox}

\InfoLabel{知识来源} 《3 Biotransportation and Biotransformation\_WXIAO》第30页；《3 Biotransportation and
Biotransformation\_WXIAO》第31页

\ContentBand{对应往年题}

\begin{pastquestion}

\QuestionLabel{【往年题1｜07级毒理口试题回忆.doc｜口试】}\\
entrohepatic recirculation的概念及毒理学意义（18号题）；肠肝循环（英语）概念及意义（27号题）

\end{pastquestion}

\begin{pastquestion}

\QuestionLabel{【往年题2｜09级口试题目.docx｜口试】}\\
肠肝循环，概念和毒理学意义（注意不是肝肠）（16）；肠肝循环（英文）概念及意义（27，备注：王琪老师问肝肠循环确实是生物半衰期延长吗）；肠肝循环的概念及意义（53）

\end{pastquestion}

\PointSeparator

\section{首过消除（first-pass
elimination/effect）的概念和毒理学意义}\label{users__fpb__26_spring__pkusph_study__tmp__pdfs__prevention-medicine-past-exams__02-ux6bd2ux7406ux5b66ux5f80ux5e74ux9898ux8003ux70b9ux6574ux7406.md__ux9996ux8fc7ux6d88ux9664first-pass-eliminationeffectux7684ux6982ux5ff5ux548cux6bd2ux7406ux5b66ux610fux4e49}

\InfoLabel{考频} 3 次

\InfoLabel{对应小节} 首过消除（First-pass elimination）/《3 Biotransportation and
Biotransformation\_WXIAO》第13页

\InfoLabel{匹配依据} ``first-pass effect的概念和毒理学意义''``first pass
effect概念和毒理学意义''直接对应本小节。

\ContentBand{知识点原文摘取}

首过消除（First-pass elimination）：外源化学物口服后，经胃肠道到达肝脏，一部分化学物在代谢酶作用下被代谢掉。A
high first-pass effect will limit exposure and minimize toxic potential。

【翻译】较强的首过效应会限制机体暴露水平，并降低潜在毒性。

\begin{correctionbox}

【勘误补充】\\
首过效应 first pass
effect：由于消化道血液循环的特点，从胃和肠吸收到局部血管的物质，都要汇入肝门静脉，再进入体循环。肝会代谢和排泄一部分外源化学物，使其未能进入体循环，这种现象称为首过效应。首过效应会使化学物的原型减少，而增加部分代谢产物，另一部分代谢产物会排入胆汁

\end{correctionbox}

\InfoLabel{知识来源} 《3 Biotransportation and Biotransformation\_WXIAO》第13页

\ContentBand{对应往年题}

\begin{pastquestion}

\QuestionLabel{【往年题1｜07级毒理口试题回忆.doc｜口试】}\\
first-pass effect的概念和毒理学意义（44号题）

\end{pastquestion}

\begin{pastquestion}

\QuestionLabel{【往年题2｜09级口试题目.docx｜口试】}\\
first pass effect概念和毒理学意义（44）

\end{pastquestion}

\PointSeparator

\section{机体对外源化学物的处置过程（ADME / ADE +
M）及其毒理学意义}\label{users__fpb__26_spring__pkusph_study__tmp__pdfs__prevention-medicine-past-exams__02-ux6bd2ux7406ux5b66ux5f80ux5e74ux9898ux8003ux70b9ux6574ux7406.md__ux673aux4f53ux5bf9ux5916ux6e90ux5316ux5b66ux7269ux7684ux5904ux7f6eux8fc7ux7a0badme-ade-mux53caux5176ux6bd2ux7406ux5b66ux610fux4e49}

\InfoLabel{考频} 6 次

\InfoLabel{对应小节} Disposition of xenobiotics（ADME）/《3 Biotransportation and
Biotransformation\_WXIAO》第3页、第34页（并见《4 生物转化》第2页）

\InfoLabel{匹配依据}
``机体对外源化学物的处置过程、意义''``ADME过程概念及意义''``简述机体对外源化学物处置过程及意义''直接对应本小节。

\ContentBand{知识点原文摘取}

ADME: Absorption, Distribution, Metabolism, and Excretion。Biotransportation: Absorption, Distribution, and
Excretion；Biotransformation: Metabolism；Elimination: Metabolism and
Excretion；Toxicokinetics/Pharmacokinetics: The Quantitative characterization of xenobiotic
disposition。Disposition of xenobiotics determines toxicity：External dose → Internal dose → Target dose →
Toxicity。

【翻译】ADME
指吸收、分布、代谢和排泄。生物转运包括吸收、分布和排泄；生物转化即代谢；消除包括代谢和排泄；毒代动力学/药代动力学是对外源化学物体内处置过程的定量描述。外源化学物的体内处置决定其毒性，即外暴露剂量影响内剂量，内剂量影响靶剂量，最终导致毒性。

（《4
生物转化》第2页）机体对外源化学物的处置：吸收、分布、与排泄（ADE）过程为生物转运（biotransportation）；代谢（M）过程为生物转化（biotransformation）。

\InfoLabel{知识来源} 《3 Biotransportation and Biotransformation\_WXIAO》第3页；《3 Biotransportation and
Biotransformation\_WXIAO》第34页；《4 生物转化》第2页

\ContentBand{对应往年题}

\begin{pastquestion}

\QuestionLabel{【往年题1｜21级毒理期末口试回忆.xlsx｜口试】}\\
简述机体对外源化学物的作用及其毒理学意义（序号26）；机体对外源化合物的处置过程、意义（序号44）

\end{pastquestion}

\begin{pastquestion}

\QuestionLabel{【往年题2｜19级毒理期末口试回忆.xlsx｜口试】}\\
简述机体对外源化学物处置过程及意义（序号84）；ADME过程概念及意义（序号84）；机体对化学物的处置及毒理学意义（序号26）

\end{pastquestion}

\begin{pastquestion}

\QuestionLabel{【往年题3｜17级毒理试题汇总.xlsx｜口试】}\\
机体对外源化学物质处置过程及毒理学意义（题号63）；机体对化学物的处置过程及毒理学意义（题号15）

\end{pastquestion}

\begin{pastquestion}

\QuestionLabel{【往年题4｜15级毒理学考试题目汇总.xls｜口试】}\\
外源化合物进入机体的四阶段（题号79）；外源化合物在体内发挥毒性的阶段（题号46）

\end{pastquestion}

\PointSeparator

\section{跨膜转运方式与主动转运的特点（含简单扩散）}\label{users__fpb__26_spring__pkusph_study__tmp__pdfs__prevention-medicine-past-exams__02-ux6bd2ux7406ux5b66ux5f80ux5e74ux9898ux8003ux70b9ux6574ux7406.md__ux8de8ux819cux8f6cux8fd0ux65b9ux5f0fux4e0eux4e3bux52a8ux8f6cux8fd0ux7684ux7279ux70b9ux542bux7b80ux5355ux6269ux6563}

\InfoLabel{考频} 3 次

\InfoLabel{对应小节} Disposition of xenobiotics: Biotransportation /《3 Biotransportation and
Biotransformation\_WXIAO》第5--6页、第9页''主动转运的特点及意义''

\InfoLabel{匹配依据} ``跨膜转运有几种方式，主动转运的特点''``简述简单扩散和主动转运的特点''直接对应本小节。

\ContentBand{知识点原文摘取}

Passive transport（follow ΔC）：Simple diffusion（最常见，3 NOs：No ATP、No carrier、No competitive and
saturability inhibition；3 Conditions：ΔC、high lipophilicity、non-ionization）、Filtration。Carrier
transport：Active transport、Facilitated transport。Cytosis：Endocytosis、Exocytosis。

【翻译】被动转运顺浓度差进行，包括简单扩散和滤过。简单扩散最常见，其特点是不需要
ATP、不需要载体、无竞争性和饱和性抑制；发生条件包括存在浓度差、高脂溶性和非离子化。载体转运包括主动转运和易化转运。胞吞/胞吐包括内吞和外排。

主动转运的特点及意义：①媒介：需有载体参与；②逆浓度：外源化学物可以逆浓度梯度转运；③耗能：需消耗能量，代谢抑制剂可阻止转运；④选择性：载体对转运的外源化学物有特异选择性；⑤饱和性：转运量存在极限；⑥竞争性：同一载体转运的两种外源化学物竞争性抑制。主动转运对已吸收的外源化学物在体内的不均匀分布和排泄具有重要意义。

\InfoLabel{知识来源} 《3 Biotransportation and Biotransformation\_WXIAO》第5页；《3 Biotransportation and
Biotransformation\_WXIAO》第6页；《3 Biotransportation and Biotransformation\_WXIAO》第9页；《3
Biotransportation and Biotransformation\_WXIAO》第33页

\ContentBand{对应往年题}

\begin{pastquestion}

\QuestionLabel{【往年题1｜19级毒理期末口试回忆.xlsx｜口试】}\\
跨膜转运有几种方式，主动转运的特点（序号15）

\end{pastquestion}

\begin{pastquestion}

\QuestionLabel{【往年题2｜3 Biotransportation and Biotransformation\_WXIAO（课件 Quiz，第33页）｜思考题】}\\
简述简单扩散和主动转运的特点（与往年口试同型，纳入复习）

\end{pastquestion}

\begin{pastquestion}

\QuestionLabel{【往年题3｜20级完整试题回忆.docx｜单项选择】}\\
有关弱有机酸经消化道吸收的叙述，不正确的是（第5题，涉及解离、脂水分配系数与跨膜吸收）

\end{pastquestion}

\PointSeparator

\section{经典毒动学模型与生理毒动学模型（PBTK）的区别及生理模型的优势；毒代动力学的概念与实验设计}\label{users__fpb__26_spring__pkusph_study__tmp__pdfs__prevention-medicine-past-exams__02-ux6bd2ux7406ux5b66ux5f80ux5e74ux9898ux8003ux70b9ux6574ux7406.md__ux7ecfux5178ux6bd2ux52a8ux5b66ux6a21ux578bux4e0eux751fux7406ux6bd2ux52a8ux5b66ux6a21ux578bpbtkux7684ux533aux522bux53caux751fux7406ux6a21ux578bux7684ux4f18ux52bfux6bd2ux4ee3ux52a8ux529bux5b66ux7684ux6982ux5ff5ux4e0eux5b9eux9a8cux8bbeux8ba1}

\InfoLabel{考频} 6 次

\InfoLabel{对应小节} Toxicokinetics / Kinetic models /《3 Biotransportation and
Biotransformation\_WXIAO》第63--69页

\InfoLabel{匹配依据}
``经典毒动学模型与生理动力学模型的区别''``生理毒动学模型相较于经典毒动学的优势''``药物毒代动力学实验目的和实验设计''直接对应本小节。

\ContentBand{知识点原文摘取}

Toxicokinetics：the quantitative study of ADME of
chemicals\ldots 研究内容------化学物在体内的ADME动态变化规律；毒理学研究方案的设计（剂量和染毒途径选择）；暴露-时间依赖的靶器官剂量，解释毒作用机制；将动物试验结果外推到人的安全性评价。Kinetic
models：经典毒代动力学模型（Classical toxicokinetic
model）与生理毒代动力学模型（PBTK）。经典模型：速率过程（零级消除、一级消除、非线性动力学）；房室模型（一/二/多房室）。构建PBTK模型所需参数：生理参数、理化参数（脂水分配系数、血气分配系数等）、热力学参数。

【翻译】毒代动力学是对化学物 ADME
过程进行定量研究的学科。动力学模型包括经典毒代动力学模型和生理毒代动力学模型。经典模型关注速率过程和房室模型；PBTK
模型的构建需要生理参数、理化参数和热力学参数。

\begin{correctionbox}

【勘误补充】\\
PBTK 的优点：\\
1. 能提供毒物分布于任何组织器官的时程\\
2. 通过改变参数，能预测毒物在组织器官中分布的改变\\
3. 通过体表面积换算，相同模型可以反映毒物在不同物种中的毒代动力学特征\\
4. 使复杂的治疗方案容易调整\\
PBTK 的局限性：需要较多的参数信息且不易定义或获取，一些数学方程难以掌握

\end{correctionbox}

\InfoLabel{知识来源} 《3 Biotransportation and Biotransformation\_WXIAO》第63页；《3 Biotransportation and
Biotransformation\_WXIAO》第64页；《3 Biotransportation and Biotransformation\_WXIAO》第65页；《3
Biotransportation and Biotransformation\_WXIAO》第67页；《3 Biotransportation and
Biotransformation\_WXIAO》第68页

\ContentBand{对应往年题}

\begin{pastquestion}

\QuestionLabel{【往年题1｜21级毒理期末口试回忆.xlsx｜口试】}\\
经典毒动学模型与生理动力学模型的区别（序号25）；生理毒动学模型的优点（序号55）；生理毒动学模型相较于经典毒动学的优势（序号87，备注：王旗老师引导我说了二者的区别，让我顺着推优势）

\end{pastquestion}

\begin{pastquestion}

\QuestionLabel{【往年题2｜17级毒理试题汇总.xlsx｜口试】}\\
生理毒代动力学相比经典毒代动力学的优势（题号65）；毒代动力学研究的原则（题号16，备注：郝老师追问毒代动力学和致敏实验常用的实验动物）

\end{pastquestion}

\begin{pastquestion}

\QuestionLabel{【往年题3｜15级毒理学考试题目汇总.xls｜口试】}\\
药物毒代动力学实验目的和实验设计（题号8）；药物毒代动力学研究的目的？毒动学研究实验设计包括哪些内容？（题号8）

\end{pastquestion}

\begin{pastquestion}

\QuestionLabel{【往年题4｜14级毒理考题汇总.xls｜口试】}\\
毒代动力学目的及实验设计（题号8）；比较说明生理毒代学和经典款区别特点
以及为什么生理的好（题号36）；毒动力学的意义
实验设计要点（题号71）；药物毒代学实验的研究目的？毒物代谢动力学的试验设计内容有哪些？（题号37）

\end{pastquestion}

\PointSeparator

\chapter{生物转化}\label{users__fpb__26_spring__pkusph_study__tmp__pdfs__prevention-medicine-past-exams__02-ux6bd2ux7406ux5b66ux5f80ux5e74ux9898ux8003ux70b9ux6574ux7406.md__ux7b2cux56dbux7ae0-4-ux751fux7269ux8f6cux5316}

\section{生物转化的概念、特点和毒理学意义}\label{users__fpb__26_spring__pkusph_study__tmp__pdfs__prevention-medicine-past-exams__02-ux6bd2ux7406ux5b66ux5f80ux5e74ux9898ux8003ux70b9ux6574ux7406.md__ux751fux7269ux8f6cux5316ux7684ux6982ux5ff5ux7279ux70b9ux548cux6bd2ux7406ux5b66ux610fux4e49}

\InfoLabel{考频} 9 次

\InfoLabel{对应小节} 第一节 生物转化和代谢酶 /《4
生物转化》第4页''生物转化''、第5--6页''生物转化的意义/结果''、第17--18页''代谢酶的基本特性''

\InfoLabel{匹配依据}
``生物转化的特点和毒理学意义''``生物转化概念、特点、毒理学意义''``生物转化特点意义化学反应类型''反复出现，对应本小节。

\ContentBand{知识点原文摘取}

生物转化（biotransformation）：外源化学物在机体内经多种酶催化的反应、发生化学结构改变的过程；机体对外源化学物处置的重要环节；机体维持稳态的主要机制；影响外源化学物的活性及其体内动力学。

一、生物转化的意义：1、水溶性显著增加，易于从尿或粪便中排泄（也有例外，如挥发性化合物的呼出排泄）；2、改变其化学结构，引起生物效应改变------代谢灭活（metabolic
inactivation，代谢产物活性降低）、代谢活化（metabolic
activation，活性代谢产物发挥药效或毒效）。生物转化的结果（毒理角度）：代谢解毒（毒性降低）、代谢活化（毒性增强，生成中间产物/终毒物）。代谢酶的基本特性：底物特异性广泛；结构酶（持续少量表达）；诱导酶（外源化学物可诱导合成）；多态性；立体选择性。生物转化反应类型分Ⅰ相反应和Ⅱ相反应。

\InfoLabel{知识来源} 《4 生物转化》第4页；《4 生物转化》第5页；《4 生物转化》第6页；《4 生物转化》第15页；《4
生物转化》第17页；《4 生物转化》第18页

\ContentBand{对应往年题}

\begin{pastquestion}

\QuestionLabel{【往年题1｜07级毒理口试题回忆.doc｜口试】}\\
生物转化的特点与毒理学意义（8号题）；生物转化的特点和毒理学意义（52号、54号题）；外源化合物的生物转化和毒理学意义（60号题）

\end{pastquestion}

\begin{pastquestion}

\QuestionLabel{【往年题2｜09级口试题目.docx｜口试】}\\
生物转化（52）；生物转化特点及意义（60）

\end{pastquestion}

\begin{pastquestion}

\QuestionLabel{【往年题3｜21级毒理期末口试回忆.xlsx｜口试】}\\
生物转化的特点以及意义（序号22）

\end{pastquestion}

\begin{pastquestion}

\QuestionLabel{【往年题4｜19级毒理期末口试回忆.xlsx｜口试】}\\
生物转化的特点和意义（序号22）；生物转化的特点及毒理学意义（序号57）

\end{pastquestion}

\begin{pastquestion}

\QuestionLabel{【往年题5｜17级毒理试题汇总.xlsx｜口试】}\\
生物转化概念、特点、毒理学意义（题号67、19）

\end{pastquestion}

\begin{pastquestion}

\QuestionLabel{【往年题6｜15级毒理学考试题目汇总.xls｜口试】}\\
生物转化特点，反应类型，毒理学意义（题号72）；生物转化的特点、包括哪些反应、毒理学意义（题号82）；生物转化的类型，特点，意义（题号82）

\end{pastquestion}

\begin{pastquestion}

\QuestionLabel{【往年题7｜14级毒理考题汇总.xls｜口试】}\\
生物转化的特点，包括什么反应，意义（题号72）；生物转化相关（题号82）；生物转化特点意义化学反应类型（题号9）

\end{pastquestion}

\PointSeparator

\section{Ⅰ相反应和Ⅱ相反应的概念及意义}\label{users__fpb__26_spring__pkusph_study__tmp__pdfs__prevention-medicine-past-exams__02-ux6bd2ux7406ux5b66ux5f80ux5e74ux9898ux8003ux70b9ux6574ux7406.md__ux76f8ux53cdux5e94ux548cux2171ux76f8ux53cdux5e94ux7684ux6982ux5ff5ux53caux610fux4e49}

\InfoLabel{考频} 5 次

\InfoLabel{对应小节} 第二节 生物转化反应 /《4 生物转化》第15页、第26页''Ⅰ相反应''、第41页''Ⅱ相反应''

\InfoLabel{匹配依据} ``Ⅰ相反应和Ⅱ相反应的概念，意义''``一相二相反应的概念及其意义''直接对应本小节。

\ContentBand{知识点原文摘取}

生物转化酶催化的反应一般分为两大类：Ⅰ相反应和Ⅱ相反应。Ⅰ相反应：氧化反应、还原反应和水解反应。Ⅱ相反应（结合作用
conjugation）：毒物原有的官能团或由Ⅰ相反应引入/暴露的官能团，与内源性辅助因子反应；除甲基化和乙酰化结合反应外，Ⅱ相反应显著增加毒物的水溶性、促进其排泄；包括葡萄糖醛酸化、硫酸化、乙酰化、甲基化，与谷胱甘肽结合以及与氨基酸结合。

\InfoLabel{知识来源} 《4 生物转化》第15页；《4 生物转化》第26页；《4 生物转化》第41页

\ContentBand{对应往年题}

\begin{pastquestion}

\QuestionLabel{【往年题1｜09级口试题目.docx｜口试】}\\
Ⅰ相反应和Ⅱ相反应的概念，意义（12）；Ⅰ相反应和Ⅱ相反应的概念及毒理学意义（40）

\end{pastquestion}

\begin{pastquestion}

\QuestionLabel{【往年题2｜19级毒理期末口试回忆.xlsx｜口试】}\\
I相反应II相反应（序号58）；I、II相反应的概念及意义（序号12）

\end{pastquestion}

\begin{pastquestion}

\QuestionLabel{【往年题3｜17级毒理试题汇总.xlsx｜口试】}\\
一相二相反应的概念及其意义（题号84）

\end{pastquestion}

\begin{pastquestion}

\QuestionLabel{【往年题4｜15级毒理学考试题目汇总.xls｜口试】}\\
代谢，12相反应，毒理学意义（题号72）

\end{pastquestion}

\begin{pastquestion}

\QuestionLabel{【往年题5｜14级毒理考题汇总.xls｜口试】}\\
二相反应包括？二相反应能代谢活化？（题号72追问）

\end{pastquestion}

\PointSeparator

\section{代谢酶的诱导（enzyme
induction）的概念和毒理学意义}\label{users__fpb__26_spring__pkusph_study__tmp__pdfs__prevention-medicine-past-exams__02-ux6bd2ux7406ux5b66ux5f80ux5e74ux9898ux8003ux70b9ux6574ux7406.md__ux4ee3ux8c22ux9176ux7684ux8bf1ux5bfcenzyme-inductionux7684ux6982ux5ff5ux548cux6bd2ux7406ux5b66ux610fux4e49}

\InfoLabel{考频} 6 次

\InfoLabel{对应小节} 第三节 生物转化的影响因素 /《4
生物转化》第58页''代谢酶的诱导、激活及抑制''、第59--60页''毒物代谢酶的诱导''

\InfoLabel{匹配依据} ``酶诱导（英文）作用的定义和意义''``enzyme
induction（酶诱导）的概念和毒理学意义''``转化酶诱导的定义和毒理学意义''直接对应本小节。

\ContentBand{知识点原文摘取}

当代谢酶被诱导和激活：对经代谢活化的毒物，则表现出毒性增强；对经代谢转化减毒的毒物，则表现为毒性降低。当代谢酶被抑制和阻遏，则得到相反的结果。酶的诱导（induction）：有些毒物可使某些代谢酶系合成增加并伴活力增强的现象。诱导剂（inducer）：具有诱导效应的毒物（如长期饮酒，吸烟）。双功能诱导剂：苯并吡
诱导Ⅰ相和Ⅱ相代谢酶。

\InfoLabel{知识来源} 《4 生物转化》第58页；《4 生物转化》第59页；《4 生物转化》第60页

\ContentBand{对应往年题}

\begin{pastquestion}

\QuestionLabel{【往年题1｜21级毒理期末口试回忆.xlsx｜口试】}\\
酶诱导（英文）作用的定义和意义（序号11）

\end{pastquestion}

\begin{pastquestion}

\QuestionLabel{【往年题2｜19级毒理期末口试回忆.xlsx｜口试】}\\
酶诱导的定义及意义（序号57）

\end{pastquestion}

\begin{pastquestion}

\QuestionLabel{【往年题3｜17级毒理试题汇总.xlsx｜口试】}\\
enzyme induction（酶诱导）的概念和毒理学意义（题号85）；酶诱导的概念和毒理学意义（题号12）

\end{pastquestion}

\begin{pastquestion}

\QuestionLabel{【往年题4｜15级毒理学考试题目汇总.xls｜口试】}\\
生物转化的影响因素？酶诱导的定义和毒理学意义（题号10）；生物转化的影响因素，酶诱导的概念和毒理学意义（题号73）

\end{pastquestion}

\begin{pastquestion}

\QuestionLabel{【往年题5｜14级毒理考题汇总.xls｜口试】}\\
生物转化的影响因素，转化酶诱导的定义和毒理学意义（题号39）；生物转化的影响因素，酶诱导的概念和意义（题号73）

\end{pastquestion}

\PointSeparator

\section{代谢酶（毒物代谢）的两重性和可饱和性的意义}\label{users__fpb__26_spring__pkusph_study__tmp__pdfs__prevention-medicine-past-exams__02-ux6bd2ux7406ux5b66ux5f80ux5e74ux9898ux8003ux70b9ux6574ux7406.md__ux4ee3ux8c22ux9176ux6bd2ux7269ux4ee3ux8c22ux7684ux4e24ux91cdux6027ux548cux53efux9971ux548cux6027ux7684ux610fux4e49}

\InfoLabel{考频} 4 次

\InfoLabel{对应小节} 第一节 生物转化和代谢酶 /《4
生物转化》第5--6页''生物转化的意义/结果''、第20页''代谢酶的种类及特点''、第62页''代谢酶的抑制（饱和/竞争）''

\InfoLabel{匹配依据}
``毒物代谢两重性和可饱和性的意义''``代谢的饱和性和两重性''``生物代谢的两重性和可饱和性，及其意义''直接对应------两重性即代谢解毒/代谢活化双重结果，可饱和性即酶催化能力有限、可被竞争抑制。

\ContentBand{知识点原文摘取}

生物转化的结果（毒理角度）：代谢解毒（metabolic
detoxication，外源化学物经生物转化使其毒性降低）；代谢活化（metabolic
activation，外源化学物经生物转化使其毒性增强，生成中间产物/终毒物）。非专一性酶（CYP450）特点：专一性低，活性有限，个体差异大，活性可受化学物影响。代谢酶的抑制可分：抑制物与酶的活性中心发生可逆或不可逆性结合；两种不同毒物在同一个酶的活性中心发生竞争性抑制；破坏酶；减少酶的合成；变构作用；缺乏辅因子。

\begin{correctionbox}

【勘误补充】\\
毒物代谢具有两重性和可饱和性。两重性是指外源化学物经生物转化后，既可能发生代谢解毒，使毒性降低、水溶性增加、易于排泄；也可能发生代谢活化，使毒性增强，生成活性中间产物或终毒物，引起组织损伤、致突变、致癌等毒效应。可饱和性是指毒物代谢依赖酶催化，而代谢酶活性和数量有限，当毒物剂量较大、暴露时间较长或多种毒物竞争同一代谢酶时，代谢能力可达到饱和，导致母体化合物或活性代谢产物蓄积，毒性增强。因此，毒物代谢的两重性和可饱和性是影响毒性大小、剂量-反应关系、个体差异、联合作用和毒性评价的重要因素。

\end{correctionbox}

\InfoLabel{知识来源} 《4 生物转化》第5页；《4 生物转化》第6页；《4 生物转化》第20页；《4 生物转化》第62页

\ContentBand{对应往年题}

\begin{pastquestion}

\QuestionLabel{【往年题1｜21级毒理期末口试回忆.xlsx｜口试】}\\
毒物代谢两重性和可饱和性的意义（序号16）；代谢的饱和性和两重性（序号62）

\end{pastquestion}

\begin{pastquestion}

\QuestionLabel{【往年题2｜19级毒理期末口试回忆.xlsx｜口试】}\\
生物代谢的两重性和可饱和性的意义（序号62）

\end{pastquestion}

\begin{pastquestion}

\QuestionLabel{【往年题3｜09级口试题目.docx｜口试】}\\
代谢酶的两重性与可饱和性（11）；代谢酶的两重性和饱和性的意义（42）

\end{pastquestion}

\begin{pastquestion}

\QuestionLabel{【往年题4｜17级毒理试题汇总.xlsx｜口试】}\\
生物代谢的两重性和可饱和性，及其意义（题号21）

\end{pastquestion}

\PointSeparator

\section{终毒物（ultimate
toxicant）的概念和分类}\label{users__fpb__26_spring__pkusph_study__tmp__pdfs__prevention-medicine-past-exams__02-ux6bd2ux7406ux5b66ux5f80ux5e74ux9898ux8003ux70b9ux6574ux7406.md__ux7ec8ux6bd2ux7269ultimate-toxicantux7684ux6982ux5ff5ux548cux5206ux7c7b}

\InfoLabel{考频} 4 次

\InfoLabel{对应小节} 第一节 生物转化和代谢酶 /《4
生物转化》第10--12页''终毒物''\,``外源化学物的活性形式和终毒物''（并见《5 毒作用机制》、《3》第36页）

\InfoLabel{匹配依据} ``终毒物的概念和分类''``终毒物的定义和类型''直接对应本小节。

\ContentBand{知识点原文摘取}

终毒物（ultimate
toxicant）：直接与内源性靶分子（如受体、酶、DNA、微丝蛋白、脂质）反应并造成机体损害的化学形态。终毒物：①外源化学物本身，如尼古丁，氰化氢；②代谢活化产物；③代谢过程激发内源毒物产生，如活性氧。毒效应的强度主要取决于终毒物在其作用位点的浓度及持续时间。外源化学物生成终毒物的几种类型：①亲电子剂；②自由基；③亲核剂；④氧化还原反应物。

\InfoLabel{知识来源} 《4 生物转化》第10页；《4 生物转化》第11页；《4 生物转化》第12页

\ContentBand{对应往年题}

\begin{pastquestion}

\QuestionLabel{【往年题1｜21级毒理期末口试回忆.xlsx｜口试】}\\
终毒物的概念和分类（序号85）

\end{pastquestion}

\begin{pastquestion}

\QuestionLabel{【往年题2｜19级毒理期末口试回忆.xlsx｜口试】}\\
终毒物的定义和类型（序号41）

\end{pastquestion}

\begin{pastquestion}

\QuestionLabel{【往年题3｜17级毒理试题汇总.xlsx｜口试】}\\
终毒物的概念与分类（题号56）；终毒物的概念和分类（题号35）

\end{pastquestion}

\PointSeparator

\section{生物转化的影响因素（机体、化学物、环境因素；代谢酶基因多态性）}\label{users__fpb__26_spring__pkusph_study__tmp__pdfs__prevention-medicine-past-exams__02-ux6bd2ux7406ux5b66ux5f80ux5e74ux9898ux8003ux70b9ux6574ux7406.md__ux751fux7269ux8f6cux5316ux7684ux5f71ux54cdux56e0ux7d20ux673aux4f53ux5316ux5b66ux7269ux73afux5883ux56e0ux7d20ux4ee3ux8c22ux9176ux57faux56e0ux591aux6001ux6027}

\InfoLabel{考频} 3 次

\InfoLabel{对应小节} 第三节 生物转化的影响因素 /《4 生物转化》第54--57页

\InfoLabel{匹配依据}
``生物转化的影响因素''作为独立设问反复出现（常与酶诱导合并），对应本小节遗传生理、环境因素及代谢酶多态性。

\ContentBand{知识点原文摘取}

生物转化是外源化学物毒作用的决定因素，受机体的遗传生理和环境两大类因素影响：遗传生理因素如物种、性别、遗传、年龄等，代谢酶种类、数量和活性差异（如代谢酶的多态性）；各种环境因素主要通过影响代谢酶和辅酶的合成及催化过程（如代谢酶的诱导和抑制）；其他因素如营养、疾病等。代谢酶基因多态性：外源化学物代谢酶的遗传差异是不同个体和种族间对外源化学物毒性和肿瘤易感性差异的原因之一；Ⅰ相和Ⅱ相代谢酶均存在多态性；代谢酶不同的表型------慢代谢型（PM）、中间代谢型（IM）、快代谢型（EM）、超快代谢型（UM）。

\InfoLabel{知识来源} 《4 生物转化》第54页；《4 生物转化》第55页；《4 生物转化》第56页

\ContentBand{对应往年题}

\begin{pastquestion}

\QuestionLabel{【往年题1｜15级毒理学考试题目汇总.xls｜口试】}\\
生物转化的影响因素？酶诱导的定义和毒理学意义（题号10）；生物转化的影响因素，酶诱导的概念和意义（题号73）

\end{pastquestion}

\begin{pastquestion}

\QuestionLabel{【往年题2｜14级毒理考题汇总.xls｜口试】}\\
生物转化的影响因素，转化酶诱导的定义和毒理学意义（题号39）

\end{pastquestion}

\begin{pastquestion}

\QuestionLabel{【往年题3｜09级口试题目.docx｜口试】}\\
姚老师提问，举一个代谢活化例子（12号备注）

\end{pastquestion}

\PointSeparator

\chapter{毒作用机制、急性毒性和局部毒性试验}\label{users__fpb__26_spring__pkusph_study__tmp__pdfs__prevention-medicine-past-exams__02-ux6bd2ux7406ux5b66ux5f80ux5e74ux9898ux8003ux70b9ux6574ux7406.md__ux7b2cux4e94ux7ae0-5-ux6bd2ux4f5cux7528ux673aux5236ux6025ux6027ux6bd2ux6027ux548cux5c40ux90e8ux6bd2ux6027ux8bd5ux9a8c}

\section{急性毒性试验的参数及其意义（毒性上限参数与下限参数）}\label{users__fpb__26_spring__pkusph_study__tmp__pdfs__prevention-medicine-past-exams__02-ux6bd2ux7406ux5b66ux5f80ux5e74ux9898ux8003ux70b9ux6574ux7406.md__ux6025ux6027ux6bd2ux6027ux8bd5ux9a8cux7684ux53c2ux6570ux53caux5176ux610fux4e49ux6bd2ux6027ux4e0aux9650ux53c2ux6570ux4e0eux4e0bux9650ux53c2ux6570}

\InfoLabel{考频} 11 次

\InfoLabel{对应小节} 第二节 急性毒性作用 /《5
毒作用机制、急性毒性和局部毒性试验》第18页''急性毒性实验及其设计要点''（并见《1
绪论》第22--23页致死/非致死参数；18级回忆文件给出完整答案）

\InfoLabel{匹配依据}
``急性毒性实验的参数有哪些及其意义''``列举急性毒性试验参数及其意义''``常见的急性毒性参数和意义''反复出现，对应急性毒性上限指标与非致死性指标。

\ContentBand{知识点原文摘取}

以死亡为终点的急性毒性实验：检测急性毒性上限指标 LD50/LC50
或近似LD50；非致死性急性毒性实验：可检测非致死性指标 LOAEL 和
NOAEL。（18级回忆答案）毒性参数可分为两类：一类为毒性上限参数（以死亡为终点）：LD100、LD50、LD01、LD0；一类为毒性下限参数（以非致死急性毒性为终点）：NOAEL、LOAEL。LD100绝对致死剂量、LD50半数致死剂量、LD01最小致死剂量、LD0/MTD最大耐受剂量；LOAEL观察到有害作用的最低水平、NOAEL未观察到有害作用水平；阈剂量为引起少数个体最轻微异常改变的最低剂量。

\InfoLabel{知识来源} 《5 毒作用机制、急性毒性和局部毒性试验》第18页；《1 绪论》第22页；《1 绪论》第23页

\ContentBand{对应往年题}

\begin{pastquestion}

\QuestionLabel{【往年题1｜07级毒理口试题回忆.doc｜口试】}\\
急性毒性实验的参数及意义（22题、58号题）；急性毒性参数，及其意义（27号题）

\end{pastquestion}

\begin{pastquestion}

\QuestionLabel{【往年题2｜09级口试题目.docx｜口试】}\\
列举急性毒性试验参数及其意义（22）；急性毒性参数及意义（27）；急性毒性试验描述指标有哪些及其定义：LD0，LD1，LD50，LD100（58）；急毒常用指标及意义（60）

\end{pastquestion}

\begin{pastquestion}

\QuestionLabel{【往年题3｜21级毒理期末口试回忆.xlsx｜口试】}\\
急性毒性实验参数（序号22、68，备注：追问LD50的意义）

\end{pastquestion}

\begin{pastquestion}

\QuestionLabel{【往年题4｜19级毒理期末口试回忆.xlsx｜口试】}\\
急性毒性实验的参数及各自的意义（序号68）；常见的急性毒性参数和意义（序号22）；急性毒性（序号74）

\end{pastquestion}

\begin{pastquestion}

\QuestionLabel{【往年题5｜17级毒理试题汇总.xlsx｜口试】}\\
急性毒性参数以及毒理学意义（题号22）；急性毒性（题号74）

\end{pastquestion}

\begin{pastquestion}

\QuestionLabel{【往年题6｜15级毒理学考试题目汇总.xls｜口试】}\\
急性毒性实验的参数有哪些？其意义？（题号22）；急毒参数及意义（题号50）；急性毒性试验的毒性参数（题号76）；急性毒性实验能够得到哪些毒理学参数？意义？（题号50）；毒性参数及意义（题号76）

\end{pastquestion}

\begin{pastquestion}

\QuestionLabel{【往年题7｜14级毒理考题汇总.xls｜口试】}\\
急性毒性参数有哪些 毒理学意义是（题号50）；急性毒性有哪些指标？其意义？（题号22）

\end{pastquestion}

\begin{pastquestion}

\QuestionLabel{【往年题8｜18级毒理回忆by连昕瑶（搬运自经验分享）.docx｜简答答案】}\\
急性毒性实验参数（上限参数 LD100/LD50/LD01/LD0 与下限参数 NOAEL/LOAEL 的定义与意义）

\end{pastquestion}

\PointSeparator

\section{经典急性致死性毒性试验的设计原则和实验步骤}\label{users__fpb__26_spring__pkusph_study__tmp__pdfs__prevention-medicine-past-exams__02-ux6bd2ux7406ux5b66ux5f80ux5e74ux9898ux8003ux70b9ux6574ux7406.md__ux7ecfux5178ux6025ux6027ux81f4ux6b7bux6027ux6bd2ux6027ux8bd5ux9a8cux7684ux8bbeux8ba1ux539fux5219ux548cux5b9eux9a8cux6b65ux9aa4}

\InfoLabel{考频} 9 次

\InfoLabel{对应小节} 三、急性毒性实验及其设计要点 /《5 毒作用机制、急性毒性和局部毒性试验》第18--28页

\InfoLabel{匹配依据}
``经典急性毒性实验的设计原则和实验步骤''``经典急性毒性试验设计原则，实验步骤''反复出现，对应本小节实验动物选择、试验设计（剂量设计）、观察、LD50计算。

\ContentBand{知识点原文摘取}

经典急性致死性毒性试验设计要点：实验动物的选择和处置（一般常用大、小鼠和狗，首选大鼠；体重大鼠180\textasciitilde240g、小鼠18\textasciitilde25g；雌雄各半；检疫5\textasciitilde7天；随机分组至少设3个剂量组）；试验设计（确定LD50计算方法；剂量设计至少3个剂量，高剂量全部/大部分死亡、中剂量50\%死亡、低剂量不引起/很少死亡；先进行预实验获得10\textasciitilde90\%致死剂量范围）；毒作用观察（染毒后一般观察14天，记录中毒体征、体重变化、非致死性指标可逆性）；LD50的计算（求回归方程
Y=a+bX，X为剂量对数值，Y为死亡率概率单位）。

\InfoLabel{知识来源} 《5 毒作用机制、急性毒性和局部毒性试验》第19页；《5
毒作用机制、急性毒性和局部毒性试验》第20页；《5 毒作用机制、急性毒性和局部毒性试验》第21页；《5
毒作用机制、急性毒性和局部毒性试验》第25页；《5 毒作用机制、急性毒性和局部毒性试验》第28页

\ContentBand{对应往年题}

\begin{pastquestion}

\QuestionLabel{【往年题1｜09级口试题目.docx｜口试】}\\
经典急性毒性试验原则及步骤（24）；经典急性毒性实验的设计原则和步骤（41）

\end{pastquestion}

\begin{pastquestion}

\QuestionLabel{【往年题2｜21级毒理期末口试回忆.xlsx｜口试】}\\
经典急性毒性实验的设计原则和实验步骤（序号16、59，备注：郝老师追问了急性毒性的替代实验有哪些）

\end{pastquestion}

\begin{pastquestion}

\QuestionLabel{【往年题3｜19级毒理期末口试回忆.xlsx｜口试】}\\
急性毒性试验的设计原则、实验步骤（序号59）

\end{pastquestion}

\begin{pastquestion}

\QuestionLabel{【往年题4｜17级毒理试题汇总.xlsx｜口试】}\\
经典急性毒性实验的设计原则并简述实验步骤（题号70）；急性毒性实验设计原则，实验步骤（题号23）

\end{pastquestion}

\begin{pastquestion}

\QuestionLabel{【往年题5｜15级毒理学考试题目汇总.xls｜口试】}\\
经典急性毒性实验设计（题号48）；经典急性毒性（题号20）；急性毒性试验设计及要点（题号46）

\end{pastquestion}

\begin{pastquestion}

\QuestionLabel{【往年题6｜14级毒理考题汇总.xls｜口试】}\\
急性毒性实验的设计及要点（题号46）；经典急性致死性毒性的试验设计要点（no 20）

\end{pastquestion}

\PointSeparator

\section{急性毒性替代试验（固定剂量法、急性毒性分级法、上-下法、限量试验）}\label{users__fpb__26_spring__pkusph_study__tmp__pdfs__prevention-medicine-past-exams__02-ux6bd2ux7406ux5b66ux5f80ux5e74ux9898ux8003ux70b9ux6574ux7406.md__ux6025ux6027ux6bd2ux6027ux66ffux4ee3ux8bd5ux9a8cux56faux5b9aux5242ux91cfux6cd5ux6025ux6027ux6bd2ux6027ux5206ux7ea7ux6cd5ux4e0a-ux4e0bux6cd5ux9650ux91cfux8bd5ux9a8c}

\InfoLabel{考频} 9 次

\InfoLabel{对应小节} 第三节 急性毒性替代试验 /《5 毒作用机制、急性毒性和局部毒性试验》第38--44页

\InfoLabel{匹配依据}
``急性毒性替代实验固定剂量法原理''``简述急毒替代实验固定剂量法的原理''``急性毒性的代替实验举1-2个，要点''反复出现，对应本小节四种替代方法。

\ContentBand{知识点原文摘取}

急性毒性替代试验：固定剂量法（fixed dose procedure，OECD
1992/2001，不以动物死亡作为观察终点，按5.0/50/300/2000
mg/kg观察存活及中毒表现判断毒性等级）；急性毒性分级法（acute toxic class method，OECD
1996/2001，以死亡为终点，少量动物，大鼠每阶段3只/同性别，剂量5、50、300、2000
mg/kg，观察14天，平均经2---4个阶段即可判断）；上、下移动法（up/down method，OECD
1998/2001，以死亡为终点，大鼠6\textasciitilde9只，由前一只反应决定下一只剂量）；限量试验（limit
test，单次染毒剂量≤10 g/kg体重，食品毒理学≤15 g/kg）。

【翻译】急性毒性替代试验包括固定剂量法、急性毒性分级法、上/下移动法和限量试验。固定剂量法不以动物死亡为观察终点，而根据存活和中毒表现判断毒性等级；急性毒性分级法和上/下移动法以死亡为终点但减少动物使用量；限量试验用于单次高剂量染毒条件下的急性毒性评价。

\InfoLabel{知识来源} 《5 毒作用机制、急性毒性和局部毒性试验》第39页；《5
毒作用机制、急性毒性和局部毒性试验》第40页；《5 毒作用机制、急性毒性和局部毒性试验》第42页；《5
毒作用机制、急性毒性和局部毒性试验》第43页；《5 毒作用机制、急性毒性和局部毒性试验》第44页

\ContentBand{对应往年题}

\begin{pastquestion}

\QuestionLabel{【往年题1｜21级毒理期末口试回忆.xlsx｜口试】}\\
急性毒理学实验固定剂量法的基本原理（序号18）；急性毒性试验固定剂量法（序号60）；急性毒性分级法简述（序号94）；急性分级毒性实验要点（序号25）

\end{pastquestion}

\begin{pastquestion}

\QuestionLabel{【往年题2｜19级毒理期末口试回忆.xlsx｜口试】}\\
急性毒性替代实验固定剂量法原理（序号18，备注：追问了其他的急毒替代实验的名称、经典急毒的局限性）；急毒固定剂量法（序号60）

\end{pastquestion}

\begin{pastquestion}

\QuestionLabel{【往年题3｜17级毒理试题汇总.xlsx｜口试】}\\
简述急性毒性分级实验（题号28）；简述急毒替代实验固定剂量法的原理和具体内容（题号73）；急毒替代实验中固定剂量法的原理（题号25，追问了还有什么替代实验）

\end{pastquestion}

\begin{pastquestion}

\QuestionLabel{【往年题4｜15级毒理学考试题目汇总.xls｜口试】}\\
急性毒性的代替实验举1-2个，要点（题号23）；急性毒性试验的替代实验和要点（题号23）；1-2种经典急性毒性实验的主要替代方法，说出实验要点（题号51，备注：上下移动法观察24小时再进行下一步判断）；急性致死性试验的替代试验设计（题号75）；急毒实验的替代法（题号48备注）

\end{pastquestion}

\begin{pastquestion}

\QuestionLabel{【往年题5｜14级毒理考题汇总.xls｜口试】}\\
一到两个急性致死性毒性试验的替代方法，说出其设计要点（无题号）；急性致死性试验的替代试验设计（题号75）

\end{pastquestion}

\PointSeparator

\section{LD50
的毒理学意义和局限性}\label{users__fpb__26_spring__pkusph_study__tmp__pdfs__prevention-medicine-past-exams__02-ux6bd2ux7406ux5b66ux5f80ux5e74ux9898ux8003ux70b9ux6574ux7406.md__ld50-ux7684ux6bd2ux7406ux5b66ux610fux4e49ux548cux5c40ux9650ux6027}

\InfoLabel{考频} 6 次

\InfoLabel{对应小节} 四、LD50的毒理学意义和局限性 /《5 毒作用机制、急性毒性和局部毒性试验》第30--31页（并见《1
绪论》第22页 LD50定义）

\InfoLabel{匹配依据}
``LD50的概念优缺点''``经典急性毒性实验（LD50）的优点和局限性''``ld50优缺点''反复出现，对应本小节。

\ContentBand{知识点原文摘取}

（一）意义：比较毒物急性毒性的大小；计算药物的治疗指数（LD50/ED50），了解药效学剂量和毒性剂量的距离；为后续重复给药毒理试验剂量选择提供参考；比较不同途径的LD50值，获得生物利用度的信息；试验结果可用来推测人类的致死剂量以及中毒后的体征。（二）局限性：LD50值给予的有效信息少，作用有限；LD50的波动性很大，无稳定的数值；物种差异对LD50的影响大，对于人急性中毒的诊疗意义很小；对于开发新药，无需得到精确的LD50值；消耗的动物数量大，经典法一次需30\textasciitilde70只动物。

\InfoLabel{知识来源} 《5 毒作用机制、急性毒性和局部毒性试验》第30页；《5
毒作用机制、急性毒性和局部毒性试验》第31页；《1 绪论》第22页

\ContentBand{对应往年题}

\begin{pastquestion}

\QuestionLabel{【往年题1｜07级毒理口试题回忆.doc｜口试】}\\
LD50的概念和优缺点（1号、3号题）

\end{pastquestion}

\begin{pastquestion}

\QuestionLabel{【往年题2｜09级口试题目.docx｜口试】}\\
ld50的概念优缺点（13）

\end{pastquestion}

\begin{pastquestion}

\QuestionLabel{【往年题3｜17级毒理试题汇总.xlsx｜口试】}\\
ld50概念意义局限（题号23）

\end{pastquestion}

\begin{pastquestion}

\QuestionLabel{【往年题4｜15级毒理学考试题目汇总.xls｜口试】}\\
ld50优缺点（题号48）；LD50的优点和局限性（题号21）；LD50实验的优点和局限性（题号77）；急毒实验ld50的优缺点（题号49）

\end{pastquestion}

\begin{pastquestion}

\QuestionLabel{【往年题5｜14级毒理考题汇总.xls｜口试】}\\
经典急性毒性实验(LD50）的优点和局限性（题号21）；急性毒性试验（LD50）的优点和局限性（题号49，备注：追问了经典急性毒性的不足）

\end{pastquestion}

\PointSeparator

\section{局部毒性试验的概念及评价方法（皮肤刺激/腐蚀、眼刺激/腐蚀等）}\label{users__fpb__26_spring__pkusph_study__tmp__pdfs__prevention-medicine-past-exams__02-ux6bd2ux7406ux5b66ux5f80ux5e74ux9898ux8003ux70b9ux6574ux7406.md__ux5c40ux90e8ux6bd2ux6027ux8bd5ux9a8cux7684ux6982ux5ff5ux53caux8bc4ux4ef7ux65b9ux6cd5ux76aeux80a4ux523aux6fc0ux8150ux8680ux773cux523aux6fc0ux8150ux8680ux7b49}

\InfoLabel{考频} 6 次

\InfoLabel{对应小节} 第四节 局部毒性作用 /《5 毒作用机制、急性毒性和局部毒性试验》第45--53页

\InfoLabel{匹配依据}
``简述局部毒性试验的概念并例举一种评价方法''``描述局部毒性实验的概念并列举1种常见评价方法''直接对应本小节。

\ContentBand{知识点原文摘取}

局部毒性为外源化学物对机体接触部位的毒效应，不包括经接触部位的吸收而引起的系统毒性。局部毒性又称短程毒性试验。一、皮肤原发性刺激试验（包括单次和多次皮肤刺激试验、完整皮肤和破损皮肤刺激试验等，观察终点为皮肤刺激和皮肤腐蚀；Draize试验：白色家兔或豚鼠，脱毛、染毒、观察红斑水肿评分）。二、眼原发性刺激试验（Draize试验：白色家兔，0.1ml受试物滴入一眼，观察结膜、角膜和虹膜反应）。皮肤刺激/腐蚀试验替代方法：皮肤组织培养、重组人体皮肤模型（EpiSkin™、EpiDerm™）；眼刺激替代方法：离体牛眼角膜浑浊度和通透性试验（BCOP）、离体鸡眼试验（ICE）、重组人角膜上皮细胞模型试验（RhCE）。

\InfoLabel{知识来源} 《5 毒作用机制、急性毒性和局部毒性试验》第45页；《5
毒作用机制、急性毒性和局部毒性试验》第46页；《5 毒作用机制、急性毒性和局部毒性试验》第47页；《5
毒作用机制、急性毒性和局部毒性试验》第49页；《5 毒作用机制、急性毒性和局部毒性试验》第50页；《5
毒作用机制、急性毒性和局部毒性试验》第52页

\ContentBand{对应往年题}

\begin{pastquestion}

\QuestionLabel{【往年题1｜21级毒理期末口试回忆.xlsx｜口试】}\\
简述局部毒性试验的概念并例举一种评价方法（序号26）；局部毒性试验概念并举例（序号72）；局部毒性的概念和评价方法（序号95）

\end{pastquestion}

\begin{pastquestion}

\QuestionLabel{【往年题2｜19级毒理期末口试回忆.xlsx｜口试】}\\
局部毒性试验并举一个方法（序号26）；描述局部毒性实验的概念并列举1种常见评价方法（序号95）

\end{pastquestion}

\begin{pastquestion}

\QuestionLabel{【往年题3｜17级毒理试题汇总.xlsx｜口试】}\\
局部毒性实验的概念，并简述1、2种评价方法（题号70）；局部毒性试验的概念并举两个例子（题号16）

\end{pastquestion}

\PointSeparator

\section{毒作用模式（MOA）、分子起始事件（MIE）、有害结局通路（AOP）、毒性通路（TP）等新概念}\label{users__fpb__26_spring__pkusph_study__tmp__pdfs__prevention-medicine-past-exams__02-ux6bd2ux7406ux5b66ux5f80ux5e74ux9898ux8003ux70b9ux6574ux7406.md__ux6bd2ux4f5cux7528ux6a21ux5f0fmoaux5206ux5b50ux8d77ux59cbux4e8bux4ef6mieux6709ux5bb3ux7ed3ux5c40ux901aux8defaopux6bd2ux6027ux901aux8deftpux7b49ux65b0ux6982ux5ff5}

\InfoLabel{考频} 5 次

\InfoLabel{对应小节} 第一节 毒作用机制 一、新概念 /《5 毒作用机制、急性毒性和局部毒性试验》第2--3页

\InfoLabel{匹配依据} ``AOP,MIE.MOA,TP名解''``MOA，AOP，MIE，TP相关概念云云''``tp
moa等四个概念''直接对应本小节新概念。

\ContentBand{知识点原文摘取}

毒作用模式（Mode of
Action，MOA）：指从化学物与生物分子相互作用开始的、证据权重支持的、可能导致毒性有关终点的一组事件或过程。分子起始事件（Molecular
Initiating
Event，MIE）：在机体内化学物-生物分子相互作用，启动通路扰动的初始点，沿被扰动通路引起进一步的关键事件，并与有害结局相关联。有害结局通路（Adverse
Outcome
Pathway，AOP）：是在个人或群体的水平，关于分子起始事件和有害结局之间相互联系的已有的知识汇总。毒性通路（Toxicity
Pathway）：在受到一定程度的干扰时会导致有害效应的细胞应答通路。

【翻译】MOA 是毒作用模式，MIE 是分子起始事件，AOP 是有害结局通路，TP
是毒性通路。它们用于把化学物与生物分子的初始相互作用、后续关键事件、细胞应答通路和最终有害结局联系起来。

\InfoLabel{知识来源} 《5 毒作用机制、急性毒性和局部毒性试验》第2页；《5
毒作用机制、急性毒性和局部毒性试验》第3页

\ContentBand{对应往年题}

\begin{pastquestion}

\QuestionLabel{【往年题1｜19级毒理期末口试回忆.xlsx｜口试】}\\
AOP,MIE.MOA,TP名解（序号84）；aop mie之类的（序号40）；tp moa等四个概念（序号84）

\end{pastquestion}

\begin{pastquestion}

\QuestionLabel{【往年题2｜17级毒理试题汇总.xlsx｜口试】}\\
moa,aop,mie,tp相关概念云云（题号57）；MOA，AOP，MIE，TP（毒作用机制概述）（题号34，备注：郝老师说这是前沿词汇）

\end{pastquestion}

\PointSeparator

\section{细胞应激与细胞损伤；机体/细胞遭受损伤时早期（最初）的防御机制及毒理学意义}\label{users__fpb__26_spring__pkusph_study__tmp__pdfs__prevention-medicine-past-exams__02-ux6bd2ux7406ux5b66ux5f80ux5e74ux9898ux8003ux70b9ux6574ux7406.md__ux7ec6ux80deux5e94ux6fc0ux4e0eux7ec6ux80deux635fux4f24ux673aux4f53ux7ec6ux80deux906dux53d7ux635fux4f24ux65f6ux65e9ux671fux6700ux521dux7684ux9632ux5fa1ux673aux5236ux53caux6bd2ux7406ux5b66ux610fux4e49}

\InfoLabel{考频} 7 次

\InfoLabel{对应小节} 第一节 毒作用机制 一、细胞应激与细胞损伤 /《5
毒作用机制、急性毒性和局部毒性试验》第9--10页

\InfoLabel{匹配依据}
``机体和细胞遭受损伤时，最初的防御机制是什么？毒理学意义''``细胞和机体应对有害因素的早期防御反应有哪些？毒理学意义''``应激反应的内容，毒理学意义''直接对应本小节适应性应激反应。

\ContentBand{知识点原文摘取}

外源化学物作用于细胞后，细胞具有一定的抗损害能力，即首先发生适应性应激反应，当外源化学物的作用强度超过细胞的抗损害能力后，就会引起细胞功能障碍（细胞损伤），最后还可导致死亡。细胞应激（Cell
Stress）：指细胞处于不利环境和遇到有害刺激时所产生的防御和适应性反应。包括热应激、氧化应激、缺氧应激、遗传损伤应激、内质网应激、线粒体应激。

\begin{correctionbox}

【勘误补充】\\
细胞应激是机体和细胞遭受有害因素刺激或损伤时最早发生的防御和适应性反应，是细胞对不利环境的初步保护机制。其主要形式包括热应激、氧化应激、缺氧应激、遗传损伤应激、内质网应激和线粒体应激等。

\end{correctionbox}

其毒理学意义在于：当外源化学物作用于细胞后，细胞首先通过适应性应激反应增强抗损伤能力，维持细胞稳态，减轻或延缓毒性损伤；但当外源化学物作用强度超过细胞的抗损害能力时，应激反应不足以维持稳态，细胞可出现功能障碍、结构损伤，甚至发生细胞死亡。因此，细胞应激既是早期毒作用的重要表现，也是判断毒性损伤发生、发展和机体防御能力的重要依据。

\InfoLabel{知识来源} 《5 毒作用机制、急性毒性和局部毒性试验》第9页；《5
毒作用机制、急性毒性和局部毒性试验》第10页

\ContentBand{对应往年题}

\begin{pastquestion}

\QuestionLabel{【往年题1｜15级毒理学考试题目汇总.xls｜口试】}\\
细胞应激（题号47）；机体和细胞遭受损伤时，最初的防御机制是什么？毒理学意义（题号47）；应激反应的内容，毒理学意义（题号19）；机体受到外界有害因素刺激时首先的防御以及毒理学意义（题号78）；应对有害刺激的首要防御机制（题号78）

\end{pastquestion}

\begin{pastquestion}

\QuestionLabel{【往年题2｜14级毒理考题汇总.xls｜口试】}\\
细胞和机体对于有害作用产生的初步防御机制是什么，毒理学意义（题号47）；细胞或机体遭受损伤时，早期的防御机制是什么？毒理学意义（题号19）；细胞和机体应对有害因素的早期防御反应有哪些？毒理学意义（题号81）；机体受到外界有害因素刺激时首先的防御以及毒理学意义（题号78）

\end{pastquestion}

\PointSeparator

\section{化学物引起毒作用的几个阶段（外源化学物对机体毒效应主要作用阶段）}\label{users__fpb__26_spring__pkusph_study__tmp__pdfs__prevention-medicine-past-exams__02-ux6bd2ux7406ux5b66ux5f80ux5e74ux9898ux8003ux70b9ux6574ux7406.md__ux5316ux5b66ux7269ux5f15ux8d77ux6bd2ux4f5cux7528ux7684ux51e0ux4e2aux9636ux6bb5ux5916ux6e90ux5316ux5b66ux7269ux5bf9ux673aux4f53ux6bd2ux6548ux5e94ux4e3bux8981ux4f5cux7528ux9636ux6bb5}

\InfoLabel{考频} 4 次

\InfoLabel{对应小节} 第一节 毒作用机制 /《5
毒作用机制、急性毒性和局部毒性试验》第4页''化学物引起毒作用的4个阶段、3种途径''

\InfoLabel{匹配依据}
``简述外源化学物对机体毒性作用的几个过程''``外源化学物对机体毒效应机制主要作用于哪几个阶段''``外源性化合物毒性作用阶段''直接对应本小节。

\ContentBand{知识点原文摘取}

化学物引起毒作用的4个阶段、3种途径：暴露 → 处置 → 与靶分子相互作用/生物环境的改变 → 细胞功能障碍 → 修复失调 →
毒作用。（1）化学物不与靶分子发生作用而直接作用于接触的部位；（2）毒物与靶分子发生作用，造成细胞功能障碍或损伤；（3）毒物造成细胞损伤后，机体出现异常修复（包括分子水平、细胞水平和组织水平的修复）。毒理学上常见的靶分子：膜脂质、核酸（特别是DNA）、蛋白质。

\InfoLabel{知识来源} 《5 毒作用机制、急性毒性和局部毒性试验》第4页；《5
毒作用机制、急性毒性和局部毒性试验》第5页

\ContentBand{对应往年题}

\begin{pastquestion}

\QuestionLabel{【往年题1｜14级毒理考题汇总.xls｜口试】}\\
简述外源化学物对机体毒性作用的几个过程（题号18）；外源性化合物毒性作用阶段（题号46）；外源化学物对机体毒效应机制主要作用于哪几个阶段（题号79）

\end{pastquestion}

\begin{pastquestion}

\QuestionLabel{【往年题2｜15级毒理学考试题目汇总.xls｜口试】}\\
外源化合物在体内发挥毒性的阶段（题号46）

\end{pastquestion}

\PointSeparator

\section{急性毒性分级（GHS
与我国食品毒理急性毒性分级）（久远题库补充）}\label{users__fpb__26_spring__pkusph_study__tmp__pdfs__prevention-medicine-past-exams__02-ux6bd2ux7406ux5b66ux5f80ux5e74ux9898ux8003ux70b9ux6574ux7406.md__ux6025ux6027ux6bd2ux6027ux5206ux7ea7ghs-ux4e0eux6211ux56fdux98dfux54c1ux6bd2ux7406ux6025ux6027ux6bd2ux6027ux5206ux7ea7ux4e45ux8fdcux9898ux5e93ux8865ux5145}

\InfoLabel{考频} 2 次

\InfoLabel{对应小节} 急性毒性分级 /《5 毒作用机制、急性毒性和局部毒性试验》第33--37页

\InfoLabel{匹配依据} 久远题库与卫生毒理学习题中多次出现急性毒性分级、LD50/LC50 与
GHS/食品毒理分级相关内容；现有整理有''急性毒性参数''和''GHS健康危害分类''，但未在急性毒性章节单列分级。

\ContentBand{知识点原文摘取}

GHS 的化学品急性毒性分级基于经口、皮肤途径和吸入途径的 LD50/LC50 将化学品急性毒性分为第 1 类至第 5
类。急性毒性危害类别和急性毒性估计（ATE）值：经口 LD50 第 1 类 ≤5 mg/kg，第 2 类 ≤50 mg/kg，第 3 类 ≤300
mg/kg，第 4 类 ≤2000 mg/kg，第 5 类 ≤5000 mg/kg；皮肤 LD50 第 1 类 ≤50 mg/kg，第 2 类 ≤200 mg/kg，第 3 类
≤1000 mg/kg，第 4 类 ≤2000 mg/kg；气体、蒸气、粉尘和烟雾吸入临界值以 4 小时试验接触为基础。

我国食品毒理急性毒性分级法（1994）：按小鼠一次经口 LD50 分为 6 级，极毒 \textless1 mg/kg，剧毒 1--50
mg/kg，中等毒 51--500 mg/kg，低毒 501--5000 mg/kg，实际无毒 5001--15000 mg/kg，无毒 \textgreater15000 mg/kg。

\InfoLabel{知识来源} 《5 毒作用机制、急性毒性和局部毒性试验》第33页；《5
毒作用机制、急性毒性和局部毒性试验》第34页；《5 毒作用机制、急性毒性和局部毒性试验》第37页

\ContentBand{对应往年题}

\begin{pastquestion}

\QuestionLabel{【往年题1｜相对久远的总结/卫生毒理学习题.doc｜单项选择/久远题库补充】}\\
急性毒性分级、LD50/LC50 与急性毒性估计值相关选择题。

\end{pastquestion}

\begin{pastquestion}

\QuestionLabel{【往年题2｜相对久远的总结/毒理题库-10年更新.doc｜久远题库补充】}\\
急性毒性参数分类；利用 LD50 比较不同外源性化学物急性毒性。

\end{pastquestion}

\PointSeparator

\chapter{亚慢性和慢性毒性试验}\label{users__fpb__26_spring__pkusph_study__tmp__pdfs__prevention-medicine-past-exams__02-ux6bd2ux7406ux5b66ux5f80ux5e74ux9898ux8003ux70b9ux6574ux7406.md__ux7b2cux516dux7ae0-6-ux4e9aux6162ux6027ux548cux6162ux6027ux6bd2ux6027ux8bd5ux9a8c}

\section{亚慢性毒性与慢性毒性的概念、目的及主要观察指标}\label{users__fpb__26_spring__pkusph_study__tmp__pdfs__prevention-medicine-past-exams__02-ux6bd2ux7406ux5b66ux5f80ux5e74ux9898ux8003ux70b9ux6574ux7406.md__ux4e9aux6162ux6027ux6bd2ux6027ux4e0eux6162ux6027ux6bd2ux6027ux7684ux6982ux5ff5ux76eeux7684ux53caux4e3bux8981ux89c2ux5bdfux6307ux6807}

\InfoLabel{考频} 10 次

\InfoLabel{对应小节} （一）Concept；The purpose；六、观察指标 /《6 亚慢性和慢性毒性试验》第8--9页、第20页

\InfoLabel{匹配依据} ``亚慢性和慢性毒理学试验的定义、目的和参考指标''``subchronic toxicity和chronic
toxicity的概念和研究目的''``慢性亚慢性毒性定义与研究目的''反复出现，对应本小节。

\ContentBand{知识点原文摘取}

Subchronic toxicity：The toxicity induced by xenobiotics under repeated exposure（1-3 months）in experimental
animals. Chronic toxicity：The toxicity induced by xenobiotics under repeated exposure（long time or life
long）in experimental animals. The purpose of sub-chronic and chronic toxicity experiment：1. New toxic
effect；2. Target organ；3. Dose-response；4. Reversibility；5. Reference to chronic；6.
Extrapolate。观察指标：1.一般性指标（动物外观、生长发育、行为功能、体征；动物体重；食物利用率）；2.实验室检查（血液学、血液生化、尿液）；3.眼科学检查；4.心脏血管检查；5.病理学检查（大体解剖脏器系数、组织病理学）；6.特异性指标（反映损伤机制、效应生物学标志）。

【翻译】亚慢性毒性是实验动物在 1-3
个月重复暴露于外源化学物后产生的毒性。慢性毒性是实验动物长期或终生重复暴露于外源化学物后产生的毒性。亚慢性和慢性毒性试验的目的包括发现新的毒性效应、确定靶器官、阐明剂量-反应关系、判断可逆性、为慢性试验提供参考以及进行外推。

\InfoLabel{知识来源} 《6 亚慢性和慢性毒性试验》第8页；《6 亚慢性和慢性毒性试验》第9页；《6
亚慢性和慢性毒性试验》第20页

\ContentBand{对应往年题}

\begin{pastquestion}

\QuestionLabel{【往年题1｜07级毒理口试题回忆.doc｜口试】}\\
亚慢性毒性和慢性毒性的概念、目的和检测指标（10号题）；亚慢性和慢性毒性的概念、目的、观测指标（44号题，观测指标老师要求说得很细）；亚慢性、慢性毒性的定义、目的、观察指标（45号、46号题）

\end{pastquestion}

\begin{pastquestion}

\QuestionLabel{【往年题2｜09级口试题目.docx｜口试】}\\
亚慢性和慢性毒理学试验的定义、目的和参考指标都有哪些（10、44）；简述亚慢性和慢性毒理试验的概念、目的、主要观察指标（11、19）

\end{pastquestion}

\begin{pastquestion}

\QuestionLabel{【往年题3｜21级毒理期末口试回忆.xlsx｜口试】}\\
subchronic toxicity和chronic toxicity的概念和研究目的（序号43、87）

\end{pastquestion}

\begin{pastquestion}

\QuestionLabel{【往年题4｜19级毒理期末口试回忆.xlsx｜口试】}\\
慢性亚慢性毒性定义与研究目的（序号99）

\end{pastquestion}

\begin{pastquestion}

\QuestionLabel{【往年题5｜17级毒理试题汇总.xlsx｜口试】}\\
subchronic toxicity \& chronic
toxicity的概念和目的（题号43）；亚慢性和慢性毒性的概念和目的（亚慢性和慢性是英文）（题号53）；亚慢性毒性和慢性毒性概念，目的和观察指标（题号31）

\end{pastquestion}

\PointSeparator

\section{慢性毒性试验的动物选择、剂量设计、结果评价原则}\label{users__fpb__26_spring__pkusph_study__tmp__pdfs__prevention-medicine-past-exams__02-ux6bd2ux7406ux5b66ux5f80ux5e74ux9898ux8003ux70b9ux6574ux7406.md__ux6162ux6027ux6bd2ux6027ux8bd5ux9a8cux7684ux52a8ux7269ux9009ux62e9ux5242ux91cfux8bbeux8ba1ux7ed3ux679cux8bc4ux4ef7ux539fux5219}

\InfoLabel{考频} 8 次

\InfoLabel{对应小节} 实验设计（动物种属、剂量设计、结果分析）/《6 亚慢性和慢性毒性试验》第14--20页、第33--36页

\InfoLabel{匹配依据}
``慢性毒性试验动物选择、剂量分组、结果评价原则''``慢性毒性实验动物选择，剂量设计及结果评价的原则''反复出现，对应本小节。

\ContentBand{知识点原文摘取}

动物种属选择：啮齿类常用大鼠，非啮齿类狗。年龄：大鼠50-70g（慢性）。动物数量：大鼠至少40只（慢性），狗至少8只（慢性）。染毒期限：慢性毒性研究------工业毒理6个月，环境、食品毒理1年以上，OECD要求至少一年。剂量设计：急性毒性阈剂量、LD50的1/5\textasciitilde1/20，组距3-10倍；高剂量出现较明显毒性、无死亡或个别死亡，中剂量出现有害作用（LOAEL），低剂量未观察到有害作用（NOAEL），并设对照（溶剂对照）。结果分析：数据特征、统计学分析、综合分析，需要同时考虑生物学意义。

\InfoLabel{知识来源} 《6 亚慢性和慢性毒性试验》第14页；《6 亚慢性和慢性毒性试验》第15页；《6
亚慢性和慢性毒性试验》第16页；《6 亚慢性和慢性毒性试验》第18页；《6 亚慢性和慢性毒性试验》第19页

\ContentBand{对应往年题}

\begin{pastquestion}

\QuestionLabel{【往年题1｜07级毒理口试题回忆.doc｜口试】}\\
慢性毒性实验的选择动物原则，剂量分组原则和结果评价原则（3号题）；慢性毒性实验的动物选择、剂量设计、结果评价（5号题）；慢性毒性试验动物选择、剂量设置、结果评价原则（33号题）

\end{pastquestion}

\begin{pastquestion}

\QuestionLabel{【往年题2｜09级口试题目.docx｜口试】}\\
慢性毒性试验动物选择、剂量设计、评判原则（3）；慢性毒性试验动物选择、剂量分组、结果评价的原则（5）；慢性毒性试验的剂量设计、动物选择和结果判定（33）

\end{pastquestion}

\begin{pastquestion}

\QuestionLabel{【往年题3｜21级毒理期末口试回忆.xlsx｜口试】}\\
慢性毒性实验的动物选择，剂量设计及结果评价的原则（序号14）；慢性毒性实验设计（序号62）

\end{pastquestion}

\begin{pastquestion}

\QuestionLabel{【往年题4｜19级毒理期末口试回忆.xlsx｜口试】}\\
慢性毒性实验动物选择，剂量设计，结果评价的原则（序号62）

\end{pastquestion}

\begin{pastquestion}

\QuestionLabel{【往年题5｜17级毒理试题汇总.xlsx｜口试】}\\
慢性毒性实验动物选择、剂量分组、结果评价原则（题号78）；慢性毒性实验动物选择，剂量设计，结果评价原则（题号31）

\end{pastquestion}

\PointSeparator

\section{重复暴露（重复毒性）试验的类型、目的及必要性}\label{users__fpb__26_spring__pkusph_study__tmp__pdfs__prevention-medicine-past-exams__02-ux6bd2ux7406ux5b66ux5f80ux5e74ux9898ux8003ux70b9ux6574ux7406.md__ux91cdux590dux66b4ux9732ux91cdux590dux6bd2ux6027ux8bd5ux9a8cux7684ux7c7bux578bux76eeux7684ux53caux5fc5ux8981ux6027}

\InfoLabel{考频} 7 次

\InfoLabel{对应小节} Single/Repeated exposure；The tendency for toxicity experiment /《6
亚慢性和慢性毒性试验》第5--12页

\InfoLabel{匹配依据}
``重复暴露毒性试验的类型、目的、必要性''``重复毒性分类目的必要性''反复出现，对应本小节单次/重复暴露、急性毒性不能预测重复暴露毒性、试验趋势。

\ContentBand{知识点原文摘取}

Acute toxicity can not predict the repeated exposure toxicity。Single exposure / Repeated
exposure：Accumulation------materials accumulation、effect accumulation。The tendency for toxicity
experiment：01 Acute toxicity；02 Sub-chronic toxicity；03 Chronic toxicity will be studied only when
necessary；04 Chronic toxicity experiment tends to be carried out combined with carcinogenesis
experiment。（112 chemicals，Rats，97.5\% of chemicals could present their toxic effects within 90 days。）

【翻译】急性毒性不能预测重复暴露毒性。单次暴露与重复暴露的差异之一在于蓄积，包括物质蓄积和效应蓄积。毒性试验的发展趋势是先进行急性毒性试验，再进行亚慢性毒性试验；慢性毒性试验仅在必要时开展，并且倾向于与致癌试验合并进行。在大鼠
112 种化学物研究中，97.5\% 的化学物可在 90 天内表现出其毒性效应。

\InfoLabel{知识来源} 《6 亚慢性和慢性毒性试验》第5页；《6 亚慢性和慢性毒性试验》第6页；《6
亚慢性和慢性毒性试验》第11页；《6 亚慢性和慢性毒性试验》第12页

\ContentBand{对应往年题}

\begin{pastquestion}

\QuestionLabel{【往年题1｜15级毒理学考试题目汇总.xls｜口试】}\\
重复暴露性实验分类、目的、必要性（题号24）；重复暴露目的，类型，意义（题号20）；重复毒性试验的类型，意义，必要性（题号82）；重复暴露研究设计有哪些，目的是什么，必要性（题号24）

\end{pastquestion}

\begin{pastquestion}

\QuestionLabel{【往年题2｜14级毒理考题汇总.xls｜口试】}\\
重复毒性分类目的必要性（题号52）；重复暴露毒性试验的类型、目的、必要性（no
20）；重复暴露研究的类型，目的和意义（题号82）；重复毒性实验的类型 目的
必要性（题号74）；重复暴露研究设计有哪些，目的是什么，必要性（题号24）

\end{pastquestion}

\begin{pastquestion}

\QuestionLabel{【往年题3｜07级毒理口试题回忆.doc｜口试】}\\
亚慢性毒性和慢性毒性的概念、目的、检测指标（10号题，含重复暴露意义）

\end{pastquestion}

\PointSeparator

\section{亚慢性和慢性毒性试验设置不同剂量组的目的；亚慢性毒性实验设计}\label{users__fpb__26_spring__pkusph_study__tmp__pdfs__prevention-medicine-past-exams__02-ux6bd2ux7406ux5b66ux5f80ux5e74ux9898ux8003ux70b9ux6574ux7406.md__ux4e9aux6162ux6027ux548cux6162ux6027ux6bd2ux6027ux8bd5ux9a8cux8bbeux7f6eux4e0dux540cux5242ux91cfux7ec4ux7684ux76eeux7684ux4e9aux6162ux6027ux6bd2ux6027ux5b9eux9a8cux8bbeux8ba1}

\InfoLabel{考频} 6 次

\InfoLabel{对应小节} 四、实验分组；五、剂量设计 /《6 亚慢性和慢性毒性试验》第13--26页（第18页实验分组）

\InfoLabel{匹配依据}
``亚慢性和慢性毒性试验设置不同剂量组的目的''``亚慢性毒性实验设计''``亚慢性试验全要素（实验设计）''反复出现，对应本小节。

\ContentBand{知识点原文摘取}

四、实验分组：高剂量------较为明显的毒性，无死亡或个别动物死亡；中剂量------出现有害作用（LOAEL）；低剂量------未观察到有害作用（NOAEL）；对照------溶剂对照。整体动物实验设计的要素：一个伦理审查；二个暴露因素（途径、期限）；三个剂量组（低、中、高）；四个动物相关要素（物种、性别、年龄、数量）；五类观察指标（一般性指标、实验室检查、仪器检查、病理学检查、特异性指标）。

\InfoLabel{知识来源} 《6 亚慢性和慢性毒性试验》第18页；《6 亚慢性和慢性毒性试验》第19页；《6
亚慢性和慢性毒性试验》第26页

\ContentBand{对应往年题}

\begin{pastquestion}

\QuestionLabel{【往年题1｜21级毒理期末口试回忆.xlsx｜口试】}\\
亚慢性毒性实验设计（序号94，备注：亚慢性毒性实验设计追问了探究可逆性时的特殊设计）；亚慢性和慢性设置剂量组的目的（序号95）；亚慢性毒性实验设计追问探究可逆性时的特殊设计（序号94备注）

\end{pastquestion}

\begin{pastquestion}

\QuestionLabel{【往年题2｜19级毒理期末口试回忆.xlsx｜口试】}\\
亚慢性和慢性毒性实验中设置不同剂量组的目的（序号15）；亚慢性毒性实验设计（序号5）；亚慢性和慢性实验设置不同剂量组的目的（序号95）

\end{pastquestion}

\begin{pastquestion}

\QuestionLabel{【往年题3｜17级毒理试题汇总.xlsx｜口试】}\\
亚慢性实验设计（题号48）；亚慢性毒性试验的实验设计和实施方案（题号96）；亚慢性毒性实验设计和实验方案（题号40、88）；慢性亚慢性实验设置剂量组的目的（题号26）

\end{pastquestion}

\begin{pastquestion}

\QuestionLabel{【往年题4｜15级毒理学考试题目汇总.xls｜口试】}\\
亚慢性毒性试验全部要点（题号53）；亚慢性毒性实验全因素（试验设计）（题号25）；亚慢性试验全要素（实验设计）（题号53）

\end{pastquestion}

\PointSeparator

\section{不同物质（化学物）在重复暴露/重复毒性试验中需要考虑的区别与注意事项}\label{users__fpb__26_spring__pkusph_study__tmp__pdfs__prevention-medicine-past-exams__02-ux6bd2ux7406ux5b66ux5f80ux5e74ux9898ux8003ux70b9ux6574ux7406.md__ux4e0dux540cux7269ux8d28ux5316ux5b66ux7269ux5728ux91cdux590dux66b4ux9732ux91cdux590dux6bd2ux6027ux8bd5ux9a8cux4e2dux9700ux8981ux8003ux8651ux7684ux533aux522bux4e0eux6ce8ux610fux4e8bux9879}

\InfoLabel{考频} 4 次

\InfoLabel{对应小节} 受试物A/B/C；接触化学物的方式；剂量设计 /《6
亚慢性和慢性毒性试验》第16--17页、第37--39页（并见《13 安全性评价》第7--8页）

\InfoLabel{匹配依据}
``不同化学物质在重复暴露实验研究中区别''``不同物质重复暴露实验设计步骤中的区别''``不同物质重复毒性试验的注意事项''对应本小节按化学物/药品/保健品差异设计的内容（与安全性评价分类原则相关）。

\ContentBand{知识点原文摘取}

三、接触化学物的方式：1.模拟人的实际接触方式；2.亚慢性与慢性毒性研究接触方式一致。消化道染毒、呼吸道染毒（工业污染物6小时，环境污染物6小时/22-24小时）、皮肤染毒4-6小时。五、剂量设计按药效学或功能学资料：临床治疗或拟用量的倍数（保健食品100、化学药品30、中药50）。受试物A（化学物，LD50为100mg/kg，脂溶性好挥发度小）、受试物B（药品，LD50为1000mg/kg，临床推荐用量每日一次每次两片，每片含300mg）、受试物C（保健品，液体，LD50大于15000mg/kg，推荐使用量10ml/天）。

\InfoLabel{知识来源} 《6 亚慢性和慢性毒性试验》第16页；《6 亚慢性和慢性毒性试验》第17页；《6
亚慢性和慢性毒性试验》第19页；《6 亚慢性和慢性毒性试验》第37页；《6 亚慢性和慢性毒性试验》第38页；《6
亚慢性和慢性毒性试验》第39页

\ContentBand{对应往年题}

\begin{pastquestion}

\QuestionLabel{【往年题1｜15级毒理学考试题目汇总.xls｜口试】}\\
不同类型物质重复染毒实验设计流程注意事项（题号54）；不同化学物在重复暴露实验设计步骤中的区别（题号54）；不同类型物质在重复暴露毒性研究中的问题（题号72）；不同物质重复暴露实验设计有什么要注意的（题号54）

\end{pastquestion}

\begin{pastquestion}

\QuestionLabel{【往年题2｜14级毒理考题汇总.xls｜口试】}\\
不同化学物质在重复暴露实验研究中区别（题号26）；不同物质重复毒性试验的注意事项（题号72，追问二相反应包括/能否代谢活化）

\end{pastquestion}

\begin{pastquestion}

\QuestionLabel{【往年题3｜17级毒理试题汇总.xlsx｜口试】}\\
不同化学物在重复暴露实验时需要考虑的事项（题号72）

\end{pastquestion}

\PointSeparator

\section{蓄积作用（accumulation）的概念、分类、意义及蓄积系数判定（久远题库补充）}\label{users__fpb__26_spring__pkusph_study__tmp__pdfs__prevention-medicine-past-exams__02-ux6bd2ux7406ux5b66ux5f80ux5e74ux9898ux8003ux70b9ux6574ux7406.md__ux84c4ux79efux4f5cux7528accumulationux7684ux6982ux5ff5ux5206ux7c7bux610fux4e49ux53caux84c4ux79efux7cfbux6570ux5224ux5b9aux4e45ux8fdcux9898ux5e93ux8865ux5145}

\InfoLabel{考频} 6 次

\InfoLabel{对应小节} Single/Repeated exposure；四阶段毒理学安全性评价 /《6 亚慢性和慢性毒性试验》第6页；《13
安全性评价》第6页

\InfoLabel{匹配依据}
现有文件在附录中暂归档''蓄积作用''为未完全匹配题；本次临时补充文件夹中的久远题库和卫生毒理学习题均提供完整概念、分类和蓄积系数判定，且《6
亚慢性和慢性毒性试验》明确把蓄积列为单次/重复暴露差异，故补入第六章。

\ContentBand{知识点原文摘取}

Single exposure / Repeated exposure：Accumulation------materials accumulation、effect
accumulation。四阶段毒理学安全性评价第二阶段包括蓄积试验、遗传毒理学试验和致畸试验。

蓄积作用（accumulation）：外源化合物连续、反复进入机体，而且吸收速度超过代谢转化和排泄速度时，化学物质在体内逐渐增加，称为化学物的蓄积作用。分类包括物质蓄积和功能蓄积。蓄积作用是发生亚慢性和慢性毒性的基础；在毒理学评价实际工作中，可根据受试物蓄积毒性强弱评估其毒性作用，也是制定卫生标准时选用安全系数大小的重要参考依据。蓄积系数
K 值常用于判断蓄积强弱：K\textless1 为高度蓄积，1--3 为明显蓄积，3--5 为中等蓄积，≥5 为轻度蓄积。

\InfoLabel{知识来源} 《6 亚慢性和慢性毒性试验》第6页；《13
安全性评价》第6页；相对久远的总结《毒理题库-10年更新.doc》第18题；相对久远的总结《卫生毒理学习题.doc》蓄积系数题

\ContentBand{对应往年题}

\begin{pastquestion}

\QuestionLabel{【往年题1｜21级毒理期末口试回忆.xlsx｜口试】}\\
accumulation的含义，分类及意义（序号42）

\end{pastquestion}

\begin{pastquestion}

\QuestionLabel{【往年题2｜19级毒理期末口试回忆.xlsx｜口试】}\\
accumulation的意义及分类，意义（序号86）；accumulation，分类，毒理学意义（序号忘了/姚碧云）

\end{pastquestion}

\begin{pastquestion}

\QuestionLabel{【往年题3｜17级毒理试题汇总.xlsx｜口试】}\\
什么是accumulation,分几类，毒理学意义是什么（题号42、55）

\end{pastquestion}

\begin{pastquestion}

\QuestionLabel{【往年题4｜20级完整试题回忆.docx｜口试/单项选择】}\\
蓄积作用的分类及毒理学意义（其他同学考题A）；发生慢性毒性的基础是 A选择毒性 B局部毒性 C特殊毒性 D急性毒性
E蓄积作用（第8题）

\end{pastquestion}

\begin{pastquestion}

\QuestionLabel{【往年题5｜相对久远的总结/毒理题库-10年更新.doc｜久远题库补充】}\\
蓄积作用的概念和意义。

\end{pastquestion}

\begin{pastquestion}

\QuestionLabel{【往年题6｜相对久远的总结/卫生毒理学习题.doc｜单项选择/久远题库补充】}\\
蓄积系数法判断化学物蓄积作用强弱。

\end{pastquestion}

\PointSeparator

\chapter{遗传毒理学（一）}\label{users__fpb__26_spring__pkusph_study__tmp__pdfs__prevention-medicine-past-exams__02-ux6bd2ux7406ux5b66ux5f80ux5e74ux9898ux8003ux70b9ux6574ux7406.md__ux7b2cux4e03ux7ae0-7-ux9057ux4f20ux6bd2ux7406ux5b66ux4e00}

\section{DNA
损伤修复的类型（途径）及其与致突变作用的关系}\label{users__fpb__26_spring__pkusph_study__tmp__pdfs__prevention-medicine-past-exams__02-ux6bd2ux7406ux5b66ux5f80ux5e74ux9898ux8003ux70b9ux6574ux7406.md__dna-ux635fux4f24ux4feeux590dux7684ux7c7bux578bux9014ux5f84ux53caux5176ux4e0eux81f4ux7a81ux53d8ux4f5cux7528ux7684ux5173ux7cfb}

\InfoLabel{考频} 11 次

\InfoLabel{对应小节} DNA损伤修复途径；DNA损伤耐受机制；DNA损伤修复与突变 /《7 遗传毒理学（一）》第32--35页

\InfoLabel{匹配依据}
``DNA损伤修复的类型以及DNA损伤修复与突变之间的关系''``基因修复和突变的关系（无误修复/有误修复，DNA修复酶的多态性）''``损伤修复与突变关系''反复出现，对应本小节。

\ContentBand{知识点原文摘取}

DNA损伤修复途径：直接修复（光复活；O6-甲基鸟嘌呤-DNA甲基转移酶修复O6-甲基鸟嘌呤）；切除修复（碱基切除修复------嘧啶二聚体、氧化损伤、辐射、烷化剂引起的损伤；核苷切除修复------底物范围宽；错配碱基修复------错配碱基）；双链断裂修复（同源重组；非同源末端连接）；交联修复------DNA链内交联。DNA损伤耐受机制：易误修复（error-prone
repair，SOS
repair，常为易误修复）；复制后修复（避免了致死性后果，但DNA损伤并未移除）。DNA损伤修复与突变：DNA损伤---DNA损伤修复---突变；DNA损伤修复能力与突变；DNA损伤修复保真性与突变；DNA损伤修复酶的多态性。

\InfoLabel{知识来源} 《7 遗传毒理学（一）》第32页；《7 遗传毒理学（一）》第33页；《7
遗传毒理学（一）》第34页；《7 遗传毒理学（一）》第35页

\ContentBand{对应往年题}

\begin{pastquestion}

\QuestionLabel{【往年题1｜07级毒理口试题回忆.doc｜口试】}\\
基因修复和突变的关系（答：无误修复有哪些------减少了突变；有误修复有哪些------增加了突变，人个体差异：dna修复酶的多态性，辩证地回答）（2号题）；DNA损伤修复与突变的关系（28号题）；突变与dna损伤的关系（46号题）；dna损伤修复与突变的关系（118）

\end{pastquestion}

\begin{pastquestion}

\QuestionLabel{【往年题2｜09级口试题目.docx｜口试】}\\
DNA损伤修复和突变的关系（23）；DNA损伤修复与突变的关系（28）；突变与dna损伤的关系（46）

\end{pastquestion}

\begin{pastquestion}

\QuestionLabel{【往年题3｜21级毒理期末口试回忆.xlsx｜口试】}\\
DNA损伤修复的类型以及DNA损伤修复与突变之间的关系（序号31）；DNA损伤修复能力和致突变作用的关系（序号42，备注：第二题问我损伤是都能修复吗）；DNA损伤修复有哪些？DNA损伤修复与致突变作用的关系（序号74）

\end{pastquestion}

\begin{pastquestion}

\QuestionLabel{【往年题4｜19级毒理期末口试回忆.xlsx｜口试】}\\
突变修复的内容及其与致突变性的关系（序号72，备注：漏了多态性这个点；核苷酸切除修复比碱基切除修复更广泛的个人理解）；DNA损伤修复的途径有哪些？简述DNA修复与致突变作用的关系（序号74）

\end{pastquestion}

\begin{pastquestion}

\QuestionLabel{【往年题5｜17级毒理试题汇总.xlsx｜口试】}\\
DNA损伤修复的类型，与致突变作用的关系（题号50）；DNA损伤修复和致突变性的关系（题号43）；修复与突变关系（题号43）

\end{pastquestion}

\begin{pastquestion}

\QuestionLabel{【往年题6｜15级毒理学考试题目汇总.xls｜口试】}\\
DNA损伤修复和致突变性的关系（题号43）；DNA修复和致突变性关系（题号24）

\end{pastquestion}

\begin{pastquestion}

\QuestionLabel{【往年题7｜14级毒理考题汇总.xls｜口试】}\\
DNA修复和致突变性的关联（题号24）；dna修复和致突变作用的关系（题号4，追问致突变都导致什么）；损失修复与突变的关系（无题号，追问从损伤到突变要经历一个什么过程）

\end{pastquestion}

\PointSeparator

\section{化学致突变作用的不良后果（致突变物的有害影响）}\label{users__fpb__26_spring__pkusph_study__tmp__pdfs__prevention-medicine-past-exams__02-ux6bd2ux7406ux5b66ux5f80ux5e74ux9898ux8003ux70b9ux6574ux7406.md__ux5316ux5b66ux81f4ux7a81ux53d8ux4f5cux7528ux7684ux4e0dux826fux540eux679cux81f4ux7a81ux53d8ux7269ux7684ux6709ux5bb3ux5f71ux54cd}

\InfoLabel{考频} 9 次

\InfoLabel{对应小节} 致突变作用的不良后果 /《7 遗传毒理学（一）》第38页（并见《8 遗传毒理学（二）》第2页）

\InfoLabel{匹配依据}
``化学致突变作用的影响或后果''``致突变物的危害和后果''``突变的不良影响和后果''反复出现，对应本小节体细胞突变与生殖细胞突变后果。

\ContentBand{知识点原文摘取}

致突变作用的不良后果：哺乳动物给予效剂量的诱变剂后诱发突变，分为：体细胞突变------肿瘤、畸胎（如果胚胎接触毒物）；生殖细胞突变------显性致死（不可遗传的改变）、可遗传的改变（基因突变、平衡易位、小的缺失）。

\InfoLabel{知识来源} 《7 遗传毒理学（一）》第38页；《8 遗传毒理学（二）》第2页

\ContentBand{对应往年题}

\begin{pastquestion}

\QuestionLabel{【往年题1｜07级毒理口试题回忆.doc｜口试】}\\
突变对机体的危害（15号题）；突变的遗传学终点并举例（25号题）；化学突变的不良后果和影响（59号题）；化学性致突变剂的不良后果（29号题）；致突变物的有害影响和后果（31号题）

\end{pastquestion}

\begin{pastquestion}

\QuestionLabel{【往年题2｜09级口试题目.docx｜口试】}\\
化学致突变作用的后果（29、142）；突变的不利影响（31）；突变的不良影响和后果（48）

\end{pastquestion}

\begin{pastquestion}

\QuestionLabel{【往年题3｜21级毒理期末口试回忆.xlsx｜口试】}\\
化学致突变作用的影响或后果（序号17）

\end{pastquestion}

\begin{pastquestion}

\QuestionLabel{【往年题4｜19级毒理期末口试回忆.xlsx｜口试】}\\
化学致突变作用的后果（序号40）；化学致突变作用的影响或后果（序号17）

\end{pastquestion}

\begin{pastquestion}

\QuestionLabel{【往年题5｜17级毒理试题汇总.xlsx｜口试】}\\
化学物致突变物所致不良后果（题号1）；致突变物的危害（题号49）

\end{pastquestion}

\begin{pastquestion}

\QuestionLabel{【往年题6｜15级毒理学考试题目汇总.xls｜口试】}\\
致突变的影响或后果（题号29）；突变的后果（题号48）

\end{pastquestion}

\begin{pastquestion}

\QuestionLabel{【往年题7｜14级毒理考题汇总.xls｜口试】}\\
致突变物的危害和后果（题号46）；化学致突变影响（题号9）

\end{pastquestion}

\PointSeparator

\section{染色体（结构）畸变的主要类型及毒理学意义}\label{users__fpb__26_spring__pkusph_study__tmp__pdfs__prevention-medicine-past-exams__02-ux6bd2ux7406ux5b66ux5f80ux5e74ux9898ux8003ux70b9ux6574ux7406.md__ux67d3ux8272ux4f53ux7ed3ux6784ux7578ux53d8ux7684ux4e3bux8981ux7c7bux578bux53caux6bd2ux7406ux5b66ux610fux4e49}

\InfoLabel{考频} 6 次

\InfoLabel{对应小节} 染色体畸变主要类型 /《7 遗传毒理学（一）》第19页

\InfoLabel{匹配依据}
``染色体结构畸变的类型和毒理学意义''``染色体畸变的类型及意义''反复出现，对应本小节裂隙、断裂、断片、缺失、环、易位、倒位等。

\ContentBand{知识点原文摘取}

染色体畸变主要类型：裂隙（gap，在一条或两条染色单体上出现无染色质的区域，但分割的两段染色体仍保持线性）；断裂（break，无染色质区域分割的两段不再保持线性）；断片（fragment，无着丝粒部分与有着丝粒部分分开）；缺失（deletion，有末端缺失和中间缺失）；环（ring，断裂后带着丝粒部分两端连接成环）；易位（translocation，两个非同源染色体断裂后互相交换片段）；倒位（inversion，两处断裂中间节段旋转180°后再连接，有臂间倒位和臂内倒位）。各畸变分染色体型和染色单体型。

\InfoLabel{知识来源} 《7 遗传毒理学（一）》第19页

\ContentBand{对应往年题}

\begin{pastquestion}

\QuestionLabel{【往年题1｜21级毒理期末口试回忆.xlsx｜口试】}\\
染色体结构畸变的类型和毒理学意义（序号68，备注：大礼包不全，要看PPT\&书且好好背）

\end{pastquestion}

\begin{pastquestion}

\QuestionLabel{【往年题2｜19级毒理期末口试回忆.xlsx｜口试】}\\
染色体畸变的类型及意义（意义里要答到肿瘤和致癌）（序号68）；染色体畸变的类型与意义（序号99）

\end{pastquestion}

\begin{pastquestion}

\QuestionLabel{【往年题3｜17级毒理试题汇总.xlsx｜口试】}\\
染色体结构畸变的类型和毒理学意义（题号85，追问染色体型畸变和染色单体型畸变的区别、环状染色体怎么形成的）；染色体结构畸变的类型和意义（题号37，备注：染色体畸变记得分成稳定和非稳定两大类回答）

\end{pastquestion}

\begin{pastquestion}

\QuestionLabel{【往年题4｜20级完整试题回忆.docx｜单项选择】}\\
DNA链上鸟嘌呤被腺嘌呤取代，此种突变称为 A移码突变 B转换 C颠换 D错义突变
E无义突变（第11题）；体细胞突变可能与以下事件有关，除了 A致畸 B致癌 C动脉粥样硬化 D衰老 E遗传性疾病（第12题）

\end{pastquestion}

\PointSeparator

\section{致突变作用的分类（基因突变、染色体突变、基因组突变）及致突变机制}\label{users__fpb__26_spring__pkusph_study__tmp__pdfs__prevention-medicine-past-exams__02-ux6bd2ux7406ux5b66ux5f80ux5e74ux9898ux8003ux70b9ux6574ux7406.md__ux81f4ux7a81ux53d8ux4f5cux7528ux7684ux5206ux7c7bux57faux56e0ux7a81ux53d8ux67d3ux8272ux4f53ux7a81ux53d8ux57faux56e0ux7ec4ux7a81ux53d8ux53caux81f4ux7a81ux53d8ux673aux5236}

\InfoLabel{考频} 3 次

\InfoLabel{对应小节} 致突变作用的分类；DNA损伤类型；致突变作用机制 /《7 遗传毒理学（一）》第11--30页

\InfoLabel{匹配依据}
20级选择题考碱基取代/移码突变的命名；07级口试''基因突变实验不能用Ames实验时用什么检测''涉及突变分类；对应本小节。

\ContentBand{知识点原文摘取}

致突变作用的分类：基因突变（gene mutation------base-pair substitution mutation、frame shift mutation、codon
mutation、fragment mutation）；染色体突变（chromosome mutation------chromosome aberration、chromatid
aberration）；基因组突变（genomic
mutation------aneuploid、euploid）。整码突变又称密码子的插入或缺失。片段突变：基因中小片段核苷酸序列发生改变。致突变作用机制：基因突变、染色体畸变------DNA损伤（碱基类似物取代、共价结合形成加合物、改变碱基结构、大分子嵌入DNA链、DNA-蛋白交联等）；非整倍体、整倍体------有丝分裂或减数分裂器损伤（如对纺锤体的毒作用）。

【翻译】致突变作用可分为基因突变、染色体突变和基因组突变。基因突变包括碱基置换突变、移码突变、密码子突变和片段突变；染色体突变包括染色体畸变和染色单体畸变；基因组突变包括非整倍体和整倍体改变。

\InfoLabel{知识来源} 《7 遗传毒理学（一）》第11页；《7 遗传毒理学（一）》第17页；《7
遗传毒理学（一）》第22页；《7 遗传毒理学（一）》第23页；《7 遗传毒理学（一）》第30页

\ContentBand{对应往年题}

\begin{pastquestion}

\QuestionLabel{【往年题1｜20级完整试题回忆.docx｜单项选择】}\\
DNA链上鸟嘌呤被腺嘌呤取代，此种突变称为（第11题，转换/颠换的判别）

\end{pastquestion}

\begin{pastquestion}

\QuestionLabel{【往年题2｜07级毒理口试题回忆.doc｜口试】}\\
致突变实验的终点（提问，基因突变实验不能用Ames实验时用什么检测，哪些不能用Ames实验）（27号题）

\end{pastquestion}

\begin{pastquestion}

\QuestionLabel{【往年题3｜14级毒理考题汇总.xls｜口试】}\\
姚老师追问免疫/突变固定（题号27、43备注：问上课有没有讲过突变固定）

\end{pastquestion}

\PointSeparator

\section{致突变作用的不良后果------遗传性疾病发生率（背景知识）}\label{users__fpb__26_spring__pkusph_study__tmp__pdfs__prevention-medicine-past-exams__02-ux6bd2ux7406ux5b66ux5f80ux5e74ux9898ux8003ux70b9ux6574ux7406.md__ux81f4ux7a81ux53d8ux4f5cux7528ux7684ux4e0dux826fux540eux679cux9057ux4f20ux6027ux75beux75c5ux53d1ux751fux7387ux80ccux666fux77e5ux8bc6}

\InfoLabel{考频} 1 次

\InfoLabel{对应小节} 遗传性疾病发生率 /《7 遗传毒理学（一）》第41页

\InfoLabel{匹配依据} 部分''突变的危害''答题需扩展到人群遗传性疾病发生率，单列以备复习；本知识点在课件中明确。

\ContentBand{知识点原文摘取}

遗传性疾病发生率：在新生儿中，1\%患有染色体显性突变的疾病；0.25\%患常染色体隐性突变的疾病；0.5\%患性连锁的疾病；0.4\%患与易位及非整倍体等染色体异常有关的疾病；3\%\textasciitilde6\%的新生儿患先天畸形。如果包括晚些时候才发生的和遗传有关的多病因疾病（如心脏病、高血压和糖尿病），受影响的人群比例可达60\%以上。

\InfoLabel{知识来源} 《7 遗传毒理学（一）》第41页

\ContentBand{对应往年题}

\begin{pastquestion}

\QuestionLabel{【往年题1｜07级毒理口试题回忆.doc｜口试】}\\
突变对机体的危害（15号题，扩展遗传性疾病发生率，归入本考点辅助复习）

\end{pastquestion}

\PointSeparator

\chapter{遗传毒理学（二）}\label{users__fpb__26_spring__pkusph_study__tmp__pdfs__prevention-medicine-past-exams__02-ux6bd2ux7406ux5b66ux5f80ux5e74ux9898ux8003ux70b9ux6574ux7406.md__ux7b2cux516bux7ae0-8-ux9057ux4f20ux6bd2ux7406ux5b66ux4e8c}

\section{Ames
试验（细菌回复突变试验）的结果判读、原理、菌株鉴定及方法学评价}\label{users__fpb__26_spring__pkusph_study__tmp__pdfs__prevention-medicine-past-exams__02-ux6bd2ux7406ux5b66ux5f80ux5e74ux9898ux8003ux70b9ux6574ux7406.md__ames-ux8bd5ux9a8cux7ec6ux83ccux56deux590dux7a81ux53d8ux8bd5ux9a8cux7684ux7ed3ux679cux5224ux8bfbux539fux7406ux83ccux682aux9274ux5b9aux53caux65b9ux6cd5ux5b66ux8bc4ux4ef7}

\InfoLabel{考频} 30 余次（各年实验非操作题/口试反复出现，``Ames 结果判读/判定''为多届高频实验非操作题）

\InfoLabel{对应小节} 细菌回复突变试验---Ames Test；Ames试验原理；标准试验菌株的基因型；菌株鉴定；方法学评价
/《8 遗传毒理学（二）》第5--16页

\InfoLabel{匹配依据}
``Ames结果''``ames实验结果判读''``叠氮钠Ames平皿掺入试验结果评定''``Ames试验原理及优缺点，菌株鉴定的方法及意义''在
07/09/14/15/17/19/21/20
各文件大量重复出现，对应本小节。判定要点（结果阳性需满足剂量反应关系、回复突变数达自发回变 2
倍及以上、结果具有可重复性，且说明有无 S9 代谢活化、直接/间接致突变物、突变类型）来自历届回忆与课件原理。

\ContentBand{知识点原文摘取}

Ames试验原理：鼠伤寒沙门氏菌原养型（his+）经正向突变成组氨酸营养缺陷型突变株（his-），受试物±代谢活化系统作用下发生回复突变。Ames测试菌株的遗传特性：rfa改变（细菌荚膜脂多糖屏障，使致突变物有更大的通透性）；ΔuvrB（切除修复系统缺失，增加对很多致突变物的敏感性）；pKM101（增强对依赖SOS系统的致突变物的敏感性）；pAQ1（pAQ1上hisG428的回复性可被测试）。S9混合液（S9mix）：S9、MgCl2、KCl、磷酸盐缓冲液、葡萄糖-6-磷酸溶液、辅酶Ⅱ。菌株鉴定：组氨酸缺陷型（加组氨酸溶液）、自发回复（在标准范围内）、阳性物反应（有/无代谢活化条件下）、rfa改变（结晶紫溶液毒性）、ΔuvrB（紫外线照射）、pKM101（青霉素抗性）、pAQ1（四环素抗性）。方法学评价：优点------致突变检测首选方法、积累资料多、简单快速；缺点------指示生物为原核生物、假阳性和假阴性、不适用某些化学物。

\InfoLabel{知识来源} 《8 遗传毒理学（二）》第6页；《8 遗传毒理学（二）》第8页；《8
遗传毒理学（二）》第9页；《8 遗传毒理学（二）》第13页；《8 遗传毒理学（二）》第15页；《8
遗传毒理学（二）》第16页

\ContentBand{对应往年题}

\begin{pastquestion}

\QuestionLabel{【往年题1｜07级毒理口试题回忆.doc｜实验/口试】}\\
Ames实验结果鉴定及意义（3号题）；AMES
结果读图（4号题）；Ames试验点试验结果判定（4号、15号、46号、54号题）；Ames试验菌株鉴定的内容和意义（36号题）；MNNT/MNNG的ames试验点试验结果判定（10号、60号题）；Ames试验和微核试验的原理及观察终点（43号题）

\end{pastquestion}

\begin{pastquestion}

\QuestionLabel{【往年题2｜09级口试题目.docx｜实验/口试】}\\
叠氮钠的Ames试验结果判定（要点：TA100菌株、呈剂量反应关系，超过自发回变数）（10、28、42）；Ames试验菌种鉴定的内容和意义（菌株组合提高灵敏性）（3）；Ames实验平板掺入评价（29、40）；NaN3的AMES结果评定（28）

\end{pastquestion}

\begin{pastquestion}

\QuestionLabel{【往年题3｜21级毒理期末口试回忆.xlsx｜实验非操作题/口试】}\\
Ames结果 / ames实验结果判读 /
Ames平板掺入实验结果判读（序号8、16、22、24、26等大量序号，备注：判断阳性结果有三点①剂量反应关系②回复突变数达自发回变2倍及以上③结果具有可重复性，第三点容易漏掉）；Ames试验的原理、优缺点，基因型鉴定的内容、意义（序号85）

\end{pastquestion}

\begin{pastquestion}

\QuestionLabel{【往年题4｜19级毒理期末口试回忆.xlsx｜实验非操作题/口试】}\\
Ames结果（序号41、86、97、18等大量序号，备注：ames结果要从阳性阴性结果、直接间接诱变物、突变类型几方面来答；TA100是碱基置换）；Ames实验原理和优缺点，菌株鉴定的方法和意义（序号27）

\end{pastquestion}

\begin{pastquestion}

\QuestionLabel{【往年题5｜17级毒理试题汇总.xlsx｜实验非操作题/口试】}\\
ames结果判断 /
叠氮钠阳性结果判断（题号9、42、34等大量题号，备注：ames要答``有无代谢活化系统，是直接/间接致突变物''，TA100是什么类型的突变）；Ames实验原理及优缺点，标准菌株鉴定的方法及意义（菌株组合提高灵敏性）（题号61）

\end{pastquestion}

\begin{pastquestion}

\QuestionLabel{【往年题6｜15级毒理学考试题目汇总.xls｜实验非操作题/口试】}\\
ames结果 / NaN3 Ames实验结果解读 /
叠氮钠平板掺入法结果解读（题号11、50、73等大量题号，备注：ames实验不要提回变数两倍以上！杨老师有毒；只答存在剂量反应关系和叠氮钠是阳性致突变物）；ames实验和微核试验的原理和遗传学重点（无题号）

\end{pastquestion}

\begin{pastquestion}

\QuestionLabel{【往年题7｜14级毒理考题汇总.xls｜实验非操作题】}\\
Ames结果判断（题号50、52、80等大量题号，备注：标准答案------随叠氮钠剂量水平增加菌落数也增加，说明存在剂量反应关系，可认为叠氮钠对TA100菌株有致突变性；不要说回变菌落数逐渐增加）

\end{pastquestion}

\begin{pastquestion}

\QuestionLabel{【往年题8｜20级完整试题回忆.docx｜实验/口试】}\\
请对受试物NaN3一次Ames平皿掺入试验（见图片）进行结果评定（序号76，0.1/0.5/1/2μg；杨老师希望听到三个关键词''剂量反应关系''\,``结果阳性''``致突变性''，并强调要完整严谨）

\end{pastquestion}

\PointSeparator

\section{突变的遗传学终点及举例（常用致突变试验小结）}\label{users__fpb__26_spring__pkusph_study__tmp__pdfs__prevention-medicine-past-exams__02-ux6bd2ux7406ux5b66ux5f80ux5e74ux9898ux8003ux70b9ux6574ux7406.md__ux7a81ux53d8ux7684ux9057ux4f20ux5b66ux7ec8ux70b9ux53caux4e3eux4f8bux5e38ux7528ux81f4ux7a81ux53d8ux8bd5ux9a8cux5c0fux7ed3}

\InfoLabel{考频} 9 次

\InfoLabel{对应小节} 常用致突变试验小结 /《8 遗传毒理学（二）》第38页（并见《7》第42页 遗传毒理学试验分类）

\InfoLabel{匹配依据}
``突变的遗传学终点及举例''``遗传毒性检测的终点并举例''``四个遗传学终点并举例''反复出现，对应本小节各试验对应的遗传学终点。

\ContentBand{知识点原文摘取}

常用致突变试验小结：Ames
Test（细菌，体外，基因突变）；哺乳动物细胞基因突变试验（哺乳动物，体外，基因突变）；染色体畸变试验（哺乳动物，体内或体外，染色体畸变）；显性致死试验（哺乳动物，体内，染色体畸变）；微核试验（哺乳动物，体内或体外，染色体畸变、非整倍体）；姐妹染色单体交换试验（哺乳动物，体内或体外，DNA损伤）；单细胞凝胶电泳试验（哺乳动物，体内或体外，DNA损伤）；程序外DNA合成试验（哺乳动物，体内或体外，DNA损伤）；转基因动物致突变试验（哺乳动物，体内，基因突变）。遗传毒理学试验分类：基因突变试验、染色体损伤试验、非整倍体试验、DNA损伤的试验。

\InfoLabel{知识来源} 《8 遗传毒理学（二）》第38页；《7 遗传毒理学（一）》第42页

\ContentBand{对应往年题}

\begin{pastquestion}

\QuestionLabel{【往年题1｜07级毒理口试题回忆.doc｜口试】}\\
突变的遗传学终点并举例（25号题）；致突变实验的终点（27号题）；遗传毒性检测的终点并举例（52号题）；遗传毒理学试验的终点（35号题）

\end{pastquestion}

\begin{pastquestion}

\QuestionLabel{【往年题2｜09级口试题目.docx｜口试】}\\
四个遗传学终点并举例（25）；遗传终点并举例（27、52）；遗传学终点，举例（52）

\end{pastquestion}

\begin{pastquestion}

\QuestionLabel{【往年题3｜21级毒理期末口试回忆.xlsx｜口试】}\\
突变的遗传学终点及举例（序号95）

\end{pastquestion}

\begin{pastquestion}

\QuestionLabel{【往年题4｜17级毒理试题汇总.xlsx｜口试】}\\
遗传毒性试验终点和举例（题号96）；遗传毒理终点，举例（题号27）；突变的遗传学终点及举例（题号95）

\end{pastquestion}

\begin{pastquestion}

\QuestionLabel{【往年题5｜15级毒理学考试题目汇总.xls｜口试】}\\
遗传毒性的主要终点，举例（题号46）；遗传毒理学实验的终点及举例（题号66）；遗传实验终点（题号27）

\end{pastquestion}

\begin{pastquestion}

\QuestionLabel{【往年题6｜14级毒理考题汇总.xls｜口试】}\\
遗传毒理学终点（题号46）；遗传学实验终点类型及举例（题号66）

\end{pastquestion}

\PointSeparator

\section{遗传毒性试验成组应用的组合原则、在预测致癌性中的作用与局限（不确定性的原因）}\label{users__fpb__26_spring__pkusph_study__tmp__pdfs__prevention-medicine-past-exams__02-ux6bd2ux7406ux5b66ux5f80ux5e74ux9898ux8003ux70b9ux6574ux7406.md__ux9057ux4f20ux6bd2ux6027ux8bd5ux9a8cux6210ux7ec4ux5e94ux7528ux7684ux7ec4ux5408ux539fux5219ux5728ux9884ux6d4bux81f4ux764cux6027ux4e2dux7684ux4f5cux7528ux4e0eux5c40ux9650ux4e0dux786eux5b9aux6027ux7684ux539fux56e0}

\InfoLabel{考频} 12 次

\InfoLabel{对应小节} 遗传毒理学试验的组合应用；灵敏性/特异性 /《8 遗传毒理学（二）》第39--47页

\InfoLabel{匹配依据}
``遗传毒理学对致癌性预测的试验的组成原则以及不确定的原因''``遗传毒性试验预测致癌性的作用和局限性''``遗传毒理学实验预测致癌性的不确定性''反复出现，对应本小节成组原则与灵敏性/特异性局限。

\ContentBand{知识点原文摘取}

灵敏性（sensitivity）：受试致癌物中，在遗传毒理学试验中呈致突变性阳性结果的比例。特异性（specificity）：受试非致癌物中，在遗传毒理学试验中呈阴性结果的比例。遗传毒性试验成组应用组合原则：①应包括多个遗传学终点；②试验指示生物包括若干进化阶段的物种；③应包括体外试验和体内试验；④在预测可遗传的危害时，应包括体细胞及性细胞的试验。Battery
concept：每个短期试验都有其固有的优势和局限，联合使用时可相互重叠和补偿。（预测不确定的原因：不能预测遗传非致癌物及非遗传致癌物，只能预测非遗传非致癌物及遗传致癌物，故不准确。）

\InfoLabel{知识来源} 《8 遗传毒理学（二）》第40页；《8 遗传毒理学（二）》第41页；《8
遗传毒理学（二）》第46页；《8 遗传毒理学（二）》第47页

\ContentBand{对应往年题}

\begin{pastquestion}

\QuestionLabel{【往年题1｜07级毒理口试题回忆.doc｜口试】}\\
遗传毒理学实验对预测致癌性的作用和局限性（45号、57号题）；遗传毒理学成组试验成组的原则、预测不肯定的原因（30号题）

\end{pastquestion}

\begin{pastquestion}

\QuestionLabel{【往年题2｜09级口试题目.docx｜口试】}\\
遗传毒理学对致癌性预测的试验的组成原则以及不确定的原因（注意不要答生殖细胞/体细胞那条）（7）；遗传毒理学成组实验的成组原则，预测不肯定的原因（34）；遗传毒理学试验成组应用原则及预测不准确的原因（12，备注：不能预测遗传非致癌物及非遗传致癌物，只能预测非遗传非致癌物及遗传致癌物）；遗传毒理学实验预测致癌性的不确定性（57）

\end{pastquestion}

\begin{pastquestion}

\QuestionLabel{【往年题3｜21级毒理期末口试回忆.xlsx｜口试】}\\
遗传毒理学试验在致癌物预测和评定遗传毒性的意义和局限性（序号39）；遗传毒性实验研究致癌的试验组成原则？我国标准的试验组合？（序号45）

\end{pastquestion}

\begin{pastquestion}

\QuestionLabel{【往年题4｜19级毒理期末口试回忆.xlsx｜口试】}\\
遗传毒理学实验在致畸物筛选和评价中的作用和局限（20级补充A/B）；遗传毒性试验致癌性判断的组合原则，以及我国的标准组合（序号4）

\end{pastquestion}

\begin{pastquestion}

\QuestionLabel{【往年题5｜17级毒理试题汇总.xlsx｜口试】}\\
遗传毒性实验预测致癌性的作用和局限性（题号51）；由遗传毒理学实验推致突变性的不确定性（题号3，备注：要答原癌基因抑癌基因、遗传致癌物非遗传非致癌物）；遗传毒性试验致癌性判断的组合原则，以及我国的标准组合（题号4，备注：要背我国的标准）

\end{pastquestion}

\begin{pastquestion}

\QuestionLabel{【往年题6｜15级毒理学考试题目汇总.xls｜口试】}\\
由遗传毒理学实验推致突变性的不确定性（题号5）；遗传毒性实验推测致癌性的不确定性（题号25）；为什么用遗传实验推测致癌实验存在不确定性（题号5）；遗传毒理学实验预测致癌性的不确定性（题号25）；遗传毒理试验预测致癌性时有不确定性的原因（题号44）

\end{pastquestion}

\begin{pastquestion}

\QuestionLabel{【往年题7｜14级毒理考题汇总.xls｜口试】}\\
遗传毒性推致癌性的不确定性（题号25）；遗传毒理学组合原则（题号26）；遗传学实验组合的原则，我国推荐的实验组合（题号65，备注：第二题让举具体实验的例子）；遗传毒理学试验的组合应用（题号45）；为什么用遗传毒理学实验预测致癌性有不肯定性（题号5）

\end{pastquestion}

\PointSeparator

\section{我国《药物遗传毒性研究技术指导原则》及食品/农药安全性评价推荐的标准试验组合}\label{users__fpb__26_spring__pkusph_study__tmp__pdfs__prevention-medicine-past-exams__02-ux6bd2ux7406ux5b66ux5f80ux5e74ux9898ux8003ux70b9ux6574ux7406.md__ux6211ux56fdux836fux7269ux9057ux4f20ux6bd2ux6027ux7814ux7a76ux6280ux672fux6307ux5bfcux539fux5219ux53caux98dfux54c1ux519cux836fux5b89ux5168ux6027ux8bc4ux4ef7ux63a8ux8350ux7684ux6807ux51c6ux8bd5ux9a8cux7ec4ux5408}

\InfoLabel{考频} 5 次

\InfoLabel{对应小节} 遗传毒性试验成组应用组合原则；我国推荐的标准试验组合 /《8
遗传毒理学（二）》第48页、第49页、第53页

\InfoLabel{匹配依据}
``遗传毒性试验预测致癌性时的试验组成原则？我国药物遗传毒性研究技术指导原则推荐的标准试验组合''``我国推荐的实验组合''直接对应本小节。

\ContentBand{知识点原文摘取}

我国《药物遗传毒性研究技术指导原则》推荐的标准试验组合：①一项体外细菌基因突变试验；②一项采用哺乳动物细胞进行的体外染色体损伤评估试验，或体外小鼠淋巴瘤tk试验；③一项采用啮齿类动物造血细胞进行的体内染色体损伤试验。对于结果为阴性的受试物，完成上述三项试验组合通常可提示其无遗传毒性。《食品安全性毒理学评价程序》（2014）组合1：细菌回复突变试验，哺乳动物红细胞微核试验或哺乳动物骨髓细胞染色体畸变试验；小鼠精原细胞或精母细胞染色体畸变试验或啮齿类动物显性致死试验。

\InfoLabel{知识来源} 《8 遗传毒理学（二）》第48页；《8 遗传毒理学（二）》第49页；《8 遗传毒理学（二）》第53页

\ContentBand{对应往年题}

\begin{pastquestion}

\QuestionLabel{【往年题1｜20级完整试题回忆.docx｜口试】}\\
遗传毒性试验预测致癌性时的试验组成原则？我国药物遗传毒性研究技术指导原则推荐的标准试验组合？（序号76的B）

\end{pastquestion}

\begin{pastquestion}

\QuestionLabel{【往年题2｜21级毒理期末口试回忆.xlsx｜口试】}\\
遗传毒性实验研究致癌的试验组成原则？我国标准的试验组合？（序号45）

\end{pastquestion}

\begin{pastquestion}

\QuestionLabel{【往年题3｜17级毒理试题汇总.xlsx｜口试】}\\
遗传毒性试验致癌性判断的组合原则，以及我国的标准组合是什么！！（题号4，备注：要背我国的标准！）

\end{pastquestion}

\begin{pastquestion}

\QuestionLabel{【往年题4｜19级毒理期末口试回忆.xlsx｜口试】}\\
遗传学实验组合的原则，我国推荐的实验组合（序号76）

\end{pastquestion}

\begin{pastquestion}

\QuestionLabel{【往年题5｜14级毒理考题汇总.xls｜口试】}\\
遗传学实验组合的原则，我国推荐的实验组合（题号65）

\end{pastquestion}

\PointSeparator

\section{微核试验（微核的概念）及骨髓多染红细胞微核试验/阅片（PCE、NCE、MN）}\label{users__fpb__26_spring__pkusph_study__tmp__pdfs__prevention-medicine-past-exams__02-ux6bd2ux7406ux5b66ux5f80ux5e74ux9898ux8003ux70b9ux6574ux7406.md__ux5faeux6838ux8bd5ux9a8cux5faeux6838ux7684ux6982ux5ff5ux53caux9aa8ux9ad3ux591aux67d3ux7ea2ux7ec6ux80deux5faeux6838ux8bd5ux9a8cux9605ux7247pcencemn}

\InfoLabel{考频} 30 余次（各年实验操作题/实验非操作题''微核制片''\,``微核阅片/找PCE NCE MN''反复出现）

\InfoLabel{对应小节} 微核试验；小鼠骨髓多染红细胞微核试验 /《8 遗传毒理学（二）》第21--24页

\InfoLabel{匹配依据}
``微核制片''``微核阅片，找PCE，NCE，MN''``镜下找出PCE、NCE和微核细胞''为多届固定实验题，对应本小节微核定义与骨髓多染红细胞微核试验。

\ContentBand{知识点原文摘取}

微核：在有丝分裂（减数分裂）末期迟滞的染色体断片或整条染色体生成的，分离于细胞主核之外的小核。小鼠骨髓多染红细胞微核试验（Micronucleus
assay in bone marrow polychromatic
erythrocytes）。微核试验遗传学终点为染色体畸变、非整倍体。（历届补充：微核的大小为细胞大小的1/20\textasciitilde1/5；微核制片为推片不是拉片，往钝角方向推；阅片需掌握PCE、NCE、MN三种细胞的判别标准。）

\InfoLabel{知识来源} 《8 遗传毒理学（二）》第21页；《8 遗传毒理学（二）》第22页；《8 遗传毒理学（二）》第38页

\ContentBand{对应往年题}

\begin{pastquestion}

\QuestionLabel{【往年题1｜07级毒理口试题回忆.doc｜实验操作】}\\
镜下找出PCE、NCE和微核细胞（8号题）；微核涂片（22号题）；微核骨髓制片（30号、44号题，老师会问小牛血清的剂量、怎样评价涂片的好坏）；辨认PCE、NCE、MN（17号、80号题）；正染、多染、微核的读片（59号题）；找多染，正染，微核（52号题）

\end{pastquestion}

\begin{pastquestion}

\QuestionLabel{【往年题2｜09级口试题目.docx｜实验操作】}\\
小鼠微核制片（14、22、44）；微核骨髓采样、制片技术（22）；微核涂片找正染，多染，微核（52）

\end{pastquestion}

\begin{pastquestion}

\QuestionLabel{【往年题3｜21级毒理期末口试回忆.xlsx｜实验操作/实验非操作】}\\
微核制片 /
微核阅片（序号3、11、14、19等大量序号，备注：微核的大小是细胞大小的1/20-1/5，要说清楚PCE、NCE和微核的判断标准）

\end{pastquestion}

\begin{pastquestion}

\QuestionLabel{【往年题4｜19级毒理期末口试回忆.xlsx｜实验操作/实验非操作】}\\
微核制片 / 微核阅片（序号41、49、5等大量序号，备注：微核阅片会问PCE和NCE的判别标准）

\end{pastquestion}

\begin{pastquestion}

\QuestionLabel{【往年题5｜17级毒理试题汇总.xlsx｜实验操作/实验非操作】}\\
微核阅片，找PCE，NCE，MN / 微核制片（题号61、85、95等大量题号，备注：微核制片注意推片方向，别弄成刮片）

\end{pastquestion}

\begin{pastquestion}

\QuestionLabel{【往年题6｜15级毒理学考试题目汇总.xls｜实验操作/实验非操作】}\\
微核试验制片 / 微核读片找微核 /
PCE、NCE、MN辨认及要点（题号11、5、76等大量题号，备注：追问了微核大小占细胞大小的比例）

\end{pastquestion}

\begin{pastquestion}

\QuestionLabel{【往年题7｜14级毒理考题汇总.xls｜实验操作/实验非操作】}\\
微核制片 / 微核阅片（题号50、46、78等大量题号）

\end{pastquestion}

\begin{pastquestion}

\QuestionLabel{【往年题8｜20级完整试题回忆.docx｜实验操作】}\\
辨认PCE NCE MN，老师会问这些单词的含义及辨别依据（其他同学考题）

\end{pastquestion}

\PointSeparator

\section{显性致死试验（dominant lethal
test）的概念和评价方法}\label{users__fpb__26_spring__pkusph_study__tmp__pdfs__prevention-medicine-past-exams__02-ux6bd2ux7406ux5b66ux5f80ux5e74ux9898ux8003ux70b9ux6574ux7406.md__ux663eux6027ux81f4ux6b7bux8bd5ux9a8cdominant-lethal-testux7684ux6982ux5ff5ux548cux8bc4ux4ef7ux65b9ux6cd5}

\InfoLabel{考频} 3 次

\InfoLabel{对应小节} 显性致死试验 /《8 遗传毒理学（二）》第25--26页

\InfoLabel{匹配依据} ``简述dominal lethal 的概念和评价方法''``显性致死概念 试验''直接对应本小节。

\ContentBand{知识点原文摘取}

显性致死：引起胚胎或胎儿死亡的相应。显性致死突变：发生于生殖细胞的突变，它不引起受精障碍，但能引起受精卵或发育中的胚胎死亡。引起显性致死的突变主要是染色体断裂。小鼠、大鼠精子分化阶段与交配周次的关系：与受试物时精子所处的分化阶段（输精管及睾丸中的精子、精细胞、精母细胞、精原细胞）对应不同交配周次。

\InfoLabel{知识来源} 《8 遗传毒理学（二）》第25页；《8 遗传毒理学（二）》第26页；《8 遗传毒理学（二）》第38页

\ContentBand{对应往年题}

\begin{pastquestion}

\QuestionLabel{【往年题1｜07级毒理口试题回忆.doc｜口试】}\\
简述dominal lethal 的概念和评价方法（51号题）

\end{pastquestion}

\begin{pastquestion}

\QuestionLabel{【往年题2｜09级口试题目.docx｜口试】}\\
显性致死概念 试验（题号未知）

\end{pastquestion}

\begin{pastquestion}

\QuestionLabel{【往年题3｜08级（07级口试题回忆.doc 引用）｜口试】}\\
显性致死的概念及评价方法（归入本考点辅助复习）

\end{pastquestion}

\PointSeparator

\section{Ames
试验与微核试验的原理和检测终点（合并对比题）}\label{users__fpb__26_spring__pkusph_study__tmp__pdfs__prevention-medicine-past-exams__02-ux6bd2ux7406ux5b66ux5f80ux5e74ux9898ux8003ux70b9ux6574ux7406.md__ames-ux8bd5ux9a8cux4e0eux5faeux6838ux8bd5ux9a8cux7684ux539fux7406ux548cux68c0ux6d4bux7ec8ux70b9ux5408ux5e76ux5bf9ux6bd4ux9898}

\InfoLabel{考频} 8 次

\InfoLabel{对应小节} 细菌回复突变试验---Ames Test；微核试验；常用致突变试验小结 /《8
遗传毒理学（二）》第6页、第21页、第38页

\InfoLabel{匹配依据}
``Ames试验和微核试验的原理和检测终点''``ames实验和微核试验的原理和遗传学终点''作为合并设问反复出现，对应两试验原理与遗传学终点。

\ContentBand{知识点原文摘取}

Ames试验原理：鼠伤寒沙门氏菌组氨酸营养缺陷型（his-）在受试物±代谢活化系统作用下发生回复突变（his- →
HIS+），遗传学终点为基因突变。微核：在有丝分裂（减数分裂）末期迟滞的染色体断片或整条染色体生成的、分离于细胞主核之外的小核；微核试验遗传学终点为染色体畸变、非整倍体。

\InfoLabel{知识来源} 《8 遗传毒理学（二）》第6页；《8 遗传毒理学（二）》第21页；《8 遗传毒理学（二）》第38页

\ContentBand{对应往年题}

\begin{pastquestion}

\QuestionLabel{【往年题1｜07级毒理口试题回忆.doc｜口试】}\\
Ames试验和微核试验的原理及观察终点（观察终点要答是哪种遗传损伤）（43号题）

\end{pastquestion}

\begin{pastquestion}

\QuestionLabel{【往年题2｜15级毒理学考试题目汇总.xls｜口试】}\\
ames实验和微核试验的原理和遗传学重点（无题号）；Ames试验和微核试验的原理和检测终点（题号8、47、67）；ames实验与微核试验（题号47）；ames试验和微核试验的原理和检测终点（题号28）

\end{pastquestion}

\begin{pastquestion}

\QuestionLabel{【往年题3｜14级毒理考题汇总.xls｜口试】}\\
Ames试验和微核试验的原理和检测终点（题号8）；Ames实验和微核试验的实验原理和实验终点（题号47）；ames和微核的原理（题号28）

\end{pastquestion}

\PointSeparator

\section{单细胞凝胶电泳（彗星试验，Comet
Assay）检测的损伤类型}\label{users__fpb__26_spring__pkusph_study__tmp__pdfs__prevention-medicine-past-exams__02-ux6bd2ux7406ux5b66ux5f80ux5e74ux9898ux8003ux70b9ux6574ux7406.md__ux5355ux7ec6ux80deux51ddux80f6ux7535ux6cf3ux5f57ux661fux8bd5ux9a8ccomet-assayux68c0ux6d4bux7684ux635fux4f24ux7c7bux578b}

\InfoLabel{考频} 1 次

\InfoLabel{对应小节} 单细胞凝胶电泳试验（Comet Assay）；常用致突变试验小结 /《8
遗传毒理学（二）》第30--31页、第38页

\InfoLabel{匹配依据} YYS 题库补充明确给出''彗星实验检验的突变类型''选择题，对应本小节单细胞凝胶电泳检测 DNA
损伤（DNA 链断裂）。

\ContentBand{知识点原文摘取}

单细胞凝胶电泳试验（SCGE，Comet
Assay）：在单细胞水平上定量检测DNA损伤；在电泳槽中，DNA断片在电场作用下由细胞核中移出向阳极游动，经荧光染色后形成如同彗星一样的彗头和彗尾，DNA迁移的程度与单链DNA断裂的水平有关；技术简单、快速、价廉、检测DNA损伤的敏感性高。常用致突变试验小结：单细胞凝胶电泳试验遗传学终点为DNA损伤。

\InfoLabel{知识来源} 《8 遗传毒理学（二）》第30页；《8 遗传毒理学（二）》第31页；《8 遗传毒理学（二）》第38页

\ContentBand{对应往年题}

\begin{pastquestion}

\QuestionLabel{【往年题1｜毒理学-来自课本的题库补充byYYS.docx｜单项选择】}\\
彗星实验检验的突变类型是：A 基因突变 B DNA损伤 C 非整倍体突变 D 整倍体突变 E 染色体损伤（2025/5/18 补充题）

\end{pastquestion}

\PointSeparator

\section{其他遗传毒性试验：SCE、UDS、转基因动物致突变试验和睾丸染色体畸变试验（久远题库补充）}\label{users__fpb__26_spring__pkusph_study__tmp__pdfs__prevention-medicine-past-exams__02-ux6bd2ux7406ux5b66ux5f80ux5e74ux9898ux8003ux70b9ux6574ux7406.md__ux5176ux4ed6ux9057ux4f20ux6bd2ux6027ux8bd5ux9a8csceudsux8f6cux57faux56e0ux52a8ux7269ux81f4ux7a81ux53d8ux8bd5ux9a8cux548cux777eux4e38ux67d3ux8272ux4f53ux7578ux53d8ux8bd5ux9a8cux4e45ux8fdcux9898ux5e93ux8865ux5145}

\InfoLabel{考频} 4 次

\InfoLabel{对应小节} 姐妹染色单体交换试验；程序外DNA合成试验；转基因动物致突变试验；常用致突变试验小结 /《8
遗传毒理学（二）》第27--38页

\InfoLabel{匹配依据} 久远题库和卫生毒理学习题除 Ames、微核、彗星、显性致死外，还反复涉及
SCE、UDS、转基因动物致突变试验、睾丸染色体畸变试验等检测终点；现有整理已有''突变遗传学终点''总述，但未单列这些替代/补充试验。

\ContentBand{知识点原文摘取}

姐妹染色单体交换试验（SCE）：染色体同源位点上 DNA 复制产物的相交换，其频率与 DNA 的断裂和修复有关。程序外 DNA
合成试验（UDS）：正常细胞仅在 S 期进行 DNA 复制合成，当 DNA 受损后，DNA 修复合成可发生在 S 期以外；可用
3H-胸腺嘧啶核苷掺入 DNA 的放射活性评价 DNA 修复合成程度，间接反映 DNA
损伤程度。转基因动物致突变试验：转基因动物是基因组中整合有实验方法导入的 DNA 并能遗传给后代的动物，如 BigBlue
Mouse 以大肠杆菌 LacI 为靶基因，Muta Mouse 以大肠杆菌 LacZ
为靶基因。常用致突变试验小结中，SCE、单细胞凝胶电泳和 UDS 的遗传学终点为 DNA
损伤，转基因动物致突变试验终点为基因突变；睾丸染色体畸变试验可检测生殖细胞染色体畸变。

\InfoLabel{知识来源} 《8 遗传毒理学（二）》第27页；《8 遗传毒理学（二）》第33页；《8
遗传毒理学（二）》第34页；《8 遗传毒理学（二）》第35页；《8 遗传毒理学（二）》第38页

\ContentBand{对应往年题}

\begin{pastquestion}

\QuestionLabel{【往年题1｜相对久远的总结/卫生毒理学习题.doc｜单项选择/久远题库补充】}\\
SCE、细菌回复突变试验、微核试验、哺乳动物细胞正向突变试验、显性致死试验、彗星试验、转基因动物致突变试验、程序外DNA合成等试验的检测终点和适用对象。

\end{pastquestion}

\begin{pastquestion}

\QuestionLabel{【往年题2｜相对久远的总结/毒理题库-10年更新.doc｜久远题库补充】}\\
遗传毒理学试验的终点：基因突变试验、染色体损伤试验、非整倍体试验和原发性 DNA 损伤试验；原发性 DNA
损伤包括姐妹染色单体交换试验、单细胞凝胶电泳、程序外 DNA 合成试验等。

\end{pastquestion}

\begin{pastquestion}

\QuestionLabel{【往年题3｜相对久远的总结/毒理题库-12年更新.doc｜久远题库补充】}\\
突变的遗传学终点并举例；基因突变实验不能用 Ames 实验时用什么检测，哪些不能用 Ames 实验。

\end{pastquestion}

\begin{pastquestion}

\QuestionLabel{【往年题4｜04级毒理整理和部分考题.docx｜口试/久远题库补充】}\\
致突变实验的终点；遗传毒理学的观察终点和相应试验。

\end{pastquestion}

\PointSeparator

\chapter{化学致癌作用}\label{users__fpb__26_spring__pkusph_study__tmp__pdfs__prevention-medicine-past-exams__02-ux6bd2ux7406ux5b66ux5f80ux5e74ux9898ux8003ux70b9ux6574ux7406.md__ux7b2cux4e5dux7ae0-9-ux5316ux5b66ux81f4ux764cux4f5cux7528}

\section{化学致癌的多阶段学说（引发、促长、进展阶段）及形态学/生物学特征}\label{users__fpb__26_spring__pkusph_study__tmp__pdfs__prevention-medicine-past-exams__02-ux6bd2ux7406ux5b66ux5f80ux5e74ux9898ux8003ux70b9ux6574ux7406.md__ux5316ux5b66ux81f4ux764cux7684ux591aux9636ux6bb5ux5b66ux8bf4ux5f15ux53d1ux4fc3ux957fux8fdbux5c55ux9636ux6bb5ux53caux5f62ux6001ux5b66ux751fux7269ux5b66ux7279ux5f81}

\InfoLabel{考频} 14 次

\InfoLabel{对应小节} 化学致癌的多阶段学说；引发剂、促长剂、进展剂；多阶段致癌的形态学和生物学特征 /《9
化学致癌作用》第54--68页

\InfoLabel{匹配依据}
``癌症发生的多阶段学说''``致癌三阶段学说介绍''``多阶段致癌的形态学和生物学特征''反复出现，对应本小节。

\ContentBand{知识点原文摘取}

化学致癌的多阶段学说：正常细胞经过遗传学改变的积累，才能转变为癌细胞，癌症的发生是多阶段过程，至少包括引发、促长和进展阶段。引发阶段（Initiation）：化学物或其活性代谢产物与DNA作用，导致体细胞突变成引发细胞的阶段，是一个遗传毒性事件；引发剂本身有致癌性，大多数是致突变物，没有可检测的阈剂量，靶主要是原癌基因和肿瘤抑制基因；引发阶段很短，引发细胞不是肿瘤细胞，引发作用不可逆并且是累积性的。促长阶段（Promotion）：引发细胞增殖成为癌前病变或良性肿瘤的过程，主要是非遗传机制；历时较长，持续给以促长剂是必需的，早期有可逆性、晚期不可逆；促长剂本身不能诱发肿瘤，常是非致突变物，通常具有阈剂量，其剂量反应关系呈S形曲线。进展阶段（Progression）：从引发细胞群（癌前病变、良性肿瘤）转变成恶性肿瘤的过程，不可逆，主要表现为核型不稳定及表型的异质性；兼有引发剂、促长剂和进展剂作用的化学致癌物称为完全致癌物。

\InfoLabel{知识来源} 《9 化学致癌作用》第54页；《9 化学致癌作用》第57页；《9 化学致癌作用》第60页；《9
化学致癌作用》第61页；《9 化学致癌作用》第63页；《9 化学致癌作用》第64页；《9 化学致癌作用》第65页；《9
化学致癌作用》第66页

\ContentBand{对应往年题}

\begin{pastquestion}

\QuestionLabel{【往年题1｜07级毒理口试题回忆.doc｜口试】}\\
致癌的三阶段学说介绍（2号题）；癌变的多阶段学说（14号题）；癌症发生的多阶段学说（5号、26号题）；多阶段研究（22号题）；解释癌症的多阶段学说（5号题）；简述致癌多阶段学说（44号题）

\end{pastquestion}

\begin{pastquestion}

\QuestionLabel{【往年题2｜09级口试题目.docx｜口试】}\\
致癌三阶段（引发和促长阶段的重点都是什么细胞？）（2）；肿瘤发生的多阶段学说（5、14）；肿瘤发生机制的三阶段学说（22）；多阶段致癌学说（26）

\end{pastquestion}

\begin{pastquestion}

\QuestionLabel{【往年题3｜21级毒理期末口试回忆.xlsx｜口试】}\\
多阶段学说（序号79）；癌症三阶段详解、组织学与生理学意义（序号78）

\end{pastquestion}

\begin{pastquestion}

\QuestionLabel{【往年题4｜17级毒理试题汇总.xlsx｜口试】}\\
癌症的多阶段学说（题号62）；肿瘤发生的多阶段学说（题号41）；多阶段致癌的形态学和生物学特征（题号36，备注：促长阶段是良性肿瘤进展阶段是恶性肿瘤）；致癌三阶段的形态学变化和生理学意义（题号84）

\end{pastquestion}

\begin{pastquestion}

\QuestionLabel{【往年题5｜15级毒理学考试题目汇总.xls｜口试】}\\
多阶段致癌的生物学和形态学特点（题号12）；致癌多阶段中的形态学和生物学特征（题号51，备注：按照书上总结的表格背一下）

\end{pastquestion}

\begin{pastquestion}

\QuestionLabel{【往年题6｜14级毒理考题汇总.xls｜口试】}\\
致癌多阶段（题号60）；多阶段致癌的形态学生物学要点（表格）（题号12）；致癌多阶段的生理学与形态学特征（无题号）

\end{pastquestion}

\PointSeparator

\section{引发剂、促长剂（与进展剂）的作用特点}\label{users__fpb__26_spring__pkusph_study__tmp__pdfs__prevention-medicine-past-exams__02-ux6bd2ux7406ux5b66ux5f80ux5e74ux9898ux8003ux70b9ux6574ux7406.md__ux5f15ux53d1ux5242ux4fc3ux957fux5242ux4e0eux8fdbux5c55ux5242ux7684ux4f5cux7528ux7279ux70b9}

\InfoLabel{考频} 10 次

\InfoLabel{对应小节} 引发剂、促长剂、进展剂；引发阶段/促长阶段/进展阶段特点 /《9 化学致癌作用》第55--65页

\InfoLabel{匹配依据}
``促长剂和引发剂的作用特点''``引发剂、促长剂的特点''反复出现，对应本小节各阶段化学物特点。

\ContentBand{知识点原文摘取}

引发剂（initiator）：本身或其代谢产物能与DNA作用，导致体细胞突变为引发细胞的化学物；引发剂本身有致癌性，大多数是致突变物，没有可检测的阈剂量，其作用的靶主要是原癌基因和肿瘤抑制基因；引发作用不可逆，并且是累积性的。促长剂（promoter）：作用于引发细胞增殖成为癌前病变或良性肿瘤的化学物；促长剂本身不能诱发肿瘤，常是非致突变物，通常具有阈剂量，其剂量反应关系呈S形曲线；持续给以促长剂是必需的，早期可逆、晚期不可逆，对生理因素（衰老、饮食和激素）调节敏感。进展剂（progressor）：作用于促长阶段的细胞转变成进展期的化学物。

\InfoLabel{知识来源} 《9 化学致癌作用》第55页；《9 化学致癌作用》第58页；《9 化学致癌作用》第60页；《9
化学致癌作用》第63页；《9 化学致癌作用》第64页；《9 化学致癌作用》第66页

\ContentBand{对应往年题}

\begin{pastquestion}

\QuestionLabel{【往年题1｜07级毒理口试题回忆.doc｜口试】}\\
促长剂和引发剂的作用特点（10号、18号、58号题）；引发剂、促长剂的作用特点（17号题）

\end{pastquestion}

\begin{pastquestion}

\QuestionLabel{【往年题2｜09级口试题目.docx｜口试】}\\
简述引发剂、促长剂的特点（10，张宝旭老师追问常见的引发剂和促长剂）；促长剂与引发剂的作用特点（58）；引发剂和促长剂的概念和特点（61，追问：促长剂s型曲线的原因）

\end{pastquestion}

\begin{pastquestion}

\QuestionLabel{【往年题3｜19级毒理期末口试回忆.xlsx｜口试】}\\
引发剂和促长剂的特点（序号78）；引发剂和促长剂的作用特点（无序号）

\end{pastquestion}

\begin{pastquestion}

\QuestionLabel{【往年题4｜17级毒理试题汇总.xlsx｜口试】}\\
引发剂和促长剂的作用特点（题号40、63）

\end{pastquestion}

\begin{pastquestion}

\QuestionLabel{【往年题5｜09级口试题目.docx｜口试备注】}\\
郝老师追问''促长剂s型曲线的原因''（61号题，呼应促长剂剂量反应关系呈S形曲线、有阈剂量）

\end{pastquestion}

\PointSeparator

\section{哺乳动物致癌试验（目的、动物选择、剂量设计、观察指标、结果判定/阳性标准）}\label{users__fpb__26_spring__pkusph_study__tmp__pdfs__prevention-medicine-past-exams__02-ux6bd2ux7406ux5b66ux5f80ux5e74ux9898ux8003ux70b9ux6574ux7406.md__ux54faux4e73ux52a8ux7269ux81f4ux764cux8bd5ux9a8cux76eeux7684ux52a8ux7269ux9009ux62e9ux5242ux91cfux8bbeux8ba1ux89c2ux5bdfux6307ux6807ux7ed3ux679cux5224ux5b9aux9633ux6027ux6807ux51c6}

\InfoLabel{考频} 9 次

\InfoLabel{对应小节} 化学物致癌性的鉴定---Long-term carcinogenicity study；致癌试验阳性的判断标准（WHO）/《9
化学致癌作用》第22--30页、第73页

\InfoLabel{匹配依据}
``哺乳动物致癌试验目的，动物要求，剂量设计原则，观察指标，结果判断''``致癌实验目的剂量动物选择结果判读''反复出现，对应本小节。

\ContentBand{知识点原文摘取}

Long-term carcinogenicity study
目的：确定受试化学物对实验动物的致癌性；剂量-反应关系；诱发肿瘤的靶器官。实验动物：常规选用大鼠和小鼠，也可用仓鼠；品系应选较敏感、自发肿瘤率低、生活力强及寿命较长的（NCI推荐Fischer344大鼠和B6C3F1小鼠）；同等数量雌雄两种性别；使用刚离乳的动物。剂量选择和动物数量：3个剂量组和1个对照组，每组至少雌雄各50只，希望出现第1例肿瘤时每组还有不少于25只；美国NCI推荐以最大耐受剂量（MTD）为高剂量。染毒途径：尽可能模拟人体可能的接触途径。试验期限：小鼠和仓鼠18个月、大鼠24个月等。肿瘤发生率（\%）=（实验结束时患肿瘤动物数/有效动物总数）×100\%。致癌试验阳性的判断标准（WHO）：对照组也出现的肿瘤试验组发生率增加；试验组发生对照组没有的肿瘤类型；试验组肿瘤发生早于对照组；与对照组比较试验组每只动物的平均肿瘤数增加。

\InfoLabel{知识来源} 《9 化学致癌作用》第23页；《9 化学致癌作用》第24页；《9 化学致癌作用》第25页；《9
化学致癌作用》第26页；《9 化学致癌作用》第27页；《9 化学致癌作用》第28页；《9 化学致癌作用》第73页

\ContentBand{对应往年题}

\begin{pastquestion}

\QuestionLabel{【往年题1｜09级口试题目.docx｜口试】}\\
哺乳动物致癌试验的动物选择（物种，品系，性别，数量，年龄），剂量设计原则，观测指标及结果评价（55）

\end{pastquestion}

\begin{pastquestion}

\QuestionLabel{【往年题2｜21级毒理期末口试回忆.xlsx｜口试】}\\
致癌试验的实验目的、动物选择原则、剂量设计原则和主要观察指标以及阳性的判定（序号11）；哺乳动物致癌实验目的，动物选择，主要观察指标，结果判定（序号65）

\end{pastquestion}

\begin{pastquestion}

\QuestionLabel{【往年题3｜19级毒理期末口试回忆.xlsx｜口试】}\\
哺乳动物致癌试验 所有（序号97）；哺乳动物致癌实验设计（序号65）

\end{pastquestion}

\begin{pastquestion}

\QuestionLabel{【往年题4｜17级毒理试题汇总.xlsx｜口试】}\\
哺乳动物致癌实验的设计（题号82）；致癌实验设计（题号34）；评价化学物致癌作用的实验有哪些（题号87，备注：哺乳动物短期致癌实验有哪些）

\end{pastquestion}

\begin{pastquestion}

\QuestionLabel{【往年题5｜15级毒理学考试题目汇总.xls｜口试】}\\
致癌试验的实验目的，实验动物，剂量设计，观察指标，结果分析（题号10）；哺乳动物致癌性试验目的，动物要求，剂量设计原则，观察指标，结果判断（题号30，备注：补充提问什么时候做化学物致癌实验）；致癌实验目的剂量动物选择结果判读（题号48）；哺乳动物致癌实验设计（目的，选动物，剂量设计，观察指标，结果判定）（题号49）

\end{pastquestion}

\begin{pastquestion}

\QuestionLabel{【往年题6｜14级毒理考题汇总.xls｜口试】}\\
哺乳动物致癌实验的目的，动物原则选择，剂量设计标准，观察指标，结果分析（一定要回答病理解剖诊断肿瘤发生率）（题号58）；哺乳动物致癌实验的目的、实验动物选择、剂量设计、观察指标、结果判断（题号49）；哺乳动物细胞致癌实验（题号30）

\end{pastquestion}

\PointSeparator

\section{IARC
基于证据权重的化学致癌物分类}\label{users__fpb__26_spring__pkusph_study__tmp__pdfs__prevention-medicine-past-exams__02-ux6bd2ux7406ux5b66ux5f80ux5e74ux9898ux8003ux70b9ux6574ux7406.md__iarc-ux57faux4e8eux8bc1ux636eux6743ux91cdux7684ux5316ux5b66ux81f4ux764cux7269ux5206ux7c7b}

\InfoLabel{考频} 8 次

\InfoLabel{对应小节} 化学致癌物的分类---IARC基于证据权重的化学致癌物分类 /《9 化学致癌作用》第16--20页

\InfoLabel{匹配依据}
``IARC化学致癌物分类''``基于IARC的化学致癌物分类''``IARC基于证据权重对化学致癌物的分类''反复出现，对应本小节。

\ContentBand{知识点原文摘取}

IARC根据对人类和实验动物致癌性资料，以及在实验系统和人类其它有关的资料进行综合评价，将环境因子和类别、混合物及暴露环境与人类癌症的关系分为：组1，对人类是致癌物（对人类致癌性证据充分者）；组2，对人类是很可能或可能致癌物，分为组2A（对人类很可能
probably 致癌，对人类致癌性证据有限、对实验动物致癌性证据充分）和组2B（对人类是可能 possible
致癌，对人类致癌性证据有限、对实验动物致癌性证据并不充分；或对人类证据不足、对实验动物证据充分）；组3，现有的证据不能对人类致癌性进行分类；组4，人类可能非致癌物。IARC
CLASSIFICATION依据：人群流行病学证据、动物实验证据、生物学机制的证据。

\InfoLabel{知识来源} 《9 化学致癌作用》第16页；《9 化学致癌作用》第18页；《9 化学致癌作用》第20页

\ContentBand{对应往年题}

\begin{pastquestion}

\QuestionLabel{【往年题1｜21级毒理期末口试回忆.xlsx｜口试】}\\
IARC基于证据权重对化学致癌物的分类（序号10，备注：追问IARC是什么，还有没有其他分类标准）；基于IARC的化学致癌物分类（序号74）

\end{pastquestion}

\begin{pastquestion}

\QuestionLabel{【往年题2｜17级毒理试题汇总.xlsx｜口试】}\\
IARC致癌物分类（题号83）；IARC根据证据权重对致癌物的分类（题号35）；基于证据的致癌物分类（题号93）

\end{pastquestion}

\begin{pastquestion}

\QuestionLabel{【往年题3｜15级毒理学考试题目汇总.xls｜口试】}\\
IARC化学致癌物分类（题号11）；IARC致癌物分类（题号50）；IRAC证据权重致癌物分级（题号No.11）；IARC基于证据权重的致癌物分类（题号50）；致癌物分级（题号31）

\end{pastquestion}

\begin{pastquestion}

\QuestionLabel{【往年题4｜14级毒理考题汇总.xls｜口试】}\\
依证据权重的致癌物分级（题号50）；IRAC（题号11）；IROC基于证据权重的化学致癌物分类（题号31，备注：每类致癌物举个例子）

\end{pastquestion}

\begin{pastquestion}

\QuestionLabel{【往年题5｜19级毒理期末口试回忆.xlsx｜口试】}\\
致癌物毒性分级（序号72）

\end{pastquestion}

\begin{pastquestion}

\QuestionLabel{【往年题6｜毒理学-来自课本的题库补充byYYS.docx｜口试】}\\
基于IARC的化学致癌物分类（居然问IARC是什么，还有哪些化学致癌物分类方法）

\end{pastquestion}

\PointSeparator

\section{癌基因（原癌基因）和抑癌基因（肿瘤抑制基因）的概念}\label{users__fpb__26_spring__pkusph_study__tmp__pdfs__prevention-medicine-past-exams__02-ux6bd2ux7406ux5b66ux5f80ux5e74ux9898ux8003ux70b9ux6574ux7406.md__ux764cux57faux56e0ux539fux764cux57faux56e0ux548cux6291ux764cux57faux56e0ux80bfux7624ux6291ux5236ux57faux56e0ux7684ux6982ux5ff5}

\InfoLabel{考频} 5 次

\InfoLabel{对应小节} 化学致癌机制---癌基因学说；癌基因；抑癌基因；TP53；致癌过程中涉及的两类基因比较 /《9
化学致癌作用》第43--48页

\InfoLabel{匹配依据} ``癌基因和抑癌基因概念''``癌基因和抑癌基因举例''反复出现，对应本小节。

\ContentBand{知识点原文摘取}

癌基因（Oncogene）：一类能引起细胞恶性转化及癌变的基因，是化学致癌物作用的主要靶分子。原癌基因（proto-oncogene）：本身是正常的细胞基因，是一类调控细胞生长分化的基因，在进化过程中高度保守，通常处于静止状态，在化学致癌物等作用下被激活而表达，促进细胞分裂增殖，引起细胞转化，最终形成肿瘤；活化的原癌基因才被称为癌基因。抑癌基因（Tumor
suppressor
gene）：作用方式与癌基因相反，在正常细胞中起抑制细胞增殖和促进分化的作用；在环境致癌因素作用下，抑癌基因发生纯合缺失或失活会引起细胞的恶性转化而致癌；抑癌基因必须一对等位基因丢失或突变后失活才能起作用，故又称隐性癌基因。野生型TP53为肿瘤抑制基因，参与G1/S关卡检查；突变型P53起癌基因作用。

\InfoLabel{知识来源} 《9 化学致癌作用》第44页；《9 化学致癌作用》第45页；《9 化学致癌作用》第46页；《9
化学致癌作用》第48页

\ContentBand{对应往年题}

\begin{pastquestion}

\QuestionLabel{【往年题1｜09级口试题目.docx｜口试】}\\
癌基因和抑癌基因的含义（42，备注：癌基因从原癌基因转化的机制）；癌基因和抑癌基因概念（50，备注：举出常见的癌基因和抑癌基因；为什么现在研究热点是癌基因而非抑癌基因）；简述癌基因和抑癌基因的概念（56）

\end{pastquestion}

\begin{pastquestion}

\QuestionLabel{【往年题2｜17级毒理试题汇总.xlsx｜口试】}\\
癌基因和抑癌基因概念（题号32，备注：郝老师死亡追问，癌基因和抑癌基因举例）

\end{pastquestion}

\begin{pastquestion}

\QuestionLabel{【往年题3｜07级毒理口试题回忆.doc｜口试】}\\
癌基因和抑癌基因（42号题）

\end{pastquestion}

\PointSeparator

\section{化学致癌物的分类（按致癌作用模式/机制；遗传毒性与非遗传毒性致癌物）及致癌机制学派}\label{users__fpb__26_spring__pkusph_study__tmp__pdfs__prevention-medicine-past-exams__02-ux6bd2ux7406ux5b66ux5f80ux5e74ux9898ux8003ux70b9ux6574ux7406.md__ux5316ux5b66ux81f4ux764cux7269ux7684ux5206ux7c7bux6309ux81f4ux764cux4f5cux7528ux6a21ux5f0fux673aux5236ux9057ux4f20ux6bd2ux6027ux4e0eux975eux9057ux4f20ux6bd2ux6027ux81f4ux764cux7269ux53caux81f4ux764cux673aux5236ux5b66ux6d3e}

\InfoLabel{考频} 3 次

\InfoLabel{对应小节} 根据致癌作用模式/机制分类；Part 3 化学致癌机制 /《9 化学致癌作用》第15页、第21页、第42页

\InfoLabel{匹配依据} 致癌物分类及''遗传机制学派/非遗传机制学派''为致癌作用机制题的组成部分，对应本小节。

\ContentBand{知识点原文摘取}

化学致癌物的分类：根据致癌作用模式/机制分类------遗传毒性致癌物（直接致癌物、间接致癌物、无机致癌物）；非遗传毒性致癌物（促长剂、激素调节剂、细胞毒剂、过氧化物酶体增殖剂、免疫抑制剂、固态物质如石棉等）；未分类（二噁烷、美沙吡林等）。化学致癌机制：遗传机制学派（亲电子剂学说，体细胞突变学说，癌基因学说，多阶段学说）；非遗传机制学派（细胞异常增生，免疫抑制，内分泌激素失衡，表观遗传改变等）。

\InfoLabel{知识来源} 《9 化学致癌作用》第15页；《9 化学致癌作用》第21页；《9 化学致癌作用》第42页

\ContentBand{对应往年题}

\begin{pastquestion}

\QuestionLabel{【往年题1｜07级毒理口试题回忆.doc｜口试】}\\
致癌的三阶段学说介绍（2号题，含致癌机制学派背景）

\end{pastquestion}

\begin{pastquestion}

\QuestionLabel{【往年题2｜09级口试题目.docx｜口试备注】}\\
张宝旭老师说致癌物与致突变物不是一回事儿（5号题备注，呼应遗传毒性/非遗传毒性致癌物分类）

\end{pastquestion}

\begin{pastquestion}

\QuestionLabel{【往年题3｜17级毒理试题汇总.xlsx｜口试】}\\
评价化学物致癌作用的实验有哪些（题号87，含致癌物筛选短期试验/体外细胞转化，呼应分类与鉴定）

\end{pastquestion}

\PointSeparator

\section{致癌物与致突变物的关系；助癌剂（cocarcinogen）的作用（久远题库补充）}\label{users__fpb__26_spring__pkusph_study__tmp__pdfs__prevention-medicine-past-exams__02-ux6bd2ux7406ux5b66ux5f80ux5e74ux9898ux8003ux70b9ux6574ux7406.md__ux81f4ux764cux7269ux4e0eux81f4ux7a81ux53d8ux7269ux7684ux5173ux7cfbux52a9ux764cux5242cocarcinogenux7684ux4f5cux7528ux4e45ux8fdcux9898ux5e93ux8865ux5145}

\InfoLabel{考频} 3 次

\InfoLabel{对应小节} 根据致癌作用模式/机制分类；化学致癌机制；多阶段致癌 /《9
化学致癌作用》第15页、第21页、第42页、第55--65页

\InfoLabel{匹配依据}
久远题库明确列出''致癌物是否都是致突变物''\,``助癌剂对癌症发生的作用''。现有整理已有致癌物分类、引发剂/促长剂与多阶段学说，但未把这两个辨析题单列。

\ContentBand{知识点原文摘取}

化学致癌物按致癌作用模式/机制可分为遗传毒性致癌物和非遗传毒性致癌物。遗传毒性致癌物包括直接致癌物、间接致癌物和无机致癌物；非遗传毒性致癌物包括促长剂、激素调节剂、细胞毒剂、过氧化物酶体增殖剂、免疫抑制剂和固态物质等。化学致癌机制包括遗传机制学派（亲电子剂学说、体细胞突变学说、癌基因学说、多阶段学说）和非遗传机制学派（细胞异常增生、免疫抑制、内分泌激素失衡、表观遗传改变等）。

补充辨析：大部分致癌物是致突变物，但致癌物不都等同于致突变物。许多遗传毒性致癌物可直接或经代谢活化后损伤
DNA，引起基因突变、染色体畸变或基因组改变；但非遗传毒性致癌物可不直接与 DNA
反应，而通过促长、激素调节、免疫抑制、细胞毒性、表观遗传改变等机制促进癌变。助癌剂（cocarcinogen）本身无致癌性，但在致癌物之前或同时应用可显著增加癌症发生，机制可涉及增加致癌物摄入、增加活化致癌物比例、耗尽竞争性亲核剂、抑制
DNA 修复或使 DNA 损伤固定为突变。

\InfoLabel{知识来源} 《9 化学致癌作用》第15页；《9 化学致癌作用》第21页；《9 化学致癌作用》第42页；《9
化学致癌作用》第55--65页；相对久远的总结《毒理题库-10年更新.doc》第36题、第41题

\ContentBand{对应往年题}

\begin{pastquestion}

\QuestionLabel{【往年题1｜相对久远的总结/毒理题库-10年更新.doc｜久远题库补充】}\\
致癌物是否都是致突变物。

\end{pastquestion}

\begin{pastquestion}

\QuestionLabel{【往年题2｜相对久远的总结/毒理题库-10年更新.doc｜久远题库补充】}\\
助癌剂对癌症发生的作用。

\end{pastquestion}

\begin{pastquestion}

\QuestionLabel{【往年题3｜09级口试题目.docx｜口试备注】}\\
张宝旭老师说致癌物与致突变物不是一回事儿（5号题备注，呼应遗传毒性/非遗传毒性致癌物分类）

\end{pastquestion}

\PointSeparator

\chapter{免疫毒理学}\label{users__fpb__26_spring__pkusph_study__tmp__pdfs__prevention-medicine-past-exams__02-ux6bd2ux7406ux5b66ux5f80ux5e74ux9898ux8003ux70b9ux6574ux7406.md__ux7b2cux5341ux7ae0-10-ux514dux75abux6bd2ux7406ux5b66}

\section{外源化学物对免疫系统作用的特点（灵敏性、复杂性、选择性、双向性）及免疫损伤的表现}\label{users__fpb__26_spring__pkusph_study__tmp__pdfs__prevention-medicine-past-exams__02-ux6bd2ux7406ux5b66ux5f80ux5e74ux9898ux8003ux70b9ux6574ux7406.md__ux5916ux6e90ux5316ux5b66ux7269ux5bf9ux514dux75abux7cfbux7edfux4f5cux7528ux7684ux7279ux70b9ux7075ux654fux6027ux590dux6742ux6027ux9009ux62e9ux6027ux53ccux5411ux6027ux53caux514dux75abux635fux4f24ux7684ux8868ux73b0}

\InfoLabel{考频} 10 次

\InfoLabel{对应小节} 外源化学物对免疫系统作用的特点 /《10 免疫毒理学》第10--15页

\InfoLabel{匹配依据}
``外源化学物对免疫系统作用的表现和特点''``毒作用对免疫系统的影响和特点（特点指灵敏、复杂、选择、双向）''反复出现，对应本小节四个特点。

\ContentBand{知识点原文摘取}

外源化学物对免疫系统作用的特点：反应的灵敏性；反应的复杂性；反应的选择性；反应的双向性（功能抑制、功能亢进）。化学物引起的免疫损伤的类型：免疫抑制；超敏反应（过敏反应、变态反应）；自身免疫反应。（灵敏性：硝酸锰土
2 mg/kg 抑制PFC，而其它毒性效应 NOAEL 为 20、200 mg/kg；选择性、复杂性、双向性见相应案例。）

\InfoLabel{知识来源} 《10 免疫毒理学》第10页；《10 免疫毒理学》第11页；《10 免疫毒理学》第12页；《10
免疫毒理学》第13页；《10 免疫毒理学》第14页；《10 免疫毒理学》第15页

\ContentBand{对应往年题}

\begin{pastquestion}

\QuestionLabel{【往年题1｜07级毒理口试题回忆.doc｜口试】}\\
外源化学物对机体免疫系统影响的表现和特点（29号题）；外源化学物对免疫系统的作用和特点（51号题）

\end{pastquestion}

\begin{pastquestion}

\QuestionLabel{【往年题2｜09级口试题目.docx｜口试】}\\
免疫毒性的表现和特点（34）；外源化学物对机体免疫系统损伤的表现和特点（需要与普通损伤相比较）（55）；外源化学物对免疫系统的作用及特点（题号未知）

\end{pastquestion}

\begin{pastquestion}

\QuestionLabel{【往年题3｜21级毒理期末口试回忆.xlsx｜口试】}\\
毒作用对免疫系统的影响和特点（序号30）；免疫系统损伤的表现和特点（序号74）

\end{pastquestion}

\begin{pastquestion}

\QuestionLabel{【往年题4｜19级毒理期末口试回忆.xlsx｜口试】}\\
简述毒物对免疫系统作用的表现和特点（特点指灵敏、复杂、选择、双向）（序号74）；外源化合物对免疫系统作用的表现和特点（序号30）

\end{pastquestion}

\begin{pastquestion}

\QuestionLabel{【往年题5｜17级毒理试题汇总.xlsx｜口试】}\\
外源性化学物质对免疫系统的影响和特点（题号68）

\end{pastquestion}

\begin{pastquestion}

\QuestionLabel{【往年题6｜14级毒理考题汇总.xls｜口试】}\\
免疫毒性的类型特点 评价方法（题号71）

\end{pastquestion}

\PointSeparator

\section{免疫毒理学的定义和研究内容}\label{users__fpb__26_spring__pkusph_study__tmp__pdfs__prevention-medicine-past-exams__02-ux6bd2ux7406ux5b66ux5f80ux5e74ux9898ux8003ux70b9ux6574ux7406.md__ux514dux75abux6bd2ux7406ux5b66ux7684ux5b9aux4e49ux548cux7814ux7a76ux5185ux5bb9}

\InfoLabel{考频} 8 次

\InfoLabel{对应小节} 概念---免疫毒理学；研究内容 /《10 免疫毒理学》第7页、第9页

\InfoLabel{匹配依据}
``免疫毒理学（英文）的定义及研究内容''``immunotoxicology的定义及研究内容''反复出现，对应本小节。

\ContentBand{知识点原文摘取}

免疫毒理学（immunotoxicology）：研究外源性的因素对人体和实验动物免疫系统产生的有害作用及其机制以及危害性风险评估的科学。研究内容：免疫毒性损伤识别（特征，剂量）；阐明免疫毒性及其机制（AOP，作用靶点）；免疫毒性风险评估（根据特定的损伤类型建立不同的模型）；完善和发展免疫毒性评价方法。

\InfoLabel{知识来源} 《10 免疫毒理学》第7页；《10 免疫毒理学》第9页

\ContentBand{对应往年题}

\begin{pastquestion}

\QuestionLabel{【往年题1｜21级毒理期末口试回忆.xlsx｜口试】}\\
免疫毒理学（英文）的定义+研究内容（序号74）；immunotoxicology的定义及研究内容（序号91）

\end{pastquestion}

\begin{pastquestion}

\QuestionLabel{【往年题2｜19级毒理期末口试回忆.xlsx｜口试】}\\
免疫毒理学的定义和研究内容（序号27、72）

\end{pastquestion}

\begin{pastquestion}

\QuestionLabel{【往年题3｜17级毒理试题汇总.xlsx｜口试】}\\
免疫毒理学的定义及研究内容（题号69）；免疫毒理学定义与研究内容（题号26）

\end{pastquestion}

\begin{pastquestion}

\QuestionLabel{【往年题4｜毒理学-来自课本的题库补充byYYS.docx｜口试】}\\
免疫毒理学（英文）的定义+研究内容

\end{pastquestion}

\PointSeparator

\section{免疫抑制的机制及评价（检测）方法}\label{users__fpb__26_spring__pkusph_study__tmp__pdfs__prevention-medicine-past-exams__02-ux6bd2ux7406ux5b66ux5f80ux5e74ux9898ux8003ux70b9ux6574ux7406.md__ux514dux75abux6291ux5236ux7684ux673aux5236ux53caux8bc4ux4ef7ux68c0ux6d4bux65b9ux6cd5}

\InfoLabel{考频} 9 次

\InfoLabel{对应小节} 免疫抑制的定义/机制；免疫抑制评价方法 /《10 免疫毒理学》第18--30页、第55--57页

\InfoLabel{匹配依据} ``免疫抑制的机制''``免疫抑制的评价方法''``不同免疫损伤的检测方法''反复出现，对应本小节。

\ContentBand{知识点原文摘取}

免疫抑制的定义：外源性的化学因素或者物理因素引起免疫系统的不同亚型细胞数量下降，细胞功能降低，导致机体免疫应答减弱或者消失。免疫抑制的机制：直接作用（功能改变------影响细胞成熟发育、增殖分化、干扰抗原递呈、干扰免疫细胞功能；结构改变；组成成分改变；直接的细胞毒作用）；间接作用（代谢活化、其它器官损伤的继发效应、对神经内分泌的影响、对营养和代谢的影响）。免疫抑制评价方法：免疫病理；免疫功能评价------体液免疫（PFC，血清溶血素）、细胞免疫（DTH，淋巴细胞增殖）、非特异免疫（巨噬细胞吞噬，NK细胞活性）；宿主抵抗力。

\InfoLabel{知识来源} 《10 免疫毒理学》第18页；《10 免疫毒理学》第20页；《10 免疫毒理学》第25页；《10
免疫毒理学》第28页；《10 免疫毒理学》第55页

\ContentBand{对应往年题}

\begin{pastquestion}

\QuestionLabel{【往年题1｜07级毒理口试题回忆.doc｜口试】}\\
不同种免疫损伤的检测方法（14号题）；不同免疫损伤的检测方法（39号题）；举例说明不同免疫损伤的检测方法（49号题）；外源化学物对免疫系统的作用和特点（51号题）

\end{pastquestion}

\begin{pastquestion}

\QuestionLabel{【往年题2｜09级口试题目.docx｜口试】}\\
不同免疫损伤的检测方法（14、39）；举例说明不同的免疫损伤的检查方法（56）

\end{pastquestion}

\begin{pastquestion}

\QuestionLabel{【往年题3｜21级毒理期末口试回忆.xlsx｜口试】}\\
免疫损伤的检测方法（序号32、76）；免疫抑制的机制（序号75、92）；免疫抑制的评价（序号54）；免疫抑制的评价方法（序号45、88）

\end{pastquestion}

\begin{pastquestion}

\QuestionLabel{【往年题4｜19级毒理期末口试回忆.xlsx｜口试】}\\
免疫损伤的检测方法（序号62、76）；免疫抑制的评价方法（序号88）

\end{pastquestion}

\begin{pastquestion}

\QuestionLabel{【往年题5｜17级毒理试题汇总.xlsx｜口试】}\\
免疫抑制的主要机制（题号42）；免疫抑制的评价指标（评价方法？）（题号44）；免疫损伤的检测方法（题号66）

\end{pastquestion}

\begin{pastquestion}

\QuestionLabel{【往年题6｜15级毒理学考试题目汇总.xls｜口试】}\\
免疫抑制的作用机制（题号57）；免疫抑制机制（题号29）；免疫毒性实验类型（题号27）；免疫的分类，特点，检测（题号55）

\end{pastquestion}

\begin{pastquestion}

\QuestionLabel{【往年题7｜14级毒理考题汇总.xls｜口试】}\\
免疫抑制机制（题号57）；免疫反应的类型，特点和评价方法（题号55）；免疫毒理，类型特点，反应（题号27）

\end{pastquestion}

\begin{pastquestion}

\QuestionLabel{【往年题8｜20级完整试题回忆.docx｜口试】}\\
举例介绍不同类型免疫损伤的检测方法（序号76的A）

\end{pastquestion}

\PointSeparator

\section{超敏反应的类型及特点（I、II、III、IV
型）}\label{users__fpb__26_spring__pkusph_study__tmp__pdfs__prevention-medicine-past-exams__02-ux6bd2ux7406ux5b66ux5f80ux5e74ux9898ux8003ux70b9ux6574ux7406.md__ux8d85ux654fux53cdux5e94ux7684ux7c7bux578bux53caux7279ux70b9iiiiiiiv-ux578b}

\InfoLabel{考频} 6 次

\InfoLabel{对应小节} 超敏反应（过敏反应、变态反应）；外源化学物引起超敏反应的类型和特点 /《10
免疫毒理学》第33--44页

\InfoLabel{匹配依据}
``超敏反应的类型及特点''``过敏反应的类型、特点''``化学物致敏的类型和机制''反复出现，对应本小节四型超敏反应。

\ContentBand{知识点原文摘取}

外源化学物引起超敏反应的类型和特点：I型（速发型）------IgE、肥大细胞、嗜碱性粒细胞，被变应原激活的肥大细胞脱颗粒释放组胺等介质，临床表现为过敏性休克、支气管哮喘、荨麻疹、鼻炎等；II型（细胞毒型）------IgG或IgM、补体、NK、K细胞，临床表现为溶血性贫血、粒细胞减少、血小板减少性紫癜、输血反应等；III型（免疫复合物型）------IgG、IgM、IgA、补体、中性粒细胞，抗原抗体复合物在组织中沉积激活补体导致组织损伤；IV型（迟发型）------T细胞和巨噬细胞，致敏T细胞激活后分泌淋巴因子介导细胞免疫反应，临床表现为接触性皮炎、湿疹、移植排斥反应等。超敏反应定义：机体针对无害抗原产生的无意义的强烈免疫应答；特点：特异质，难预测。

\InfoLabel{知识来源} 《10 免疫毒理学》第39页；《10 免疫毒理学》第40页；《10 免疫毒理学》第44页

\ContentBand{对应往年题}

\begin{pastquestion}

\QuestionLabel{【往年题1｜21级毒理期末口试回忆.xlsx｜口试】}\\
超敏反应的类型及特点（序号46）；超敏反应的分类及特点（序号100，备注：追问了II型超敏反应和自身免疫病的区别）

\end{pastquestion}

\begin{pastquestion}

\QuestionLabel{【往年题2｜19级毒理期末口试回忆.xlsx｜口试】}\\
超敏反应的类型和特点（序号46）

\end{pastquestion}

\begin{pastquestion}

\QuestionLabel{【往年题3｜17级毒理试题汇总.xlsx｜口试】}\\
过敏反应的类型、特点（题号50）；超敏反应类型和机制（题号28）

\end{pastquestion}

\begin{pastquestion}

\QuestionLabel{【往年题4｜15级毒理学考试题目汇总.xls｜口试】}\\
致敏反应的类型和原理（题号28）；化学物致敏的类型和机制（题号28）；致敏反应的类型和机制（题号28）

\end{pastquestion}

\begin{pastquestion}

\QuestionLabel{【往年题5｜14级毒理考题汇总.xls｜口试】}\\
化学致敏反应的类型及作用机制（题号56）；超敏反应类型和机制（题号28）

\end{pastquestion}

\begin{pastquestion}

\QuestionLabel{【往年题6｜20级完整试题回忆.docx｜单项选择】}\\
外来化合物对免疫系统作用特点，下述哪项不正确（）A反应的复杂性 B反应的选择性 C反应的均衡性 D反应的灵敏性
E反应的双重性（第15题）

\end{pastquestion}

\PointSeparator

\section{自身免疫反应的定义、后果及产生机制}\label{users__fpb__26_spring__pkusph_study__tmp__pdfs__prevention-medicine-past-exams__02-ux6bd2ux7406ux5b66ux5f80ux5e74ux9898ux8003ux70b9ux6574ux7406.md__ux81eaux8eabux514dux75abux53cdux5e94ux7684ux5b9aux4e49ux540eux679cux53caux4ea7ux751fux673aux5236}

\InfoLabel{考频} 3 次

\InfoLabel{对应小节} 自身免疫反应；自身免疫产生的机制 /《10 免疫毒理学》第47--50页

\InfoLabel{匹配依据}
``自身免疫损伤的定义及后果''``化学物引起自身免疫损伤的机制''对应本小节（多并入''免疫损伤类型/机制''题）。

\ContentBand{知识点原文摘取}

自身免疫反应：化学物接触后，诱导机体对自身抗原产生免疫应答，由于自身识别障碍，免疫球蛋白及T细胞受体与自身抗原发生反应，导致组织损伤。后果：出现局部或者全身性组织损伤，表现为器官的炎症或者全身性的炎症损伤（系统性红斑狼疮）。自身免疫产生的机制：干扰中枢耐受；改变外围耐受（IL-10,
TGF-β；组织损伤、细胞应激）；改变抗原结构（化学物氧化自身蛋白的巯基）；遗传易感性；交叉反应（化学物抗原与自身抗原具有相同的抗原决定簇）。

\InfoLabel{知识来源} 《10 免疫毒理学》第47页；《10 免疫毒理学》第49页；《10 免疫毒理学》第50页

\ContentBand{对应往年题}

\begin{pastquestion}

\QuestionLabel{【往年题1｜07级毒理口试题回忆.doc｜口试】}\\
外源性化学物对机体免疫系统影响的表现和特点（含免疫损伤类型，29号题，追问免疫学试验：细胞免疫、体液免疫、非特异性免疫、超敏反应）

\end{pastquestion}

\begin{pastquestion}

\QuestionLabel{【往年题2｜09级口试题目.docx｜口试】}\\
免疫学毒性有哪些（29，含自身免疫反应）

\end{pastquestion}

\begin{pastquestion}

\QuestionLabel{【往年题3｜21级毒理期末口试回忆.xlsx｜口试】}\\
化学物引起的免疫损伤的类型（含自身免疫反应，序号30/74等并入免疫损伤类型）

\end{pastquestion}

\PointSeparator

\chapter{化学致畸作用、生殖毒性、发育毒性2021 / 化学致畸作用和发育毒性2022}\label{users__fpb__26_spring__pkusph_study__tmp__pdfs__prevention-medicine-past-exams__02-ux6bd2ux7406ux5b66ux5f80ux5e74ux9898ux8003ux70b9ux6574ux7406.md__ux7b2cux5341ux4e00ux7ae0-11-ux5316ux5b66ux81f4ux7578ux4f5cux7528ux751fux6b96ux6bd2ux6027ux53d1ux80b2ux6bd2ux6027202111-ux5316ux5b66ux81f4ux7578ux4f5cux7528ux548cux53d1ux80b2ux6bd2ux60272022}

\section{生殖毒性和发育毒性的定义与主要表现（含发育毒性观察终点）}\label{users__fpb__26_spring__pkusph_study__tmp__pdfs__prevention-medicine-past-exams__02-ux6bd2ux7406ux5b66ux5f80ux5e74ux9898ux8003ux70b9ux6574ux7406.md__ux751fux6b96ux6bd2ux6027ux548cux53d1ux80b2ux6bd2ux6027ux7684ux5b9aux4e49ux4e0eux4e3bux8981ux8868ux73b0ux542bux53d1ux80b2ux6bd2ux6027ux89c2ux5bdfux7ec8ux70b9}

\InfoLabel{考频} 12 次

\InfoLabel{对应小节} 生殖毒性与发育毒性；发育毒性的主要表现 /《11
化学致畸作用、生殖毒性、发育毒性2021》第10--13页（并见《11 化学致畸作用和发育毒性2022》第4--5页
发育毒性观察终点）

\InfoLabel{匹配依据} ``生殖毒性和发育毒性的定义与主要表现''``development
toxicity的概念及主要表现''``发育毒性的概念+表现''反复出现，对应本小节。

\ContentBand{知识点原文摘取}

生殖毒性（Reproductive
Toxicity）：外源化学物对生殖过程的不良影响，主要指对生殖细胞的发生、卵细胞受精、胚胎形成、妊娠、分娩和哺乳过程的损害作用；着重研究外源化学物对亲代的生殖功能的有害作用。发育毒性（Developmental
Toxicity）：外源化学物对发育过程的不良影响，包括对胚胎发育、胎仔发育及出生幼仔发育的损害作用；指出生前后接触有害因素，子代个体发育为成体之前诱发的任何有害影响；着重研究外源化学物对子代的发育过程的有害作用。发育毒性的主要表现：发育生物体死亡（death
of the developing organism）；生长改变（altered growth，生长迟缓）；结构异常（structural
abnormality，外观、内脏、骨骼畸形）；功能缺陷（functional deficiency，生理、生化、免疫、行为、智力）。

【翻译】生殖毒性是外源化学物对生殖过程造成的不良影响。发育毒性是外源化学物对发育过程造成的不良影响。发育毒性的主要表现包括发育中个体死亡、生长改变、结构异常和功能缺陷。

\InfoLabel{知识来源} 《11 化学致畸作用、生殖毒性、发育毒性2021》第10页；《11
化学致畸作用、生殖毒性、发育毒性2021》第11页；《11 化学致畸作用、生殖毒性、发育毒性2021》第13页；《11
化学致畸作用和发育毒性2022》第4页；《11 化学致畸作用和发育毒性2022》第5页

\ContentBand{对应往年题}

\begin{pastquestion}

\QuestionLabel{【往年题1｜07级毒理口试题回忆.doc｜口试】}\\
发育毒性的定义及表现（46号题）；developmental toxicity内容和表现（60号题）；生殖毒性概念及表现（21号题）

\end{pastquestion}

\begin{pastquestion}

\QuestionLabel{【往年题2｜09级口试题目.docx｜口试】}\\
发育毒性概念表现（19、24）；生殖毒性概念及表现（21）；developmental toxicity定义及表现（60）

\end{pastquestion}

\begin{pastquestion}

\QuestionLabel{【往年题3｜21级毒理期末口试回忆.xlsx｜口试】}\\
发育毒性的概念和主要表现（序号3）；生殖毒性和发育毒性的定义与主要表现（序号44、78、88）；生殖毒性及其表现（序号49）

\end{pastquestion}

\begin{pastquestion}

\QuestionLabel{【往年题4｜19级毒理期末口试回忆.xlsx｜口试】}\\
发育毒性的概念+表现（序号49）；生殖毒性与发育毒性的定义与主要表现（序号88）

\end{pastquestion}

\begin{pastquestion}

\QuestionLabel{【往年题5｜17级毒理试题汇总.xlsx｜口试】}\\
生殖毒性和发育毒性（英文)的定义和主要表现（题号44）；development toxicity的概念及主要表现（题号8）

\end{pastquestion}

\begin{pastquestion}

\QuestionLabel{【往年题6｜15级毒理学考试题目汇总.xls｜口试】}\\
生殖毒性发育毒性概念，发育毒性表现，生殖毒性发育毒性特点（题号61）；生殖发育毒性概念、发育毒性表现，特点（题号13）；生殖和发育毒性的概念、表现、特点（题号33）

\end{pastquestion}

\begin{pastquestion}

\QuestionLabel{【往年题7｜14级毒理考题汇总.xls｜口试】}\\
生殖发育毒性 定义
特点（题号52）；生殖和发育毒性概念，特点，发育毒性的表现（题号33）；生殖发育毒性概念，表现，特点（题号13）

\end{pastquestion}

\begin{pastquestion}

\QuestionLabel{【往年题8｜20级完整试题回忆.docx｜单项选择】}\\
对于不同类型发育毒性毒物的剂量-反应关系模式，说法正确的是（第17题）

\end{pastquestion}

\PointSeparator

\section{三段生殖毒性试验的要点、内容及异同}\label{users__fpb__26_spring__pkusph_study__tmp__pdfs__prevention-medicine-past-exams__02-ux6bd2ux7406ux5b66ux5f80ux5e74ux9898ux8003ux70b9ux6574ux7406.md__ux4e09ux6bb5ux751fux6b96ux6bd2ux6027ux8bd5ux9a8cux7684ux8981ux70b9ux5185ux5bb9ux53caux5f02ux540c}

\InfoLabel{考频} 12 次

\InfoLabel{对应小节} 生殖与发育毒性试验---三段生殖毒性试验 /《11
化学致畸作用、生殖毒性、发育毒性2021》第56--62页（并见《11 化学致畸作用和发育毒性2022》第15页）

\InfoLabel{匹配依据}
``三段生殖毒性试验的要点''``三段致畸试验的要点''``三段生殖试验的异同点（试验设计课上ppt的表格）''反复出现，对应本小节三段试验目的/动物/染毒期/处死时间/观察项目。

\ContentBand{知识点原文摘取}

三段生殖毒性试验：生育力和早期胚胎发育毒性试验（一般生殖毒性试验）；胚体-胎体毒性试验（致畸试验）；出生前和出生后发育毒性试验（围生期毒性试验）。主要为发育毒性检测，属于化学物毒理学安全性评价试验四个阶段中的第二阶段。一般生殖毒性试验：评价配子发育及着床前期接触化学物对生育力及早期胚胎发育的影响，雄鼠交配前28天、雌鼠交配前14天到着床染毒。致畸试验（胚体-胎体毒性试验）：评价致畸敏感期接触化学物对妊娠、胎体死亡、生长及结构异常的影响，通常两种动物（大鼠和家兔），胎鼠器官形成期（大鼠孕6-15天，家兔孕6-18天）染毒，分娩前一天处死。围生期毒性试验：评价母体自着床至断乳期间接触化学物对母鼠妊娠、仔鼠生长发育、形态结构及神经行为功能的影响。

\InfoLabel{知识来源} 《11 化学致畸作用、生殖毒性、发育毒性2021》第56页；《11
化学致畸作用、生殖毒性、发育毒性2021》第59页；《11 化学致畸作用、生殖毒性、发育毒性2021》第60页；《11
化学致畸作用、生殖毒性、发育毒性2021》第62页；《11 化学致畸作用和发育毒性2022》第15页

\ContentBand{对应往年题}

\begin{pastquestion}

\QuestionLabel{【往年题1｜07级毒理口试题回忆.doc｜口试】}\\
三段发育毒性实验特点（未知题号）；三段式实验的注意事项（15号题）；三段生殖毒性实验的要点（58号题）

\end{pastquestion}

\begin{pastquestion}

\QuestionLabel{【往年题2｜09级口试题目.docx｜口试】}\\
三段生殖试验的要点 内容很多（58）；三阶段生殖发育毒性（题号未知）

\end{pastquestion}

\begin{pastquestion}

\QuestionLabel{【往年题3｜21级毒理期末口试回忆.xlsx｜口试】}\\
三阶段毒性试验（序号37）；三段生殖实验的要点（序号46）

\end{pastquestion}

\begin{pastquestion}

\QuestionLabel{【往年题4｜19级毒理期末口试回忆.xlsx｜口试】}\\
简述三段生殖毒性实验要点（序号37，备注：追问三段生殖毒性实验适用于什么------药品）；简要介绍生殖毒理学实验设计原则（序号24）

\end{pastquestion}

\begin{pastquestion}

\QuestionLabel{【往年题5｜17级毒理试题汇总.xlsx｜口试】}\\
三段致畸试验的要点（题号53）；三段生殖毒性试验要点（题号5）

\end{pastquestion}

\begin{pastquestion}

\QuestionLabel{【往年题6｜15级毒理学考试题目汇总.xls｜口试】}\\
三段生殖毒性实验（题号34）；三段生殖毒性试验全部要点（题号53）；三段生殖试验异同点表格（题号34）

\end{pastquestion}

\begin{pastquestion}

\QuestionLabel{【往年题7｜14级毒理考题汇总.xls｜口试】}\\
三段生殖毒性试验的各方面简述（题号14）；三段生殖试验的异同点，也是试验设计课上ppt的表格（题号53）；三段生殖毒性实验，目的，染毒时间处死时间，观察项目和动物品系（题号62）

\end{pastquestion}

\PointSeparator

\section{致畸作用的影响因素（含发育毒性的影响因素）}\label{users__fpb__26_spring__pkusph_study__tmp__pdfs__prevention-medicine-past-exams__02-ux6bd2ux7406ux5b66ux5f80ux5e74ux9898ux8003ux70b9ux6574ux7406.md__ux81f4ux7578ux4f5cux7528ux7684ux5f71ux54cdux56e0ux7d20ux542bux53d1ux80b2ux6bd2ux6027ux7684ux5f71ux54cdux56e0ux7d20}

\InfoLabel{考频} 10 次

\InfoLabel{对应小节} 致畸作用的影响因素 /《11 化学致畸作用、生殖毒性、发育毒性2021》第19--34页（并见《11
化学致畸作用和发育毒性2022》第7--10页）

\InfoLabel{匹配依据}
``影响致畸作用的因素''``致畸作用的影响因素，其对致畸试验有哪些指导意义''``发育毒性的影响因素''反复出现，对应本小节。

\ContentBand{知识点原文摘取}

致畸作用的影响因素：接触致畸原的时间（致畸试验的染毒时间必须安排在器官形成期，即致畸作用关键期 critical
period）；接触致畸原的剂量（致畸作用具有阈剂量）；物种差异（某一化学物在某一物种中是否有致畸性取决于化学物能否在这一物种体内转化成具有致畸活性的产物）；母体因素和母体毒性（母体因素：遗传背景、营养状态、疾病状态、应激、年龄、产次、代谢状态、其他暴露；母体毒性：胎盘毒性、母体效应）；其他因素（化学物的理化性质、染毒途径等；父源性因素）。

\InfoLabel{知识来源} 《11 化学致畸作用、生殖毒性、发育毒性2021》第19页；《11
化学致畸作用、生殖毒性、发育毒性2021》第24页；《11 化学致畸作用、生殖毒性、发育毒性2021》第29页；《11
化学致畸作用、生殖毒性、发育毒性2021》第33页；《11 化学致畸作用、生殖毒性、发育毒性2021》第34页；《11
化学致畸作用和发育毒性2022》第7页

\ContentBand{对应往年题}

\begin{pastquestion}

\QuestionLabel{【往年题1｜07级毒理口试题回忆.doc｜口试】}\\
影响致畸作用的因素（35号题，备注：还要补充染毒途径）

\end{pastquestion}

\begin{pastquestion}

\QuestionLabel{【往年题2｜21级毒理期末口试回忆.xlsx｜口试】}\\
影响致畸作用的主要因素（序号50）

\end{pastquestion}

\begin{pastquestion}

\QuestionLabel{【往年题3｜19级毒理期末口试回忆.xlsx｜口试】}\\
发育毒性的影响因素（序号51、1，备注：器官形成期暴露会导致结构畸形，之后暴露则主要导致功能畸形）；发育毒性的影响因素（序号79）

\end{pastquestion}

\begin{pastquestion}

\QuestionLabel{【往年题4｜17级毒理试题汇总.xlsx｜口试】}\\
影响致畸作用的因素（题号91）；致畸作用的影响因素（题号35，备注：以此谈实验设计注意问题）；发育毒性的影响因素（题号43）

\end{pastquestion}

\begin{pastquestion}

\QuestionLabel{【往年题5｜15级毒理学考试题目汇总.xls｜口试】}\\
致畸实验影响因素和对实验设计的指导（题号54）；致畸作用的影响因素，并以此阐明实验设计的要点（题号15）；致畸影响因素，基于影响因素的实验设计（题号35）

\end{pastquestion}

\begin{pastquestion}

\QuestionLabel{【往年题6｜14级毒理考题汇总.xls｜口试】}\\
致畸作用的影响因素，其对致畸试验有哪些指导意义（题号35）；致畸作用的影响因素及致畸实验的设计（无题号）；致畸试验的影响因素，并结合影响因素谈实验设计（题号54）

\end{pastquestion}

\PointSeparator

\section{致畸试验的设计要点}\label{users__fpb__26_spring__pkusph_study__tmp__pdfs__prevention-medicine-past-exams__02-ux6bd2ux7406ux5b66ux5f80ux5e74ux9898ux8003ux70b9ux6574ux7406.md__ux81f4ux7578ux8bd5ux9a8cux7684ux8bbeux8ba1ux8981ux70b9}

\InfoLabel{考频} 10 次

\InfoLabel{对应小节} 致畸试验 /《11 化学致畸作用、生殖毒性、发育毒性2021》第36--37页（并见《11
化学致畸作用和发育毒性2022》第14页）

\InfoLabel{匹配依据}
``致畸试验设计要点''``致畸试验的要点''``畸形试验要点''反复出现，对应本小节''一个关键期、两个物种、三个对照组、四个剂量、五个结果观察''。

\ContentBand{知识点原文摘取}

致畸试验：试验动物------大鼠、小鼠、家兔；剂量------高剂量应在母体中产生轻度的毒性（LD50的1/5-1/3），中剂量应引起最小的毒作用（LOAEL），低剂量NOAEL或LD50的1/100-1/30；暴露途径与染毒时机------应与人的暴露途径相同，器官形成期；对照组------阴性对照（溶剂或赋形剂）、阳性对照（乙酰水杨酸、Vit-A等）；结果观察------母体观察、胎仔检查（一般检查、外观检查、内脏检查、骨骼检查）。试验设计概括：一个关键期（致畸作用关键期）；两个物种（啮齿类、非啮齿类哺乳动物）；三个对照组（阴性、空白、阳性）；四个剂量（零、低、中、高）；五个结果观察（胎仔一般、外观、内脏、骨骼，母体）。

\InfoLabel{知识来源} 《11 化学致畸作用、生殖毒性、发育毒性2021》第36页；《11
化学致畸作用、生殖毒性、发育毒性2021》第37页；《11 化学致畸作用和发育毒性2022》第14页

\ContentBand{对应往年题}

\begin{pastquestion}

\QuestionLabel{【往年题1｜07级毒理口试题回忆.doc｜口试】}\\
致畸试验的要点（2号题）；致畸试验的设计（61号题）；畸形试验要点（37号题）

\end{pastquestion}

\begin{pastquestion}

\QuestionLabel{【往年题2｜09级口试题目.docx｜口试】}\\
致畸试验的要点（2、61，追问大鼠/小鼠致畸敏感期时间）；致畸试验影响因素（16）

\end{pastquestion}

\begin{pastquestion}

\QuestionLabel{【往年题3｜21级毒理期末口试回忆.xlsx｜口试】}\\
致畸试验设计要点（序号24、52，追问大小鼠器官形成期是什么时间、结果观察什么）；致畸试验的要点（序号86，备注：别忘了致畸指数）

\end{pastquestion}

\begin{pastquestion}

\QuestionLabel{【往年题4｜19级毒理期末口试回忆.xlsx｜口试】}\\
致畸试验设计要点（序号忘了/姚碧云）；致畸性实验设计（序号52）

\end{pastquestion}

\begin{pastquestion}

\QuestionLabel{【往年题5｜17级毒理试题汇总.xlsx｜口试】}\\
致畸实验设计要点（题号12）

\end{pastquestion}

\begin{pastquestion}

\QuestionLabel{【往年题6｜09级口试题目.docx｜口试】}\\
三阶段致畸试验的要点（题号未知，含致畸试验设计）

\end{pastquestion}

\PointSeparator

\section{致畸作用关键期（critical
period）/靶窗的概念和毒理学意义}\label{users__fpb__26_spring__pkusph_study__tmp__pdfs__prevention-medicine-past-exams__02-ux6bd2ux7406ux5b66ux5f80ux5e74ux9898ux8003ux70b9ux6574ux7406.md__ux81f4ux7578ux4f5cux7528ux5173ux952eux671fcritical-periodux9776ux7a97ux7684ux6982ux5ff5ux548cux6bd2ux7406ux5b66ux610fux4e49}

\InfoLabel{考频} 6 次

\InfoLabel{对应小节} 接触致畸原的时间---关键期/靶窗 /《11
化学致畸作用、生殖毒性、发育毒性2021》第20--23页（并见《11 化学致畸作用和发育毒性2022》第7页）

\InfoLabel{匹配依据} ``致畸作用的critical
period的概念、毒理学意义''``关键期定义，毒理学的意义''反复出现，对应本小节。

\ContentBand{知识点原文摘取}

关键期（critical
period）：胚胎发育过程中，易于发生形态结构异常的整个时期。通常略长于胚胎的器官形成期（organogenesis
period）。靶窗（Target
window）：不同器官对致畸作用的特殊敏感期。致畸试验的染毒时间必须安排在器官形成期，即致畸作用关键期。不同动物胚胎发育时间（妊娠天数）：人着床8-13、器官形成21-56；大鼠器官形成6-15；小鼠器官形成6-15。

\InfoLabel{知识来源} 《11 化学致畸作用、生殖毒性、发育毒性2021》第20页；《11
化学致畸作用、生殖毒性、发育毒性2021》第21页；《11 化学致畸作用、生殖毒性、发育毒性2021》第22页；《11
化学致畸作用和发育毒性2022》第7页

\ContentBand{对应往年题}

\begin{pastquestion}

\QuestionLabel{【往年题1｜07级毒理口试题回忆.doc｜口试】}\\
致畸作用的critical period的概念、毒理学意义（30号、31号题）；致畸作用的critical
period的含义，毒理学意义（31号题）

\end{pastquestion}

\begin{pastquestion}

\QuestionLabel{【往年题2｜09级口试题目.docx｜口试】}\\
致畸试验的critical period的定义和毒理学意义（31）；致畸作用关键期的概念及意义（53）；致畸Critical
period概念，毒理学意义，小鼠、大鼠的具体时间段（41）

\end{pastquestion}

\begin{pastquestion}

\QuestionLabel{【往年题3｜19级毒理期末口试回忆.xlsx｜口试】}\\
关键期定义，毒理学的意义（序号5）

\end{pastquestion}

\PointSeparator

\section{致畸作用的剂量反应关系及其毒理学意义（含三种剂量反应模式）}\label{users__fpb__26_spring__pkusph_study__tmp__pdfs__prevention-medicine-past-exams__02-ux6bd2ux7406ux5b66ux5f80ux5e74ux9898ux8003ux70b9ux6574ux7406.md__ux81f4ux7578ux4f5cux7528ux7684ux5242ux91cfux53cdux5e94ux5173ux7cfbux53caux5176ux6bd2ux7406ux5b66ux610fux4e49ux542bux4e09ux79cdux5242ux91cfux53cdux5e94ux6a21ux5f0f}

\InfoLabel{考频} 6 次

\InfoLabel{对应小节} 接触致畸原的剂量；致畸作用的剂量反应关系 /《11
化学致畸作用、生殖毒性、发育毒性2021》第25--26页、第35页

\InfoLabel{匹配依据}
``外源化学物致畸作用的剂量反应关系特点及意义''``简述致畸曲线及其毒理学意义''``致畸效应的剂量反应关系和毒理学意义''反复出现，对应本小节。

\ContentBand{知识点原文摘取}

接触致畸原的剂量：剂量增加与毒作用表现的关系------随剂量增加依次出现无作用带、致畸带、胚胎死亡带、母体死亡带；致畸作用具有阈剂量。致畸试验剂量设计应适中：过低不足以显示致畸性；过高引起大量胚胎死亡、畸胎率反而降低，或母体毒性过强不能判断。致畸作用的剂量反应关系三种模式（一次染毒，剂量水平不导致母体严重毒性）：不引起胚胎死亡导致整窝胚胎畸形（高度致畸性）；一窝中正常胎、结构畸形、生长迟缓、死亡共存（常见）；只有胚胎死亡和生长迟缓没有畸形（具有胚胎毒性但非致畸性）。意义：了解机制。

\InfoLabel{知识来源} 《11 化学致畸作用、生殖毒性、发育毒性2021》第25页；《11
化学致畸作用、生殖毒性、发育毒性2021》第26页；《11 化学致畸作用、生殖毒性、发育毒性2021》第35页

\ContentBand{对应往年题}

\begin{pastquestion}

\QuestionLabel{【往年题1｜07级毒理口试题回忆.doc｜口试】}\\
简述致畸曲线及其毒理学意义（意义是了解机制）（33号题）；外源化学物致畸作用的剂量反应关系特点及意义（57号题）

\end{pastquestion}

\begin{pastquestion}

\QuestionLabel{【往年题2｜09级口试题目.docx｜口试】}\\
致畸效应的剂量反应关系和毒理学意义（33）；致畸剂量反应曲线和毒理学意义（57）

\end{pastquestion}

\begin{pastquestion}

\QuestionLabel{【往年题3｜17级毒理试题汇总.xlsx｜口试】}\\
致畸作用的剂量反应关系和毒理学意义（题号98）

\end{pastquestion}

\begin{pastquestion}

\QuestionLabel{【往年题4｜20级完整试题回忆.docx｜单项选择】}\\
对于不同类型发育毒性毒物的剂量-反应关系模式，说法正确的是（第17题）

\end{pastquestion}

\PointSeparator

\section{致畸物、致畸性的概念与致畸指数；生殖与发育毒性的特点；雄/雌性生殖毒性靶点}\label{users__fpb__26_spring__pkusph_study__tmp__pdfs__prevention-medicine-past-exams__02-ux6bd2ux7406ux5b66ux5f80ux5e74ux9898ux8003ux70b9ux6574ux7406.md__ux81f4ux7578ux7269ux81f4ux7578ux6027ux7684ux6982ux5ff5ux4e0eux81f4ux7578ux6307ux6570ux751fux6b96ux4e0eux53d1ux80b2ux6bd2ux6027ux7684ux7279ux70b9ux96c4ux96ccux6027ux751fux6b96ux6bd2ux6027ux9776ux70b9}

\InfoLabel{考频} 5 次

\InfoLabel{对应小节}
致畸作用---致畸物/致畸性；致畸试验结果评定---致畸指数；生殖与发育毒性的特点；雄/雌性生殖毒性的作用靶点 /《11
化学致畸作用、生殖毒性、发育毒性2021》第14、16、46页（并见《11 化学致畸作用和发育毒性2022》第10、13、14页）

\InfoLabel{匹配依据}
``致畸指数的概念及毒理学意义''``雄性生殖毒性的主要靶点''``生殖与发育毒性的特点''对应本小节。

\ContentBand{知识点原文摘取}

致畸物（Teratogen）：在出生前暴露可诱发永久性结构异常的物质；能引起畸形的环境因子。致畸性（Teratogenicity）：外源化学物作用于妊娠母体，干扰胚胎及胎儿的正常发育过程，使胎儿出现形态结构及功能异常的性质。致畸指数=化合物对雌性动物LD50/最小致畸剂量；\textless10不具致畸作用，10-100有致畸作用，\textgreater100强致畸作用。生殖与发育毒性的特点：生殖过程对某些化学物的毒作用更为敏感（在成体系统毒性的NOAEL，胚胎即可受到影响）；损害作用不仅表现在暴露化学物质的机体本身，还可影响其后代。雄性生殖毒性作用靶点：Leydig细胞、Sertoli细胞、生精上皮、附睾上皮；雌性生殖毒性作用靶点：卵巢、输卵管、子宫、激素分泌。

\InfoLabel{知识来源} 《11 化学致畸作用、生殖毒性、发育毒性2021》第14页；《11
化学致畸作用、生殖毒性、发育毒性2021》第16页；《11 化学致畸作用、生殖毒性、发育毒性2021》第46页；《11
化学致畸作用和发育毒性2022》第10页；《11 化学致畸作用和发育毒性2022》第13页；《11
化学致畸作用和发育毒性2022》第14页

\ContentBand{对应往年题}

\begin{pastquestion}

\QuestionLabel{【往年题1｜19级毒理期末口试回忆.xlsx｜口试备注】}\\
别忘了致畸指数朋友们（序号86 备注，呼应致畸指数概念及意义）

\end{pastquestion}

\begin{pastquestion}

\QuestionLabel{【往年题2｜21级毒理期末口试回忆.xlsx｜口试】}\\
生殖毒性和发育毒性的定义与主要表现（序号44，追问生殖毒性的表现、贮存库的危害；序号49
追问生殖毒性的起点与终点）

\end{pastquestion}

\begin{pastquestion}

\QuestionLabel{【往年题3｜15级毒理学考试题目汇总.xls｜口试】}\\
生殖毒性发育毒性概念，发育毒性表现，生殖毒性发育毒性特点（题号61）

\end{pastquestion}

\begin{pastquestion}

\QuestionLabel{【往年题4｜14级毒理考题汇总.xls｜口试】}\\
生殖发育毒性 定义 特点（题号52）；生殖发育毒性概念，表现，特点（题号13）

\end{pastquestion}

\begin{pastquestion}

\QuestionLabel{【往年题5｜11 化学致畸作用和发育毒性2022（课件思考题，第11页）｜思考题】}\\
致畸指数的概念及毒理学意义（与往年口试同型，纳入复习）

\end{pastquestion}

\PointSeparator

\section{致畸试验/发育毒物的体外预筛（替代）试验}\label{users__fpb__26_spring__pkusph_study__tmp__pdfs__prevention-medicine-past-exams__02-ux6bd2ux7406ux5b66ux5f80ux5e74ux9898ux8003ux70b9ux6574ux7406.md__ux81f4ux7578ux8bd5ux9a8cux53d1ux80b2ux6bd2ux7269ux7684ux4f53ux5916ux9884ux7b5bux66ffux4ee3ux8bd5ux9a8c}

\InfoLabel{考频} 2 次

\InfoLabel{对应小节} 致畸试验替代方法（EST、MM、WEC）/《11
化学致畸作用、生殖毒性、发育毒性2021》第49--54页（并见《11 化学致畸作用和发育毒性2022》第17--19页）

\InfoLabel{匹配依据}
课件思考题与往年''体外替代方法能否取代传统致畸试验''对应本小节胚胎干细胞试验、细胞微团试验、全胚胎培养。

\ContentBand{知识点原文摘取}

致畸试验替代方法：胚胎干细胞试验（embryonic stem cell test,
EST，使用小鼠胚胎干细胞D3和成纤维细胞3T3，比较受试物抑制干细胞分化的浓度和抑制细胞生长的浓度评价胚胎毒性）；细胞微团（micromass,
MM）培养试验（高密度胚胎原代细胞共培养，评价对细胞增殖和分化的影响）；哺乳动物全胚胎培养（whole embryo
culture,
WEC）试验（9.5日龄大鼠胚胎体外培养，对器官形成期胚胎器官分化和生长发育的影响判定胚胎毒性和致畸性）。发育毒物体外预筛试验：小鼠胚体干细胞培养、肢芽培养、中脑微团培养、全胚胎培养。

\InfoLabel{知识来源} 《11 化学致畸作用、生殖毒性、发育毒性2021》第49页；《11
化学致畸作用、生殖毒性、发育毒性2021》第50页；《11 化学致畸作用、生殖毒性、发育毒性2021》第52页；《11
化学致畸作用、生殖毒性、发育毒性2021》第54页；《11 化学致畸作用和发育毒性2022》第17页

\ContentBand{对应往年题}

\begin{pastquestion}

\QuestionLabel{【往年题1｜11 化学致畸作用、生殖毒性、发育毒性2021（课件思考，第67页）｜思考题】}\\
要对一个化学物的致畸性进行预筛，只用一种体外替代方法够吗？为什么？体外替代方法能完全取代传统致畸试验吗？为什么？（与往年口试同型，纳入复习）

\end{pastquestion}

\begin{pastquestion}

\QuestionLabel{【往年题2｜11 化学致畸作用和发育毒性2022（课件思考题，第18页）｜思考题】}\\
发育毒物体外预筛试验（与往年口试同型，纳入复习）

\end{pastquestion}

\PointSeparator

\chapter{神经和神经行为毒理学}\label{users__fpb__26_spring__pkusph_study__tmp__pdfs__prevention-medicine-past-exams__02-ux6bd2ux7406ux5b66ux5f80ux5e74ux9898ux8003ux70b9ux6574ux7406.md__ux7b2cux5341ux4e8cux7ae0-12-ux795eux7ecfux548cux795eux7ecfux884cux4e3aux6bd2ux7406ux5b66}

\section{迟发性神经毒性的定义和表现（及作用机理、实验动物）}\label{users__fpb__26_spring__pkusph_study__tmp__pdfs__prevention-medicine-past-exams__02-ux6bd2ux7406ux5b66ux5f80ux5e74ux9898ux8003ux70b9ux6574ux7406.md__ux8fdfux53d1ux6027ux795eux7ecfux6bd2ux6027ux7684ux5b9aux4e49ux548cux8868ux73b0ux53caux4f5cux7528ux673aux7406ux5b9eux9a8cux52a8ux7269}

\InfoLabel{考频} 8 次

\InfoLabel{对应小节} 三、迟发性神经毒性 /《12 神经和神经行为毒理学》第30--33页

\InfoLabel{匹配依据}
``迟发性神经毒性的定义和表现''``迟发性神经毒性的概念和表现''反复出现，对应本小节定义、机理及试验动物。

\ContentBand{知识点原文摘取}

迟发性神经毒性（delayed
toxicity）定义：在中毒反应发生后8～14天，再次出现的较持久的神经中毒反应，主要表现为迟缓性麻痹或轻瘫，而后出现脊髓损伤体征，如共济失调或强直。多见于有机磷农药、CO中毒。中毒作用机理（OPIDN）：第一步神经毒靶酯老化；第二步钙稳态的失衡，引起依赖钙激活蛋白酶（CANP）的活化；第三步CANP引起远端轴突的降解。常见引起迟发性神经毒性的有机磷农药：TOCP（三磷甲苯磷酸酯）、丙胺氟磷、敌敌畏和敌百虫等。有机磷农药的毒性试验中应该有迟发性神经毒性试验，实验动物：母鸡；剂量：LD50；试验期限：21天。

\InfoLabel{知识来源} 《12 神经和神经行为毒理学》第30页；《12 神经和神经行为毒理学》第31页；《12
神经和神经行为毒理学》第32页；《12 神经和神经行为毒理学》第33页

\ContentBand{对应往年题}

\begin{pastquestion}

\QuestionLabel{【往年题1｜07级毒理口试题回忆.doc｜口试】}\\
迟发性神经毒性的定义和表现（25号题）；迟发型神经毒性的定义及变现（54号题，郝东东追问可能的机制是什么？用什么动物？可能导致迟发神毒的毒物有啥？）；迟发性神经毒性的概念和表现（59号题）

\end{pastquestion}

\begin{pastquestion}

\QuestionLabel{【往年题2｜09级口试题目.docx｜口试】}\\
简述迟发型神经毒性的定义和表现（11）；迟发性神经毒性定义及表现（25）

\end{pastquestion}

\begin{pastquestion}

\QuestionLabel{【往年题3｜21级毒理期末口试回忆.xlsx｜口试】}\\
迟发性神经毒性的定义和表现（序号19，备注：wxt老师追问迟发神经毒性用什么实验动物）；迟发性神经毒性的定义及表现（序号65）

\end{pastquestion}

\begin{pastquestion}

\QuestionLabel{【往年题4｜19级毒理期末口试回忆.xlsx｜口试】}\\
迟发性神经毒性（序号65）

\end{pastquestion}

\begin{pastquestion}

\QuestionLabel{【往年题5｜17级毒理试题汇总.xlsx｜口试】}\\
迟发性神经毒性的定义和表现（题号20）

\end{pastquestion}

\begin{pastquestion}

\QuestionLabel{【往年题6｜20级完整试题回忆.docx｜口试】}\\
迟发型神经毒性的定义和表现（没答上来老师还追问了组织损伤特点）（其他同学考题）

\end{pastquestion}

\PointSeparator

\section{神经毒性的研究方法（神经系统毒性研究方法）}\label{users__fpb__26_spring__pkusph_study__tmp__pdfs__prevention-medicine-past-exams__02-ux6bd2ux7406ux5b66ux5f80ux5e74ux9898ux8003ux70b9ux6574ux7406.md__ux795eux7ecfux6bd2ux6027ux7684ux7814ux7a76ux65b9ux6cd5ux795eux7ecfux7cfbux7edfux6bd2ux6027ux7814ux7a76ux65b9ux6cd5}

\InfoLabel{考频} 5 次

\InfoLabel{对应小节} 四、神经毒性的研究方法 /《12 神经和神经行为毒理学》第34页

\InfoLabel{匹配依据}
``神经毒性的研究方法''``神经系统毒性研究方法''反复出现，对应本小节人体研究、动物研究、细胞组织研究。

\ContentBand{知识点原文摘取}

四、神经毒性的研究方法：1.人体的研究------临床观察或者检查（神经学检查、电生理检查、生化检查）；2.动物研究------（1）功能试验组合（functional
observational FOB，US EPA
1994）：自主体征、感觉、运动、生命体征；（2）神经学检查；（3）形态学检查：病理；（4）电生理；（5）生化检查：相应的酶、靶蛋白；（6）分子生物学检测：信号通路，表观遗传改变等；3.细胞组织研究------神经元细胞的体外培养、干细胞的定向诱导和神经元干细胞的应用、中脑微团的培养、脑片的体外培养。

\InfoLabel{知识来源} 《12 神经和神经行为毒理学》第34页

\ContentBand{对应往年题}

\begin{pastquestion}

\QuestionLabel{【往年题1｜15级毒理学考试题目汇总.xls｜口试】}\\
神经毒性的研究方法（题号17）；神经系统毒性的研究方法（题号44）；神经毒性的作用部位，机制（题号79）

\end{pastquestion}

\begin{pastquestion}

\QuestionLabel{【往年题2｜14级毒理考题汇总.xls｜口试】}\\
神经系统毒性研究方法（题号80）；神经系统毒性研究方法有哪些（题号17）；神经毒性的机制和效应（题号45）

\end{pastquestion}

\begin{pastquestion}

\QuestionLabel{【往年题3｜12 神经和神经行为毒理学（课件内容，第34页）｜复习】}\\
神经毒性的研究方法（人体/动物/细胞组织三层次，纳入复习）

\end{pastquestion}

\PointSeparator

\section{神经毒性的损伤类型/机制（按作用部位和机制分类）及外源化学物对神经系统毒作用的临床表现}\label{users__fpb__26_spring__pkusph_study__tmp__pdfs__prevention-medicine-past-exams__02-ux6bd2ux7406ux5b66ux5f80ux5e74ux9898ux8003ux70b9ux6574ux7406.md__ux795eux7ecfux6bd2ux6027ux7684ux635fux4f24ux7c7bux578bux673aux5236ux6309ux4f5cux7528ux90e8ux4f4dux548cux673aux5236ux5206ux7c7bux53caux5916ux6e90ux5316ux5b66ux7269ux5bf9ux795eux7ecfux7cfbux7edfux6bd2ux4f5cux7528ux7684ux4e34ux5e8aux8868ux73b0}

\InfoLabel{考频} 6 次

\InfoLabel{对应小节} 二、神经毒性类型及机制；外源化学物对神经系统毒作用的表现 /《12
神经和神经行为毒理学》第17--27页

\InfoLabel{匹配依据}
``按机制和部位的神经毒性分类''``神经毒效应按照作用部位和机制可分为哪几类''``神经中毒的临床表现分类''反复出现，对应本小节。

\ContentBand{知识点原文摘取}

二、神经毒性类型及机制：1.血－脑及血－神经屏障与神经毒性损伤；2.神经元损伤（neuronopathy，有关化学物如氨基糖甙类抗生素、铅、汞、锰、铊等）；3.轴病（axonopathy，如丙烯酰胺、二硫化碳、有机磷、正己烷等）；4.微管相关的神经毒作用---周围神经病（如秋水仙素、紫杉醇）；5.髓鞘病（myelinopathies）；6.神经传递与神经毒性（transmission
toxicity，如苯丙胺、阿托品、可卡因、烟碱等）。外源化学物对神经系统毒作用的表现：1.中毒性神经官能症（以脑功能失调和精神障碍为主，可逆）；2.中毒性脑病（中枢神经系统器质性病变，急性/慢性）；3.中毒性神经炎（损害周围神经导致单神经炎或多神经炎）。

\InfoLabel{知识来源} 《12 神经和神经行为毒理学》第17页；《12 神经和神经行为毒理学》第19页；《12
神经和神经行为毒理学》第20页；《12 神经和神经行为毒理学》第25页；《12 神经和神经行为毒理学》第26页；《12
神经和神经行为毒理学》第27页

\ContentBand{对应往年题}

\begin{pastquestion}

\QuestionLabel{【往年题1｜15级毒理学考试题目汇总.xls｜口试】}\\
神经毒性临床反应类型（题号78）；神经中毒的临床表现分类（题号15）；按机制和部位的神经毒性分类（题号16）；外源化学物导致神经毒性的类型（题号43）

\end{pastquestion}

\begin{pastquestion}

\QuestionLabel{【往年题2｜14级毒理考题汇总.xls｜口试】}\\
神经毒性的临床表现（记得举例子）（题号78）；神经毒性临床类型（题号43）；根据神经毒作用机制与部位，简述神经毒效应（题号16）；神经毒效应按照作用部位和机制可分为哪几类（题号79）；按机制和部位的神经毒性分类（题号16）

\end{pastquestion}

\begin{pastquestion}

\QuestionLabel{【往年题3｜12 神经和神经行为毒理学（课件内容，第16页）｜复习】}\\
神经毒理学定义：研究外源化学物引起神经系统结构和功能损害作用及其损伤机制的一门学科（纳入复习）

\end{pastquestion}

\PointSeparator

\section{神经行为毒理学的定义及研究方法}\label{users__fpb__26_spring__pkusph_study__tmp__pdfs__prevention-medicine-past-exams__02-ux6bd2ux7406ux5b66ux5f80ux5e74ux9898ux8003ux70b9ux6574ux7406.md__ux795eux7ecfux884cux4e3aux6bd2ux7406ux5b66ux7684ux5b9aux4e49ux53caux7814ux7a76ux65b9ux6cd5}

\InfoLabel{考频} 1 次

\InfoLabel{对应小节} 神经行为毒理学 /《12 神经和神经行为毒理学》第45--53页

\InfoLabel{匹配依据}
课件单列神经行为毒理学定义与研究方法，往年神经毒性研究方法题常需扩展至神经行为学方法，单列以备复习。

\ContentBand{知识点原文摘取}

神经行为毒理学（neurobehavioral
toxicology）定义：主要是研究外源性化学物特别是低剂量长期接触对人神经行为的毒性效应，亦即人的心理功能的毒性效应。特点：敏感、能定量、可早期检出亚临床表现、是整体试验。研究方法：（1）动物试验------主要检测学习记忆，空间辨别试验（Y迷宫、Morris水迷宫）、厌恶学习（穿梭箱、避暗、跳台）、自主活动度测定、焦虑测定（高架十字迷宫）、药物成瘾性（条件位置偏好）；（2）人的观察------心理功能测试量表（韦氏智力/记忆量表），NCTB（WHO
1983），计算机化神经行为评价系统（NES）。

\InfoLabel{知识来源} 《12 神经和神经行为毒理学》第45页；《12 神经和神经行为毒理学》第46页；《12
神经和神经行为毒理学》第53页

\ContentBand{对应往年题}

\begin{pastquestion}

\QuestionLabel{【往年题1｜15级毒理学考试题目汇总.xls｜口试】}\\
神经毒性的研究方法（题号17，含神经行为学方法扩展，归入本考点辅助复习）

\end{pastquestion}

\PointSeparator

\chapter{安全性评价}\label{users__fpb__26_spring__pkusph_study__tmp__pdfs__prevention-medicine-past-exams__02-ux6bd2ux7406ux5b66ux5f80ux5e74ux9898ux8003ux70b9ux6574ux7406.md__ux7b2cux5341ux4e09ux7ae0-13-ux5b89ux5168ux6027ux8bc4ux4ef7}

\section{化学品安全性评价的概念和主要内容；四阶段毒理学安全性评价及为什么要分阶段进行}\label{users__fpb__26_spring__pkusph_study__tmp__pdfs__prevention-medicine-past-exams__02-ux6bd2ux7406ux5b66ux5f80ux5e74ux9898ux8003ux70b9ux6574ux7406.md__ux5316ux5b66ux54c1ux5b89ux5168ux6027ux8bc4ux4ef7ux7684ux6982ux5ff5ux548cux4e3bux8981ux5185ux5bb9ux56dbux9636ux6bb5ux6bd2ux7406ux5b66ux5b89ux5168ux6027ux8bc4ux4ef7ux53caux4e3aux4ec0ux4e48ux8981ux5206ux9636ux6bb5ux8fdbux884c}

\InfoLabel{考频} 14 次

\InfoLabel{对应小节} 化学物对生物体健康风险的安全性评价；四阶段毒理学安全性评价；实验结果判定原则 /《13
安全性评价》第4--8页

\InfoLabel{匹配依据}
``毒理学安全性评价的概念和基本内容''``安全性评价为什么要分阶段进行？各阶段的具体内容''``化学品安全性评价的概念和内容''反复出现，对应本小节四阶段及分阶段原因（含历届备注：经济成本因素，非食入化学物只需前三阶段，食品和药品须完整四阶段）。

\ContentBand{知识点原文摘取}

安全性评价（safety
evaluation）：按照一定的程序要求对外源化学物的毒作用进行检测，以评价化学物的毒作用特点、毒作用的剂量－反应关系，经足够的试验之后，提出NOAEL和LOAEL，再除以一定的安全系数，制订出安全限值。开展安全性评价必备的几个基本要素：GLP（good
laboratory practice，实验室认证认可）、SOP（standard operation procedures）、QAU（quality assurance
unit）。四阶段毒理学安全性评价：第一阶段------急性毒性试验，局部毒性试验；第二阶段------蓄积试验、遗传毒理学试验和致畸试验；第三阶段------亚慢性毒性试验、生殖试验和代谢试验；第四阶段------慢性毒性试验和致癌试验。实验结果判定原则：急性毒性小于人摄入量100倍放弃；遗传毒性两项阳性弃用，一项阳性再选两个备选试验（其中一个为体内试验）；90天经口毒性100、300倍；生殖毒性100、300倍；慢性毒性和致癌试验50、100倍。

\InfoLabel{知识来源} 《13 安全性评价》第4页；《13 安全性评价》第5页；《13 安全性评价》第6页；《13
安全性评价》第7页；《13 安全性评价》第8页

\ContentBand{对应往年题}

\begin{pastquestion}

\QuestionLabel{【往年题1｜07级毒理口试题回忆.doc｜口试】}\\
安全评价的主要内容（还要答人群流行病实验）（未知题号）；安全性评价的概念，几个阶段的内容（4号题）；安全性评价为什么要分阶段进行？各阶段的具体内容？（8号题）；化学品安全性评价程序为什么要分阶段进行？分几个阶段？各阶段包括哪些项目（10号题）；安全性评价的四阶段以及为什么要分阶段进行（17号题）；毒理学安全性评价的四个阶段？为什么要分阶段进行？（36号题）；安全性评价为啥要分阶段？各阶段有啥？（54号题，郝东东追问什么时候用三阶段试验------药物，非药物用啥试验）

\end{pastquestion}

\begin{pastquestion}

\QuestionLabel{【往年题2｜09级口试题目.docx｜口试】}\\
毒理学安全性评价的概念和基本内容（7、34）；安全性评价的主要内容（9）；毒理学安全性评价为什么要分阶段进行？都有哪几个步骤？（10，备注：补充经济成本因素，非食入化学物只需前三阶段，食品和药品必须完整四阶段）；安全性评价的定义及内容（23，备注：针对新的、未知的化学物，已知化学物多用危险性评定）；为什么毒理学安全性评价要分阶段进行？分为哪几个阶段，各阶段有什么内容（12）

\end{pastquestion}

\begin{pastquestion}

\QuestionLabel{【往年题3｜21级毒理期末口试回忆.xlsx｜口试】}\\
化学品安全性评价的概念和内容（序号31）；化学物安全评价的概念和内容（序号75）；安全性评价分哪几个阶段以及为什么要分阶段（序号64，备注：以食品为例会做哪些实验）

\end{pastquestion}

\begin{pastquestion}

\QuestionLabel{【往年题4｜17级毒理试题汇总.xlsx｜口试】}\\
安全性评价分阶段原因，阶段内容，具体项目（题号32，备注：郝老师死亡追问，安全性评价阶段的''具体项目''，如新药生殖发育毒性要做哪些实验）；安全性评价概念，内容（题号31）；化学品安全性评价概念、内容（题号67）

\end{pastquestion}

\begin{pastquestion}

\QuestionLabel{【往年题5｜15级毒理学考试题目汇总.xls｜口试】}\\
化学品毒理学安全性评价的内容（题号21）；化学品安全性评价的主要内容（题号40）

\end{pastquestion}

\begin{pastquestion}

\QuestionLabel{【往年题6｜14级毒理考题汇总.xls｜口试】}\\
毒理学安全性评价的内容（题号1）；化学物安全性评价的内容（题号21）；化学品安全性评价的内容（题号40）；安全性评价的方法（四阶段）（题号60，追问试验的四个阶段是否都需要完成）

\end{pastquestion}

\PointSeparator

\section{生物药物与化学药物在安全性评价上的异同（侧重点/特别关注点）}\label{users__fpb__26_spring__pkusph_study__tmp__pdfs__prevention-medicine-past-exams__02-ux6bd2ux7406ux5b66ux5f80ux5e74ux9898ux8003ux70b9ux6574ux7406.md__ux751fux7269ux836fux7269ux4e0eux5316ux5b66ux836fux7269ux5728ux5b89ux5168ux6027ux8bc4ux4ef7ux4e0aux7684ux5f02ux540cux4fa7ux91cdux70b9ux7279ux522bux5173ux6ce8ux70b9}

\InfoLabel{考频} 9 次

\InfoLabel{对应小节} 生物药物的特点；生物药物的安全性评价特殊性；安全评估试验-生物药物与化学药物 /《13
安全性评价》第14--36页

\InfoLabel{匹配依据}
``生物药物和化学药物的安全性评价的异同点''``生药和化药安全性评价异同''``生物药安全性评价的关注点''反复出现，对应本小节对照表与特别关注点。

\ContentBand{知识点原文摘取}

生物药物的特点（区别于化药）：（1）空间结构很复杂致使结构确定存在不完整性；（2）生理结构差异导致生物活性差异；（3）生物药物受体广泛决定其具有多功能性；（4）生物药物是异源性大分子，具有免疫原性。安全评估试验-生物药物（ICH
S6）与化学药物（ICH
M3）：普通毒理学（生物药物啮齿和非啮齿动物通常6个月，可只用一种相关动物；化药啮齿至多6个月、非啮齿9-12个月，需2种动物）；组织交叉反应（生物药物单克隆抗体及衍生物需要，化药不需要）；遗传毒理学（生物药物通常不需要，化药需要）；安全药理学（生物药物通常不需要，可与普通毒理学合并）；生殖毒理学；致癌性试验（生物药物个案考虑；化药需大鼠和小鼠致癌试验）；光毒性试验；代谢特征鉴定（生物药物不做，化药通常需要）。生物药物特别关注点：生殖毒性、致癌性、安全药理学、组织交叉反应性、细胞因子释放、免疫原性。

\InfoLabel{知识来源} 《13 安全性评价》第14页；《13 安全性评价》第16页；《13 安全性评价》第19页；《13
安全性评价》第21页；《13 安全性评价》第22页；《13 安全性评价》第23页

\ContentBand{对应往年题}

\begin{pastquestion}

\QuestionLabel{【往年题1｜21级毒理期末口试回忆.xlsx｜口试】}\\
（生物药物相关，序号84''机体对外源化合物的处置过程''等并行考查）；生物药物安全性评价相关（结合分阶段评价）

\end{pastquestion}

\begin{pastquestion}

\QuestionLabel{【往年题2｜19级毒理期末口试回忆.xlsx｜口试】}\\
生物药（序号84相关）；机体对外源化学物处置及生物药安全性评价（结合）

\end{pastquestion}

\begin{pastquestion}

\QuestionLabel{【往年题3｜17级毒理试题汇总.xlsx｜口试】}\\
生药和化药安全性评价的侧重点（题号76）；生物药安全性评价的关注点（题号48）；生药
化药安全性评价各自的区别和要点（题号97）；生物药物和化学药物在安全性评价上的侧重点（题号88）；生物药物安全性评价（题号49，备注：郝老师问临床生物性药物的毒理实验期限）；生物药（题号10）

\end{pastquestion}

\begin{pastquestion}

\QuestionLabel{【往年题4｜15级毒理学考试题目汇总.xls｜口试】}\\
化学药物和生物药物的安全性评价的异同点（题号23）；生药和化药在安全性评价中的异同（题号42）；生物药物和化学药物安全性试验的异同（题号23）；生物药物和化学药物在安全性评价中的异同点（题号3、72）

\end{pastquestion}

\begin{pastquestion}

\QuestionLabel{【往年题5｜14级毒理考题汇总.xls｜口试】}\\
生药和化药安全性评价异同（题号42）；生物药物和化学药物的安全性评价；什么时候需要做化学药物的致癌实验（题号62）

\end{pastquestion}

\PointSeparator

\section{不同受试物选择开展试验的原则；不同物质安全性评价需注意的特殊点}\label{users__fpb__26_spring__pkusph_study__tmp__pdfs__prevention-medicine-past-exams__02-ux6bd2ux7406ux5b66ux5f80ux5e74ux9898ux8003ux70b9ux6574ux7406.md__ux4e0dux540cux53d7ux8bd5ux7269ux9009ux62e9ux5f00ux5c55ux8bd5ux9a8cux7684ux539fux5219ux4e0dux540cux7269ux8d28ux5b89ux5168ux6027ux8bc4ux4ef7ux9700ux6ce8ux610fux7684ux7279ux6b8aux70b9}

\InfoLabel{考频} 6 次

\InfoLabel{对应小节} 不同受试物选择开展试验的原则；化学品A/药物B/转基因水稻案例 /《13
安全性评价》第7页、第9--11页

\InfoLabel{匹配依据}
``不同物质安全性评价需要注意哪些内容''``不同化学物安全性评价的特点''``不同物质安全性评价的不同要点''反复出现，对应本小节按受试物类型选择试验的原则。

\ContentBand{知识点原文摘取}

不同受试物选择开展试验的原则：首创物质，构效分析（潜在毒性），使用范围广，摄入量大------急性至致癌（完整程序）；已知食品中相关物质的类似物，或者部分地区有食用历史------急毒、遗传毒性、28天经口毒性，根据具体情况确定是否继续。案例：化学品A（硫酸盐，溶于水，具有作为食品添加剂的特性）；药物B（单克隆抗体，用于癌症治疗）；转cry-1Ab蛋白水稻（抗虫，用于常规粮食）------请设计安全性评价方案。

\InfoLabel{知识来源} 《13 安全性评价》第7页；《13 安全性评价》第9页；《13 安全性评价》第10页；《13
安全性评价》第11页

\ContentBand{对应往年题}

\begin{pastquestion}

\QuestionLabel{【往年题1｜15级毒理学考试题目汇总.xls｜口试】}\\
不同化学物安全性评价的特点（题号61）；不同物质安全性评价需要注意的特殊点（题号22，备注：追问转基因食品的安全性评价问题）；不同物质安全性评价特点（题号41）；不同化学物在重复暴露实验时需要考虑的事项（题号72）

\end{pastquestion}

\begin{pastquestion}

\QuestionLabel{【往年题2｜14级毒理考题汇总.xls｜口试】}\\
不同物质的安全性评价需要注意哪些内容？（题号22）；不同物质安全性评价的不同要点（无题号）；不同物质安全性评价的特点（题号61）

\end{pastquestion}

\begin{pastquestion}

\QuestionLabel{【往年题3｜21级毒理期末口试回忆.xlsx｜口试】}\\
安全性评价分哪几个阶段以及为什么要分阶段（序号64，备注：以食品为例会做哪些实验，呼应不同受试物原则）

\end{pastquestion}

\PointSeparator

\section{对化学物进行毒理学试验的程序和项目}\label{users__fpb__26_spring__pkusph_study__tmp__pdfs__prevention-medicine-past-exams__02-ux6bd2ux7406ux5b66ux5f80ux5e74ux9898ux8003ux70b9ux6574ux7406.md__ux5bf9ux5316ux5b66ux7269ux8fdbux884cux6bd2ux7406ux5b66ux8bd5ux9a8cux7684ux7a0bux5e8fux548cux9879ux76ee}

\InfoLabel{考频} 4 次

\InfoLabel{对应小节} 四阶段毒理学安全性评价（结合《2 毒理学实验基础》第24页 受试样品的准备）/《13
安全性评价》第6--8页

\InfoLabel{匹配依据}
``对化学物进行毒理学试验的程序和项目''``化学物质毒理学实验的程序、项目''``化学物毒理实验的程序和项目''反复出现，对应安全性评价程序（含前期准备阶段对理化性质/结构类似物毒性的了解，再进入分阶段试验）。

\ContentBand{知识点原文摘取}

四阶段毒理学安全性评价：第一阶段急性毒性试验、局部毒性试验；第二阶段蓄积试验、遗传毒理学试验和致畸试验；第三阶段亚慢性毒性试验、生殖试验和代谢试验；第四阶段慢性毒性试验和致癌试验。（《2
毒理学实验基础》第24页
受试样品的准备：①纯度及杂质成分；②化学结构和理化性质；③相似化合物的毒性参考资料；④同种同一批号；⑤成分和配方必须固定；⑥稳定性受试物应一次备齐全部实验的用量------即试验程序前的前期准备阶段。）

\InfoLabel{知识来源} 《13 安全性评价》第6页；《2 毒理学实验基础》第24页

\ContentBand{对应往年题}

\begin{pastquestion}

\QuestionLabel{【往年题1｜21级毒理期末口试回忆.xlsx｜口试】}\\
外源化学物毒理学试验的程序和方法（序号3，备注：从安全性评价的角度的毒理学试验设计进行回答）；对化学物进行毒理学试验的程序和项目（序号72，备注：不要一上来就答实验设计，不要忘了前期的准备阶段，对化学物的理化性质等进行理解，查找与其结构类似的化学物的毒性）

\end{pastquestion}

\begin{pastquestion}

\QuestionLabel{【往年题2｜17级毒理试题汇总.xlsx｜口试】}\\
化学物质毒理学实验的程序、项目（题号41）；化学物毒理实验的程序和项目（题号89）

\end{pastquestion}

\PointSeparator

\chapter{风险评估、管理毒理学}\label{users__fpb__26_spring__pkusph_study__tmp__pdfs__prevention-medicine-past-exams__02-ux6bd2ux7406ux5b66ux5f80ux5e74ux9898ux8003ux70b9ux6574ux7406.md__ux7b2cux5341ux56dbux7ae0-14-ux98ceux9669ux8bc4ux4f30ux7ba1ux7406ux6bd2ux7406ux5b66}

\section{健康风险评估（health risk
assessment）的概念及主要步骤}\label{users__fpb__26_spring__pkusph_study__tmp__pdfs__prevention-medicine-past-exams__02-ux6bd2ux7406ux5b66ux5f80ux5e74ux9898ux8003ux70b9ux6574ux7406.md__ux5065ux5eb7ux98ceux9669ux8bc4ux4f30health-risk-assessmentux7684ux6982ux5ff5ux53caux4e3bux8981ux6b65ux9aa4}

\InfoLabel{考频} 9 次

\InfoLabel{对应小节} 风险评估；风险评估的步骤；危害鉴定/危害表征/暴露评定/风险表征 /《14
风险评估、管理毒理学》第2--16页

\InfoLabel{匹配依据} ``risk assessment的定义和步骤''``Health risk assessment
的概念及主要步骤''``健康风险评估的步骤和内容''反复出现，对应本小节四步骤。

\ContentBand{知识点原文摘取}

风险评估（risk
assessment）：对不良结果发生的机率进行描述和定量，包括环境风险评估和健康风险评估。健康风险评估：基于人群流行病学资料、毒理学试验资料、环境化学物的接触资料等科学数据的分析，确定接触外源化学物后对公众健康危害的可能性，包括对有害作用性质、强度的定性描述，接触水平与可能出现损害的危险性水平的定量评定及对评价结论和评定的不确定性的分析和描述。风险评估的步骤：危害识别（hazard
identification）；危害表征（hazard characterization）；暴露评定（exposure assessment）；风险表征（risk
characterization）。危害鉴定：分析确定化学物与机体接触后产生不良作用的性质和特点。暴露评定是风险评估中最为不确定的部分。风险表征：将前三步分析和结论综合，对人体危害的性质及危险度的大小做出估计，为危险度管理提供依据。

\InfoLabel{知识来源} 《14 风险评估、管理毒理学》第2页；《14 风险评估、管理毒理学》第3页；《14
风险评估、管理毒理学》第5页；《14 风险评估、管理毒理学》第6页；《14 风险评估、管理毒理学》第13页；《14
风险评估、管理毒理学》第16页

\ContentBand{对应往年题}

\begin{pastquestion}

\QuestionLabel{【往年题1｜07级毒理口试题回忆.doc｜口试】}\\
risk assessment的定义和步骤（1号、28号题）；risk assessment
的概念和具体步骤（18号题）；危险性评价（英文）的步骤（59号题）。备注：危险度评价老师问剂量反应关系评定这个步骤怎么做的，要分零阈值和非零阈值答。

\end{pastquestion}

\begin{pastquestion}

\QuestionLabel{【往年题2｜09级口试题目.docx｜口试】}\\
危险评估的概念和步骤（16）；risk assessment的定义和步骤（28）；risk assessment的意义与步骤（48）

\end{pastquestion}

\begin{pastquestion}

\QuestionLabel{【往年题3｜19级毒理期末口试回忆.xlsx｜口试】}\\
Health risk assessment 的概念及主要步骤（序号51，备注：魏老师可能会对一些细节进行补充提问）

\end{pastquestion}

\begin{pastquestion}

\QuestionLabel{【往年题4｜17级毒理试题汇总.xlsx｜口试】}\\
health risk assessment
概念主要步骤（题号4）；健康风险评估的步骤和内容（题号65，备注：追问健康风险评估的概念和与安全性评价的区别）

\end{pastquestion}

\begin{pastquestion}

\QuestionLabel{【往年题5｜15级毒理学考试题目汇总.xls｜口试】}\\
健康风险评价的步骤和内容（题号65，备注：追问健康风险评估的概念和与安全性评价的区别）

\end{pastquestion}

\begin{pastquestion}

\QuestionLabel{【往年题6｜14级毒理考题汇总.xls｜口试】}\\
健康风险评估的内容及步骤（题号56，备注：追问危害识别的应用方法即剂量反应关系研究方法）；健康风险评价的步骤和内容（题号65，备注：老师说还有危险管理和危险交流）；风险评估的步骤和内容（题号37）

\end{pastquestion}

\PointSeparator

\section{危害表征中的阈值法与非阈值法（有阈/无阈化学物的剂量-反应关系评定）}\label{users__fpb__26_spring__pkusph_study__tmp__pdfs__prevention-medicine-past-exams__02-ux6bd2ux7406ux5b66ux5f80ux5e74ux9898ux8003ux70b9ux6574ux7406.md__ux5371ux5bb3ux8868ux5f81ux4e2dux7684ux9608ux503cux6cd5ux4e0eux975eux9608ux503cux6cd5ux6709ux9608ux65e0ux9608ux5316ux5b66ux7269ux7684ux5242ux91cf-ux53cdux5e94ux5173ux7cfbux8bc4ux5b9a}

\InfoLabel{考频} 8 次

\InfoLabel{对应小节} 危害表征；有阈化学物的剂量---反应关系评定（RfD、BMD）；无阈化学物的剂量---反应关系评定
/《14 风险评估、管理毒理学》第6--12页

\InfoLabel{匹配依据}
``危害表征的阈值法和非阈值法''``阈值法与非阈值法''``健康风险评估中，阈值法和非阈值法的应用''反复出现，对应本小节（致癌、致突变用非阈值法，致畸是有阈值，见历届备注）。

\ContentBand{知识点原文摘取}

危害表征（hazard
characterization）：确定外源化学物接触水平与有害效应发生频率之间的关系。有阈化学物的剂量---反应关系评定：参考剂量（RfD）------日平均接触剂量的估计值，RfD=NOAEL或LOAEL/(UFs×MF)；关键效应------一般为具有最低NOAEL的有害效应；基准剂量（BMD）------依据动物试验剂量-反应关系结果，用统计学模式求得引起一定比例（通常1-10\%）动物出现阳性反应剂量的95\%可信限区间的下限值。无阈化学物的剂量---反应关系评定：概率分布模型；机制模型（一次打击、多次打击、多阶段、线性多阶段、随机两阶段模型）；PBPK模型；BBDR模型。（致癌、致突变用非阈值法，致畸是有阈值法。）

\InfoLabel{知识来源} 《14 风险评估、管理毒理学》第6页；《14 风险评估、管理毒理学》第7页；《14
风险评估、管理毒理学》第8页；《14 风险评估、管理毒理学》第9页；《14 风险评估、管理毒理学》第11页

\ContentBand{对应往年题}

\begin{pastquestion}

\QuestionLabel{【往年题1｜07级毒理口试题回忆.doc｜口试】}\\
有阈毒性和无阈毒性的评价方法（37号题）

\end{pastquestion}

\begin{pastquestion}

\QuestionLabel{【往年题2｜09级口试题目.docx｜口试】}\\
有阈和无阈化学物的剂量反应关系（21）；有阈和无阈的剂量反应关系评价（46）

\end{pastquestion}

\begin{pastquestion}

\QuestionLabel{【往年题3｜21级毒理期末口试回忆.xlsx｜口试】}\\
阈值法与非阈值法（序号32）；危害表征的有阈值法和无阈值法（序号78）

\end{pastquestion}

\begin{pastquestion}

\QuestionLabel{【往年题4｜19级毒理期末口试回忆.xlsx｜口试】}\\
健康风险评估中，阈值法和非阈值法的应用（序号48，备注：致癌、致突变用非阈值法，致畸是可阈值；RfC的概念、计算公式）；危害表征的有阈值法和无阈值法（序号78）；阈值法和非阈值法（序号62，备注：参考剂量的估计方法、不确定性的来源）

\end{pastquestion}

\begin{pastquestion}

\QuestionLabel{【往年题5｜17级毒理试题汇总.xlsx｜口试】}\\
危害表征的阈值法和非阈值法（题号66、98）；阈值法与非阈值法（题号56，备注：一定要看书，不是学长总结的那个）；在健康风险评估中，危害表征的阈值法和非阈值法（题号8）

\end{pastquestion}

\begin{pastquestion}

\QuestionLabel{【往年题6｜15级毒理学考试题目汇总.xls｜口试】}\\
阈值法与非阈值法（题号57）；危害表征的阈值法和非阈值法（题号66）

\end{pastquestion}

\begin{pastquestion}

\QuestionLabel{【往年题7｜14级毒理考题汇总.xls｜口试】}\\
简述危害表征中的阈值法与非阈值法（题号18，备注：追问两种方法用于评价哪类化学品）；阈值法非阈值法（题号57，备注：非阈值法有一个百万分之一的系数）；危害表征的阈值法和非阈值法（题号66）

\end{pastquestion}

\PointSeparator

\section{安全性评价和风险评估的联系和区别及各自应用领域}\label{users__fpb__26_spring__pkusph_study__tmp__pdfs__prevention-medicine-past-exams__02-ux6bd2ux7406ux5b66ux5f80ux5e74ux9898ux8003ux70b9ux6574ux7406.md__ux5b89ux5168ux6027ux8bc4ux4ef7ux548cux98ceux9669ux8bc4ux4f30ux7684ux8054ux7cfbux548cux533aux522bux53caux5404ux81eaux5e94ux7528ux9886ux57df}

\InfoLabel{考频} 3 次

\InfoLabel{对应小节} 安全性评价和风险评估的联系和区别；研究/风险评估/风险管理的关系 /《14
风险评估、管理毒理学》第19--20页

\InfoLabel{匹配依据}
``安全性评价和危险评估的关系及各自的应用领域''直接对应本小节（部分并入风险评估题作追问）。

\ContentBand{知识点原文摘取}

安全性评价和风险评估的联系和区别：健康风险评估是在安全性评价的基础上发展起来的；安全性评价和风险评估的危害识别阶段所用毒性测试的实验方法基本相同；安全性评价表示为确定安全的程序，具有预警性质，以NOAEL/LOAEL作为外推的起始点，并考虑变异性和不确定性，制定安全限值或暴露指导值；安全性评价常用于①暴露可能是受控制的化学物安全评价情形中，②新化学物或新产品生产、使用的许可和管理；风险评估是风险分析决策程序的一部分。

\InfoLabel{知识来源} 《14 风险评估、管理毒理学》第19页；《14 风险评估、管理毒理学》第20页

\ContentBand{对应往年题}

\begin{pastquestion}

\QuestionLabel{【往年题1｜17级毒理试题汇总.xlsx｜口试】}\\
安全性评价和危险评估的关系及各自的应用领域（题号62）；健康风险评估的概念和与安全性评价的区别（题号65 追问）

\end{pastquestion}

\begin{pastquestion}

\QuestionLabel{【往年题2｜15级毒理学考试题目汇总.xls｜口试】}\\
健康风险评估的概念和与安全性评价的区别（题号65 追问）

\end{pastquestion}

\begin{pastquestion}

\QuestionLabel{【往年题3｜09级口试题目.docx｜口试】}\\
``安全性评价''针对新的、未知的化学物，已知的化学物多用危险性评定（23号题备注）

\end{pastquestion}

\PointSeparator

\section{管理毒理学；化学品管理过程中毒理学家的作用}\label{users__fpb__26_spring__pkusph_study__tmp__pdfs__prevention-medicine-past-exams__02-ux6bd2ux7406ux5b66ux5f80ux5e74ux9898ux8003ux70b9ux6574ux7406.md__ux7ba1ux7406ux6bd2ux7406ux5b66ux5316ux5b66ux54c1ux7ba1ux7406ux8fc7ux7a0bux4e2dux6bd2ux7406ux5b66ux5bb6ux7684ux4f5cux7528}

\InfoLabel{考频} 5 次

\InfoLabel{对应小节} 管理毒理学；毒理学研究和化学物质管理的相互关系；化学品管理过程中毒理学家的作用 /《14
风险评估、管理毒理学》第21--30页

\InfoLabel{匹配依据}
``毒理学家在安全性管理方面的工作''``作为毒理学家，能在安全性管理中做什么''``在化学品管理中毒理学家的作用''反复出现，对应本小节。

\ContentBand{知识点原文摘取}

管理毒理学（Regulatory
Toxicology）：依据描述和机制毒理学对毒作用规律的研究成果，协助政府部门制订相关法规条例和管理措施并付诸实施，以确保对化学品、药品、食品等进行有效的管理，保障接触人群的健康和环境安全。化学品管理过程中毒理学家的作用：（1）参与有关法律、法规的制订；（2）对化学品分类、分级、标签管理；（3）提供技术咨询和技术支持；（4）卫生标准和环境标准制订；（5）对新化学物质进行毒理学安全性评价，并参与其专业技术评审；（6）对重要的环境污染物和化学品进行健康危险度评价；（7）参与化学事故的应急救援。毒理学研究和化学物质管理的相互关系：总体目标一致；化学品的管理是以行政法规为依据采取行政措施，而不是直接依据毒理学的科研结论。

\InfoLabel{知识来源} 《14 风险评估、管理毒理学》第21页；《14 风险评估、管理毒理学》第23页；《14
风险评估、管理毒理学》第24页；《14 风险评估、管理毒理学》第30页

\ContentBand{对应往年题}

\begin{pastquestion}

\QuestionLabel{【往年题1｜14级毒理考题汇总.xls｜口试】}\\
毒理学家在安全性管理方面的工作（题号36）；安全性管理毒理学工作者能干啥（题号36）；毒理工作者可以做些什么（题号16）

\end{pastquestion}

\begin{pastquestion}

\QuestionLabel{【往年题2｜15级毒理学考试题目汇总.xls｜口试】}\\
作为毒理学家，能在安全性管理中做什么（题号55）；在化学品管理中毒理学家的作用（无题号）；毒理工作者的任务（题号16）

\end{pastquestion}

\begin{pastquestion}

\QuestionLabel{【往年题3｜07级毒理口试题回忆.doc｜口试】}\\
（管理毒理学相关，结合安全性评价/风险管理体系考查）

\end{pastquestion}

\PointSeparator

\section{GHS 的概念及健康危害分类（10
类）}\label{users__fpb__26_spring__pkusph_study__tmp__pdfs__prevention-medicine-past-exams__02-ux6bd2ux7406ux5b66ux5f80ux5e74ux9898ux8003ux70b9ux6574ux7406.md__ghs-ux7684ux6982ux5ff5ux53caux5065ux5eb7ux5371ux5bb3ux5206ux7c7b10-ux7c7b}

\InfoLabel{考频} 5 次

\InfoLabel{对应小节} GHS；GHS危险性分类；GHS健康危害性分类 /《14 风险评估、管理毒理学》第37--43页（并见《5
毒作用机制、急性毒性和局部毒性试验》第33--36页 急性毒性分级 GHS）

\InfoLabel{匹配依据}
``GHS的概念解释；健康危害10类''``GHS健康危害的分类''``GHS健康危害主要内容（10项那个）''反复出现，对应本小节；YYS题库补充亦明确
GHS+10大健康危害。

\ContentBand{知识点原文摘取}

GHS：《全球化学品统一分类和标签制度》（Globally Harmonized System of Classification and Labelling of
Chemicals）是在联合国有关机构的协调下，以世界各国现行的主要化学品分类制度为基础，创建的一套科学的、统一标准化的化学品分类和标签制度。GHS分类适用于所有化学品（含农药），既适用于纯化学物质，也适用于混合物；但不适用于医药品、食品添加剂、食品中农药残留物或消费者使用的化妆品。GHS危险性分类：理化危害16类、健康危害10类、环境危害2类。GHS健康危害性分类（10类）：急性毒性；皮肤腐蚀/刺激性；严重眼损伤/眼刺激性；呼吸或皮肤致敏性；生殖细胞致突变性；致癌性；生殖毒性；特定靶器官毒性（单次接触）；特定靶器官毒性（反复接触）；吸入危害性。

\InfoLabel{知识来源} 《14 风险评估、管理毒理学》第37页；《14 风险评估、管理毒理学》第38页；《14
风险评估、管理毒理学》第39页；《14 风险评估、管理毒理学》第40页；《5
毒作用机制、急性毒性和局部毒性试验》第33页

\ContentBand{对应往年题}

\begin{pastquestion}

\QuestionLabel{【往年题1｜21级毒理期末口试回忆.xlsx｜口试】}\\
GHS的概念解释；健康危害10类（序号39）

\end{pastquestion}

\begin{pastquestion}

\QuestionLabel{【往年题2｜17级毒理试题汇总.xlsx｜口试】}\\
GHS的概念和健康危害分类（题号58）

\end{pastquestion}

\begin{pastquestion}

\QuestionLabel{【往年题3｜15级毒理学考试题目汇总.xls｜口试】}\\
GHS健康危害主要内容（10项那个）（无题号）；GHS10种毒性（题号19）；健康危害性分类（题号67）

\end{pastquestion}

\begin{pastquestion}

\QuestionLabel{【往年题4｜14级毒理考题汇总.xls｜口试】}\\
ghs健康危害的分类（题号39）；GHS的危害性分类：重点是健康危害的10项（题号58）；GHS健康危害的分类（题号19）

\end{pastquestion}

\begin{pastquestion}

\QuestionLabel{【往年题5｜毒理学-来自课本的题库补充byYYS.docx｜补充考点】}\\
GHS分类+10大健康危害（含 GHS 定义、适用范围与健康危害10类）

\end{pastquestion}

\PointSeparator

\section{GLP、SOP、QAU
的概念和意义}\label{users__fpb__26_spring__pkusph_study__tmp__pdfs__prevention-medicine-past-exams__02-ux6bd2ux7406ux5b66ux5f80ux5e74ux9898ux8003ux70b9ux6574ux7406.md__glpsopqau-ux7684ux6982ux5ff5ux548cux610fux4e49}

\InfoLabel{考频} 4 次

\InfoLabel{对应小节} GLP（Good Laboratory Practice）；GLP的主要内容 /《14
风险评估、管理毒理学》第31--32页（并见《13 安全性评价》第5页、《1 绪论》第32页）

\InfoLabel{匹配依据} ``GLP、SOP和QA的概念和意义''``GLP SOP QA概念 意义''反复出现，对应本小节。

\ContentBand{知识点原文摘取}

GLP（Good Laboratory
Practice）：为保证试验数据的准确、可靠，对于非临床的实验室研究在研究计划制定、实施、监督、记录及报告等各项工作的过程和条件提出的要求和指导。GLP的主要内容：关键-人员素质；基础-实验设施；手段-SOP（Standard
Operation Procedure）；保证-QAU（Quality Assurance Unit）。

【翻译】GLP 指良好实验室规范，SOP 指标准操作规程，QAU 指质量保证部门。GLP
通过人员素质、实验设施、标准操作规程和质量保证体系，确保非临床实验室研究数据准确可靠。

\InfoLabel{知识来源} 《14 风险评估、管理毒理学》第31页；《14 风险评估、管理毒理学》第32页；《13
安全性评价》第5页；《1 绪论》第32页

\ContentBand{对应往年题}

\begin{pastquestion}

\QuestionLabel{【往年题1｜07级毒理口试题回忆.doc｜口试】}\\
简述GLP、SOP和QA的概念和意义（51号题）；GLP SOP QA概念 意义（题号未知）

\end{pastquestion}

\begin{pastquestion}

\QuestionLabel{【往年题2｜21级毒理期末口试回忆.xlsx｜口试】}\\
GLP SOP QA意义与解释（序号52）

\end{pastquestion}

\begin{pastquestion}

\QuestionLabel{【往年题3｜19级毒理期末口试回忆.xlsx｜口试】}\\
GLP，SOP，QA的概念和意义（序号52）

\end{pastquestion}

\begin{pastquestion}

\QuestionLabel{【往年题4｜17级毒理试题汇总.xlsx｜口试】}\\
glp,sop,qa概念和意义（题号93）

\end{pastquestion}

\PointSeparator

\section{风险管理与风险交流；可接受风险、实用安全剂量；动物外推不确定性系数（背景与延伸）}\label{users__fpb__26_spring__pkusph_study__tmp__pdfs__prevention-medicine-past-exams__02-ux6bd2ux7406ux5b66ux5f80ux5e74ux9898ux8003ux70b9ux6574ux7406.md__ux98ceux9669ux7ba1ux7406ux4e0eux98ceux9669ux4ea4ux6d41ux53efux63a5ux53d7ux98ceux9669ux5b9eux7528ux5b89ux5168ux5242ux91cfux52a8ux7269ux5916ux63a8ux4e0dux786eux5b9aux6027ux7cfbux6570ux80ccux666fux4e0eux5ef6ux4f38}

\InfoLabel{考频} 2 次

\InfoLabel{对应小节} 风险管理；风险交流；可接受风险/实用安全剂量；推导慢性RfD时不确定性系数和修正系数 /《14
风险评估、管理毒理学》第4页、第10页、第17--18页

\InfoLabel{匹配依据}
历届''健康风险评价的步骤和内容''被追问''还有危险管理和危险交流''\,``不确定性的来源''，对应本小节风险管理/交流及不确定性系数。

\ContentBand{知识点原文摘取}

风险管理（risk
management）：管理部门依据危险度评价的结果，综合技术、社会、经济及政治等其他因素，确定可接受的危险度，制订容许限量标准。风险交流（risk
communication）：在风险评估者、风险管理者、消费者和其他有关各方之间进行有关风险和风险相关因素的信息和观点的交流过程，应贯穿风险分析的各个阶段。可接受风险（acceptable
risk）：公众及社会在精神及心理等方面可以接受的风险。实用安全剂量（VSD）：相近于可接受风险水平的化学物接触剂量。推导慢性RfD时的不确定性系数：B（人群个体敏感性的变异）、A（使用动物资料外推到人）、S（亚慢性研究外推到慢性研究）、L（由LOAEL代替NOAEL）、D（数据库不完整的不确定性）、MF（修正系数）。

\InfoLabel{知识来源} 《14 风险评估、管理毒理学》第4页；《14 风险评估、管理毒理学》第10页；《14
风险评估、管理毒理学》第17页；《14 风险评估、管理毒理学》第18页

\ContentBand{对应往年题}

\begin{pastquestion}

\QuestionLabel{【往年题1｜14级毒理考题汇总.xls｜口试】}\\
健康风险评价的步骤和内容（题号65，备注：第一题老师说还有危险管理和危险交流）

\end{pastquestion}

\begin{pastquestion}

\QuestionLabel{【往年题2｜19级毒理期末口试回忆.xlsx｜口试】}\\
阈值法和非阈值法（序号62，备注：还问了不确定性的来源------呼应不确定性系数 A/B/S/L/D/MF）

\end{pastquestion}

\PointSeparator

\chapter*{附：实验操作题与未完全匹配到课件的题目}\label{users__fpb__26_spring__pkusph_study__tmp__pdfs__prevention-medicine-past-exams__02-ux6bd2ux7406ux5b66ux5f80ux5e74ux9898ux8003ux70b9ux6574ux7406.md__ux9644ux5b9eux9a8cux64cdux4f5cux9898ux4e0eux672aux5b8cux5168ux5339ux914dux5230ux8bfeux4ef6ux7684ux9898ux76ee}
\addcontentsline{toc}{chapter}{附：实验操作题与未完全匹配到课件的题目}

\begin{quote}
以下题目（主要为实验操作题及个别选择题）在 14
节课件知识点正文中未找到完全对应的原文表述，按要求保留并根据题意暂归入最合适章节，单独建立考点条目。
\end{quote}

\section{实验操作题：小鼠经口灌胃及受试物配液计算（1\% 原液配制 50 mg/kg
剂量）（未完全匹配到）}\label{users__fpb__26_spring__pkusph_study__tmp__pdfs__prevention-medicine-past-exams__02-ux6bd2ux7406ux5b66ux5f80ux5e74ux9898ux8003ux70b9ux6574ux7406.md__ux5b9eux9a8cux64cdux4f5cux9898ux5c0fux9f20ux7ecfux53e3ux704cux80c3ux53caux53d7ux8bd5ux7269ux914dux6db2ux8ba1ux7b971-ux539fux6db2ux914dux5236-50-mgkg-ux5242ux91cfux672aux5b8cux5168ux5339ux914dux5230}

\InfoLabel{考频} 30 余次（各年实验操作题反复出现）

\InfoLabel{对应小节} 未完全匹配到课件知识点正文；根据题意暂归入：第二章 /《2
毒理学实验基础》（七、实验动物的染毒和处置------经口（胃肠道）染毒：灌胃）

\InfoLabel{匹配依据} 该题为实验操作技能（灌胃手法、回抽看负压、落空感，及''1\%原液配制50mg/kg剂量''的配液计算
c=d/v、倍比稀释），在课件中仅有''经口（胃肠道）染毒：灌胃、喂饲和吞咽胶囊''的简略表述，未给出操作细节与计算方法，故标记为未完全匹配到，归入实验基础章节。

\ContentBand{知识点原文摘取}

（《2
毒理学实验基础》第41--42页）（二）染毒途径：1.经口（胃肠道）染毒：灌胃、喂饲和吞咽胶囊。其余操作细节（称重、计算
ml 数、抽、灌、回抽看负压、落空感、染毒后继续禁食2\textasciitilde4小时）及配液计算未完全匹配到课件原文。

\InfoLabel{知识来源} 《2 毒理学实验基础》第41页（操作细节与配液计算未完全匹配到）

\ContentBand{对应往年题}

\begin{pastquestion}

\QuestionLabel{【往年题1｜07级毒理口试题回忆.doc｜实验操作】}\\
灌胃（称重，算ml数，抽，灌，问几个很简单的问题）（1号题）；1\%的原液配置50mg/kg
剂量的实验试剂（2号、3号、5号、26号、45号、58号题，答出计算过程，倍比稀释）；小鼠灌胃（考点：落空感、回抽）（25号、35号题）

\end{pastquestion}

\begin{pastquestion}

\QuestionLabel{【往年题2｜09级口试题目.docx｜实验操作】}\\
计算（都是那个1\%的题）（2）；剂量计算，1\%原液配制50mg/kg剂量（3、5、24、53）；演示小鼠灌胃操作（7、9、12、23、46）；药物配比（26）；1\%配5mg/kg（13）

\end{pastquestion}

\begin{pastquestion}

\QuestionLabel{【往年题3｜21级毒理期末口试回忆.xlsx｜实验操作】}\\
灌胃（序号3、11、19、43等大量序号，备注：灌胃口头询问剂量；灌胃操作第一步需要计算受试物配置）

\end{pastquestion}

\begin{pastquestion}

\QuestionLabel{【往年题4｜19级毒理期末口试回忆.xlsx｜实验操作】}\\
灌胃（序号51、84、27等大量序号，备注：灌胃要计算浓度 c=d/v，操作时记得回抽看负压再灌）

\end{pastquestion}

\begin{pastquestion}

\QuestionLabel{【往年题5｜17级毒理试题汇总.xlsx｜实验操作】}\\
灌胃 / 小鼠灌胃（题号78、62、34等大量题号）

\end{pastquestion}

\begin{pastquestion}

\QuestionLabel{【往年题6｜15级毒理学考试题目汇总.xls｜实验操作】}\\
灌胃（包括称重和灌胃，记得0.1ml/10g）（题号10、57、61等大量题号）

\end{pastquestion}

\begin{pastquestion}

\QuestionLabel{【往年题7｜14级毒理考题汇总.xls｜实验操作】}\\
灌胃（题号14、22、57等大量题号）

\end{pastquestion}

\begin{pastquestion}

\QuestionLabel{【往年题8｜20级完整试题回忆.docx｜实验操作】}\\
请演示经口灌胃操作（序号76）

\end{pastquestion}

\PointSeparator

\section{实验操作题：小鼠采血技术（内眦/眼眶静脉取血、鼠尾取血、内眦静脉丛取血）（未完全匹配到）}\label{users__fpb__26_spring__pkusph_study__tmp__pdfs__prevention-medicine-past-exams__02-ux6bd2ux7406ux5b66ux5f80ux5e74ux9898ux8003ux70b9ux6574ux7406.md__ux5b9eux9a8cux64cdux4f5cux9898ux5c0fux9f20ux91c7ux8840ux6280ux672fux5185ux7726ux773cux7736ux9759ux8109ux53d6ux8840ux9f20ux5c3eux53d6ux8840ux5185ux7726ux9759ux8109ux4e1bux53d6ux8840ux672aux5b8cux5168ux5339ux914dux5230}

\InfoLabel{考频} 30 余次（各年实验操作题反复出现）

\InfoLabel{对应小节} 未完全匹配到课件知识点正文；根据题意暂归入：第二章 /《2
毒理学实验基础》（七、实验动物处死及生物标本采集------血液采集）

\InfoLabel{匹配依据}
该题为采血实验操作技能（内眦取血需将老鼠提起、沿玻璃管滴血、棉球按压止血；鼠尾取血酒精消毒后干棉球擦干、形成大血滴、转移到含抗凝剂EP管、压迫止血），课件仅列''血液采集''条目，无操作步骤细节，故标记为未完全匹配到。

\ContentBand{知识点原文摘取}

（《2
毒理学实验基础》第47页）（三）实验动物处死及生物标本采集：1.实验动物处死方法；2.血液采集；3.尿液采集；4.病理解剖和标本留取；5.大体解剖\ldots\ldots 具体采血操作（内眦/眼眶静脉取血、鼠尾取血、内眦静脉丛取血）的步骤与注意事项未完全匹配到课件原文。

\InfoLabel{知识来源} 《2 毒理学实验基础》第47页（采血操作步骤未完全匹配到）

\ContentBand{对应往年题}

\begin{pastquestion}

\QuestionLabel{【往年题1｜07级毒理口试题回忆.doc｜实验操作】}\\
内眦静脉取血（1号、18号、36号、39号、61号题）；鼠尾取血10ul（27号题，注意第一滴血要弃掉，堵住胶头滴管上的洞）；鼠尾取血（31号、33号、49号、51号、57号题）

\end{pastquestion}

\begin{pastquestion}

\QuestionLabel{【往年题2｜09级口试题目.docx｜实验操作】}\\
内眦取血（11、16、34、55）；鼠尾取血（21、27、31、33、48）；小鼠内眦静脉丛取血（41）；小鼠内眦静脉取血（39，备注：记得用棉花球按压鼠眼止血）

\end{pastquestion}

\begin{pastquestion}

\QuestionLabel{【往年题3｜21级毒理期末口试回忆.xlsx｜实验操作】}\\
鼠尾采血 / 内眦取血 /
内眦采血（序号8、10、17等大量序号，备注：鼠尾采血最好边做边讲每步作用，最后将取的血转移到有抗凝剂的ep管中，干棉球按压止血；内眦取血需把老鼠提起来沿玻璃管滴血）

\end{pastquestion}

\begin{pastquestion}

\QuestionLabel{【往年题4｜19级毒理期末口试回忆.xlsx｜实验操作】}\\
内眦取血 / 鼠尾取血（序号41、86、97等大量序号，备注：取完血记得用棉球按压止血）

\end{pastquestion}

\begin{pastquestion}

\QuestionLabel{【往年题5｜17级毒理试题汇总.xlsx｜实验操作】}\\
内眦取血 / 鼠尾取血（题号61、68、85等大量题号）

\end{pastquestion}

\begin{pastquestion}

\QuestionLabel{【往年题6｜15级毒理学考试题目汇总.xls｜实验操作】}\\
眼眶取血 / 内眦取血 /
鼠尾采血（题号5、29、47等大量题号，备注：鼠尾取血要形成大血滴再取；酒精消毒后用干棉球擦干）

\end{pastquestion}

\begin{pastquestion}

\QuestionLabel{【往年题7｜14级毒理考题汇总.xls｜实验操作】}\\
内眦静脉取血 / 鼠尾取血（题号52、80、55等大量题号）

\end{pastquestion}

\begin{pastquestion}

\QuestionLabel{【往年题8｜20级完整试题回忆.docx｜实验操作】}\\
内眦取血（其他同学考题）

\end{pastquestion}

\begin{pastquestion}

\QuestionLabel{【往年题9｜毒理学-来自课本的题库补充byYYS.docx｜实验操作】}\\
球内眦取血，需要把老鼠提起来，沿着玻璃管将血滴到平板上

\end{pastquestion}

\PointSeparator

\section{实验操作题：小鼠骨髓取样/微核制片（取股骨、小牛血清用量）（未完全匹配到）}\label{users__fpb__26_spring__pkusph_study__tmp__pdfs__prevention-medicine-past-exams__02-ux6bd2ux7406ux5b66ux5f80ux5e74ux9898ux8003ux70b9ux6574ux7406.md__ux5b9eux9a8cux64cdux4f5cux9898ux5c0fux9f20ux9aa8ux9ad3ux53d6ux6837ux5faeux6838ux5236ux7247ux53d6ux80a1ux9aa8ux5c0fux725bux8840ux6e05ux7528ux91cfux672aux5b8cux5168ux5339ux914dux5230}

\InfoLabel{考频} 多次（各年实验操作题，常与微核试验合并）

\InfoLabel{对应小节} 未完全匹配到课件操作细节；知识点对应：第八章 /《8
遗传毒理学（二）》第22页''小鼠骨髓多染红细胞微核试验''

\InfoLabel{匹配依据}
微核试验的原理与遗传学终点已在第八章考点5中收录；本条专指''小鼠骨髓取法/制片''实验操作细节（取股骨、股骨头和膝关节处操作、抽多少
ml 小牛血清、推片方向、评价片子好坏），课件无操作步骤原文，故标记为未完全匹配到。

\ContentBand{知识点原文摘取}

（《8 遗传毒理学（二）》第22页）小鼠骨髓多染红细胞微核试验（Micronucleus assay in bone marrow polychromatic
erythrocytes）。具体骨髓取样、制片操作（小牛血清用量、取股骨方法、推片方向、片子优劣评价）未完全匹配到课件原文。

\InfoLabel{知识来源} 《8 遗传毒理学（二）》第22页（骨髓取样/制片操作细节未完全匹配到）

\ContentBand{对应往年题}

\begin{pastquestion}

\QuestionLabel{【往年题1｜07级毒理口试题回忆.doc｜实验操作】}\\
小鼠骨髓取法，制片（14号题）；微核骨髓采样、制片（44号题，老师会问小牛血清的剂量，怎样评价涂片的好坏）

\end{pastquestion}

\begin{pastquestion}

\QuestionLabel{【往年题2｜09级口试题目.docx｜实验操作】}\\
小鼠骨髓取样、涂片制作（40，备注：吴双老师问小牛血清的用量以及取股骨时股骨头和膝关节处的操作方法）

\end{pastquestion}

\begin{pastquestion}

\QuestionLabel{【往年题3｜14级毒理考题汇总.xls｜实验操作】}\\
微核制片（题号55备注：实验的孟老师\ldots\ldots 一开始就问要取多少血清，中间问推片手法注意事项和冲洗骨髓时冲干净的标志，最后让你评价自己的片子好不好）

\end{pastquestion}

\begin{pastquestion}

\QuestionLabel{【往年题4｜21级毒理期末口试回忆.xlsx｜实验操作】}\\
微核制片（序号94备注：微核制片会问抽多少ml的牛血清）

\end{pastquestion}

\PointSeparator

\section{蓄积作用（accumulation）的含义、分类及毒理学意义（未完全匹配到）}\label{users__fpb__26_spring__pkusph_study__tmp__pdfs__prevention-medicine-past-exams__02-ux6bd2ux7406ux5b66ux5f80ux5e74ux9898ux8003ux70b9ux6574ux7406.md__ux84c4ux79efux4f5cux7528accumulationux7684ux542bux4e49ux5206ux7c7bux53caux6bd2ux7406ux5b66ux610fux4e49ux672aux5b8cux5168ux5339ux914dux5230}

\InfoLabel{考频} 6 次

\InfoLabel{对应小节} 未完全匹配到课件知识点正文；根据题意暂归入：第六章 /《6
亚慢性和慢性毒性试验》第6页（并见《13 安全性评价》第6页 蓄积试验）

\InfoLabel{匹配依据}
该题在课件原文中未找到完整对应内容，课件仅在《6》第6页以英文列出''Accumulation------materials
accumulation、effect
accumulation''，《13》第6页第二阶段含''蓄积试验''；根据题干关键词和知识逻辑（物质蓄积/功能蓄积及其在重复暴露毒性、安全性评价中的意义），推测属于亚慢性和慢性毒性章节内容。

\ContentBand{知识点原文摘取}

（《6 亚慢性和慢性毒性试验》第6页）Single exposure / Repeated exposure：Accumulation------materials
accumulation、effect accumulation。（《13 安全性评价》第6页）四阶段毒理学安全性评价
第二阶段：蓄积试验、遗传毒理学试验和致畸试验。蓄积作用的完整分类（物质蓄积、功能蓄积）及其毒理学意义未完全匹配到课件原文。

【翻译】单次暴露和重复暴露的区别之一是蓄积作用，包括物质蓄积和效应蓄积。

\InfoLabel{知识来源} 《6 亚慢性和慢性毒性试验》第6页；《13 安全性评价》第6页（完整分类与意义未完全匹配到）

\ContentBand{对应往年题}

\begin{pastquestion}

\QuestionLabel{【往年题1｜21级毒理期末口试回忆.xlsx｜口试】}\\
accumulation的含义，分类及意义（序号42）

\end{pastquestion}

\begin{pastquestion}

\QuestionLabel{【往年题2｜19级毒理期末口试回忆.xlsx｜口试】}\\
accumulation的意义及分类，意义（序号86）；accumulation，分类，毒理学意义（序号忘了/姚碧云）

\end{pastquestion}

\begin{pastquestion}

\QuestionLabel{【往年题3｜17级毒理试题汇总.xlsx｜口试】}\\
什么是accumulation,分几类，毒理学意义是什么（题号42、55）

\end{pastquestion}

\begin{pastquestion}

\QuestionLabel{【往年题4｜20级完整试题回忆.docx｜口试】}\\
蓄积作用的分类及毒理学意义（其他同学考题A）

\end{pastquestion}

\begin{pastquestion}

\QuestionLabel{【往年题5｜20级完整试题回忆.docx｜单项选择】}\\
发生慢性毒性的基础是 A选择毒性 B局部毒性 C特殊毒性 D急性毒性 E蓄积作用（第8题）

\end{pastquestion}

\PointSeparator

\section{卫生毒理学的概念（单项选择题）（未完全匹配到）}\label{users__fpb__26_spring__pkusph_study__tmp__pdfs__prevention-medicine-past-exams__02-ux6bd2ux7406ux5b66ux5f80ux5e74ux9898ux8003ux70b9ux6574ux7406.md__ux536bux751fux6bd2ux7406ux5b66ux7684ux6982ux5ff5ux5355ux9879ux9009ux62e9ux9898ux672aux5b8cux5168ux5339ux914dux5230}

\InfoLabel{考频} 1 次

\InfoLabel{对应小节} 未完全匹配到课件原文；根据题意暂归入：第一章 /《1 绪论》（毒理学定义及范畴）

\InfoLabel{匹配依据} 该题在课件原文中未找到''卫生毒理学''专门定义；课件《1
绪论》给出的是''毒理学''的定义与范畴（描述/机制/管理）。根据题干关键词（卫生毒理学属预防医学范畴、是毒理学分支学科、研究生活和生产活动中外来化合物对机体损害及机制），推测属于绪论章节内容。

\ContentBand{知识点原文摘取}

未完全匹配到课件原文。（最接近的课件表述：《1
绪论》第4页''毒理学：从生物学角度研究外源因素和内源因素对生物体和生态系统的损害作用与机制，并进行安全性评价、危险度评定/风险评估、危险性管理与交流的科学''；第10页基础毒理学与环境和职业毒理学、食品毒理学、药物毒理学等分支的关系。）

\InfoLabel{知识来源} 未完全匹配到（最接近《1 绪论》第4页、第10页）

\ContentBand{对应往年题}

\begin{pastquestion}

\QuestionLabel{【往年题1｜20级完整试题回忆.docx｜单项选择】}\\
关于卫生毒理学的概念，不正确的是 A、属于预防医学的范畴 B、是毒理学的一个分支学科
C、研究生活和生产活动中可能接触的外来化合物对机体的损害作用及其机制 D、人体用药安全是其主要研究内容之一
E、环境毒理学、食品毒理学和工业毒理学与卫生毒理学关系最密切（第1题）

\end{pastquestion}

\PointSeparator

\section{关于毒性的说法及生物学标志（单项选择题，零散考点归并）（未完全匹配到/部分匹配）}\label{users__fpb__26_spring__pkusph_study__tmp__pdfs__prevention-medicine-past-exams__02-ux6bd2ux7406ux5b66ux5f80ux5e74ux9898ux8003ux70b9ux6574ux7406.md__ux5173ux4e8eux6bd2ux6027ux7684ux8bf4ux6cd5ux53caux751fux7269ux5b66ux6807ux5fd7ux5355ux9879ux9009ux62e9ux9898ux96f6ux6563ux8003ux70b9ux5f52ux5e76ux672aux5b8cux5168ux5339ux914dux5230ux90e8ux5206ux5339ux914d}

\InfoLabel{考频} 3 次

\InfoLabel{对应小节} 部分匹配：第九章《9 化学致癌作用》第3页''毒性-毒效应''、第一章《1
绪论》第18页''毒性''、第29--30页''生物学标志''

\InfoLabel{匹配依据}
20级选择题''关于毒性的说法错误的是''、``生物学标志包括''、``3Rs原则是指''等零散选择题，知识点已分别在第一章/第九章相关考点收录，此处归并保留题干原文，确保不遗漏。

\ContentBand{知识点原文摘取}

（《9
化学致癌作用》第3页）毒性（toxicity）：化学物引起有害作用的固有的能力。毒性是物质一种内在的、不变的分子性质，取决于其化学结构；毒效应（toxic
effect）：化学物对机体健康引起的有害作用。（《1 绪论》第29--30页 生物学标志含接触、效应、易感性三类；第39页
3Rs：减少、替代、优化。）

\InfoLabel{知识来源} 《9 化学致癌作用》第3页；《1 绪论》第18页；《1 绪论》第29页；《1 绪论》第30页；《1
绪论》第39页

\ContentBand{对应往年题}

\begin{pastquestion}

\QuestionLabel{【往年题1｜20级完整试题回忆.docx｜单项选择】}\\
关于毒性的说法错误的是 A毒性是化学物引起有害作用的能力 B毒性很强说明需要很大剂量才能显示毒性
C\&D\&E\ldots\ldots（第16题）

\end{pastquestion}

\begin{pastquestion}

\QuestionLabel{【往年题2｜20级完整试题回忆.docx｜单项选择】}\\
生物学标志包括（第2题）；3Rs原则是指（第3题）

\end{pastquestion}

\begin{pastquestion}

\QuestionLabel{【往年题3｜20级完整试题回忆.docx｜单项选择】}\\
18、\ldots\ldots（试题回忆原文残缺，仅记录题号，按要求保留不予删除）

\end{pastquestion}

\PointSeparator

\SetSubjectAccent{B16A2B}{流行病学}

\part{流行病学}

\chapter*{说明}\label{users__fpb__26_spring__pkusph_study__tmp__pdfs__prevention-medicine-past-exams__03-ux6d41ux884cux75c5ux5b66ux5f80ux5e74ux9898ux8003ux70b9ux6574ux7406.md__ux8bf4ux660e}
\addcontentsline{toc}{chapter}{说明}

\begin{itemize}
\tightlist
\item
  本文件根据《流行病学》第9版课本知识点（24 章 Markdown 文件）和大礼包所有的历年往年题文件整理。
\item
  知识点正文均摘自 24 章课本知识点 Markdown 文件。
\item
  往年题均摘自用户提供的往年题文件夹（04、05、07 级回忆，10、11、12 级口试题库，14、17、19、20、21
  级口试回忆）。
\item
  最终排序依据为真实授课顺序，而不是课本 24 章原始顺序。
\item
  同一授课章节内，考点按考频从高到低排列；同频时，按课本原始出现顺序排列。
\item
  一道考题若同时考查多个知识点，分别归入对应知识点，并各计 1 次。
\item
  仅纳入在往年题中匹配到的知识点；如果某个授课章节未在往年题中匹配到明确考点，则保留该章节标题并标注``本章未在提供的往年题中匹配到明确考点''。
\item
  对于无法明确匹配但高度相关的题目，在``匹配依据''中标注``疑似对应''。
\item
  对于课本知识点文件中未检索到明确原文的题目，不自行补充原文，只保留题目并说明未检索到。
\end{itemize}

\textbf{关于考频的统计口径（重要）：} 10---21
级多为``口试回忆题库''，同一份文件（同一年级）内由多位同学回忆，同一道题卡会被反复记录。为避免重复计数失真，本文件约定：\textbf{考频
= 该考点在所提供的 11
个年级文件（04/05/07/10/11/12/14/17/19/20/21）中出现的年级数；同一年级文件内的重复回忆合并计 1
次（即``按届计数''）。} 因此考频最大为
11。口试题库中同一考点在同一年级内通常出现多次，下方``对应往年题''仅摘取代表性题干，并标注所属年级文件。

\PointSeparator

\chapter{绪论}\label{users__fpb__26_spring__pkusph_study__tmp__pdfs__prevention-medicine-past-exams__03-ux6d41ux884cux75c5ux5b66ux5f80ux5e74ux9898ux8003ux70b9ux6574ux7406.md__ux7b2cux4e00ux7ae0-ux7eeaux8bba}

\section{流行病学的定义}\label{users__fpb__26_spring__pkusph_study__tmp__pdfs__prevention-medicine-past-exams__03-ux6d41ux884cux75c5ux5b66ux5f80ux5e74ux9898ux8003ux70b9ux6574ux7406.md__ux6d41ux884cux75c5ux5b66ux7684ux5b9aux4e49}

\FrequencyPill{2 次}

\InfoLabel{对应小节} 第一章 绪论 / 第二节 流行病学的定义 / 一、流行病学定义的演变（书中页码 第5页）

\InfoLabel{匹配依据} 04 级名解直接考``流行病学（英文）''；20/21 级题库总论部分将``绪论''列为必背范围。题干考查
epidemiology 的规范定义。

\ContentBand{知识点原文摘取}

自预防医学专业规划教材《流行病学》第3版起，我国采用如下定义：流行病学是研究人群中疾病与健康状况的分布及其影响因素，并研究防制疾病及促进健康的策略和措施的科学（连志浩，1992）。

1983年，Last主编的《流行病学词典》中定义：``流行病学研究人群中与健康有关的状态或事件的分布及其决定因素，并应用这些研究来维持和促进健康。''

``流行病学''一词的英文源自希腊语，由epi（在\ldots\ldots 之中、之上）、demos（人群）和logos（研究）组成，直译为``研究人群中发生之事的学问''。

\ContentBand{对应往年题}

\begin{pastquestion}

\QuestionLabel{【往年题1｜04级流病考题回忆.docx｜名解】}\\
流行病学（英文）

\end{pastquestion}

\begin{pastquestion}

\QuestionLabel{【往年题2｜20级回忆试题库.xlsx｜总论复习范围】}\\
总论：绪论、三间分布、描述性、队列、病例对照、RCT、筛检、病因与因果推断、偏倚及其控制

\end{pastquestion}

\PointSeparator

\chapter{疾病分布}\label{users__fpb__26_spring__pkusph_study__tmp__pdfs__prevention-medicine-past-exams__03-ux6d41ux884cux75c5ux5b66ux5f80ux5e74ux9898ux8003ux70b9ux6574ux7406.md__ux7b2cux4e8cux7ae0-ux75beux75c5ux5206ux5e03}

\section{横断面分析与出生队列分析（疾病年龄分布的两种分析方法）}\label{users__fpb__26_spring__pkusph_study__tmp__pdfs__prevention-medicine-past-exams__03-ux6d41ux884cux75c5ux5b66ux5f80ux5e74ux9898ux8003ux70b9ux6574ux7406.md__ux6a2aux65adux9762ux5206ux6790ux4e0eux51faux751fux961fux5217ux5206ux6790ux75beux75c5ux5e74ux9f84ux5206ux5e03ux7684ux4e24ux79cdux5206ux6790ux65b9ux6cd5}

\FrequencyPill{7 次}

\InfoLabel{对应小节} 第二章 疾病的分布 / 第三节 疾病的分布特征 / 一、人群分布 /（一）年龄 /
2.疾病年龄分布的分析方法（书中页码 第18---19页）

\InfoLabel{匹配依据} 10、11、12、17、19、20、21
级反复出现``横断面分析（研究）和出生队列分析（研究）的异同点/优点''\,``横断面分析的缺点（肺癌案例）应该用什么方法''等简答与综合分析题，直接对应课本两种年龄分布分析方法及其优缺点。

\ContentBand{知识点原文摘取}

（1）横断面分析（cross-sectional
analysis）：主要分析同一时期不同年龄组或不同时期同一年龄组的发病率、患病率或死亡率的变化，多用于某时期传染病或潜伏期较短疾病的年龄分布分析。对于慢性病，由于暴露时间距发病时间可能很长，致病因素在不同时间的强度也可能发生变化，横断面分析难以正确显示疾病与年龄的关系。

（2）出生队列分析（birth cohort analysis）：同一时期出生的一组人群称为出生队列（birth
cohort），对其随访若干年，以观察发病情况，这种利用出生队列资料将疾病年龄分布和时间分布结合描述的方法称出生队列分析。该方法在评价疾病年龄分布的长期趋势及提供病因线索等方面发挥重要作用。它可以直观地呈现致病因素与年龄的关系，有助于探索年龄、所处时代暴露特点及经历在疾病频率变化中的作用。相比横断面分析，出生队列分析更合理地展示了年龄与发病的关系，有助于正确分辨出年龄、时间、暴露经历对疾病的作用。

\ContentBand{对应往年题}

\begin{pastquestion}

\QuestionLabel{【往年题1｜10级考题回忆.pdf｜简答】}\\
横断面研究和出生队列研究异同点

\end{pastquestion}

\begin{pastquestion}

\QuestionLabel{【往年题2｜10级考题回忆.pdf｜大题】}\\
美国肺癌死亡率的年龄分布特点：横断面分析和出生队列分析的对比和各自的作用

\end{pastquestion}

\begin{pastquestion}

\QuestionLabel{【往年题3｜17级流病口试回忆.xlsx｜简答/大题（题号3、1）】}\\
横断面分析和出生队列分析的异同；横断面分析的缺点（给了一个肺癌案例），应该用什么分析方法（出生队列分析），为什么？

\end{pastquestion}

\begin{pastquestion}

\QuestionLabel{【往年题4｜19级流行病学期末口试回忆.xlsx｜简答（题号3）】}\\
出生队列和横断面异同

\end{pastquestion}

\begin{pastquestion}

\QuestionLabel{【往年题5｜21级流病考题回忆汇总.xlsx｜简答（题号8、11、48）】}\\
研究疾病的时间分布的两种方法，横断面研究和出生队列分析的异同点；出生队列分析的作用；出生队列研究与横断面研究相比有什么优点

\end{pastquestion}

\PointSeparator

\section{罹患率}\label{users__fpb__26_spring__pkusph_study__tmp__pdfs__prevention-medicine-past-exams__03-ux6d41ux884cux75c5ux5b66ux5f80ux5e74ux9898ux8003ux70b9ux6574ux7406.md__ux7f79ux60a3ux7387}

\FrequencyPill{7 次}

\InfoLabel{对应小节} 第二章 疾病的分布 / 第一节 疾病频率测量指标 / 一、发病频率指标 /（二）罹患率（书中页码
第12页）

\InfoLabel{匹配依据} 04、07、10、11、12、19、20 级均考``罹患率 / attack
rate''名词解释，部分追问``与发病率的区别、适用场景''。

\ContentBand{知识点原文摘取}

罹患率（attack
rate）通常是指在某一局限范围短时间内的发病率，也是测量某人群某病新病例发生频率的指标。其计算公式与发病率相同，但它的观察时间较短，常以日、周、旬、月为单位，使用比较灵活。它的优点是能根据暴露程度较精确地测量发病频率，在食物中毒、职业中毒或传染病的暴发及流行中，经常使用该指标。

\[罹患率=\frac{观察期间发病人数}{同期暴露人口数}\times K\]

\ContentBand{对应往年题}

\begin{pastquestion}

\QuestionLabel{【往年题1｜04级流病考题回忆.docx｜名解】}\\
罹患率

\end{pastquestion}

\begin{pastquestion}

\QuestionLabel{【往年题2｜11级流行病学口试考题回忆.pdf｜名解】}\\
attack rate

\end{pastquestion}

\begin{pastquestion}

\QuestionLabel{【往年题3｜20级回忆试题库.xlsx｜名解（题号15）】}\\
attack rate（罹患率）

\end{pastquestion}

\PointSeparator

\section{续发率（二代发病率）}\label{users__fpb__26_spring__pkusph_study__tmp__pdfs__prevention-medicine-past-exams__03-ux6d41ux884cux75c5ux5b66ux5f80ux5e74ux9898ux8003ux70b9ux6574ux7406.md__ux7eedux53d1ux7387ux4e8cux4ee3ux53d1ux75c5ux7387}

\FrequencyPill{6 次}

\InfoLabel{对应小节} 第二章 疾病的分布 / 第一节 疾病频率测量指标 / 一、发病频率指标 /（三）续发率（书中页码
第12页）

\InfoLabel{匹配依据} 04、07、11、12、19、20 级均考``续发率 / secondary attack rate / 二代发病率''名词解释。

\ContentBand{知识点原文摘取}

续发率（secondary attack
rate，SAR）也称二代发病率，指在传染病最短潜伏期到最长潜伏期之间，易感接触者中发病人数占所有易感接触者总数的百分比。

\[续发率=\frac{潜伏期内易感接触者中发病人数}{易感接触者总人数}\times100\%\]

第一个病例发生后，在该病最短与最长潜伏期之间出现的病例称续发病例，又称二代病例。续发率可用于比较传染病传染力的强弱、分析传染病流行因素及评价卫生防疫措施的效果。

\ContentBand{对应往年题}

\begin{pastquestion}

\QuestionLabel{【往年题1｜11级流行病学口试考题回忆.pdf｜名解】}\\
续发率；二代发病率

\end{pastquestion}

\begin{pastquestion}

\QuestionLabel{【往年题2｜19级流行病学期末口试回忆.xlsx｜名解（题号12）】}\\
续发率

\end{pastquestion}

\begin{pastquestion}

\QuestionLabel{【往年题3｜20级回忆试题库.xlsx｜名解（题号21）】}\\
secondary attack rate（续发率）

\end{pastquestion}

\PointSeparator

\section{期间患病率（及患病率）}\label{users__fpb__26_spring__pkusph_study__tmp__pdfs__prevention-medicine-past-exams__03-ux6d41ux884cux75c5ux5b66ux5f80ux5e74ux9898ux8003ux70b9ux6574ux7406.md__ux671fux95f4ux60a3ux75c5ux7387ux53caux60a3ux75c5ux7387}

\FrequencyPill{7 次}

\InfoLabel{对应小节} 第二章 疾病的分布 / 第一节 疾病频率测量指标 / 二、患病频率指标 /（一）患病率（书中页码
第12---13页）

\InfoLabel{匹配依据} 04、05、10、12、14、17、20 级考``期间患病率 / period prevalence rate''名词解释；19、20
级另考``患病率 prevalence rate''。

\ContentBand{知识点原文摘取}

患病率（prevalence）也称现患率，是指某特定时间内总人口中某病新旧病例所占的比例。患病率可按观察时间的不同分为时点患病率和期间患病率。时点患病率的观察时间一般不超过一个月，而期间患病率所对应的是特定的一段时间，通常为几个月到几年。

\[时点患病率=\frac{某一时点某人群中某病新旧病例数}{该时点人口数}\times K\]

\[期间患病率=\frac{观察期间某人群中某病的新旧病例数}{同期的平均人口数}\times K\]

\ContentBand{对应往年题}

\begin{pastquestion}

\QuestionLabel{【往年题1｜04级流病考题回忆.docx｜名解】}\\
period prevalence rate

\end{pastquestion}

\begin{pastquestion}

\QuestionLabel{【往年题2｜12级流病口试考题回忆.pdf｜名解（题号12）】}\\
期间患病率

\end{pastquestion}

\begin{pastquestion}

\QuestionLabel{【往年题3｜17级流病口试回忆.xlsx｜名解（题号3）】}\\
期间患病率

\end{pastquestion}

\PointSeparator

\section{发病率}\label{users__fpb__26_spring__pkusph_study__tmp__pdfs__prevention-medicine-past-exams__03-ux6d41ux884cux75c5ux5b66ux5f80ux5e74ux9898ux8003ux70b9ux6574ux7406.md__ux53d1ux75c5ux7387}

\FrequencyPill{6 次}

\InfoLabel{对应小节} 第二章 疾病的分布 / 第一节 疾病频率测量指标 / 一、发病频率指标 /（一）发病率（书中页码
第11---12页）

\InfoLabel{匹配依据} 04、07、10、14、19、20 级考``发病率 / incidence
rate''名词解释；部分追问与患病率、死亡率、病死率的区别。

\ContentBand{知识点原文摘取}

发病率（incidence rate）是指一定期间内，一定范围人群中某病新发生病例出现的频率。

\[发病率=\frac{一定时期内某人群中某病新发病例数}{同期该人群暴露人口数}\times K\]

发病率是反映疾病流行强度的指标，衡量疾病对人群健康影响的程度，发病率高的疾病对人群健康的危害大。通过发病率的比较，可了解疾病流行特征，探讨病因因素，提出病因假说，评价防制措施的效果。

\ContentBand{对应往年题}

\begin{pastquestion}

\QuestionLabel{【往年题1｜07级考题回忆.docx｜名解】}\\
incidence rate（追问了患病率、死亡率、病死率）

\end{pastquestion}

\begin{pastquestion}

\QuestionLabel{【往年题2｜14级流病考题回忆汇总.xls｜名解（题号14）】}\\
发病率

\end{pastquestion}

\begin{pastquestion}

\QuestionLabel{【往年题3｜19级流行病学期末口试回忆.xlsx｜名解（题号3）】}\\
发病率

\end{pastquestion}

\PointSeparator

\section{死亡专率与死亡率（含死亡率、病死率的比较）}\label{users__fpb__26_spring__pkusph_study__tmp__pdfs__prevention-medicine-past-exams__03-ux6d41ux884cux75c5ux5b66ux5f80ux5e74ux9898ux8003ux70b9ux6574ux7406.md__ux6b7bux4ea1ux4e13ux7387ux4e0eux6b7bux4ea1ux7387ux542bux6b7bux4ea1ux7387ux75c5ux6b7bux7387ux7684ux6bd4ux8f83}

\FrequencyPill{7 次}

\InfoLabel{对应小节} 第二章 疾病的分布 / 第一节 疾病频率测量指标 / 三、死亡与生存频率指标
/（一）死亡率、（二）病死率（书中页码 第14---15页）

\InfoLabel{匹配依据} 05、10、11、12、17、20 级考``死亡专率 / specific mortality rate''名词解释；17、20、21
级简答考``病死率和死亡率的区别和联系''；多年计算题考粗死亡率、死亡专率、病死率。

\ContentBand{知识点原文摘取}

死亡率（mortality
rate）指在一定期间内，某人群中总死亡人数在该人群中所占的比例，是测量人群死亡危险最常用的指标。根据公式计算所得的死亡率也称粗死亡率（crude
death rate）。死亡率可按不同的人群属性（如年龄、性别、职业、民族）、疾病种类等特征分别计算，此即死亡专率。

\[死亡率=\frac{某人群某年总死亡人数}{该人群同年平均人口数}\times K\]

病死率（case fatality rate）表示一定时期内因某病死亡者占患该病人数的比例，表示某病病人因该病死亡的危险性。

\[病死率=\frac{某时期内因某病死亡人数}{同期患某病人数}\times100\%\]

病死率表示确诊某病者的死亡概率，它可反映疾病的严重程度，也可反映医疗水平和诊治能力，常用于急性传染病。死亡率用于衡量某一时期、某一地区人群死亡危险性的大小。

\ContentBand{对应往年题}

\begin{pastquestion}

\QuestionLabel{【往年题1｜12级流病口试考题回忆.pdf｜名解（题号6）】}\\
死亡专率

\end{pastquestion}

\begin{pastquestion}

\QuestionLabel{【往年题2｜17级流病口试回忆.xlsx｜简答（题号10）】}\\
病死率和死亡率的区别和联系（追问病程稳定时能否换算）

\end{pastquestion}

\begin{pastquestion}

\QuestionLabel{【往年题3｜20级回忆试题库.xlsx｜计算（题号22）】}\\
某市某年有10万人，该年内死亡1000人，有心脏病患者300人（男200女100），心脏病患者死亡数为60（男50女10）。计算死亡率、心脏病死亡专率、病死率、男性心脏病死亡率、心脏病归因死亡率五个数值。

\end{pastquestion}

\PointSeparator

\section{长期趋势（长期变动）}\label{users__fpb__26_spring__pkusph_study__tmp__pdfs__prevention-medicine-past-exams__03-ux6d41ux884cux75c5ux5b66ux5f80ux5e74ux9898ux8003ux70b9ux6574ux7406.md__ux957fux671fux8d8bux52bfux957fux671fux53d8ux52a8}

\FrequencyPill{6 次}

\InfoLabel{对应小节} 第二章 疾病的分布 / 第三节 疾病的分布特征 / 三、时间分布 /（四）长期趋势（书中页码
第26---27页）

\InfoLabel{匹配依据} 04、10、11、12、14、20 级考``长期趋势 / secular trend / 疾病长期趋势''名词解释。

\ContentBand{知识点原文摘取}

长期趋势（secular trend）也称长期变动（secular
change），是指在一个比较长的时间内，通常为几年或几十年，疾病的临床表现、分布特征和流行强度等方面发生变化的现象。有些疾病可表现出经过几年或几十年发病率或死亡率持续上升或下降的趋势。长期趋势的主要原因有病因或致病因素的变化、病原体的变异、机体免疫状况的改变、医疗和防制水平的提高、报告及登记制度完善程度等。

\ContentBand{对应往年题}

\begin{pastquestion}

\QuestionLabel{【往年题1｜04级流病考题回忆.docx｜名解】}\\
疾病长期趋势

\end{pastquestion}

\begin{pastquestion}

\QuestionLabel{【往年题2｜12级流病口试考题回忆.pdf｜名解（题号16、47）】}\\
长期趋势

\end{pastquestion}

\begin{pastquestion}

\QuestionLabel{【往年题3｜20级回忆试题库.xlsx｜名解（题号3、26）】}\\
secular trend（长期趋势）

\end{pastquestion}

\PointSeparator

\section{大流行（疾病流行强度：散发、暴发、流行、大流行）}\label{users__fpb__26_spring__pkusph_study__tmp__pdfs__prevention-medicine-past-exams__03-ux6d41ux884cux75c5ux5b66ux5f80ux5e74ux9898ux8003ux70b9ux6574ux7406.md__ux5927ux6d41ux884cux75beux75c5ux6d41ux884cux5f3aux5ea6ux6563ux53d1ux66b4ux53d1ux6d41ux884cux5927ux6d41ux884c}

\FrequencyPill{4 次}

\InfoLabel{对应小节} 第二章 疾病的分布 / 第二节 疾病流行强度 /
一、散发、二、暴发、三、流行、四、大流行（书中页码 第17页）

\InfoLabel{匹配依据} 04、11、12、17 级考``大流行（pandemic）/ 散发''等流行强度名词解释；07 级考``散发''。

\ContentBand{知识点原文摘取}

散发（sporadic）指发病率呈历年的一般水平，各病例间在发病时间和地点上无明显联系，表现为散在发生的现象。

暴发（outbreak）是指局部地区或集体单位，短时间内突然出现很多症状相似病人的现象。大多数病人在疾病的最短潜伏期到最长潜伏期之间发病，通常具有相同的传染源或传播途径。

流行（epidemic）是指在某地区某病的发病率显著超过该病历年发病率水平的现象，各病例之间呈现明显的时间和空间联系。

大流行（pandemic）是指某病发病率显著超过该病历年发病率水平，疾病蔓延迅速，涉及地区广，在短期内可跨越省界、国界甚至洲界的现象。

\ContentBand{对应往年题}

\begin{pastquestion}

\QuestionLabel{【往年题1｜04级流病考题回忆.docx｜名解】}\\
大流行；流行

\end{pastquestion}

\begin{pastquestion}

\QuestionLabel{【往年题2｜11级流行病学口试考题回忆.pdf｜名解】}\\
大流行，匹配

\end{pastquestion}

\begin{pastquestion}

\QuestionLabel{【往年题3｜17级流病口试回忆.xlsx｜名解（题号1）】}\\
pandemic（大流行）

\end{pastquestion}

\PointSeparator

\section{移民流行病学}\label{users__fpb__26_spring__pkusph_study__tmp__pdfs__prevention-medicine-past-exams__03-ux6d41ux884cux75c5ux5b66ux5f80ux5e74ux9898ux8003ux70b9ux6574ux7406.md__ux79fbux6c11ux6d41ux884cux75c5ux5b66}

\FrequencyPill{6 次}

\InfoLabel{对应小节} 第二章 疾病的分布 / 第三节 疾病的分布特征 / 四、疾病分布的综合描述
/（二）移民流行病学（书中页码 第27---29页）

\InfoLabel{匹配依据} 04、05、07、10、12、21
级考``移民流行病学''名词解释及``举例说明移民流行病学（在病因研究中的作用）''问答/综合分析题，部分追问二代移民发病情况变化。

\ContentBand{知识点原文摘取}

移民流行病学（migrant
epidemiology）是通过观察疾病在移民、移居地当地居民及原居地人群间的发病率或死亡率的差异，探讨疾病的发生与遗传因素或环境因素关系的一种观察性研究方法。移民流行病学常用于慢性非传染性疾病及某些遗传病的病因和流行因素的探讨。移民流行病学研究应遵循的原则如下：

\begin{enumerate}
\def\labelenumi{\arabic{enumi}.}
\item
  若某病发病率或死亡率的差别主要是环境因素作用的结果，则该病在移民人群中的发病率或死亡率与原居国（地区）人群不同，而接近移居国（地区）当地人群的发病率或死亡率。
\item
  若该病发病率或死亡率的差别主要与遗传因素有关，则移民人群与原居国（地区）人群的发病率或死亡率近似，而不同于移居国（地区）当地人群。
\end{enumerate}

移民在新环境居住的世代数越多，越接近移居地居民的发病水平。

\ContentBand{对应往年题}

\begin{pastquestion}

\QuestionLabel{【往年题1｜04级流病考题回忆.docx｜问答】}\\
举例说明移民流行病学

\end{pastquestion}

\begin{pastquestion}

\QuestionLabel{【往年题2｜07级考题回忆.docx｜名解】}\\
移民流行病学（追问二代移民发病情况变化）

\end{pastquestion}

\begin{pastquestion}

\QuestionLabel{【往年题3｜21级流病考题回忆汇总.xlsx｜综合分析（题号3）】}\\
移民流行病学

\end{pastquestion}

\PointSeparator

\section{影响疾病年龄分布的因素}\label{users__fpb__26_spring__pkusph_study__tmp__pdfs__prevention-medicine-past-exams__03-ux6d41ux884cux75c5ux5b66ux5f80ux5e74ux9898ux8003ux70b9ux6574ux7406.md__ux5f71ux54cdux75beux75c5ux5e74ux9f84ux5206ux5e03ux7684ux56e0ux7d20}

\FrequencyPill{1 次}

\InfoLabel{对应小节} 第二章 疾病的分布 / 第三节 疾病的分布特征 / 一、人群分布 /（一）年龄 /
1.疾病年龄分布特征（书中页码 第18页）

\InfoLabel{匹配依据} 21 级简答两次考``举例说明不同年龄疾病分布不同的原因 /
影响疾病年龄分布的因素''，对应课本急性传染病与慢性病的年龄分布特征及疫苗接种、致病因素变化对年龄模式的影响。

\ContentBand{知识点原文摘取}

（1）急性传染病一般有发病率随年龄的增长而下降的趋势。出生6个月内的婴儿因具有从母体获得的抗体，一般不易患传染病，但随着从母体获得的抗体逐渐减少或消失，儿童、青少年期易患某些急性呼吸道传染病如麻疹、百日咳、流行性腮腺炎等，至成年期免疫力增强，易感性逐渐降低。

（2）慢性病发病率常随年龄增长而上升。这是由于慢性病致病因素需要较长时间积累，才可至疾病发生，如大多数的恶性肿瘤和心血管疾病多发生于45岁以后。

（3）接种疫苗可改变一些急性传染病感染的年龄模式。如麻疹既往主要发生于幼儿及学龄前儿童，在推行免疫规划后，发病高峰后延。

（4）随着致病因素的变化，疾病的年龄分布也在动态变化。某些恶性肿瘤有年轻化趋势，如肺癌、乳腺癌等。

\ContentBand{对应往年题}

\begin{pastquestion}

\QuestionLabel{【往年题1｜21级流病考题回忆汇总.xlsx｜简答（题号13）】}\\
举例说明不同年龄疾病分布不同的原因

\end{pastquestion}

\begin{pastquestion}

\QuestionLabel{【往年题2｜21级流病考题回忆汇总.xlsx｜简答（题号39）】}\\
疾病年龄分布差异的原因

\end{pastquestion}

\PointSeparator

\section{潜在减寿年数（PYLL）}\label{users__fpb__26_spring__pkusph_study__tmp__pdfs__prevention-medicine-past-exams__03-ux6d41ux884cux75c5ux5b66ux5f80ux5e74ux9898ux8003ux70b9ux6574ux7406.md__ux6f5cux5728ux51cfux5bffux5e74ux6570pyll}

\FrequencyPill{2 次}

\InfoLabel{对应小节} 第二章 疾病的分布 / 第一节 疾病频率测量指标 / 四、疾病负担指标
/（一）潜在减寿年数（书中页码 第15---16页）

\InfoLabel{匹配依据} 19、21 级考``potential years of life lost（潜在损失寿命）/ 潜在寿命损失年 /
潜在损寿年数''名词解释。

\ContentBand{知识点原文摘取}

潜在减寿年数（potential years of life
lost，PYLL）是某病某年龄组人群死亡者的期望寿命与实际死亡年龄之差的总和，即死亡所造成的寿命损失。

\[PYLL=\sum_{i=1}^{e} a_i d_i\]

PYLL是疾病负担测量的一个直接指标，也是评价人群健康水平的一个重要指标，是在考虑死亡数量的基础上，以期望寿命为基准，进一步衡量死亡造成的寿命损失，强调了早死对人群健康的损害。

\ContentBand{对应往年题}

\begin{pastquestion}

\QuestionLabel{【往年题1｜19级流行病学期末口试回忆.xlsx｜名解（题号12）】}\\
potentialy years of life lost（潜在损失寿命）

\end{pastquestion}

\begin{pastquestion}

\QuestionLabel{【往年题2｜21级流病考题回忆汇总.xlsx｜名解（题号30）】}\\
潜在损寿年数

\end{pastquestion}

\PointSeparator

\section{感染率}\label{users__fpb__26_spring__pkusph_study__tmp__pdfs__prevention-medicine-past-exams__03-ux6d41ux884cux75c5ux5b66ux5f80ux5e74ux9898ux8003ux70b9ux6574ux7406.md__ux611fux67d3ux7387}

\FrequencyPill{1 次}

\InfoLabel{对应小节} 第二章 疾病的分布 / 第一节 疾病频率测量指标 / 二、患病频率指标 /（二）感染率（书中页码
第13---14页）

\InfoLabel{匹配依据} 21 级名解考``感染率''。

\ContentBand{知识点原文摘取}

感染率（prevalence of
infection）是指在某时间内受检人群中某病原体现有的感染人数所占的比例，通常用百分率表示。感染率的性质与患病率相似。

\[感染率=\frac{受检者中感染人数}{受检人数}\times100\%\]

感染率在流行病学工作中应用较广泛，特别是对隐性感染、病原携带者及轻型和不典型病例的调查较为常用。

\ContentBand{对应往年题}

\begin{pastquestion}

\QuestionLabel{【往年题1｜21级流病考题回忆汇总.xlsx｜名解（题号9）】}\\
感染率

\end{pastquestion}

\PointSeparator

\section{生存率}\label{users__fpb__26_spring__pkusph_study__tmp__pdfs__prevention-medicine-past-exams__03-ux6d41ux884cux75c5ux5b66ux5f80ux5e74ux9898ux8003ux70b9ux6574ux7406.md__ux751fux5b58ux7387}

\FrequencyPill{1 次}

\InfoLabel{对应小节} 第二章 疾病的分布 / 第一节 疾病频率测量指标 / 三、死亡与生存频率指标
/（三）生存率（书中页码 第15页）

\InfoLabel{匹配依据} 19 级名解考``生存率''。

\ContentBand{知识点原文摘取}

生存率（survival rate）指接受某种治疗的病人或某病病人中，经 \(n\) 年随访后，尚存活的病例数所占的比例。

\[生存率=\frac{随访满n年尚存活的病例数}{随访满n年的病例数}\times100\%\]

生存率反映疾病对生命的危害程度，可用于评价某些病程较长疾病的远期疗效，常用于癌症、心血管疾病、结核病等慢性疾病的预后研究。

\ContentBand{对应往年题}

\begin{pastquestion}

\QuestionLabel{【往年题1｜19级流行病学期末口试回忆.xlsx｜名解（题号20）】}\\
生存率

\end{pastquestion}

\PointSeparator

\chapter{描述性研究}\label{users__fpb__26_spring__pkusph_study__tmp__pdfs__prevention-medicine-past-exams__03-ux6d41ux884cux75c5ux5b66ux5f80ux5e74ux9898ux8003ux70b9ux6574ux7406.md__ux7b2cux4e09ux7ae0-ux63cfux8ff0ux6027ux7814ux7a76}

\section{描述性研究的种类与特征}\label{users__fpb__26_spring__pkusph_study__tmp__pdfs__prevention-medicine-past-exams__03-ux6d41ux884cux75c5ux5b66ux5f80ux5e74ux9898ux8003ux70b9ux6574ux7406.md__ux63cfux8ff0ux6027ux7814ux7a76ux7684ux79cdux7c7bux4e0eux7279ux5f81}

\FrequencyPill{6 次}

\InfoLabel{对应小节} 第三章 描述性研究 / 第一节 概述 / 二、种类；及第3页``主要特点''（书中页码 第30---32页）

\InfoLabel{匹配依据} 04、05、10、11、12、20
级反复考``描述性（流行病）研究的分类和特征/种类与特点''\,``举例说明描述性研究的特点''。注意 05、04
级标注``要答胡永华课件上的分类，书上的三分法不答''------以课程课件分类为准，课本原文列于下供参考。

\ContentBand{知识点原文摘取}

描述性研究常见的类型主要有：现况研究、病例报告、病例系列分析、个案研究、历史资料分析、随访研究、生态学研究等。

它的主要特点包括：\\
1.
描述性研究以观察为主要研究手段，不对研究对象采取任何干预措施，仅通过观察和收集相关资料，分析和总结研究对象或事件的特点。\\
2. 描述性研究中，其暴露因素的分配不是随机的，且在研究开始时一般不设立对照组。\\
3.
暴露与结局的时序关系无法确定，对于暴露与结局间关系的因果推断存在一定的局限性，仅可做一些初步的比较性分析，但可为后续的分析性或实验研究提供线索。

\ContentBand{对应往年题}

\begin{pastquestion}

\QuestionLabel{【往年题1｜05流病考题回忆.doc｜简答】}\\
描述性流行病学的分类和特征（要答胡永华课件上的分类，书上的三分法不答）

\end{pastquestion}

\begin{pastquestion}

\QuestionLabel{【往年题2｜10级考题回忆.pdf｜简答】}\\
描述性研究的特点和分类

\end{pastquestion}

\begin{pastquestion}

\QuestionLabel{【往年题3｜20级回忆试题库.xlsx｜简答（题号6、46、50）】}\\
描述研究的分类和特征；描述性研究特点；举例说明描述性研究的特点

\end{pastquestion}

\PointSeparator

\section{抽样调查（及随机抽样方法）}\label{users__fpb__26_spring__pkusph_study__tmp__pdfs__prevention-medicine-past-exams__03-ux6d41ux884cux75c5ux5b66ux5f80ux5e74ux9898ux8003ux70b9ux6574ux7406.md__ux62bdux6837ux8c03ux67e5ux53caux968fux673aux62bdux6837ux65b9ux6cd5}

\FrequencyPill{8 次}

\InfoLabel{对应小节} 第三章 描述性研究 / 第二节 现况研究 /
一、概述（三）类型；二、研究设计与实施（三）确定样本量和抽样方法（书中页码 第33---39页）

\InfoLabel{匹配依据} 名解层面，05、10、11、12、14、17、20、21 级考``抽样调查 sampling survey / 系统抽样
systematic sampling / 分层抽样 stratified sampling / 整群抽样 cluster sampling / 多阶段抽样''；简答层面
05、11、12、14、20 级考``现况研究的随机抽样方法及特点''。合并为一个考点。

\ContentBand{知识点原文摘取}

抽样调查（sampling
survey）指通过随机抽样的方法，对特定时点、特定范围内人群的一个代表性样本进行调查，以样本的统计量来估计总体参数所在范围。抽样必须随机化，样本量要足够。

（1）单纯随机抽样（simple random sampling）：按等概率原则从含有 \(N\) 个观察单位的总体中抽取 \(n\)
个组成样本。总体中每个观察单位被抽到的概率相等（均为 \(n/N\)）。

（2）系统抽样（systematic
sampling）：又称机械抽样，按照某种顺序给总体中的个体编号，然后机械地每隔若干单位抽取一个单位。优点是样本分布均匀、代表性较好；缺点是若总体分布有周期性且抽样间隔与之一致，会产生偏性。

（3）分层抽样（stratified
sampling）：先将总体按某种特征分为若干层，再从每层内单纯随机抽样。每一层内个体变异越小越好，层间变异则越大越好。分层抽样比单纯随机抽样精确度更高。

（4）整群抽样（cluster
sampling）：将总体分成若干群组，以群组为抽样单位抽取样本，对抽中群组的全部个体进行调查。特点：易于组织、节省人力物力；抽样误差较大，故通常在单纯随机抽样样本量估算的基础上再增加1/2。

（5）多阶段抽样（multistage
sampling）：将抽样过程分为几个阶段，每个阶段使用的抽样方法往往不同，是大型流行病学调查常用的抽样方法。

\ContentBand{对应往年题}

\begin{pastquestion}

\QuestionLabel{【往年题1｜11级流行病学口试考题回忆.pdf｜简答】}\\
现况研究的随机抽样方法及抽样方法以及三间分布的分析

\end{pastquestion}

\begin{pastquestion}

\QuestionLabel{【往年题2｜17级流病口试回忆.xlsx｜名解（题号11、2）】}\\
系统抽样；抽样调查

\end{pastquestion}

\begin{pastquestion}

\QuestionLabel{【往年题3｜21级流病考题回忆汇总.xlsx｜名解（题号45、25）】}\\
分层抽样 stratified sampling；整群抽样

\end{pastquestion}

\begin{pastquestion}

\QuestionLabel{【往年题4｜14级流病考题回忆汇总.xls｜简答（题号12）】}\\
现况研究中的随机抽样方法，各自的特点

\end{pastquestion}

\PointSeparator

\section{普查}\label{users__fpb__26_spring__pkusph_study__tmp__pdfs__prevention-medicine-past-exams__03-ux6d41ux884cux75c5ux5b66ux5f80ux5e74ux9898ux8003ux70b9ux6574ux7406.md__ux666eux67e5}

\FrequencyPill{6 次}

\InfoLabel{对应小节} 第三章 描述性研究 / 第二节 现况研究 / 一、概述（三）类型 / 1.普查（书中页码 第33页）

\InfoLabel{匹配依据} 04、05、10、12、17、20 级考``普查 census''名词解释；19
级简答考``普查和抽样调查的优缺点''。

\ContentBand{知识点原文摘取}

普查（census）即全面调查，是指将特定时点或时期内、特定范围内的全部人群（总体）作为研究对象的调查。

普查的目的主要包括：①早期发现、早期诊断和早期治疗病人；②了解慢性病的患病情况及急性传染性疾病的疫情分布；③了解人群健康水平；④建立某些生理生化指标的正常值范围。

普查的优点有：①调查对象为全体目标人群，不存在抽样误差；②可以同时调查多种疾病或健康状况；③能发现目标人群中的全部病例，实现``三早''。局限性包括：①不适用于患病率低且诊断技术复杂的疾病；②工作量大，难免漏查；③调查质量不易控制；④耗费人力物力大，费用高。

\ContentBand{对应往年题}

\begin{pastquestion}

\QuestionLabel{【往年题1｜10级考题回忆.pdf｜名解】}\\
census 普查

\end{pastquestion}

\begin{pastquestion}

\QuestionLabel{【往年题2｜12级流病口试考题回忆.pdf｜名解（题号3）】}\\
普查（普查是为了普治）

\end{pastquestion}

\begin{pastquestion}

\QuestionLabel{【往年题3｜19级流行病学期末口试回忆.xlsx｜简答（题号23）】}\\
普查和抽样调查优缺点

\end{pastquestion}

\PointSeparator

\section{生态学研究（含生态比较研究、生态趋势研究）}\label{users__fpb__26_spring__pkusph_study__tmp__pdfs__prevention-medicine-past-exams__03-ux6d41ux884cux75c5ux5b66ux5f80ux5e74ux9898ux8003ux70b9ux6574ux7406.md__ux751fux6001ux5b66ux7814ux7a76ux542bux751fux6001ux6bd4ux8f83ux7814ux7a76ux751fux6001ux8d8bux52bfux7814ux7a76}

\FrequencyPill{6 次}

\InfoLabel{对应小节} 第三章 描述性研究 / 第三节 生态学研究 / 一、概述、二、类型（书中页码 第43---45页）

\InfoLabel{匹配依据} 10、11、12、14、19、20 级考``生态学研究 ecological study / 生态比较研究 ecological
comparison study / 生态趋势研究 ecological trend study''名词解释。

\ContentBand{知识点原文摘取}

生态学研究（ecological
study）是描述性研究的一种类型，它是在群体的水平上研究某种暴露因素与疾病之间的关系，以群体为观察和分析的单位，通过描述不同人群中某因素的暴露情况与疾病的频率，分析该暴露因素与疾病之间的关系。

生态比较研究（ecological comparison
study）：观察不同人群或地区某种疾病的分布，然后根据疾病分布的差异，提出病因假设；更常用来比较不同人群中某因素的平均暴露水平和某疾病频率之间的关系。

生态趋势研究（ecological trend
study）：连续观察人群中某因素平均暴露水平的改变与某种疾病的发病率、死亡率变化的关系，通过比较暴露水平变化前后疾病频率的变化情况，判断某因素与某疾病的联系。

\ContentBand{对应往年题}

\begin{pastquestion}

\QuestionLabel{【往年题1｜12级流病口试考题回忆.pdf｜名解（题号13、39）】}\\
生态学研究

\end{pastquestion}

\begin{pastquestion}

\QuestionLabel{【往年题2｜20级回忆试题库.xlsx｜名解（题号23、47）】}\\
生态比较研究；生态趋势研究（ecological trend study）

\end{pastquestion}

\begin{pastquestion}

\QuestionLabel{【往年题3｜19级流行病学期末口试回忆.xlsx｜名解（题号23）】}\\
ecological study

\end{pastquestion}

\PointSeparator

\section{描述性研究与分析性研究的区别}\label{users__fpb__26_spring__pkusph_study__tmp__pdfs__prevention-medicine-past-exams__03-ux6d41ux884cux75c5ux5b66ux5f80ux5e74ux9898ux8003ux70b9ux6574ux7406.md__ux63cfux8ff0ux6027ux7814ux7a76ux4e0eux5206ux6790ux6027ux7814ux7a76ux7684ux533aux522b}

\FrequencyPill{4 次}

\InfoLabel{对应小节} 第三章 描述性研究 / 第一节 概述（与第四、五章分析性研究对照）（书中页码 第30---32页）

\InfoLabel{匹配依据} 07、10、12、19 级考``描述性研究和分析性研究的区别/不同点''，19
级进一步要求``简述描述性、分析性、实验性流行病学的定义及其异同''。

\ContentBand{知识点原文摘取}

描述性研究主要用来描述人群中疾病或健康状况及暴露因素的分布情况，目的是提出病因假设，为进一步调查研究提供线索；其暴露因素的分配不是随机的，且在研究开始时一般不设立对照组，暴露与结局的时序关系无法确定，因果推断存在局限性。（分析流行病学主要负责检验科研假设，设立对照组比较暴露与疾病的关联，见第四、五章。）

（课本知识点文件中未检索到逐条对照``描述性 vs
分析性''的成表原文，以上为两类研究定位的原文摘取；建议结合课件回答。）

\ContentBand{对应往年题}

\begin{pastquestion}

\QuestionLabel{【往年题1｜07级考题回忆.docx｜简答】}\\
描述性研究和分析性研究的区别

\end{pastquestion}

\begin{pastquestion}

\QuestionLabel{【往年题2｜10级考题回忆.pdf｜简答】}\\
描述性研究和分析性研究的不同点

\end{pastquestion}

\begin{pastquestion}

\QuestionLabel{【往年题3｜19级流行病学期末口试回忆.xlsx｜简答（题号2）】}\\
简述描述性、分析性、实验性流行病学的定义及其异同

\end{pastquestion}

\PointSeparator

\section{现况研究的常见偏倚}\label{users__fpb__26_spring__pkusph_study__tmp__pdfs__prevention-medicine-past-exams__03-ux6d41ux884cux75c5ux5b66ux5f80ux5e74ux9898ux8003ux70b9ux6574ux7406.md__ux73b0ux51b5ux7814ux7a76ux7684ux5e38ux89c1ux504fux501a}

\FrequencyPill{4 次}

\InfoLabel{对应小节} 第三章 描述性研究 / 第二节 现况研究 / 四、常见偏倚及其控制（书中页码 第42页）

\InfoLabel{匹配依据} 10、11、14、20 级考``现况研究的偏倚/现况研究中存在的偏倚/现况研究的常见偏倚''。

\ContentBand{知识点原文摘取}

现况研究中，偏倚产生的原因主要有：①主观选择研究对象，将随意抽样当作随机抽样；②无应答偏倚（无应答者的患病与暴露情况与应答者可能不同；应答率低于70\%难以用调查结果估计总体）；③幸存者偏倚（所调查对象均为幸存者）；④报告偏倚或回忆偏倚；⑤调查偏倚；⑥测量偏倚。前三种属于选择偏倚，后三种属于信息偏倚；此外混杂因素的存在可导致混杂偏倚。

\ContentBand{对应往年题}

\begin{pastquestion}

\QuestionLabel{【往年题1｜11级流行病学口试考题回忆.pdf｜简答】}\\
现况研究的偏倚

\end{pastquestion}

\begin{pastquestion}

\QuestionLabel{【往年题2｜14级流病考题回忆汇总.xls｜简答（题号38）】}\\
现况研究中存在的偏倚

\end{pastquestion}

\begin{pastquestion}

\QuestionLabel{【往年题3｜20级回忆试题库.xlsx｜简答（题号23）】}\\
现况研究的常见偏倚

\end{pastquestion}

\PointSeparator

\section{生态学谬误}\label{users__fpb__26_spring__pkusph_study__tmp__pdfs__prevention-medicine-past-exams__03-ux6d41ux884cux75c5ux5b66ux5f80ux5e74ux9898ux8003ux70b9ux6574ux7406.md__ux751fux6001ux5b66ux8c2cux8bef}

\FrequencyPill{2 次}

\InfoLabel{对应小节} 第三章 描述性研究 / 第三节 生态学研究 / 三、优点和局限性（二）局限性 /
1.生态学谬误（书中页码 第45---46页）

\InfoLabel{匹配依据} 19、20 级考``生态学谬误 ecological fallacy''名词解释，部分追问能否举实际例子。

\ContentBand{知识点原文摘取}

生态学谬误（ecological
fallacy）是生态学研究中最主要的局限性。它是由于生态学研究以各个不同情况的个体``集合''而成的群体（组）为观察和分析的单位，以及存在的混杂因素等原因而造成研究结果与真实情况不符。在群体水平上发现的某因素与某疾病分布上的一致性与该人群中个体的真实情况不符时，就发生了``生态学谬误''。

\ContentBand{对应往年题}

\begin{pastquestion}

\QuestionLabel{【往年题1｜19级流行病学期末口试回忆.xlsx｜名解（题号10、20、27）】}\\
ecological fallacy / 生态学谬误

\end{pastquestion}

\begin{pastquestion}

\QuestionLabel{【往年题2｜20级回忆试题库.xlsx｜名解（题号19、26）】}\\
ecological fallacy（生态学谬误）；生态学研究举例

\end{pastquestion}

\PointSeparator

\section{现况研究与生态学研究的异同点}\label{users__fpb__26_spring__pkusph_study__tmp__pdfs__prevention-medicine-past-exams__03-ux6d41ux884cux75c5ux5b66ux5f80ux5e74ux9898ux8003ux70b9ux6574ux7406.md__ux73b0ux51b5ux7814ux7a76ux4e0eux751fux6001ux5b66ux7814ux7a76ux7684ux5f02ux540cux70b9}

\FrequencyPill{1 次}

\InfoLabel{对应小节} 第三章 描述性研究 / 第二节 现况研究、第三节 生态学研究（思考题3）（书中页码 第30---46页）

\InfoLabel{匹配依据} 21
级综合分析题考``现况研究和生态学研究的异同点''，对应课本思考题``现况研究与生态学研究的异同点及优缺点比较''。

\ContentBand{知识点原文摘取}

现况研究以个体为观察和分析单位，可同时描述疾病分布并探讨多个暴露因素与多种疾病的关系；生态学研究以群体为观察和分析的单位（这是其最基本特征），无法得知个体的暴露与效应间的因果关系，是一种粗线条的描述性研究，最主要局限性是生态学谬误。两者同属描述性研究，均能提供病因线索、产生病因假设，但均不能据此作因果推断。

\ContentBand{对应往年题}

\begin{pastquestion}

\QuestionLabel{【往年题1｜21级流病考题回忆汇总.xlsx｜简答（题号5）】}\\
现况研究和生态学研究的异同点

\end{pastquestion}

\PointSeparator

\chapter{队列研究}\label{users__fpb__26_spring__pkusph_study__tmp__pdfs__prevention-medicine-past-exams__03-ux6d41ux884cux75c5ux5b66ux5f80ux5e74ux9898ux8003ux70b9ux6574ux7406.md__ux7b2cux56dbux7ae0-ux961fux5217ux7814ux7a76}

\section{相对危险度（RR）、归因危险度（AR）、AR\%、PAR、PAR\%
的概念、计算与意义}\label{users__fpb__26_spring__pkusph_study__tmp__pdfs__prevention-medicine-past-exams__03-ux6d41ux884cux75c5ux5b66ux5f80ux5e74ux9898ux8003ux70b9ux6574ux7406.md__ux76f8ux5bf9ux5371ux9669ux5ea6rrux5f52ux56e0ux5371ux9669ux5ea6ararparpar-ux7684ux6982ux5ff5ux8ba1ux7b97ux4e0eux610fux4e49}

\FrequencyPill{11 次}

\InfoLabel{对应小节} 第四章 队列研究 / 第三节 资料的整理与分析 / 四、效应估计（书中页码 第57---59页）

\InfoLabel{匹配依据} 04、05、07、10、11、12、14、17、19、20、21
级\textbf{每一届}都考队列研究效应指标的计算与解释------给暴露组/非暴露组（及全人群）发病率或死亡率，要求计算
RR、AR、AR\%、PAR、PAR\% 并解释含义、流行病学意义及影响因素，常追问``RR 的病因学意义 vs AR/PAR
的公共卫生学意义''。这是本课程考频最高的计算考点。

\ContentBand{知识点原文摘取}

队列研究的主要效应测量指标是相对危险度（relative risk，RR）与归因危险度（attributable
risk，AR），即暴露组与对照组之间的危险度比和危险度差。

\begin{enumerate}
\def\labelenumi{\arabic{enumi}.}
\item
  相对危险度（RR）=
  \(I_e/I_0\)（暴露组率与对照组率之比）。RR表明暴露组发病或死亡的风险是对照组的多少倍。RR=1表示暴露与疾病无关；RR\textgreater1表示暴露是危险因素；RR\textless1表示暴露是保护因素。RR值越远离1，暴露与结局关联的强度越大。
\item
  归因危险度（AR）=
  \(I_e-I_0=I_0(RR-1)\)，又称特异危险度、危险度差（RD）和超额危险度，是暴露组发病率与对照组发病率之差的绝对值，表示危险特异地归因于暴露因素的程度。
\item
  归因危险度百分比（AR\%）=
  \(\frac{I_e-I_0}{I_e}\times100\%=\frac{RR-1}{RR}\times100\%\)，又称病因分值（etiologic
  fraction，EF），指暴露人群中的发病或死亡归因于暴露的部分占全部发病或死亡的百分比。
\item
  人群归因危险度（PAR）= \(I_t-I_0\)；人群归因危险度百分比（PAR\%）=
  \(\frac{I_t-I_0}{I_t}\times100\%=\frac{P_e(RR-1)}{P_e(RR-1)+1}\times100\%\)（\(I_t\) 为全人群率，\(P_e\)
  为人群中暴露者比例）。PAR、PAR\% 既与 RR、AR 有关，又与人群中暴露者比例有关。
\end{enumerate}

RR与AR的区别：RR说明暴露者发生疾病的危险是非暴露者的多少倍（侧重病因学意义）；AR是指暴露人群与非暴露人群比较所增加的疾病发生数量，若消除暴露可减少这个数量的疾病发生（侧重公共卫生学意义）。如吸烟与肺癌、心血管疾病：从RR看吸烟对肺癌作用大（病因联系强），从AR看吸烟对心血管疾病作用大（控烟的社会收益更大）。

\ContentBand{对应往年题}

\begin{pastquestion}

\QuestionLabel{【往年题1｜10级考题回忆.pdf｜计算/大题】}\\
算了几个 rr 和 ar，唐迅问我为啥 rr 具有病因学意义、ar 具有公共卫生学意义，然后要评价这个队列研究

\end{pastquestion}

\begin{pastquestion}

\QuestionLabel{【往年题2｜19级流行病学期末口试回忆.xlsx｜综合分析（题号23、4、29）】}\\
吸烟肺癌死亡关联队列研究的吸烟组、非吸烟组和全人群的肺癌死亡率，求RR、AR、AR\%、PAR、PAR\%并解释流行病学意义；给RR、Pe，算PAR\%，解释指标含义，与RR关系

\end{pastquestion}

\begin{pastquestion}

\QuestionLabel{【往年题3｜21级流病考题回忆汇总.xlsx｜综合分析（题号10、13、39）】}\\
队列研究计算RR、AR、AR\%、PAR、PAR\%（具体概念需要说出）；RR，AR，AR\%

\end{pastquestion}

\begin{pastquestion}

\QuestionLabel{【往年题4｜04级流病考题回忆.docx｜计算】}\\
若干AR、RR然后评价队列研究结果；队列研究的RR和AR

\end{pastquestion}

\begin{pastquestion}

\QuestionLabel{【往年题5｜11级流行病学口试考题回忆.pdf｜计算】}\\
队列效应指标的计算和解释（被王涛追问 par\% 的实际应用，被孙凤追问 RD 和 AR 的关系，以及 RD 英文）

\end{pastquestion}

\PointSeparator

\section{暴露人年（人时）及其计算}\label{users__fpb__26_spring__pkusph_study__tmp__pdfs__prevention-medicine-past-exams__03-ux6d41ux884cux75c5ux5b66ux5f80ux5e74ux9898ux8003ux70b9ux6574ux7406.md__ux66b4ux9732ux4ebaux5e74ux4ebaux65f6ux53caux5176ux8ba1ux7b97}

\FrequencyPill{7 次}

\InfoLabel{对应小节} 第四章 队列研究 / 第三节 资料的整理与分析 / 二、人时的计算（书中页码 第55---56页）

\InfoLabel{匹配依据} 04、07、10、11、12、14、19 级考``暴露人年（exposure person
year）''名词解释及``人月/人时的计算''计算题，部分追问暴露人年的优缺点、适用研究（动态队列）。

\ContentBand{知识点原文摘取}

队列研究由于时间跨度一般较长，观察人口不稳定。较合理的办法是考虑时间因素，引入人时（person
time）的概念来描述观察对象的暴露经历，人时即观察人数与观察时间的乘积。常用的人时计算方法有三种：

\begin{enumerate}
\def\labelenumi{\arabic{enumi}.}
\tightlist
\item
  精确法计算人时：将每个人的精确的观察时间相加。结果精确，但资料处理麻烦。
\item
  近似法计算人时：用每年观察的平均人数（取相邻两年的年初人口的平均数或年中人口数）作为该年的观察人年数，再将各年相加。计算简单，精确性略差。
\item
  寿命表法计算人时：常将观察当年内进入队列的个体作1/2人年计算；失访或出现终点结局的个体也作1/2人年计算。
\end{enumerate}

对于动态队列，多以观察人时为分母计算发病率，称为发病密度（incidence density）。

\ContentBand{对应往年题}

\begin{pastquestion}

\QuestionLabel{【往年题1｜12级流病口试考题回忆.pdf｜计算（题号7、11）】}\\
人月的计算；暴露人年

\end{pastquestion}

\begin{pastquestion}

\QuestionLabel{【往年题2｜10级考题回忆.pdf｜名解/计算】}\\
exposed person year（暴露人年）；人月计算

\end{pastquestion}

\begin{pastquestion}

\QuestionLabel{【往年题3｜14级流病考题回忆汇总.xls｜名解/计算（题号43、50）】}\\
暴露人年（追问局限性、什么情况下适用------动态队列；约登指数等）；暴露人月计算

\end{pastquestion}

\PointSeparator

\section{队列研究对照人群的选择（内对照、外对照、总人口对照）}\label{users__fpb__26_spring__pkusph_study__tmp__pdfs__prevention-medicine-past-exams__03-ux6d41ux884cux75c5ux5b66ux5f80ux5e74ux9898ux8003ux70b9ux6574ux7406.md__ux961fux5217ux7814ux7a76ux5bf9ux7167ux4ebaux7fa4ux7684ux9009ux62e9ux5185ux5bf9ux7167ux5916ux5bf9ux7167ux603bux4ebaux53e3ux5bf9ux7167}

\FrequencyPill{5 次}

\InfoLabel{对应小节} 第四章 队列研究 / 第二节 研究设计与实施 /
三、确定研究现场与研究人群（二）研究人群（书中页码 第51---52页）

\InfoLabel{匹配依据} 05、07、10、11 级考``队列研究的对照（组）类型及优缺点/暴露人群选择原则''简答；14、16、20
级考``内对照（internal control）/外对照（external control）''名词解释，部分追问内对照适用场景及与外对照区别。

\ContentBand{知识点原文摘取}

（1）内对照（internal
control）：依据与暴露无关的因素确定一个研究人群，将其中暴露于所研究因素的对象作为暴露组，未暴露者为对照组。好处是选取对照比较容易，可比性较好。当暴露变量是定量或等级变量时，可按暴露剂量分组，以最低暴露水平人群为对照组。

（2）外对照（external
control）：当选择职业人群或特殊暴露人群作为暴露组时，常需在该人群之外去寻找对照组。优点是随访时可免受暴露组的``沾染''，缺点是需另外组织一项人群工作。

（3）总人口对照（total population
control）：利用整个地区现成的发病或死亡统计资料，以全人群为对照。优点是对比资料容易得到，缺点是资料比较粗糙、可比性较差，一般用于总人群中暴露者比例很小的情形（通过计算标化比指标）。

\ContentBand{对应往年题}

\begin{pastquestion}

\QuestionLabel{【往年题1｜05流病考题回忆.doc｜简答】}\\
队列研究的对照类型及其优缺点

\end{pastquestion}

\begin{pastquestion}

\QuestionLabel{【往年题2｜11级流行病学口试考题回忆.pdf｜简答】}\\
队列研究暴露人群选择原则

\end{pastquestion}

\begin{pastquestion}

\QuestionLabel{【往年题3｜20级回忆试题库.xlsx｜名解（题号8、40）】}\\
internal control（内对照）（追问适用场景及与外对照的区别）

\end{pastquestion}

\PointSeparator

\section{队列研究的优点与局限性}\label{users__fpb__26_spring__pkusph_study__tmp__pdfs__prevention-medicine-past-exams__03-ux6d41ux884cux75c5ux5b66ux5f80ux5e74ux9898ux8003ux70b9ux6574ux7406.md__ux961fux5217ux7814ux7a76ux7684ux4f18ux70b9ux4e0eux5c40ux9650ux6027}

\FrequencyPill{4 次}

\InfoLabel{对应小节} 第四章 队列研究 / 第五节 优缺点及其他实践类型 / 一、优点、二、局限性（书中页码
第60---61页）

\InfoLabel{匹配依据} 11、19、21 级考``队列研究的优劣/优缺点''简答；13 级（14 级文件）等亦涉及。

\ContentBand{知识点原文摘取}

优点：①暴露资料在结局发生之前收集，资料完整可靠，回忆偏倚相对较小；②可直接获得发病率或死亡率，计算RR和AR等关联强度指标，直接反映暴露的病因作用；③暴露在前、结局在后，因果时间顺序合理，检验病因假设能力较强；④有助于了解疾病的自然史；⑤能分析一种暴露因素与多种结局之间的关系。

局限性：①不适于发病率很低的疾病的病因研究；②随访时间较长，容易产生失访偏倚；③前瞻性队列研究耗费人力、物力、财力和时间较多，不易实施。

\ContentBand{对应往年题}

\begin{pastquestion}

\QuestionLabel{【往年题1｜19级流行病学期末口试回忆.xlsx｜简答（题号13）】}\\
队列研究的优劣

\end{pastquestion}

\begin{pastquestion}

\QuestionLabel{【往年题2｜21级流病考题回忆汇总.xlsx｜简答（题号7）】}\\
队列研究优缺点

\end{pastquestion}

\PointSeparator

\section{历史性（回顾性）队列研究}\label{users__fpb__26_spring__pkusph_study__tmp__pdfs__prevention-medicine-past-exams__03-ux6d41ux884cux75c5ux5b66ux5f80ux5e74ux9898ux8003ux70b9ux6574ux7406.md__ux5386ux53f2ux6027ux56deux987eux6027ux961fux5217ux7814ux7a76}

\FrequencyPill{3 次}

\InfoLabel{对应小节} 第四章 队列研究 / 第一节 概述 / 四、研究类型（二）历史性队列研究（书中页码 第49---50页）

\InfoLabel{匹配依据} 05、19、20 级考``历史性队列研究 / historical cohort study /
回顾性队列''名词解释，部分追问其特点（是否前瞻性）。

\ContentBand{知识点原文摘取}

历史性队列研究也称回顾性队列研究（retrospective cohort
study），是依据研究开始时研究者已掌握的有关研究对象过去某个时点暴露状况的历史材料进行分组，研究结局在研究开始时已经发生，不需要前瞻性观察。虽然收集暴露与结局资料的方法是回顾性的，但究其本质仍然是先有暴露后有结局，研究方向依然是从因到果。该类研究具有省时、省力、出结果快的特点，仅在具备详细、准确的历史资料的条件下适用。

\ContentBand{对应往年题}

\begin{pastquestion}

\QuestionLabel{【往年题1｜05流病考题回忆.doc｜名解】}\\
历史性队列研究（追问：历史性队列研究是不是前瞻性的？）

\end{pastquestion}

\begin{pastquestion}

\QuestionLabel{【往年题2｜19级流行病学期末口试回忆.xlsx｜名解（题号10）】}\\
historical cohort study / 回顾性队列

\end{pastquestion}

\begin{pastquestion}

\QuestionLabel{【往年题3｜20级回忆试题库.xlsx｜综合分析（题号10）】}\\
回顾性队列研究；相关指标计算

\end{pastquestion}

\PointSeparator

\section{标化死亡比（SMR）}\label{users__fpb__26_spring__pkusph_study__tmp__pdfs__prevention-medicine-past-exams__03-ux6d41ux884cux75c5ux5b66ux5f80ux5e74ux9898ux8003ux70b9ux6574ux7406.md__ux6807ux5316ux6b7bux4ea1ux6bd4smr}

\FrequencyPill{2 次}

\InfoLabel{对应小节} 第四章 队列研究 / 第三节 资料的整理与分析 / 三、率的计算（一）常用指标 /
3.标化比（书中页码 第56页）

\InfoLabel{匹配依据} 07 级（颜力组）名解考``SMR''，并比较实验流行病学与队列研究异同。SMR
在职业流行病学队列研究中常用。

\ContentBand{知识点原文摘取}

当采用总人口作对照或研究对象数量较少、结局事件发生率较低时，常以全人口发病（死亡）率作为标准，计算该观察人群的预期发病（死亡）人数，以观察人群实际发病（死亡）人数除以预期发病（死亡）人数来计算标化发病/死亡比（standardized
morbidity/mortality ratio，SMR）。

\[SMR=\frac{研究人群观察死亡数（O）}{以标准人口死亡率计算的预期死亡数（E）}\]

标化比实际上不是率，而是以全人口的发病（死亡）率作为对照计算出来的比，其流行病学意义与效应指标RR类似。

\ContentBand{对应往年题}

\begin{pastquestion}

\QuestionLabel{【往年题1｜07级考题回忆.docx｜名解（颜力组，记录于05级回忆汇总）】}\\
SMR

\end{pastquestion}

\PointSeparator

\chapter{病例对照研究}\label{users__fpb__26_spring__pkusph_study__tmp__pdfs__prevention-medicine-past-exams__03-ux6d41ux884cux75c5ux5b66ux5f80ux5e74ux9898ux8003ux70b9ux6574ux7406.md__ux7b2cux4e94ux7ae0-ux75c5ux4f8bux5bf9ux7167ux7814ux7a76}

\section{比值比（OR）的概念、计算与意义（含非匹配四格表 OR、1∶1 配对
OR）}\label{users__fpb__26_spring__pkusph_study__tmp__pdfs__prevention-medicine-past-exams__03-ux6d41ux884cux75c5ux5b66ux5f80ux5e74ux9898ux8003ux70b9ux6574ux7406.md__ux6bd4ux503cux6bd4orux7684ux6982ux5ff5ux8ba1ux7b97ux4e0eux610fux4e49ux542bux975eux5339ux914dux56dbux683cux8868-or11-ux914dux5bf9-or}

\FrequencyPill{11 次}

\InfoLabel{对应小节} 第五章 病例对照研究 / 第三节 资料的整理与分析 / 二、资料的分析（书中页码 第70---73页）

\InfoLabel{匹配依据} 04、05、07、10、11、12、14、17、19、20、21 级\textbf{每一届}都考``OR（odds ratio /
比值比）''名词解释及计算------非匹配四格表（OR=ad/bc）、1∶1
配对（OR=c/b）、雌激素与子宫内膜癌的配对病例对照、口服避孕药与心肌梗死等经典例题，常追问结果解释、能否据此判断因果关系、95\%CI
含义。本课程最高频计算考点之一。

\ContentBand{知识点原文摘取}

比值比又称比数比、优势比，为病例组与对照组两组暴露比值之比。非匹配资料的四格表 OR：

\[OR=\frac{ad}{bc}\]

OR值的含义与RR值相同，均指暴露者疾病的危险性是非暴露者的多少倍。OR\textgreater1说明暴露可增加疾病的危险性（危险因素）；OR\textless1说明暴露可降低疾病的危险性（保护因素）；OR=1表明暴露与疾病无统计学联系。当所研究疾病的发病率很低（如小于5\%）时，OR≈RR。

非匹配四格表暴露与疾病关联性检验用 \(\chi^2=\dfrac{(ad-bc)^2N}{(a+b)(c+d)(a+c)(b+d)}\)。

1∶1 配对资料（表内 \(a,b,c,d\) 为对子数）用 McNemar \(\chi^2\) 检验：\(\chi^2=\dfrac{(b-c)^2}{b+c}\)，当
\((b+c)<40\) 时用校正公式 \(\chi^2=\dfrac{(|b-c|-1)^2}{b+c}\)；配对 OR：

\[OR=\frac{c}{b}\ (b\ne0)\]

OR 95\%CI 常用 Miettinen 法：\(OR^{(1\pm1.96/\sqrt{\chi^2})}\)。若 95\%CI 不包括
1，可认为研究因素与疾病有关联。

（外源性雌激素与子宫内膜癌配对资料：校正
\(\chi^2=19.53\)，\(OR=29/3=9.67\)，提示外源性雌激素是子宫内膜癌的危险因素。）

\ContentBand{对应往年题}

\begin{pastquestion}

\QuestionLabel{【往年题1｜05流病考题回忆.doc｜计算】}\\
1：1配对的卡方、OR计算，根据结果进行解释，能不能据此确定暴露和疾病的关系

\end{pastquestion}

\begin{pastquestion}

\QuestionLabel{【往年题2｜12级流病口试考题回忆.pdf｜计算（题号34、40）】}\\
雌激素和子宫内膜癌 计算卡方和OR，这个病例对照可能存在的偏倚；口服雌激素与子宫内膜癌1∶1配对，计算OR及卡方

\end{pastquestion}

\begin{pastquestion}

\QuestionLabel{【往年题3｜20级回忆试题库.xlsx｜综合分析（题号13）】}\\
1比1配对病例对照研究：第一问计算OR值和卡方值；第二问根据结果是否能推断两变量间的因果关系

\end{pastquestion}

\begin{pastquestion}

\QuestionLabel{【往年题4｜21级流病考题回忆汇总.xlsx｜综合分析（题号29）】}\\
1：1配对的病例对照研究，有无高血压病史和心梗发病率的关系，计算OR值并结合题目给出的卡方值解释

\end{pastquestion}

\PointSeparator

\section{匹配（配比、频数匹配、个体匹配、配对、过度匹配）及匹配的目的}\label{users__fpb__26_spring__pkusph_study__tmp__pdfs__prevention-medicine-past-exams__03-ux6d41ux884cux75c5ux5b66ux5f80ux5e74ux9898ux8003ux70b9ux6574ux7406.md__ux5339ux914dux914dux6bd4ux9891ux6570ux5339ux914dux4e2aux4f53ux5339ux914dux914dux5bf9ux8fc7ux5ea6ux5339ux914dux53caux5339ux914dux7684ux76eeux7684}

\FrequencyPill{10 次}

\InfoLabel{对应小节} 第五章 病例对照研究 / 第二节 研究设计与实施 / 三、确定研究对象（二）对照的选择 /
3.选择对照的方法（书中页码 第66页）

\InfoLabel{匹配依据} 04、05、07、10、11、12、14、17、19、20、21 级广泛考``匹配 matching / 配对 pair matching /
频数匹配 frequency matching / 个体匹配 individual matching / 过度匹配
overmatching''名词解释，以及``病例对照研究为什么要匹配、匹配不当会怎样、匹配的目的、如何控制''问答。

\ContentBand{知识点原文摘取}

匹配或称配比，是要求对照在某些特征或因素上与病例保持一致，保证对照与病例具有可比性，以便对两组进行比较时排除匹配因素的干扰。匹配的目的主要是提高研究效率，其次是控制混杂因素的干扰。

匹配变量必须是确定或有充分理由怀疑为混杂因素，否则不应匹配：①疾病因果链上的中间变量不应匹配；②只与可疑病因有关而与疾病无关的因素不应匹配。把不起混杂作用的因素作为匹配变量，试图使对照组与病例组在多方面都一致，结果导致所研究的因素也趋于一致，反而降低了研究效率，这种情况称为匹配过度（overmatching）。

根据匹配方式不同，分为频数匹配（frequency
matching，对照组某特征所占比例与病例组一致或相近）和个体匹配（individual
matching，以个体为单位匹配）。1个病例匹配1个对照叫配对（pair matching），也可1∶r
匹配；病例与对照之比为1∶1时统计学效率最高。

\ContentBand{对应往年题}

\begin{pastquestion}

\QuestionLabel{【往年题1｜04级流病考题回忆.docx｜问答】}\\
为啥配比，会出现哪些问题，怎样纠正？；病例对照研究为什么要匹配，匹配不当会怎样，如何控制

\end{pastquestion}

\begin{pastquestion}

\QuestionLabel{【往年题2｜10级考题回忆.pdf｜名解】}\\
个体匹配；过度匹配（over-matching）；频数匹配（frequency matching）

\end{pastquestion}

\begin{pastquestion}

\QuestionLabel{【往年题3｜12级流病口试考题回忆.pdf｜名解（题号29）】}\\
over matching 过度匹配

\end{pastquestion}

\PointSeparator

\section{以医院为基础与以社区为基础的病例对照研究的优缺点}\label{users__fpb__26_spring__pkusph_study__tmp__pdfs__prevention-medicine-past-exams__03-ux6d41ux884cux75c5ux5b66ux5f80ux5e74ux9898ux8003ux70b9ux6574ux7406.md__ux4ee5ux533bux9662ux4e3aux57faux7840ux4e0eux4ee5ux793eux533aux4e3aux57faux7840ux7684ux75c5ux4f8bux5bf9ux7167ux7814ux7a76ux7684ux4f18ux7f3aux70b9}

\FrequencyPill{8 次}

\InfoLabel{对应小节} 第五章 病例对照研究 / 第二节 研究设计与实施 /
三、确定研究对象（一）病例的选择---3.病例的来源、（二）对照的选择---2.对照的来源（书中页码 第65页）

\InfoLabel{匹配依据} 04、05、10、11、12、17、19、21
级反复考``以医院为基础和以社区为基础的病例对照研究的优缺点''，常作为大题第一问，与 OR 计算配套出现。

\ContentBand{知识点原文摘取}

病例来源：一类从医院选择病例，可节省费用、合作性好、资料容易得到且信息较完整准确，但不同医院接收的病人特征不同，仅从一所医院选择代表性较差，应尽量选自不同水平、不同种类的医院。另一类从社区人群中选择病例，代表性好，结果推及人群的可信程度较高，但调查工作比较困难，花费人力物力较多。

对照来源（医院对照）：易于选取、合作性好、可利用档案资料，但暴露分布常不同于病例的源人群（医院对照的暴露水平可能高于源人群）；为避免选择偏倚，①因已知与所研究暴露有关的病种入院的病人不能作为对照，②对照应由尽可能多的病种的病人组成。社区/团体人群对照不易出现上述选择偏倚问题，但实施难度大、费用高、对照不易配合。

\ContentBand{对应往年题}

\begin{pastquestion}

\QuestionLabel{【往年题1｜04级流病考题回忆.docx｜简答】}\\
以医院为基础和以社区为基础的病例对照研究的优缺点

\end{pastquestion}

\begin{pastquestion}

\QuestionLabel{【往年题2｜05流病考题回忆.doc｜计算（蔡赐河）】}\\
以医院为基础的病例对照研究的优缺点（追问以社区为基础的研究的优缺点），计算OR并解释结果

\end{pastquestion}

\begin{pastquestion}

\QuestionLabel{【往年题3｜21级流病考题回忆汇总.xlsx｜综合分析（题号29、30、4）】}\\
以医院为基础的病例对照研究的优缺点是什么，计算并解释OR值

\end{pastquestion}

\PointSeparator

\section{病例对照研究与队列研究的原理及异同}\label{users__fpb__26_spring__pkusph_study__tmp__pdfs__prevention-medicine-past-exams__03-ux6d41ux884cux75c5ux5b66ux5f80ux5e74ux9898ux8003ux70b9ux6574ux7406.md__ux75c5ux4f8bux5bf9ux7167ux7814ux7a76ux4e0eux961fux5217ux7814ux7a76ux7684ux539fux7406ux53caux5f02ux540c}

\FrequencyPill{7 次}

\InfoLabel{对应小节} 第五章 病例对照研究 / 第一节 概述、第五节 优点与局限性（表5-11）；第四章
队列研究（书中页码 第62---63、81页）

\InfoLabel{匹配依据} 04、05、10、11、12、14、19 级反复考``病例对照研究与队列研究的原理并比较两者异同/区别''。

\ContentBand{知识点原文摘取}

病例对照研究是以当前已确诊患某病的一组病人作为病例组，以不患该病但有可比性的一组个体作为对照组，搜集既往暴露史并比较两组暴露比例差异，是一种由``果''及``因''的分析性研究，因果论证强度不及队列研究，只能用
OR 估计 RR。

队列研究是将研究人群按是否暴露分为不同亚组，随访观察各组结局发生率的差异，是由``因''及``果''的分析性研究，可直接计算发病率、RR
和 AR，因果时间顺序明确，检验病因假说能力较强。

两者均为分析性观察研究、均设对照组；区别在于：研究方向（由果溯因 vs 由因及果）、时间顺序、效应指标（OR vs
RR/AR）、适用疾病（病例对照适于罕见病、省时省力 vs
队列适于罕见暴露、可研究多结局）、论证因果能力（队列更强）。

\ContentBand{对应往年题}

\begin{pastquestion}

\QuestionLabel{【往年题1｜12级流病口试考题回忆.pdf｜简答（题号6、23）】}\\
队列研究和病例对照研究的原理并比较异同点；病例对照研究和队列研究的基本原理，并比较两者异同

\end{pastquestion}

\begin{pastquestion}

\QuestionLabel{【往年题2｜14级流病考题回忆汇总.xls｜简答（题号1）】}\\
病例对照和队列的原理及异同点比较

\end{pastquestion}

\begin{pastquestion}

\QuestionLabel{【往年题3｜19级流行病学期末口试回忆.xlsx｜简答（题号14）】}\\
病例对照和队列研究异同

\end{pastquestion}

\PointSeparator

\section{病例对照研究样本量的影响因素（及与队列研究样本量异同）}\label{users__fpb__26_spring__pkusph_study__tmp__pdfs__prevention-medicine-past-exams__03-ux6d41ux884cux75c5ux5b66ux5f80ux5e74ux9898ux8003ux70b9ux6574ux7406.md__ux75c5ux4f8bux5bf9ux7167ux7814ux7a76ux6837ux672cux91cfux7684ux5f71ux54cdux56e0ux7d20ux53caux4e0eux961fux5217ux7814ux7a76ux6837ux672cux91cfux5f02ux540c}

\FrequencyPill{5 次}

\InfoLabel{对应小节} 第五章 病例对照研究 / 第二节 研究设计与实施 /
四、确定样本量（一）影响样本量的因素（书中页码 第66页）

\InfoLabel{匹配依据} 10、11、12、19、20
级考``影响病例对照研究样本量的因素''及``病例对照与队列研究样本量影响因素的异同''。

\ContentBand{知识点原文摘取}

病例对照研究的样本含量与下列四个因素有关：①研究因素在对照组或人群中的暴露率（\(P_0\)）；②研究因素与疾病关联强度，即比值比（OR）；③希望达到的检验假设的显著性水平，即第Ⅰ类错误概率（\(\alpha\)），一般取
0.05；④希望达到的检验效能或把握度（\(1-\beta\)），\(\beta\) 为第Ⅱ类错误概率，一般取
0.1。如果采取匹配设计，估计样本量时还要考虑病例和对照的比例。

\ContentBand{对应往年题}

\begin{pastquestion}

\QuestionLabel{【往年题1｜19级流行病学期末口试回忆.xlsx｜简答（题号19）】}\\
病例对照样本量计算的影响因素（追问 P0、P1 的含义）

\end{pastquestion}

\begin{pastquestion}

\QuestionLabel{【往年题2｜20级回忆试题库.xlsx｜简答（题号5）】}\\
病例对照和队列研究的样本量影响因素

\end{pastquestion}

\begin{pastquestion}

\QuestionLabel{【往年题3｜10级考题回忆.pdf｜简答（题号25）】}\\
病例对照与队列样本量影响因素的异同

\end{pastquestion}

\PointSeparator

\section{分级（剂量-反应）资料的 OR
计算与趋势卡方（吸烟与肺癌）}\label{users__fpb__26_spring__pkusph_study__tmp__pdfs__prevention-medicine-past-exams__03-ux6d41ux884cux75c5ux5b66ux5f80ux5e74ux9898ux8003ux70b9ux6574ux7406.md__ux5206ux7ea7ux5242ux91cf-ux53cdux5e94ux8d44ux6599ux7684-or-ux8ba1ux7b97ux4e0eux8d8bux52bfux5361ux65b9ux5438ux70dfux4e0eux80baux764c}

\FrequencyPill{4 次}

\InfoLabel{对应小节} 第五章 病例对照研究 / 第三节 资料的整理与分析 /（四）剂量-反应关系的分析（书中页码
第76---77页）

\InfoLabel{匹配依据} 04、05、10、12 级考``吸烟与肺癌分级病例对照资料，计算各暴露水平
OR、趋势卡方，分析剂量-反应关系''，追问为什么要分级（剂量-反应关系作为因果推断证据）。

\ContentBand{知识点原文摘取}

如果能够获得某些暴露因素不同暴露水平的资料（分级资料），可将不同暴露等级分别与无暴露或最低水平暴露作比较，分析剂量-反应关系，增加因果关系推断的依据。以不暴露或最低水平暴露组为参照，分别计算各暴露水平的
OR。以 Doll 和 Hill 男性吸烟与肺癌资料为例，每日吸烟 1～4、5～14、15～ 三组 OR 分别为
8.10、11.52、17.93，随吸烟量增加而递增，呈明显剂量-反应关系，再经 \(\chi^2\)
趋势检验（\(\chi^2=40.01\)，\(P<0.001\)）判明该剂量-反应关系有统计学意义。

\ContentBand{对应往年题}

\begin{pastquestion}

\QuestionLabel{【往年题1｜05流病考题回忆.doc｜计算（老绵羊/李宇杰）】}\\
吸烟和肺癌关系分层计算OR还有趋势卡方；分成6层，计算各层OR并对危险度做出分析

\end{pastquestion}

\begin{pastquestion}

\QuestionLabel{【往年题2｜04级流病考题回忆.docx｜计算】}\\
卡方趋势检验（从不吸烟 轻度中度重度的病例对照）计算各种指标和意义

\end{pastquestion}

\begin{pastquestion}

\QuestionLabel{【往年题3｜10级考题回忆.pdf｜计算/大题】}\\
吸烟与肺癌分层 or 的计算，注意是剂量反应不是剂量效应关系

\end{pastquestion}

\PointSeparator

\section{巢式病例对照研究}\label{users__fpb__26_spring__pkusph_study__tmp__pdfs__prevention-medicine-past-exams__03-ux6d41ux884cux75c5ux5b66ux5f80ux5e74ux9898ux8003ux70b9ux6574ux7406.md__ux5de2ux5f0fux75c5ux4f8bux5bf9ux7167ux7814ux7a76}

\FrequencyPill{2 次}

\InfoLabel{对应小节} 第五章 病例对照研究 / 第一节 概述 / 三、研究类型（衍生类型）（书中页码 第62---63页）

\InfoLabel{匹配依据} 04、19 级考``巢式病例对照研究（nested case-control study）''名词解释及优缺点。

\ContentBand{知识点原文摘取}

（课本第五章正文重点介绍非匹配与匹配两种基本类型，并提及随研究方法发展产生了多种改进的衍生类型；巢式病例对照研究即为在队列基础上的衍生设计。课本知识点文件中未检索到``巢式病例对照研究''的独立定义原文，暂不自行补充，建议结合课件回答其``在队列内、由因及果、省时省力、回忆偏倚小''等特点。）

\ContentBand{对应往年题}

\begin{pastquestion}

\QuestionLabel{【往年题1｜04级流病考题回忆.docx｜问答】}\\
nested case control study 优缺点

\end{pastquestion}

\begin{pastquestion}

\QuestionLabel{【往年题2｜19级流行病学期末口试回忆.xlsx｜名解（题号11）】}\\
巢式病例对照

\end{pastquestion}

\PointSeparator

\chapter{实验研究（实验流行病学）}\label{users__fpb__26_spring__pkusph_study__tmp__pdfs__prevention-medicine-past-exams__03-ux6d41ux884cux75c5ux5b66ux5f80ux5e74ux9898ux8003ux70b9ux6574ux7406.md__ux7b2cux516dux7ae0-ux5b9eux9a8cux7814ux7a76ux5b9eux9a8cux6d41ux884cux75c5ux5b66}

\section{疫苗保护率（PR）与效果指数（IE）的概念、计算与评价}\label{users__fpb__26_spring__pkusph_study__tmp__pdfs__prevention-medicine-past-exams__03-ux6d41ux884cux75c5ux5b66ux5f80ux5e74ux9898ux8003ux70b9ux6574ux7406.md__ux75abux82d7ux4fddux62a4ux7387prux4e0eux6548ux679cux6307ux6570ieux7684ux6982ux5ff5ux8ba1ux7b97ux4e0eux8bc4ux4ef7}

\FrequencyPill{9 次}

\InfoLabel{对应小节} 第六章 实验流行病学 / 第三节 资料的整理与分析 / 三、评价指标 /
2.评价预防措施效果的主要指标（书中页码 第96---97页）

\InfoLabel{匹配依据} 10、11、12、14、17、19、20、21 级反复考``疫苗保护率（protective rate）、效果指数（index
of
effectiveness）''的名词解释、公式与计算，并常作为疫苗现场试验大题的计算部分（如老年组/非老年组保护率比较）。20
级名解亦考``protective rate / index of effectiveness''。

\ContentBand{知识点原文摘取}

（1）保护率（protective rate，PR）=
\(\dfrac{对照组发病（或死亡）率-试验组发病（或死亡）率}{对照组发病（或死亡）率}\times100\%\)

（2）效果指数（index of effectiveness，IE）= \(\dfrac{对照组发病（或死亡）率}{试验组发病（或死亡）率}\)

预防措施效果的考核还可用抗体阳转率、抗体滴度几何平均数、病情轻重变化等指标评价。（RCT
危险度评价指标中：相对危险度 RR=EER/CER；效果指数 IE=CER/EER；相对危险度减少值/保护率
RRR=\textbar CER−EER\textbar/CER。）

\ContentBand{对应往年题}

\begin{pastquestion}

\QuestionLabel{【往年题1｜17级流病口试回忆.xlsx｜简答（题号7）】}\\
归因危险度百分比、疫苗保护率概念和公式

\end{pastquestion}

\begin{pastquestion}

\QuestionLabel{【往年题2｜20级回忆试题库.xlsx｜综合分析（题号2、9、30、41）】}\\
疫苗接种的有效率、效果指数；EV71疫苗的现场试验：计算和评价；疫苗保护率和治疗指数计算

\end{pastquestion}

\begin{pastquestion}

\QuestionLabel{【往年题3｜21级流病考题回忆汇总.xlsx｜综合分析（题号24、44）】}\\
疫苗的干预实验，分老年组和非老年组，计算两组的保护率，解释差别和可能的原因

\end{pastquestion}

\PointSeparator

\section{高血压两种药物（针剂/片剂）实验设计的评价与重新设计}\label{users__fpb__26_spring__pkusph_study__tmp__pdfs__prevention-medicine-past-exams__03-ux6d41ux884cux75c5ux5b66ux5f80ux5e74ux9898ux8003ux70b9ux6574ux7406.md__ux9ad8ux8840ux538bux4e24ux79cdux836fux7269ux9488ux5242ux7247ux5242ux5b9eux9a8cux8bbeux8ba1ux7684ux8bc4ux4ef7ux4e0eux91cdux65b0ux8bbeux8ba1}

\FrequencyPill{8 次}

\InfoLabel{对应小节} 第六章 实验流行病学 / 第二节
研究设计与实施（随机、对照、盲法、依从性）；第三节（不依从、ITT 分析）（书中页码 第86---95页）

\InfoLabel{匹配依据} 04、05、07、10、11、12、17、19、20
级反复出现这道经典大题：``研究两种高血压药（一种片剂、一种针剂，针剂组病人自行注射），6 周后针剂组约 50\%
自行停药、片剂组约 15\%
停药，要求①评价该实验设计②自己重新设计。''答题要点为实验流行病学的随机、对照、盲法、前瞻原则，以及减少不依从、控制混杂等。

\ContentBand{知识点原文摘取}

实验流行病学研究具有以下基本特征：①属于前瞻性研究；②必须施加一种或多种干预处理；③研究对象来自同一总体、随机分配；④必须有平行的试验组和对照组，开始时两组基线特征近似或可比。

随机化：将研究对象随机分配到试验组和对照组，平衡已知和未知的混杂因素，提高可比性。盲法分单盲、双盲、三盲。设立对照可排除不能预知的结局、向均数回归、霍桑效应、安慰剂效应、潜在未知因素的影响。

不依从（noncompliance）：试验组成员不遵守干预规程相当于退出试验组；对照组成员私下接受干预相当于加入试验组。为减少不依从应进行宣传教育、注意设计合理性、试验期限不宜过长、简化干预措施。对不依从者不能剔除，应采用意向治疗分析（ITT）。

\ContentBand{对应往年题}

\begin{pastquestion}

\QuestionLabel{【往年题1｜05流病考题回忆.doc｜开放题】}\\
高血压药物研究，片剂、针剂两组，针剂50\%没坚持，片剂15\%没坚持。对上述研究评价，并提出自己的方法

\end{pastquestion}

\begin{pastquestion}

\QuestionLabel{【往年题2｜11级流行病学口试考题回忆.pdf｜计算/设计】}\\
高血压针剂片剂，评价原设计，自己重新设计（答案要按照随机、对照、干预、前瞻来说，而且一定要说这几个字）

\end{pastquestion}

\begin{pastquestion}

\QuestionLabel{【往年题3｜17级流病口试回忆.xlsx｜综合分析（题号11）】}\\
高血压片剂和针剂，重新实验设计（给出了原先设计当中两组依从性，原实验问题在于依从性太差）；追问意向治疗分析、依从性分析和符合方案分析的区别

\end{pastquestion}

\PointSeparator

\section{疫苗现场试验的设计与评价（现场与研究对象的选择原则）}\label{users__fpb__26_spring__pkusph_study__tmp__pdfs__prevention-medicine-past-exams__03-ux6d41ux884cux75c5ux5b66ux5f80ux5e74ux9898ux8003ux70b9ux6574ux7406.md__ux75abux82d7ux73b0ux573aux8bd5ux9a8cux7684ux8bbeux8ba1ux4e0eux8bc4ux4ef7ux73b0ux573aux4e0eux7814ux7a76ux5bf9ux8c61ux7684ux9009ux62e9ux539fux5219}

\FrequencyPill{8 次}

\InfoLabel{对应小节} 第六章 实验流行病学 / 第二节 研究设计与实施 /
二、确定试验现场、三、选择研究对象（书中页码 第87页）

\InfoLabel{匹配依据} 10、11、12、17、19、20、21
级反复考``评价某疫苗现场试验的设计/自己设计一个疫苗效果评价实验''，要求回答现场试验的种类、现场和人群的选择原则、是否合适、指标计算与评价、得出结论（如
EV71 疫苗、四价流感疫苗）。

\ContentBand{知识点原文摘取}

确定试验现场通常考虑：①人口相对稳定、流动性小且数量足够；②试验疾病在该地区有较高而稳定的发病率；③评价疫苗免疫学效果时应选择近期内未发生该病流行的地区；④有较好的医疗卫生条件、健全的登记报告制度；⑤领导重视、群众愿意接受、协作配合好。

选择研究对象的主要原则：①选择干预措施可能起效的人群（如疫苗评价应选易感人群）；②选择预期发病率较高的人群（疾病高发区/高危人群）；③选择干预对其无害的人群；④选择能将试验坚持到底的人群；⑤选择依从性好的人群。

现场试验以尚未患病的人为研究对象，以个体为单位随机分配措施，多为预防性试验，需较多研究对象，尽可能应用盲法。

\ContentBand{对应往年题}

\begin{pastquestion}

\QuestionLabel{【往年题1｜20级回忆试题库.xlsx｜综合分析（题号9、30）】}\\
EV71疫苗的现场试验：现场和对象选择的原则有哪些？本题干描述是否符合？计算和评价；结论？

\end{pastquestion}

\begin{pastquestion}

\QuestionLabel{【往年题2｜11级流行病学口试考题回忆.pdf｜计算/设计】}\\
疫苗的评价一个现场研究，选择的现场和人群是否合适？（考实验研究现场和人群的选择原则）计算各种评价疫苗的指标并分析

\end{pastquestion}

\begin{pastquestion}

\QuestionLabel{【往年题3｜21级流病考题回忆汇总.xlsx｜综合分析（题号24）】}\\
四价流感疫苗，实验组新疫苗，对照组旧疫苗：算老年组和非老年组的保护率；讨论保护率差异可能的原因

\end{pastquestion}

\PointSeparator

\section{实验研究与观察性研究、前瞻性队列研究的异同}\label{users__fpb__26_spring__pkusph_study__tmp__pdfs__prevention-medicine-past-exams__03-ux6d41ux884cux75c5ux5b66ux5f80ux5e74ux9898ux8003ux70b9ux6574ux7406.md__ux5b9eux9a8cux7814ux7a76ux4e0eux89c2ux5bdfux6027ux7814ux7a76ux524dux77bbux6027ux961fux5217ux7814ux7a76ux7684ux5f02ux540c}

\FrequencyPill{8 次}

\InfoLabel{对应小节} 第六章 实验流行病学 / 第一节 概述 / 三、基本特征和用途（书中页码 第82---84页）

\InfoLabel{匹配依据} 05、10、11、12、14、17、21
级考``观察性研究和实验性研究的区别/不同点''；10、11、12、14、17、21
级考``实验研究和前瞻性队列研究的异同''。两类对比合并为一个考点。

\ContentBand{知识点原文摘取}

``观察''是在不干预、自然的情况下认识现象；``实验''则是在研究者的控制下，对研究对象人为施加或去除某种因素，观察其改变，评价这些人为措施的效果。

实验流行病学的基本特征：前瞻性、人为施加干预、随机分组、设平行对照组。其与前瞻性队列研究的关键区别在于：实验研究的干预是人为给予并随机分配的，而队列研究的暴露是自然存在、非随机的（观察法）。由于干预由研究者控制、随机分组，实验研究检验因果关系假设的能力强于分析性（观察性）研究，可作为确证因果关系的最终手段。两者相同点：均为前瞻性、由因及果、设对照组、可计算发病率与
RR 等。

\ContentBand{对应往年题}

\begin{pastquestion}

\QuestionLabel{【往年题1｜05流病考题回忆.doc｜简答】}\\
观察性研究和实验性研究的区别？（追问：观察性研究和实验性研究的种类？实验性研究最大的两个特点是什么？）

\end{pastquestion}

\begin{pastquestion}

\QuestionLabel{【往年题2｜10级考题回忆.pdf｜简答】}\\
比较实验研究和前瞻性队列研究的异同；实验性研究观察性研究的不同

\end{pastquestion}

\begin{pastquestion}

\QuestionLabel{【往年题3｜21级流病考题回忆汇总.xlsx｜简答（题号28、46）】}\\
实验研究和前瞻性队列研究的异同；RCT的设计原则？RCT在实验流行病学中的优势？

\end{pastquestion}

\PointSeparator

\section{实验流行病学的类型与特点用途（RCT、临床试验、现场试验、社区试验、类实验）}\label{users__fpb__26_spring__pkusph_study__tmp__pdfs__prevention-medicine-past-exams__03-ux6d41ux884cux75c5ux5b66ux5f80ux5e74ux9898ux8003ux70b9ux6574ux7406.md__ux5b9eux9a8cux6d41ux884cux75c5ux5b66ux7684ux7c7bux578bux4e0eux7279ux70b9ux7528ux9014rctux4e34ux5e8aux8bd5ux9a8cux73b0ux573aux8bd5ux9a8cux793eux533aux8bd5ux9a8cux7c7bux5b9eux9a8c}

\FrequencyPill{8 次}

\InfoLabel{对应小节} 第六章 实验流行病学 / 第一节 概述 / 三、基本特征和用途、四、主要类型（书中页码
第83---86页）

\InfoLabel{匹配依据} 名解层面，04、05、10、11、12、14、20、21
级广泛考``RCT（随机对照试验）、社区试验（community
trial）、类实验（quasi-experiment）、实验流行病学''；简答层面考``实验流行病学的特点和主要用途''。合并为一个考点。

\ContentBand{知识点原文摘取}

通常把实验流行病学研究分为临床试验、现场试验和社区试验三类。

临床试验（clinical
trial）：以病人个体为单位分组；随机化临床试验（RCT，又称随机对照试验）应用最广。现场试验（field
trial）：以尚未患病的人为研究对象，以个体为单位分配措施，多为预防性试验。社区试验（community
trial）：以人群作为整体进行实验，一般采用整群随机分配。

类实验（quasi-experiment）：缺乏严格的随机分组或平行对照的实验，又称自然实验，可分为不设平行对照组（如自身前后对照）和设对照组但非随机分组两类，本质上仍是观察性研究。

实验流行病学主要用于评价疾病防治效果（如疫苗预防效果、生活方式干预、保健策略评价、治疗方案评价）。基本特征：前瞻性、施加干预、随机分组、设平行对照组。

\ContentBand{对应往年题}

\begin{pastquestion}

\QuestionLabel{【往年题1｜11级流行病学口试考题回忆.pdf｜名解】}\\
RCT；community trial（社区试验）；类实验

\end{pastquestion}

\begin{pastquestion}

\QuestionLabel{【往年题2｜14级流病考题回忆汇总.xls｜简答（题号9、61）】}\\
流行病学实验特点用途；实验流行病学的特点和主要用途

\end{pastquestion}

\begin{pastquestion}

\QuestionLabel{【往年题3｜12级流病口试考题回忆.pdf｜名解（题号13、36）】}\\
RCT；类实验

\end{pastquestion}

\PointSeparator

\section{不依从与意向治疗分析（ITT）、符合方案分析、接受治疗分析}\label{users__fpb__26_spring__pkusph_study__tmp__pdfs__prevention-medicine-past-exams__03-ux6d41ux884cux75c5ux5b66ux5f80ux5e74ux9898ux8003ux70b9ux6574ux7406.md__ux4e0dux4f9dux4eceux4e0eux610fux5411ux6cbbux7597ux5206ux6790ittux7b26ux5408ux65b9ux6848ux5206ux6790ux63a5ux53d7ux6cbbux7597ux5206ux6790}

\FrequencyPill{6 次}

\InfoLabel{对应小节} 第六章 实验流行病学 / 第三节 资料的整理与分析 /
一、资料的整理（二）退出；二、资料的分析（一）意向治疗分析（书中页码 第94---96页）

\InfoLabel{匹配依据} 10、11、12、17、19、20 级考``不依从（noncompliance）''名解及``ITT
意向性分析/符合方案（依从者）分析/接受治疗分析的概念、区别、缺点''。

\ContentBand{知识点原文摘取}

意向治疗分析（intention-to-treat
analysis，ITT）：所有病人被随机分入任意一组，不管他们是否完成试验或是否真正接受了该组治疗，都保留在原组进行结果分析。目的在于避免选择偏倚，并使各治疗组之间保持可比性。

效力分析（即依从者分析/符合方案分析）：只比较真正完成各自治疗的病人，忽略未完成或转组者。接受治疗分析：按研究对象实际接受的治疗分组比较。

三种分析方法结果常不一致：ITT 反映实际临床应用后的效果（若试验确实有效，ITT
可能低估治疗效果）；依从者分析和接受治疗分析在不依从不均衡时会高估治疗效果。评价试验效力时建议同时使用三种分析。

\ContentBand{对应往年题}

\begin{pastquestion}

\QuestionLabel{【往年题1｜19级流行病学期末口试回忆.xlsx｜简答/综合分析（题号10、14）】}\\
意向性治疗分析的概念、意义、优缺点；ITT的意义、应用

\end{pastquestion}

\begin{pastquestion}

\QuestionLabel{【往年题2｜20级回忆试题库.xlsx｜综合分析（题号19、27、35）】}\\
ITT那三种分析方法是什么，缺点分别是什么；ITT、依从者分析、接受治疗分析计算

\end{pastquestion}

\begin{pastquestion}

\QuestionLabel{【往年题3｜12级流病口试考题回忆.pdf｜简答（题号47）】}\\
不依从的三种分析方法：意向性分析、遵循研究方案分析、接受干预措施分析

\end{pastquestion}

\PointSeparator

\section{盲法（单盲、双盲、三盲）及其优缺点}\label{users__fpb__26_spring__pkusph_study__tmp__pdfs__prevention-medicine-past-exams__03-ux6d41ux884cux75c5ux5b66ux5f80ux5e74ux9898ux8003ux70b9ux6574ux7406.md__ux76f2ux6cd5ux5355ux76f2ux53ccux76f2ux4e09ux76f2ux53caux5176ux4f18ux7f3aux70b9}

\FrequencyPill{6 次}

\InfoLabel{对应小节} 第六章 实验流行病学 / 第二节 研究设计与实施 / 七、盲法的应用（书中页码 第91---92页）

\InfoLabel{匹配依据} 04、05、11、12、17、19、21
级考``单盲、双盲、三盲''名词解释及``三种盲法的特点及优缺点''，部分追问各种盲法想消除的偏倚。

\ContentBand{知识点原文摘取}

根据盲法程度分为单盲、双盲和三盲。单盲（single blind）指研究对象不知道自己是试验组还是对照组。双盲（double
blind）指研究对象和研究观察者都不了解试验分组情况，由研究设计者来安排和控制全部试验。三盲（triple
blind）指不但研究者和研究对象不了解分组情况，而且负责资料收集和分析的人员也不了解分组情况，从而较好地避免了偏倚。

未用盲法的试验称为开放性试验（open trial），优点是易设计和实施，主要缺点是容易产生偏倚。

\ContentBand{对应往年题}

\begin{pastquestion}

\QuestionLabel{【往年题1｜12级流病口试考题回忆.pdf｜简答（题号18）】}\\
单盲、双盲、三盲的特点及优缺点

\end{pastquestion}

\begin{pastquestion}

\QuestionLabel{【往年题2｜17级流病口试回忆.xlsx｜简答（题号12）】}\\
三种盲法及其优缺点（秦老师补充：分析资料人员也不知道分组）

\end{pastquestion}

\begin{pastquestion}

\QuestionLabel{【往年题3｜21级流病考题回忆汇总.xlsx｜名解（题号13、39）】}\\
三盲（Three Blind Method）

\end{pastquestion}

\PointSeparator

\section{安慰剂效应}\label{users__fpb__26_spring__pkusph_study__tmp__pdfs__prevention-medicine-past-exams__03-ux6d41ux884cux75c5ux5b66ux5f80ux5e74ux9898ux8003ux70b9ux6574ux7406.md__ux5b89ux6170ux5242ux6548ux5e94}

\FrequencyPill{6 次}

\InfoLabel{对应小节} 第六章 实验流行病学 / 第二节 研究设计与实施 / 六、设立对照（一）设立对照的必要性 /
4.安慰剂效应（书中页码 第91页）

\InfoLabel{匹配依据} 05、10、11、12、19、20 级考``安慰剂效应（placebo effect）''名词解释。

\ContentBand{知识点原文摘取}

安慰剂效应（placebo
effect）是指某些疾病的病人由于依赖医药而表现出的一种正向心理效应，因此当以主观症状的改善情况作为疗效评价指标时，其``效应''中可能包括有安慰剂效应在内。

安慰剂（placebo）是一种无论在外观、颜色、味觉、嗅觉上均与积极治疗的药品无从辨别的物品，但没有特定已知的治疗成分。目前已知的安慰剂可使三分之一的病人增强信心、减轻病情、减少不适症状，这一现象称为安慰剂效应。

\ContentBand{对应往年题}

\begin{pastquestion}

\QuestionLabel{【往年题1｜12级流病口试考题回忆.pdf｜名解（题号4、43）】}\\
安慰剂效应

\end{pastquestion}

\begin{pastquestion}

\QuestionLabel{【往年题2｜20级回忆试题库.xlsx｜名解（题号46）】}\\
placebo effects（安慰剂效应）

\end{pastquestion}

\begin{pastquestion}

\QuestionLabel{【往年题3｜19级流行病学期末口试回忆.xlsx｜名解（题号12、22）】}\\
安慰剂效应 / placebo

\end{pastquestion}

\PointSeparator

\section{需治疗人数（NNT）}\label{users__fpb__26_spring__pkusph_study__tmp__pdfs__prevention-medicine-past-exams__03-ux6d41ux884cux75c5ux5b66ux5f80ux5e74ux9898ux8003ux70b9ux6574ux7406.md__ux9700ux6cbbux7597ux4ebaux6570nnt}

\FrequencyPill{4 次}

\InfoLabel{对应小节} 第六章 实验流行病学 / 第三节 资料的整理与分析 / 三、评价指标 / 3.需治疗人数（书中页码
第97---98页）

\InfoLabel{匹配依据} 10、11、12、19 级考``需治疗人数（number needed to
treat，NNT）''名词解释，部分追问其公共卫生学意义。

\ContentBand{知识点原文摘取}

需治疗人数（number needed to
treat，NNT）指为预防1例不良事件发生，临床医师在一段时间内应用某一疗法需治疗的病人数。从数学关系上讲，NNT等于绝对危险度减少值（ARR）的倒数，即
NNT=1/\textbar CER−EER\textbar。NNT 具有直观易懂、操作方便、可指导个体病人临床决策等优点。

\ContentBand{对应往年题}

\begin{pastquestion}

\QuestionLabel{【往年题1｜11级流行病学口试考题回忆.pdf｜名解】}\\
需要治疗人数（NNT）

\end{pastquestion}

\begin{pastquestion}

\QuestionLabel{【往年题2｜19级流行病学期末口试回忆.xlsx｜名解（题号9、28）】}\\
NNT / 需治疗人数（英文）（追问 NNT 的意义------公共卫生学意义，预防）

\end{pastquestion}

\PointSeparator

\section{实验研究设立对照的原因（必要性）及对照类型}\label{users__fpb__26_spring__pkusph_study__tmp__pdfs__prevention-medicine-past-exams__03-ux6d41ux884cux75c5ux5b66ux5f80ux5e74ux9898ux8003ux70b9ux6574ux7406.md__ux5b9eux9a8cux7814ux7a76ux8bbeux7acbux5bf9ux7167ux7684ux539fux56e0ux5fc5ux8981ux6027ux53caux5bf9ux7167ux7c7bux578b}

\FrequencyPill{3 次}

\InfoLabel{对应小节} 第六章 实验流行病学 / 第二节 研究设计与实施 / 六、设立对照（书中页码 第90---92页）

\InfoLabel{匹配依据} 04、14、21
级考``实验研究设对照的原因及方式/对照的设置原因和方式''，对应课本设立对照的必要性（不能预知的结局、向均数回归、霍桑效应、安慰剂效应、潜在未知因素）及对照类型。

\ContentBand{知识点原文摘取}

合理的对照能将干预措施的真实效应客观、充分地暴露出来。设立对照的必要性源于干预效应受以下因素影响：①不能预知的结局；②向均数回归；③霍桑效应；④安慰剂效应；⑤潜在的未知因素的影响。

对照类型：①标准疗法对照（有效对照）；②安慰剂对照；③自身对照；④交叉对照（研究对象随机分两组，先后互换干预与对照，每个对象均兼作试验组和对照组，所需样本量小于平行对照，但要求第一阶段干预不影响第二阶段）。

\ContentBand{对应往年题}

\begin{pastquestion}

\QuestionLabel{【往年题1｜21级流病考题回忆汇总.xlsx｜简答（题号42）】}\\
实验研究的对照组选取原因和方法

\end{pastquestion}

\begin{pastquestion}

\QuestionLabel{【往年题2｜14级流病考题回忆汇总.xls｜简答（题号57）】}\\
实验研究设对照的原因

\end{pastquestion}

\PointSeparator

\chapter{筛检}\label{users__fpb__26_spring__pkusph_study__tmp__pdfs__prevention-medicine-past-exams__03-ux6d41ux884cux75c5ux5b66ux5f80ux5e74ux9898ux8003ux70b9ux6574ux7406.md__ux7b2cux4e03ux7ae0-ux7b5bux68c0}

\section{真实性指标的计算：灵敏度、特异度、约登指数、预测值、假阳性率、似然比}\label{users__fpb__26_spring__pkusph_study__tmp__pdfs__prevention-medicine-past-exams__03-ux6d41ux884cux75c5ux5b66ux5f80ux5e74ux9898ux8003ux70b9ux6574ux7406.md__ux771fux5b9eux6027ux6307ux6807ux7684ux8ba1ux7b97ux7075ux654fux5ea6ux7279ux5f02ux5ea6ux7ea6ux767bux6307ux6570ux9884ux6d4bux503cux5047ux9633ux6027ux7387ux4f3cux7136ux6bd4}

\FrequencyPill{11 次}

\InfoLabel{对应小节} 第七章 筛检 / 第二节 筛检试验的评价 /
二、筛检试验的评价（一）真实性、（三）预测值（书中页码 第104---108页）

\InfoLabel{匹配依据} 04、05、07、10、11、12、14、17、19、20、21
级\textbf{每一届}都考四格表筛检计算------灵敏度、特异度、约登指数（正确指数）、阳性预测值、阴性预测值、假阳性率（误诊率）、假阴性率（漏诊率）、一致率，常追问各指标含义、如何提高灵敏度/特异度、选哪种方法用于人群筛查。本课程最高频计算考点之一。

\ContentBand{知识点原文摘取}

灵敏度（sensitivity，真阳性率）=\(\dfrac{TP}{TP+FN}\times100\%\)，反映筛检发现病人的能力；假阴性率（漏诊率）=1−灵敏度。

特异度（specificity，真阴性率）=\(\dfrac{TN}{FP+TN}\times100\%\)，反映鉴别排除非病例的能力；假阳性率（误诊率）=1−特异度。

正确指数（约登指数，Youden's index）=（灵敏度+特异度）−1，范围0～1，指数越大真实性越高。

阳性似然比 +LR=灵敏度/(1−特异度)；阴性似然比 −LR=(1−灵敏度)/特异度。

阳性预测值（Pr+）=\(\dfrac{TP}{TP+FP}\times100\%\)；阴性预测值（Pr−）=\(\dfrac{TN}{TN+FN}\times100\%\)。

\ContentBand{对应往年题}

\begin{pastquestion}

\QuestionLabel{【往年题1｜19级流行病学期末口试回忆.xlsx｜计算（题号25）】}\\
四格表计算灵敏度、特异度、预测值

\end{pastquestion}

\begin{pastquestion}

\QuestionLabel{【往年题2｜20级回忆试题库.xlsx｜综合分析（题号21）】}\\
四格表计算灵敏度、特异度、约登指数、预测值

\end{pastquestion}

\begin{pastquestion}

\QuestionLabel{【往年题3｜10级考题回忆.pdf｜计算】}\\
计算 sen/spe/youden index/ppv/npv 并解释各个指标的含义

\end{pastquestion}

\begin{pastquestion}

\QuestionLabel{【往年题4｜04级流病考题回忆.docx｜计算】}\\
灵敏度、特异度及联合应用；阳性预测值、阴性预测值的计算

\end{pastquestion}

\PointSeparator

\section{串联试验与并联试验的灵敏度、特异度计算及选择}\label{users__fpb__26_spring__pkusph_study__tmp__pdfs__prevention-medicine-past-exams__03-ux6d41ux884cux75c5ux5b66ux5f80ux5e74ux9898ux8003ux70b9ux6574ux7406.md__ux4e32ux8054ux8bd5ux9a8cux4e0eux5e76ux8054ux8bd5ux9a8cux7684ux7075ux654fux5ea6ux7279ux5f02ux5ea6ux8ba1ux7b97ux53caux9009ux62e9}

\FrequencyPill{10 次}

\InfoLabel{对应小节} 第七章 筛检 / 第三节 筛检效果的评价 / 一、收益（三）联合试验对收益的影响（书中页码
第112---113页）

\InfoLabel{匹配依据} 04、05、07、10、11、12、14、19、20、21 级反复考``两种筛检方法（如尿糖与血糖、FOB 与
TF、大肠癌两种检测）的串联/并联试验的灵敏度、特异度计算，及如何在此基础上提高灵敏度、特异度''，并追问临床上应选哪种。

\ContentBand{知识点原文摘取}

串联试验（serial
test，序列试验）：一组筛检试验按顺序相连，全部结果均为阳性者才定为阳性。相较单一方法，串联试验可提高特异度和阳性检出率，但会损失一定的灵敏度。适用于筛检人群基数大，但发病率和患病率较低的疾病。

并联试验（parallel
test，平行试验）：全部筛检试验同时平行开展，任何一项阳性即判为阳性。优点是弥补两种方法灵敏度都不足的问题，提高筛检整体的灵敏度，但会降低特异度。

\ContentBand{对应往年题}

\begin{pastquestion}

\QuestionLabel{【往年题1｜05流病考题回忆.doc｜计算】}\\
算尿糖、血糖筛查方法各自的灵敏度、特异度，如果要提高灵敏度、特异度该怎么办，结果如何？（串、并联）

\end{pastquestion}

\begin{pastquestion}

\QuestionLabel{【往年题2｜20级回忆试题库.xlsx｜综合分析（题号23）】}\\
串联、并联试验灵敏度和特异度的计算，大肠癌选哪一种（联合试验从临床角度考虑）

\end{pastquestion}

\begin{pastquestion}

\QuestionLabel{【往年题3｜11级流行病学口试考题回忆.pdf｜计算】}\\
大肠癌便潜血和隐蛋白两个筛检试验，计算串并联实验的灵敏度和特异度，根据计算结果能否选择试验方法，并说明原因

\end{pastquestion}

\PointSeparator

\section{筛检试验的评价（真实性、可靠性、预测值三个方面）}\label{users__fpb__26_spring__pkusph_study__tmp__pdfs__prevention-medicine-past-exams__03-ux6d41ux884cux75c5ux5b66ux5f80ux5e74ux9898ux8003ux70b9ux6574ux7406.md__ux7b5bux68c0ux8bd5ux9a8cux7684ux8bc4ux4ef7ux771fux5b9eux6027ux53efux9760ux6027ux9884ux6d4bux503cux4e09ux4e2aux65b9ux9762}

\FrequencyPill{8 次}

\InfoLabel{对应小节} 第七章 筛检 / 第二节 筛检试验的评价 / 二、筛检试验的评价（书中页码 第104页）

\InfoLabel{匹配依据} 05、10、14、17、19、20、21
级考``如何评价筛检试验/筛检试验评价包括哪些方面/筛检实验的核心指标及其意义''，对应课本评价的三大方面。

\ContentBand{知识点原文摘取}

筛检试验方法是否有效是开展筛检项目的基础，评价内容包含：真实性、可靠性和预测值。

真实性（validity，效度）：测量值与实际值相符合的程度，评价指标有灵敏度与假阴性率、特异度与假阳性率、正确指数和似然比。

可靠性（reliability，信度/精确度/可重复性）：在相同条件下重复测量同一批受试者时结果的一致程度，指标有符合率、Kappa
值等。

预测值（predictive
value）：应用筛检结果的阳性或阴性来估计受检者为病人或非病人可能性的指标，分阳性预测值和阴性预测值，反映筛检实际应用到人群后的收益大小。

\ContentBand{对应往年题}

\begin{pastquestion}

\QuestionLabel{【往年题1｜14级流病考题回忆汇总.xls｜简答（题号5）】}\\
筛检方法的评价（真实性、可靠性、预测度）

\end{pastquestion}

\begin{pastquestion}

\QuestionLabel{【往年题2｜19级流行病学期末口试回忆.xlsx｜简答（题号8）】}\\
筛检实验的核心指标及其意义

\end{pastquestion}

\begin{pastquestion}

\QuestionLabel{【往年题3｜20级回忆试题库.xlsx｜简答（题号27、44）】}\\
筛检试验的指标和评价方法；如何评价筛检方法

\end{pastquestion}

\PointSeparator

\section{影响筛检试验可靠性的因素}\label{users__fpb__26_spring__pkusph_study__tmp__pdfs__prevention-medicine-past-exams__03-ux6d41ux884cux75c5ux5b66ux5f80ux5e74ux9898ux8003ux70b9ux6574ux7406.md__ux5f71ux54cdux7b5bux68c0ux8bd5ux9a8cux53efux9760ux6027ux7684ux56e0ux7d20}

\FrequencyPill{5 次}

\InfoLabel{对应小节} 第七章 筛检 / 第二节 筛检试验的评价 / 二、筛检试验的评价（二）可靠性 /
2.影响筛检试验可靠性的因素（书中页码 第107页）

\InfoLabel{匹配依据} 10、11、12、14、21
级考``影响筛检（实验）可靠性的因素''。注意此处易答偏，应紧扣可靠性（重复测量一致性）而非真实性。

\ContentBand{知识点原文摘取}

影响筛检试验可靠性的因素：\\
（1）受试对象生物学变异：由于个体的生物周期等生物学变异，同一受试对象在不同时间获得的临床测量值有所波动（如血压在一天内不同时间的测量值存在变异）。\\
（2）观察者之间差异：测量者之间技术水平不一致，同一测量者在不同时间判断尺度不同，均可导致重复测量结果不一致（如不同阅片者报告的
X 线片结果不同）。\\
（3）实验室条件不一致：测量仪器不稳定、试验方法本身不稳定、不同厂家/批号试剂盒差异等可引起测量误差。

\ContentBand{对应往年题}

\begin{pastquestion}

\QuestionLabel{【往年题1｜11级流行病学口试考题回忆.pdf｜简答】}\\
影响筛检实验可靠性的因素

\end{pastquestion}

\begin{pastquestion}

\QuestionLabel{【往年题2｜12级流病口试考题回忆.pdf｜简答（题号30）】}\\
影响筛检可靠性的因素

\end{pastquestion}

\begin{pastquestion}

\QuestionLabel{【往年题3｜21级流病考题回忆汇总.xlsx｜简答（题号16）】}\\
影响筛检可靠性的因素

\end{pastquestion}

\PointSeparator

\section{灵敏度、特异度、患病率与预测值的关系（间接计算法）}\label{users__fpb__26_spring__pkusph_study__tmp__pdfs__prevention-medicine-past-exams__03-ux6d41ux884cux75c5ux5b66ux5f80ux5e74ux9898ux8003ux70b9ux6574ux7406.md__ux7075ux654fux5ea6ux7279ux5f02ux5ea6ux60a3ux75c5ux7387ux4e0eux9884ux6d4bux503cux7684ux5173ux7cfbux95f4ux63a5ux8ba1ux7b97ux6cd5}

\FrequencyPill{5 次}

\InfoLabel{对应小节} 第七章 筛检 / 第二节 筛检试验的评价 / 二（三）预测值 /
2.间接计算法、3.预测值与真实性指标、患病率的关系（书中页码 第108页）

\InfoLabel{匹配依据} 04、05、19、20
级考``给灵敏度特异度和患病率，计算阳性/阴性预测值；改变患病率后预测值如何变化；社区人群与高危人群的预测值差异''，对应课本
Bayes 公式及预测值随患病率变化的规律。

\ContentBand{知识点原文摘取}

间接计算法（Bayes 公式）：

\[阳性预测值=\frac{灵敏度\times{}患病率}{灵敏度\times{}患病率+(1-患病率)\times(1-特异度)}\]

\[阴性预测值=\frac{特异度\times(1-患病率)}{特异度\times(1-患病率)+(1-灵敏度)\times{}患病率}\]

关系规律：①当灵敏度与特异度一定时，疾病患病率增加，阳性预测值增加，阴性预测值降低；②当人群患病率不变时，灵敏度升高、特异度降低，阳性预测值下降，阴性预测值增加。因此在高危（高患病率）人群中筛检可提高阳性预测值。

\ContentBand{对应往年题}

\begin{pastquestion}

\QuestionLabel{【往年题1｜05流病考题回忆.doc｜计算】}\\
筛检计算阳性和阴性预测值，又给了个发病率高的问阳阴预测值如何变

\end{pastquestion}

\begin{pastquestion}

\QuestionLabel{【往年题2｜19级流行病学期末口试回忆.xlsx｜综合分析（题号29）】}\\
社区人群和高危人群的阴性预测值和阳性预测值计算；患病率计算（社区人群需带入患病率、灵敏度、特异度那个间接算法公式）

\end{pastquestion}

\begin{pastquestion}

\QuestionLabel{【往年题3｜04级流病考题回忆.docx｜计算】}\\
阳性预测值、阴性预测值的计算；改变患病率以后算阳性预测值

\end{pastquestion}

\PointSeparator

\section{筛检试验与诊断试验的区别}\label{users__fpb__26_spring__pkusph_study__tmp__pdfs__prevention-medicine-past-exams__03-ux6d41ux884cux75c5ux5b66ux5f80ux5e74ux9898ux8003ux70b9ux6574ux7406.md__ux7b5bux68c0ux8bd5ux9a8cux4e0eux8bcaux65adux8bd5ux9a8cux7684ux533aux522b}

\FrequencyPill{4 次}

\InfoLabel{对应小节} 第七章 筛检 / 第二节 筛检试验的评价 / 一、筛检试验的定义（表7-1）（书中页码 第101页）

\InfoLabel{匹配依据} 10、11、19、20 级考``筛检试验和诊断试验的区别''。

\ContentBand{知识点原文摘取}

筛检试验和诊断试验的目的、对象、结果判读及后续处理都不相同：筛检试验目的是区分可能患病的个体与可能未患病者，对象是健康的人或无明显临床症状的病人，要求快速、简便、无创、有高灵敏度，阳性者须进一步通过诊断试验确诊；诊断试验目的是明确判断病人与未患病的人，对象是病人或筛检阳性者，要求灵敏度和特异度高、结果更准确权威，阳性者要及时治疗。筛检的准确性要求不如诊断性试验高，允许存在一定比例的错判。

\ContentBand{对应往年题}

\begin{pastquestion}

\QuestionLabel{【往年题1｜11级流行病学口试考题回忆.pdf｜简答】}\\
筛检实验和诊断实验区别

\end{pastquestion}

\begin{pastquestion}

\QuestionLabel{【往年题2｜20级回忆试题库.xlsx｜简答（题号33）】}\\
筛检试验和诊断试验的区别

\end{pastquestion}

\begin{pastquestion}

\QuestionLabel{【往年题3｜19级流行病学期末口试回忆.xlsx｜综合分析（题号20）】}\\
筛检实验的定义，和诊断试验的区别

\end{pastquestion}

\PointSeparator

\section{筛检效果评价中的偏倚（领先时间偏倚、病程偏倚、志愿者偏倚、过度诊断偏倚）}\label{users__fpb__26_spring__pkusph_study__tmp__pdfs__prevention-medicine-past-exams__03-ux6d41ux884cux75c5ux5b66ux5f80ux5e74ux9898ux8003ux70b9ux6574ux7406.md__ux7b5bux68c0ux6548ux679cux8bc4ux4ef7ux4e2dux7684ux504fux501aux9886ux5148ux65f6ux95f4ux504fux501aux75c5ux7a0bux504fux501aux5fd7ux613fux8005ux504fux501aux8fc7ux5ea6ux8bcaux65adux504fux501a}

\FrequencyPill{5 次}

\InfoLabel{对应小节} 第七章 筛检 / 第三节 筛检效果的评价 / 二、生物学效果评价（三）偏倚及其控制（书中页码
第114---115页）

\InfoLabel{匹配依据} 04、10、19、20、21 级考``筛检的偏倚有哪些/筛检试验的常见偏倚和解释''，名解层面 19、21
级考``领先时间偏倚（lead time bias）''。

\ContentBand{知识点原文摘取}

\begin{enumerate}
\def\labelenumi{\arabic{enumi}.}
\item
  领先时间偏倚（lead time
  bias）：领先时间是临床前筛检诊断的时点至常规临床诊断时点之间的时间间隔。如果筛检只提前了发现疾病的时点，而并未改变死亡时点，也会观察到筛检人群比不筛检人群生存时间更长的假象。以生命年评价筛检效果时，应扣除领先时间，否则会高估筛检效果。
\item
  病程偏倚（length
  bias）：疾病被检出的可能性和疾病的进展速度有关。如果筛检组中疾病进展缓慢的病人占较大比例，可能观察到筛检组生存概率更高或生存时间更长，从而高估筛检效果。
\item
  志愿者偏倚（volunteer
  bias）：参加筛检者往往受教育程度更高、更关注健康、不良行为较少，总的发病或死亡风险可能低于不参加者，导致筛检效果被高估。
\item
  过度诊断偏倚（over diagnosis
  bias）：筛检发现的``惰性病例''等被计入病人总体，导致经筛检发现的病人有较多生存者或较长平均生存期的假象，高估筛检效果。
\end{enumerate}

\ContentBand{对应往年题}

\begin{pastquestion}

\QuestionLabel{【往年题1｜10级考题回忆.pdf｜简答】}\\
筛检偏倚的类型

\end{pastquestion}

\begin{pastquestion}

\QuestionLabel{【往年题2｜21级流病考题回忆汇总.xlsx｜名解/简答（题号44、4、24）】}\\
领先时间偏倚（外暴露标志物）；筛检中的偏倚；筛检试验的常见偏倚和解释

\end{pastquestion}

\begin{pastquestion}

\QuestionLabel{【往年题3｜19级流行病学期末口试回忆.xlsx｜简答（题号22）】}\\
筛检的偏倚

\end{pastquestion}

\PointSeparator

\section{筛检的概念、目的、类型与实施原则}\label{users__fpb__26_spring__pkusph_study__tmp__pdfs__prevention-medicine-past-exams__03-ux6d41ux884cux75c5ux5b66ux5f80ux5e74ux9898ux8003ux70b9ux6574ux7406.md__ux7b5bux68c0ux7684ux6982ux5ff5ux76eeux7684ux7c7bux578bux4e0eux5b9eux65bdux539fux5219}

\FrequencyPill{5 次}

\InfoLabel{对应小节} 第七章 筛检 / 第一节 概述 /
一、筛检的概念、二、筛检的目的及类型、三、筛检的实施原则（书中页码 第101---103页）

\InfoLabel{匹配依据} 04、05、12、19、20 级考``筛检（screening）/选择性筛检（selective
screening）/整群筛检（mass screening）''名解，以及``筛检的条件/原则/目的''。

\ContentBand{知识点原文摘取}

筛检（screening）是针对临床前期或早期的疾病阶段，运用快速、简便的试验、检查或其他方法，将未察觉或未诊断疾病的人群中那些可能有病或缺陷、但表面健康的个体，同那些可能无病者鉴别开来的一系列医疗卫生服务措施。

按筛检对象范围分为整群筛检（mass screening，对一定范围内人群全体进行无差异普通筛检）和选择性筛检（selective
screening，又称高危人群筛检，选择疾病的高危人群进行筛检）。

Wilson 和 Jungner（1968）提出实施筛检项目的 10
条原则，核心包括：被筛检疾病是重要的卫生问题、检出的病人有相应治疗方法、具备诊断治疗设施、疾病处于潜伏期或早期、有适宜且为人群接受的检查方法、对疾病自然史有足够了解、成本与医疗保健支出总体持平等。

\ContentBand{对应往年题}

\begin{pastquestion}

\QuestionLabel{【往年题1｜04级流病考题回忆.docx｜问答】}\\
筛选的条件；筛检的原则

\end{pastquestion}

\begin{pastquestion}

\QuestionLabel{【往年题2｜20级回忆试题库.xlsx｜名解（题号37、49）】}\\
mass screening（整群筛检）；selective screening（选择性筛检）

\end{pastquestion}

\begin{pastquestion}

\QuestionLabel{【往年题3｜12级流病口试考题回忆.pdf｜名解（题号63）】}\\
筛检

\end{pastquestion}

\PointSeparator

\chapter{偏倚及其控制}\label{users__fpb__26_spring__pkusph_study__tmp__pdfs__prevention-medicine-past-exams__03-ux6d41ux884cux75c5ux5b66ux5f80ux5e74ux9898ux8003ux70b9ux6574ux7406.md__ux7b2cux516bux7ae0-ux504fux501aux53caux5176ux63a7ux5236}

\begin{quote}
说明：课本将三大类偏倚分散于第三章（现况研究）、第四章（队列研究）、第五章（病例对照研究）等章节，本章按授课顺序集中整理。
\end{quote}

\section{入院率偏倚（伯克森偏倚）及其控制（含社区/医院病例对照 OR
差异大题）}\label{users__fpb__26_spring__pkusph_study__tmp__pdfs__prevention-medicine-past-exams__03-ux6d41ux884cux75c5ux5b66ux5f80ux5e74ux9898ux8003ux70b9ux6574ux7406.md__ux5165ux9662ux7387ux504fux501aux4f2fux514bux68eeux504fux501aux53caux5176ux63a7ux5236ux542bux793eux533aux533bux9662ux75c5ux4f8bux5bf9ux7167-or-ux5deeux5f02ux5927ux9898}

\FrequencyPill{8 次}

\InfoLabel{对应小节} 第五章 病例对照研究 / 第四节 常见偏倚及其控制 / 一、选择偏倚（一）入院率偏倚（书中页码
第79页）

\InfoLabel{匹配依据} 10、12、14、17、19、20、21 级反复考``入院率偏倚（admission rate
bias）/伯克森偏倚（Berkson's bias）''名解；并以大题形式出现：``社区来源和医院来源的两个病例对照研究 OR
值有差异，原因是什么？如何避免入院率偏倚？计算偏倚的大小与方向。''

\ContentBand{知识点原文摘取}

入院率偏倚（admission rate bias）也称伯克森偏倚（Berkson's
bias），在以医院为基础的病例对照研究中常发生，即当选择医院病人作为病例和对照时，病例只是该医院或某些医院的特定病例而非全体病人的随机样本，对照是医院的某一部分病人而非全体目标人群的随机样本，由于各种疾病的入院率不同可导致病例组与对照组在某些特征上的系统误差。

控制：尽可能在社区人群中选择病例和对照；如进行以医院为基础的病例对照研究，最好在多个不同级别、不同种类的医院选择病例，在与病例相同的多个医院、多个科室、多病种的病人中选择对照，因已知与所研究暴露因素有关的病种而就诊的病人不宜作为对照。

\ContentBand{对应往年题}

\begin{pastquestion}

\QuestionLabel{【往年题1｜17级流病口试回忆.xlsx｜综合分析（题号12）】}\\
病例对照社区和医院来源的 OR 差异，原因，计算偏倚大小方向，如何控制

\end{pastquestion}

\begin{pastquestion}

\QuestionLabel{【往年题2｜19级流行病学期末口试回忆.xlsx｜综合分析（题号8）】}\\
两个病例对照研究（一个社区一个医院），OR 值有差异，差异的原因是什么？如何避免入院率偏移，计算偏移的大小与方向

\end{pastquestion}

\begin{pastquestion}

\QuestionLabel{【往年题3｜20级回忆试题库.xlsx｜名解（题号9）】}\\
admission rate bias（入院率偏倚）

\end{pastquestion}

\PointSeparator

\section{混杂偏倚与混杂因素的判断及控制}\label{users__fpb__26_spring__pkusph_study__tmp__pdfs__prevention-medicine-past-exams__03-ux6d41ux884cux75c5ux5b66ux5f80ux5e74ux9898ux8003ux70b9ux6574ux7406.md__ux6df7ux6742ux504fux501aux4e0eux6df7ux6742ux56e0ux7d20ux7684ux5224ux65adux53caux63a7ux5236}

\FrequencyPill{8 次}

\InfoLabel{对应小节} 第四章 队列研究 / 第四节 三、混杂偏倚；第五章 病例对照研究 /
第三节（三）非匹配资料的分层分析（书中页码 第60、73---76页）

\InfoLabel{匹配依据} 04、05、10、12、14、19、21 级考``混杂偏倚（confounding
bias）''名解、``举例说明混杂偏倚如何控制''、``如何判断某因素为混杂因素（书上91页/通过列四格表确定）''及混杂因素判断计算题。

\ContentBand{知识点原文摘取}

混杂偏倚（confounding
bias）是指当研究某个因素与某种疾病的关联时，某个既与疾病有关、又与所研究的暴露因素有关的外来因素的影响，掩盖或夸大了所研究的暴露因素与疾病的联系所造成的偏倚。该外来因素称为混杂因素（confounding
factor）。

混杂因素的特征：①与疾病有关（是疾病的危险因素）；②与暴露因素有关；③不是暴露与疾病因果链上的中间环节。只有在暴露组与对照组中分布不均衡时混杂才会产生。

控制：研究设计阶段可采取限制、匹配（配比）等方式控制；资料分析阶段可采用分层分析、标准化或多因素分析等方法控制。（如口服避孕药与心肌梗死研究中，按年龄分层后两层
OR（2.80、2.78）均大于不分层时的 OR（2.20），\(OR_{MH}=2.79\)，说明年龄是混杂因素。）

\ContentBand{对应往年题}

\begin{pastquestion}

\QuestionLabel{【往年题1｜04级流病考题回忆.docx｜问答】}\\
分析性研究混杂的解释及其控制方法，要求举例；举例说明混杂偏倚，如何控制

\end{pastquestion}

\begin{pastquestion}

\QuestionLabel{【往年题2｜05流病考题回忆.doc｜计算】}\\
如何判断某因素为混杂因素（书上91页）

\end{pastquestion}

\begin{pastquestion}

\QuestionLabel{【往年题3｜21级流病考题回忆汇总.xlsx｜综合分析（题号16、41、50）】}\\
病例对照、判断混杂因素；非匹配分层病例对照研究，探究性别是否为混杂因素，分别计算总 OR 和分层 OR，根据 OR
判断是不是有混杂

\end{pastquestion}

\PointSeparator

\section{错分偏倚（无差异错分、有差异错分）及偏倚大小和方向的计算}\label{users__fpb__26_spring__pkusph_study__tmp__pdfs__prevention-medicine-past-exams__03-ux6d41ux884cux75c5ux5b66ux5f80ux5e74ux9898ux8003ux70b9ux6574ux7406.md__ux9519ux5206ux504fux501aux65e0ux5deeux5f02ux9519ux5206ux6709ux5deeux5f02ux9519ux5206ux53caux504fux501aux5927ux5c0fux548cux65b9ux5411ux7684ux8ba1ux7b97}

\FrequencyPill{7 次}

\InfoLabel{对应小节} 第四章 队列研究 / 第四节 二、信息偏倚；第三章（偏倚分类）（书中页码 第60、42页）

\InfoLabel{匹配依据} 10、11、12、14、19、20、21 级反复考``无差异错分（non-differential
misclassification）/有差异错分（differential
misclassification）''名解，以及``病例对照偏倚计算：无差异错分下计算校正
OR、偏倚的大小和方向''（多源自实习指导原题）。

\ContentBand{知识点原文摘取}

（课本将信息偏倚定义为``在收集信息过程中由于测量暴露与结局的方法有缺陷造成的系统误差''。课本知识点文件中未检索到``错分偏倚/无差异错分/有差异错分''的独立定义原文，该内容主要来自课件与实习指导，故不自行补充课本原文。建议结合实习指导回答：错分偏倚是信息偏倚的一种，指对暴露或疾病状态的错误分类；若错分在比较组间是均衡的（无差异错分），通常使关联趋于无效值（OR
趋向 1）；若错分在组间不均衡（有差异错分），可使关联被高估或低估，偏倚方向不定。偏倚大小可由``真实 OR 与观察
OR 之差/比''衡量。）

\ContentBand{对应往年题}

\begin{pastquestion}

\QuestionLabel{【往年题1｜12级流病口试考题回忆.pdf｜名解/计算（题号8、15）】}\\
无差异错分；无差异性错分的偏倚计算

\end{pastquestion}

\begin{pastquestion}

\QuestionLabel{【往年题2｜19级流行病学期末口试回忆.xlsx｜综合分析（题号26）】}\\
课件中无差异错分分类计算原题（计算校正 OR、偏倚大小等）

\end{pastquestion}

\begin{pastquestion}

\QuestionLabel{【往年题3｜20级回忆试题库.xlsx｜综合分析（题号38）】}\\
实习手册病例对照那个无差异错分的原题，计算校正 OR、偏倚大小等

\end{pastquestion}

\PointSeparator

\section{检出症候偏倚（暴露偏倚/揭露伪装偏倚）}\label{users__fpb__26_spring__pkusph_study__tmp__pdfs__prevention-medicine-past-exams__03-ux6d41ux884cux75c5ux5b66ux5f80ux5e74ux9898ux8003ux70b9ux6574ux7406.md__ux68c0ux51faux75c7ux5019ux504fux501aux66b4ux9732ux504fux501aux63edux9732ux4f2aux88c5ux504fux501a}

\FrequencyPill{7 次}

\InfoLabel{对应小节} 第五章 病例对照研究 / 第四节 一、选择偏倚（三）检出症候偏倚（书中页码 第79页）

\InfoLabel{匹配依据} 04、10、11、12、14、19、21 级考``检出症候偏倚（detection signal bias / unmasking
bias）''名词解释。

\ContentBand{知识点原文摘取}

某因素虽然不是所研究疾病的病因，但有该因素的个体容易出现某些症状或体征，并常因此而就医，从而提高了所研究疾病早期病例的检出率。如果病例对照研究中病例组包括了较多的这种早期病例，致使过高地估计了病例组的暴露程度，由此而产生的系统误差即为检出症候偏倚（detection
signal bias）。因此，在医院中收集病例时，最好包括不同来源的早、中、晚期病人，以减少这种偏倚。

\ContentBand{对应往年题}

\begin{pastquestion}

\QuestionLabel{【往年题1｜12级流病口试考题回忆.pdf｜名解（题号3）】}\\
检出症候偏倚

\end{pastquestion}

\begin{pastquestion}

\QuestionLabel{【往年题2｜21级流病考题回忆汇总.xlsx｜名解（题号50）】}\\
unmasking bias（检出症候偏倚）

\end{pastquestion}

\begin{pastquestion}

\QuestionLabel{【往年题3｜10级考题回忆.pdf｜名解】}\\
detection signal bias（检出症候偏倚）

\end{pastquestion}

\PointSeparator

\section{现患病例-新发病例偏倚（奈曼偏倚）}\label{users__fpb__26_spring__pkusph_study__tmp__pdfs__prevention-medicine-past-exams__03-ux6d41ux884cux75c5ux5b66ux5f80ux5e74ux9898ux8003ux70b9ux6574ux7406.md__ux73b0ux60a3ux75c5ux4f8b-ux65b0ux53d1ux75c5ux4f8bux504fux501aux5948ux66fcux504fux501a}

\FrequencyPill{6 次}

\InfoLabel{对应小节} 第五章 病例对照研究 / 第四节 一、选择偏倚（二）现患病例-新发病例偏倚（书中页码 第79页）

\InfoLabel{匹配依据} 04、05、10、19、20 级考``奈曼偏倚（Neyman
bias）/现患病例-新发病例偏倚（prevalence-incidence bias）''名词解释，常追问开放性分析。

\ContentBand{知识点原文摘取}

现患病例-新发病例偏倚（prevalence-incidence bias）也称奈曼偏倚（Neyman
bias），即如果调查对象选自现患病例（存活病例），特别是病程较长的现患病例，得到的一些暴露信息可能只与存活有关，而未必与该病的发病有关，从而错误地估计这些因素的病因作用；另一种情况是，某病的幸存者由于疾病而改变了原有的一些暴露特征，当他们被调查时容易误将这些改变了的暴露特征当作疾病前的状况。因此，选择新发病例作为研究对象可避免或减少此类偏倚。

\ContentBand{对应往年题}

\begin{pastquestion}

\QuestionLabel{【往年题1｜04级流病考题回忆.docx｜名解】}\\
奈曼偏倚（英文）（追问奈曼偏倚的开放性分析）

\end{pastquestion}

\begin{pastquestion}

\QuestionLabel{【往年题2｜19级流行病学期末口试回忆.xlsx｜名解（题号20）】}\\
prevalence-incidence bias（现患新发偏倚）

\end{pastquestion}

\begin{pastquestion}

\QuestionLabel{【往年题3｜20级回忆试题库.xlsx｜名解（题号17）】}\\
Neyman bias（奈曼偏倚）

\end{pastquestion}

\PointSeparator

\section{选择偏倚及其控制}\label{users__fpb__26_spring__pkusph_study__tmp__pdfs__prevention-medicine-past-exams__03-ux6d41ux884cux75c5ux5b66ux5f80ux5e74ux9898ux8003ux70b9ux6574ux7406.md__ux9009ux62e9ux504fux501aux53caux5176ux63a7ux5236}

\FrequencyPill{6 次}

\InfoLabel{对应小节} 第四章 队列研究 / 第四节 一、选择偏倚；第五章 第四节 一、选择偏倚（书中页码
第59---60、79页）

\InfoLabel{匹配依据} 10、11、12、14、17、20 级考``选择偏倚（selection bias）''名解及``选择偏倚的控制方法''。

\ContentBand{知识点原文摘取}

选择偏倚（selection
bias）是由于研究对象的选择不当，即入选的研究对象与未入选的研究对象在某些特征上存在差异而引起的系统误差，常发生在研究的设计阶段。病例对照研究中常见的选择偏倚包括入院率偏倚、现患病例-新发病例偏倚等；队列研究中失访偏倚本质上也是一种选择偏倚。

控制：选择偏倚一旦产生很难精确估计和有效处理，重在预防------严格按照标准选择研究对象，尽量从源人群中随机选取病例和对照，随访时尽量提高研究对象的应答率和依从性，减少失访。

\ContentBand{对应往年题}

\begin{pastquestion}

\QuestionLabel{【往年题1｜11级流行病学口试考题回忆.pdf｜简答】}\\
如何控制选择偏倚

\end{pastquestion}

\begin{pastquestion}

\QuestionLabel{【往年题2｜20级回忆试题库.xlsx｜简答（题号30、35）】}\\
选择偏倚的控制；控制选择偏移的方法

\end{pastquestion}

\begin{pastquestion}

\QuestionLabel{【往年题3｜14级流病考题回忆汇总.xls｜简答（题号11）】}\\
选择偏倚的控制方法

\end{pastquestion}

\PointSeparator

\section{霍桑效应}\label{users__fpb__26_spring__pkusph_study__tmp__pdfs__prevention-medicine-past-exams__03-ux6d41ux884cux75c5ux5b66ux5f80ux5e74ux9898ux8003ux70b9ux6574ux7406.md__ux970dux6851ux6548ux5e94}

\FrequencyPill{6 次}

\InfoLabel{对应小节} 第六章 实验流行病学 / 第二节 六、设立对照（一）设立对照的必要性 / 3.霍桑效应（书中页码
第91页）

\InfoLabel{匹配依据} 04、10、12、17、19、20 级考``霍桑效应（Hawthorne effect）''名词解释，10
级追问``霍桑效应为什么叫霍桑''。

\ContentBand{知识点原文摘取}

霍桑效应（Hawthorne
effect）：在试验研究（干预研究）中，受试者由于知道自己成为特殊被关注的对象后，所出现的改变自己行为或状态的一种倾向，与他们接受的干预措施的特异性作用无关，是由于病人渴望取悦于他们的医师，使医师感到其医疗活动是成功的。这是病人的一种心理、生理效应，会对疗效产生正向影响。当然，有时也可因病人厌恶医师或不信任某医院而产生负向效应。

\ContentBand{对应往年题}

\begin{pastquestion}

\QuestionLabel{【往年题1｜10级考题回忆.pdf｜名解】}\\
霍桑效应（老师追问霍桑效应为什么叫霍桑）

\end{pastquestion}

\begin{pastquestion}

\QuestionLabel{【往年题2｜17级流病口试回忆.xlsx｜名解（题号12）】}\\
霍桑效应

\end{pastquestion}

\begin{pastquestion}

\QuestionLabel{【往年题3｜19级流行病学期末口试回忆.xlsx｜名解（题号17）】}\\
霍桑效应

\end{pastquestion}

\PointSeparator

\section{信息偏倚及其控制（回忆偏倚、报告偏倚、调查偏倚、测量偏倚）}\label{users__fpb__26_spring__pkusph_study__tmp__pdfs__prevention-medicine-past-exams__03-ux6d41ux884cux75c5ux5b66ux5f80ux5e74ux9898ux8003ux70b9ux6574ux7406.md__ux4fe1ux606fux504fux501aux53caux5176ux63a7ux5236ux56deux5fc6ux504fux501aux62a5ux544aux504fux501aux8c03ux67e5ux504fux501aux6d4bux91cfux504fux501a}

\FrequencyPill{5 次}

\InfoLabel{对应小节} 第四章 队列研究 / 第四节 二、信息偏倚；第五章 第四节 二、信息偏倚（书中页码
第60、79---80页）

\InfoLabel{匹配依据} 04、11、12、14、20 级考``信息偏倚（information bias）及其控制方法''及``回忆偏倚（recall
bias）''名解。

\ContentBand{知识点原文摘取}

信息偏倚（information
bias）又称观察偏倚或测量偏倚，是在收集信息过程中由于测量暴露与结局的方法有缺陷造成的系统误差，主要发生在研究的实施阶段。病例对照研究中常见回忆偏倚、调查偏倚等。

回忆偏倚（recall
bias）：由研究对象对暴露史或既往史回忆的准确性和完整性存在系统误差而引起，是病例对照研究中最常见的信息偏倚。

控制：选择精确稳定的测量方法并在研究期间保持不变；校准仪器、严格遵守实验流程；做好调查员培训、统一调查口径；充分利用客观记录资料，问卷调查时重视提问方式；选择新发病例作为调查对象可减少回忆偏倚。

\ContentBand{对应往年题}

\begin{pastquestion}

\QuestionLabel{【往年题1｜12级流病口试考题回忆.pdf｜简答（题号8、33）】}\\
无差异错分偏倚的计算；控制信息偏倚措施

\end{pastquestion}

\begin{pastquestion}

\QuestionLabel{【往年题2｜14级流病考题回忆汇总.xls｜简答（题号1、54）】}\\
信息偏倚控制；信息偏倚的控制

\end{pastquestion}

\begin{pastquestion}

\QuestionLabel{【往年题3｜04级流病考题回忆.docx｜名解】}\\
information bias（信息偏倚）；Recall bias（回忆偏倚）

\end{pastquestion}

\PointSeparator

\section{其他偏倚名解（无应答偏倚、易感性偏倚、时间效应偏倚、排除）}\label{users__fpb__26_spring__pkusph_study__tmp__pdfs__prevention-medicine-past-exams__03-ux6d41ux884cux75c5ux5b66ux5f80ux5e74ux9898ux8003ux70b9ux6574ux7406.md__ux5176ux4ed6ux504fux501aux540dux89e3ux65e0ux5e94ux7b54ux504fux501aux6613ux611fux6027ux504fux501aux65f6ux95f4ux6548ux5e94ux504fux501aux6392ux9664}

\FrequencyPill{5 次}

\InfoLabel{对应小节} 第三章 现况研究 / 四、常见偏倚；第六章 第三节 一、资料的整理（一）排除（书中页码
第42、94页）

\InfoLabel{匹配依据} 04、10、12、19、20 级考``无应答偏倚（non-response bias）、易感性偏倚（susceptibility
bias）、时间效应偏倚（time effect bias）、排除（exclusion）''等名词解释。合并为一个考点。

\ContentBand{知识点原文摘取}

无应答偏倚（non-response
bias）：调查对象因不愿合作或其他原因没有应答，无应答者的患病状况与暴露情况与应答者可能不同，从而导致系统误差。若应答率低于70\%就难以用调查结果来估计整个研究总体的状况。

排除（exclusion）：指在随机分组前研究对象因各种原因没有被纳入。排除对研究结果的内部真实性不会产生影响，但可能影响研究结果的外推；被排除的研究对象愈多，结果推广的面愈小。

（``易感性偏倚''\,``时间效应偏倚''在课本 24
章中未检索到明确独立定义原文，多见于课件，暂不自行补充。按课件，易感性偏倚指不同特征人群对疾病的易感性不同所致的偏倚；时间效应偏倚指从暴露到发病的潜隐过程中，处于早期或临床前期者被错分所致的偏倚。）

\ContentBand{对应往年题}

\begin{pastquestion}

\QuestionLabel{【往年题1｜20级回忆试题库.xlsx｜名解（题号24）】}\\
non-response bias（无应答偏倚）

\end{pastquestion}

\begin{pastquestion}

\QuestionLabel{【往年题2｜19级流行病学期末口试回忆.xlsx｜名解（题号5）】}\\
无应答偏倚

\end{pastquestion}

\begin{pastquestion}

\QuestionLabel{【往年题3｜12级流病口试考题回忆.pdf｜名解（题号37、24）】}\\
易感性偏倚；时间效应偏倚

\end{pastquestion}

\PointSeparator

\chapter{病因及病因推断}\label{users__fpb__26_spring__pkusph_study__tmp__pdfs__prevention-medicine-past-exams__03-ux6d41ux884cux75c5ux5b66ux5f80ux5e74ux9898ux8003ux70b9ux6574ux7406.md__ux7b2cux4e5dux7ae0-ux75c5ux56e0ux53caux75c5ux56e0ux63a8ux65ad}

\section{病因推断的准则（希尔 Hill
九条）及综合分析（反应停、类固醇/雌激素与畸形等能否判断因果）}\label{users__fpb__26_spring__pkusph_study__tmp__pdfs__prevention-medicine-past-exams__03-ux6d41ux884cux75c5ux5b66ux5f80ux5e74ux9898ux8003ux70b9ux6574ux7406.md__ux75c5ux56e0ux63a8ux65adux7684ux51c6ux5219ux5e0cux5c14-hill-ux4e5dux6761ux53caux7efcux5408ux5206ux6790ux53cdux5e94ux505cux7c7bux56faux9187ux96ccux6fc0ux7d20ux4e0eux7578ux5f62ux7b49ux80fdux5426ux5224ux65adux56e0ux679c}

\FrequencyPill{10 次}

\InfoLabel{对应小节} 第八章 病因及其发现和推断 / 第五节 因果关系推论 /
三、综合所有证据的推论：希尔的病因推断准则（书中页码 第132---134页）

\InfoLabel{匹配依据} 04、07、10、12、14、17、19、20、21
级反复以简答和大题形式考查病因推断------``给一道病例对照/生态学/描述性资料（如反应停与海豹肢畸形、类固醇/雌激素与畸胎、A
药与先天性心脏病），能否据此判断因果关系？为什么？进一步需做什么研究？''要求按希尔准则逐条分析（病例对照不能确证因果，需队列/RCT
等）。

\ContentBand{知识点原文摘取}

希尔病因推断的9条标准：①时间顺序；②关联强度；③剂量-反应关系；④结果的一致性；⑤实验证据；⑥生物学合理性；⑦生物学一致性；⑧特异性；⑨相似性。（1991年
Susser 增加``预测力''一项。）

以上标准中，存在关联（包括剂量-反应关系）以及关联的时间特征是判断因果关系的必要条件和特异条件；其余7项是非特异、非必要的条件，其中结果的一致性最为重要。所有标准都不是充分条件，即使满足所有条件，也不能百分百肯定是因果关系。

时间顺序的可信度上，临床试验、队列研究、病例对照研究和横断面研究依次降低。病例对照研究由果及因，论证因果能力较弱，不能确证因果关系，需通过队列研究或随机对照试验进一步验证（队列研究是在人群中验证病因最可靠的方法）。

\ContentBand{对应往年题}

\begin{pastquestion}

\QuestionLabel{【往年题1｜12级流病口试考题回忆.pdf｜综合分析（题号33、58）】}\\
反应停与短肢畸形：属于哪类调查方法，怎么验证因果关系；一项孕产妇怀孕期间服用类固醇药物与新生儿先天性畸形的研究，问能否得出结论，如果不能应该怎么办（主要考病因及因果推断的
Hill 标准）

\end{pastquestion}

\begin{pastquestion}

\QuestionLabel{【往年题2｜21级流病考题回忆汇总.xlsx｜简答/综合分析（题号1、5、43）】}\\
简述希尔法则 /
病因确定的原则；2015年新上市A药用于治疗孕期抑郁症，某医院临床药师根据病例系列发现65\%先天性心脏病患儿母亲孕期吃过A药------能否判断A药是病因？为什么？对可疑病因应使用什么流行病学方法研究？为什么？

\end{pastquestion}

\begin{pastquestion}

\QuestionLabel{【往年题3｜20级回忆试题库.xlsx｜综合分析（题号1、28）】}\\
以一例病例对照研究分析，病因因果推断10条；围产儿先天畸形的母亲使用类固醇药物的比例高于对照------能否判断使用药物是导致死胎和畸形的原因？

\end{pastquestion}

\PointSeparator

\section{病因模型（三角模型、轮状模型、病因网）及其公共卫生学意义}\label{users__fpb__26_spring__pkusph_study__tmp__pdfs__prevention-medicine-past-exams__03-ux6d41ux884cux75c5ux5b66ux5f80ux5e74ux9898ux8003ux70b9ux6574ux7406.md__ux75c5ux56e0ux6a21ux578bux4e09ux89d2ux6a21ux578bux8f6eux72b6ux6a21ux578bux75c5ux56e0ux7f51ux53caux5176ux516cux5171ux536bux751fux5b66ux610fux4e49}

\FrequencyPill{5 次}

\InfoLabel{对应小节} 第八章 病因及其发现和推断 / 第二节 病因学说与病因模型 /
一、传染病病因的三角模型、二、病因的轮状模型、五、病因网（书中页码 第120---123页）

\InfoLabel{匹配依据} 10、12、14、21 级考``流行病学三角（三角模型）及其公共卫生学意义''等病因模型相关简答。

\ContentBand{知识点原文摘取}

传染病病因的三角模型（epidemiologic
triad）：影响传染病发生发展的因素归为宿主（host）、病原体（agent）和环境（environment）三方面，三者对传染病流行缺一不可，以等边三角形的平衡关系描述，保持动态平衡时发病率维持常态，一旦一个或多个因素变化破坏平衡，发病率就会变化。其公共卫生学意义：揭示了在病原体之外存在可以用来预防和控制传染病的因素，寻找切断三角中任何一条（或多条）链条的措施即可控制疾病流行。

病因网（web of
causation）：同一疾病不同病因链相互连接、相互交错形成的网状结构。意义：去除一条病因链中任何一个因素即可切断整个病因链；有效的预防应切断主要的病因链；同时切断多个病因链可预防更多病例。

\ContentBand{对应往年题}

\begin{pastquestion}

\QuestionLabel{【往年题1｜12级流病口试考题回忆.pdf｜简答（题号1、64）】}\\
三角模型公共卫生学意义；流行病学三角，公卫意义

\end{pastquestion}

\begin{pastquestion}

\QuestionLabel{【往年题2｜21级流病考题回忆汇总.xlsx｜简答（题号38）】}\\
流行病学三角及其公共卫生学意义

\end{pastquestion}

\begin{pastquestion}

\QuestionLabel{【往年题3｜14级流病考题回忆汇总.xls｜简答（题号52）】}\\
三角模型及其公共卫生学意义

\end{pastquestion}

\PointSeparator

\section{病因（cause）与必要病因（necessary
cause）}\label{users__fpb__26_spring__pkusph_study__tmp__pdfs__prevention-medicine-past-exams__03-ux6d41ux884cux75c5ux5b66ux5f80ux5e74ux9898ux8003ux70b9ux6574ux7406.md__ux75c5ux56e0causeux4e0eux5fc5ux8981ux75c5ux56e0necessary-cause}

\FrequencyPill{6 次}

\InfoLabel{对应小节} 第八章 病因及其发现和推断 / 第一节 概述 / 一、病因与因果关系；第三节
充分病因-组分病因模型 / 二、充分病因和组分病因（书中页码 第117---118、125---126页）

\InfoLabel{匹配依据} 04、05、07、12、19、21 级考``病因（cause of disease）/必要病因（necessary
cause）''名词解释，05 级追问公卫与临床上病因的区别（与书上定义不一样）。

\ContentBand{知识点原文摘取}

病因（cause of
disease）是引起疾病发生的原因，即能够影响未来疾病发生的因素或事件。（Rothman：病因就是那些对疾病发生起着核心作用的事件、特征和条件。）

必要病因（necessary
cause）：是一个疾病发生必需的组分病因，是该疾病所有充分病因都需要的组分病因；若该病因不存在，疾病就不会发生，因此所有病人都具有该病因。最典型的必要病因就是传染病的病原体。

\ContentBand{对应往年题}

\begin{pastquestion}

\QuestionLabel{【往年题1｜21级流病考题回忆汇总.xlsx｜名解（题号23、45）】}\\
必要病因（necessary cause / necessary causation）

\end{pastquestion}

\begin{pastquestion}

\QuestionLabel{【往年题2｜05流病考题回忆.doc｜名解】}\\
病因（李老师给详细解释了一下，和书上的定义不一样）

\end{pastquestion}

\begin{pastquestion}

\QuestionLabel{【往年题3｜19级流行病学期末口试回忆.xlsx｜名解（题号5）】}\\
病因（追问公卫与临床上病因的区别）

\end{pastquestion}

\PointSeparator

\section{穆勒法则（求同法等）与因果关系的三个条件}\label{users__fpb__26_spring__pkusph_study__tmp__pdfs__prevention-medicine-past-exams__03-ux6d41ux884cux75c5ux5b66ux5f80ux5e74ux9898ux8003ux70b9ux6574ux7406.md__ux7a46ux52d2ux6cd5ux5219ux6c42ux540cux6cd5ux7b49ux4e0eux56e0ux679cux5173ux7cfbux7684ux4e09ux4e2aux6761ux4ef6}

\FrequencyPill{2 次}

\InfoLabel{对应小节} 第八章 病因及其发现和推断 / 第四节 发现和验证因果关系 /
一、提出因果假设、二、因果推断的逻辑方法（书中页码 第127---129页）

\InfoLabel{匹配依据} 04
级名解考``求同法''；病因推断综合题中常需说明``依据穆勒法则/因果关系三条件（时间顺序、关联关系、因变性）''。

\ContentBand{知识点原文摘取}

因果关系必须同时满足三个基本条件：①时间顺序；②关联关系；③因变性。

穆勒法则（Mill's canons）包括求同法、求异法、同异共求法、共变法和剩余法。求同法（method of
agreement）：如果所有患同一疾病的病人都具有某一共同的因素，而其他因素并非每个病人都有，那么该因素可能是该疾病的病因。共求法奠定了病例对照研究、队列研究和随机对照试验的理论基础；共变法强调对剂量-反应关系的考量。

\ContentBand{对应往年题}

\begin{pastquestion}

\QuestionLabel{【往年题1｜04级流病考题回忆.docx｜名解】}\\
求同法

\end{pastquestion}

\PointSeparator

\chapter{疾病预防与监测}\label{users__fpb__26_spring__pkusph_study__tmp__pdfs__prevention-medicine-past-exams__03-ux6d41ux884cux75c5ux5b66ux5f80ux5e74ux9898ux8003ux70b9ux6574ux7406.md__ux7b2cux5341ux7ae0-ux75beux75c5ux9884ux9632ux4e0eux76d1ux6d4b}

\section{疾病的三级预防（一级、二级、三级预防）}\label{users__fpb__26_spring__pkusph_study__tmp__pdfs__prevention-medicine-past-exams__03-ux6d41ux884cux75c5ux5b66ux5f80ux5e74ux9898ux8003ux70b9ux6574ux7406.md__ux75beux75c5ux7684ux4e09ux7ea7ux9884ux9632ux4e00ux7ea7ux4e8cux7ea7ux4e09ux7ea7ux9884ux9632}

\FrequencyPill{9 次}

\InfoLabel{对应小节} 第九章 预防策略 / 第二节 预防策略与措施 / 二、疾病预防（二）疾病的三级预防（书中页码
第140---142页）

\InfoLabel{匹配依据} 04、07、10、12、14、17、19、20、21 级反复考``一级预防（primary
prevention）/二级预防（secondary prevention）/三级预防''名词解释及``三级预防及其意义、举例说明''。

\ContentBand{知识点原文摘取}

第一级预防（primary
prevention），又称病因预防，是指在疾病和残疾尚未发生时针对病因或危险因素采取措施，降低有害暴露的水平，增强个体对抗有害暴露的能力，预防疾病的发生，或至少推迟疾病的发生。措施如免疫接种、戒烟限酒、合理膳食、增加身体活动等。

第二级预防（secondary
prevention），又称``三早''预防，即早发现、早诊断、早治疗。在疾病早期，通过及早发现并诊断疾病、及时治疗，阻止疾病发展。早发现可通过筛检、定期体检等实现。

第三级预防（tertiary
prevention），又称临床预防或疾病管理。在疾病的症状体征明显表现出来之后，通过适当治疗缓解症状、预防恶化、预防合并症和残疾，对已发生的残疾进行康复治疗，提高生活质量、延长寿命。

\ContentBand{对应往年题}

\begin{pastquestion}

\QuestionLabel{【往年题1｜17级流病口试回忆.xlsx｜简答（题号5）】}\\
三级预防，举例说明

\end{pastquestion}

\begin{pastquestion}

\QuestionLabel{【往年题2｜20级回忆试题库.xlsx｜名解（题号13、48）】}\\
primary prevention（第一级预防）；secondary prevention（第二级预防）

\end{pastquestion}

\begin{pastquestion}

\QuestionLabel{【往年题3｜12级流病口试考题回忆.pdf｜简答（题号45、5）】}\\
三级预防及其意义；二级预防

\end{pastquestion}

\PointSeparator

\section{全人群策略与高危人群策略及其优劣}\label{users__fpb__26_spring__pkusph_study__tmp__pdfs__prevention-medicine-past-exams__03-ux6d41ux884cux75c5ux5b66ux5f80ux5e74ux9898ux8003ux70b9ux6574ux7406.md__ux5168ux4ebaux7fa4ux7b56ux7565ux4e0eux9ad8ux5371ux4ebaux7fa4ux7b56ux7565ux53caux5176ux4f18ux52a3}

\FrequencyPill{6 次}

\InfoLabel{对应小节} 第九章 预防策略 / 第二节 二、疾病预防（二）1.第一级预防
/（1）高危人群策略、（2）全人群策略（书中页码 第140---141页）

\InfoLabel{匹配依据} 10、11、12、14、20、21 级考``全人群策略（population-based
strategy）/高危人群策略（high-risk strategy）''名解及``全人群策略与高危策略的优劣''简答。

\ContentBand{知识点原文摘取}

高危人群策略（high-risk
strategy）：以临床医学思维为导向，对未来发病风险高的一小部分个体，针对致病危险因素采取有针对性的措施。优点是对资源的利用更符合成本效益原则；局限是要求少数人在行为上与众不同，效果受限，且当病因的原因波及整个人群时是治标不治本的策略。

全人群策略（population-based
strategy）：以公共卫生思维为导向，不需确定个体风险高低，而是通过降低整个人群有害暴露的水平（针对病因的原因），降低人群总的疾病负担。根据
Rose
的风险（预防）悖论，大部分病例出自低/中等暴露水平人群，全人群策略使整个人群风险分布曲线向低风险方向平移，总健康收益可观，但平均每个个体所得甚少。

两种策略各有优势和不足，并非非此即彼，而是互为补充、协同作用。

\ContentBand{对应往年题}

\begin{pastquestion}

\QuestionLabel{【往年题1｜11级流行病学口试考题回忆.pdf｜简答】}\\
全人群策略的优劣；人群策略

\end{pastquestion}

\begin{pastquestion}

\QuestionLabel{【往年题2｜12级流病口试考题回忆.pdf｜简答（题号21、37）】}\\
高危策略及其优劣；全人群策略的优劣

\end{pastquestion}

\begin{pastquestion}

\QuestionLabel{【往年题3｜20级回忆试题库.xlsx｜名解（题号50）】}\\
population-based strategy（全人群策略）

\end{pastquestion}

\PointSeparator

\section{主动监测与被动监测}\label{users__fpb__26_spring__pkusph_study__tmp__pdfs__prevention-medicine-past-exams__03-ux6d41ux884cux75c5ux5b66ux5f80ux5e74ux9898ux8003ux70b9ux6574ux7406.md__ux4e3bux52a8ux76d1ux6d4bux4e0eux88abux52a8ux76d1ux6d4b}

\FrequencyPill{8 次}

\InfoLabel{对应小节} 第十章 公共卫生监测 / 第一节 一、基本概念（二）监测相关的基本概念（书中页码
第146---147页）

\InfoLabel{匹配依据} 04、05、10、11、12、14、17、19、20、21 级广泛考``主动监测（active
surveillance）/被动监测（passive surveillance）''名词解释。

\ContentBand{知识点原文摘取}

被动监测（passive
surveillance）是指下级单位常规地向上级机构报告监测资料，而上级单位被动地接受。主要依据相关法律法规要求进行，我国的法定传染病报告信息系统、突发公共卫生事件报告系统、药品不良反应监测自发报告系统等多属于被动监测。

主动监测（active
surveillance）是指根据公共卫生问题的特殊需要，由上级单位专门组织调查或统一部署监测软件工具主动识别收集资料。我国的传染病智能监测预警前置软件、免疫接种率监测、传染病漏报调查以及对某些重点疾病（如不明原因发热）的监测等，多属于主动监测。

\ContentBand{对应往年题}

\begin{pastquestion}

\QuestionLabel{【往年题1｜04级流病考题回忆.docx｜名解】}\\
主动监测

\end{pastquestion}

\begin{pastquestion}

\QuestionLabel{【往年题2｜20级回忆试题库.xlsx｜名解（题号35）】}\\
passive surveillance（被动监测）

\end{pastquestion}

\begin{pastquestion}

\QuestionLabel{【往年题3｜14级流病考题回忆汇总.xls｜名解】}\\
被动监测；主动监测

\end{pastquestion}

\PointSeparator

\section{公共卫生监测的概念、目的及种类与内容}\label{users__fpb__26_spring__pkusph_study__tmp__pdfs__prevention-medicine-past-exams__03-ux6d41ux884cux75c5ux5b66ux5f80ux5e74ux9898ux8003ux70b9ux6574ux7406.md__ux516cux5171ux536bux751fux76d1ux6d4bux7684ux6982ux5ff5ux76eeux7684ux53caux79cdux7c7bux4e0eux5185ux5bb9}

\FrequencyPill{6 次}

\InfoLabel{对应小节} 第十章 公共卫生监测 / 第一节 一、基本概念、二、目的和意义；第二节 种类与内容（书中页码
第146---150页）

\InfoLabel{匹配依据} 10、11、12、17、19、21 级考``公共卫生监测（public health
surveillance）''名解及``疾病监测/公共卫生监测的种类和内容并举例说明''。

\ContentBand{知识点原文摘取}

公共卫生监测（public health
surveillance）是指长期、连续、系统地收集人群中有关公共卫生问题的资料，经过科学分析和解释后，获得重要的公共卫生信息，并及时反馈给需要这些信息的个人或机构，用以指导制定、完善和评价公共卫生干预策略与措施的过程。

种类主要包括：疾病监测（传染病监测、慢性非传染性疾病监测、医院感染监测、死因监测）、症状监测、行为危险因素监测，以及环境、食品与营养、药物不良反应等其他公共卫生监测。

\ContentBand{对应往年题}

\begin{pastquestion}

\QuestionLabel{【往年题1｜10级考题回忆.pdf｜简答】}\\
公共卫生监测的种类和内容并举例说明

\end{pastquestion}

\begin{pastquestion}

\QuestionLabel{【往年题2｜17级流病口试回忆.xlsx｜简答（题号11）】}\\
介绍疾病监测的内容和种类（追问传染病和 NCD、死因监测的区别，监测的目的）

\end{pastquestion}

\begin{pastquestion}

\QuestionLabel{【往年题3｜21级流病考题回忆汇总.xlsx｜名解（题号9、20）】}\\
公共卫生监测 / public health surveillance

\end{pastquestion}

\PointSeparator

\section{哨点监测}\label{users__fpb__26_spring__pkusph_study__tmp__pdfs__prevention-medicine-past-exams__03-ux6d41ux884cux75c5ux5b66ux5f80ux5e74ux9898ux8003ux70b9ux6574ux7406.md__ux54e8ux70b9ux76d1ux6d4b}

\FrequencyPill{5 次}

\InfoLabel{对应小节} 第十章 公共卫生监测 / 第一节 一、基本概念（二）4.哨点监测（书中页码 第147页）

\InfoLabel{匹配依据} 10、11、12、14、21 级考``哨点监测（sentinel
surveillance）''名词解释，部分追问用于哪类疾病。

\ContentBand{知识点原文摘取}

哨点监测（sentinel
surveillance）是为了更清楚地了解某些疾病在不同地区、不同人群的分布以及相应的影响因素等，根据被监测疾病的流行特点，选择若干有代表性的地区和/或人群，按统一的监测方案连续地开展监测。最典型的是艾滋病哨点监测，选择有代表性的地区和艾滋病相关高危人群，开展
HIV 抗体检测，同时收集监测人群的高危行为信息。

\ContentBand{对应往年题}

\begin{pastquestion}

\QuestionLabel{【往年题1｜10级考题回忆.pdf｜名解】}\\
哨点监测（sentinel surveillance）

\end{pastquestion}

\begin{pastquestion}

\QuestionLabel{【往年题2｜11级流行病学口试考题回忆.pdf｜名解】}\\
哨点监测（英文）

\end{pastquestion}

\begin{pastquestion}

\QuestionLabel{【往年题3｜21级流病考题回忆汇总.xlsx｜名解（题号24）】}\\
哨点监测

\end{pastquestion}

\PointSeparator

\section{健康促进与健康保护}\label{users__fpb__26_spring__pkusph_study__tmp__pdfs__prevention-medicine-past-exams__03-ux6d41ux884cux75c5ux5b66ux5f80ux5e74ux9898ux8003ux70b9ux6574ux7406.md__ux5065ux5eb7ux4fc3ux8fdbux4e0eux5065ux5eb7ux4fddux62a4}

\FrequencyPill{5 次}

\InfoLabel{对应小节} 第九章 预防策略 / 第二节 三、健康保护与健康促进（书中页码 第142---143页）

\InfoLabel{匹配依据} 10、11、12、14、21 级考``健康促进（health promotion）/健康保护（health
protection）''名词解释及健康促进的策略。

\ContentBand{知识点原文摘取}

健康保护（health
protection）又称健康防护，即采取有针对性的措施保护个体或人群免受来自外界环境的有害物质（生物、物理、化学类有害物质）对健康的威胁。

健康促进（health
promotion）是增强人们控制健康影响因素、改善自身健康的能力的过程。它是一个综合的社会和政治活动过程，不仅包括直接改善个体行为和加强生活技能的健康教育，还包括通过政策、立法、经济手段和环境工程改善社会、经济和环境条件的社会行动。Tannahill
模型将健康促进置于健康保护和健康教育的策略之上，涵盖预防、健康教育、健康保护三种策略。

\ContentBand{对应往年题}

\begin{pastquestion}

\QuestionLabel{【往年题1｜11级流行病学口试考题回忆.pdf｜名解】}\\
健康促进

\end{pastquestion}

\begin{pastquestion}

\QuestionLabel{【往年题2｜21级流病考题回忆汇总.xlsx｜名解（题号18、43、46）】}\\
health protection（健康保护）；健康促进

\end{pastquestion}

\begin{pastquestion}

\QuestionLabel{【往年题3｜14级流病考题回忆汇总.xls｜简答（题号35）】}\\
健康促进策略

\end{pastquestion}

\PointSeparator

\section{症状监测}\label{users__fpb__26_spring__pkusph_study__tmp__pdfs__prevention-medicine-past-exams__03-ux6d41ux884cux75c5ux5b66ux5f80ux5e74ux9898ux8003ux70b9ux6574ux7406.md__ux75c7ux72b6ux76d1ux6d4b}

\FrequencyPill{4 次}

\InfoLabel{对应小节} 第十章 公共卫生监测 / 第二节 种类与内容 / 二、症状监测（书中页码 第149页）

\InfoLabel{匹配依据} 10、12、19、21 级考``症状监测（syndromic surveillance）''名词解释。

\ContentBand{知识点原文摘取}

症状监测（syndromic
surveillance）又称为综合征监测或症候群监测，是指通过长期、连续、系统地收集特定临床症候群或与疾病相关现象的发生频率，从而对某类疾病的发生或流行进行早期探查、预警和作出快速反应的监测方法。尤其适用于新发疾病（病因未明、临床上尚无明确诊断方法）。常用主题有流感样病例监测、发热监测、腹泻监测等。由于症状的出现总是先于疾病的确诊，可提高监测系统的敏感性。

\ContentBand{对应往年题}

\begin{pastquestion}

\QuestionLabel{【往年题1｜12级流病口试考题回忆.pdf｜名解（题号59）】}\\
症状监测

\end{pastquestion}

\begin{pastquestion}

\QuestionLabel{【往年题2｜19级流行病学期末口试回忆.xlsx｜名解（题号20）】}\\
症状监测

\end{pastquestion}

\begin{pastquestion}

\QuestionLabel{【往年题3｜21级流病考题回忆汇总.xlsx｜名解（题号20）】}\\
症状监测

\end{pastquestion}

\PointSeparator

\chapter{传染病流行病学}\label{users__fpb__26_spring__pkusph_study__tmp__pdfs__prevention-medicine-past-exams__03-ux6d41ux884cux75c5ux5b66ux5f80ux5e74ux9898ux8003ux70b9ux6574ux7406.md__ux7b2cux5341ux4e00ux7ae0-ux4f20ux67d3ux75c5ux6d41ux884cux75c5ux5b66}

\section{传播途径（各种途径及举例；经空气/水/食物/节肢动物的流行特征）}\label{users__fpb__26_spring__pkusph_study__tmp__pdfs__prevention-medicine-past-exams__03-ux6d41ux884cux75c5ux5b66ux5f80ux5e74ux9898ux8003ux70b9ux6574ux7406.md__ux4f20ux64adux9014ux5f84ux5404ux79cdux9014ux5f84ux53caux4e3eux4f8bux7ecfux7a7aux6c14ux6c34ux98dfux7269ux8282ux80a2ux52a8ux7269ux7684ux6d41ux884cux7279ux5f81}

\FrequencyPill{9 次}

\InfoLabel{对应小节} 第十一章 传染病流行病学 / 第三节 流行过程 / 一、基本环节（二）传播途径（书中页码
第160---162页）

\InfoLabel{匹配依据} 04、05、07、10、11、12、14、19、20、21 级广泛考``传播途径（route of
transmission）有哪些、举例说明''及各种传播途径名解（直接接触传播、间接接触传播、医源性传播、机械携带、生物媒介、水平/垂直传播），并考``经空气/经水/经食物/经节肢动物传播的传染病流行特征''。

\ContentBand{知识点原文摘取}

传播途径（route of
transmission）是指病原体从传染源排出至侵入新的易感宿主前，在外环境中所经历的全过程。传播方式分水平传播和垂直传播（母婴传播）。

主要传播途径：①经空气传播（飞沫、飞沫核、尘埃）------流行特征：传播易实现、发病率高、有冬春季季节性、无免疫预防人群发病呈周期性、人口密度大地区高发；②经水传播（饮用水、疫水接触）；③经食物传播；④经接触传播（直接接触、间接接触/日常生活接触）；⑤经节肢动物传播（机械携带
mechanical vector、生物学传播 biological
vector）------流行特征：有地区性、季节性、职业及年龄差异；⑥经土壤传播；⑦医源性传播（iatrogenic）；⑧母婴传播（经胎盘、上行性、分娩时）。

\ContentBand{对应往年题}

\begin{pastquestion}

\QuestionLabel{【往年题1｜12级流病口试考题回忆.pdf｜简答（题号3）】}\\
节肢动物传播的流行病学特点

\end{pastquestion}

\begin{pastquestion}

\QuestionLabel{【往年题2｜21级流病考题回忆汇总.xlsx｜名解（题号25、43、11）】}\\
医源性传播；生物媒介；生物学媒介

\end{pastquestion}

\begin{pastquestion}

\QuestionLabel{【往年题3｜04级流病考题回忆.docx｜问答】}\\
传播途径有哪些，举例；空气传染病的流行病学特点

\end{pastquestion}

\PointSeparator

\section{人群易感性升高与降低的原因}\label{users__fpb__26_spring__pkusph_study__tmp__pdfs__prevention-medicine-past-exams__03-ux6d41ux884cux75c5ux5b66ux5f80ux5e74ux9898ux8003ux70b9ux6574ux7406.md__ux4ebaux7fa4ux6613ux611fux6027ux5347ux9ad8ux4e0eux964dux4f4eux7684ux539fux56e0}

\FrequencyPill{8 次}

\InfoLabel{对应小节} 第十一章 传染病流行病学 / 第三节 流行过程 / 一、基本环节（三）易感人群（书中页码
第162---163页）

\InfoLabel{匹配依据} 10、11、12、14、17、19、20、21
级反复考``使人群易感性升高/降低的因素（分升高和降低作答）''。

\ContentBand{知识点原文摘取}

引起人群易感性升高的主要因素：①新生儿增加（出生6个月以上婴儿源自母体的抗体逐渐消失）；②易感人口迁入；③免疫人口减少（免疫水平随时间消退或免疫人口死亡）；④新型病原体出现或病原体变异。

导致人群易感性降低的主要因素：①预防接种（降低人群易感性最主要的因素）；②传染病流行（流行后相当数量易感者因患病或隐性感染获得免疫力）。

\ContentBand{对应往年题}

\begin{pastquestion}

\QuestionLabel{【往年题1｜19级流行病学期末口试回忆.xlsx｜简答（题号9）】}\\
使人群易感性升高/降低的因素（分升高和降低作答）

\end{pastquestion}

\begin{pastquestion}

\QuestionLabel{【往年题2｜20级回忆试题库.xlsx｜简答（题号2）】}\\
导致人群易感性降低的原因

\end{pastquestion}

\begin{pastquestion}

\QuestionLabel{【往年题3｜10级考题回忆.pdf｜简答】}\\
人群易感性增加的原因；人群易感性降低的原因

\end{pastquestion}

\PointSeparator

\section{传染过程、感染谱、隐性感染与冰山现象}\label{users__fpb__26_spring__pkusph_study__tmp__pdfs__prevention-medicine-past-exams__03-ux6d41ux884cux75c5ux5b66ux5f80ux5e74ux9898ux8003ux70b9ux6574ux7406.md__ux4f20ux67d3ux8fc7ux7a0bux611fux67d3ux8c31ux9690ux6027ux611fux67d3ux4e0eux51b0ux5c71ux73b0ux8c61}

\FrequencyPill{6 次}

\InfoLabel{对应小节} 第十一章 传染病流行病学 / 第二节 传染过程 / 三、感染谱（书中页码 第158页）

\InfoLabel{匹配依据} 04、05、10、11、12、20 级考``传染过程（infection process）/隐性感染（latent/covert
infection）/传染病的冰山现象（iceberg phenomenon）''名词解释。

\ContentBand{知识点原文摘取}

传染过程（infection process）指病原体进入宿主机体后，与机体相互作用、相互斗争的过程，是个体现象。

感染谱（spectrum of
infection）又称感染梯度，是指机体感染病原体后所呈现出的各种临床表现的集合。出现临床症状和体征的感染称显性感染（临床感染）；不引起症状或体征、但可通过实验室方法发现的感染称隐性感染（亚临床感染、无症状感染）。

``冰山现象''：以隐性感染为主的传染病，临床上能发现的显性病例只是其中一小部分，犹如浮出水面的冰山一角，大量隐性感染者隐藏在水面之下。隔离病人对以隐性感染为主的传染病作用甚微，而对以显性感染为主的传染病则较为有效。

\ContentBand{对应往年题}

\begin{pastquestion}

\QuestionLabel{【往年题1｜05流病考题回忆.doc｜名解】}\\
冰山现象

\end{pastquestion}

\begin{pastquestion}

\QuestionLabel{【往年题2｜20级回忆试题库.xlsx｜名解（题号21、11）】}\\
iceberg phenomenon（冰山现象，见传染病）；latent infection（潜伏期感染/潜隐期）

\end{pastquestion}

\begin{pastquestion}

\QuestionLabel{【往年题3｜12级流病口试考题回忆.pdf｜名解（题号63、57）】}\\
隐性感染

\end{pastquestion}

\PointSeparator

\section{消毒（随时消毒、终末消毒）}\label{users__fpb__26_spring__pkusph_study__tmp__pdfs__prevention-medicine-past-exams__03-ux6d41ux884cux75c5ux5b66ux5f80ux5e74ux9898ux8003ux70b9ux6574ux7406.md__ux6d88ux6bd2ux968fux65f6ux6d88ux6bd2ux7ec8ux672bux6d88ux6bd2}

\FrequencyPill{6 次}

\InfoLabel{对应小节} 第十一章 传染病流行病学 / 第四节 预防策略与措施（针对传播途径的措施）（书中页码 第166页）

\InfoLabel{匹配依据} 10、11、12、14、20、21 级考``随时消毒（current disinfection）/终末消毒（terminal
disinfection）''名词解释。

\ContentBand{知识点原文摘取}

（课本第十一章正文将``用化学、物理、生物等方法杀灭或消除外界环境中的致病性微生物''作为切断传播途径的主要措施。课本知识点文件中未检索到``随时消毒/终末消毒''的独立定义原文，按课件：随时消毒指对现有传染源存在的疫源地进行的、随时杀灭其排出病原体的消毒；终末消毒指传染源因住院、死亡或痊愈而离开疫源地后所进行的一次彻底消毒，以完全清除传染源所遗留的病原体。建议结合课件作答。）

\ContentBand{对应往年题}

\begin{pastquestion}

\QuestionLabel{【往年题1｜11级流行病学口试考题回忆.pdf｜名解】}\\
随时消毒；终末消毒

\end{pastquestion}

\begin{pastquestion}

\QuestionLabel{【往年题2｜20级回忆试题库.xlsx｜名解（题号5）】}\\
terminal disinfection（终末消毒）

\end{pastquestion}

\begin{pastquestion}

\QuestionLabel{【往年题3｜21级流病考题回忆汇总.xlsx｜名解（题号31、47）】}\\
随时消毒（current disinfection）

\end{pastquestion}

\PointSeparator

\section{主动免疫、被动免疫与免疫规划（计划免疫）}\label{users__fpb__26_spring__pkusph_study__tmp__pdfs__prevention-medicine-past-exams__03-ux6d41ux884cux75c5ux5b66ux5f80ux5e74ux9898ux8003ux70b9ux6574ux7406.md__ux4e3bux52a8ux514dux75abux88abux52a8ux514dux75abux4e0eux514dux75abux89c4ux5212ux8ba1ux5212ux514dux75ab}

\FrequencyPill{6 次}

\InfoLabel{对应小节} 第十一章 传染病流行病学 / 第五节 免疫规划及其人群评价 /
一、预防接种、二、免疫规划（书中页码 第168---169页）

\InfoLabel{匹配依据} 04、10、11、12、14、20 级考``计划免疫/免疫规划、主动免疫（active
immunization）、人工被动免疫（passive immunization）''名词解释。

\ContentBand{知识点原文摘取}

预防接种（immunization/vaccination）又称人工免疫或免疫接种，包括主动免疫和被动免疫。

主动免疫（active immunization）指将疫苗接种到机体，使之产生特异性免疫，是最常见的预防接种形式。

被动免疫（passive
immunization）是人体接种含特异性抗体的血清或制剂，使机体被动地获得特异性免疫力而受到保护。接种后即产生免疫力，见效快，但维持时间较短，不会产生免疫记忆，常用于疫情发生时的紧急预防或治疗。

免疫规划（1981年起逐渐取代``计划免疫''一词）：根据国家传染病防制规划，使用有效疫苗对易感人群进行预防接种所制定的规划、计划和策略，按确定的疫苗品种、免疫程序在人群中有计划地进行预防接种。

\ContentBand{对应往年题}

\begin{pastquestion}

\QuestionLabel{【往年题1｜04级流病考题回忆.docx｜名解】}\\
计划免疫

\end{pastquestion}

\begin{pastquestion}

\QuestionLabel{【往年题2｜20级回忆试题库.xlsx｜名解（题号12、45）】}\\
active immunization（主动免疫）；artificial passive immunization（人工被动免疫）

\end{pastquestion}

\begin{pastquestion}

\QuestionLabel{【往年题3｜12级流病口试考题回忆.pdf｜名解（题号10、65）】}\\
主动免疫；人工被动免疫

\end{pastquestion}

\PointSeparator

\section{潜伏期及其流行病学意义}\label{users__fpb__26_spring__pkusph_study__tmp__pdfs__prevention-medicine-past-exams__03-ux6d41ux884cux75c5ux5b66ux5f80ux5e74ux9898ux8003ux70b9ux6574ux7406.md__ux6f5cux4f0fux671fux53caux5176ux6d41ux884cux75c5ux5b66ux610fux4e49}

\FrequencyPill{5 次}

\InfoLabel{对应小节} 第十一章 传染病流行病学 / 第三节 流行过程 /
一、基本环节（一）传染源---1.病人（1）潜伏期（书中页码 第158---159页）

\InfoLabel{匹配依据} 04、05、17、19、21 级考``潜伏期（incubation period）''名解及其流行病学意义，17
级追问``一定要答出追溯传染源和传播途径''。

\ContentBand{知识点原文摘取}

潜伏期（incubation period）是指从病原体侵入机体到最早出现临床症状或体征的一段时间。

潜伏期的流行病学意义及用途：①判断病人被感染的时间，追溯传染源，确定传播途径；②确定接触者的隔离和医学观察期限，一般为平均潜伏期加1～2天，危害严重的传染病可按最长潜伏期处理；③确定预防接种时间；④评价防制措施的效果（发病数经过一个潜伏期明显下降可认为措施可能有效）；⑤潜伏期的长短会影响疾病的流行特征，潜伏期短的常以暴发形式出现，潜伏期长的多表现为散发。

\ContentBand{对应往年题}

\begin{pastquestion}

\QuestionLabel{【往年题1｜17级流病口试回忆.xlsx｜名解（题号10、9）】}\\
incubation period / 潜伏期（追问意义：一定要答出追溯传染源和传播途径）

\end{pastquestion}

\begin{pastquestion}

\QuestionLabel{【往年题2｜04级流病考题回忆.docx｜名解】}\\
潜伏期（含意义）

\end{pastquestion}

\begin{pastquestion}

\QuestionLabel{【往年题3｜21级流病考题回忆汇总.xlsx｜名解（题号29、50）】}\\
潜伏期 incubation period（追问潜隐期的概念）

\end{pastquestion}

\PointSeparator

\section{疫源地及其消灭条件}\label{users__fpb__26_spring__pkusph_study__tmp__pdfs__prevention-medicine-past-exams__03-ux6d41ux884cux75c5ux5b66ux5f80ux5e74ux9898ux8003ux70b9ux6574ux7406.md__ux75abux6e90ux5730ux53caux5176ux6d88ux706dux6761ux4ef6}

\FrequencyPill{5 次}

\InfoLabel{对应小节} 第十一章 传染病流行病学 / 第三节 流行过程 / 三、疫源地与流行过程（书中页码
第163---164页）

\InfoLabel{匹配依据} 04、07、14、19、20 级考``疫源地（epidemic focus）''名解及``疫源地的消灭（消除）条件''。

\ContentBand{知识点原文摘取}

疫源地（epidemic
focus）是指传染源及其排出的病原体向周围播散所能波及的范围，即新病例或新感染可能发生的范围。范围较小的疫源地称疫点，范围较大的称疫区。

疫源地的消灭必须同时具备下列条件：①传染源被移除（住院、死亡或移至他处）或不再排出病原体（治愈）；②通过各种措施消灭了传染源排到外环境的病原体或消除了传播途径（如彻底灭蚊）；③传染源周围的所有易感接触者经过该病最长潜伏期没有出现新病例或新感染者。

\ContentBand{对应往年题}

\begin{pastquestion}

\QuestionLabel{【往年题1｜04级流病考题回忆.docx｜简答】}\\
疫源地及消除条件

\end{pastquestion}

\begin{pastquestion}

\QuestionLabel{【往年题2｜20级回忆试题库.xlsx｜名解（题号16）】}\\
epidemic focus（疫源地）

\end{pastquestion}

\begin{pastquestion}

\QuestionLabel{【往年题3｜19级流行病学期末口试回忆.xlsx｜名解（题号21）】}\\
疫源地

\end{pastquestion}

\PointSeparator

\section{传染源与病原携带者}\label{users__fpb__26_spring__pkusph_study__tmp__pdfs__prevention-medicine-past-exams__03-ux6d41ux884cux75c5ux5b66ux5f80ux5e74ux9898ux8003ux70b9ux6574ux7406.md__ux4f20ux67d3ux6e90ux4e0eux75c5ux539fux643aux5e26ux8005}

\FrequencyPill{5 次}

\InfoLabel{对应小节} 第十一章 传染病流行病学 / 第三节 流行过程 / 一、基本环节（一）传染源（书中页码
第158---160页）

\InfoLabel{匹配依据} 04、10、12、14、20 级考``传染源（source of infection）/病原携带者（carrier）''名词解释。

\ContentBand{知识点原文摘取}

传染源（source of
infection）是指体内有病原体生长、繁殖，并能排出病原体的人和动物，包括传染病病人、病原携带者和受感染的动物。

病原携带者（carrier）是指感染病原体后无临床症状但能排出病原体的人，有时又称无症状感染者，包括隐性感染者和处于潜伏期或恢复期的显性感染者，可分为潜伏期病原携带者、恢复期病原携带者和健康病原携带者。在饮食服务、供水、托幼等机构工作的病原携带者对人群健康威胁非常大。

\ContentBand{对应往年题}

\begin{pastquestion}

\QuestionLabel{【往年题1｜10级考题回忆.pdf｜名解】}\\
source of infection（传染源）

\end{pastquestion}

\begin{pastquestion}

\QuestionLabel{【往年题2｜20级回忆试题库.xlsx｜名解（题号10）】}\\
carrier（病原携带者）

\end{pastquestion}

\begin{pastquestion}

\QuestionLabel{【往年题3｜04级流病考题回忆.docx｜名解】}\\
病原携带者

\end{pastquestion}

\PointSeparator

\section{传染病的预防控制措施（策略与措施）}\label{users__fpb__26_spring__pkusph_study__tmp__pdfs__prevention-medicine-past-exams__03-ux6d41ux884cux75c5ux5b66ux5f80ux5e74ux9898ux8003ux70b9ux6574ux7406.md__ux4f20ux67d3ux75c5ux7684ux9884ux9632ux63a7ux5236ux63aaux65bdux7b56ux7565ux4e0eux63aaux65bd}

\FrequencyPill{5 次}

\InfoLabel{对应小节} 第十一章 传染病流行病学 / 第四节 预防策略与措施 / 二、措施（书中页码 第165---167页）

\InfoLabel{匹配依据} 10、11、12、14、20 级考``传染病的预防控制措施/策略''。

\ContentBand{知识点原文摘取}

预防控制策略可概括为``预防为主、加强监测、群防群控、联防联控''。

措施包括：①经常性的传染病防制措施（传染病报告、改善卫生条件、健康教育、国境卫生检疫、预防接种）；②针对传染源的措施（对病人做到早发现、早诊断、早报告、早隔离、早治疗；对病原携带者隔离治疗、限制职业行为；对接触者医学观察、应急接种、药物预防；对动物传染源无害化处理等）；③针对传播途径的措施（消毒、杀虫等切断传播途径）；④针对易感人群的措施（预防接种、药物预防、个人防护）。

\ContentBand{对应往年题}

\begin{pastquestion}

\QuestionLabel{【往年题1｜12级流病口试考题回忆.pdf｜简答（题号8）】}\\
谈一下传染病的预防控制措施

\end{pastquestion}

\begin{pastquestion}

\QuestionLabel{【往年题2｜20级回忆试题库.xlsx｜简答（题号28）】}\\
传染病的预防控制措施 / 传染病的防控策略

\end{pastquestion}

\begin{pastquestion}

\QuestionLabel{【往年题3｜11级流行病学口试考题回忆.pdf｜简答】}\\
传染病的预防控制措施

\end{pastquestion}

\PointSeparator

\section{人畜共患病（zoonosis）}\label{users__fpb__26_spring__pkusph_study__tmp__pdfs__prevention-medicine-past-exams__03-ux6d41ux884cux75c5ux5b66ux5f80ux5e74ux9898ux8003ux70b9ux6574ux7406.md__ux4ebaux755cux5171ux60a3ux75c5zoonosis}

\FrequencyPill{4 次}

\InfoLabel{对应小节} 第十一章 传染病流行病学 / 第三节 流行过程 /
一、基本环节（一）传染源---3.受感染的动物（书中页码 第160页）

\InfoLabel{匹配依据} 10、11、12、14 级考``人畜共患病（zoonosis）''名词解释，常追问举例。

\ContentBand{知识点原文摘取}

在脊椎动物与人类之间自然传播感染的疫病称为人兽共患病或人畜共患病（zoonosis），如鼠疫、狂犬病等，可分为四类：①以动物为主（如狂犬病、森林脑炎、禽流感）；②以人为主（如人型结核、阿米巴痢疾）；③人和动物并重（如血吸虫病）；④真性共患病（病原体必须以人和动物分别作为终宿主和中间宿主，如牛绦虫病）。

\ContentBand{对应往年题}

\begin{pastquestion}

\QuestionLabel{【往年题1｜10级考题回忆.pdf｜名解】}\\
zoonosis（人畜共患病）

\end{pastquestion}

\begin{pastquestion}

\QuestionLabel{【往年题2｜12级流病口试考题回忆.pdf｜名解（题号33、34）】}\\
zoonosis；人畜共患病

\end{pastquestion}

\begin{pastquestion}

\QuestionLabel{【往年题3｜11级流行病学口试考题回忆.pdf｜名解】}\\
人畜共患疾病

\end{pastquestion}

\PointSeparator

\section{暴发调查的步骤}\label{users__fpb__26_spring__pkusph_study__tmp__pdfs__prevention-medicine-past-exams__03-ux6d41ux884cux75c5ux5b66ux5f80ux5e74ux9898ux8003ux70b9ux6574ux7406.md__ux66b4ux53d1ux8c03ux67e5ux7684ux6b65ux9aa4}

\FrequencyPill{3 次}

\InfoLabel{对应小节} 第十一章 传染病流行病学（结合实习课``暴发调查''）（书中页码
见第十一章流行过程及实习指导）

\InfoLabel{匹配依据} 14、20、21 级（尤其 20 级，多人）考``暴发调查的步骤''，与实习课``暴发调查''直接对应。

\ContentBand{知识点原文摘取}

（课本 24
章知识点文件中未检索到``暴发调查步骤''的成条原文；该内容主要来自实习课与课件，故不自行补充课本原文。建议按实习指导回答暴发调查的一般步骤：①核实诊断、确认暴发；②建立病例定义、核实病例并计算病例数；③描述三间分布（核心是绘制流行曲线、判断暴露时间）；④建立并检验病因假设（必要时开展病例对照或队列分析）；⑤采取控制措施并评价效果；⑥书面报告与总结。）

\ContentBand{对应往年题}

\begin{pastquestion}

\QuestionLabel{【往年题1｜20级回忆试题库.xlsx｜简答（题号9、22、26、47 等多处）】}\\
暴发调查的步骤（应该就是那10个）；暴发的调查步骤

\end{pastquestion}

\begin{pastquestion}

\QuestionLabel{【往年题2｜21级流病考题回忆汇总.xlsx｜简答（题号45、49）】}\\
暴发调查的步骤

\end{pastquestion}

\PointSeparator

\section{法定传染病的分类及举例}\label{users__fpb__26_spring__pkusph_study__tmp__pdfs__prevention-medicine-past-exams__03-ux6d41ux884cux75c5ux5b66ux5f80ux5e74ux9898ux8003ux70b9ux6574ux7406.md__ux6cd5ux5b9aux4f20ux67d3ux75c5ux7684ux5206ux7c7bux53caux4e3eux4f8b}

\FrequencyPill{2 次}

\InfoLabel{对应小节} 第十一章 传染病流行病学 / 第四节 二、措施（一）1.传染病报告（书中页码 第165---166页）

\InfoLabel{匹配依据} 19、21 级考``法定传染病的分类和举例/三类法定传染病每类举例''。19
级提示``法定传染病是三类40种（39+新冠）\ldots\ldots 要翻 PPT''。

\ContentBand{知识点原文摘取}

依据2025年修订实施的《传染病防治法》，我国将法定报告传染病分为甲、乙、丙三类共计40种，以及突发原因不明的传染病等。

甲类（2种）：鼠疫、霍乱。\\
乙类（27种）：新型冠状病毒感染、传染性非典型肺炎、艾滋病、病毒性肝炎、脊髓灰质炎、人感染新亚型流感、麻疹、流行性出血热、狂犬病、流行性乙型脑炎、登革热、猴痘、炭疽、细菌性和阿米巴性痢疾、肺结核、伤寒和副伤寒、流行性脑脊髓膜炎、百日咳、白喉、新生儿破伤风、猩红热、布鲁氏菌病、淋病、梅毒、钩端螺旋体病、血吸虫病、疟疾。\\
丙类（11种）：流行性感冒、流行性腮腺炎、风疹、急性出血性结膜炎、麻风病、流行性和地方性斑疹伤寒、黑热病、包虫病、丝虫病、手足口病，以及除霍乱、细菌性和阿米巴性痢疾、伤寒和副伤寒以外的感染性腹泻病。

\ContentBand{对应往年题}

\begin{pastquestion}

\QuestionLabel{【往年题1｜21级流病考题回忆汇总.xlsx｜简答（题号9、29、41）】}\\
法定传染病的分类和举例；三类法定传染病，每类举例2种以上

\end{pastquestion}

\begin{pastquestion}

\QuestionLabel{【往年题2｜19级流行病学期末口试回忆.xlsx｜简答（题号8）】}\\
（疾病监测追问）法定传染病是三类40种

\end{pastquestion}

\PointSeparator

\chapter{病毒性肝炎、感染性腹泻}\label{users__fpb__26_spring__pkusph_study__tmp__pdfs__prevention-medicine-past-exams__03-ux6d41ux884cux75c5ux5b66ux5f80ux5e74ux9898ux8003ux70b9ux6574ux7406.md__ux7b2cux5341ux4e8cux7ae0-ux75c5ux6bd2ux6027ux809dux708eux611fux67d3ux6027ux8179ux6cfb}

\section{丙型肝炎的流行过程（传播途径）与预防措施}\label{users__fpb__26_spring__pkusph_study__tmp__pdfs__prevention-medicine-past-exams__03-ux6d41ux884cux75c5ux5b66ux5f80ux5e74ux9898ux8003ux70b9ux6574ux7406.md__ux4e19ux578bux809dux708eux7684ux6d41ux884cux8fc7ux7a0bux4f20ux64adux9014ux5f84ux4e0eux9884ux9632ux63aaux65bd}

\FrequencyPill{8 次}

\InfoLabel{对应小节} 第十九章 病毒性肝炎 / 第二节 流行过程 / 三、丙型肝炎；第四节
二、措施（二）乙型、丙型和丁型肝炎的预防措施（书中页码 第19章相应小节）

\InfoLabel{匹配依据} 04、05、07、10、12、14、20、21
级反复考``丙肝的流行病学特征/流行过程/防治措施''，常追问与乙肝的区别（乙肝可治疗维持但目前不能治愈，丙肝可治愈但药物昂贵）。

\ContentBand{知识点原文摘取}

丙肝的主要传播途径为经血液传播，其次为性传播，而母婴传播概率较低。

预防措施：丙肝系非消化道传播的病毒性肝炎，以经血液传播为主，主要预防措施是高危人群筛查、诊断和治疗。①管理传染源：按乙类传染病管理，所有成人进行抗-HCV
筛查，高危人群定期筛查，抗-HCV 阳性者检测 HCV
RNA，动员病人``应治尽治''、规范抗病毒治疗。②切断传播途径：严格筛选献血者、安全注射、加强文身文眉修脚等行业工具消毒、避免共用剃须刀牙具；阻断母婴传播；预防性传播。③目前尚无特异性预防
HCV 感染的疫苗。

\ContentBand{对应往年题}

\begin{pastquestion}

\QuestionLabel{【往年题1｜04级流病考题回忆.docx｜问答】}\\
丙肝的防治措施；丙肝的流行病学特征

\end{pastquestion}

\begin{pastquestion}

\QuestionLabel{【往年题2｜21级流病考题回忆汇总.xlsx｜简答（题号31）】}\\
丙肝流行过程

\end{pastquestion}

\begin{pastquestion}

\QuestionLabel{【往年题3｜12级流病口试考题回忆.pdf｜简答（题号19、52）】}\\
丙肝的预防策略；丙肝的流行过程

\end{pastquestion}

\PointSeparator

\section{乙型肝炎的传播途径与预防措施}\label{users__fpb__26_spring__pkusph_study__tmp__pdfs__prevention-medicine-past-exams__03-ux6d41ux884cux75c5ux5b66ux5f80ux5e74ux9898ux8003ux70b9ux6574ux7406.md__ux4e59ux578bux809dux708eux7684ux4f20ux64adux9014ux5f84ux4e0eux9884ux9632ux63aaux65bd}

\FrequencyPill{7 次}

\InfoLabel{对应小节} 第十九章 病毒性肝炎 / 第二节 流行过程 / 二、乙型肝炎；第四节 二、措施（二）（书中页码
第19章相应小节）

\InfoLabel{匹配依据} 10、11、12、14、20、21
级考``乙肝的传播途径和预防措施''，部分追问乙肝预防中最重要的是什么（新生儿乙肝疫苗接种）。

\ContentBand{知识点原文摘取}

乙肝的主要传播途径包括经血液传播、母婴传播和性传播。

预防措施：乙肝和丁肝的主要预防措施是乙肝疫苗接种，同时加强乙肝的筛查、诊断和治疗。①管理传染源：按乙类传染病管理，筛查（HBsAg、抗-HBs、抗-HBc），筛查后管理与抗病毒治疗。②切断传播途径：预防经血液传播（严格筛选献血者、安全注射、器械消毒）、阻断母婴传播（HBsAg
阳性孕妇所生新生儿出生 12 小时内注射 HBIG 并接种乙肝疫苗）、预防性传播。③保护易感人群：乙肝疫苗接种是预防 HBV
感染最有效的方法，接种对象主要是新生儿，按 0、1、6 个月程序全程接种 3 剂；母亲 HBsAg
阳性或不详的新生儿采用人工自动免疫联合人工被动免疫（HBIG）。

\ContentBand{对应往年题}

\begin{pastquestion}

\QuestionLabel{【往年题1｜20级回忆试题库.xlsx｜综合分析（题号8）】}\\
简述乙肝的传播途径和预防措施

\end{pastquestion}

\begin{pastquestion}

\QuestionLabel{【往年题2｜12级流病口试考题回忆.pdf｜简答（题号15、60）】}\\
乙肝病毒的传播途径和预防措施（注意乙肝病毒的传播途径不包括日常接触传播，书上有误）

\end{pastquestion}

\begin{pastquestion}

\QuestionLabel{【往年题3｜10级考题回忆.pdf｜大题】}\\
乙肝的传播途径和预防

\end{pastquestion}

\PointSeparator

\section{甲型肝炎与戊型肝炎的传播途径及预防措施}\label{users__fpb__26_spring__pkusph_study__tmp__pdfs__prevention-medicine-past-exams__03-ux6d41ux884cux75c5ux5b66ux5f80ux5e74ux9898ux8003ux70b9ux6574ux7406.md__ux7532ux578bux809dux708eux4e0eux620aux578bux809dux708eux7684ux4f20ux64adux9014ux5f84ux53caux9884ux9632ux63aaux65bd}

\FrequencyPill{6 次}

\InfoLabel{对应小节} 第十九章 病毒性肝炎 / 第二节 流行过程 / 一、甲型肝炎、五、戊型肝炎；第四节
二、措施（一）甲型和戊型肝炎的预防措施（书中页码 第19章相应小节）

\InfoLabel{匹配依据} 04、05、07、12、14、20
级考``甲肝的传播途径和预防措施/甲型和戊型肝炎的防治措施''，常作为大题。

\ContentBand{知识点原文摘取}

甲肝的主要传播途径是消化道传播，易感人群会因摄入被 HAV
感染者粪便污染的食物或水而感染，也可能通过与感染者直接或间接接触感染。戊肝的主要传播途径是消化道传播，经水传播造成流行多见。

预防措施：甲肝和戊肝应采取以切断消化道传播和预防接种相结合的综合性预防措施。①管理传染源：按乙类传染病管理，24
小时内网络直报，治愈前不得从事直接为顾客服务的工作。②切断传播途径：疫源地消毒；做好食品卫生、饮用水消毒、餐具消毒和粪便无害化处理，防止``病从口入''。③保护易感人群：人工自动免疫（甲肝疫苗已纳入国家免疫规划；戊肝有重组戊肝疫苗，推荐对畜牧养殖者、疫区旅行者、餐饮业人群等高危人群接种）；人工被动免疫（人免疫球蛋白用于甲肝接触者的紧急预防）。

\ContentBand{对应往年题}

\begin{pastquestion}

\QuestionLabel{【往年题1｜07级考题回忆.docx｜问答】}\\
甲肝与戊肝的预防措施；甲肝的传播途径及预防措施

\end{pastquestion}

\begin{pastquestion}

\QuestionLabel{【往年题2｜04级流病考题回忆.docx｜大题】}\\
甲型和戊型肝炎的防治措施

\end{pastquestion}

\begin{pastquestion}

\QuestionLabel{【往年题3｜12级流病口试考题回忆.pdf｜大题（题号35）】}\\
甲肝乙肝流行过程；甲肝和戊肝的预防措施

\end{pastquestion}

\PointSeparator

\section{乙肝抗原抗体系统（乙肝五项）及生物标志的流行病学意义}\label{users__fpb__26_spring__pkusph_study__tmp__pdfs__prevention-medicine-past-exams__03-ux6d41ux884cux75c5ux5b66ux5f80ux5e74ux9898ux8003ux70b9ux6574ux7406.md__ux4e59ux809dux6297ux539fux6297ux4f53ux7cfbux7edfux4e59ux809dux4e94ux9879ux53caux751fux7269ux6807ux5fd7ux7684ux6d41ux884cux75c5ux5b66ux610fux4e49}

\FrequencyPill{4 次}

\InfoLabel{对应小节} 第十九章 病毒性肝炎 / 第一节 病原学和临床特征 /
二、临床特征（二）乙型肝炎---2.实验室检测指标及意义（书中页码 第4页）

\InfoLabel{匹配依据} 10、11、14、21 级考``乙肝抗原抗体系统（乙肝五项）及生物标志的流行病学意义''。

\ContentBand{知识点原文摘取}

（1）HBsAg：是现症感染的标志。\\
（2）抗-HBs：为乙肝恢复期或接种乙肝疫苗后产生的保护性抗体，是机体获得免疫力的标志。\\
（3）HBeAg：是病毒复制活跃和传染性较强的标志，与 HBV DNA 具有良好的相关性。\\
（4）抗-HBe：非保护性抗体，为病毒复制水平和传染性降低的标志。\\
（5）HBcAg：与 HBV DNA 呈正相关，是病毒复制和具有传染性的标志。\\
（6）抗-HBc：非保护性抗体，分 IgM 和 IgG；抗-HBc IgM 是急性感染或慢性感染病情活动的标志，抗-HBc IgG
是慢性感染或既往感染的标志。\\
（7）HBV DNA：是病毒复制和具有传染性的直接标志。

\ContentBand{对应往年题}

\begin{pastquestion}

\QuestionLabel{【往年题1｜21级流病考题回忆汇总.xlsx｜简答（题号40）】}\\
乙肝抗原抗体系统及生物标志流行病学意义；乙肝五项

\end{pastquestion}

\begin{pastquestion}

\QuestionLabel{【往年题2｜10级考题回忆.pdf｜简答】}\\
乙肝抗原抗体系统流行病学意义

\end{pastquestion}

\begin{pastquestion}

\QuestionLabel{【往年题3｜14级流病考题回忆汇总.xls｜简答（题号29）】}\\
简述乙肝抗原抗体系统及生物学标志物的流行病学意义

\end{pastquestion}

\PointSeparator

\section{甲型肝炎与乙型肝炎流行过程的比较}\label{users__fpb__26_spring__pkusph_study__tmp__pdfs__prevention-medicine-past-exams__03-ux6d41ux884cux75c5ux5b66ux5f80ux5e74ux9898ux8003ux70b9ux6574ux7406.md__ux7532ux578bux809dux708eux4e0eux4e59ux578bux809dux708eux6d41ux884cux8fc7ux7a0bux7684ux6bd4ux8f83}

\FrequencyPill{4 次}

\InfoLabel{对应小节} 第十九章 病毒性肝炎 / 第二节 流行过程 / 一、甲型肝炎、二、乙型肝炎（书中页码
第19章相应小节）

\InfoLabel{匹配依据} 11、12、14、20 级考``甲肝和乙肝流行过程的比较''。

\ContentBand{知识点原文摘取}

甲肝：传染源为甲肝病人和隐性感染者；主要传播途径是消化道传播（粪-口途径）；人群普遍易感，感染后可获持久免疫力。

乙肝：传染源为急、慢性乙肝病人和 HBV
携带者（健康病原携带者数量多，是非常重要的传染源）；主要传播途径为经血液传播、母婴传播和性传播；人群普遍易感，感染
HBV 时的年龄是影响慢性化的主要因素（新生儿及1岁以下婴幼儿感染慢性化风险高达90\%，成人低于5\%）。

\ContentBand{对应往年题}

\begin{pastquestion}

\QuestionLabel{【往年题1｜11级流行病学口试考题回忆.pdf｜简答】}\\
甲肝乙肝流行过程比较

\end{pastquestion}

\begin{pastquestion}

\QuestionLabel{【往年题2｜20级回忆试题库.xlsx｜综合分析（题号40）】}\\
甲肝和乙肝的流行过程比较

\end{pastquestion}

\PointSeparator

\section{乙型肝炎与丁型肝炎的预防措施}\label{users__fpb__26_spring__pkusph_study__tmp__pdfs__prevention-medicine-past-exams__03-ux6d41ux884cux75c5ux5b66ux5f80ux5e74ux9898ux8003ux70b9ux6574ux7406.md__ux4e59ux578bux809dux708eux4e0eux4e01ux578bux809dux708eux7684ux9884ux9632ux63aaux65bd}

\FrequencyPill{3 次}

\InfoLabel{对应小节} 第十九章 病毒性肝炎 / 第二节 流行过程 / 四、丁型肝炎；第四节 二、措施（二）（书中页码
第19章相应小节）

\InfoLabel{匹配依据} 04、14、21 级考``乙肝和丁肝的预防措施''。

\ContentBand{知识点原文摘取}

丁肝的主要传播途径为经血液传播和性传播，很少发生母婴传播。

乙肝和丁肝的主要预防措施是乙肝疫苗接种。对 HBV 有免疫力者也不再感染 HDV，因此接种乙肝疫苗可通过预防 HBV
感染而间接达到预防 HDV（丁肝）感染的目的。在出生时进行人工自动免疫和人工被动免疫，也可预防 HDV 从 HBV/HDV
同时感染或重叠感染的母亲传播给子女。

\ContentBand{对应往年题}

\begin{pastquestion}

\QuestionLabel{【往年题1｜21级流病考题回忆汇总.xlsx｜简答（题号50）】}\\
乙肝和丁肝的预防措施

\end{pastquestion}

\begin{pastquestion}

\QuestionLabel{【往年题2｜04级流病考题回忆.docx｜问答】}\\
乙肝和丁肝的预防措施

\end{pastquestion}

\PointSeparator

\section{戊型肝炎的流行过程与特征}\label{users__fpb__26_spring__pkusph_study__tmp__pdfs__prevention-medicine-past-exams__03-ux6d41ux884cux75c5ux5b66ux5f80ux5e74ux9898ux8003ux70b9ux6574ux7406.md__ux620aux578bux809dux708eux7684ux6d41ux884cux8fc7ux7a0bux4e0eux7279ux5f81}

\FrequencyPill{3 次}

\InfoLabel{对应小节} 第十九章 病毒性肝炎 / 第二节 流行过程 / 五、戊型肝炎（书中页码 第19章相应小节）

\InfoLabel{匹配依据} 04、12、20 级考``戊型肝炎的流行病学特征/流行过程''。

\ContentBand{知识点原文摘取}

戊肝的主要传播途径是消化道传播，经水传播造成流行多见，易感人群通过饮用被 HEV
污染的水而感染，多在雨季和洪水后出现戊肝暴发或流行。食用被 HEV
污染的食物，如未加工的贝类以及被感染动物未煮熟的肉、肝脏等，可引起散发或暴发。（由 3 型和 4 型 HEV
引起的戊肝可能存在动物宿主。）

\ContentBand{对应往年题}

\begin{pastquestion}

\QuestionLabel{【往年题1｜12级流病口试考题回忆.pdf｜综合分析（题号58）】}\\
戊型肝炎的流行过程

\end{pastquestion}

\begin{pastquestion}

\QuestionLabel{【往年题2｜20级回忆试题库.xlsx｜简答（题号36）】}\\
戊肝的流行过程

\end{pastquestion}

\PointSeparator

\chapter{呼吸系统传染病}\label{users__fpb__26_spring__pkusph_study__tmp__pdfs__prevention-medicine-past-exams__03-ux6d41ux884cux75c5ux5b66ux5f80ux5e74ux9898ux8003ux70b9ux6574ux7406.md__ux7b2cux5341ux4e09ux7ae0-ux547cux5438ux7cfbux7edfux4f20ux67d3ux75c5}

\section{流感病毒抗原变异（抗原漂移、抗原转换）及其与流行强度的关系}\label{users__fpb__26_spring__pkusph_study__tmp__pdfs__prevention-medicine-past-exams__03-ux6d41ux884cux75c5ux5b66ux5f80ux5e74ux9898ux8003ux70b9ux6574ux7406.md__ux6d41ux611fux75c5ux6bd2ux6297ux539fux53d8ux5f02ux6297ux539fux6f02ux79fbux6297ux539fux8f6cux6362ux53caux5176ux4e0eux6d41ux884cux5f3aux5ea6ux7684ux5173ux7cfb}

\FrequencyPill{7 次}

\InfoLabel{对应小节} 第二十三章 流行性感冒 / 第一节 病原学 / 二、抗原变异（书中页码 第23章相应小节）

\InfoLabel{匹配依据} 04、10、11、12、14、20、21
级反复考``流感病毒的抗原变异类型（抗原漂移、抗原转换）及其与流行强度/流行规模的关系''。

\ContentBand{知识点原文摘取}

流感病毒可引起季节性流行和流感大流行，主要是其 HA 和 NA
的抗原性容易发生变异所致，其中甲型流感病毒的抗原变异最容易发生，它同流感大流行密切相关。流感病毒抗原变异幅度的大小直接影响流行规模。

\begin{enumerate}
\def\labelenumi{\arabic{enumi}.}
\item
  抗原漂移（antigenic
  drift）：是指流感病毒亚型内部经常发生的小幅度的变异，属于量变。漂移的结果常引起流感的季节性流行（主要原因是病毒抗原基因突变）。
\item
  抗原转换（antigenic shift）：是指流感病毒抗原变异幅度大，形成新的亚型，即新毒株的 HA 和/或 NA
  与前次流行株不同，属于质变。如 H1N1 转换成
  H2N2，转换的结果常引起流感大流行（机制主要是不同病毒表面抗原发生基因重配）。
\end{enumerate}

\ContentBand{对应往年题}

\begin{pastquestion}

\QuestionLabel{【往年题1｜04级流病考题回忆.docx｜简答】}\\
流感病毒的抗原变异和流行强度范围的关系

\end{pastquestion}

\begin{pastquestion}

\QuestionLabel{【往年题2｜20级回忆试题库.xlsx｜简答（题号18）】}\\
抗原变异主要类型（抗原转换和抗原漂移）

\end{pastquestion}

\begin{pastquestion}

\QuestionLabel{【往年题3｜11级流行病学口试考题回忆.pdf｜简答】}\\
流感病毒抗原变异及流行强度的关系（抗原漂变和转变）

\end{pastquestion}

\PointSeparator

\section{WHO
针对流感大流行的措施与对策}\label{users__fpb__26_spring__pkusph_study__tmp__pdfs__prevention-medicine-past-exams__03-ux6d41ux884cux75c5ux5b66ux5f80ux5e74ux9898ux8003ux70b9ux6574ux7406.md__who-ux9488ux5bf9ux6d41ux611fux5927ux6d41ux884cux7684ux63aaux65bdux4e0eux5bf9ux7b56}

\FrequencyPill{2 次}

\InfoLabel{对应小节} 第二十三章 流行性感冒 / 第二节 流行过程、第三节 流行特征及防控（书中页码 第23章相应小节）

\InfoLabel{匹配依据} 10、11 级考``WHO 针对流感（大流行/暴发）的措施和对策/全球大流行的策略''。

\ContentBand{知识点原文摘取}

（课本第二十三章重点阐述流感的病原学、抗原变异、流行过程与流行特征，并指出流感是第一个实行全球监测的传染病，中国是全球流感监测的重要哨点。关于``WHO
针对流感大流行的具体措施与对策''的成条原文，课本知识点文件中未检索到完整列举，建议结合课件作答：以全球流感监测网络为核心，加强病原学和疫情监测、疫苗研发与接种、抗病毒药物储备、非药物干预（隔离、减少聚集）及大流行应急预案等综合措施。）

\ContentBand{对应往年题}

\begin{pastquestion}

\QuestionLabel{【往年题1｜10级考题回忆.pdf｜简答】}\\
WHO 针对流感爆发的措施和对策

\end{pastquestion}

\begin{pastquestion}

\QuestionLabel{【往年题2｜11级流行病学口试考题回忆.pdf｜简答】}\\
who 流感大流行的策略

\end{pastquestion}

\PointSeparator

\section{耐多药结核病（MDR-TB）}\label{users__fpb__26_spring__pkusph_study__tmp__pdfs__prevention-medicine-past-exams__03-ux6d41ux884cux75c5ux5b66ux5f80ux5e74ux9898ux8003ux70b9ux6574ux7406.md__ux8010ux591aux836fux7ed3ux6838ux75c5mdr-tb}

\FrequencyPill{3 次}

\InfoLabel{对应小节} 第二十章 结核病 / 第二节（相应小节）三、耐药结核病（书中页码 第20章相应小节）

\InfoLabel{匹配依据} 14、21 级考``耐多药结核病（multidrug-resistant
tuberculosis，MDR-TB）/多耐药性结核菌''名词解释。

\ContentBand{知识点原文摘取}

耐多药结核病（multidrug-resistant
tuberculosis，MDR-TB）是指结核患者感染的结核分枝杆菌在体外被证实至少对异烟肼、利福平耐药。异烟肼和利福平是一线抗结核药物中效力最强的两种药物，一旦结核分枝杆菌发生异烟肼和利福平耐药，常规化疗就很难发挥治疗作用。MDR-TB
治疗需要采用二线抗结核药物，治疗成本高昂。

\ContentBand{对应往年题}

\begin{pastquestion}

\QuestionLabel{【往年题1｜21级流病考题回忆汇总.xlsx｜名解（题号38）】}\\
耐多药结核病（英文）

\end{pastquestion}

\begin{pastquestion}

\QuestionLabel{【往年题2｜14级流病考题回忆汇总.xls｜名解（题号17）】}\\
意向性分析，多耐药性结核菌

\end{pastquestion}

\PointSeparator

\section{短程督导治疗（DOTS
策略）}\label{users__fpb__26_spring__pkusph_study__tmp__pdfs__prevention-medicine-past-exams__03-ux6d41ux884cux75c5ux5b66ux5f80ux5e74ux9898ux8003ux70b9ux6574ux7406.md__ux77edux7a0bux7763ux5bfcux6cbbux7597dots-ux7b56ux7565}

\FrequencyPill{1 次}

\InfoLabel{对应小节} 第二十章 结核病 / 第一节 概述（化疗发展）及第四节 预防控制策略（DOTS 策略）（书中页码
第20章相应小节）

\InfoLabel{匹配依据} 04 级考``短程督导治疗''名词解释，对应课本现代结核病控制策略 DOTS。

\ContentBand{知识点原文摘取}

1995 年起，结核病高负担国家和地区全面实施了以直接面视下的短程化疗（directly observed treatment, short
course；DOTS）为核心的现代结核病控制策略。结核病化疗的原则是：早期、联合、适量、规律和全程用药。一旦确定化疗方案并开始治疗，卫生服务提供者就要采取有效的健康教育和经常性的督导等管理措施，以确保患者连续不间断地治疗，直至完成规定的疗程。DOTS
策略的推行大幅度提高了患者的治愈率和发现率，同时可以防止耐药菌株的产生。

\ContentBand{对应往年题}

\begin{pastquestion}

\QuestionLabel{【往年题1｜04级流病考题回忆.docx｜名解】}\\
短程督导治疗

\end{pastquestion}

\PointSeparator

\section{人感染高致病性禽流感}\label{users__fpb__26_spring__pkusph_study__tmp__pdfs__prevention-medicine-past-exams__03-ux6d41ux884cux75c5ux5b66ux5f80ux5e74ux9898ux8003ux70b9ux6574ux7406.md__ux4ebaux611fux67d3ux9ad8ux81f4ux75c5ux6027ux79bdux6d41ux611f}

\FrequencyPill{2 次}

\InfoLabel{对应小节} 第二十三章 流行性感冒（动物流感病毒跨种属传播）（书中页码 第23章相应小节）

\InfoLabel{匹配依据} 12、14 级考``人感染高致病性禽流感/人高致病性禽流感''名词解释。

\ContentBand{知识点原文摘取}

（课本第二十三章指出：某些动物流感病毒偶尔跨越种属屏障引起人类感染，对公共卫生安全造成威胁，流感病毒的跨物种传播正日益引起关注。课本知识点文件中未检索到``人感染高致病性禽流感''的独立成条定义原文，建议结合课件作答：人感染高致病性禽流感是指由高致病性禽流感病毒（如
H5N1、H7N9
等亚型）跨种属感染人类所引起的急性呼吸道传染病，病情重、病死率高，属乙类传染病中按甲类管理/重点监测的人感染新亚型流感范畴。）

\ContentBand{对应往年题}

\begin{pastquestion}

\QuestionLabel{【往年题1｜14级流病考题回忆汇总.xls｜名解（题号9、12）】}\\
人感染高致病性禽流感

\end{pastquestion}

\begin{pastquestion}

\QuestionLabel{【往年题2｜12级流病口试考题回忆.pdf｜名解（题号49）】}\\
人高致病性禽流感

\end{pastquestion}

\PointSeparator

\chapter{性传播疾病}\label{users__fpb__26_spring__pkusph_study__tmp__pdfs__prevention-medicine-past-exams__03-ux6d41ux884cux75c5ux5b66ux5f80ux5e74ux9898ux8003ux70b9ux6574ux7406.md__ux7b2cux5341ux56dbux7ae0-ux6027ux4f20ux64adux75beux75c5}

\section{性传播疾病（STD）的预防措施}\label{users__fpb__26_spring__pkusph_study__tmp__pdfs__prevention-medicine-past-exams__03-ux6d41ux884cux75c5ux5b66ux5f80ux5e74ux9898ux8003ux70b9ux6574ux7406.md__ux6027ux4f20ux64adux75beux75c5stdux7684ux9884ux9632ux63aaux65bd}

\FrequencyPill{5 次}

\InfoLabel{对应小节} 第二十一章 性传播疾病 / 第三节 预防策略与措施 / 二、措施（书中页码 第21章相应小节）

\InfoLabel{匹配依据} 10、11、14、20、21 级考``性传播疾病的预防（控制）措施''，21
级老师喜欢追问举例（如不安全性行为指什么）。

\ContentBand{知识点原文摘取}

和其他传染病一样，STDs 的预防控制包括针对传染源、传播途径、易感人群和影响因素的多种措施以及 STDs 的监测。

①针对传染源：及时发现、规范诊断和治疗 STDs
感染者（大量未被发现的感染者是最重要的传染源）。②针对传播途径：改变不安全性行为------通过健康教育和健康促进（同伴教育、流动宣传、电话咨询等）传播安全性行为知识，减少性伴数，促进安全套的使用；防止母婴传播（规范化使用药物或疫苗）。③针对易感人群：健康教育是控制
STDs 最经济有效的方法；持续正确使用安全套可有效预防包括 HIV
感染在内的性传播感染；着重对青少年开展早期性教育；乙肝疫苗和人乳头瘤病毒（HPV）疫苗是预防 STDs
的重大突破。④监测。

\ContentBand{对应往年题}

\begin{pastquestion}

\QuestionLabel{【往年题1｜21级流病考题回忆汇总.xlsx｜简答（题号25、30）】}\\
性传播疾病的预防措施（老师喜欢举例，如不安全的性行为是什么）

\end{pastquestion}

\begin{pastquestion}

\QuestionLabel{【往年题2｜11级流行病学口试考题回忆.pdf｜简答】}\\
STD 预防措施

\end{pastquestion}

\begin{pastquestion}

\QuestionLabel{【往年题3｜14级流病考题回忆汇总.xls｜简答（题号28）】}\\
性传播疾病预防措施

\end{pastquestion}

\PointSeparator

\section{影响性传播疾病流行的社会因素}\label{users__fpb__26_spring__pkusph_study__tmp__pdfs__prevention-medicine-past-exams__03-ux6d41ux884cux75c5ux5b66ux5f80ux5e74ux9898ux8003ux70b9ux6574ux7406.md__ux5f71ux54cdux6027ux4f20ux64adux75beux75c5ux6d41ux884cux7684ux793eux4f1aux56e0ux7d20}

\FrequencyPill{4 次}

\InfoLabel{对应小节} 第二十一章 性传播疾病 / 第二节 流行过程 / 四、影响流行过程的因素（二）社会因素（书中页码
第6页）

\InfoLabel{匹配依据} 10、19、20 级考``影响性传播疾病的社会因素/性传播疾病的社会影响因素''。

\ContentBand{知识点原文摘取}

社会因素：①经济快速发展和城市化加速，人口流动性增强，性活跃人群高度集中与频繁迁徙为 STDs
传播扩散提供条件；②健康教育不足（许多青少年从未接受过 STDs 教育）；③卫生资源分配不均衡地区 STDs
相关知识匮乏，无保护的商业性行为成为部分地区传播的重要途径；④社会舆论和心理压力使许多患者不愿就医或选择非正规渠道治疗，导致病情迁延、耐药；⑤短期内与多个性伴发生性行为显著提高感染风险；⑥吸毒和贩毒问题也是不可忽视的社会因素。

\ContentBand{对应往年题}

\begin{pastquestion}

\QuestionLabel{【往年题1｜19级流行病学期末口试回忆.xlsx｜简答（题号19）】}\\
影响性传播疾病的社会因素

\end{pastquestion}

\begin{pastquestion}

\QuestionLabel{【往年题2｜20级回忆试题库.xlsx｜简答（题号37）】}\\
性传染病的社会因素

\end{pastquestion}

\begin{pastquestion}

\QuestionLabel{【往年题3｜10级考题回忆.pdf｜简答】}\\
性传播疾病的社会影响因素

\end{pastquestion}

\PointSeparator

\section{艾滋病窗口期}\label{users__fpb__26_spring__pkusph_study__tmp__pdfs__prevention-medicine-past-exams__03-ux6d41ux884cux75c5ux5b66ux5f80ux5e74ux9898ux8003ux70b9ux6574ux7406.md__ux827eux6ecbux75c5ux7a97ux53e3ux671f}

\FrequencyPill{3 次}

\InfoLabel{对应小节} 第二十一章 性传播疾病 / 第四节 艾滋病 / 一、概述（二）疾病自然史---1.急性期（书中页码
第8页）

\InfoLabel{匹配依据} 10、11、12 级考``艾滋病窗口期''名词解释。

\ContentBand{知识点原文摘取}

（课本知识点文件中未检索到``窗口期''的独立成条定义原文。课本在 HIV 感染急性期中描述：``此期血液中可检测到 HIV
RNA 和 P24 抗原，HIV 抗体在 2～3 周内逐渐由阴转阳。''据此，窗口期通常指从 HIV 感染到血液中能检测出 HIV
抗体之前的这段时期------此期感染者体内已有病毒复制并具有传染性，但抗体检测可呈阴性。建议结合课件作答其概念及流行病学意义------窗口期内抗体检测阴性者仍可能传播
HIV，是输血等防控的重要隐患。）

\ContentBand{对应往年题}

\begin{pastquestion}

\QuestionLabel{【往年题1｜11级流行病学口试考题回忆.pdf｜名解】}\\
艾滋病的窗口期

\end{pastquestion}

\begin{pastquestion}

\QuestionLabel{【往年题2｜12级流病口试考题回忆.pdf｜名解（题号57）】}\\
艾滋病窗口期

\end{pastquestion}

\begin{pastquestion}

\QuestionLabel{【往年题3｜10级考题回忆.pdf｜名解】}\\
艾滋病窗口期、causation of disease

\end{pastquestion}

\PointSeparator

\section{艾滋病的传播途径、危害与预防控制}\label{users__fpb__26_spring__pkusph_study__tmp__pdfs__prevention-medicine-past-exams__03-ux6d41ux884cux75c5ux5b66ux5f80ux5e74ux9898ux8003ux70b9ux6574ux7406.md__ux827eux6ecbux75c5ux7684ux4f20ux64adux9014ux5f84ux5371ux5bb3ux4e0eux9884ux9632ux63a7ux5236}

\FrequencyPill{4 次}

\InfoLabel{对应小节} 第二十一章 性传播疾病 / 第四节 艾滋病 / 一、概述（三）危害；三、预防策略与措施（书中页码
第8---9页及相应小节）

\InfoLabel{匹配依据} 04、05、10、11 级考``艾滋病的传播途径/主要危险因素/危害/预防控制措施/我国流行特征''，04
级追问我国最可能的传播途径（性传播）。

\ContentBand{知识点原文摘取}

AIDS
由人类免疫缺陷病毒（HIV）引起。传播途径：不安全的性行为、注射吸毒时共用针具、输入受污染的血液或血液制品、母婴传播是
AIDS 的主要传播途径（全球 90\% 以上婴儿和儿童的 HIV 感染是通过母婴传播）。

危害：①影响人口质量和人口结构；②加剧贫穷和不平等；③危害家庭；④偏见和歧视；⑤卫生服务压力；⑥社会经济和政治影响。

预防控制：针对传染源（发现、随访、治疗管理）、针对传播途径（安全性行为、安全套、清洁针具、阻断母婴传播、血液安全）、针对易感人群（健康教育，尤其青少年；自愿性男性包皮环切术；HIV
暴露前预防 PrEP 与暴露后预防 PEP）。我国近年学生 HIV 感染疫情上升较快，传播途径以男男性行为传播为主。

\ContentBand{对应往年题}

\begin{pastquestion}

\QuestionLabel{【往年题1｜04级流病考题回忆.docx｜问答】}\\
艾滋病的传播途径，危害，控制方法。我国最可能的传播途径是什么（性传播）

\end{pastquestion}

\begin{pastquestion}

\QuestionLabel{【往年题2｜05流病考题回忆.doc｜问答】}\\
艾滋病的主要危险因素和预防控制措施

\end{pastquestion}

\begin{pastquestion}

\QuestionLabel{【往年题3｜10级考题回忆.pdf｜简答】}\\
我国现阶段艾滋病的流行特征

\end{pastquestion}

\PointSeparator

\section{STD 监测的目的及意义 /
性传播疾病（名解）}\label{users__fpb__26_spring__pkusph_study__tmp__pdfs__prevention-medicine-past-exams__03-ux6d41ux884cux75c5ux5b66ux5f80ux5e74ux9898ux8003ux70b9ux6574ux7406.md__std-ux76d1ux6d4bux7684ux76eeux7684ux53caux610fux4e49-ux6027ux4f20ux64adux75beux75c5ux540dux89e3}

\FrequencyPill{3 次}

\InfoLabel{对应小节} 第二十一章 性传播疾病 / 概述、第三节（四）监测、第四节（四）监测工作（书中页码
第21章相应小节）

\InfoLabel{匹配依据} 14、20 级考``STD 监测的目的及意义/简述性传播疾病的目的和意义''；21
级名解考``性传播疾病''。

\ContentBand{知识点原文摘取}

性传播疾病（sexually transmitted
diseases，STDs）：以性行为接触或类似性行为接触为主要传播途径的、可引起泌尿生殖器官及附属淋巴系统病变的一类疾病，也可导致全身主要器官的病变（WHO
定义）。有临床症状的性传播感染（STIs）称为 STDs。

监测意义：HIV 感染/AIDS
病例报告系统通过规范病例登记、数据收集与管理，实现对疫情的实时监测和动态评估，整合病例基本信息、传播途径、随访治疗情况等多维度数据，为流行病学研究、政策制定及公共卫生干预提供科学依据。

\ContentBand{对应往年题}

\begin{pastquestion}

\QuestionLabel{【往年题1｜14级流病考题回忆汇总.xls｜简答（题号35）】}\\
STD 监测的目的及意义

\end{pastquestion}

\begin{pastquestion}

\QuestionLabel{【往年题2｜20级回忆试题库.xlsx｜综合分析（题号13）】}\\
简述性传播疾病的目的和意义

\end{pastquestion}

\begin{pastquestion}

\QuestionLabel{【往年题3｜21级流病考题回忆汇总.xlsx｜名解（题号4、33、40）】}\\
性传播疾病

\end{pastquestion}

\PointSeparator

\chapter{慢性病流行病学}\label{users__fpb__26_spring__pkusph_study__tmp__pdfs__prevention-medicine-past-exams__03-ux6d41ux884cux75c5ux5b66ux5f80ux5e74ux9898ux8003ux70b9ux6574ux7406.md__ux7b2cux5341ux4e94ux7ae0-ux6162ux6027ux75c5ux6d41ux884cux75c5ux5b66}

\section{慢性非传染性疾病（NCD）的主要危险因素（影响因素）}\label{users__fpb__26_spring__pkusph_study__tmp__pdfs__prevention-medicine-past-exams__03-ux6d41ux884cux75c5ux5b66ux5f80ux5e74ux9898ux8003ux70b9ux6574ux7406.md__ux6162ux6027ux975eux4f20ux67d3ux6027ux75beux75c5ncdux7684ux4e3bux8981ux5371ux9669ux56e0ux7d20ux5f71ux54cdux56e0ux7d20}

\FrequencyPill{8 次}

\InfoLabel{对应小节} 第十四章 慢性病流行病学 / 第一节 概述 / 三、慢性病的主要影响因素；第二节
二、主要危险因素的流行特征（书中页码 第14章相应小节）

\InfoLabel{匹配依据} 04、10、11、12、14、19、20、21
级反复考``NCD（慢性非传染性疾病）的危险因素/病因/影响因素''，10
级提示``除了那四条（吸烟、饮酒、不健康饮食、缺乏身体活动），还要答生物因素、环境因素，往病因链上达''。

\ContentBand{知识点原文摘取}

慢性病的主要可改变影响因素可以分为社会驱动因素、环境驱动因素、行为因素和代谢因素四大类；还包括不可改变因素，如年龄、性别、种族、遗传等。

\begin{itemize}
\tightlist
\item
  社会驱动因素：政策、经济发展水平、科技进步、社会文化习俗等。
\item
  环境驱动因素：如空气污染等。
\item
  行为因素：高盐、高脂肪、高糖饮食，蔬果摄入不足，身体活动不足，吸烟，过量饮酒，长期睡眠不足等。
\item
  代谢因素：高血压、高血糖、高血脂（尤其 LDL-C 升高）、肥胖、肾功能减退等。
\end{itemize}

主要危险因素及其流行特征包括：高血压、吸烟、不健康的饮食习惯、糖尿病、超重和肥胖、空气污染、过量饮酒等。

\ContentBand{对应往年题}

\begin{pastquestion}

\QuestionLabel{【往年题1｜10级考题回忆.pdf｜简答】}\\
我国 NCD 发病的危险因素（除了那四条，还要答生物因素、环境因素（自然\&社会）、医疗保健机制，往病因链上达）

\end{pastquestion}

\begin{pastquestion}

\QuestionLabel{【往年题2｜12级流病口试考题回忆.pdf｜简答（题号2、40）】}\\
NCD 的危险因素（不只吸烟饮酒运动饮食，要答全面，根据病因链回答）；简述 NCD 的病因

\end{pastquestion}

\begin{pastquestion}

\QuestionLabel{【往年题3｜21级流病考题回忆汇总.xlsx｜简答（题号23、4）】}\\
NCD 的病因；NCD 发病的危险因素

\end{pastquestion}

\PointSeparator

\section{心血管疾病一级预防的策略和措施}\label{users__fpb__26_spring__pkusph_study__tmp__pdfs__prevention-medicine-past-exams__03-ux6d41ux884cux75c5ux5b66ux5f80ux5e74ux9898ux8003ux70b9ux6574ux7406.md__ux5fc3ux8840ux7ba1ux75beux75c5ux4e00ux7ea7ux9884ux9632ux7684ux7b56ux7565ux548cux63aaux65bd}

\FrequencyPill{4 次}

\InfoLabel{对应小节} 第十四章 慢性病流行病学 / 第三节 预防策略与措施；第九章
预防策略（一级预防、全人群策略与高危策略）（书中页码 第14章、第9章相应小节）

\InfoLabel{匹配依据} 04、05、20 级考``心血管疾病（高血压）一级预防的策略和措施''。

\ContentBand{知识点原文摘取}

慢性病的预防应特别强调第一级预防中的根本预防（在人群水平上预防危险因素的流行），从生命早期开始、贯穿生命全过程（全生命周期策略），整合以个体为基础的高危策略和以人群为基础的全人群策略。慢性病防控应是政府主导、多部门协作、全社会参与的系统工程，围绕导致疾病负担的主要慢性病的共同和可改变的危险因素来进行：一方面建立支持性的环境，使个体有机会作出健康选择；另一方面提高个体的健康素养。

心血管疾病一级预防针对其主要危险因素（高血压、吸烟、不健康饮食/高盐、缺乏身体活动、超重肥胖、糖尿病、高血脂、过量饮酒、空气污染等），采取人群策略（控烟、减盐、改善膳食、促进身体活动等）与高危策略（对高危个体进行危险因素评估与干预，如戒烟限酒、适量运动、控制体重、控制血压血脂血糖）相结合的综合措施。

\ContentBand{对应往年题}

\begin{pastquestion}

\QuestionLabel{【往年题1｜20级回忆试题库.xlsx｜简答（题号16）】}\\
心血管疾病的一级预防策略和措施

\end{pastquestion}

\begin{pastquestion}

\QuestionLabel{【往年题2｜05流病考题回忆.doc｜问答】}\\
心血管疾病的一级预防策略及实施

\end{pastquestion}

\begin{pastquestion}

\QuestionLabel{【往年题3｜04级流病考题回忆.docx｜简答】}\\
高血压的一级预防；心血管疾病一级预防的策略和措施

\end{pastquestion}

\PointSeparator

\section{心血管疾病的社区综合防治}\label{users__fpb__26_spring__pkusph_study__tmp__pdfs__prevention-medicine-past-exams__03-ux6d41ux884cux75c5ux5b66ux5f80ux5e74ux9898ux8003ux70b9ux6574ux7406.md__ux5fc3ux8840ux7ba1ux75beux75c5ux7684ux793eux533aux7efcux5408ux9632ux6cbb}

\FrequencyPill{3 次}

\InfoLabel{对应小节} 第十四章 慢性病流行病学 / 第三节 预防策略与措施、第四节 一、心血管疾病（书中页码
第14章相应小节）

\InfoLabel{匹配依据} 04、05、07 级考``心血管疾病的社区综合防治''名词解释。

\ContentBand{知识点原文摘取}

（课本 24
章知识点文件中未检索到``心血管疾病的社区综合防治''的独立成条定义原文，该提法多见于课件。建议结合课件作答其概念：心血管疾病的社区综合防治是以社区为单位、以全人群策略与高危策略相结合为核心，通过政府主导、多部门协作、全社会参与，针对心血管疾病的共同可改变危险因素（如高血压、吸烟、高盐饮食、缺乏运动等）开展健康教育、健康促进、危险因素监测与干预、高危人群管理和病人规范治疗的综合性防治模式，如
WHO 的 MONICA 方案、社区高血压综合干预等。）

\ContentBand{对应往年题}

\begin{pastquestion}

\QuestionLabel{【往年题1｜05流病考题回忆.doc｜名解】}\\
心血管疾病的社区综合防治

\end{pastquestion}

\begin{pastquestion}

\QuestionLabel{【往年题2｜04级流病考题回忆.docx｜名解】}\\
心血管病的社区综合防治

\end{pastquestion}

\begin{pastquestion}

\QuestionLabel{【往年题3｜07级考题回忆.docx｜名解】}\\
心血管疾病的社区综合防治

\end{pastquestion}

\PointSeparator

\section{控烟策略}\label{users__fpb__26_spring__pkusph_study__tmp__pdfs__prevention-medicine-past-exams__03-ux6d41ux884cux75c5ux5b66ux5f80ux5e74ux9898ux8003ux70b9ux6574ux7406.md__ux63a7ux70dfux7b56ux7565}

\FrequencyPill{2 次}

\InfoLabel{对应小节} 第十四章 慢性病流行病学 / 第三节 预防策略与措施 /
二、措施（WHO``最合算''干预措施）（书中页码 第14章相应小节）

\InfoLabel{匹配依据} 12、14 级考``从慢病防控角度谈控烟策略''，14 级提示``控烟策略 WHO 有一个，要按那个答''（指
WHO MPOWER 系列控烟措施）。

\ContentBand{知识点原文摘取}

吸烟（包括主动吸烟、被动吸烟及其他形式的烟草暴露）是导致多种慢性病（肺癌、心血管疾病和
COPD）的主要原因。控烟是 WHO 推荐的针对危险因素、在人群水平上的``最合算''（best
buy）干预措施之一，属于慢性病第一级预防中的根本预防（通过控烟法规等政策直接影响健康行为的选择）。

（课本未逐条列出 WHO MPOWER 控烟策略原文；按 WHO《烟草控制框架公约》MPOWER
系列：监测烟草使用与预防政策、保护人们免受烟草烟雾危害、提供戒烟帮助、警示烟草危害、确保禁止烟草广告促销和赞助、提高烟税。建议结合课件作答。）

\ContentBand{对应往年题}

\begin{pastquestion}

\QuestionLabel{【往年题1｜12级流病口试考题回忆.pdf｜简答（题号23）】}\\
结合学过的知识，从慢病防控角度谈控烟策略

\end{pastquestion}

\begin{pastquestion}

\QuestionLabel{【往年题2｜14级流病考题回忆汇总.xls｜简答（题号51）】}\\
控烟策略（李老师说 WHO 有一个，要按照那个答）

\end{pastquestion}

\PointSeparator

\chapter{伤害流行病学}\label{users__fpb__26_spring__pkusph_study__tmp__pdfs__prevention-medicine-past-exams__03-ux6d41ux884cux75c5ux5b66ux5f80ux5e74ux9898ux8003ux70b9ux6574ux7406.md__ux7b2cux5341ux516dux7ae0-ux4f24ux5bb3ux6d41ux884cux75c5ux5b66}

\section{伤害（injury）的定义}\label{users__fpb__26_spring__pkusph_study__tmp__pdfs__prevention-medicine-past-exams__03-ux6d41ux884cux75c5ux5b66ux5f80ux5e74ux9898ux8003ux70b9ux6574ux7406.md__ux4f24ux5bb3injuryux7684ux5b9aux4e49}

\FrequencyPill{5 次}

\InfoLabel{对应小节} 第十五章 伤害流行病学 / 第一节 概述 / 一、伤害的定义（书中页码 第15章相应小节）

\InfoLabel{匹配依据} 10、11、12、14、20 级考``伤害（injury）''名词解释。

\ContentBand{知识点原文摘取}

美国疾病预防控制中心（CDC）给伤害下的定义为：由于运动、热量、化学、电或放射线的能量交换，在机体组织无法耐受的水平上所造成的组织损伤或由窒息而引起的缺氧称为伤害。

当前伤害的可操作性定义为``因为能量（机械能、热能、化学能等）的传递或干扰超过人体的耐受性造成组织损伤，或窒息导致缺氧，影响了正常活动，需要医治或者看护''。我国流行病学界定标准为``经医疗单位诊断为某一类损伤或因损伤请假（休工、休学、休息）一日以上''。

\ContentBand{对应往年题}

\begin{pastquestion}

\QuestionLabel{【往年题1｜20级回忆试题库.xlsx｜名解（题号11、47、28）】}\\
injury（伤害）

\end{pastquestion}

\begin{pastquestion}

\QuestionLabel{【往年题2｜14级流病考题回忆汇总.xls｜名解（题号8、12）】}\\
伤害

\end{pastquestion}

\begin{pastquestion}

\QuestionLabel{【往年题3｜12级流病口试考题回忆.pdf｜名解（题号9）】}\\
伤害

\end{pastquestion}

\PointSeparator

\section{伤害的能量交换（转移）模型与伤害发生的原因及影响因素}\label{users__fpb__26_spring__pkusph_study__tmp__pdfs__prevention-medicine-past-exams__03-ux6d41ux884cux75c5ux5b66ux5f80ux5e74ux9898ux8003ux70b9ux6574ux7406.md__ux4f24ux5bb3ux7684ux80fdux91cfux4ea4ux6362ux8f6cux79fbux6a21ux578bux4e0eux4f24ux5bb3ux53d1ux751fux7684ux539fux56e0ux53caux5f71ux54cdux56e0ux7d20}

\FrequencyPill{4 次}

\InfoLabel{对应小节} 第十五章 伤害流行病学 / 第一节 概述 / 三、伤害发生的原因及影响因素（书中页码
第15章相应小节）

\InfoLabel{匹配依据} 10、11、12、19 级考``根据伤害能量交换（转移）模型，解释伤害发生的原因和影响因素''。

\ContentBand{知识点原文摘取}

所有伤害共享一个致病机制------伤害的能量交换模型。该模型源自传染病流行病学的三角模型，其中传染病发生的动因是病原体，而伤害发生的动因是能量。从该模型出发，伤害发生的原因包括致伤因子、宿主和环境三个方面：

①致伤因子：能量（energy），能量的异常交换或在短时间内暴露于大剂量能量都会导致伤害，常见的致伤能量有动能（机械能）、热能、电能、辐射能、化学能等。\\
②宿主：受伤害的个体，宿主的条件和耐受性影响伤害的发生与否及严重程度；既有个人内在因素（性别、年龄等），也有外在因素（疲劳、酗酒等）。\\
③环境：包括物理环境和社会环境等。

\ContentBand{对应往年题}

\begin{pastquestion}

\QuestionLabel{【往年题1｜11级流行病学口试考题回忆.pdf｜简答】}\\
伤害的能量交换模型，解释伤害发生的病因和影响因素

\end{pastquestion}

\begin{pastquestion}

\QuestionLabel{【往年题2｜19级流行病学期末口试回忆.xlsx｜简答（题号19）】}\\
根据伤害能量交换模型，解释伤害发生的原因和影响因素

\end{pastquestion}

\begin{pastquestion}

\QuestionLabel{【往年题3｜12级流病口试考题回忆.pdf｜简答（题号9、31）】}\\
伤害模型的原因和影响因素

\end{pastquestion}

\PointSeparator

\section{伤害预防的''5E''综合策略}\label{users__fpb__26_spring__pkusph_study__tmp__pdfs__prevention-medicine-past-exams__03-ux6d41ux884cux75c5ux5b66ux5f80ux5e74ux9898ux8003ux70b9ux6574ux7406.md__ux4f24ux5bb3ux9884ux9632ux76845eux7efcux5408ux7b56ux7565}

\FrequencyPill{3 次}

\InfoLabel{对应小节} 第十五章 伤害流行病学 / 第四节 预防策略与措施 /
一、策略（四）``5E''伤害预防综合策略（书中页码 第15章相应小节）

\InfoLabel{匹配依据} 11、12、20 级考``伤害预防的 5E 原则/5E 干预策略及说明''。

\ContentBand{知识点原文摘取}

``5E''伤害预防综合策略是目前国际公认的伤害预防综合策略，包括五个方面：\\
①教育预防策略（education strategy）：提高公众对伤害预防的认识和知识水平，改变危险行为和习惯。\\
②环境改善策略（environmental modification strategy）：通过改善物理环境减少伤害发生的可能性和严重性。\\
③工程策略（engineering strategy）：通过安全设计和技术改进消除或减低环境中的危险因素。\\
④强化执法策略（enforcement strategy）：通过制定和执行法律及公安部门的措施，维持某些行为和规范的实施。\\
⑤评估策略（evaluation strategy）：通过监测和评估衡量哪些干预措施、项目和政策对预防伤害最有效。

\ContentBand{对应往年题}

\begin{pastquestion}

\QuestionLabel{【往年题1｜20级回忆试题库.xlsx｜简答（题号11）】}\\
伤害预防的 5E 原则

\end{pastquestion}

\begin{pastquestion}

\QuestionLabel{【往年题2｜11级流行病学口试考题回忆.pdf｜简答】}\\
伤害 5E 干预策略及说明

\end{pastquestion}

\begin{pastquestion}

\QuestionLabel{【往年题3｜12级流病口试考题回忆.pdf｜简答（题号24）】}\\
5E

\end{pastquestion}

\PointSeparator

\section{Haddon
伤害预防的十大策略（哈顿十策）}\label{users__fpb__26_spring__pkusph_study__tmp__pdfs__prevention-medicine-past-exams__03-ux6d41ux884cux75c5ux5b66ux5f80ux5e74ux9898ux8003ux70b9ux6574ux7406.md__haddon-ux4f24ux5bb3ux9884ux9632ux7684ux5341ux5927ux7b56ux7565ux54c8ux987fux5341ux7b56}

\FrequencyPill{2 次}

\InfoLabel{对应小节} 第十五章 伤害流行病学 / 第四节 预防策略与措施 / 一、策略（三）Haddon
伤害预防的十大策略（书中页码 第215---216页）

\InfoLabel{匹配依据} 10、12 级考``Haddon 十策/哈顿十策/伤害预防十大策略''。

\ContentBand{知识点原文摘取}

Haddon
提出的预防与控制伤害发生和减少死亡的十大策略：①预防危险因素的形成；②降低危险因素（减少或控制其数量/强度，如限车速）；③预防危险因素的释放；④改变危险因素的空间分布（如使用安全带）；⑤将危险因素从时间、空间上与受保护者分开（如行人走人行道、戴安全帽）；⑥用屏障将危险因素与受保护者分开；⑦改变危险因素的基本性质（如车内尖锐件改钝角）；⑧增加人体对危险因素的抵抗力；⑨对已造成的损伤提供有针对性的预防与控制措施（如急救）；⑩对伤者采取有效治疗及康复措施。

\ContentBand{对应往年题}

\begin{pastquestion}

\QuestionLabel{【往年题1｜12级流病口试考题回忆.pdf｜简答（题号54、10）】}\\
哈顿十策；伤害预防十大策略

\end{pastquestion}

\begin{pastquestion}

\QuestionLabel{【往年题2｜10级考题回忆.pdf｜简答】}\\
haddon 十策

\end{pastquestion}

\PointSeparator

\section{伤害的测量指标}\label{users__fpb__26_spring__pkusph_study__tmp__pdfs__prevention-medicine-past-exams__03-ux6d41ux884cux75c5ux5b66ux5f80ux5e74ux9898ux8003ux70b9ux6574ux7406.md__ux4f24ux5bb3ux7684ux6d4bux91cfux6307ux6807}

\FrequencyPill{2 次}

\InfoLabel{对应小节} 第十五章 伤害流行病学 / 第三节 伤害的流行病学研究 / 一、测量指标（书中页码
第15章相应小节）

\InfoLabel{匹配依据} 07 级考``伤害的指标''，结合伤害测量指标内容。

\ContentBand{知识点原文摘取}

伤害发生频率的测量指标包括伤害发生率、伤害死亡率等。伤害发生率指单位时间（通常是年）内伤害发生的人数（或人次数）与同期该人群的平均人口数之比，是进行伤害研究与监测常用的指标。反映伤害造成的损失程度还可以采用潜在减寿年数（PYLL）和伤残调整寿命年（DALY）等测量指标。

\ContentBand{对应往年题}

\begin{pastquestion}

\QuestionLabel{【往年题1｜07级考题回忆.docx｜简答】}\\
伤害的指标

\end{pastquestion}

\PointSeparator

\chapter{分子、遗传流行病学}\label{users__fpb__26_spring__pkusph_study__tmp__pdfs__prevention-medicine-past-exams__03-ux6d41ux884cux75c5ux5b66ux5f80ux5e74ux9898ux8003ux70b9ux6574ux7406.md__ux7b2cux5341ux4e03ux7ae0-ux5206ux5b50ux9057ux4f20ux6d41ux884cux75c5ux5b66}

\section{生物标志（biomarker）的概念与分类}\label{users__fpb__26_spring__pkusph_study__tmp__pdfs__prevention-medicine-past-exams__03-ux6d41ux884cux75c5ux5b66ux5f80ux5e74ux9898ux8003ux70b9ux6574ux7406.md__ux751fux7269ux6807ux5fd7biomarkerux7684ux6982ux5ff5ux4e0eux5206ux7c7b}

\FrequencyPill{7 次}

\InfoLabel{对应小节} 第十七章 分子流行病学 / 第一节 概述 / 一、定义；第二节 生物标志（书中页码
第17章相应小节）

\InfoLabel{匹配依据} 07、10、11、12、14、20、21
级广泛考``生物标志（物）的概念和种类（分类）''及``外暴露标志/内暴露标志/生物有效剂量标志''名词解释，常追问举例。

\ContentBand{知识点原文摘取}

生物标志（biomarker）主要指从暴露到疾病这个连续过程中可测量的、能反映功能或结构变化的细胞、亚细胞、分子水平的物质。

分子流行病学研究中的生物标志总体上可分为三类：\\
①暴露标志（exposure
biomarker）：与疾病或健康状态有关的暴露因素的生物标志，主要包括外暴露标志（暴露因素进入机体之前的标志和剂量）、内暴露剂量标志（被宿主吸收的外源性暴露物质的量）和生物有效剂量标志（最终与靶组织细胞内
DNA 或蛋白质相互作用的外源性物质或其反应产物的含量，如 DNA 加合物）。\\
②效应标志（effect
biomarker）：宿主暴露后产生功能性或结构性变化的生物标志，包括早期生物效应标志、结构和/或功能改变标志、临床疾病标志。\\
③易感性标志（susceptibility
biomarker）：在暴露因素作用下，宿主对疾病发生、发展易感程度的生物标志，可潜在地修饰从暴露到疾病发生及预后的每一步骤。

\ContentBand{对应往年题}

\begin{pastquestion}

\QuestionLabel{【往年题1｜20级回忆试题库.xlsx｜简答（题号14、15）】}\\
什么是生物标志？生物标志的种类有哪些？（追问各类的举例）；生物标志的概念和分类

\end{pastquestion}

\begin{pastquestion}

\QuestionLabel{【往年题2｜11级流行病学口试考题回忆.pdf｜名解/简答】}\\
生物学标志（把分类也说一说）；外暴露标志（external exposure marker）；内暴露标志剂量

\end{pastquestion}

\begin{pastquestion}

\QuestionLabel{【往年题3｜21级流病考题回忆汇总.xlsx｜名解（题号24、44）】}\\
生物标志（物）；外暴露标志物

\end{pastquestion}

\PointSeparator

\section{生物标志的选择原则}\label{users__fpb__26_spring__pkusph_study__tmp__pdfs__prevention-medicine-past-exams__03-ux6d41ux884cux75c5ux5b66ux5f80ux5e74ux9898ux8003ux70b9ux6574ux7406.md__ux751fux7269ux6807ux5fd7ux7684ux9009ux62e9ux539fux5219}

\FrequencyPill{4 次}

\InfoLabel{对应小节} 第十七章 分子流行病学 / 第三节 主要研究方法 /
三、生物标志的选择和检测（一）生物标志的选择（书中页码 第241页）

\InfoLabel{匹配依据} 10、12、14、20 级考``生物标志（物）的选择原则/标准''。

\ContentBand{知识点原文摘取}

在分子流行病学研究中，应根据研究目的的不同选择相应的生物标志：研究暴露及其水平与疾病的关系选暴露标志；探讨暴露引起的生物学效应选效应标志；研究易感性选易感性标志；综合研究选多项（类）生物标志。生物标志的选择还需符合下列原则：

（1）生物标志应特异、稳定。\\
（2）标本采集、储存方便。\\
（3）检测方法简单、实用，而且操作规范，便于与同类研究结果比较。\\
（4）检测方法的灵敏度和特异度高。

\ContentBand{对应往年题}

\begin{pastquestion}

\QuestionLabel{【往年题1｜12级流病口试考题回忆.pdf｜简答（题号6）】}\\
生物标志的选择原则

\end{pastquestion}

\begin{pastquestion}

\QuestionLabel{【往年题2｜20级回忆试题库.xlsx｜简答（题号45）】}\\
生物标志物选择原则

\end{pastquestion}

\begin{pastquestion}

\QuestionLabel{【往年题3｜14级流病考题回忆汇总.xls｜简答】}\\
生物标志的选择原则

\end{pastquestion}

\PointSeparator

\section{分子流行病学在慢性非传染性疾病（NCD）防控中的应用}\label{users__fpb__26_spring__pkusph_study__tmp__pdfs__prevention-medicine-past-exams__03-ux6d41ux884cux75c5ux5b66ux5f80ux5e74ux9898ux8003ux70b9ux6574ux7406.md__ux5206ux5b50ux6d41ux884cux75c5ux5b66ux5728ux6162ux6027ux975eux4f20ux67d3ux6027ux75beux75c5ncdux9632ux63a7ux4e2dux7684ux5e94ux7528}

\FrequencyPill{4 次}

\InfoLabel{对应小节} 第十七章 分子流行病学 / 第一节 概述 / 二、发展简史（应用现状）（书中页码 第17章相应小节）

\InfoLabel{匹配依据} 10、12、14、17 级考``分子流行病学在慢性非传染性疾病（NCD）预防（和控制）中的应用''。

\ContentBand{知识点原文摘取}

分子流行病学应用分子生物学技术从分子和基因水平阐明生物标志在人群中的分布及其与疾病/健康的关系和影响因素，可以揭示疾病自然史的分子事件，打开传统流行病学中的``黑匣子''，为病因研究和防治措施评价开辟新的途径。

在慢性非传染性疾病，如肿瘤、心血管疾病、糖尿病等的病因和发病机制以及个体易感性等的研究中，分子流行病学作出了巨大贡献（如全基因组关联分析）。具体可用于：疾病和健康状态相关生物标志的分布及影响因素、人群易感性研究、防治效果评价、预后分析等。

\ContentBand{对应往年题}

\begin{pastquestion}

\QuestionLabel{【往年题1｜17级流病口试回忆.xlsx｜简答（题号4）】}\\
分子流行病学在 NCD 的预防和控制方面的应用

\end{pastquestion}

\begin{pastquestion}

\QuestionLabel{【往年题2｜10级考题回忆.pdf｜简答】}\\
分子流行病学在慢性非传染性疾病预防中的应用

\end{pastquestion}

\begin{pastquestion}

\QuestionLabel{【往年题3｜14级流病考题回忆汇总.xls｜简答（题号24）】}\\
分子生物学在慢性非传染性疾病的应用

\end{pastquestion}

\PointSeparator

\section{系统流行病学}\label{users__fpb__26_spring__pkusph_study__tmp__pdfs__prevention-medicine-past-exams__03-ux6d41ux884cux75c5ux5b66ux5f80ux5e74ux9898ux8003ux70b9ux6574ux7406.md__ux7cfbux7edfux6d41ux884cux75c5ux5b66}

\FrequencyPill{2 次}

\InfoLabel{对应小节} 第十七章 分子流行病学 / 第一节 概述 /
二、发展简史（应用现状---系统流行病学正在兴起）（书中页码 第17章相应小节）

\InfoLabel{匹配依据} 10、19 级考``系统流行病学（systems epidemiology）的意义/概念''，19
级标注``考得好偏，系统流行病学是分子流行病学里很偏的内容''。

\ContentBand{知识点原文摘取}

系统流行病学（systems
epidemiology）的概念应运而生，已经逐渐成为一个由流行病学人群观察、干预研究与系统生物学概念交叉形成的新兴学科。这种新的研究模式不再局限于单一的分子标志与疾病的关系，而是将传统流行病学理论和多维度的组学数据（基因组学、表观组学、转录组学、蛋白组学、代谢组学等）相结合，形成一个与疾病或表型相关的数据网络系统，整合遗传易感性、表观遗传改变、基因表达、代谢、肠道微生物等信息，为深入理解疾病发生的内在分子机制提供全面、系统的全新视角。

\ContentBand{对应往年题}

\begin{pastquestion}

\QuestionLabel{【往年题1｜19级流行病学期末口试回忆.xlsx｜名解（题号29）】}\\
系统流行病学（追问：分子流行病学里面很偏的系统流行病学）

\end{pastquestion}

\begin{pastquestion}

\QuestionLabel{【往年题2｜10级考题回忆.pdf｜简答】}\\
系统流行病学意义

\end{pastquestion}

\PointSeparator

\section{遗传度（heritability）}\label{users__fpb__26_spring__pkusph_study__tmp__pdfs__prevention-medicine-past-exams__03-ux6d41ux884cux75c5ux5b66ux5f80ux5e74ux9898ux8003ux70b9ux6574ux7406.md__ux9057ux4f20ux5ea6heritability}

\FrequencyPill{1 次}

\InfoLabel{对应小节} 第十七章 分子流行病学（遗传易感性）；遗传流行病学相关内容（书中页码 第17章相应小节）

\InfoLabel{匹配依据} 07 级名解考``遗传度（heritability）''，追问``遗传度 60\%
说明什么''。属授课内容中相对偏门的遗传流行病学考点。

\ContentBand{知识点原文摘取}

（课本第十七章分子流行病学中讲述了遗传易感性研究------由个体遗传背景差异所导致的不同个体对同一疾病易感程度或治疗反应的强弱，包括候选基因策略和全基因组关联分析策略。课本知识点文件中未检索到``遗传度''的独立成条定义原文，该概念属遗传流行病学，建议结合课件作答：遗传度（heritability）是指在多基因遗传病中，遗传因素所起作用的相对大小，一般以百分数表示；遗传度越高，说明遗传因素在该病发生中的作用越大。如遗传度
60\% 表示该病的发生中遗传因素的贡献约占 60\%，环境因素约占 40\%。）

\ContentBand{对应往年题}

\begin{pastquestion}

\QuestionLabel{【往年题1｜07级考题回忆.docx｜名解】}\\
遗传度（追问：遗传度 60\% 说明什么）

\end{pastquestion}

\PointSeparator

\chapter{研究设计与方案学应用}\label{users__fpb__26_spring__pkusph_study__tmp__pdfs__prevention-medicine-past-exams__03-ux6d41ux884cux75c5ux5b66ux5f80ux5e74ux9898ux8003ux70b9ux6574ux7406.md__ux7b2cux5341ux516bux7ae0-ux7814ux7a76ux8bbeux8ba1ux4e0eux65b9ux6848ux5b66ux5e94ux7528}

\begin{quote}
说明：本章为综合应用类大题，整合第三至六章的研究设计原理，课本无单独对应章节，``对应小节''指向相关方法学章节。
\end{quote}

\section{探索某病危险因素的综合研究设计（工业城市肺癌、高血压危险因素等）}\label{users__fpb__26_spring__pkusph_study__tmp__pdfs__prevention-medicine-past-exams__03-ux6d41ux884cux75c5ux5b66ux5f80ux5e74ux9898ux8003ux70b9ux6574ux7406.md__ux63a2ux7d22ux67d0ux75c5ux5371ux9669ux56e0ux7d20ux7684ux7efcux5408ux7814ux7a76ux8bbeux8ba1ux5de5ux4e1aux57ceux5e02ux80baux764cux9ad8ux8840ux538bux5371ux9669ux56e0ux7d20ux7b49}

\FrequencyPill{8 次}

\InfoLabel{对应小节} 第五章 病例对照研究（设计与实施）；第三章 描述性研究；第四章 队列研究（书中页码
相关方法学章节）

\InfoLabel{匹配依据} 04、05、07、10、12、17、19、20
级反复出现这道经典综合设计大题：``已知吸烟是肺癌的危险因素，现研究表明大气污染、家具污染、职业因素等也可能与肺癌有关，请在某工业城市设计一项流行病学研究探索肺癌的（除吸烟外的）危险因素，写出设计要点''；或``某地从未进行过高血压危险因素调查，设计一个探索性研究''。答题首选病例对照研究（可同时探索多个危险因素、由果及因、适合罕见/潜伏期长的疾病），由于吸烟已知为危险因素，分析阶段需分层分析以控制其混杂。

\ContentBand{知识点原文摘取}

病例对照研究最常被用于疾病病因或危险因素的研究，可以广泛探索病因或危险因素，也可在描述性研究初步形成病因假说的基础上检验病因假说。设计与实施要点包括：确定研究目的与类型→选择病例（统一明确的诊断标准，首选新发病例；来源于医院或社区）与对照（与病例同源、有可比性，必要时匹配）→确定样本量（与
\(P_0\)、OR、\(\alpha\)、\(\beta\)
有关）→确定研究因素（明确暴露定义与测量，纳入可能的混杂因素）→资料收集（可比地收集病例与对照信息，减少回忆/调查偏倚）→分析（计算
OR 及 95\%CI，分层分析控制混杂如吸烟，必要时剂量-反应分析）。

（综合分析题答题应包含：研究类型选择及理由、研究对象的选择、暴露与结局的测量、样本量、资料收集方法、偏倚控制、统计分析与因果推断。）

\ContentBand{对应往年题}

\begin{pastquestion}

\QuestionLabel{【往年题1｜20级回忆试题库.xlsx｜综合分析（题号34、50）】}\\
肺炎/肺癌患病危险因素除吸烟还有大气污染、家具污染、职业因素等，让你在某工业城市进行一项流行病学研究，探索其他危险因素（送分题，可套用实习指导石棉工人那道题）

\end{pastquestion}

\begin{pastquestion}

\QuestionLabel{【往年题2｜19级流行病学期末口试回忆.xlsx｜综合分析（题号16、41）】}\\
实验设计/研究设计：探究一个工业城市肺癌除吸烟以外的危险因素

\end{pastquestion}

\begin{pastquestion}

\QuestionLabel{【往年题3｜05流病考题回忆.doc｜大题】}\\
某工业城市从未进行过肺癌的流行病学调查，请你设计一个研究计划，开展概括病因研究设计并写出设计要点

\end{pastquestion}

\PointSeparator

\section{现况（横断面）研究设计应用（如某地区某时期某病患病率调查）}\label{users__fpb__26_spring__pkusph_study__tmp__pdfs__prevention-medicine-past-exams__03-ux6d41ux884cux75c5ux5b66ux5f80ux5e74ux9898ux8003ux70b9ux6574ux7406.md__ux73b0ux51b5ux6a2aux65adux9762ux7814ux7a76ux8bbeux8ba1ux5e94ux7528ux5982ux67d0ux5730ux533aux67d0ux65f6ux671fux67d0ux75c5ux60a3ux75c5ux7387ux8c03ux67e5}

\FrequencyPill{1 次}

\InfoLabel{对应小节} 第三章 描述性研究 / 第二节 现况研究 / 二、研究设计与实施（书中页码 第32---40页）

\InfoLabel{匹配依据} 21 级综合分析题考``调查某时期某地 18
岁以下人群的结核病患病率，说一下研究设计、研究方法、样本量计算、疾病诊断等''（现况研究/横断面研究设计）。

\ContentBand{知识点原文摘取}

现况研究设计与实施步骤：确定研究目的与类型（普查或抽样调查）→确定研究对象（明确人群特征、地域范围、时间点）→确定样本量和抽样方法（样本量与预期现患率
\(p\)、容许误差 \(d\)、显著性水平 \(\alpha\)
有关；抽样方法包括单纯随机、系统、分层、整群、多阶段抽样）→资料收集（暴露定义和疾病标准明确统一，调查员培训）→整理与分析（计算现患率/患病率，按人群、时间、地区描述分布）。设计中需特别重视抽样调查中研究对象的代表性，应随机抽取足够样本并避免选择偏倚。

\ContentBand{对应往年题}

\begin{pastquestion}

\QuestionLabel{【往年题1｜21级流病考题回忆汇总.xlsx｜综合分析（题号1、19、48）】}\\
某研究者想分析某时期某地区 18
岁以下人群的结核病发病率/患病率，请设计一项研究，从研究类型、调查方法、样本量和诊断标准等方面进行回答（现况研究）

\end{pastquestion}

\PointSeparator

\section{队列研究设计应用（如核废料暴露与恶性肿瘤死亡）}\label{users__fpb__26_spring__pkusph_study__tmp__pdfs__prevention-medicine-past-exams__03-ux6d41ux884cux75c5ux5b66ux5f80ux5e74ux9898ux8003ux70b9ux6574ux7406.md__ux961fux5217ux7814ux7a76ux8bbeux8ba1ux5e94ux7528ux5982ux6838ux5e9fux6599ux66b4ux9732ux4e0eux6076ux6027ux80bfux7624ux6b7bux4ea1}

\FrequencyPill{1 次}

\InfoLabel{对应小节} 第四章 队列研究 / 第二节 研究设计与实施（书中页码 第50---54页）

\InfoLabel{匹配依据} 21 级综合分析题考``某地有一个放置 50
年的核工厂/核废料，政府想了解核废料对当地居民健康的影响（以恶性肿瘤死亡为结局），请设计一个流行病学研究方案''（队列研究）。

\ContentBand{知识点原文摘取}

队列研究设计与实施要点：确定研究因素（暴露的定义与测量，如核废料/辐射暴露）→确定研究结局（明确客观的结局标准，如恶性肿瘤死亡）→确定研究现场与研究人群（暴露人群可选特殊暴露人群；对照可选内对照、外对照或总人口对照）→确定样本量（与对照组发病率
\(p_0\)、两组发病率之差 \(d\)、\(\alpha\)、把握度 \(1-\beta\)
有关，并考虑失访率适当扩大）→资料收集与随访→效应估计（计算 RR、AR
等）。由于辐射等特殊暴露危险一旦认识到即会采取防护，常采用历史性队列研究或双向性队列研究。

\ContentBand{对应往年题}

\begin{pastquestion}

\QuestionLabel{【往年题1｜21级流病考题回忆汇总.xlsx｜综合分析（题号28、46）】}\\
某地有一个露天放置 50
多年的核废料，政府想知道核废料对当地居民健康（以恶性肿瘤死亡为结局指标）的影响，请设计一个合适的研究设计（研究类型、思路、分析指标等）

\end{pastquestion}

\PointSeparator

\chapter{附录：未在授课顺序中明确出现但往年题涉及的考点}\label{users__fpb__26_spring__pkusph_study__tmp__pdfs__prevention-medicine-past-exams__03-ux6d41ux884cux75c5ux5b66ux5f80ux5e74ux9898ux8003ux70b9ux6574ux7406.md__ux9644ux5f55ux672aux5728ux6388ux8bfeux987aux5e8fux4e2dux660eux786eux51faux73b0ux4f46ux5f80ux5e74ux9898ux6d89ux53caux7684ux8003ux70b9}

\begin{quote}
以下考点对应课本中未列入 18 个授课主题的章节（主要是第十六章
突发公共卫生事件流行病学），但在近年往年题中考频较高，需结合老师要求确认是否纳入复习范围。
\end{quote}

\section{突发公共卫生事件的定义和分类}\label{users__fpb__26_spring__pkusph_study__tmp__pdfs__prevention-medicine-past-exams__03-ux6d41ux884cux75c5ux5b66ux5f80ux5e74ux9898ux8003ux70b9ux6574ux7406.md__ux7a81ux53d1ux516cux5171ux536bux751fux4e8bux4ef6ux7684ux5b9aux4e49ux548cux5206ux7c7b}

\FrequencyPill{3 次}

\InfoLabel{对应小节} 第十六章 突发公共卫生事件流行病学（该章未列入 18 个授课主题）（书中页码 第16章相应小节）

\InfoLabel{匹配依据} 19、20、21 级反复考``突发公共卫生事件（emergency public health
events）的定义和分类''名解及简答，追问能否举例其他类型。该章在课本中存在，但不在 prompt
列出的真实授课顺序内，故归入附录，\textbf{需结合老师要求确认是否纳入复习范围}。

\ContentBand{知识点原文摘取}

（课本第十六章为``突发公共卫生事件流行病学''，但该章未列入本课程 18
个授课主题，且未作为本次整理的重点阅读章节，故此处不摘录课本原文。建议结合课件/《突发公共卫生事件应急条例》作答：突发公共卫生事件是指突然发生、造成或者可能造成社会公众健康严重损害的重大传染病疫情、群体性不明原因疾病、重大食物和职业中毒以及其他严重影响公众健康的事件；通常按性质、危害程度、涉及范围分为特别重大（Ⅰ级）、重大（Ⅱ级）、较大（Ⅲ级）和一般（Ⅳ级）四级。）

\ContentBand{对应往年题}

\begin{pastquestion}

\QuestionLabel{【往年题1｜19级流行病学期末口试回忆.xlsx｜简答（题号27、7、25）】}\\
突发公共卫生事件的定义和分类；公共卫生应急事件的定义和分类

\end{pastquestion}

\begin{pastquestion}

\QuestionLabel{【往年题2｜21级流病考题回忆汇总.xlsx｜名解（题号8、41）】}\\
emergency public health events（突发公共卫生事件）

\end{pastquestion}

\begin{pastquestion}

\QuestionLabel{【往年题3｜20级回忆试题库.xlsx｜名解（题号25、26 等）】}\\
emergency public health events（突发公共卫生事件）

\end{pastquestion}

\PointSeparator

\section{治疗指数、ROC
曲线等其他零星考点}\label{users__fpb__26_spring__pkusph_study__tmp__pdfs__prevention-medicine-past-exams__03-ux6d41ux884cux75c5ux5b66ux5f80ux5e74ux9898ux8003ux70b9ux6574ux7406.md__ux6cbbux7597ux6307ux6570roc-ux66f2ux7ebfux7b49ux5176ux4ed6ux96f6ux661fux8003ux70b9}

\FrequencyPill{2 次}

\InfoLabel{对应小节} 第七章 筛检（ROC 曲线，书中页码 第109---110页）；治疗指数（相关章节）

\InfoLabel{匹配依据} 19 级名解考``ROC（曲线）/治疗指数''；20 级名解考``疫苗保护率和治疗指数计算''。其中 ROC
曲线属第七章筛检内容，治疗指数为零星考点。

\ContentBand{知识点原文摘取}

ROC 曲线（receiver operator characteristic curve，受试者工作特征曲线）：将病人和非病人的测量值排序，以
1−特异度为 X 轴、灵敏度为 Y 轴绘制坐标点连成曲线。作用：①选取正确指数最大的点作为最佳截断值；②用 ROC
曲线下面积（AUC）评价筛检方法整体准确度，AUC 越接近 1.0 准确性越高，等于 0.5 时无应用价值。

（``治疗指数''在课本 24 章知识点文件中未检索到明确独立定义原文，建议结合课件作答其概念。）

\ContentBand{对应往年题}

\begin{pastquestion}

\QuestionLabel{【往年题1｜19级流行病学期末口试回忆.xlsx｜名解（题号19、27）】}\\
ROC 曲线（追问横纵坐标）；治疗指数

\end{pastquestion}

\begin{pastquestion}

\QuestionLabel{【往年题2｜20级回忆试题库.xlsx｜综合分析（题号2）】}\\
疫苗保护率和治疗指数计算

\end{pastquestion}

\PointSeparator

\chapter{处理备注}\label{users__fpb__26_spring__pkusph_study__tmp__pdfs__prevention-medicine-past-exams__03-ux6d41ux884cux75c5ux5b66ux5f80ux5e74ux9898ux8003ux70b9ux6574ux7406.md__ux5904ux7406ux5907ux6ce8}

\begin{itemize}
\item
  已读取的课本知识点文件（24 章汇总）：

  \begin{itemize}
  \tightlist
  \item
    01 第一章 绪论
  \item
    02 第二章 疾病的分布
  \item
    03 第三章 描述性研究
  \item
    04 第四章 队列研究
  \item
    05 第五章 病例对照研究
  \item
    06 第六章 实验流行病学
  \item
    07 第七章 筛检
  \item
    08 第八章 病因及其发现和推断
  \item
    09 第九章 预防策略
  \item
    10 第十章 公共卫生监测
  \item
    11 第十一章 传染病流行病学
  \item
    14 第十四章 慢性病流行病学
  \item
    15 第十五章 伤害流行病学
  \item
    17 第十七章 分子流行病学
  \item
    19 第十九章 病毒性肝炎
  \item
    20 第二十章 结核病（耐药结核病、DOTS 相关小节）
  \item
    21 第二十一章 性传播疾病
  \item
    23 第二十三章 流行性感冒（抗原变异等小节）
  \item
    未作为重点阅读的章节：12 医院感染、13 消毒与病媒生物防制、16 突发公共卫生事件、18 循证医学与系统综述、22
    感染性腹泻、24 地方病（其中 16、22 涉及的考点已分别归入``附录''和``第十二章''说明）。
  \end{itemize}
\item
  已读取的往年题文件（7 往年题）：

  \begin{itemize}
  \tightlist
  \item
    04级流病考题回忆.docx（已用 textutil 抽取）
  \item
    05流病考题回忆.doc（已用 textutil 抽取）
  \item
    07级考题回忆.docx（已用 textutil 抽取）
  \item
    10级考题回忆.pdf（已用 pdftotext 抽取，口试题库网格版）
  \item
    11级流行病学口试考题回忆.pdf（已用 pdftotext 抽取，口试题库网格版）
  \item
    12级流病口试考题回忆.pdf（已用 pdftotext 抽取，题号 1--65）
  \item
    14级流病考题回忆汇总.xls（已用 pandas/xlrd 抽取，问卷星版 62 条）
  \item
    17级流病口试回忆.xlsx（已用 openpyxl 抽取，新冠主题）
  \item
    19级流行病学期末口试回忆.xlsx（已用 openpyxl 抽取）
  \item
    20级回忆试题库.xlsx（已用 openpyxl 抽取，原版+英文翻译版）
  \item
    21级流病考题回忆汇总.xlsx（已用 openpyxl 抽取）
  \end{itemize}
\item
  无法读取或内容不完整的文件：无。所有 11 个往年题文件均成功抽取为文本（含 .doc/.xls 旧格式）。PDF（10/11/12
  级）为口试题库网格，pdftotext 抽取后存在列错位，已结合上下文人工校读，个别题干文字可能有错位或漏字。
\item
  关于考频统计：本文件考频按``出现的年级文件数（届数）''计，同一年级文件内的重复回忆合并计 1
  次（详见开头``说明''）。口试回忆题库中同一考点在同一年级内通常被多人重复记录，故``对应往年题''仅摘取代表性题干。最高频考点（如
  RR/AR/OR/灵敏度特异度计算、各级预防、人群易感性、传播途径等）几乎每届必考。
\item
  疑似对应但未能完全确认的考点：

  \begin{itemize}
  \tightlist
  \item
    ``霍桑效应''归入第八章偏倚，但课本原文位于第六章（设立对照的必要性），属跨章对应。
  \item
    ``冰山现象''\,``移民流行病学''等跨章概念已按课本原文所在章归类，并在对应小节注明。
  \end{itemize}
\item
  课本知识点文件中未检索到明确原文的往年题（已按 prompt 要求保留题目、未自行编造课本原文，仅据课件方向说明）：

  \begin{itemize}
  \tightlist
  \item
    错分偏倚（无差异错分/有差异错分）及偏倚大小和方向的计算（第八章）------主要来自课件与实习指导。
  \item
    随时消毒、终末消毒（第十一章）------课本未单列定义。
  \item
    暴发调查的步骤（第十一章）------主要来自实习课。
  \item
    巢式病例对照研究的独立定义（第五章衍生类型）。
  \item
    易感性偏倚、时间效应偏倚（第八章）。
  \item
    WHO 针对流感大流行的具体措施、人感染高致病性禽流感定义（第二十三章）。
  \item
    艾滋病窗口期（第二十一章，课本仅描述抗体 2～3 周阴转阳）。
  \item
    心血管疾病的社区综合防治、控烟策略（MPOWER）（第十四章）。
  \item
    遗传度（第十七章，属遗传流行病学）、治疗指数（附录）。
  \item
    突发公共卫生事件的定义和分类（第十六章，未列入授课主题）。
  \end{itemize}
\end{itemize}

\SetSubjectAccent{3D8A4F}{环境健康学}

\part{环境健康学}

\chapter*{说明}\label{users__fpb__26_spring__pkusph_study__tmp__pdfs__prevention-medicine-past-exams__04-ux73afux5883ux5065ux5eb7ux5b66ux5f80ux5e74ux9898ux8003ux70b9ux6574ux7406.md__ux8bf4ux660e}
\addcontentsline{toc}{chapter}{说明}

\begin{itemize}
\tightlist
\item
  本文件根据《环境健康学------笔记by徐庆松》PDF 和用户提供的往年题文件整理。
\item
  本次将往年题拆分为 280 个题目条目；每个条目均已归入至少一个考点。
\item
  知识点正文均摘自知识总结文件/课件 Markdown/PDF
  抽取文本原文；未找到明确原文的题目已保留并标记为``未完全匹配到''。
\item
  往年题均摘自用户提供的往年题文件。
\item
  排序依据为：先按知识总结文件/课件章节顺序，再在同一章节内按考频从高到低排列；同频时按原文出现顺序排列。
\item
  一道考题若同时考查多个知识点，可以分别归入多个考点，并各计 1 次。
\item
  如果某道题未在知识总结文件/课件中完全匹配到对应知识点，则标记为``未完全匹配到''，并根据题意暂归入最合适章节。
\item
  已反向核对题库拆分条目，确保最终文件中没有题目条目遗漏。
\end{itemize}

\PointSeparator

\chapter{绪论}\label{users__fpb__26_spring__pkusph_study__tmp__pdfs__prevention-medicine-past-exams__04-ux73afux5883ux5065ux5eb7ux5b66ux5f80ux5e74ux9898ux8003ux70b9ux6574ux7406.md__ux7b2cux4e00ux7ae0-ux7eeaux8bba}

\section{环境暴露、暴露测量、暴露重建与环境流行病学（未完全匹配到完整原文）}\label{users__fpb__26_spring__pkusph_study__tmp__pdfs__prevention-medicine-past-exams__04-ux73afux5883ux5065ux5eb7ux5b66ux5f80ux5e74ux9898ux8003ux70b9ux6574ux7406.md__ux73afux5883ux66b4ux9732ux66b4ux9732ux6d4bux91cfux66b4ux9732ux91cdux5efaux4e0eux73afux5883ux6d41ux884cux75c5ux5b66ux672aux5b8cux5168ux5339ux914dux5230ux5b8cux6574ux539fux6587}

\InfoLabel{考频} 20 次

\InfoLabel{对应小节} 未完全匹配到完整原文；根据题意暂归入：第一章 绪论，并参考第十二章/第十三章相关片段

\InfoLabel{匹配依据}
题干集中于环境暴露全过程、时间-活动模式、暴露重建、暴露组学、定组研究、环境流行病学内容特点等。资料中仅有暴露效应关系、暴露评价和环境流行病学一句定义，未找到这些术语的完整原文。

\ContentBand{知识点原文摘取}

未完全匹配到完整原文。\\
可参考原文片段：\\
暴露评估（Exposure Assessment）：人们每天通过各种途径（如饮食、呼吸等）摄入的这种物质的量是多少？\\
大气污染对健康影响的人群健康调查中，暴露评价包括：大气监测资料、调查问卷、个体暴露测定、生物材料监测。\\
环境流行病学是研究疾病在人群中的分布以及影响和决定这种分布的环境因素的科学。

\InfoLabel{知识来源}
《环境健康学------笔记by徐庆松》第25页；《环境健康学------笔记by徐庆松》第94页；《环境健康学------笔记by徐庆松》第95页

\ContentBand{对应往年题}

\begin{pastquestion}

\QuestionLabel{【往年题1｜03级环境卫生重点考题总结-曾.doc｜选择】}\\
评价环境暴露

\end{pastquestion}

\begin{pastquestion}

\QuestionLabel{【往年题2｜03级环境卫生重点考题总结-曾.doc｜选择】}\\
环境流病要点

\end{pastquestion}

\begin{pastquestion}

\QuestionLabel{【往年题3｜03级环境卫生重点考题总结-曾.doc｜大题】}\\
环境暴露的完整定义（4）

\end{pastquestion}

\begin{pastquestion}

\QuestionLabel{【往年题4｜03级环境卫生重点考题总结-曾.doc｜口试题】}\\
暴露定义，环境暴露全过程

\end{pastquestion}

\begin{pastquestion}

\QuestionLabel{【往年题5｜03级环境卫生重点考题总结-曾.doc｜口试题】}\\
环境流行病学设计原则和注意事项

\end{pastquestion}

\begin{pastquestion}

\QuestionLabel{【往年题6｜03级环境卫生重点考题总结-曾.doc｜口试题】}\\
环境流行病学的内容特点

\end{pastquestion}

\begin{pastquestion}

\QuestionLabel{【往年题7｜04级环卫考题回忆.doc｜名词解释】}\\
environmental response gene

\end{pastquestion}

\begin{pastquestion}

\QuestionLabel{【往年题8｜04级环卫考题回忆.doc｜名词解释】}\\
时间-活动模式

\end{pastquestion}

\begin{pastquestion}

\QuestionLabel{【往年题9｜04级环卫考题回忆.doc｜简答题】}\\
简述环境流行病学研究的主要内容

\end{pastquestion}

\begin{pastquestion}

\QuestionLabel{【往年题10｜04级环卫考题回忆.doc｜选择】}\\
环境流行病学定组研究一般是指：在不同的时间点对同一组研究对象进行连续的观测并进行分析的方法

\end{pastquestion}

\begin{pastquestion}

\QuestionLabel{【往年题11｜04级环卫考题回忆.doc｜选择】}\\
接触点测量和暴露情景模拟？

\end{pastquestion}

\begin{pastquestion}

\QuestionLabel{【往年题12｜08级环卫考题回忆.docx｜名词解释】}\\
时间-活动模式

\end{pastquestion}

\begin{pastquestion}

\QuestionLabel{【往年题13｜08级环卫考题回忆.docx｜简答】}\\
暴露组学的含义

\end{pastquestion}

\begin{pastquestion}

\QuestionLabel{【往年题14｜10级环卫考题回忆.docx｜名词解释】}\\
暴露重建

\end{pastquestion}

\begin{pastquestion}

\QuestionLabel{【往年题15｜10级环卫考题回忆.docx｜简答】}\\
环境暴露的定义及内容

\end{pastquestion}

\begin{pastquestion}

\QuestionLabel{【往年题16｜12级环境卫生学考试题回忆版（不全）.docx｜名词解释】}\\
情景评估（英文，潘老师课件）

\end{pastquestion}

\begin{pastquestion}

\QuestionLabel{【往年题17｜14级环卫考题回忆by 小川川.pdf｜名词解释】}\\
Exposure reconstruction

\end{pastquestion}

\begin{pastquestion}

\QuestionLabel{【往年题18｜14级环卫考题回忆by 小川川.pdf｜简答】}\\
环境暴露测量的要素有哪些，分别有什么意义

\end{pastquestion}

\begin{pastquestion}

\QuestionLabel{【往年题19｜2021级环境健康学考题yys.docx｜名词解释】}\\
生物有效剂量

\end{pastquestion}

\begin{pastquestion}

\QuestionLabel{【往年题20｜21级本科预防二班+胡润泽+环境卫生+主观题总结.docx｜名词解释】}\\
生物有效剂量

\end{pastquestion}

\PointSeparator

\section{人与环境的统一关系、原生环境与次生环境}\label{users__fpb__26_spring__pkusph_study__tmp__pdfs__prevention-medicine-past-exams__04-ux73afux5883ux5065ux5eb7ux5b66ux5f80ux5e74ux9898ux8003ux70b9ux6574ux7406.md__ux4ebaux4e0eux73afux5883ux7684ux7edfux4e00ux5173ux7cfbux539fux751fux73afux5883ux4e0eux6b21ux751fux73afux5883}

\InfoLabel{考频} 5 次

\InfoLabel{对应小节} 二、环境与健康 / （一）人与环境的统一关系；（二）环境的分类；（八）原生环境的健康问题

\InfoLabel{匹配依据}
题干关键词为``人与环境统一关系''\,``原生环境''\,``次生环境''\,``气候/适应''，与绪论中人与环境关系、环境分类和气象因素定义相对应。

\ContentBand{知识点原文摘取}

人与环境在物质上的统一性：机体通过新陈代谢与环境不断进行着物质、能量和信息的交换和转移，使机体与周围环境之间保持着动态平衡。\\
机体对环境的适应性：机体随外界环境条件的改变而改变自身的特性或生活方式。\\
机体与环境的相互作用：人类的健康、疾病、寿命都是外界环境因素与机体内在因素相互作用的结果。\\
按照是否受到人类活动的影响，环境又可分为原生环境(primary environment)和次生环境(secondary
environment)。原生环境：天然形成的基本上未受人为活动影响的环境；次生环境：受到人为活动影响的环境。\\
气象因素是由大气的气温、气湿、气流、气压等多种要素组成。在一定地区在一定时间内多种气象要素的综合状态称为天气(weather)。某地区长期天气变化情况的总和则称作气候(climate)。

\InfoLabel{知识来源}
《环境健康学------笔记by徐庆松》第3页；《环境健康学------笔记by徐庆松》第4页；《环境健康学------笔记by徐庆松》第6页

\ContentBand{对应往年题}

\begin{pastquestion}

\QuestionLabel{【往年题1｜08级环卫考题回忆.docx｜名词解释】}\\
原生环境

\end{pastquestion}

\begin{pastquestion}

\QuestionLabel{【往年题2｜12级环境卫生学考试题回忆版（不全）.docx｜简答】}\\
简要说明原生环境与次生环境的区别

\end{pastquestion}

\begin{pastquestion}

\QuestionLabel{【往年题3｜12级环境卫生学考试题回忆版（不全）.docx｜应用】}\\
举例说明人与环境的统一关系

\end{pastquestion}

\begin{pastquestion}

\QuestionLabel{【往年题4｜18级环境健康期末考题回忆by崔浩亮.docx｜名词解释】}\\
气候

\end{pastquestion}

\begin{pastquestion}

\QuestionLabel{【往年题5｜18级环境健康期末考题回忆by崔浩亮.docx｜简答题】}\\
简述气候实应。

\end{pastquestion}

\PointSeparator

\section{健康效应谱、易感人群与环境污染健康危害特点}\label{users__fpb__26_spring__pkusph_study__tmp__pdfs__prevention-medicine-past-exams__04-ux73afux5883ux5065ux5eb7ux5b66ux5f80ux5e74ux9898ux8003ux70b9ux6574ux7406.md__ux5065ux5eb7ux6548ux5e94ux8c31ux6613ux611fux4ebaux7fa4ux4e0eux73afux5883ux6c61ux67d3ux5065ux5eb7ux5371ux5bb3ux7279ux70b9}

\InfoLabel{考频} 5 次

\InfoLabel{对应小节} 二、环境与健康 / （六）健康效应谱；（七）易感人群；（九）环境污染

\InfoLabel{匹配依据}
题干涉及高危/易感人群、健康效应指标选择、环境污染物健康影响特点和长期影响，匹配绪论中的健康效应谱、易感人群和污染危害特点。

\ContentBand{知识点原文摘取}

健康效应谱：环境因素对机体的效应是一个连续的多个阶段的过程。\\
在人群中，由于个体实际的暴露水平和暴露时间存在着差异，以及易感人群的存在，同样的外环境暴露水平下，人群中的不同个体可出现不同的反应，结果表现为不同效应在人群中有不同的发生频率。这种不同水平的效应在人群中的分布称之为健康效应谱。\\
易感人群：由于易感性(年龄、性别、生理状况、健康状况、遗传因素等)的差异，有少数人可能出现病理改变甚至死亡，这类易受环境损伤的人群称为敏感人群或易感人群。\\
环境污染健康危害的特点：环境污染物可通过大气、土壤、水和食物等多种介质进入人体产生危害，其作用对象是整个人群；一般生活环境中的环境污染物的水平很低，但人群长期生活在这样的环境中，累积暴露量大，会出现慢性中毒。

\InfoLabel{知识来源}
《环境健康学------笔记by徐庆松》第6页；《环境健康学------笔记by徐庆松》第7页；《环境健康学------笔记by徐庆松》第8页

\ContentBand{对应往年题}

\begin{pastquestion}

\QuestionLabel{【往年题1｜03级环境卫生重点考题总结-曾.doc｜名词解释】}\\
高危人群

\end{pastquestion}

\begin{pastquestion}

\QuestionLabel{【往年题2｜03级环境卫生重点考题总结-曾.doc｜口试题】}\\
环境污染物对人群健康影响的特点

\end{pastquestion}

\begin{pastquestion}

\QuestionLabel{【往年题3｜03级环境卫生重点考题总结-曾.doc｜口试题】}\\
高危人群的概念及其主要特征

\end{pastquestion}

\begin{pastquestion}

\QuestionLabel{【往年题4｜03级环境卫生重点考题总结-曾.doc｜口试题】}\\
效应指标及其选择

\end{pastquestion}

\begin{pastquestion}

\QuestionLabel{【往年题5｜21级本科预防二班+胡润泽+环境卫生+主观题总结.docx｜名词解释】}\\
易感人群

\end{pastquestion}

\PointSeparator

\section{暴露效应关系、暴露反应关系与剂量反应}\label{users__fpb__26_spring__pkusph_study__tmp__pdfs__prevention-medicine-past-exams__04-ux73afux5883ux5065ux5eb7ux5b66ux5f80ux5e74ux9898ux8003ux70b9ux6574ux7406.md__ux66b4ux9732ux6548ux5e94ux5173ux7cfbux66b4ux9732ux53cdux5e94ux5173ux7cfbux4e0eux5242ux91cfux53cdux5e94}

\InfoLabel{考频} 3 次

\InfoLabel{对应小节} 二、环境与健康 / （五）暴露效应关系；第十二章风险评估 / 建立剂量反应关系

\InfoLabel{匹配依据}
题干要求区分剂量-效应、剂量-反应或考查剂量反应关系，绪论中有暴露效应/暴露反应关系，第十二章进一步说明剂量反应评估。

\ContentBand{知识点原文摘取}

环境因素的暴露不同，会使机体产生不同的效应。效应可以从轻微的生理或生化改变到严重的疾病甚至死亡。暴露量越大，效应越严重，环境因素的暴露与机体所呈现出的生物效应强度间的关系，称为暴露效应关系。\\
在某一群体中，相同暴露量的环境因素对不同的个体有不同的效应，从无健康损害→代偿性损伤→亚临床状态→疾病→死亡。环境因素的暴露量不同时，各种效应的比例也就相应地改变。这种随暴露量不同在人群中某种效应发生率不同的关系称为暴露反应关系。\\
暴露效应关系和暴露反应关系是制订环境健康和环境质量标准的理论基础。\\
建立剂量反应关系（Dose-response
Assessment）：如果该物质有害，那么人们摄入的剂量与患病的可能性有怎样的关系（剂量反应关系）？

\InfoLabel{知识来源} 《环境健康学------笔记by徐庆松》第6页；《环境健康学------笔记by徐庆松》第94页

\ContentBand{对应往年题}

\begin{pastquestion}

\QuestionLabel{【往年题1｜03级环境卫生重点考题总结-曾.doc｜口试题】}\\
剂量反应关系，剂量效应关系

\end{pastquestion}

\begin{pastquestion}

\QuestionLabel{【往年题2｜08级环卫考题回忆.docx｜简答】}\\
剂量-效应关系与剂量-反应关系之间的区别和联系

\end{pastquestion}

\begin{pastquestion}

\QuestionLabel{【往年题3｜14级环卫考题回忆by 小川川.pdf｜选择回忆】}\\
什么是剂量反应关系

\end{pastquestion}

\PointSeparator

\section{环境污染、公害事件与公害病（未完全匹配到完整原文）}\label{users__fpb__26_spring__pkusph_study__tmp__pdfs__prevention-medicine-past-exams__04-ux73afux5883ux5065ux5eb7ux5b66ux5f80ux5e74ux9898ux8003ux70b9ux6574ux7406.md__ux73afux5883ux6c61ux67d3ux516cux5bb3ux4e8bux4ef6ux4e0eux516cux5bb3ux75c5ux672aux5b8cux5168ux5339ux914dux5230ux5b8cux6574ux539fux6587}

\InfoLabel{考频} 3 次

\InfoLabel{对应小节} 二、环境与健康 / （九）环境污染；第二章 / 大气污染急性危害

\InfoLabel{匹配依据}
题干涉及公害病定义、国外公害事件、米糠油事件等。资料中有环境污染定义、环境污染阶段和若干污染事件，但未完整给出``公害病''定义。

\ContentBand{知识点原文摘取}

环境污染：污染物进入环境后，对人群的机体和精神状态产生了直接的、间接的或者潜在的有害影响；或者在很大范围内妨碍了各种生物的生活，使环境条件恶化，影响了生态平衡，称为环境污染(environmental
pollution)。\\
世界范围环境污染的几个阶段：发生期；发展期；泛滥期；转移和扩散期；全球影响期。\\
环境污染物引起的急性中毒以大气污染事件比较多见。\\
生产事故：印度博帕尔毒气泄漏事件；前苏联切尔诺贝利核电站爆炸事件；重庆市开县特大天然气井喷事件；日本福岛核泄漏事件；天津港``8·12''火灾爆炸事件。

\InfoLabel{知识来源}
《环境健康学------笔记by徐庆松》第7页；《环境健康学------笔记by徐庆松》第8页；《环境健康学------笔记by徐庆松》第19页

\ContentBand{对应往年题}

\begin{pastquestion}

\QuestionLabel{【往年题1｜10级环卫考题回忆.docx｜大题】}\\
就公害病的定义和国外公害病的历史事件谈谈国内目前公害病的现状和特点

\end{pastquestion}

\begin{pastquestion}

\QuestionLabel{【往年题2｜12级环境卫生学考试题回忆版（不全）.docx｜名词解释】}\\
环境污染

\end{pastquestion}

\begin{pastquestion}

\QuestionLabel{【往年题3｜20级完整试题回忆.docx｜单选题3】}\\
3、下列哪项不属于历史上著名的大气污染公害事件\\
A、多诺拉事件\\
B、洛杉矶光化学烟雾事件\\
C、四日市哮喘事件\\
D、米糠油事件\\
E、马斯河谷事件

\end{pastquestion}

\PointSeparator

\section{环境内分泌干扰物与污染物联合作用}\label{users__fpb__26_spring__pkusph_study__tmp__pdfs__prevention-medicine-past-exams__04-ux73afux5883ux5065ux5eb7ux5b66ux5f80ux5e74ux9898ux8003ux70b9ux6574ux7406.md__ux73afux5883ux5185ux5206ux6cccux5e72ux6270ux7269ux4e0eux6c61ux67d3ux7269ux8054ux5408ux4f5cux7528}

\InfoLabel{考频} 3 次

\InfoLabel{对应小节} 二、环境与健康 / （九）环境污染；5、环境内分泌干扰物

\InfoLabel{匹配依据} 题干涉及 EED、二噁英干扰内分泌以及联合作用类型，匹配绪论和二噁英章节中的原文。

\ContentBand{知识点原文摘取}

环境中的污染物种类很多，可同时进入人体，产生联合作用。环境污染物的联合作用可表现为相加作用、增强作用、拮抗作用或无关作用。\\
改变健康生物及其子孙或者其群体的内分泌功能，对它们的健康产生不良影响的外源性物质或混合物，称为环境内分泌干扰物(environmental
endocrine disruptors, EEDs)。\\
二噁英类是环境内分泌干扰物的代表。

\InfoLabel{知识来源} 《环境健康学------笔记by徐庆松》第8页；《环境健康学------笔记by徐庆松》第25页

\ContentBand{对应往年题}

\begin{pastquestion}

\QuestionLabel{【往年题1｜03级环境卫生重点考题总结-曾.doc｜口试题】}\\
环境内分泌干扰物

\end{pastquestion}

\begin{pastquestion}

\QuestionLabel{【往年题2｜11级环卫考题回忆.docx｜名词解释】}\\
EED

\end{pastquestion}

\begin{pastquestion}

\QuestionLabel{【往年题3｜20级完整试题回忆.docx｜单选题2】}\\
2、以下是联合作用的类型，哪项除外\\
A、独立作用\\
B、相加作用\\
C、拮抗作用\\
D、相乘作用\\
E、协同作用

\end{pastquestion}

\PointSeparator

\section{环境健康学的定义、研究对象与研究内容}\label{users__fpb__26_spring__pkusph_study__tmp__pdfs__prevention-medicine-past-exams__04-ux73afux5883ux5065ux5eb7ux5b66ux5f80ux5e74ux9898ux8003ux70b9ux6574ux7406.md__ux73afux5883ux5065ux5eb7ux5b66ux7684ux5b9aux4e49ux7814ux7a76ux5bf9ux8c61ux4e0eux7814ux7a76ux5185ux5bb9}

\InfoLabel{考频} 2 次

\InfoLabel{对应小节} 一、定义 / （一）环境健康学；（三）环境健康学的研究对象；（四）环境健康学的内容

\InfoLabel{匹配依据}
题干直接涉及``环境健康学研究内容''\,``环境健康学和环境健康工作''等概念性内容，匹配绪论中对学科定义、对象和内容的表述。

\ContentBand{知识点原文摘取}

环境健康学：是研究环境中的物理、化学、生物、社会以及心理社会因素与人体健康，包括生活质量的关系，揭示环境因素对健康影响的发生、发展规律，为充分利用对人群健康有利的环境因素，消除和改善不利的环境因素提出卫生要求和预防措施的学科。\\
环境健康学的研究对象：人类及其周围的环境。\\
环境健康学的内容：大气、水体、土壤与健康；饮水卫生与健康；住宅及室内环境与健康；公共场所卫生人居环境与健康；家用化学物品、个人用品与健康；环境质量和健康危险度评价；环境与健康管理体系；灾害卫生；全球环境变化与健康。

\InfoLabel{知识来源} 《环境健康学------笔记by徐庆松》第2页；《环境健康学------笔记by徐庆松》第3页

\ContentBand{对应往年题}

\begin{pastquestion}

\QuestionLabel{【往年题1｜14级环卫考题回忆by 小川川.pdf｜选择回忆】}\\
环境卫生学研究的主要内容

\end{pastquestion}

\begin{pastquestion}

\QuestionLabel{【往年题2｜14级环卫考题回忆by 小川川.pdf｜简答】}\\
环境健康学和环境健康工作的区别和联系

\end{pastquestion}

\PointSeparator

\section{环境自净、二次污染物、迁移转化与食物链}\label{users__fpb__26_spring__pkusph_study__tmp__pdfs__prevention-medicine-past-exams__04-ux73afux5883ux5065ux5eb7ux5b66ux5f80ux5e74ux9898ux8003ux70b9ux6574ux7406.md__ux73afux5883ux81eaux51c0ux4e8cux6b21ux6c61ux67d3ux7269ux8fc1ux79fbux8f6cux5316ux4e0eux98dfux7269ux94fe}

\InfoLabel{考频} 2 次

\InfoLabel{对应小节} 二、环境与健康 / （三）环境的自净作用；（四）二次污染物或次生污染物

\InfoLabel{匹配依据}
题干涉及自净作用、二次污染物、食物链及迁移转化，均对应绪论中环境自净和环境化学物迁移转化内容。

\ContentBand{知识点原文摘取}

自然界依靠自身的能力，而不是依靠人为的力量，将一些有害因子消除到无害程度，这种作用称为环境的自净作用(self
purification)。\\
环境中的化学污染物可通过多种途径在环境中迁移、转化。一些污染物在环境中由于物理、化学、生物学的作用形成与原来污染物的理化性质和毒性不同的新型污染物，称为二次污染物。\\
迁移：单一介质内的迁移；不同介质间的迁移；生物性转移：通过食物链、食物网；生物放大作用（营养级越高化学物质浓度越高）。\\
转化：化学转化；生物转化。

\InfoLabel{知识来源} 《环境健康学------笔记by徐庆松》第4页；《环境健康学------笔记by徐庆松》第5页

\ContentBand{对应往年题}

\begin{pastquestion}

\QuestionLabel{【往年题1｜03级环境卫生重点考题总结-曾.doc｜口试题】}\\
食物链及其卫生学意义

\end{pastquestion}

\begin{pastquestion}

\QuestionLabel{【往年题2｜04级环卫考题回忆.doc｜名词解释】}\\
自净作用

\end{pastquestion}

\PointSeparator

\section{环境健康标准与基准（未完全匹配到）}\label{users__fpb__26_spring__pkusph_study__tmp__pdfs__prevention-medicine-past-exams__04-ux73afux5883ux5065ux5eb7ux5b66ux5f80ux5e74ux9898ux8003ux70b9ux6574ux7406.md__ux73afux5883ux5065ux5eb7ux6807ux51c6ux4e0eux57faux51c6ux672aux5b8cux5168ux5339ux914dux5230}

\InfoLabel{考频} 1 次

\InfoLabel{对应小节} 未完全匹配到原文；根据题意暂归入：第一章 绪论/环境与健康

\InfoLabel{匹配依据}
题干直接问环境健康标准和基准的区别联系。资料中未找到``环境健康基准''完整定义，仅有标准制定理论基础和环境质量评价标准相关表述。

\ContentBand{知识点原文摘取}

未完全匹配到原文。\\
可参考原文片段：暴露效应关系和暴露反应关系是制订环境健康和环境质量标准的理论基础。

\InfoLabel{知识来源} 《环境健康学------笔记by徐庆松》第6页

\ContentBand{对应往年题}

\begin{pastquestion}

\QuestionLabel{【往年题1｜20级完整试题回忆.docx｜简答题】}\\
试述环境健康标准和基准的区别和联系

\end{pastquestion}

\PointSeparator

\section{题干信息过少无法判断的条目（未完全匹配到）}\label{users__fpb__26_spring__pkusph_study__tmp__pdfs__prevention-medicine-past-exams__04-ux73afux5883ux5065ux5eb7ux5b66ux5f80ux5e74ux9898ux8003ux70b9ux6574ux7406.md__ux9898ux5e72ux4fe1ux606fux8fc7ux5c11ux65e0ux6cd5ux5224ux65adux7684ux6761ux76eeux672aux5b8cux5168ux5339ux914dux5230}

\InfoLabel{考频} 1 次

\InfoLabel{对应小节} 未完全匹配到原文；根据题库原文暂存于第一章

\InfoLabel{匹配依据} 题库中只有``来源？''等无法判明具体知识点的残缺题干，不能删除，故单列保留。

\ContentBand{知识点原文摘取}

未完全匹配到原文。

\InfoLabel{知识来源} 未完全匹配到

\ContentBand{对应往年题}

\begin{pastquestion}

\QuestionLabel{【往年题1｜04级环卫考题回忆.doc｜选择】}\\
来源？

\end{pastquestion}

\PointSeparator

\chapter{大气与健康}\label{users__fpb__26_spring__pkusph_study__tmp__pdfs__prevention-medicine-past-exams__04-ux73afux5883ux5065ux5eb7ux5b66ux5f80ux5e74ux9898ux8003ux70b9ux6574ux7406.md__ux7b2cux4e8cux7ae0-ux5927ux6c14ux4e0eux5065ux5eb7}

\section{风玫瑰图、逆温、大气稳定度、有效排出高度及浓度影响因素}\label{users__fpb__26_spring__pkusph_study__tmp__pdfs__prevention-medicine-past-exams__04-ux73afux5883ux5065ux5eb7ux5b66ux5f80ux5e74ux9898ux8003ux70b9ux6574ux7406.md__ux98ceux73abux7470ux56feux9006ux6e29ux5927ux6c14ux7a33ux5b9aux5ea6ux6709ux6548ux6392ux51faux9ad8ux5ea6ux53caux6d53ux5ea6ux5f71ux54cdux56e0ux7d20}

\InfoLabel{考频} 14 次

\InfoLabel{对应小节} 三、大气污染 / （四）影响大气污染物浓度的因素

\InfoLabel{匹配依据}
题干涉及风向频率图/风玫瑰图、逆温、γ、γd、有效排出高度、大气稳定度、影响大气污染物浓度因素、对流层等。

\ContentBand{知识点原文摘取}

污染物的排放量是决定大气污染程度的最基本的因素。\\
烟囱的有效排出高度(effective height of
emission)，即烟囱本身的高度和烟气抬升高度之和，可以用烟波中心轴到地面的距离表示。在其它条件相同时，排出高度越高，烟波断面越大，污染物的稀释程度就越大，烟波着陆点的浓度就越低。\\
将一定时期内各个风向出现的频率按比例标在罗盘坐标上，可以绘制成风向频率图（风玫瑰图，wind
rose）。风向频率图能够反映某地区一定时期内的主导风向，从而能够指示该地区受某一污染源影响的主要方位。\\
大气温度垂直递减率（γ）：高度每增加 100m 气温下降的度数，通常为 0.65℃。\\
气温下低上高，这种大气温度随着距地面高度的增加而增加的现象称为逆温(temperature inversion)，此时 γ＜0。\\
大气稳定度是指大气垂直运动的程度。

\InfoLabel{知识来源}
《环境健康学------笔记by徐庆松》第10页；《环境健康学------笔记by徐庆松》第14页；《环境健康学------笔记by徐庆松》第15页；《环境健康学------笔记by徐庆松》第16页；《环境健康学------笔记by徐庆松》第17页

\ContentBand{对应往年题}

\begin{pastquestion}

\QuestionLabel{【往年题1｜03级环境卫生重点考题总结-曾.doc｜口试题】}\\
逆温及r、rd

\end{pastquestion}

\begin{pastquestion}

\QuestionLabel{【往年题2｜03级环境卫生重点考题总结-曾.doc｜口试题】}\\
大气垂直递减率

\end{pastquestion}

\begin{pastquestion}

\QuestionLabel{【往年题3｜03级环境卫生重点考题总结-曾.doc｜口试题】}\\
风向频率图定义、卫生学用途

\end{pastquestion}

\begin{pastquestion}

\QuestionLabel{【往年题4｜03级环境卫生重点考题总结-曾.doc｜口试题】}\\
大气稳定度对空气污染影响

\end{pastquestion}

\begin{pastquestion}

\QuestionLabel{【往年题5｜03级环境卫生重点考题总结-曾.doc｜口试题】}\\
有效排放高度

\end{pastquestion}

\begin{pastquestion}

\QuestionLabel{【往年题6｜03级环境卫生重点考题总结-曾.doc｜口试题】}\\
影响大气污染浓度的因素

\end{pastquestion}

\begin{pastquestion}

\QuestionLabel{【往年题7｜10级环卫考题回忆.docx｜选择】}\\
逆温时γ\textless0

\end{pastquestion}

\begin{pastquestion}

\QuestionLabel{【往年题8｜10级环卫考题回忆.docx｜名词解释】}\\
风向频率图

\end{pastquestion}

\begin{pastquestion}

\QuestionLabel{【往年题9｜11级环卫考题回忆.docx｜名词解释】}\\
有效排出高度

\end{pastquestion}

\begin{pastquestion}

\QuestionLabel{【往年题10｜14级环卫考题回忆by 小川川.pdf｜名词解释】}\\
风玫瑰图

\end{pastquestion}

\begin{pastquestion}

\QuestionLabel{【往年题11｜18级环境健康期末考题回忆by崔浩亮.docx｜名词解释】}\\
逆温

\end{pastquestion}

\begin{pastquestion}

\QuestionLabel{【往年题12｜20级完整试题回忆.docx｜单选题4】}\\
4、对大气污染物浓度影响最大的为\\
A、对流层\\
B、平流层\\
C、热成层\\
D、逸散层\\
E、边界层

\end{pastquestion}

\begin{pastquestion}

\QuestionLabel{【往年题13｜20级完整试题回忆.docx｜名词解释】}\\
风向玫瑰图

\end{pastquestion}

\begin{pastquestion}

\QuestionLabel{【往年题14｜21级本科预防二班+胡润泽+环境卫生+主观题总结.docx｜名词解释】}\\
风玫瑰图

\end{pastquestion}

\PointSeparator

\section{煤烟型烟雾事件与光化学烟雾}\label{users__fpb__26_spring__pkusph_study__tmp__pdfs__prevention-medicine-past-exams__04-ux73afux5883ux5065ux5eb7ux5b66ux5f80ux5e74ux9898ux8003ux70b9ux6574ux7406.md__ux7164ux70dfux578bux70dfux96feux4e8bux4ef6ux4e0eux5149ux5316ux5b66ux70dfux96fe}

\InfoLabel{考频} 13 次

\InfoLabel{对应小节} 四、大气污染对健康的影响 / （一）直接危害 / 烟雾事件

\InfoLabel{匹配依据} 题干直接问光化学烟雾名解、形成条件、成分、危害，或比较煤烟型与光化学型烟雾事件。

\ContentBand{知识点原文摘取}

煤烟型烟雾事件：主要由燃煤产生的大量污染物排入大气，在不良气象条件下不能充分扩散所致。引起人群健康危害得到主要大气污染物是烟尘、SO2
以及硫酸雾。\\
光化学烟雾：污染源排入大气中的氮氧化物（NOx）和挥发性有机物（VOC）在太阳紫外线的作用下发生光化学反应所产生的刺激性很强的浅蓝色烟雾所致。属于二次污染物。\\
著名事件：洛杉矶光化学烟雾事件。\\
主要污染源：城市机动车尾气排放。\\
主要成分：臭氧、醛类和各种过氧酰基硝酸酯。称为光化学氧化剂。

\InfoLabel{知识来源} 《环境健康学------笔记by徐庆松》第18页；《环境健康学------笔记by徐庆松》第19页

\ContentBand{对应往年题}

\begin{pastquestion}

\QuestionLabel{【往年题1｜03级环境卫生重点考题总结-曾.doc｜名词解释】}\\
光化学烟雾

\end{pastquestion}

\begin{pastquestion}

\QuestionLabel{【往年题2｜03级环境卫生重点考题总结-曾.doc｜口试题】}\\
``伦敦烟雾事件''和``洛山基烟雾事件''的比较

\end{pastquestion}

\begin{pastquestion}

\QuestionLabel{【往年题3｜03级环境卫生重点考题总结-曾.doc｜口试题】}\\
光化学烟雾产生的条件及对身体的影响

\end{pastquestion}

\begin{pastquestion}

\QuestionLabel{【往年题4｜03级环境卫生重点考题总结-曾.doc｜口试题】}\\
光化学烟雾的成分及危害

\end{pastquestion}

\begin{pastquestion}

\QuestionLabel{【往年题5｜04级环卫考题回忆.doc｜名词解释】}\\
photochemical smog

\end{pastquestion}

\begin{pastquestion}

\QuestionLabel{【往年题6｜04级环卫考题回忆.doc｜选择】}\\
光化学烟雾的特点？

\end{pastquestion}

\begin{pastquestion}

\QuestionLabel{【往年题7｜08级环卫考题回忆.docx｜名词解释】}\\
光化学烟雾（phtochemical smog)

\end{pastquestion}

\begin{pastquestion}

\QuestionLabel{【往年题8｜10级环卫考题回忆.docx｜简答】}\\
光化学烟雾的发生条件，引发的症状？

\end{pastquestion}

\begin{pastquestion}

\QuestionLabel{【往年题9｜11级环卫考题回忆.docx｜简答】}\\
燃煤型和光化学烟雾型烟雾事件的形成和健康危害

\end{pastquestion}

\begin{pastquestion}

\QuestionLabel{【往年题10｜12级环境卫生学考试题回忆版（不全）.docx｜名词解释】}\\
光化学烟雾（英文）

\end{pastquestion}

\begin{pastquestion}

\QuestionLabel{【往年题11｜18级环境健康期末考题回忆by崔浩亮.docx｜名词解释】}\\
Photochemical smog

\end{pastquestion}

\begin{pastquestion}

\QuestionLabel{【往年题12｜20级完整试题回忆.docx｜简答题】}\\
试述煤烟型烟雾事件和光化学型烟雾事件的区别

\end{pastquestion}

\begin{pastquestion}

\QuestionLabel{【往年题13｜21级本科预防二班+胡润泽+环境卫生+主观题总结.docx｜简答题】}\\
煤烟型烟雾事件和光化学烟雾之间的区别

\end{pastquestion}

\PointSeparator

\section{气溶胶、大气颗粒物、PM2.5/PM10及健康影响}\label{users__fpb__26_spring__pkusph_study__tmp__pdfs__prevention-medicine-past-exams__04-ux73afux5883ux5065ux5eb7ux5b66ux5f80ux5e74ux9898ux8003ux70b9ux6574ux7406.md__ux6c14ux6eb6ux80f6ux5927ux6c14ux9897ux7c92ux7269pm2.5pm10ux53caux5065ux5eb7ux5f71ux54cd}

\InfoLabel{考频} 8 次

\InfoLabel{对应小节} 二、大气的自然特征 / 气溶胶；三、大气污染 /
大气颗粒物；六、大气中主要污染物对人体健康的影响 / 颗粒物

\InfoLabel{匹配依据} 题干涉及 fine particulate
matter、PM2.5、PM10、气溶胶、颗粒物健康影响和影响颗粒物生物学作用因素。

\ContentBand{知识点原文摘取}

气溶胶：是固体或液体微粒均匀地分散在气体中形成的相对稳定的悬浮体系。\\
气溶胶态的大气污染物，常常被称作大气颗粒物(particulate matter, PM)。\\
PM2.5 的定义：空气动力学直径小于或等于 2.5μm 的大气颗粒物。\\
颗粒物对呼吸系统的影响：大量的颗粒物进入肺部对局部组织有堵塞作用；吸附着有害气体的颗粒物可以刺激或腐蚀肺泡壁；颗粒物还可以增加动物对病原微生物的敏感性。\\
影响颗粒物生物学作用的因素：颗粒物的粒径；颗粒物的成分；呼吸道对颗粒物的清除作用；其他。

\InfoLabel{知识来源}
《环境健康学------笔记by徐庆松》第11页；《环境健康学------笔记by徐庆松》第13页；《环境健康学------笔记by徐庆松》第21页；《环境健康学------笔记by徐庆松》第22页

\ContentBand{对应往年题}

\begin{pastquestion}

\QuestionLabel{【往年题1｜03级环境卫生重点考题总结-曾.doc｜大题】}\\
影像颗粒物生物学作用的因素（10）

\end{pastquestion}

\begin{pastquestion}

\QuestionLabel{【往年题2｜03级环境卫生重点考题总结-曾.doc｜口试题】}\\
PM10的来源及危害

\end{pastquestion}

\begin{pastquestion}

\QuestionLabel{【往年题3｜04级环卫考题回忆.doc｜名词解释】}\\
细颗粒物

\end{pastquestion}

\begin{pastquestion}

\QuestionLabel{【往年题4｜10级环卫考题回忆.docx｜名词解释】}\\
气溶胶

\end{pastquestion}

\begin{pastquestion}

\QuestionLabel{【往年题5｜11级环卫考题回忆.docx｜名词解释】}\\
PM\_(2.5)

\end{pastquestion}

\begin{pastquestion}

\QuestionLabel{【往年题6｜12级环境卫生学考试题回忆版（不全）.docx｜简答】}\\
大气颗粒物的健康影响

\end{pastquestion}

\begin{pastquestion}

\QuestionLabel{【往年题7｜14级环卫考题回忆by 小川川.pdf｜简答】}\\
大气颗粒物对人体健康的影响有哪些

\end{pastquestion}

\begin{pastquestion}

\QuestionLabel{【往年题8｜20级完整试题回忆.docx｜名词解释】}\\
fine particulate matter

\end{pastquestion}

\PointSeparator

\section{大气污染健康影响调查、监测与调查方案设计}\label{users__fpb__26_spring__pkusph_study__tmp__pdfs__prevention-medicine-past-exams__04-ux73afux5883ux5065ux5eb7ux5b66ux5f80ux5e74ux9898ux8003ux70b9ux6574ux7406.md__ux5927ux6c14ux6c61ux67d3ux5065ux5eb7ux5f71ux54cdux8c03ux67e5ux76d1ux6d4bux4e0eux8c03ux67e5ux65b9ux6848ux8bbeux8ba1}

\InfoLabel{考频} 8 次

\InfoLabel{对应小节} 七、大气污染对健康影响的调查和监测；八、大气污染控制措施

\InfoLabel{匹配依据}
题干要求设计工厂废气、汽车尾气、铅污染或雾霾健康影响调查方案，或问暴露评价内容，匹配本节污染源调查、污染状况监测和人群健康调查。

\ContentBand{知识点原文摘取}

污染源的调查：点源污染，即对一个工厂或一座烟囱对周围大气影响的调查；面源污染，即对整个城市或工业区的大气污染源进行调查；线源污染，主要应调查该线路上交通工具的种类、流量和行驶状态，燃料的种类和燃烧情况，废气的成分等。\\
污染状况监测：采样点选择；采样时间；监测指标；采样记录；监测结果分析与评价。\\
人群健康调查：暴露评价；健康效应测定；资料统计。\\
暴露评价：大气监测资料；调查问卷；个体暴露测定；生物材料监测。\\
健康效应测定：疾病资料；体检；生物材料监测。\\
大气污染控制措施：合理安排工业布局，调整工业结构；完善城市绿化系统；加强居住区内局部污染源的管理；改善能源结构，大力节约能耗；控制机动车尾气污染；改进生产工艺，减少废气排放。

\InfoLabel{知识来源} 《环境健康学------笔记by徐庆松》第25页；《环境健康学------笔记by徐庆松》第26页

\ContentBand{对应往年题}

\begin{pastquestion}

\QuestionLabel{【往年题1｜03级环境卫生重点考题总结-曾.doc｜口试题】}\\
大气污染对健康的影响

\end{pastquestion}

\begin{pastquestion}

\QuestionLabel{【往年题2｜03级环境卫生重点考题总结-曾.doc｜口试题】}\\
大气卫生监测的目的、内容、方法

\end{pastquestion}

\begin{pastquestion}

\QuestionLabel{【往年题3｜04级环卫考题回忆.doc｜问答题】}\\
某地有一座工厂，常年经过一烟囱向周围排放废气。经初步测定，废气的主要成分有烟尘、二氧化硫、苯并芘和一氧化碳等。\\
现拟调查该工厂引起的大气污染对周围居民健康的影响，请回答：\\
如何判断受该污染影响最大的方位？应收集哪些资料？\\
如何调查该污染源对周围居民健康的影响？选择哪些指标？

\end{pastquestion}

\begin{pastquestion}

\QuestionLabel{【往年题4｜08级环卫考题回忆.docx｜应用题】}\\
评价某工厂废气排放情况及对周围居民健康的影

\end{pastquestion}

\begin{pastquestion}

\QuestionLabel{【往年题5｜12级环境卫生学考试题回忆版（不全）.docx｜应用】}\\
工厂废气中铅的对环境及居民健康效应的影响调查方案设计

\end{pastquestion}

\begin{pastquestion}

\QuestionLabel{【往年题6｜14级环卫考题回忆by 小川川.pdf｜应用】}\\
为了解汽车尾气污染对人群的急性健康影响，请设计一个调查计划，至少包含以下几方面：\\
1）研究地点；2）研究人群；3）研究设计；4）研究内容；5）研究方法。

\end{pastquestion}

\begin{pastquestion}

\QuestionLabel{【往年题7｜18级环境健康期末考题回忆by崔浩亮.docx｜论述题】}\\
大气颗粒物是我国大部分城市和地区的首要污染物。自2013年我国实施《大气污染防治行动计划》以来包括颗粒物在内的多数大气污染物水平呈明显下降趋势，但大气PM2.s和0,污染仍是我国突出的环境问题2020年中国生态环境状況公报显示，全国337个地级及以上城市以PM2.s和0，为首要污染物的超标天数分别占总超标天数的51.0\%和37.1\%。针对当前我国大气污染状况，你认为应该采取哪些措施，才能进一步改善空气质量和减少大气污染所致人群健康危害？

\end{pastquestion}

\begin{pastquestion}

\QuestionLabel{【往年题8｜20级完整试题回忆.docx｜单选题6】}\\
6、在大气污染对健康影响的调查和监测中，暴露评价不包括以下哪一项\\
A、疾病诊断\\
B、生物材料监测\\
C、个体暴露测定\\
D、大气监测资料\\
E、调查问卷

\end{pastquestion}

\PointSeparator

\section{SO2、NO2、CO、铅、多环芳烃与二噁英健康影响}\label{users__fpb__26_spring__pkusph_study__tmp__pdfs__prevention-medicine-past-exams__04-ux73afux5883ux5065ux5eb7ux5b66ux5f80ux5e74ux9898ux8003ux70b9ux6574ux7406.md__so2no2coux94c5ux591aux73afux82b3ux70c3ux4e0eux4e8cux5641ux82f1ux5065ux5eb7ux5f71ux54cd}

\InfoLabel{考频} 4 次

\InfoLabel{对应小节} 六、大气中主要污染物对人体健康的影响

\InfoLabel{匹配依据} 题干涉及 SO2、NO2、CO、BaP、二噁英等具体大气污染物的来源或健康危害。

\ContentBand{知识点原文摘取}

SO2 是水溶性的刺激性气体，易被上呼吸道和支气管黏膜的富水性黏液所吸收；SO2 还有促癌作用，可增强
BaP（苯并芘）的致癌作用。\\
NO2 的毒性比 NO 高 4～5 倍；NO2
较难溶于水，故对上呼吸道和眼睛的刺激作用较小，主要作用于深部呼吸道、细支气管及肺泡。\\
一氧化碳（carbon monoxide，CO）是含碳物质不完全燃烧的产物，无色、无臭、无刺激性。CO
暴露与人群心血管疾病的发病率和死亡率增加有关。\\
多环芳烃有致癌性。\\
二噁英类是环境内分泌干扰物的代表；二噁英类有明显的免疫毒性；二噁英类还能引起皮肤损害；有极强的致癌性。

\InfoLabel{知识来源}
《环境健康学------笔记by徐庆松》第23页；《环境健康学------笔记by徐庆松》第24页；《环境健康学------笔记by徐庆松》第25页

\ContentBand{对应往年题}

\begin{pastquestion}

\QuestionLabel{【往年题1｜03级环境卫生重点考题总结-曾.doc｜口试题】}\\
SO2的主要来源及其对健康的影响

\end{pastquestion}

\begin{pastquestion}

\QuestionLabel{【往年题2｜03级环境卫生重点考题总结-曾.doc｜口试题】}\\
NO2的主要来源及其健康危害

\end{pastquestion}

\begin{pastquestion}

\QuestionLabel{【往年题3｜03级环境卫生重点考题总结-曾.doc｜口试题】}\\
CO的知识

\end{pastquestion}

\begin{pastquestion}

\QuestionLabel{【往年题4｜20级完整试题回忆.docx｜单选题5】}\\
5、二噁英类产生广泛的健康影响主要是通过干扰机体的（）\\
A、心血管功能\\
B、呼吸功能\\
C、内分泌功能\\
D、肝肾功能\\
E、免疫功能

\end{pastquestion}

\PointSeparator

\section{大气质量标准、AQI/API与空气质量指数（部分未完全匹配）}\label{users__fpb__26_spring__pkusph_study__tmp__pdfs__prevention-medicine-past-exams__04-ux73afux5883ux5065ux5eb7ux5b66ux5f80ux5e74ux9898ux8003ux70b9ux6574ux7406.md__ux5927ux6c14ux8d28ux91cfux6807ux51c6aqiapiux4e0eux7a7aux6c14ux8d28ux91cfux6307ux6570ux90e8ux5206ux672aux5b8cux5168ux5339ux914d}

\InfoLabel{考频} 4 次

\InfoLabel{对应小节} 五、我国的大气质量标准；第十三章 / 大气质量评价

\InfoLabel{匹配依据} 题干涉及 API、AQI、空气质量指数、日平均浓度和一次最高容许浓度。资料中有大气质量标准和 AQI
缺点，但 API 分级及日平均/一次最高容许浓度未完整匹配。

\ContentBand{知识点原文摘取}

大气环境质量标准是大气中有害物质的法定最高限值，是控制大气污染、保护居民健康、保护生态环境、评价污染程度、制订防护措施的法定依据。\\
现行的《环境空气质量标准》(GB 3095---2012)将我国全国范围分为两类不同的环境空气质功能区。\\
空气质量指数(AQI)的主要缺点：最后对各污染物指数没有综合，仅给出最大分指数的污染物及其污染指数，不能反映其它污染物的污染状况；淡化了其它污染物对大气环境质量的综合影响。

\InfoLabel{知识来源}
《环境健康学------笔记by徐庆松》第20页；《环境健康学------笔记by徐庆松》第21页；《环境健康学------笔记by徐庆松》第102页

\ContentBand{对应往年题}

\begin{pastquestion}

\QuestionLabel{【往年题1｜03级环境卫生重点考题总结-曾.doc｜选择】}\\
API250（轻重中度污染）

\end{pastquestion}

\begin{pastquestion}

\QuestionLabel{【往年题2｜03级环境卫生重点考题总结-曾.doc｜口试题】}\\
日平均浓度，一次最高容许浓度的定义、目的、意义、举例

\end{pastquestion}

\begin{pastquestion}

\QuestionLabel{【往年题3｜08级环卫考题回忆.docx｜名词解释】}\\
API

\end{pastquestion}

\begin{pastquestion}

\QuestionLabel{【往年题4｜14级环卫考题回忆by 小川川.pdf｜选择回忆】}\\
AQI 的相关描述

\end{pastquestion}

\PointSeparator

\section{大气一次污染物、二次污染物与臭氧}\label{users__fpb__26_spring__pkusph_study__tmp__pdfs__prevention-medicine-past-exams__04-ux73afux5883ux5065ux5eb7ux5b66ux5f80ux5e74ux9898ux8003ux70b9ux6574ux7406.md__ux5927ux6c14ux4e00ux6b21ux6c61ux67d3ux7269ux4e8cux6b21ux6c61ux67d3ux7269ux4e0eux81edux6c27}

\InfoLabel{考频} 3 次

\InfoLabel{对应小节} 三、大气污染 / （三）大气污染物种类；六、大气中主要污染物对人体健康的影响 / 臭氧

\InfoLabel{匹配依据}
题干问``大气二次污染物''\,``光化学烟雾主要二次污染物''等，匹配大气污染物按形成过程分类及臭氧说明。

\ContentBand{知识点原文摘取}

大气污染物按照其形成过程分为：一次污染物；二次污染物。\\
一次污染物：由污染源直接排入大气环境中，其物理和化学性质均为发生改变。\\
二次污染物：排入大气的污染物在物理、化学等因素的作用下发生变化，或与环境中的其他物质发生反应所形成的理化性质不同于一次污染物的新的污染物。如：SO2
在大气中氧化遇水形成硫酸；氮氧化物和挥发性有机物经过光化学反应生成其它污染物。\\
臭氧（ozone, O3）是光化学烟雾主要成分，其刺激性强并有强氧化性，属于二次污染物。

\InfoLabel{知识来源} 《环境健康学------笔记by徐庆松》第14页；《环境健康学------笔记by徐庆松》第24页

\ContentBand{对应往年题}

\begin{pastquestion}

\QuestionLabel{【往年题1｜10级环卫考题回忆.docx｜选择】}\\
光化学烟雾的主要二次污染物

\end{pastquestion}

\begin{pastquestion}

\QuestionLabel{【往年题2｜14级环卫考题回忆by 小川川.pdf｜选择回忆】}\\
哪个是二次污染物

\end{pastquestion}

\begin{pastquestion}

\QuestionLabel{【往年题3｜20级完整试题回忆.docx｜单选题1】}\\
1、下列属于大气二次污染物的是\\
A、SO₂\\
B、NO₂\\
C、O₃\\
D、CO\\
E、花粉

\end{pastquestion}

\PointSeparator

\section{酸雨、空气离子与大气物理性状}\label{users__fpb__26_spring__pkusph_study__tmp__pdfs__prevention-medicine-past-exams__04-ux73afux5883ux5065ux5eb7ux5b66ux5f80ux5e74ux9898ux8003ux70b9ux6574ux7406.md__ux9178ux96e8ux7a7aux6c14ux79bbux5b50ux4e0eux5927ux6c14ux7269ux7406ux6027ux72b6}

\InfoLabel{考频} 2 次

\InfoLabel{对应小节} 二、大气的自然特征 / 大气的物理性状；四、大气污染对健康的影响 / 间接危害

\InfoLabel{匹配依据} 题干涉及酸雨名解、污染空气中阳离子和阴离子，匹配大气物理性状和酸雨内容。

\ContentBand{知识点原文摘取}

空气离子：大气中带电荷物质。近地表大气空气离子分为：轻离子、重离子。污染空气中轻离子浓度低。空气阴离子对机体有镇静、催眠、镇痛、镇咳、降压等作用，而阳离子相反，引起失眠、头痛、烦躁、血压升高。\\
酸雨：降水 pH 小于 5.6。对土壤和植物产生危害；影响水生生态系统；对人类健康产生影响。

\InfoLabel{知识来源} 《环境健康学------笔记by徐庆松》第11页；《环境健康学------笔记by徐庆松》第20页

\ContentBand{对应往年题}

\begin{pastquestion}

\QuestionLabel{【往年题1｜04级环卫考题回忆.doc｜名词解释】}\\
酸雨

\end{pastquestion}

\begin{pastquestion}

\QuestionLabel{【往年题2｜10级环卫考题回忆.docx｜选择】}\\
污染空气中的阳离子和阴离子

\end{pastquestion}

\PointSeparator

\chapter{水与健康}\label{users__fpb__26_spring__pkusph_study__tmp__pdfs__prevention-medicine-past-exams__04-ux73afux5883ux5065ux5eb7ux5b66ux5f80ux5e74ux9898ux8003ux70b9ux6574ux7406.md__ux7b2cux4e09ux7ae0-ux6c34ux4e0eux5065ux5eb7}

\section{水质性状与评价指标：硬度、三氮、DO、BOD、COD、粪大肠菌群}\label{users__fpb__26_spring__pkusph_study__tmp__pdfs__prevention-medicine-past-exams__04-ux73afux5883ux5065ux5eb7ux5b66ux5f80ux5e74ux9898ux8003ux70b9ux6574ux7406.md__ux6c34ux8d28ux6027ux72b6ux4e0eux8bc4ux4ef7ux6307ux6807ux786cux5ea6ux4e09ux6c2edobodcodux7caaux5927ux80a0ux83ccux7fa4}

\InfoLabel{考频} 19 次

\InfoLabel{对应小节} 二、水质的性状和评价指标

\InfoLabel{匹配依据}
题干集中考查暂时硬度、三氮、溶解氧、五日生化需氧量、化学耗氧量/高锰酸盐指数、粪大肠菌群及水质微生物指标。

\ContentBand{知识点原文摘取}

暂时硬度：水经煮沸后，水中重碳酸盐分解形成碳酸盐而沉淀所除去的硬度。\\
水中三氮：氨氮-新鲜污染、亚硝酸盐氮-自净中，硝酸盐氮-自净完成。\\
溶解氧（DO）：溶解在水中的氧含量；可作为有机物污染及自净程度的间接指标。\\
化学耗氧量：一定条件下，使用强氧化剂（高锰酸钾或重铬酸钾等）氧化水中有机物所消耗的量；测定水中有机物含量的间接指标。\\
生化需氧量：水中有机物在有氧条件下被需氧微生物分解时消耗的溶解氧的量；实际工作中规定以 20℃培养 5 日后，1L
水中减少的溶解氧量为 5 日生化需氧量。\\
总大肠菌群与粪大肠菌群：来自人及动物粪便内的大肠菌群主要属粪大肠菌群。\\
粪大肠菌群：是标准中唯一的微生物学指标，表征水体受病原体污染的程度。

\InfoLabel{知识来源}
《环境健康学------笔记by徐庆松》第28页；《环境健康学------笔记by徐庆松》第29页；《环境健康学------笔记by徐庆松》第30页；《环境健康学------笔记by徐庆松》第40页

\ContentBand{对应往年题}

\begin{pastquestion}

\QuestionLabel{【往年题1｜03级环境卫生重点考题总结-曾.doc｜大题】}\\
粪大肠杆菌及其卫生学意义（4）

\end{pastquestion}

\begin{pastquestion}

\QuestionLabel{【往年题2｜03级环境卫生重点考题总结-曾.doc｜口试题】}\\
水体物理性状指标

\end{pastquestion}

\begin{pastquestion}

\QuestionLabel{【往年题3｜03级环境卫生重点考题总结-曾.doc｜口试题】}\\
生化需养量

\end{pastquestion}

\begin{pastquestion}

\QuestionLabel{【往年题4｜03级环境卫生重点考题总结-曾.doc｜口试题】}\\
``三氮''的卫生学意义

\end{pastquestion}

\begin{pastquestion}

\QuestionLabel{【往年题5｜04级环卫考题回忆.doc｜名词解释】}\\
temporary hardness

\end{pastquestion}

\begin{pastquestion}

\QuestionLabel{【往年题6｜04级环卫考题回忆.doc｜简答题】}\\
试述三氮的卫生学意义

\end{pastquestion}

\begin{pastquestion}

\QuestionLabel{【往年题7｜04级环卫考题回忆.doc｜选择】}\\
暂时硬度、细颗粒物

\end{pastquestion}

\begin{pastquestion}

\QuestionLabel{【往年题8｜08级环卫考题回忆.docx｜名词解释】}\\
暂时硬度

\end{pastquestion}

\begin{pastquestion}

\QuestionLabel{【往年题9｜08级环卫考题回忆.docx｜简答】}\\
三氮的卫生学意义

\end{pastquestion}

\begin{pastquestion}

\QuestionLabel{【往年题10｜10级环卫考题回忆.docx｜选择】}\\
暂时硬度

\end{pastquestion}

\begin{pastquestion}

\QuestionLabel{【往年题11｜10级环卫考题回忆.docx｜选择】}\\
粪大肠杆菌的检测意义

\end{pastquestion}

\begin{pastquestion}

\QuestionLabel{【往年题12｜11级环卫考题回忆.docx｜简答】}\\
水中三氮的卫生学意义

\end{pastquestion}

\begin{pastquestion}

\QuestionLabel{【往年题13｜12级环境卫生学考试题回忆版（不全）.docx｜名词解释】}\\
五日生化需氧量

\end{pastquestion}

\begin{pastquestion}

\QuestionLabel{【往年题14｜14级环卫考题回忆by 小川川.pdf｜选择回忆】}\\
哪一个是水质评价的微生物指标

\end{pastquestion}

\begin{pastquestion}

\QuestionLabel{【往年题15｜14级环卫考题回忆by 小川川.pdf｜选择回忆】}\\
化学耗氧量的相关描述

\end{pastquestion}

\begin{pastquestion}

\QuestionLabel{【往年题16｜2021级环境健康学考题yys.docx｜简答】}\\
溶解氧的定义及卫生学意义

\end{pastquestion}

\begin{pastquestion}

\QuestionLabel{【往年题17｜20级完整试题回忆.docx｜单选题7】}\\
7、水的暂时硬度是指\\
A、悬浮性固体和溶解性固体\\
B、钙、镁的重碳酸盐\\
C、钙、镁的硫酸盐和氧化物\\
D、水中不能除去的钙、镁盐类\\
E、溶解于水中的钙、镁盐类总和

\end{pastquestion}

\begin{pastquestion}

\QuestionLabel{【往年题18｜20级完整试题回忆.docx｜名词解释】}\\
高锰酸盐指数

\end{pastquestion}

\begin{pastquestion}

\QuestionLabel{【往年题19｜21级本科预防二班+胡润泽+环境卫生+主观题总结.docx｜简答题】}\\
溶解氧的定义和卫生学意义

\end{pastquestion}

\PointSeparator

\section{介水传染病及流行特点}\label{users__fpb__26_spring__pkusph_study__tmp__pdfs__prevention-medicine-past-exams__04-ux73afux5883ux5065ux5eb7ux5b66ux5f80ux5e74ux9898ux8003ux70b9ux6574ux7406.md__ux4ecbux6c34ux4f20ux67d3ux75c5ux53caux6d41ux884cux7279ux70b9}

\InfoLabel{考频} 5 次

\InfoLabel{对应小节} 三、水体污染 / （七）水体生物性污染的危害

\InfoLabel{匹配依据} 题干直接问介水传染病名解、特点或选择介水传染病，匹配水体生物性污染章节。

\ContentBand{知识点原文摘取}

介水传染病是通过饮用或接触受病原体污染的水而传播的疾病，又称水性传染病。\\
介水传染病流行特点：水源一次严重污染后，可呈暴发流行，短期内突然出现大量病人，且多数患者发病日期集中在同一潜伏期内；病例分布与供水范围一致。大多数患者都有饮用或接触同一水源的历史；一旦对污染源采取治理措施，并加强饮用水的净化和消毒后，疾病的流行能迅速得到控制。

\InfoLabel{知识来源} 《环境健康学------笔记by徐庆松》第37页；《环境健康学------笔记by徐庆松》第38页

\ContentBand{对应往年题}

\begin{pastquestion}

\QuestionLabel{【往年题1｜10级环卫考题回忆.docx｜选择】}\\
介水传染病的特点

\end{pastquestion}

\begin{pastquestion}

\QuestionLabel{【往年题2｜11级环卫考题回忆.docx｜名词解释】}\\
介水传染病

\end{pastquestion}

\begin{pastquestion}

\QuestionLabel{【往年题3｜18级环境健康期末考题回忆by崔浩亮.docx｜简答题】}\\
什么是介水传染病？其流行特点是什么？

\end{pastquestion}

\begin{pastquestion}

\QuestionLabel{【往年题4｜20级完整试题回忆.docx｜单选题8】}\\
8、下列属于介水传染病的是\\
A、水俣病\\
B、地方性甲状腺肿\\
C、氟斑牙\\
D、甲型肝炎\\
E、克汀病

\end{pastquestion}

\begin{pastquestion}

\QuestionLabel{【往年题5｜21级本科预防二班+胡润泽+环境卫生+主观题总结.docx｜简答题】}\\
介水传染病的特点

\end{pastquestion}

\PointSeparator

\section{地表水环境质量标准制定原则与水体监测}\label{users__fpb__26_spring__pkusph_study__tmp__pdfs__prevention-medicine-past-exams__04-ux73afux5883ux5065ux5eb7ux5b66ux5f80ux5e74ux9898ux8003ux70b9ux6574ux7406.md__ux5730ux8868ux6c34ux73afux5883ux8d28ux91cfux6807ux51c6ux5236ux5b9aux539fux5219ux4e0eux6c34ux4f53ux76d1ux6d4b}

\InfoLabel{考频} 4 次

\InfoLabel{对应小节} 四、水环境 / 水环境质量标准；三、水体污染 / 水体污染的调查与监测

\InfoLabel{匹配依据}
题干涉及《地表水环境质量标准》制定原则、江河水体监测、地面水有害物质最高容许浓度制定步骤等，匹配水环境标准和水体监测内容；``最高容许浓度制定步骤''未完整匹配。

\ContentBand{知识点原文摘取}

《地表水环境质量标准》制订原则：防止通过地表水传播疾病；防止通过地表水引起急性或慢性中毒及远期危害；保证地表水感官性状良好；保证地表水自净过程能正常进行。\\
水体污染的调查：污染源调查；水体污染的调查；水体污染对居民健康影响的调查。\\
水体污染的监测：确定水体中污染物的时空分布，追溯污染物的来源和污染途径，了解污染物的迁移转化规律，预测水污染的发展态势。\\
江河水系的监测：采样断面（至少三个，对照断面、污染断面、自净断面）与采样点选择；采样时间的频率；水质检测项目；水体底质的监测；水生生物的监测。

\InfoLabel{知识来源} 《环境健康学------笔记by徐庆松》第38页；《环境健康学------笔记by徐庆松》第39页

\ContentBand{对应往年题}

\begin{pastquestion}

\QuestionLabel{【往年题1｜03级环境卫生重点考题总结-曾.doc｜口试题】}\\
江河水体如何监测

\end{pastquestion}

\begin{pastquestion}

\QuestionLabel{【往年题2｜03级环境卫生重点考题总结-曾.doc｜口试题】}\\
地面水有害物质最高容许浓度的制定步骤

\end{pastquestion}

\begin{pastquestion}

\QuestionLabel{【往年题3｜14级环卫考题回忆by 小川川.pdf｜简答】}\\
地表水环境质量标准制定原则有哪些

\end{pastquestion}

\begin{pastquestion}

\QuestionLabel{【往年题4｜20级完整试题回忆.docx｜简答题】}\\
简述《地表水环境质量标准》制定的原则

\end{pastquestion}

\PointSeparator

\section{水体污染来源、自净机制、转归与水体类型}\label{users__fpb__26_spring__pkusph_study__tmp__pdfs__prevention-medicine-past-exams__04-ux73afux5883ux5065ux5eb7ux5b66ux5f80ux5e74ux9898ux8003ux70b9ux6574ux7406.md__ux6c34ux4f53ux6c61ux67d3ux6765ux6e90ux81eaux51c0ux673aux5236ux8f6cux5f52ux4e0eux6c34ux4f53ux7c7bux578b}

\InfoLabel{考频} 3 次

\InfoLabel{对应小节} 三、水体污染 / 水体污染来源；水体自净；水体污染物的转归

\InfoLabel{匹配依据}
题干涉及水体污染源、自净机制、江河/湖泊/地下水污染特点、四种水体污染自净危害、江河水体监测等。

\ContentBand{知识点原文摘取}

水体污染：人类活动排放的污染物进入水体，其数量超过了水体的自净能力，使水和水体底质的理化特性和水环境中的生物特性、组成等发生改变，从而影响水的有效利用，造成水质恶化，乃至危害人体健康或破坏生态环境的现象。\\
水体污染的主要来源：工业废水；生活污水；农业污水；医院污水；废物不适当堆放和填埋、意外事故等。\\
水体自净：受污染的水体通过自身物理、化学、生物学的作用，使污染成分沉入水底或不断稀释、扩散、分解破坏，污染物的浓度逐步降低，使水质最终恢复到污染前的状态。\\
水体自净的机制包括稀释、混合、吸附沉淀等物理作用，氧化还原、分解化合等化学作用，以及生物分解、生物转化和生物富集等生物学作用。\\
水体污染物的转归：迁移；转化。

\InfoLabel{知识来源}
《环境健康学------笔记by徐庆松》第30页；《环境健康学------笔记by徐庆松》第31页；《环境健康学------笔记by徐庆松》第32页；《环境健康学------笔记by徐庆松》第33页

\ContentBand{对应往年题}

\begin{pastquestion}

\QuestionLabel{【往年题1｜03级环境卫生重点考题总结-曾.doc｜口试题】}\\
四种水体的污染、自净、危害

\end{pastquestion}

\begin{pastquestion}

\QuestionLabel{【往年题2｜03级环境卫生重点考题总结-曾.doc｜口试题】}\\
水体自净机制

\end{pastquestion}

\begin{pastquestion}

\QuestionLabel{【往年题3｜03级环境卫生重点考题总结-曾.doc｜口试题】}\\
水体的污染源和主要污染物

\end{pastquestion}

\PointSeparator

\section{水体富营养化}\label{users__fpb__26_spring__pkusph_study__tmp__pdfs__prevention-medicine-past-exams__04-ux73afux5883ux5065ux5eb7ux5b66ux5f80ux5e74ux9898ux8003ux70b9ux6574ux7406.md__ux6c34ux4f53ux5bccux8425ux517bux5316}

\InfoLabel{考频} 3 次

\InfoLabel{对应小节} 三、水体污染 / 水体富营养化

\InfoLabel{匹配依据} 题干为水体富营养化名词解释或相关污染来源，匹配水体污染章节定义和危害。

\ContentBand{知识点原文摘取}

水体富营养化：受含磷、氮等污水污染造成水体中藻类大量繁殖，使水中有机物增加、溶解氧下降，水质恶化的现象。\\
含氮、磷等污水导致水体中藻类大量繁殖，使水中有机物增加、溶解氧下降、水质恶化的现象。\\
藻类浮于水面可影响水的感观性状，使水质出现异臭异味；藻类大量繁殖死亡后，在细菌分解过程中不断消耗水中的溶解氧，使水中溶解氧含量急剧降低。

\InfoLabel{知识来源} 《环境健康学------笔记by徐庆松》第30页；《环境健康学------笔记by徐庆松》第38页

\ContentBand{对应往年题}

\begin{pastquestion}

\QuestionLabel{【往年题1｜03级环境卫生重点考题总结-曾.doc｜名词解释】}\\
水体富营养化

\end{pastquestion}

\begin{pastquestion}

\QuestionLabel{【往年题2｜14级环卫考题回忆by 小川川.pdf｜名词解释】}\\
水体富营养化

\end{pastquestion}

\begin{pastquestion}

\QuestionLabel{【往年题3｜18级环境健康期末考题回忆by崔浩亮.docx｜名词解释】}\\
水体富营养化

\end{pastquestion}

\PointSeparator

\section{水体化学性污染物：甲基汞、水俣病、多氯联苯等}\label{users__fpb__26_spring__pkusph_study__tmp__pdfs__prevention-medicine-past-exams__04-ux73afux5883ux5065ux5eb7ux5b66ux5f80ux5e74ux9898ux8003ux70b9ux6574ux7406.md__ux6c34ux4f53ux5316ux5b66ux6027ux6c61ux67d3ux7269ux7532ux57faux6c5eux6c34ux4fe3ux75c5ux591aux6c2fux8054ux82efux7b49}

\InfoLabel{考频} 2 次

\InfoLabel{对应小节} 三、水体污染 / （六）水体化学性污染的危害

\InfoLabel{匹配依据} 题干出现 Minamata
disease（水俣病）、米糠油事件、几种水体污染物特点及用途，匹配汞与甲基汞、多氯联苯等内容。

\ContentBand{知识点原文摘取}

汞污染：工业废水及含汞农药的使用。结合于悬浮颗粒物上，然后慢慢地沉积于底泥在微生物的作用下，可转化为甲基汞和二甲基汞。\\
水产品中的汞主要以甲基汞的形式存在。生物对甲基汞具有浓缩作用，水生食物链的生物浓缩作用远高于陆生生物。\\
日本水俣病事件，美国墨西哥州婴儿的甲基汞中毒和伊拉克的甲基汞中毒事件，均证明甲基汞可通过胎盘影响后代。\\
多氯联苯在水环境中稳定性极高，残留时间长，具有生物蓄积性和高毒性；最典型例子是 1968
年发生在日本的``米糠油中毒事件''和 1979 年发生在我国台湾彰化县的``油症事件''。

\InfoLabel{知识来源}
《环境健康学------笔记by徐庆松》第34页；《环境健康学------笔记by徐庆松》第36页；《环境健康学------笔记by徐庆松》第37页

\ContentBand{对应往年题}

\begin{pastquestion}

\QuestionLabel{【往年题1｜03级环境卫生重点考题总结-曾.doc｜口试题】}\\
几种水体污染物的特点及研究的用途

\end{pastquestion}

\begin{pastquestion}

\QuestionLabel{【往年题2｜04级环卫考题回忆.doc｜名词解释】}\\
Minamata disease(水俣病)

\end{pastquestion}

\PointSeparator

\section{地表水与地下水特征}\label{users__fpb__26_spring__pkusph_study__tmp__pdfs__prevention-medicine-past-exams__04-ux73afux5883ux5065ux5eb7ux5b66ux5f80ux5e74ux9898ux8003ux70b9ux6574ux7406.md__ux5730ux8868ux6c34ux4e0eux5730ux4e0bux6c34ux7279ux5f81}

\InfoLabel{考频} 1 次

\InfoLabel{对应小节} 一、水资源 / 水资源的种类

\InfoLabel{匹配依据} 题干问地面水与地下水区别，匹配水资源种类中地表水和地下水的水质、水量和防护特征。

\ContentBand{知识点原文摘取}

地表水：主要来自降水的汇集和地下水的补充，水量较大；有季节变化；水质一般较软，含盐类较少；浑浊度大，细菌含量多；不易防护，易受污染。\\
地下水：降水、地表水经土壤地层渗透进地面以下形成。浅层地下水是我国广大农村最常用的水源，水质物理性状较好，细菌数较地表水少；深层地下水水质透明无色，水温恒定，细菌数很少，但盐类含量高，硬度大。

\InfoLabel{知识来源} 《环境健康学------笔记by徐庆松》第27页

\ContentBand{对应往年题}

\begin{pastquestion}

\QuestionLabel{【往年题1｜03级环境卫生重点考题总结-曾.doc｜口试题】}\\
地面水与地下水区别

\end{pastquestion}

\PointSeparator

\section{中水回用}\label{users__fpb__26_spring__pkusph_study__tmp__pdfs__prevention-medicine-past-exams__04-ux73afux5883ux5065ux5eb7ux5b66ux5f80ux5e74ux9898ux8003ux70b9ux6574ux7406.md__ux4e2dux6c34ux56deux7528}

\InfoLabel{考频} 1 次

\InfoLabel{对应小节} 四、水环境 / 水体卫生防护 / 生活污水

\InfoLabel{匹配依据} 题干为中水回用名词解释，匹配生活污水处理和利用中的中水回用定义。

\ContentBand{知识点原文摘取}

中水回用：采用物理、化学以及生物化学方法将城市污水或生活污水进行处理，使之达到一定水质要求，可在一定范围内重复使用，主要用于冲洗地面、厕所、绿化、喷洒及景观用水等。

\InfoLabel{知识来源} 《环境健康学------笔记by徐庆松》第40页

\ContentBand{对应往年题}

\begin{pastquestion}

\QuestionLabel{【往年题1｜18级环境健康期末考题回忆by崔浩亮.docx｜名词解释】}\\
中水回用

\end{pastquestion}

\PointSeparator

\chapter{饮用水卫生}\label{users__fpb__26_spring__pkusph_study__tmp__pdfs__prevention-medicine-past-exams__04-ux73afux5883ux5065ux5eb7ux5b66ux5f80ux5e74ux9898ux8003ux70b9ux6574ux7406.md__ux7b2cux56dbux7ae0-ux996eux7528ux6c34ux536bux751f}

\section{混凝沉淀、过滤、饮用水消毒与氯化消毒}\label{users__fpb__26_spring__pkusph_study__tmp__pdfs__prevention-medicine-past-exams__04-ux73afux5883ux5065ux5eb7ux5b66ux5f80ux5e74ux9898ux8003ux70b9ux6574ux7406.md__ux6df7ux51ddux6c89ux6dc0ux8fc7ux6ee4ux996eux7528ux6c34ux6d88ux6bd2ux4e0eux6c2fux5316ux6d88ux6bd2}

\InfoLabel{考频} 11 次

\InfoLabel{对应小节} 二、水质处理 / 混凝沉淀、过滤、消毒；三、饮用水与健康问题 / 饮水氯化消毒副产物

\InfoLabel{匹配依据}
题干涉及混凝沉淀原理和影响因素、氯化消毒原理/影响因素、加氯量/余氯/有效氯、消毒方法比较、消毒副产物。

\ContentBand{知识点原文摘取}

生活饮用水的常规净化工艺过程包括混凝沉淀(或澄清)→过滤→消毒。目的是除去原水中的悬浮物质、胶体颗粒和细菌等。\\
混凝沉淀：通过添加混凝剂，将难以自然沉淀的细小颗粒，特别是胶体微粒去除的过程称为混凝沉淀。混凝原理：压缩双电层作用；电性中和作用；吸附架桥作用。\\
消毒：是指杀灭水中病原微生物的措施。目前我国用于饮用水消毒的方法主要有氯化消毒、二氧化氯消毒、紫外线消毒和臭氧消毒等。\\
氯化消毒：用氯或氯制剂进行饮水消毒的一种方法。有效氯：指含氯化合物中具有杀菌能力的有效成分。\\
影响氯化消毒效果的因素：加氯量和接触时间；水的 pH 值；水温；水的浑浊度；水中微生物的种类和数量。\\
饮水氯化消毒副产物：在氯化消毒过程中氯与水中的有机物反应所产生的卤代烃类化合物。

\InfoLabel{知识来源}
《环境健康学------笔记by徐庆松》第41页；《环境健康学------笔记by徐庆松》第42页；《环境健康学------笔记by徐庆松》第43页；《环境健康学------笔记by徐庆松》第44页；《环境健康学------笔记by徐庆松》第45页

\ContentBand{对应往年题}

\begin{pastquestion}

\QuestionLabel{【往年题1｜03级环境卫生重点考题总结-曾.doc｜大题】}\\
氯化消毒的原理和影响因素（10）

\end{pastquestion}

\begin{pastquestion}

\QuestionLabel{【往年题2｜03级环境卫生重点考题总结-曾.doc｜口试题】}\\
混凝沉淀的原理和影响因素

\end{pastquestion}

\begin{pastquestion}

\QuestionLabel{【往年题3｜03级环境卫生重点考题总结-曾.doc｜口试题】}\\
氯胺消毒优缺点

\end{pastquestion}

\begin{pastquestion}

\QuestionLabel{【往年题4｜03级环境卫生重点考题总结-曾.doc｜口试题】}\\
氯消毒的原理及其影响因素

\end{pastquestion}

\begin{pastquestion}

\QuestionLabel{【往年题5｜03级环境卫生重点考题总结-曾.doc｜口试题】}\\
加氯量、余氯、有效氯

\end{pastquestion}

\begin{pastquestion}

\QuestionLabel{【往年题6｜04级环卫考题回忆.doc｜简答题】}\\
试述氯化消毒的原理及影响消毒效果的因素

\end{pastquestion}

\begin{pastquestion}

\QuestionLabel{【往年题7｜08级环卫考题回忆.docx｜简答】}\\
氯化消毒的影响因素

\end{pastquestion}

\begin{pastquestion}

\QuestionLabel{【往年题8｜10级环卫考题回忆.docx｜简答】}\\
影响氯化消毒的因素

\end{pastquestion}

\begin{pastquestion}

\QuestionLabel{【往年题9｜11级环卫考题回忆.docx｜名词解释】}\\
氯化消毒副产物

\end{pastquestion}

\begin{pastquestion}

\QuestionLabel{【往年题10｜12级环境卫生学考试题回忆版（不全）.docx｜简答】}\\
自来水消毒的方法比较

\end{pastquestion}

\begin{pastquestion}

\QuestionLabel{【往年题11｜18级环境健康期末考题回忆by崔浩亮.docx｜简答题】}\\
请从消毒力、持续消毒能力、对pH值敏感性、消毒副产物及费用几个方面对常用的饮用水消毒方法进行简单比较。

\end{pastquestion}

\PointSeparator

\section{二次供水与涉水产品}\label{users__fpb__26_spring__pkusph_study__tmp__pdfs__prevention-medicine-past-exams__04-ux73afux5883ux5065ux5eb7ux5b66ux5f80ux5e74ux9898ux8003ux70b9ux6574ux7406.md__ux4e8cux6b21ux4f9bux6c34ux4e0eux6d89ux6c34ux4ea7ux54c1}

\InfoLabel{考频} 5 次

\InfoLabel{对应小节} 三、饮用水与健康问题 / 高层建筑二次供水污染；六、涉及饮用水卫生安全产品

\InfoLabel{匹配依据} 题干为二次供水、高层建筑二次供水、涉水产品名词解释，匹配第四章对应定义。

\ContentBand{知识点原文摘取}

高层建筑二次供水：又称高层建筑二次加压供水，它是指供水单位将来自集中式供水或自备水源的生活饮用水，贮存于水箱或贮水池中，再通过机械加压或凭借高层建筑形成的自然压差，二次输送至水站或用户的供水系统。\\
涉水产品：是指在饮用水生产和供水过程中与饮用水接触的连接止水材料、塑料及有机合成管材、管件、防护涂料、水处理剂、除垢剂、水质处理器及其他新材料和化学物质。

\InfoLabel{知识来源} 《环境健康学------笔记by徐庆松》第46页；《环境健康学------笔记by徐庆松》第48页

\ContentBand{对应往年题}

\begin{pastquestion}

\QuestionLabel{【往年题1｜10级环卫考题回忆.docx｜名词解释】}\\
涉水产品

\end{pastquestion}

\begin{pastquestion}

\QuestionLabel{【往年题2｜14级环卫考题回忆by 小川川.pdf｜名词解释】}\\
高层建筑的二次供水

\end{pastquestion}

\begin{pastquestion}

\QuestionLabel{【往年题3｜2021级环境健康学考题yys.docx｜名词解释】}\\
二次供水

\end{pastquestion}

\begin{pastquestion}

\QuestionLabel{【往年题4｜20级完整试题回忆.docx｜名词解释】}\\
二次供水

\end{pastquestion}

\begin{pastquestion}

\QuestionLabel{【往年题5｜21级本科预防二班+胡润泽+环境卫生+主观题总结.docx｜名词解释】}\\
二次供水

\end{pastquestion}

\PointSeparator

\section{集中式供水与水源选择原则}\label{users__fpb__26_spring__pkusph_study__tmp__pdfs__prevention-medicine-past-exams__04-ux73afux5883ux5065ux5eb7ux5b66ux5f80ux5e74ux9898ux8003ux70b9ux6574ux7406.md__ux96c6ux4e2dux5f0fux4f9bux6c34ux4e0eux6c34ux6e90ux9009ux62e9ux539fux5219}

\InfoLabel{考频} 3 次

\InfoLabel{对应小节} 一、集中式给水 / 水源选择和卫生防护

\InfoLabel{匹配依据} 题干涉及集中式供水、水源选择原则、集中给水水源选择，匹配第四章集中式给水内容。

\ContentBand{知识点原文摘取}

集中式给水：由水源集中取水，经统一净化处理和消毒后，通过输水管和配水管网送到用户的供水方式。\\
水源选择的原则：水量充足；水质良好；便于防护；技术经济合理。\\
为保护水源，取水点周围应设置保护区；保护区内严禁修建任何可能危害水源水质卫生的设施及一切有碍水源水质卫生的行为。

\InfoLabel{知识来源} 《环境健康学------笔记by徐庆松》第41页

\ContentBand{对应往年题}

\begin{pastquestion}

\QuestionLabel{【往年题1｜03级环境卫生重点考题总结-曾.doc｜口试题】}\\
集中式供水选择原则

\end{pastquestion}

\begin{pastquestion}

\QuestionLabel{【往年题2｜11级环卫考题回忆.docx｜简答】}\\
集中给水的水源选择

\end{pastquestion}

\begin{pastquestion}

\QuestionLabel{【往年题3｜20级完整试题回忆.docx｜单选题9】}\\
9、水源选择的原则不包括\\
A、水量充足\\
B、靠近核心生活区\\
C、水质良好\\
D、便于防护\\
E、技术经济合理

\end{pastquestion}

\PointSeparator

\section{生活饮用水卫生标准制定原则与检测指标}\label{users__fpb__26_spring__pkusph_study__tmp__pdfs__prevention-medicine-past-exams__04-ux73afux5883ux5065ux5eb7ux5b66ux5f80ux5e74ux9898ux8003ux70b9ux6574ux7406.md__ux751fux6d3bux996eux7528ux6c34ux536bux751fux6807ux51c6ux5236ux5b9aux539fux5219ux4e0eux68c0ux6d4bux6307ux6807}

\InfoLabel{考频} 2 次

\InfoLabel{对应小节} 四、饮用水标准 / 生活饮用水卫生标准

\InfoLabel{匹配依据} 题干涉及生活饮用水标准制定原则、饮用水检测指标不包含哪个，匹配饮用水标准章节。

\ContentBand{知识点原文摘取}

制定标准的原则：水中不得含有病原微生物；所含化学物质及放射性物质对人体健康无害；水的感官性状良好，不改变被洗物的理化性状；在选择指标和确定标准限量值时要考虑经济技术上的可行性。\\
常规指标：微生物指标；毒理学指标；感官性状及一般化学指标；放射性指标；消毒剂指标。\\
非常规指标：微生物指标、毒理学指标、感官性状及一般化学指标。

\InfoLabel{知识来源} 《环境健康学------笔记by徐庆松》第46页；《环境健康学------笔记by徐庆松》第47页

\ContentBand{对应往年题}

\begin{pastquestion}

\QuestionLabel{【往年题1｜03级环境卫生重点考题总结-曾.doc｜口试题】}\\
生活饮用水标准制定原则

\end{pastquestion}

\begin{pastquestion}

\QuestionLabel{【往年题2｜2021级环境健康学考题yys.docx｜选择】}\\
根据国家指南，饮用水检测指标不包含哪个：毒理学指标 流行病学指标 一般指标 细菌学指标 消毒剂指标

\end{pastquestion}

\PointSeparator

\section{集中式供水卫生问题调查与监测（未完全匹配到完整原文）}\label{users__fpb__26_spring__pkusph_study__tmp__pdfs__prevention-medicine-past-exams__04-ux73afux5883ux5065ux5eb7ux5b66ux5f80ux5e74ux9898ux8003ux70b9ux6574ux7406.md__ux96c6ux4e2dux5f0fux4f9bux6c34ux536bux751fux95eeux9898ux8c03ux67e5ux4e0eux76d1ux6d4bux672aux5b8cux5168ux5339ux914dux5230ux5b8cux6574ux539fux6587}

\InfoLabel{考频} 1 次

\InfoLabel{对应小节} 未完全匹配到完整调查流程；根据题意暂归入：第四章 饮用水卫生

\InfoLabel{匹配依据}
题干要求集中式供水出现卫生问题时的调查与监测思路。资料有供水定义、水源防护和二次供水污染原因，但未给出完整调查流程。

\ContentBand{知识点原文摘取}

未完全匹配到完整原文。\\
可参考原文片段：集中式给水是指由水源集中取水，经统一净化处理和消毒后，通过输水管和配水管网送到用户的供水方式；为保护水源，取水点周围应设置保护区；二次供水水质污染对健康的影响取决于污染的来源及污染物的性质。

\InfoLabel{知识来源} 《环境健康学------笔记by徐庆松》第41页；《环境健康学------笔记by徐庆松》第46页

\ContentBand{对应往年题}

\begin{pastquestion}

\QuestionLabel{【往年题1｜20级完整试题回忆.docx｜应用题】}\\
某市某小区的集中式供水出现卫生问题，假如你是调查员，请说说你的调查与监测思路

\end{pastquestion}

\PointSeparator

\chapter{生物地球化学性疾病}\label{users__fpb__26_spring__pkusph_study__tmp__pdfs__prevention-medicine-past-exams__04-ux73afux5883ux5065ux5eb7ux5b66ux5f80ux5e74ux9898ux8003ux70b9ux6574ux7406.md__ux7b2cux4e94ux7ae0-ux751fux7269ux5730ux7403ux5316ux5b66ux6027ux75beux75c5}

\section{生物地球化学性疾病概念、流行特征与控制措施}\label{users__fpb__26_spring__pkusph_study__tmp__pdfs__prevention-medicine-past-exams__04-ux73afux5883ux5065ux5eb7ux5b66ux5f80ux5e74ux9898ux8003ux70b9ux6574ux7406.md__ux751fux7269ux5730ux7403ux5316ux5b66ux6027ux75beux75c5ux6982ux5ff5ux6d41ux884cux7279ux5f81ux4e0eux63a7ux5236ux63aaux65bd}

\InfoLabel{考频} 7 次

\InfoLabel{对应小节} 一、概述 / 生物地球化学性疾病

\InfoLabel{匹配依据} 题干直接考生物地球化学性疾病名解、特点或简述，匹配第五章概述。

\ContentBand{知识点原文摘取}

生物地球化学性疾病：系指在自然环境中，由于地壳中元素分布不均衡，有的地区土壤和水中某些元素含量过高或过少，超出了人类和其他生物的适应范围，致使当地动物、植物以及人群中发生某种特异性的疾病。\\
本病往往明显局限于一定地区故亦称为地方病。\\
流行特征：明显的地区性分布；与环境中元素水平相关。\\
影响因素：营养条件；生活习惯；多种元素的联合作用。\\
控制措施：组织措施；技术措施（限制摄入；适量补充）。

\InfoLabel{知识来源} 《环境健康学------笔记by徐庆松》第48页；《环境健康学------笔记by徐庆松》第49页

\ContentBand{对应往年题}

\begin{pastquestion}

\QuestionLabel{【往年题1｜03级环境卫生重点考题总结-曾.doc｜口试题】}\\
举例说明生物地球化学性疾病的特点

\end{pastquestion}

\begin{pastquestion}

\QuestionLabel{【往年题2｜04级环卫考题回忆.doc｜简答题】}\\
简述生物地球化学性疾病

\end{pastquestion}

\begin{pastquestion}

\QuestionLabel{【往年题3｜04级环卫考题回忆.doc｜选择】}\\
生物地球化学疾病

\end{pastquestion}

\begin{pastquestion}

\QuestionLabel{【往年题4｜08级环卫考题回忆.docx｜名词解释】}\\
生物地球化学性疾病

\end{pastquestion}

\begin{pastquestion}

\QuestionLabel{【往年题5｜10级环卫考题回忆.docx｜名词解释】}\\
生物地球化学疾病

\end{pastquestion}

\begin{pastquestion}

\QuestionLabel{【往年题6｜14级环卫考题回忆by 小川川.pdf｜名词解释】}\\
Biogeochemical disease

\end{pastquestion}

\begin{pastquestion}

\QuestionLabel{【往年题7｜18级环境健康期末考题回忆by崔浩亮.docx｜名词解释】}\\
生物地球化学性疾病

\end{pastquestion}

\PointSeparator

\section{碘缺乏病与补碘措施}\label{users__fpb__26_spring__pkusph_study__tmp__pdfs__prevention-medicine-past-exams__04-ux73afux5883ux5065ux5eb7ux5b66ux5f80ux5e74ux9898ux8003ux70b9ux6574ux7406.md__ux7898ux7f3aux4e4fux75c5ux4e0eux8865ux7898ux63aaux65bd}

\InfoLabel{考频} 2 次

\InfoLabel{对应小节} 二、碘缺乏病

\InfoLabel{匹配依据} 题干涉及碘缺乏病或碘缺乏最有效措施，匹配碘缺乏病章节。

\ContentBand{知识点原文摘取}

碘缺乏病：是指从胚胎发育到成人期由于摄入碘不足所引起的一系列病症，它对人的损害是连续性的。包括地方性甲状腺肿、地方性克汀病、地方性亚临床克汀病等。\\
碘缺乏病预防措施：补碘。碘盐；碘油；其他。\\
1994 年国家正式颁布《食盐加碘消除碘缺乏危害管理条例》，食盐加碘作为一项国策固定下来；2012
年起，因地制宜、分类指导、科学补碘。

\InfoLabel{知识来源} 《环境健康学------笔记by徐庆松》第49页；《环境健康学------笔记by徐庆松》第50页

\ContentBand{对应往年题}

\begin{pastquestion}

\QuestionLabel{【往年题1｜03级环境卫生重点考题总结-曾.doc｜名词解释】}\\
碘缺乏病

\end{pastquestion}

\begin{pastquestion}

\QuestionLabel{【往年题2｜14级环卫考题回忆by 小川川.pdf｜选择回忆】}\\
碘缺乏的最有效措施

\end{pastquestion}

\PointSeparator

\section{地方性氟中毒、氟斑牙与氟骨症}\label{users__fpb__26_spring__pkusph_study__tmp__pdfs__prevention-medicine-past-exams__04-ux73afux5883ux5065ux5eb7ux5b66ux5f80ux5e74ux9898ux8003ux70b9ux6574ux7406.md__ux5730ux65b9ux6027ux6c1fux4e2dux6bd2ux6c1fux6591ux7259ux4e0eux6c1fux9aa8ux75c7}

\InfoLabel{考频} 2 次

\InfoLabel{对应小节} 四、地方性氟中毒

\InfoLabel{匹配依据} 题干涉及氟斑牙、氟骨症、地方性氟中毒调查，匹配地方性氟中毒章节。

\ContentBand{知识点原文摘取}

地方性氟中毒：由于一定地区的环境中氟元素过多，而致生活在该环境中居民经饮水、食物和空气等途径长期摄入过量氟而发生的一种慢性全身性疾病，主要表现为氟斑牙(dental
fluorosis)和氟骨症(skeletal fluorosis)。\\
病区类型与分布：饮水型；燃煤型；饮茶型。\\
氟斑牙：釉面光泽度改变；釉面着色；釉面缺损。\\
氟骨症：疼痛（最常见自觉症状）；神经症状；肢体变形。预防措施：改换水源；饮水除氟；改良炉灶；减少食物氟污染；不用或少用高氟劣质煤。

\InfoLabel{知识来源}
《环境健康学------笔记by徐庆松》第53页；《环境健康学------笔记by徐庆松》第54页；《环境健康学------笔记by徐庆松》第55页

\ContentBand{对应往年题}

\begin{pastquestion}

\QuestionLabel{【往年题1｜03级环境卫生重点考题总结-曾.doc｜口试题】}\\
氟骨症、氟斑牙

\end{pastquestion}

\begin{pastquestion}

\QuestionLabel{【往年题2｜03级环境卫生重点考题总结-曾.doc｜口试题】}\\
发现牙龋病，如何展开调查确定地方性氟中毒

\end{pastquestion}

\PointSeparator

\chapter{土壤与健康}\label{users__fpb__26_spring__pkusph_study__tmp__pdfs__prevention-medicine-past-exams__04-ux73afux5883ux5065ux5eb7ux5b66ux5f80ux5e74ux9898ux8003ux70b9ux6574ux7406.md__ux7b2cux516dux7ae0-ux571fux58e4ux4e0eux5065ux5eb7}

\section{土壤物理特征、腐殖质与土壤环境容量}\label{users__fpb__26_spring__pkusph_study__tmp__pdfs__prevention-medicine-past-exams__04-ux73afux5883ux5065ux5eb7ux5b66ux5f80ux5e74ux9898ux8003ux70b9ux6574ux7406.md__ux571fux58e4ux7269ux7406ux7279ux5f81ux8150ux6b96ux8d28ux4e0eux571fux58e4ux73afux5883ux5bb9ux91cf}

\InfoLabel{考频} 11 次

\InfoLabel{对应小节} 一、土壤环境特征 / 土壤组成、物理学特征、化学特征

\InfoLabel{匹配依据} 题干涉及腐殖质、土壤环境容量、建筑潮湿与毛细管作用、土壤卫生学特点等，匹配土壤环境特征。

\ContentBand{知识点原文摘取}

土壤颗粒的大小和排列状态决定着土壤的孔隙率、透气性、渗水性、容水性和土壤的毛细管现象等许多物理特性，影响土壤的卫生状态。\\
土壤的毛细管作用：土壤中的水分沿着孔隙上升的作用；颗粒↓ 孔隙↓毛细管作用↑。\\
土壤环境容量：又称土壤负载容量，是一定土壤环境单元在一定时限内遵循环境质量标准，维持土壤生态系统的正常结构与功能，保证农产品的生物学产量与质量，在不使环境系统污染的前提下，土壤环境所能容纳污染物的最大负荷量。\\
腐殖质（humus）：是土壤特有的有机物质，是进入土壤的植物、动物等死亡残体经微生物分解转化形成的黑色胶体物质，约占土壤有机成分总量的
85\%～90\%。当土壤中有机物大部分转变为腐殖质时，病原体已经死亡，故土壤中腐殖质含量越高，卫生学上越安全。

\InfoLabel{知识来源} 《环境健康学------笔记by徐庆松》第58页；《环境健康学------笔记by徐庆松》第59页

\ContentBand{对应往年题}

\begin{pastquestion}

\QuestionLabel{【往年题1｜03级环境卫生重点考题总结-曾.doc｜名词解释】}\\
腐殖质

\end{pastquestion}

\begin{pastquestion}

\QuestionLabel{【往年题2｜08级环卫考题回忆.docx｜简答】}\\
腐殖质的概念和卫生学意义

\end{pastquestion}

\begin{pastquestion}

\QuestionLabel{【往年题3｜11级环卫考题回忆.docx｜名词解释】}\\
土壤环境容量

\end{pastquestion}

\begin{pastquestion}

\QuestionLabel{【往年题4｜11级环卫考题回忆.docx｜简答】}\\
土壤卫生学特点

\end{pastquestion}

\begin{pastquestion}

\QuestionLabel{【往年题5｜12级环境卫生学考试题回忆版（不全）.docx｜名词解释】}\\
土壤环境容量

\end{pastquestion}

\begin{pastquestion}

\QuestionLabel{【往年题6｜18级环境健康期末考题回忆by崔浩亮.docx｜名词解释】}\\
腐殖质

\end{pastquestion}

\begin{pastquestion}

\QuestionLabel{【往年题7｜2021级环境健康学考题yys.docx｜名词解释】}\\
腐殖质

\end{pastquestion}

\begin{pastquestion}

\QuestionLabel{【往年题8｜20级完整试题回忆.docx｜单选题10】}\\
10、建筑物地面和墙壁的潮湿现象，与土壤的哪一项卫生学特征相关\\
A、容水量\\
B、渗水性\\
C、通透性\\
D、毛细管作用\\
E、孔隙度

\end{pastquestion}

\begin{pastquestion}

\QuestionLabel{【往年题9｜20级完整试题回忆.docx｜名词解释】}\\
土壤环境容量

\end{pastquestion}

\begin{pastquestion}

\QuestionLabel{【往年题10｜20级完整试题回忆.docx｜名词解释】}\\
腐殖质

\end{pastquestion}

\begin{pastquestion}

\QuestionLabel{【往年题11｜21级本科预防二班+胡润泽+环境卫生+主观题总结.docx｜名词解释】}\\
腐殖质

\end{pastquestion}

\PointSeparator

\section{土壤自净、有机物无机化与腐殖质化}\label{users__fpb__26_spring__pkusph_study__tmp__pdfs__prevention-medicine-past-exams__04-ux73afux5883ux5065ux5eb7ux5b66ux5f80ux5e74ux9898ux8003ux70b9ux6574ux7406.md__ux571fux58e4ux81eaux51c0ux6709ux673aux7269ux65e0ux673aux5316ux4e0eux8150ux6b96ux8d28ux5316}

\InfoLabel{考频} 7 次

\InfoLabel{对应小节} 三、土壤自净作用

\InfoLabel{匹配依据} 题干直接考土壤自净作用、方式、微生物作用、有机物无机化和腐殖质化。

\ContentBand{知识点原文摘取}

土壤自净作用（soil
self-purification）：指受污染的土壤通过物理、化学和生物学等自然因素的作用，使病原体死灭，各种有害物质转化到无害的程度，土壤可逐渐恢复到污染前的状态。\\
物理净化：机械阻留、吸附及滤过作用。\\
化学净化：污染物进入土壤后，可发生一系列化学反应（包括氧化还原反应、水解、分解、酸碱中和反应等）。\\
生物净化：土壤中存在大量依靠有机物生存的微生物，它们具有氧化分解有机物的巨大能力，是土壤自净作用中最重要的途径之一，包括病原体死灭、有机物净化。\\
土壤中的有机污染物在微生物的作用下，使得有机物逐步无机化或腐殖质化。

\InfoLabel{知识来源} 《环境健康学------笔记by徐庆松》第60页；《环境健康学------笔记by徐庆松》第61页

\ContentBand{对应往年题}

\begin{pastquestion}

\QuestionLabel{【往年题1｜03级环境卫生重点考题总结-曾.doc｜大题】}\\
土壤污染物的来源和自净（6）

\end{pastquestion}

\begin{pastquestion}

\QuestionLabel{【往年题2｜03级环境卫生重点考题总结-曾.doc｜口试题】}\\
有机物无机化、腐殖质化

\end{pastquestion}

\begin{pastquestion}

\QuestionLabel{【往年题3｜03级环境卫生重点考题总结-曾.doc｜口试题】}\\
土壤的污染来源及其自净

\end{pastquestion}

\begin{pastquestion}

\QuestionLabel{【往年题4｜03级环境卫生重点考题总结-曾.doc｜口试题】}\\
微生物对土壤自净过程中的作用

\end{pastquestion}

\begin{pastquestion}

\QuestionLabel{【往年题5｜04级环卫考题回忆.doc｜选择】}\\
土壤的自净

\end{pastquestion}

\begin{pastquestion}

\QuestionLabel{【往年题6｜14级环卫考题回忆by 小川川.pdf｜名词解释】}\\
土壤的自净作用

\end{pastquestion}

\begin{pastquestion}

\QuestionLabel{【往年题7｜21级本科预防二班+胡润泽+环境卫生+主观题总结.docx｜简答题】}\\
土壤自净的定义和方式

\end{pastquestion}

\PointSeparator

\section{土壤污染定义、来源、特点与土壤环境质量标准}\label{users__fpb__26_spring__pkusph_study__tmp__pdfs__prevention-medicine-past-exams__04-ux73afux5883ux5065ux5eb7ux5b66ux5f80ux5e74ux9898ux8003ux70b9ux6574ux7406.md__ux571fux58e4ux6c61ux67d3ux5b9aux4e49ux6765ux6e90ux7279ux70b9ux4e0eux571fux58e4ux73afux5883ux8d28ux91cfux6807ux51c6}

\InfoLabel{考频} 4 次

\InfoLabel{对应小节} 二、土壤污染；五、土壤污染的防控 / 土壤环境质量标准

\InfoLabel{匹配依据}
题干涉及土壤污染名解、土壤污染特点、来源、土壤卫生标准制定原则，匹配第六章土壤污染和防控。

\ContentBand{知识点原文摘取}

土壤污染（soil
pollution）：指在人类生产和生活活动中排出的有害物质进入土壤中，超过一定限量，直接或间接地危害人畜健康的现象。\\
污染来源：工业污染；农业污染；生活污染；交通污染；灾害污染；电子垃圾污染。\\
污染特点：隐蔽性；累积性；不可逆转性；长期性。\\
土壤环境质量标准：国家为防止土壤污染、保护生态系统、维护人体健康所制订的土壤中污染物在一定的时间和空间范围内的容许含量值。\\
制订土壤环境质量标准的原则：保护陆地生态安全；保护人体健康。

\InfoLabel{知识来源}
《环境健康学------笔记by徐庆松》第59页；《环境健康学------笔记by徐庆松》第60页；《环境健康学------笔记by徐庆松》第65页；《环境健康学------笔记by徐庆松》第66页

\ContentBand{对应往年题}

\begin{pastquestion}

\QuestionLabel{【往年题1｜03级环境卫生重点考题总结-曾.doc｜口试题】}\\
土壤卫生学指定原则、方法

\end{pastquestion}

\begin{pastquestion}

\QuestionLabel{【往年题2｜10级环卫考题回忆.docx｜名词解释】}\\
土壤污染

\end{pastquestion}

\begin{pastquestion}

\QuestionLabel{【往年题3｜10级环卫考题回忆.docx｜简答】}\\
土壤环境卫生标准和大气、水环境的指定标准的区别，以及制订土壤卫生标准的原则

\end{pastquestion}

\begin{pastquestion}

\QuestionLabel{【往年题4｜14级环卫考题回忆by 小川川.pdf｜选择回忆】}\\
土壤污染的特点

\end{pastquestion}

\PointSeparator

\section{农药、持久性有机污染物与电子垃圾污染}\label{users__fpb__26_spring__pkusph_study__tmp__pdfs__prevention-medicine-past-exams__04-ux73afux5883ux5065ux5eb7ux5b66ux5f80ux5e74ux9898ux8003ux70b9ux6574ux7406.md__ux519cux836fux6301ux4e45ux6027ux6709ux673aux6c61ux67d3ux7269ux4e0eux7535ux5b50ux5783ux573eux6c61ux67d3}

\InfoLabel{考频} 4 次

\InfoLabel{对应小节} 四、土壤污染的健康影响 / 农药污染、持久性有机污染物；二、土壤污染 / 电子垃圾污染

\InfoLabel{匹配依据} 题干涉及农药在土壤中的转归、持续性有机污染物、电子产品拆解污染物健康影响，匹配农药、POPs
和电子垃圾污染内容。

\ContentBand{知识点原文摘取}

化学农药在土壤中的迁移转化：土壤对农药的吸附；化学农药在土壤中挥发、扩散、迁移；农药降解过程：光化学降解、化学降解、生物降解。\\
持久性有机污染物（persistent organic
pollutants，POPs）：指能持久存在于环境中，并可借助大气、水、生物体等环境介质进行远距离迁移，通过食物链富集，对环境和人类造成严重危害的天然或人工合成的有机污染物质。\\
持久性有机污染物的特性：持久性；蓄积性；迁移性；高毒性。\\
持久性有机污染物的危害：免疫系统；神经系统；内分泌系统；生殖系统；致癌、致畸、致突变。

\InfoLabel{知识来源}
《环境健康学------笔记by徐庆松》第60页；《环境健康学------笔记by徐庆松》第61页；《环境健康学------笔记by徐庆松》第64页；《环境健康学------笔记by徐庆松》第65页

\ContentBand{对应往年题}

\begin{pastquestion}

\QuestionLabel{【往年题1｜12级环境卫生学考试题回忆版（不全）.docx｜名词解释】}\\
持续性有机污染物（英文）

\end{pastquestion}

\begin{pastquestion}

\QuestionLabel{【往年题2｜12级环境卫生学考试题回忆版（不全）.docx｜简答】}\\
农药在土壤中的转归

\end{pastquestion}

\begin{pastquestion}

\QuestionLabel{【往年题3｜14级环卫考题回忆by 小川川.pdf｜选择回忆】}\\
持续性有机污染物的特点

\end{pastquestion}

\begin{pastquestion}

\QuestionLabel{【往年题4｜14级环卫考题回忆by 小川川.pdf｜应用】}\\
上世纪80年代末期和90年代初，贵屿开始涉及旧五金电器的拆解生意。由于获利丰厚，整个行业规模逐渐扩大。国外的电子废物通过深圳、广州和南海的转运点，开始大规模进入贵屿。大面积的土地开始抛荒，贵屿镇80\%的家庭参与到这个行业中来，并通过这个行业迅速积累财富。为了节省成本，贵屿的家庭作坊往往采用最直接和最原始的方式进行电子废物的拆解。已知电子拆解后产生的主要污染物是重金属和持续性有机污染物，请设计调查计划，了解电子产品拆解产生的污染物对当地居民健康的影响。

\end{pastquestion}

\PointSeparator

\section{土壤重金属迁移转化、危害与监测}\label{users__fpb__26_spring__pkusph_study__tmp__pdfs__prevention-medicine-past-exams__04-ux73afux5883ux5065ux5eb7ux5b66ux5f80ux5e74ux9898ux8003ux70b9ux6574ux7406.md__ux571fux58e4ux91cdux91d1ux5c5eux8fc1ux79fbux8f6cux5316ux5371ux5bb3ux4e0eux76d1ux6d4b}

\InfoLabel{考频} 3 次

\InfoLabel{对应小节} 三、土壤自净作用 / 污染物转归；四、土壤污染的健康影响 /
重金属污染；五、土壤卫生监督与监测

\InfoLabel{匹配依据}
题干涉及重金属离子迁移转化因素、重金属污染健康影响、土壤重金属监测方案、铅/镉等，匹配第六章重金属相关内容。

\ContentBand{知识点原文摘取}

重金属在土壤中的转化：土壤胶体、腐殖质的吸附和螯合作用；土壤 pH 的影响；土壤氧化还原状态的影响。\\
重金属不能被土壤微生物所分解，易于累积，或转化为毒性更大的化合物。有的重金属可通过食物链传递在人体内蓄积，严重危害人体健康。\\
土壤卫生监测的任务是要查明土壤的卫生状况，阐明其对环境的污染和对居民健康可能产生的影响，为保证生态环境和保障人体健康提出卫生要求和防护措施的依据。\\
化学污染的调查监测：不仅要调查监测土壤中化学物质的含量，还要监测当地各种农作物的含量，观察该污染物在农作物中的富集情况。

\InfoLabel{知识来源}
《环境健康学------笔记by徐庆松》第61页；《环境健康学------笔记by徐庆松》第62页；《环境健康学------笔记by徐庆松》第69页

\ContentBand{对应往年题}

\begin{pastquestion}

\QuestionLabel{【往年题1｜03级环境卫生重点考题总结-曾.doc｜口试题】}\\
土壤中重金属迁移转化的影响因素

\end{pastquestion}

\begin{pastquestion}

\QuestionLabel{【往年题2｜2021级环境健康学考题yys.docx｜选择】}\\
哪个不是影响重金属离子在土壤中迁徙转化的因素：PH 太阳辐射 腐殖质嵌合 胶体固着 氧化还原价态

\end{pastquestion}

\begin{pastquestion}

\QuestionLabel{【往年题3｜20级完整试题回忆.docx｜单选题11】}\\
11、重金属在土壤中的转化过程主要为\\
A、土壤胶体吸附\\
B、土壤微生物降解\\
C、化学降解\\
D、光化学降解\\
E、土壤微生物吸附

\end{pastquestion}

\PointSeparator

\section{痛痛病（未完全匹配到）}\label{users__fpb__26_spring__pkusph_study__tmp__pdfs__prevention-medicine-past-exams__04-ux73afux5883ux5065ux5eb7ux5b66ux5f80ux5e74ux9898ux8003ux70b9ux6574ux7406.md__ux75dbux75dbux75c5ux672aux5b8cux5168ux5339ux914dux5230}

\InfoLabel{考频} 2 次

\InfoLabel{对应小节} 未完全匹配到原文；根据题意暂归入：第六章 土壤与健康/重金属污染

\InfoLabel{匹配依据}
题库多次出现``痛痛病''名词。资料中有镉污染和慢性镉中毒防治原则，但未检索到``痛痛病''原文。

\ContentBand{知识点原文摘取}

未完全匹配到原文。

\InfoLabel{知识来源} 未完全匹配到

\ContentBand{对应往年题}

\begin{pastquestion}

\QuestionLabel{【往年题1｜03级环境卫生重点考题总结-曾.doc｜名词解释】}\\
痛痛病

\end{pastquestion}

\begin{pastquestion}

\QuestionLabel{【往年题2｜08级环卫考题回忆.docx｜名词解释】}\\
痛痛病

\end{pastquestion}

\PointSeparator

\section{土壤生物性污染危害}\label{users__fpb__26_spring__pkusph_study__tmp__pdfs__prevention-medicine-past-exams__04-ux73afux5883ux5065ux5eb7ux5b66ux5f80ux5e74ux9898ux8003ux70b9ux6574ux7406.md__ux571fux58e4ux751fux7269ux6027ux6c61ux67d3ux5371ux5bb3}

\InfoLabel{考频} 1 次

\InfoLabel{对应小节} 四、土壤污染的健康影响 / 生物性污染的危害

\InfoLabel{匹配依据} 题干问土壤生物性污染危害，匹配第六章土壤生物性污染。

\ContentBand{知识点原文摘取}

生物性污染的危害：引起肠道传染病和寄生虫病；引起钩端螺旋体病和炭疽病；引起破伤风和肉毒中毒。\\
人体排出的含有病原体的粪便污染土壤，人生吃在这种土壤中种植的蔬菜瓜果等而感染得病（人---土壤---人）。\\
含有病原体的动物粪便污染土壤后，病原体通过皮肤或黏膜进入人体而得病（动物---土壤---人）。

\InfoLabel{知识来源} 《环境健康学------笔记by徐庆松》第65页

\ContentBand{对应往年题}

\begin{pastquestion}

\QuestionLabel{【往年题1｜03级环境卫生重点考题总结-曾.doc｜口试题】}\\
土壤生物性污染危害

\end{pastquestion}

\PointSeparator

\section{土壤卫生监督与监测方案}\label{users__fpb__26_spring__pkusph_study__tmp__pdfs__prevention-medicine-past-exams__04-ux73afux5883ux5065ux5eb7ux5b66ux5f80ux5e74ux9898ux8003ux70b9ux6574ux7406.md__ux571fux58e4ux536bux751fux76d1ux7763ux4e0eux76d1ux6d4bux65b9ux6848}

\InfoLabel{考频} 1 次

\InfoLabel{对应小节} 五、土壤污染的防控 / 土壤卫生监督与监测

\InfoLabel{匹配依据}
题干要求电子工业园区土壤重金属污染卫生监测方案，或土壤卫生学制定/监测方法，匹配土壤卫生监督与监测。

\ContentBand{知识点原文摘取}

土壤卫生监测的任务是要查明土壤的卫生状况，阐明其对环境的污染和对居民健康可能产生的影响，为保证生态环境和保障人体健康提出卫生要求和防护措施的依据。\\
污染源的调查：查清污染来源和特点，要调查污染源的性质、数量、生产过程、净化设施、污染物的排放规律以及影响因素等。\\
土壤污染状况调查与监测：采样点的选择和采样方法；土壤环境背景调查监测；化学污染的调查监测；生物性污染的调查监测。\\
土壤污染对居民健康影响的调查：患病率和死亡率调查；居民询问调查；居民健康检查；有害物质在居民体内蓄积水平的调查。

\InfoLabel{知识来源}
《环境健康学------笔记by徐庆松》第68页；《环境健康学------笔记by徐庆松》第69页；《环境健康学------笔记by徐庆松》第70页

\ContentBand{对应往年题}

\begin{pastquestion}

\QuestionLabel{【往年题1｜18级环境健康期末考题回忆by崔浩亮.docx｜论述题】}\\
某市要对其郊区一处电子工业园区的土壤重金属污染问题进行监测，请提出你的卫生监测方案。

\end{pastquestion}

\PointSeparator

\chapter{住宅与健康}\label{users__fpb__26_spring__pkusph_study__tmp__pdfs__prevention-medicine-past-exams__04-ux73afux5883ux5065ux5eb7ux5b66ux5f80ux5e74ux9898ux8003ux70b9ux6574ux7406.md__ux7b2cux4e03ux7ae0-ux4f4fux5b85ux4e0eux5065ux5eb7}

\section{室内空气污染来源、种类、特点与清洁度指标}\label{users__fpb__26_spring__pkusph_study__tmp__pdfs__prevention-medicine-past-exams__04-ux73afux5883ux5065ux5eb7ux5b66ux5f80ux5e74ux9898ux8003ux70b9ux6574ux7406.md__ux5ba4ux5185ux7a7aux6c14ux6c61ux67d3ux6765ux6e90ux79cdux7c7bux7279ux70b9ux4e0eux6e05ux6d01ux5ea6ux6307ux6807}

\InfoLabel{考频} 12 次

\InfoLabel{对应小节} 四、室内空气污染与健康；三、住宅小气候与健康 / 不良小气候健康影响

\InfoLabel{匹配依据}
题干涉及室内空气污染物种类及来源、室内空气污染特点、室内空气清洁程度和燃烧产物指标、霉菌控制、室内生物污染危害、改善室内空气清洁状态措施。部分具体指标未完整匹配。

\ContentBand{知识点原文摘取}

室内空气污染：由于室内本身存在或外来进入能释放有害物质的污染源，导致室内空气中有害物质浓度和种类在有限空间不断增加，从而对人群健康产生直接或间接，近期或远期，或者潜在的有害影响。\\
室内空气污染的来源：室外源：室外空气、建筑物自身、人为带入室内、相邻住宅污染、建筑用水污染；室内源：室内燃烧加热、室内活动、室内装饰材料及家具、室内生物性污染。\\
主要种类、来源及健康影响：化学性污染物：CO2、燃烧产物、烹调油烟、甲醛等挥发性有机化合物
VOC；物理性污染物：噪声、非电离辐射、电离辐射；生物性污染物：微生物在空气中传播-尘螨；放射性污染物：氡。\\
室内温湿度高：温暖潮湿→蚊子、蟑螂、尘螨等孳生↑；气流↑→污染物排出室外↑。

\InfoLabel{知识来源}
《环境健康学------笔记by徐庆松》第76页；《环境健康学------笔记by徐庆松》第77页；《环境健康学------笔记by徐庆松》第78页

\ContentBand{对应往年题}

\begin{pastquestion}

\QuestionLabel{【往年题1｜03级环境卫生重点考题总结-曾.doc｜大题】}\\
室内空气污染物的种类及其来源（6）

\end{pastquestion}

\begin{pastquestion}

\QuestionLabel{【往年题2｜03级环境卫生重点考题总结-曾.doc｜口试题】}\\
室内空气污染物的种类及其来源

\end{pastquestion}

\begin{pastquestion}

\QuestionLabel{【往年题3｜03级环境卫生重点考题总结-曾.doc｜口试题】}\\
室内空气清洁度指标

\end{pastquestion}

\begin{pastquestion}

\QuestionLabel{【往年题4｜03级环境卫生重点考题总结-曾.doc｜口试题】}\\
室内生物污染的危害

\end{pastquestion}

\begin{pastquestion}

\QuestionLabel{【往年题5｜03级环境卫生重点考题总结-曾.doc｜口试题】}\\
指出三种室内空气污染物及其危害

\end{pastquestion}

\begin{pastquestion}

\QuestionLabel{【往年题6｜03级环境卫生重点考题总结-曾.doc｜口试题】}\\
改善室内空气清洁状态的措施

\end{pastquestion}

\begin{pastquestion}

\QuestionLabel{【往年题7｜03级环境卫生重点考题总结-曾.doc｜口试题】}\\
建筑材料污染物及其危害

\end{pastquestion}

\begin{pastquestion}

\QuestionLabel{【往年题8｜04级环卫考题回忆.doc｜简答题】}\\
反应室内空气清洁程度和燃烧产物的常用指标分别有哪些？

\end{pastquestion}

\begin{pastquestion}

\QuestionLabel{【往年题9｜10级环卫考题回忆.docx｜简答】}\\
室内空气污染的特点

\end{pastquestion}

\begin{pastquestion}

\QuestionLabel{【往年题10｜14级环卫考题回忆by 小川川.pdf｜简答】}\\
室内空气污染的主要来源有哪些

\end{pastquestion}

\begin{pastquestion}

\QuestionLabel{【往年题11｜20级完整试题回忆.docx｜单选题12】}\\
12、减少霉菌生长，可以通过\\
A、增加室内温度\\
B、增加室内湿度\\
C、加强通风\\
D、使用空气清新剂\\
E、在室内种植物

\end{pastquestion}

\begin{pastquestion}

\QuestionLabel{【往年题12｜20级完整试题回忆.docx｜名词解释】}\\
挥发性有机化合物

\end{pastquestion}

\PointSeparator

\section{不良建筑综合征（SBS）与建筑相关疾病（BRI）}\label{users__fpb__26_spring__pkusph_study__tmp__pdfs__prevention-medicine-past-exams__04-ux73afux5883ux5065ux5eb7ux5b66ux5f80ux5e74ux9898ux8003ux70b9ux6574ux7406.md__ux4e0dux826fux5efaux7b51ux7efcux5408ux5f81sbsux4e0eux5efaux7b51ux76f8ux5173ux75beux75c5bri}

\InfoLabel{考频} 9 次

\InfoLabel{对应小节} 七、住宅环境引发的健康问题

\InfoLabel{匹配依据} 题干涉及
SBS、不良建筑综合征、病态建筑综合征、建筑相关疾病以及两者区别，匹配第七章住宅环境健康问题。

\ContentBand{知识点原文摘取}

不良建筑综合征（Sick building Syndrome,
SBS）：建筑内不良因素，使得在建筑内活动的人群产生一系列症状（眼、上呼吸道刺激，头晕、头痛、皮肤干燥、记忆力衰退、全身不适等
30 多种症状）。特点：在建筑内出现症状；离开建筑，症状自行消失。\\
建筑相关疾病（Building related illness,
BRI）：由建筑内有害因素引起的疾病，如：呼吸道感染、哮喘、心血管病、肺癌等。\\
BRI 与 SBS
区别：患者的症状在临床上可以明确诊断；病因可以鉴别确认；患者即使离开致病现场，症状也不会很快消失，必须进行治疗才能恢复健康。

\InfoLabel{知识来源} 《环境健康学------笔记by徐庆松》第79页；《环境健康学------笔记by徐庆松》第80页

\ContentBand{对应往年题}

\begin{pastquestion}

\QuestionLabel{【往年题1｜03级环境卫生重点考题总结-曾.doc｜名词解释】}\\
SBS

\end{pastquestion}

\begin{pastquestion}

\QuestionLabel{【往年题2｜03级环境卫生重点考题总结-曾.doc｜口试题】}\\
与建筑相关的疾病

\end{pastquestion}

\begin{pastquestion}

\QuestionLabel{【往年题3｜04级环卫考题回忆.doc｜名词解释】}\\
BRI

\end{pastquestion}

\begin{pastquestion}

\QuestionLabel{【往年题4｜08级环卫考题回忆.docx｜名词解释】}\\
不良建筑综合症

\end{pastquestion}

\begin{pastquestion}

\QuestionLabel{【往年题5｜11级环卫考题回忆.docx｜名词解释】}\\
不良建筑物综合征

\end{pastquestion}

\begin{pastquestion}

\QuestionLabel{【往年题6｜18级环境健康期末考题回忆by崔浩亮.docx｜简答题】}\\
简述病态建筑综合征和建筑相关疾病的区别。

\end{pastquestion}

\begin{pastquestion}

\QuestionLabel{【往年题7｜2021级环境健康学考题yys.docx｜名词解释】}\\
病态建筑综合征

\end{pastquestion}

\begin{pastquestion}

\QuestionLabel{【往年题8｜20级完整试题回忆.docx｜单选题16】}\\
16、下列哪项属于建筑相关疾病\\
A、在室内使用化妆品后出现皮肤过敏\\
B、搬入新装修的新房子后引发过敏性鼻炎\\
C、由于室内污染导致的记忆力减退\\
D、室内空气干燥所致的嗓子干痒\\
E、\ldots\ldots{}

\end{pastquestion}

\begin{pastquestion}

\QuestionLabel{【往年题9｜21级本科预防二班+胡润泽+环境卫生+主观题总结.docx｜名词解释】}\\
病态建筑综合征

\end{pastquestion}

\PointSeparator

\section{室内小气候、热环境、有效温度与校正有效温度}\label{users__fpb__26_spring__pkusph_study__tmp__pdfs__prevention-medicine-past-exams__04-ux73afux5883ux5065ux5eb7ux5b66ux5f80ux5e74ux9898ux8003ux70b9ux6574ux7406.md__ux5ba4ux5185ux5c0fux6c14ux5019ux70edux73afux5883ux6709ux6548ux6e29ux5ea6ux4e0eux6821ux6b63ux6709ux6548ux6e29ux5ea6}

\InfoLabel{考频} 6 次

\InfoLabel{对应小节} 三、住宅小气候与健康

\InfoLabel{匹配依据} 题干涉及室内热环境定义、主要影响因素、有效温度、校正有效温度、微小气候卫生学要求。

\ContentBand{知识点原文摘取}

室内小气候（热环境）：住宅的室内由于屋顶、地板、门窗和墙壁等围护结构以及室内的人工空气调节设备等综合作用，形成了与室外不同的室内气候。\\
四个热环境要素：空气温度、空气湿度、气流、热辐射。\\
有效温度(effective temperature, ET)：在不同温度湿度和风速的综合作用下人体产生冷热感觉的指标。\\
校正有效温度（CET）：在有效温度的基础上，综合考虑热辐射的影响，将干球温度改为黑球温度，所得的有效温度。\\
室内小气候的卫生要求：为了保证大多数居民机体的热平衡，有良好的温热感觉，不使机体长期处于紧张状态，保持人体各项生理指标在正常范围内，有正常的学习、工作、休息和睡眠效率。

\InfoLabel{知识来源}
《环境健康学------笔记by徐庆松》第73页；《环境健康学------笔记by徐庆松》第75页；《环境健康学------笔记by徐庆松》第77页

\ContentBand{对应往年题}

\begin{pastquestion}

\QuestionLabel{【往年题1｜03级环境卫生重点考题总结-曾.doc｜口试题】}\\
校正有效温度

\end{pastquestion}

\begin{pastquestion}

\QuestionLabel{【往年题2｜03级环境卫生重点考题总结-曾.doc｜口试题】}\\
室内微小气候构成因素及其对机体健康影响

\end{pastquestion}

\begin{pastquestion}

\QuestionLabel{【往年题3｜03级环境卫生重点考题总结-曾.doc｜口试题】}\\
微小气候卫生学要求

\end{pastquestion}

\begin{pastquestion}

\QuestionLabel{【往年题4｜10级环卫考题回忆.docx｜名词解释】}\\
校正有效温度

\end{pastquestion}

\begin{pastquestion}

\QuestionLabel{【往年题5｜18级环境健康期末考题回忆by崔浩亮.docx｜名词解释】}\\
有效温度

\end{pastquestion}

\begin{pastquestion}

\QuestionLabel{【往年题6｜21级本科预防二班+胡润泽+环境卫生+主观题总结.docx｜简答题】}\\
室内热环境的定义和主要影响因素

\end{pastquestion}

\PointSeparator

\section{室内环境影响健康的基本特点与住宅基本卫生要求}\label{users__fpb__26_spring__pkusph_study__tmp__pdfs__prevention-medicine-past-exams__04-ux73afux5883ux5065ux5eb7ux5b66ux5f80ux5e74ux9898ux8003ux70b9ux6574ux7406.md__ux5ba4ux5185ux73afux5883ux5f71ux54cdux5065ux5eb7ux7684ux57faux672cux7279ux70b9ux4e0eux4f4fux5b85ux57faux672cux536bux751fux8981ux6c42}

\InfoLabel{考频} 2 次

\InfoLabel{对应小节} 一、概述；二、住宅设计与健康

\InfoLabel{匹配依据}
题干涉及住宅基本卫生要求、室内环境影响健康基本特点、住宅设计健康学意义，匹配住宅与健康概述和设计。

\ContentBand{知识点原文摘取}

室内环境对健康影响的基本特点：室内小气候适宜，也有利于其它生物繁殖；空间狭小，有害因素较难扩散稀释；多因素综合作用；慢性长期影响。\\
创造健康室内环境的基本原则：充分引进室外有利因素；充分发挥室内有利因素；尽量避免室外有害因素进入室内；尽量避免室内产生有害因素。\\
住宅的健康学意义：良好的住宅环境有利于人体健康；不良住宅环境不利于人体健康；住宅卫生状况可影响众多家庭成员甚至数代人的健康；住宅环境对健康的影响具有长期性和复杂性。

\InfoLabel{知识来源} 《环境健康学------笔记by徐庆松》第70页

\ContentBand{对应往年题}

\begin{pastquestion}

\QuestionLabel{【往年题1｜03级环境卫生重点考题总结-曾.doc｜口试题】}\\
住宅的基本卫生要求

\end{pastquestion}

\begin{pastquestion}

\QuestionLabel{【往年题2｜20级完整试题回忆.docx｜简答题】}\\
室内环境影响健康的基本特点

\end{pastquestion}

\PointSeparator

\section{采光系数与室内光环境}\label{users__fpb__26_spring__pkusph_study__tmp__pdfs__prevention-medicine-past-exams__04-ux73afux5883ux5065ux5eb7ux5b66ux5f80ux5e74ux9898ux8003ux70b9ux6574ux7406.md__ux91c7ux5149ux7cfbux6570ux4e0eux5ba4ux5185ux5149ux73afux5883}

\InfoLabel{考频} 2 次

\InfoLabel{对应小节} 二、住宅设计与健康 / 居室进深与采光照明；六、室内光环境

\InfoLabel{匹配依据} 题干为采光系数名词解释，匹配住宅采光照明评价指标。

\ContentBand{知识点原文摘取}

采光系数(daylight factor)：是指室内工作水平面上(或距窗 1
米处)散射光的照度与室外相同时间空旷无遮光物的地方接受整个天空散射光(全阴天,见不到太阳,但不是雾天)的水平面上照度的百分比(\%)。\\
能反应当地光气候，采光口大小、位置、朝向的情况，以及室外遮光物等有关影响因素，是比较全面的客观指标。

\InfoLabel{知识来源} 《环境健康学------笔记by徐庆松》第72页

\ContentBand{对应往年题}

\begin{pastquestion}

\QuestionLabel{【往年题1｜03级环境卫生重点考题总结-曾.doc｜口试题】}\\
采光系数

\end{pastquestion}

\begin{pastquestion}

\QuestionLabel{【往年题2｜08级环卫考题回忆.docx｜名词解释】}\\
采光系数

\end{pastquestion}

\PointSeparator

\section{住宅声环境与住宅噪声标准}\label{users__fpb__26_spring__pkusph_study__tmp__pdfs__prevention-medicine-past-exams__04-ux73afux5883ux5065ux5eb7ux5b66ux5f80ux5e74ux9898ux8003ux70b9ux6574ux7406.md__ux4f4fux5b85ux58f0ux73afux5883ux4e0eux4f4fux5b85ux566aux58f0ux6807ux51c6}

\InfoLabel{考频} 1 次

\InfoLabel{对应小节} 五、住宅声环境与健康

\InfoLabel{匹配依据} 题干涉及住宅区噪声限值或噪声影响休息睡眠，匹配住宅声环境章节。

\ContentBand{知识点原文摘取}

噪声：不论是物理学角度的乐声或噪声，只要人们主观上不需要的声音都认为是噪声。\\
噪声的影响早期多为可逆性、生理性改变。长期接触强噪声，可导致不可逆性、病理性损伤。\\
卧室、起居室在关窗状态下白天允许噪声级为 50dB（A 声级），夜间允许噪声级为 40dB（A 声级）。

\InfoLabel{知识来源} 《环境健康学------笔记by徐庆松》第78页

\ContentBand{对应往年题}

\begin{pastquestion}

\QuestionLabel{【往年题1｜2021级环境健康学考题yys.docx｜选择】}\\
居民居住区规定噪音不能超过?分贝

\end{pastquestion}

\PointSeparator

\chapter{公共场所与健康}\label{users__fpb__26_spring__pkusph_study__tmp__pdfs__prevention-medicine-past-exams__04-ux73afux5883ux5065ux5eb7ux5b66ux5f80ux5e74ux9898ux8003ux70b9ux6574ux7406.md__ux7b2cux516bux7ae0-ux516cux5171ux573aux6240ux4e0eux5065ux5eb7}

\section{公共场所概念、分类、健康学特点与基本卫生要求}\label{users__fpb__26_spring__pkusph_study__tmp__pdfs__prevention-medicine-past-exams__04-ux73afux5883ux5065ux5eb7ux5b66ux5f80ux5e74ux9898ux8003ux70b9ux6574ux7406.md__ux516cux5171ux573aux6240ux6982ux5ff5ux5206ux7c7bux5065ux5eb7ux5b66ux7279ux70b9ux4e0eux57faux672cux536bux751fux8981ux6c42}

\InfoLabel{考频} 4 次

\InfoLabel{对应小节} 一、公共场所定义及分类；三、公共场所的健康学要求

\InfoLabel{匹配依据} 题干涉及公共场所名解、分类、卫生学特点、基本卫生要求，匹配第八章公共场所定义及卫生要求。

\ContentBand{知识点原文摘取}

公共场所（public
place）：根据公众生活活动和社会活动的需要，人工建成的具有多种服务功能的公共建筑设施，供公众进行学习、休息、文体活动、交流、交际、购物、美容等活动之用。\\
公共场所的健康学特点：人群密集，流动性大，传播途径多，容易感染；设备及物品易被污染；涉及面广、各类场所特征不同；从业人员流动性大，素质参差不齐。\\
基本卫生要求：良好的环境；良好的小气候；良好的空气质量；公共用品用具清洁卫生，卫生设施运转正常；从业人员必须身体健康并具备基本卫生知识。

\InfoLabel{知识来源}
《环境健康学------笔记by徐庆松》第80页；《环境健康学------笔记by徐庆松》第81页；《环境健康学------笔记by徐庆松》第82页

\ContentBand{对应往年题}

\begin{pastquestion}

\QuestionLabel{【往年题1｜03级环境卫生重点考题总结-曾.doc｜口试题】}\\
公共场所的卫生学特点及分类

\end{pastquestion}

\begin{pastquestion}

\QuestionLabel{【往年题2｜08级环卫考题回忆.docx｜简答】}\\
公共卫生场所的基本卫生要求

\end{pastquestion}

\begin{pastquestion}

\QuestionLabel{【往年题3｜11级环卫考题回忆.docx｜简答】}\\
公共场所卫生学特点

\end{pastquestion}

\begin{pastquestion}

\QuestionLabel{【往年题4｜14级环卫考题回忆by 小川川.pdf｜名词解释】}\\
公共场所

\end{pastquestion}

\PointSeparator

\section{公共场所从业人员服务禁忌疾病（未完全匹配到）}\label{users__fpb__26_spring__pkusph_study__tmp__pdfs__prevention-medicine-past-exams__04-ux73afux5883ux5065ux5eb7ux5b66ux5f80ux5e74ux9898ux8003ux70b9ux6574ux7406.md__ux516cux5171ux573aux6240ux4eceux4e1aux4ebaux5458ux670dux52a1ux7981ux5fccux75beux75c5ux672aux5b8cux5168ux5339ux914dux5230}

\InfoLabel{考频} 2 次

\InfoLabel{对应小节} 未完全匹配到原文；根据题意暂归入：第八章 公共场所与健康

\InfoLabel{匹配依据}
题干问``哪些疾病不得从事直接为顾客服务的工作''。资料中仅有``从业人员必须身体健康并具备基本卫生知识''，未找到具体疾病清单。

\ContentBand{知识点原文摘取}

未完全匹配到原文。

\InfoLabel{知识来源} 未完全匹配到

\ContentBand{对应往年题}

\begin{pastquestion}

\QuestionLabel{【往年题1｜10级环卫考题回忆.docx｜选择】}\\
公共场所卫生管理条件实施细则中，哪些疾病不得从事直接为顾客服务的工作错误的是：乙肝

\end{pastquestion}

\begin{pastquestion}

\QuestionLabel{【往年题2｜14级环卫考题回忆by 小川川.pdf｜选择回忆】}\\
哪些人不能从事直接为客人服务的工作

\end{pastquestion}

\PointSeparator

\section{公共场所环境污染、公共用品污染与军团病}\label{users__fpb__26_spring__pkusph_study__tmp__pdfs__prevention-medicine-past-exams__04-ux73afux5883ux5065ux5eb7ux5b66ux5f80ux5e74ux9898ux8003ux70b9ux6574ux7406.md__ux516cux5171ux573aux6240ux73afux5883ux6c61ux67d3ux516cux5171ux7528ux54c1ux6c61ux67d3ux4e0eux519bux56e2ux75c5}

\InfoLabel{考频} 1 次

\InfoLabel{对应小节} 二、公共场所影响健康的环境因素

\InfoLabel{匹配依据}
题干涉及公共场所环境因素、室内/公共场所空气水污染、军团菌病，匹配公共场所影响健康的环境因素。

\ContentBand{知识点原文摘取}

公共场所空气污染：化学性：颗粒物、二氧化碳、甲醛、苯、总挥发性有机物、烟草烟雾、氨；生物性：细菌、病毒、真菌、病媒生物、致病性微生物；放射性：氡；物理性：气温、气湿、气流、热辐射、新风量。\\
公共场所水污染：公共浴池和一些温泉浴池水易受污染；在游泳过程中，游泳者汗液、尿液的排出和皮肤污垢进入池水，导致水质污染，引起介水传染病。\\
集中空调通风系统污染：风管表面积尘；空调系统冷却塔污染------军团病 (Legionnaires' disease)。\\
公共用品用具污染：公共场所人群密集，流动性大，设备和物品供公众长期反复使用，极易造成致病微生物污染。

\InfoLabel{知识来源} 《环境健康学------笔记by徐庆松》第82页

\ContentBand{对应往年题}

\begin{pastquestion}

\QuestionLabel{【往年题1｜03级环境卫生重点考题总结-曾.doc｜口试题】}\\
军团菌病

\end{pastquestion}

\PointSeparator

\chapter{城乡规划与健康}\label{users__fpb__26_spring__pkusph_study__tmp__pdfs__prevention-medicine-past-exams__04-ux73afux5883ux5065ux5eb7ux5b66ux5f80ux5e74ux9898ux8003ux70b9ux6574ux7406.md__ux7b2cux4e5dux7ae0-ux57ceux4e61ux89c4ux5212ux4e0eux5065ux5eb7}

\section{城市环境容量（未完全匹配到）}\label{users__fpb__26_spring__pkusph_study__tmp__pdfs__prevention-medicine-past-exams__04-ux73afux5883ux5065ux5eb7ux5b66ux5f80ux5e74ux9898ux8003ux70b9ux6574ux7406.md__ux57ceux5e02ux73afux5883ux5bb9ux91cfux672aux5b8cux5168ux5339ux914dux5230}

\InfoLabel{考频} 3 次

\InfoLabel{对应小节} 未完全匹配到原文；根据题意暂归入：第九章 城乡规划与健康

\InfoLabel{匹配依据}
题库出现``城市环境容量''。资料中有土壤环境容量定义，但未找到城市环境容量定义，故保留并标注未完全匹配。

\ContentBand{知识点原文摘取}

未完全匹配到原文。

\InfoLabel{知识来源} 未完全匹配到

\ContentBand{对应往年题}

\begin{pastquestion}

\QuestionLabel{【往年题1｜04级环卫考题回忆.doc｜名词解释】}\\
环境容量

\end{pastquestion}

\begin{pastquestion}

\QuestionLabel{【往年题2｜2021级环境健康学考题yys.docx｜名词解释】}\\
城市环境容量

\end{pastquestion}

\begin{pastquestion}

\QuestionLabel{【往年题3｜21级本科预防二班+胡润泽+环境卫生+主观题总结.docx｜名词解释】}\\
城市环境容量

\end{pastquestion}

\PointSeparator

\section{城乡规划、城市功能分区、人居环境与卫生防护距离}\label{users__fpb__26_spring__pkusph_study__tmp__pdfs__prevention-medicine-past-exams__04-ux73afux5883ux5065ux5eb7ux5b66ux5f80ux5e74ux9898ux8003ux70b9ux6574ux7406.md__ux57ceux4e61ux89c4ux5212ux57ceux5e02ux529fux80fdux5206ux533aux4ebaux5c45ux73afux5883ux4e0eux536bux751fux9632ux62a4ux8dddux79bb}

\InfoLabel{考频} 2 次

\InfoLabel{对应小节} 一、概述；二、城市规划与健康

\InfoLabel{匹配依据}
题干涉及人居环境建设原则、卫生防护距离、城市功能分区、城市规划与健康，匹配第九章城乡规划。

\ContentBand{知识点原文摘取}

城乡规划：为了实现一定时期内城市、村庄和集镇的经济和社会发展目标，确定城市、村庄和集镇的性质、规模和发展方向，合理利用城乡土地，协调城乡空间布局和各项建设的综合部署和具体安排。\\
城乡规划与健康：在城乡规划中贯彻可持续发展战略和以人为本的指导思想，以维持和恢复生态系统为宗旨，以人类与自然环境的和谐共处为目标，建立优良的人居环境，以获得人类生存所需的最佳环境质量。\\
城市规划卫生的基本原则：确定城市性质，控制城市规模；远期规划与近期规划相结合，总体规划与详细规划相结合；保护城市生态环境；维护城市文脉，改善景观环境；加强安全防患，促进人际交往。\\
城市功能分区的卫生原则：合理配置各功能区；居住用地选择；工业用地选择；预留发展余地；分区选择同时进行。\\
卫生防护距离：是指产生有害因素车间的边界至居住区边界的最小距离。

\InfoLabel{知识来源} 《环境健康学------笔记by徐庆松》第83页；《环境健康学------笔记by徐庆松》第84页

\ContentBand{对应往年题}

\begin{pastquestion}

\QuestionLabel{【往年题1｜20级完整试题回忆.docx｜单选题14】}\\
14、人居环境建设的基本原则不包括\\
A、生态原则\\
B、经济原则\\
C、技术原则\\
D、舒适原则\\
E、社会原则

\end{pastquestion}

\begin{pastquestion}

\QuestionLabel{【往年题2｜21级本科预防二班+胡润泽+环境卫生+主观题总结.docx｜名词解释】}\\
卫生防护距离

\end{pastquestion}

\PointSeparator

\section{村镇规划原则与功能分区卫生学要求}\label{users__fpb__26_spring__pkusph_study__tmp__pdfs__prevention-medicine-past-exams__04-ux73afux5883ux5065ux5eb7ux5b66ux5f80ux5e74ux9898ux8003ux70b9ux6574ux7406.md__ux6751ux9547ux89c4ux5212ux539fux5219ux4e0eux529fux80fdux5206ux533aux536bux751fux5b66ux8981ux6c42}

\InfoLabel{考频} 1 次

\InfoLabel{对应小节} 三、村镇规划与健康

\InfoLabel{匹配依据} 题干直接问乡村/村镇各功能分区卫生学要求，匹配村镇规划章节。

\ContentBand{知识点原文摘取}

村镇规划的原则：从农村实际出发，尊重村民意愿，体现地方和农村特色；全面规划、合理布局、节约用地、统筹安排、有利于可持续发展。\\
功能分区的卫生学要求：居住区：村镇自然条件和卫生条件最好的地段；工业副业区：位于当地主导风向的下风侧、河流的下游，与居住区之间设置卫生防护距离；饲养区：配置在居民点的外围，居住区的下风侧和河流下游；农业生产区：避免对居住区以及学校和医院等形成干扰，避免农业生产过程造成的污染。

\InfoLabel{知识来源} 《环境健康学------笔记by徐庆松》第86页

\ContentBand{对应往年题}

\begin{pastquestion}

\QuestionLabel{【往年题1｜20级完整试题回忆.docx｜简答题】}\\
乡村各功能分区的卫生学要求

\end{pastquestion}

\PointSeparator

\chapter{物理因素与健康}\label{users__fpb__26_spring__pkusph_study__tmp__pdfs__prevention-medicine-past-exams__04-ux73afux5883ux5065ux5eb7ux5b66ux5f80ux5e74ux9898ux8003ux70b9ux6574ux7406.md__ux7b2cux5341ux7ae0-ux7269ux7406ux56e0ux7d20ux4e0eux5065ux5eb7}

\section{噪声定义、健康影响与防控}\label{users__fpb__26_spring__pkusph_study__tmp__pdfs__prevention-medicine-past-exams__04-ux73afux5883ux5065ux5eb7ux5b66ux5f80ux5e74ux9898ux8003ux70b9ux6574ux7406.md__ux566aux58f0ux5b9aux4e49ux5065ux5eb7ux5f71ux54cdux4e0eux9632ux63a7}

\InfoLabel{考频} 1 次

\InfoLabel{对应小节} 二、噪声与健康

\InfoLabel{匹配依据} 题干考查噪声定义、健康影响、噪声说法正误和防控，匹配第十章噪声与健康。

\ContentBand{知识点原文摘取}

噪声：不论是物理学上称为乐声或噪声的声响，只要是人们主观上不需要的声音统称噪声。\\
健康影响：听觉系统和非听觉系统。\\
暂时性听阈位移：听觉适应；听觉疲劳。\\
永久性听阈位移：听力损失；噪声性耳聋。\\
神经系统：头痛、头晕、失眠、多梦、全身乏力、记忆力减退、易怒、烦躁等。\\
对心血管系统的影响：心率加快、心输出量增加、血压升高、心电图 ST 段改变；高血压、心肌梗塞、缺血性心脏病等。\\
防控：控制声源；减弱传播；个体防护；相关法律、法规、标准的制定和实施。

\InfoLabel{知识来源}
《环境健康学------笔记by徐庆松》第87页；《环境健康学------笔记by徐庆松》第88页；《环境健康学------笔记by徐庆松》第89页；《环境健康学------笔记by徐庆松》第90页

\ContentBand{对应往年题}

\begin{pastquestion}

\QuestionLabel{【往年题1｜20级完整试题回忆.docx｜单选题13】}\\
13、关于噪声的说法错误的是\\
A、人们主观上不需要的声音统称噪声\\
B、噪声只会引起听觉受损\\
C、超过50dB（A）的噪声就会影响正常休息和睡眠，65dB（A）以上噪声可干扰普通谈话\\
D\&E、\ldots\ldots{}

\end{pastquestion}

\PointSeparator

\chapter{日用化妆品与健康}\label{users__fpb__26_spring__pkusph_study__tmp__pdfs__prevention-medicine-past-exams__04-ux73afux5883ux5065ux5eb7ux5b66ux5f80ux5e74ux9898ux8003ux70b9ux6574ux7406.md__ux7b2cux5341ux4e00ux7ae0-ux65e5ux7528ux5316ux5986ux54c1ux4e0eux5065ux5eb7}

本章未在提供的往年题中匹配到明确考点。

\chapter{环境健康危险度评价}\label{users__fpb__26_spring__pkusph_study__tmp__pdfs__prevention-medicine-past-exams__04-ux73afux5883ux5065ux5eb7ux5b66ux5f80ux5e74ux9898ux8003ux70b9ux6574ux7406.md__ux7b2cux5341ux4e8cux7ae0-ux73afux5883ux5065ux5eb7ux5371ux9669ux5ea6ux8bc4ux4ef7}

\section{健康风险评估/健康危险度评价概念、应用与四个核心步骤}\label{users__fpb__26_spring__pkusph_study__tmp__pdfs__prevention-medicine-past-exams__04-ux73afux5883ux5065ux5eb7ux5b66ux5f80ux5e74ux9898ux8003ux70b9ux6574ux7406.md__ux5065ux5eb7ux98ceux9669ux8bc4ux4f30ux5065ux5eb7ux5371ux9669ux5ea6ux8bc4ux4ef7ux6982ux5ff5ux5e94ux7528ux4e0eux56dbux4e2aux6838ux5fc3ux6b65ux9aa4}

\InfoLabel{考频} 5 次

\InfoLabel{对应小节} 一、相关概念；二、环境健康危险度评估

\InfoLabel{匹配依据} 题干涉及
HRA、健康风险评价、健康危险度评价基本步骤、可以解决哪些问题、危险度名解，匹配第十二章。

\ContentBand{知识点原文摘取}

健康风险评估：对人群暴露在有害环境里面可能产生的健康危害进行量化分析和总结。\\
风险评估的四个核心步骤：有害物质识别（Hazard Identification）；建立剂量反应关系（Dose-response
Assessment）；暴露评估（Exposure Assessment）；风险表征（Risk Characterization）。\\
风险评估的日益广泛应用：回顾性研究；预测性研究；政策评估。\\
健康危险度评价 Health Risk
Assessment，HRA：按照一定的准则，对有害环境因素作用于特定人群的有害健康效应进行综合定性、定量评价的过程。\\
健康危险度评价的主要特点：健康保护观念的转变；量化环境污染对人体健康的影响。

\InfoLabel{知识来源} 《环境健康学------笔记by徐庆松》第94页；《环境健康学------笔记by徐庆松》第95页

\ContentBand{对应往年题}

\begin{pastquestion}

\QuestionLabel{【往年题1｜03级环境卫生重点考题总结-曾.doc｜口试题】}\\
健康危险度评价可以解决哪些问题？

\end{pastquestion}

\begin{pastquestion}

\QuestionLabel{【往年题2｜03级环境卫生重点考题总结-曾.doc｜口试题】}\\
健康危险度评价的概念、内容

\end{pastquestion}

\begin{pastquestion}

\QuestionLabel{【往年题3｜04级环卫考题回忆.doc｜简答题】}\\
试述环境健康危险度评价的基本步骤

\end{pastquestion}

\begin{pastquestion}

\QuestionLabel{【往年题4｜10级环卫考题回忆.docx｜名词解释】}\\
健康风险评价

\end{pastquestion}

\begin{pastquestion}

\QuestionLabel{【往年题5｜11级环卫考题回忆.docx｜名词解释】}\\
危险度

\end{pastquestion}

\PointSeparator

\section{剂量反应评估、RfD、致癌风险与终生暴露风险}\label{users__fpb__26_spring__pkusph_study__tmp__pdfs__prevention-medicine-past-exams__04-ux73afux5883ux5065ux5eb7ux5b66ux5f80ux5e74ux9898ux8003ux70b9ux6574ux7406.md__ux5242ux91cfux53cdux5e94ux8bc4ux4f30rfdux81f4ux764cux98ceux9669ux4e0eux7ec8ux751fux66b4ux9732ux98ceux9669}

\InfoLabel{考频} 4 次

\InfoLabel{对应小节} 二、环境健康危险度评估 / 建立剂量反应关系；风险表征

\InfoLabel{匹配依据} 题干涉及致癌强度系数、致癌物健康风险、BaP
终身健康风险计算，匹配剂量反应和风险表征公式；年龄敏感度加权未在资料中完整出现。

\ContentBand{知识点原文摘取}

非致癌性物质，存在阈值。\\
最小有害作用剂量（LOAEL）：在剂量反应关系实验中，对照组和暴露组的健康效应有统计意义上的差距的最低剂量。\\
最大无害作用剂量（NOAEL）：在剂量反应关系实验中，对照组和暴露组的健康效应没有统计意义上的差距的最高剂量。\\
非致癌物质的健康效应阈值，称为``参考剂量（Reference Dose,简称 RfD）''。\\
对于致癌物质来说，普遍认为不存在安全阈值。\\
风险表征：估算暴露人群/个体的风险（危险度）。致癌物质：平均日摄入剂量（Average Daily Dose，简称
ADD）；个体终生暴露风险（Lifetime Individual Risk，LIR）：LIR=β×LADD。

\InfoLabel{知识来源}
《环境健康学------笔记by徐庆松》第96页；《环境健康学------笔记by徐庆松》第97页；《环境健康学------笔记by徐庆松》第98页

\ContentBand{对应往年题}

\begin{pastquestion}

\QuestionLabel{【往年题1｜14级环卫考题回忆by 小川川.pdf｜名词解释】}\\
致癌强度系数

\end{pastquestion}

\begin{pastquestion}

\QuestionLabel{【往年题2｜2021级环境健康学考题yys.docx｜应用题】}\\
环境中测得BAP浓度1.67ng/m3,限值是1ng/m3,环境致癌因子是1e\^{}-6，不同年龄的敏感度分别是2岁、2-16岁、\textgreater16岁
10 3 1，评估对于不同年龄人群的终身健康风险。

\end{pastquestion}

\begin{pastquestion}

\QuestionLabel{【往年题3｜20级完整试题回忆.docx｜单选题18】}\\
18、致癌物质的健康风险相关，给了苯并芘的浓度，给了一个什么``ng\ldots\ldots1.1e-6''\\
A、有一定的致癌风险\\
B、有较高的致癌风险\\
C、没有致癌风险\\
D、可接受的健康风险\\
E、无法确定

\end{pastquestion}

\begin{pastquestion}

\QuestionLabel{【往年题4｜21级本科预防二班+胡润泽+环境卫生+主观题总结.docx｜应用题】}\\
给出浓度、单位致癌因子和分年龄别的敏感度，计算终身健康风险，计算流程如下。（注意年龄加权，如果忘记年龄加权，污染物浓度比限值高一点，结果算出来终身风险数量级10-4，就很吓人了）

\end{pastquestion}

\PointSeparator

\section{危险认知/风险认知（未完全匹配到）}\label{users__fpb__26_spring__pkusph_study__tmp__pdfs__prevention-medicine-past-exams__04-ux73afux5883ux5065ux5eb7ux5b66ux5f80ux5e74ux9898ux8003ux70b9ux6574ux7406.md__ux5371ux9669ux8ba4ux77e5ux98ceux9669ux8ba4ux77e5ux672aux5b8cux5168ux5339ux914dux5230}

\InfoLabel{考频} 1 次

\InfoLabel{对应小节} 未完全匹配到原文；根据题意暂归入：第十二章 环境健康危险度评价

\InfoLabel{匹配依据} 题库出现``危险认知''。资料中未找到该术语定义，故保留并标注未完全匹配。

\ContentBand{知识点原文摘取}

未完全匹配到原文。

\InfoLabel{知识来源} 未完全匹配到

\ContentBand{对应往年题}

\begin{pastquestion}

\QuestionLabel{【往年题1｜08级环卫考题回忆.docx｜名词解释】}\\
危险认知

\end{pastquestion}

\PointSeparator

\chapter{环境质量与影响评价}\label{users__fpb__26_spring__pkusph_study__tmp__pdfs__prevention-medicine-past-exams__04-ux73afux5883ux5065ux5eb7ux5b66ux5f80ux5e74ux9898ux8003ux70b9ux6574ux7406.md__ux7b2cux5341ux4e09ux7ae0-ux73afux5883ux8d28ux91cfux4e0eux5f71ux54cdux8bc4ux4ef7}

\section{环境质量、环境质量评价的目的、类型、内容与基本要素}\label{users__fpb__26_spring__pkusph_study__tmp__pdfs__prevention-medicine-past-exams__04-ux73afux5883ux5065ux5eb7ux5b66ux5f80ux5e74ux9898ux8003ux70b9ux6574ux7406.md__ux73afux5883ux8d28ux91cfux73afux5883ux8d28ux91cfux8bc4ux4ef7ux7684ux76eeux7684ux7c7bux578bux5185ux5bb9ux4e0eux57faux672cux8981ux7d20}

\InfoLabel{考频} 9 次

\InfoLabel{对应小节} 一、环境质量评价概述

\InfoLabel{匹配依据}
题干涉及环境质量、环境质量标准、环境质量评价类型/内容/要素、环境质量调查方法要素，匹配第十三章环境质量评价概述。

\ContentBand{知识点原文摘取}

环境质量：以健康为准绳衡量环境各要素的优劣程度。一般是指在一个具体的环境内，环境总体或环境的某些要素对人群的生存和繁衍以及社会经济发展的适宜程度。\\
环境质量评价：是认识研究环境的一种科学方法，是对环境质量优劣的定量描述。\\
环境质量评价的目的：掌握和比较环境质量状况及其变化趋势；寻找污染治理重点；为环境综合整治和城市规划及环境规划提供依据；研究环境质量与人群健康的关系；预测和评价规划或建设项目对周围环境可能产生的影响，即环境影响评价。\\
环境质量评价基本要素：监测数据；评价参数（监测指标，评价因子）；评价标准；评价权重；评价模型。

\InfoLabel{知识来源}
《环境健康学------笔记by徐庆松》第98页；《环境健康学------笔记by徐庆松》第100页；《环境健康学------笔记by徐庆松》第101页

\ContentBand{对应往年题}

\begin{pastquestion}

\QuestionLabel{【往年题1｜03级环境卫生重点考题总结-曾.doc｜口试题】}\\
环境质量评价的种类和内容

\end{pastquestion}

\begin{pastquestion}

\QuestionLabel{【往年题2｜10级环卫考题回忆.docx｜简答】}\\
环境质量影响依据区域或范围的分类

\end{pastquestion}

\begin{pastquestion}

\QuestionLabel{【往年题3｜11级环卫考题回忆.docx｜简答】}\\
环境质量调查方法的要素

\end{pastquestion}

\begin{pastquestion}

\QuestionLabel{【往年题4｜12级环境卫生学考试题回忆版（不全）.docx｜名词解释】}\\
环境质量标准

\end{pastquestion}

\begin{pastquestion}

\QuestionLabel{【往年题5｜14级环卫考题回忆by 小川川.pdf｜简答】}\\
环境质量调查评价的主要原理

\end{pastquestion}

\begin{pastquestion}

\QuestionLabel{【往年题6｜18级环境健康期末考题回忆by崔浩亮.docx｜名词解释】}\\
环境质量

\end{pastquestion}

\begin{pastquestion}

\QuestionLabel{【往年题7｜18级环境健康期末考题回忆by崔浩亮.docx｜简答题】}\\
环境质量评价有哪几大类型？

\end{pastquestion}

\begin{pastquestion}

\QuestionLabel{【往年题8｜20级完整试题回忆.docx｜简答题】}\\
环境质量评价方法的基本要素有哪些

\end{pastquestion}

\begin{pastquestion}

\QuestionLabel{【往年题9｜21级本科预防二班+胡润泽+环境卫生+主观题总结.docx｜简答题】}\\
环境质量评价方法的基本要素

\end{pastquestion}

\PointSeparator

\section{环境质量指数、等标污染负荷与排毒系数}\label{users__fpb__26_spring__pkusph_study__tmp__pdfs__prevention-medicine-past-exams__04-ux73afux5883ux5065ux5eb7ux5b66ux5f80ux5e74ux9898ux8003ux70b9ux6574ux7406.md__ux73afux5883ux8d28ux91cfux6307ux6570ux7b49ux6807ux6c61ux67d3ux8d1fux8377ux4e0eux6392ux6bd2ux7cfbux6570}

\InfoLabel{考频} 5 次

\InfoLabel{对应小节} 一、环境质量评价概述 / 环境质量评价内容与方法

\InfoLabel{匹配依据} 题干涉及环境质量指数、等标污染率/等标污染负荷、排毒系数、污染源综合评价方法、AQI
计算，匹配环境质量指数和污染物综合评价。

\ContentBand{知识点原文摘取}

单一污染物的评价：采用污染物排放的相对含量、绝对含量、超标率、超标倍数、检出率、标准差等，来评价污染物和污染源的强度。\\
污染物综合评价：等标污染负荷法/排毒系数法。\\
等标污染负荷：可评价各污染源和各污染物的相对危害程度。物理概念是：把 i
污染物的排放量稀释到其相应排放标准时所需的介质量。\\
等标污染负荷分担率：污染源等标污染负荷分担率 =
某污染源的等标污染负荷/工厂或区域所有污染物总等标污染负荷；污染物的等标污染负荷分担率 = 某污染物的等标污染负荷
/ 工厂或区域所有污染物总等标污染负荷。\\
排毒系数是表示污染源排放的各种污染物对人群健康潜在危害程度的相对指标，它体现了污染物的排放数量和毒性。\\
环境质量指数法：将大量监测数据经统计处理后求得其代表值，以环境卫生标准（或环境质量标准）作为评价标准，换算成定量和客观地评价环境质量的无量纲数值。

\InfoLabel{知识来源}
《环境健康学------笔记by徐庆松》第99页；《环境健康学------笔记by徐庆松》第100页；《环境健康学------笔记by徐庆松》第101页；《环境健康学------笔记by徐庆松》第102页

\ContentBand{对应往年题}

\begin{pastquestion}

\QuestionLabel{【往年题1｜03级环境卫生重点考题总结-曾.doc｜名词解释】}\\
等标污染率

\end{pastquestion}

\begin{pastquestion}

\QuestionLabel{【往年题2｜03级环境卫生重点考题总结-曾.doc｜口试题】}\\
排毒系数、分担率

\end{pastquestion}

\begin{pastquestion}

\QuestionLabel{【往年题3｜04级环卫考题回忆.doc｜选择】}\\
环境质量指数是指：环境监测数据与环境标准的比值

\end{pastquestion}

\begin{pastquestion}

\QuestionLabel{【往年题4｜20级完整试题回忆.docx｜单选题15】}\\
15、下列哪项为对污染源的综合评价方法\\
A、等标污染负荷\\
B、超标率\\
C、超标倍数\\
D、检出率\\
E、标准差

\end{pastquestion}

\begin{pastquestion}

\QuestionLabel{【往年题5｜20级完整试题回忆.docx｜单选题17】}\\
17、给出PM10，PM2.5，SO₂，NO₂，O₃的空气质量分指数，问该区域或城市的空气质量指数AQI

\end{pastquestion}

\PointSeparator

\section{环境影响评价、人群健康影响评价与研究方法}\label{users__fpb__26_spring__pkusph_study__tmp__pdfs__prevention-medicine-past-exams__04-ux73afux5883ux5065ux5eb7ux5b66ux5f80ux5e74ux9898ux8003ux70b9ux6574ux7406.md__ux73afux5883ux5f71ux54cdux8bc4ux4ef7ux4ebaux7fa4ux5065ux5eb7ux5f71ux54cdux8bc4ux4ef7ux4e0eux7814ux7a76ux65b9ux6cd5}

\InfoLabel{考频} 5 次

\InfoLabel{对应小节} 二、环境质量的影响评价

\InfoLabel{匹配依据}
题干涉及环境健康影响评价基本方法、时间序列研究、定组研究、雾霾健康问题研究设计等，匹配人群健康影响评价；``定组研究/panel
study''未完整匹配，但时间序列等方法出现于原文。

\ContentBand{知识点原文摘取}

环境影响评价：指对规划和建设项目实施后可能造成的环境影响进行分析、预测和评估，提出预防或者减轻不良环境影响的对策和措施，进行跟踪监测的方法与制度。\\
人群健康影响评价：人群选择：暴露人群和对照人群；健康效应指标选择：敏感的生理、生化及免疫指标；疾病或死亡指标；方法选择：横断面调查、病例对照研究、队列研究、时间序列分析、Meta
分析。\\
环境影响评价报告书的基本内容：拟建项目概况；拟建项目周围地区的自然、社会、经济以及人群健康等状况；拟建项目对周围地区环境及人群健康影响的预测和评价。

\InfoLabel{知识来源} 《环境健康学------笔记by徐庆松》第102页；《环境健康学------笔记by徐庆松》第103页

\ContentBand{对应往年题}

\begin{pastquestion}

\QuestionLabel{【往年题1｜11级环卫考题回忆.docx｜应用】}\\
设计一个时序研究，目的是发现雾霾和健康问题的关系

\end{pastquestion}

\begin{pastquestion}

\QuestionLabel{【往年题2｜12级环境卫生学考试题回忆版（不全）.docx｜简答】}\\
环境健康影响评价的基本方法

\end{pastquestion}

\begin{pastquestion}

\QuestionLabel{【往年题3｜18级环境健康期末考题回忆by崔浩亮.docx｜名词解释】}\\
定组研究

\end{pastquestion}

\begin{pastquestion}

\QuestionLabel{【往年题4｜2021级环境健康学考题yys.docx｜选择】}\\
选择不正确说法：A时间序列收集常见检测资料 E 环境卫生学中时间序列研究可用于长时间暴露

\end{pastquestion}

\begin{pastquestion}

\QuestionLabel{【往年题5｜20级完整试题回忆.docx｜名词解释】}\\
panel study

\end{pastquestion}

\PointSeparator

\SetSubjectAccent{C25A26}{营养与食品卫生学}

\part{营养与食品卫生学}

\chapter*{说明}\label{users__fpb__26_spring__pkusph_study__tmp__pdfs__prevention-medicine-past-exams__05-ux8425ux517bux5b66ux5f80ux5e74ux9898ux8003ux70b9ux6574ux7406.md__ux8bf4ux660e}
\addcontentsline{toc}{chapter}{说明}

\begin{itemize}
\tightlist
\item
  本文件根据《营养与食品卫生学·18 柯雅蕾笔记》Markdown 和用户提供的「5 往年题」文件夹中全部往年题整理。
\item
  往年题来源共 7 份文件：04 级考题回忆.doc（含「04 营养真题」与「其他年份回忆」两部分）、05
  级考题回忆.docx、08 级考题回忆.doc、09 级考题回忆.docx、14 级营养考题.pdf、20 级完整试题回忆.docx、21
  级营养期末考题.docx。
\item
  本次将往年题拆分为 175 个题目条目；每个条目均已归入至少一个考点（含标记为「未完全匹配到」者）。
\item
  知识点正文均摘自《营养与食品卫生学·18
  柯雅蕾笔记》原文；笔记中「老年营养」「临床营养」「食品安全风险分析与控制」三章在原文件中重复出现三次，本文统一采用内容最完整的版本。未找到明确原文的题目已保留并标记为「未完全匹配到」。
\item
  排序依据为：先按知识总结文件章节顺序，再在同一章节内按考频从高到低排列；同频时按原文出现顺序排列。
\item
  一道考题若同时考查多个知识点，可分别归入多个考点，并各计 1 次。
\item
  如果某道题未在知识总结文件中完全匹配到对应知识点，则标记为「未完全匹配到」，并根据题意暂归入最合适章节；个别完全无对应考点者归入文末「附」。
\item
  已反向核对题库拆分条目，确保最终文件中没有题目条目遗漏（详见文末「题量核对」）。
\end{itemize}

\PointSeparator

\chapter{绪论}\label{users__fpb__26_spring__pkusph_study__tmp__pdfs__prevention-medicine-past-exams__05-ux8425ux517bux5b66ux5f80ux5e74ux9898ux8003ux70b9ux6574ux7406.md__ux7b2cux4e00ux7ae0-ux7eeaux8bba}

\section{食品安全与食品卫生的概念}\label{users__fpb__26_spring__pkusph_study__tmp__pdfs__prevention-medicine-past-exams__05-ux8425ux517bux5b66ux5f80ux5e74ux9898ux8003ux70b9ux6574ux7406.md__ux98dfux54c1ux5b89ux5168ux4e0eux98dfux54c1ux536bux751fux7684ux6982ux5ff5}

\InfoLabel{考频} 1 次

\InfoLabel{对应小节} 三、学科发展 /（一）食品安全与食品卫生

\InfoLabel{匹配依据} 题干为名词解释「food safety（食品安全）」，与绪论中食品安全的定义直接对应。

\ContentBand{知识点原文摘取}

食品安全（food
safety）：指食品无毒、无害，符合应当有的营养要求，对人体健康不造成任何急性、亚急性或者慢性危害。\\
食品安全是结果安全与过程安全的完整统一，包括食物的种植（养殖）、加工、包装、储存、运输、销售、消费等环节的安全；食品安全包含食品卫生。

\InfoLabel{知识来源} 《营养与食品卫生学·18 柯雅蕾笔记》第 2-3 页

\ContentBand{对应往年题}

\begin{pastquestion}

\QuestionLabel{【往年题1｜14级营养考题.pdf｜名词解释】}\\
food safety

\end{pastquestion}

\PointSeparator

\chapter{公共营养}\label{users__fpb__26_spring__pkusph_study__tmp__pdfs__prevention-medicine-past-exams__05-ux8425ux517bux5b66ux5f80ux5e74ux9898ux8003ux70b9ux6574ux7406.md__ux7b2cux4e8cux7ae0-ux516cux5171ux8425ux517b}

\section{膳食营养素参考摄入量（DRIs）及其指标体系}\label{users__fpb__26_spring__pkusph_study__tmp__pdfs__prevention-medicine-past-exams__05-ux8425ux517bux5b66ux5f80ux5e74ux9898ux8003ux70b9ux6574ux7406.md__ux81b3ux98dfux8425ux517bux7d20ux53c2ux8003ux6444ux5165ux91cfdrisux53caux5176ux6307ux6807ux4f53ux7cfb}

\InfoLabel{考频} 3 次

\InfoLabel{对应小节} 二、膳食营养素参考摄入量 /（三）膳食营养素参考摄入量（DRIs）定义及概念

\InfoLabel{匹配依据} 题干涉及 DRIs（膳食营养素参考摄入量）、RNI 的定义及 EAR/AI/RNI/UL
的相互关系，与本节直接对应。

\ContentBand{知识点原文摘取}

膳食营养素参考摄入量（DRIs）是一组每日平均膳食营养素摄入量的参考值。\\
平均需要量（EAR）：是群体中各个体需要量的平均值，只能满足该群体中 50\%成员的需要。\\
推荐摄入量（RNI）：可以满足某一特定性别、年龄及生理状况群体中绝大多数（97\%-98\%）个体需要量的摄入水平。\\
适宜摄入量（AI）：通过观察或实验获得的健康人群某种营养素的摄入量。\\
可耐受最高摄入量（UL）：平均每日可以摄入该营养素的最高量。

\InfoLabel{知识来源} 《营养与食品卫生学·18 柯雅蕾笔记》第 6 页

\ContentBand{对应往年题}

\begin{pastquestion}

\QuestionLabel{【往年题1｜04级考题回忆.doc（附：其他年份回忆）｜名词解释】}\\
RNI

\end{pastquestion}

\begin{pastquestion}

\QuestionLabel{【往年题2｜08级考题回忆.doc｜名词解释】}\\
膳食营养参考摄入量

\end{pastquestion}

\begin{pastquestion}

\QuestionLabel{【往年题3｜21级营养期末考题.docx｜选择】}\\
选择正确的一项：B「RNI
推荐摄入量：可以满足某一特定性别、年龄及生理状况群体中绝大多数（97\%-98\%）个体需要量的摄入水平」；E「EAR\textless AL\textless RNI\textless UL」

\end{pastquestion}

\PointSeparator

\section{我国公共营养现状与主要问题}\label{users__fpb__26_spring__pkusph_study__tmp__pdfs__prevention-medicine-past-exams__05-ux8425ux517bux5b66ux5f80ux5e74ux9898ux8003ux70b9ux6574ux7406.md__ux6211ux56fdux516cux5171ux8425ux517bux73b0ux72b6ux4e0eux4e3bux8981ux95eeux9898}

\InfoLabel{考频} 2 次

\InfoLabel{对应小节} 公共营养·一、概述；绪论·一、概述（部分匹配）

\InfoLabel{匹配依据}
题干为「我国公众营养现状和存在的主要问题」。笔记给出公共营养的定义、工作内容及膳食因素疾病负担，但未列出系统的「现状---问题」清单，故为部分匹配。

\ContentBand{知识点原文摘取}

公共营养：基于人群营养状况，有针对性地提出解决营养问题的措施\ldots\ldots 如营养不良的消除策略、政策与措施等。\\
2017 年全球疾病负担分析显示，膳食因素分别占全死因的 22\%和伤残调整生命年（DALYs）的
15\%，主要是高钠摄入、全谷物不足和水果摄入不足；我国的膳食因素「贡献」更大，分别为 30.2\%（死亡）和
21.3\%（DALYs）。\\
营养学正面临着营养缺乏和营养过剩的双重挑战。

\InfoLabel{知识来源} 《营养与食品卫生学·18 柯雅蕾笔记》第 1 页、第 5 页

\ContentBand{对应往年题}

\begin{pastquestion}

\QuestionLabel{【往年题1｜04级考题回忆.doc｜简答】}\\
公众营养现状和存在的主要问题

\end{pastquestion}

\begin{pastquestion}

\QuestionLabel{【往年题2｜09级考题回忆.docx｜简答】}\\
我国公众营养现状

\end{pastquestion}

\PointSeparator

\section{食物营养改善与食品（营养）强化}\label{users__fpb__26_spring__pkusph_study__tmp__pdfs__prevention-medicine-past-exams__05-ux8425ux517bux5b66ux5f80ux5e74ux9898ux8003ux70b9ux6574ux7406.md__ux98dfux7269ux8425ux517bux6539ux5584ux4e0eux98dfux54c1ux8425ux517bux5f3aux5316}

\InfoLabel{考频} 4 次

\InfoLabel{对应小节} 五、食物营养规划与营养改善；食品添加剂·营养强化剂（GB14880）

\InfoLabel{匹配依据} 题干涉及「食品强化（food
fortification）」名词及「公共营养不良的迅速改善方法＝食品强化」。笔记以「营养强化剂」与「营养干预措施」相对应，无单列「食品强化」定义，故为部分匹配。

\ContentBand{知识点原文摘取}

营养强化剂（GB14880）：为增强营养成分而加入食品中的天然或人工合成的属于天然营养素范畴的食品添加剂。\\
营养干预项目（FAO 列举）：儿童辅助食品、孕妇辅助食品、营养教育、社区营养监测、营养康复中心等。\\
选择具体营养干预措施的步骤：①确定营养不良的高危人群；②优先选择干预措施的标准\ldots\ldots{}

\InfoLabel{知识来源} 《营养与食品卫生学·18 柯雅蕾笔记》第 8-9 页、第 76-77 页

\ContentBand{对应往年题}

\begin{pastquestion}

\QuestionLabel{【往年题1｜04级考题回忆.doc｜名词解释】}\\
食品强化（food fortification）

\end{pastquestion}

\begin{pastquestion}

\QuestionLabel{【往年题2｜09级考题回忆.docx｜名词解释】}\\
食物强化

\end{pastquestion}

\begin{pastquestion}

\QuestionLabel{【往年题3｜14级营养考题.pdf｜名词解释】}\\
Food fortification

\end{pastquestion}

\begin{pastquestion}

\QuestionLabel{【往年题4｜08级考题回忆.doc｜选择】}\\
公共营养不良的迅速改变方法：食品强化？

\end{pastquestion}

\PointSeparator

\section{膳食指南（《黄帝内经》最早的「膳食指南」）}\label{users__fpb__26_spring__pkusph_study__tmp__pdfs__prevention-medicine-past-exams__05-ux8425ux517bux5b66ux5f80ux5e74ux9898ux8003ux70b9ux6574ux7406.md__ux81b3ux98dfux6307ux5357ux9ec4ux5e1dux5185ux7ecfux6700ux65e9ux7684ux81b3ux98dfux6307ux5357}

\InfoLabel{考频} 1 次

\InfoLabel{对应小节} 三、膳食指南（未完全匹配到）

\InfoLabel{匹配依据}
题干考查《黄帝内经》「五谷为养、五果为助、五畜为益、五菜为充」这一最早的膳食指南。笔记膳食指南部分自 1989
年中国居民膳食指南起述，未收录《黄帝内经》内容，故标记为未完全匹配到。

\ContentBand{知识点原文摘取}

未完全匹配到完整原文。\\
可参考原文片段：膳食指南是根据营养学原则，结合国情，教育人民群众采用平衡膳食，以达到合理营养促进健康目的的指导性意见。1989
年中国营养学会提出我国居民膳食指南\ldots\ldots{}

\InfoLabel{知识来源} 《营养与食品卫生学·18 柯雅蕾笔记》第 7 页

\ContentBand{对应往年题}

\begin{pastquestion}

\QuestionLabel{【往年题1｜20级完整试题回忆.docx｜单选】}\\
《黄帝内经》中提出了（）这一最早的「膳食指南」：A
五谷为益、五果为充、五畜为养、五菜为助\ldots\ldots（正确为「五谷为养、五果为助、五畜为益、五菜为充」）

\end{pastquestion}

\PointSeparator

\chapter{蛋白质与健康}\label{users__fpb__26_spring__pkusph_study__tmp__pdfs__prevention-medicine-past-exams__05-ux8425ux517bux5b66ux5f80ux5e74ux9898ux8003ux70b9ux6574ux7406.md__ux7b2cux4e09ux7ae0-ux86cbux767dux8d28ux4e0eux5065ux5eb7}

\section{必需氨基酸（及氨基酸模式、参考蛋白、限制氨基酸）}\label{users__fpb__26_spring__pkusph_study__tmp__pdfs__prevention-medicine-past-exams__05-ux8425ux517bux5b66ux5f80ux5e74ux9898ux8003ux70b9ux6574ux7406.md__ux5fc5ux9700ux6c28ux57faux9178ux53caux6c28ux57faux9178ux6a21ux5f0fux53c2ux8003ux86cbux767dux9650ux5236ux6c28ux57faux9178}

\InfoLabel{考频} 4 次

\InfoLabel{对应小节} 一、概述 / 2. 构成人体蛋白质的氨基酸；3. 氨基酸模式

\InfoLabel{匹配依据} 题干为「必需氨基酸」名词解释及「哪个不是必需氨基酸」的辨析，与本节直接对应。

\ContentBand{知识点原文摘取}

必需氨基酸（essential amino
acid）：是指人体内不能合成或合成速度不能满足机体需要，必须从食物中直接获得的氨基酸。组氨酸为婴儿必需氨基酸。\\
条件必需氨基酸：自身可合成，但不能满足需要。半胱氨酸由蛋氨酸转变而来；酪氨酸由苯丙氨酸转变而来。\\
氨基酸模式：将该蛋白质中的色氨酸含量定为
1，分别计算出其他必需氨基酸的相应比值。鸡蛋蛋白质的氨基酸模式与人体最接近，常作为参考蛋白。

\InfoLabel{知识来源} 《营养与食品卫生学·18 柯雅蕾笔记》第 9-10 页

\ContentBand{对应往年题}

\begin{pastquestion}

\QuestionLabel{【往年题1｜04级考题回忆.doc｜名词解释】}\\
必需氨基酸（全英文）

\end{pastquestion}

\begin{pastquestion}

\QuestionLabel{【往年题2｜08级考题回忆.doc｜名词解释】}\\
必需氨基酸（全英文）

\end{pastquestion}

\begin{pastquestion}

\QuestionLabel{【往年题3｜09级考题回忆.docx｜名词解释】}\\
必需氨基酸（全英文）

\end{pastquestion}

\begin{pastquestion}

\QuestionLabel{【往年题4｜14级营养考题.pdf｜选择】}\\
哪个不是必需氨基酸

\end{pastquestion}

\PointSeparator

\section{蛋白质的营养学评价（消化率、生物价、氨基酸评分）}\label{users__fpb__26_spring__pkusph_study__tmp__pdfs__prevention-medicine-past-exams__05-ux8425ux517bux5b66ux5f80ux5e74ux9898ux8003ux70b9ux6574ux7406.md__ux86cbux767dux8d28ux7684ux8425ux517bux5b66ux8bc4ux4ef7ux6d88ux5316ux7387ux751fux7269ux4ef7ux6c28ux57faux9178ux8bc4ux5206}

\InfoLabel{考频} 4 次

\InfoLabel{对应小节} 四、营养学评价 /（二）蛋白质消化率；（三）蛋白质利用率

\InfoLabel{匹配依据}
题干涉及食物蛋白质营养学评价指标、评价蛋白质利用率的指标、氨基酸评分名词及真消化率计算，与本节直接对应。

\ContentBand{知识点原文摘取}

真消化率（\%）＝{[}食物氮－（粪氮－粪代谢氮）{]}/食物氮×100\%；表观消化率（\%）＝（食物氮－粪氮）/食物氮×100\%。\\
生物价 BV（\%）＝储留氮/吸收氮×100\%；蛋白质净利用率 NPU＝消化率×生物价。\\
蛋白质功效比值 PER＝动物体重增加（g）/摄入食物蛋白质（g）。\\
氨基酸评分（AAS）：用被测食物蛋白质的必需氨基酸评分模式和理想模式（或参考蛋白）进行比较；最低的必需氨基酸（第一限制氨基酸）评分值即为该蛋白质的氨基酸评分。

\InfoLabel{知识来源} 《营养与食品卫生学·18 柯雅蕾笔记》第 11-12 页

\ContentBand{对应往年题}

\begin{pastquestion}

\QuestionLabel{【往年题1｜04级考题回忆.doc（附：其他年份回忆）｜简答】}\\
食物蛋白质的营养学评价指标

\end{pastquestion}

\begin{pastquestion}

\QuestionLabel{【往年题2｜14级营养考题.pdf｜选择】}\\
评价蛋白质利用率的指标

\end{pastquestion}

\begin{pastquestion}

\QuestionLabel{【往年题3｜21级营养期末考题.docx｜名词解释】}\\
Amino acid score（氨基酸评分）

\end{pastquestion}

\begin{pastquestion}

\QuestionLabel{【往年题4｜20级完整试题回忆.docx｜单选/计算】}\\
某人摄入 90g 氮，粪氮为 3.7g，粪代谢氮为 1.2g，计算其蛋白质真消化率

\end{pastquestion}

\PointSeparator

\section{蛋白质生理功能与蛋白质-能量营养不良（PEM）}\label{users__fpb__26_spring__pkusph_study__tmp__pdfs__prevention-medicine-past-exams__05-ux8425ux517bux5b66ux5f80ux5e74ux9898ux8003ux70b9ux6574ux7406.md__ux86cbux767dux8d28ux751fux7406ux529fux80fdux4e0eux86cbux767dux8d28-ux80fdux91cfux8425ux517bux4e0dux826fpem}

\InfoLabel{考频} 4 次

\InfoLabel{对应小节} 二、生理功能；五、营养不良与评价 / 1. 蛋白质营养不良

\InfoLabel{匹配依据}
题干涉及蛋白质生理功能、蛋白质营养不良表现、「营养不良」名词及「Kwashiorker（水肿型/夸希奥科病）」，与本节直接对应。

\ContentBand{知识点原文摘取}

蛋白质生理功能：①人体组织的构成成分；②构成体内各种重要生理活性物质（酶、激素、载体蛋白、抗体等）；③供给能量（16.7kJ／4.0kcal/g）；④生物活性肽功能。\\
蛋白质-能量营养不良（PEM）是由蛋白质和/或能量摄入不足引起的营养缺乏病。\\
水肿型、夸希奥科病（Kwashiorker）：能量摄入基本满足但蛋白质严重不足，表现为腹部、腿部水肿，生长迟缓等。\\
消瘦型（Marasmus）：能量与蛋白质摄入均严重不足，多见于 2 岁以下婴幼儿，患儿多表现为消瘦无力。

\InfoLabel{知识来源} 《营养与食品卫生学·18 柯雅蕾笔记》第 10-14 页

\ContentBand{对应往年题}

\begin{pastquestion}

\QuestionLabel{【往年题1｜14级营养考题.pdf｜论述】}\\
阐述蛋白质的生理功能，蛋白质营养不良的表现

\end{pastquestion}

\begin{pastquestion}

\QuestionLabel{【往年题2｜04级考题回忆.doc｜名词解释】}\\
营养不良（全英文）

\end{pastquestion}

\begin{pastquestion}

\QuestionLabel{【往年题3｜09级考题回忆.docx｜名词解释】}\\
营养不良（全英文）

\end{pastquestion}

\begin{pastquestion}

\QuestionLabel{【往年题4｜04级考题回忆.doc（附：其他年份回忆）｜名词解释】}\\
Kwashiorker

\end{pastquestion}

\PointSeparator

\section{必要氮损失与氮平衡}\label{users__fpb__26_spring__pkusph_study__tmp__pdfs__prevention-medicine-past-exams__05-ux8425ux517bux5b66ux5f80ux5e74ux9898ux8003ux70b9ux6574ux7406.md__ux5fc5ux8981ux6c2eux635fux5931ux4e0eux6c2eux5e73ux8861}

\InfoLabel{考频} 1 次

\InfoLabel{对应小节} 三、消化吸收（氨基酸池、必要的氮损失、氮平衡）

\InfoLabel{匹配依据} 题干为「Obligatory nitrogen losses（必要氮损失）」名词解释，与本节直接对应。

\ContentBand{知识点原文摘取}

必要的氮损失（obligatory nitrogen
losses）：指机体内每天由于皮肤、毛发和黏膜脱落，妇女月经期失血以及肠道菌体死亡排出而造成的不可避免的氮损失，在
20g 以上。\\
氮平衡（nitrogen balance）：机体摄入氮和排出氮的关系，B＝I－（U＋F＋S）。

\InfoLabel{知识来源} 《营养与食品卫生学·18 柯雅蕾笔记》第 11 页

\ContentBand{对应往年题}

\begin{pastquestion}

\QuestionLabel{【往年题1｜20级完整试题回忆.docx｜名词解释】}\\
Obligatory nitrogen losses

\end{pastquestion}

\PointSeparator

\section{蛋白质含量测定（凯氏定氮法）与三聚氰胺案例}\label{users__fpb__26_spring__pkusph_study__tmp__pdfs__prevention-medicine-past-exams__05-ux8425ux517bux5b66ux5f80ux5e74ux9898ux8003ux70b9ux6574ux7406.md__ux86cbux767dux8d28ux542bux91cfux6d4bux5b9aux51efux6c0fux5b9aux6c2eux6cd5ux4e0eux4e09ux805aux6c30ux80faux6848ux4f8b}

\InfoLabel{考频} 2 次

\InfoLabel{对应小节} 四、营养学评价 /（一）蛋白质的含量（氮×6.25）

\InfoLabel{匹配依据}
题干为三聚氰胺造成蛋白质升高假象的案例，考查凯氏定氮法测氮折算蛋白质的原理、缺乏蛋白质的表现、如何从食物获得及如何全面评价该营养素，核心为蛋白质含量测定与营养学评价。

\ContentBand{知识点原文摘取}

蛋白质的含量＝氮的含量×6.25（凯氏定氮法以测得的总氮量折算蛋白质含量）。\\
蛋白质营养学评价还需结合消化率、生物价、氨基酸评分等综合指标。\\
食物蛋白质来源：奶、蛋、鱼、肉中的蛋白质属完全蛋白质，所含必需氨基酸种类齐全、数量充足、比例适当。

\InfoLabel{知识来源} 《营养与食品卫生学·18 柯雅蕾笔记》第 11-12 页

\ContentBand{对应往年题}

\begin{pastquestion}

\QuestionLabel{【往年题1｜04级考题回忆.doc｜应用题】}\\
三聚氰胺造成蛋白质升高假象，其原理是凯氏定氮法测得氮含量高，认为蛋白质含量高，如何全面评价食物中该营养素的营养价值？缺乏该营养素的表现，如何在食物中获得该营养素？

\end{pastquestion}

\begin{pastquestion}

\QuestionLabel{【往年题2｜09级考题回忆.docx｜论述】}\\
三聚氰胺案例：1）加入三聚氰胺增加了哪种营养素升高的假象，原理是什么；2）缺乏这种营养素的表现以及如何从食物中获得；3）如何评价食物中这种营养素

\end{pastquestion}

\PointSeparator

\chapter{脂质与健康}\label{users__fpb__26_spring__pkusph_study__tmp__pdfs__prevention-medicine-past-exams__05-ux8425ux517bux5b66ux5f80ux5e74ux9898ux8003ux70b9ux6574ux7406.md__ux7b2cux56dbux7ae0-ux8102ux8d28ux4e0eux5065ux5eb7}

\section{必需脂肪酸（EFA）}\label{users__fpb__26_spring__pkusph_study__tmp__pdfs__prevention-medicine-past-exams__05-ux8425ux517bux5b66ux5f80ux5e74ux9898ux8003ux70b9ux6574ux7406.md__ux5fc5ux9700ux8102ux80aaux9178efa}

\InfoLabel{考频} 2 次

\InfoLabel{对应小节} 二、分类 / 必需脂肪酸（essential fatty acid，EFA）

\InfoLabel{匹配依据} 题干为「必需脂肪酸／Essential fatty acid」名词解释，与本节直接对应。

\ContentBand{知识点原文摘取}

必需脂肪酸（essential fatty
acid，EFA）：指人体不可缺少同时自身不能合成，必须通过食物供给的脂肪酸，包括亚油酸和 α-亚麻酸。\\
作用：磷脂组成部分、前列腺素合成的前体、参与胆固醇代谢。每日摄入量不少于总能量 3\%。\\
缺乏可引起生长迟缓、皮肤损伤（如鳞屑样皮炎、湿疹）等。

\InfoLabel{知识来源} 《营养与食品卫生学·18 柯雅蕾笔记》第 18 页

\ContentBand{对应往年题}

\begin{pastquestion}

\QuestionLabel{【往年题1｜08级考题回忆.doc｜名词解释】}\\
必需脂肪酸（全英文）

\end{pastquestion}

\begin{pastquestion}

\QuestionLabel{【往年题2｜14级营养考题.pdf｜名词解释】}\\
Essential fatty acid

\end{pastquestion}

\PointSeparator

\section{脂类的分类、生理功能及参考摄入量}\label{users__fpb__26_spring__pkusph_study__tmp__pdfs__prevention-medicine-past-exams__05-ux8425ux517bux5b66ux5f80ux5e74ux9898ux8003ux70b9ux6574ux7406.md__ux8102ux7c7bux7684ux5206ux7c7bux751fux7406ux529fux80fdux53caux53c2ux8003ux6444ux5165ux91cf}

\InfoLabel{考频} 2 次

\InfoLabel{对应小节} 二、分类；三、生理功能；六、参考摄入量

\InfoLabel{匹配依据}
题干为「脂类的分类及功能」简答与「脂肪每日能量比例」选择，与本章分类、生理功能及参考摄入量直接对应。

\ContentBand{知识点原文摘取}

脂类包括脂肪（甘油三酯，占脂类总量 95\%）和类脂（磷脂、固醇类，约占 5\%）。\\
脂肪酸按碳链长度（长/中/短链）、饱和程度（饱和、单不饱和、多不饱和）、空间结构（顺式、反式）分类。\\
生理功能：储存和提供能量、维持正常体温、保护器官、维持细胞结构和功能、内分泌调节；食物脂肪可增加饱腹感、改善感官、提供脂溶性维生素。\\
参考摄入量：成年人脂肪摄入量应控制在总能量的
20\%～30\%；饱和脂肪酸＜10\%；反式脂肪酸＜1\%；必需脂肪酸不少于总能量 3\%。

\InfoLabel{知识来源} 《营养与食品卫生学·18 柯雅蕾笔记》第 17-21 页

\ContentBand{对应往年题}

\begin{pastquestion}

\QuestionLabel{【往年题1｜21级营养期末考题.docx｜简答】}\\
脂类的分类及功能

\end{pastquestion}

\begin{pastquestion}

\QuestionLabel{【往年题2｜08级考题回忆.doc｜选择】}\\
脂肪每日能量比例

\end{pastquestion}

\PointSeparator

\section{反式脂肪酸的来源及对健康的危害}\label{users__fpb__26_spring__pkusph_study__tmp__pdfs__prevention-medicine-past-exams__05-ux8425ux517bux5b66ux5f80ux5e74ux9898ux8003ux70b9ux6574ux7406.md__ux53cdux5f0fux8102ux80aaux9178ux7684ux6765ux6e90ux53caux5bf9ux5065ux5eb7ux7684ux5371ux5bb3}

\InfoLabel{考频} 1 次

\InfoLabel{对应小节} 二、分类（反式脂肪酸）；食品安全风险评估·反式不饱和脂肪酸（TFA）

\InfoLabel{匹配依据} 题干为「反式脂肪酸的主要来源及其对人体健康的危害」，与本节及风险评估案例直接对应。

\ContentBand{知识点原文摘取}

反式脂肪酸（TFA）：指氢原子在双键的相对两侧的脂肪酸。氢化植物油和人造黄油中反式不饱和脂肪酸含量很高，可达
40\%左右：人造奶油、蛋糕、饼干、油炸食品、乳酪食品、花生酱等。\\
自然界中存在的不饱和脂肪酸大部分是顺式结构，牛脂和奶脂中分别有 8\%和 20\%的不饱和脂肪酸是反式的。\\
反式不饱和脂肪酸可以增加动脉粥样硬化性心脏病发生的潜在危险。中国（DRIs，2013）建议 TFA ＜1\%膳食能量。

\InfoLabel{知识来源} 《营养与食品卫生学·18 柯雅蕾笔记》第 18 页、第 89 页

\ContentBand{对应往年题}

\begin{pastquestion}

\QuestionLabel{【往年题1｜20级完整试题回忆.docx｜简答】}\\
试述反式脂肪酸的主要来源及其对人体健康的危害

\end{pastquestion}

\PointSeparator

\section{动脉粥样硬化的膳食控制原则}\label{users__fpb__26_spring__pkusph_study__tmp__pdfs__prevention-medicine-past-exams__05-ux8425ux517bux5b66ux5f80ux5e74ux9898ux8003ux70b9ux6574ux7406.md__ux52a8ux8109ux7ca5ux6837ux786cux5316ux7684ux81b3ux98dfux63a7ux5236ux539fux5219}

\InfoLabel{考频} 1 次

\InfoLabel{对应小节} 五、脂肪与健康（未完全匹配到）

\InfoLabel{匹配依据}
题干为「动脉粥样硬化的膳食控制原则」。笔记述及高脂肪摄入增加肥胖风险、反式脂肪酸与动脉粥样硬化的关系及
DASH/地中海膳食，但未列出系统的「动脉粥样硬化膳食控制原则」，故标记为未完全匹配到。

\ContentBand{知识点原文摘取}

未完全匹配到完整原文。\\
可参考原文片段：高脂肪摄入可增加肥胖风险，减少总脂肪摄入有助于减少体重。\\
饱和脂肪酸摄入应占总能量的 10\%以下；尽量减少反式不饱和脂肪酸，控制在总能量的 1\%以下。\\
DASH
膳食模式：通过增加蔬菜、水果、鱼和低脂食物摄入，减少红肉、饱和脂肪酸和甜食摄入而进行高血压防治的膳食模式。

\InfoLabel{知识来源} 《营养与食品卫生学·18 柯雅蕾笔记》第 20-21 页、第 50 页

\ContentBand{对应往年题}

\begin{pastquestion}

\QuestionLabel{【往年题1｜04级考题回忆.doc（附：其他年份回忆）｜简答】}\\
动脉粥样硬化的膳食控制原则

\end{pastquestion}

\PointSeparator

\chapter{碳水化合物、能量平衡与健康}\label{users__fpb__26_spring__pkusph_study__tmp__pdfs__prevention-medicine-past-exams__05-ux8425ux517bux5b66ux5f80ux5e74ux9898ux8003ux70b9ux6574ux7406.md__ux7b2cux4e94ux7ae0-ux78b3ux6c34ux5316ux5408ux7269ux80fdux91cfux5e73ux8861ux4e0eux5065ux5eb7}

\section{膳食纤维的营养学意义及与慢性病的关系}\label{users__fpb__26_spring__pkusph_study__tmp__pdfs__prevention-medicine-past-exams__05-ux8425ux517bux5b66ux5f80ux5e74ux9898ux8003ux70b9ux6574ux7406.md__ux81b3ux98dfux7ea4ux7ef4ux7684ux8425ux517bux5b66ux610fux4e49ux53caux4e0eux6162ux6027ux75c5ux7684ux5173ux7cfb}

\InfoLabel{考频} 5 次

\InfoLabel{对应小节} 三、生理功能 / 膳食纤维的营养学意义

\InfoLabel{匹配依据} 题干涉及膳食纤维的生理意义/作用及其与哪些慢性病相关，与本节直接对应。

\ContentBand{知识点原文摘取}

膳食纤维的营养学意义：①增加饱腹感（吸水膨胀、延缓胃排空），有利于肥胖和糖尿病患者控制体重；②促进排便、预防便秘；③降低血糖和血胆固醇（减少小肠对糖的吸收、吸附胆固醇和胆汁酸）；④改善肠道菌群（结肠发酵产短链脂肪酸）；⑤降低某些肿瘤（如结肠癌、乳腺癌）的发生风险。\\
膳食纤维每日 25～30g 为宜。

\InfoLabel{知识来源} 《营养与食品卫生学·18 柯雅蕾笔记》第 23-24 页

\ContentBand{对应往年题}

\begin{pastquestion}

\QuestionLabel{【往年题1｜04级考题回忆.doc｜简答】}\\
膳食纤维与哪些慢性病有关

\end{pastquestion}

\begin{pastquestion}

\QuestionLabel{【往年题2｜04级考题回忆.doc（附：其他年份回忆）｜简答】}\\
膳食纤维的生理作用

\end{pastquestion}

\begin{pastquestion}

\QuestionLabel{【往年题3｜09级考题回忆.docx｜简答】}\\
膳食纤维相关疾病

\end{pastquestion}

\begin{pastquestion}

\QuestionLabel{【往年题4｜14级营养考题.pdf｜简答】}\\
膳食纤维的生理意义

\end{pastquestion}

\begin{pastquestion}

\QuestionLabel{【往年题5｜20级完整试题回忆.docx｜单选】}\\
膳食纤维的主要生理意义是：D 防止便秘，降低血胆固醇

\end{pastquestion}

\PointSeparator

\section{碳水化合物的分类与性质（双糖、寡糖、淀粉、抗性淀粉、果胶）}\label{users__fpb__26_spring__pkusph_study__tmp__pdfs__prevention-medicine-past-exams__05-ux8425ux517bux5b66ux5f80ux5e74ux9898ux8003ux70b9ux6574ux7406.md__ux78b3ux6c34ux5316ux5408ux7269ux7684ux5206ux7c7bux4e0eux6027ux8d28ux53ccux7cd6ux5be1ux7cd6ux6dc0ux7c89ux6297ux6027ux6dc0ux7c89ux679cux80f6}

\InfoLabel{考频} 3 次

\InfoLabel{对应小节} 二、分类（单糖/双糖/寡糖/多糖、抗性淀粉、可溶性纤维）

\InfoLabel{匹配依据}
题干涉及「抗性淀粉」「果胶（Pectin）」名词及碳水化合物性质综合辨析（直链淀粉遇碘变蓝、蔗糖无还原性、寡糖被益生菌利用、葡萄糖吸收方式），与本节直接对应。

\ContentBand{知识点原文摘取}

蔗糖：无还原性（D-G+D-F，α-1,2 糖苷键）。直链淀粉遇碘变蓝，易老化，难消化；支链淀粉遇碘变棕，易糊化。\\
寡糖（3-9 单糖）：肠道菌群代谢，造成胀气，如低聚果糖、低聚半乳糖。\\
抗性淀粉（RS）：健康人小肠内剩余的不被消化吸收的淀粉及其降解产物，在大肠内发酵并吸收，具有类似膳食纤维的生理特性。\\
果胶（pectin）：是被甲酯化到一定程度的半乳糖醛酸多聚体，属可溶性膳食纤维，在食品加工中被用作增稠剂，通常存在于水果和蔬菜中。

\InfoLabel{知识来源} 《营养与食品卫生学·18 柯雅蕾笔记》第 21-23 页

\ContentBand{对应往年题}

\begin{pastquestion}

\QuestionLabel{【往年题1｜08级考题回忆.doc｜名词解释】}\\
抗性淀粉（全英文）

\end{pastquestion}

\begin{pastquestion}

\QuestionLabel{【往年题2｜21级营养期末考题.docx｜名词解释】}\\
Pectin（果胶）

\end{pastquestion}

\begin{pastquestion}

\QuestionLabel{【往年题3｜21级营养期末考题.docx｜选择】}\\
选择正确的一项：A 直链淀粉遇碘变蓝，易老化成难消化抗性淀粉；B 蔗糖是唯一没有还原性的双糖；C
寡糖可以被肠道益生菌利用；D 葡萄糖可以被主动转运、被动转运及通过细胞间隙吸收

\end{pastquestion}

\PointSeparator

\section{血糖生成指数（GI）}\label{users__fpb__26_spring__pkusph_study__tmp__pdfs__prevention-medicine-past-exams__05-ux8425ux517bux5b66ux5f80ux5e74ux9898ux8003ux70b9ux6574ux7406.md__ux8840ux7cd6ux751fux6210ux6307ux6570gi}

\InfoLabel{考频} 1 次

\InfoLabel{对应小节} 五、碳水化合物参考摄入量 / 血糖生成指数 GI

\InfoLabel{匹配依据} 题干为「Glycemic index（血糖生成指数）」名词解释，与本节直接对应。

\ContentBand{知识点原文摘取}

血糖生成指数 GI＝某食物（50g 碳水化合物）在进食后 2h 血糖曲线下面积／相当含量（50g）葡萄糖在进食后 2h
血糖曲线下面积×100。\\
GI≤55 为低 GI 食物，55～70 为中等 GI 食物，＞70 为高 GI 食物。\\
GI 是评价一种食物引起人体餐后血糖应答的非常有意义的生理学指标。

\InfoLabel{知识来源} 《营养与食品卫生学·18 柯雅蕾笔记》第 24-25 页

\ContentBand{对应往年题}

\begin{pastquestion}

\QuestionLabel{【往年题1｜20级完整试题回忆.docx｜名词解释】}\\
Glycemic index

\end{pastquestion}

\PointSeparator

\section{食物热效应（食物特殊动力作用，SDA/TEF）}\label{users__fpb__26_spring__pkusph_study__tmp__pdfs__prevention-medicine-past-exams__05-ux8425ux517bux5b66ux5f80ux5e74ux9898ux8003ux70b9ux6574ux7406.md__ux98dfux7269ux70edux6548ux5e94ux98dfux7269ux7279ux6b8aux52a8ux529bux4f5cux7528sdatef}

\InfoLabel{考频} 2 次

\InfoLabel{对应小节} 能量平衡与健康 /（一）人体的能量消耗 / 2. 食物热效应

\InfoLabel{匹配依据} 题干为「SDA／Special dynamic action（食物特殊动力作用）」名词解释，与本节直接对应。

\ContentBand{知识点原文摘取}

食物热效应（thermic effect of food，TEF），又称为食物特殊动力学作用（specific dynamic
action，SDA），指人体在摄食过程中所引起的额外能量消耗，与食物营养成分、进食量和进食速度有关。\\
脂肪的 TEF 为本身产生能量的 0\%～5\%，碳水化合物为 5\%～10\%，而蛋白质为 20\%～30\%。中国人混合膳食 TEF 约为
BEE 的 10\%。

\InfoLabel{知识来源} 《营养与食品卫生学·18 柯雅蕾笔记》第 27 页

\ContentBand{对应往年题}

\begin{pastquestion}

\QuestionLabel{【往年题1｜04级考题回忆.doc｜名词解释】}\\
SDA

\end{pastquestion}

\begin{pastquestion}

\QuestionLabel{【往年题2｜14级营养考题.pdf｜名词解释】}\\
Special dynamic action

\end{pastquestion}

\PointSeparator

\section{能量平衡与能量需要量（BMI、PAL、能量摄入与消耗、减重）}\label{users__fpb__26_spring__pkusph_study__tmp__pdfs__prevention-medicine-past-exams__05-ux8425ux517bux5b66ux5f80ux5e74ux9898ux8003ux70b9ux6574ux7406.md__ux80fdux91cfux5e73ux8861ux4e0eux80fdux91cfux9700ux8981ux91cfbmipalux80fdux91cfux6444ux5165ux4e0eux6d88ux8017ux51cfux91cd}

\InfoLabel{考频} 2 次

\InfoLabel{对应小节} 能量平衡与健康 /（一）人体的能量消耗；（二）人体的能量需要量

\InfoLabel{匹配依据} 两道论述题给出身高、体重、能量摄入及活动情况，要求计算
BMI、判断能量摄入与消耗是否平衡（PAL＝TEE/BEE）、并提出膳食与锻炼建议，与本节直接对应。

\ContentBand{知识点原文摘取}

能量消耗包括基础代谢（60\%～70\%）、身体活动（15\%～30\%）和食物热效应。\\
身体活动水平 PAL＝TEE/BEE；能量需要量 EER＝BEE×PAL。\\
人体能量的摄入应与机体需要保持平衡；长期能量过剩→异常脂肪堆积→肥胖→心血管疾病、部分癌症、糖尿病等风险升高。\\
膳食能量供能比：蛋白质 10～15\%、脂肪 20～30\%、碳水化合物 50～65\%。

\InfoLabel{知识来源} 《营养与食品卫生学·18 柯雅蕾笔记》第 26-28 页

\ContentBand{对应往年题}

\begin{pastquestion}

\QuestionLabel{【往年题1｜20级完整试题回忆.docx｜论述】}\\
张三，28 岁男性，身高 175cm，体重 80kg\ldots\ldots 每天能量摄入 3200kcal，8 小时睡眠，16
小时轻体力活动。问：（1）BMI
是多少，是否需要管理体重；（2）能量摄入与能量消耗是否平衡；（3）如减重请提出具体措施

\end{pastquestion}

\begin{pastquestion}

\QuestionLabel{【往年题2｜21级营养期末考题.docx｜论述】}\\
男性，给身高、体重、腰围、低强度体力活动、每日三大产能营养素摄入量、睡眠及身体活动消耗能量及基础代谢公式。问：（1）能量需要量（PAL=1.4）；（2）能量消耗和摄入是否平衡；（3）给出膳食和锻炼建议

\end{pastquestion}

\PointSeparator

\section{运动营养与水分补充（运动与健康、补水、补糖）}\label{users__fpb__26_spring__pkusph_study__tmp__pdfs__prevention-medicine-past-exams__05-ux8425ux517bux5b66ux5f80ux5e74ux9898ux8003ux70b9ux6574ux7406.md__ux8fd0ux52a8ux8425ux517bux4e0eux6c34ux5206ux8865ux5145ux8fd0ux52a8ux4e0eux5065ux5eb7ux8865ux6c34ux8865ux7cd6}

\InfoLabel{考频} 5 次

\InfoLabel{对应小节} 能量平衡·身体活动（PA）；膳食指南·身体活动和饮水（未完全匹配到）

\InfoLabel{匹配依据}
题干涉及体力活动中水代谢改变及补水原则、运动降低疾病风险、运动营养补充误区、饮水说法及运动时补糖的意义与策略。笔记述及身体活动（PA）及饮水量推荐，但缺少运动营养（补水、补糖、运动与疾病）的系统内容，故标记为未完全匹配到。

\ContentBand{知识点原文摘取}

未完全匹配到完整原文。\\
可参考原文片段：身体活动是能量平衡和保持身体健康的重要手段；推荐成年人每天进行至少相当于快步走 6000
步以上的身体活动，每周最好进行 150 分钟中等强度的运动。\\
水是一切生命活动必需的物质；低身体活动水平的成年人每天至少饮水 1500～1700ml（7～8
杯），在高温或高身体活动水平条件下应适当增加饮水量。

\InfoLabel{知识来源} 《营养与食品卫生学·18 柯雅蕾笔记》第 16-17 页、第 27-28 页

\ContentBand{对应往年题}

\begin{pastquestion}

\QuestionLabel{【往年题1｜05级考题回忆.docx｜简答】}\\
体力活动中人体水代谢会发生什么改变？补水的原则是什么？

\end{pastquestion}

\begin{pastquestion}

\QuestionLabel{【往年题2｜05级考题回忆.docx｜选择】}\\
目前的科学证据还不能确定运动降低下列疾病发生风险的是（D 卵巢癌）

\end{pastquestion}

\begin{pastquestion}

\QuestionLabel{【往年题3｜05级考题回忆.docx｜选择】}\\
下列说法哪一项不是运动营养补充的误区（A）

\end{pastquestion}

\begin{pastquestion}

\QuestionLabel{【往年题4｜05级考题回忆.docx｜选择】}\\
下列关于饮水的说法正确的是（D 一般情况下，口渴时进行饮水能够满足机体的需要）

\end{pastquestion}

\begin{pastquestion}

\QuestionLabel{【往年题5｜08级考题回忆.doc｜简答】}\\
运动时补糖的意义和策略

\end{pastquestion}

\PointSeparator

\chapter{矿物质与健康}\label{users__fpb__26_spring__pkusph_study__tmp__pdfs__prevention-medicine-past-exams__05-ux8425ux517bux5b66ux5f80ux5e74ux9898ux8003ux70b9ux6574ux7406.md__ux7b2cux516dux7ae0-ux77ffux7269ux8d28ux4e0eux5065ux5eb7}

\section{钙的概述、含量与参考摄入量（混溶钙池）}\label{users__fpb__26_spring__pkusph_study__tmp__pdfs__prevention-medicine-past-exams__05-ux8425ux517bux5b66ux5f80ux5e74ux9898ux8003ux70b9ux6574ux7406.md__ux9499ux7684ux6982ux8ff0ux542bux91cfux4e0eux53c2ux8003ux6444ux5165ux91cfux6df7ux6eb6ux9499ux6c60}

\InfoLabel{考频} 3 次

\InfoLabel{对应小节} 二、钙 /（一）概述；（五）膳食参考摄入量

\InfoLabel{匹配依据} 题干涉及「混溶钙池」名词、「人体含量最多且最易缺乏的矿物质（钙）」及「50
岁以上老年人钙摄入量」，与本节直接对应。

\ContentBand{知识点原文摘取}

钙是人体内含量最多的无机元素，成年人总量为 1000g\textasciitilde1200g，占体重的
1.5\%\textasciitilde2.0\%。99\%存在于骨骼和牙齿中，1\%为混溶钙池（miscible calcium pool）。\\
我国人群比较容易缺乏的矿物质是钙、锌、铁、碘、硒等。\\
膳食参考摄入量（AI）：成年人 800mg/d，老年人（＞50 岁）1000 mg/d；UL 为 2000mg/d。

\InfoLabel{知识来源} 《营养与食品卫生学·18 柯雅蕾笔记》第 30、32-33 页

\ContentBand{对应往年题}

\begin{pastquestion}

\QuestionLabel{【往年题1｜04级考题回忆.doc（附：其他年份回忆）｜名词解释】}\\
混溶钙池

\end{pastquestion}

\begin{pastquestion}

\QuestionLabel{【往年题2｜20级完整试题回忆.docx｜单选】}\\
人体含量最多的且是我国人民最易缺乏的矿物质元素是（D 钙）

\end{pastquestion}

\begin{pastquestion}

\QuestionLabel{【往年题3｜08级考题回忆.doc｜选择】}\\
50 岁以上老年人 Ca 摄入量：1000mg

\end{pastquestion}

\PointSeparator

\section{钙的吸收及影响因素}\label{users__fpb__26_spring__pkusph_study__tmp__pdfs__prevention-medicine-past-exams__05-ux8425ux517bux5b66ux5f80ux5e74ux9898ux8003ux70b9ux6574ux7406.md__ux9499ux7684ux5438ux6536ux53caux5f71ux54cdux56e0ux7d20}

\InfoLabel{考频} 1 次

\InfoLabel{对应小节} 二、钙 /（三）吸收与代谢（有利/不利钙吸收的膳食因素）

\InfoLabel{匹配依据}
题干为「影响钙吸收的膳食因素及机理」（原卷此处另标注「铁」），与本节有利、不利钙吸收的膳食因素直接对应。

\ContentBand{知识点原文摘取}

有利钙吸收的膳食因素：维生素 D（促进肠黏膜对钙的吸收）、乳糖（与钙形成可溶性低分子物质、降低肠内
pH）、低磷膳食、氨基酸（赖、色、精、组）、酪蛋白磷酸肽。\\
不利钙吸收的膳食因素：草酸、植酸（与钙形成不可溶性复合物）、膳食纤维、未被吸收的脂肪酸（形成钙皂）、高钠饮食、乙醇。

\InfoLabel{知识来源} 《营养与食品卫生学·18 柯雅蕾笔记》第 30-31 页

\ContentBand{对应往年题}

\begin{pastquestion}

\QuestionLabel{【往年题1｜14级营养考题.pdf｜简答】}\\
影响钙吸收的膳食因素及机理　铁

\end{pastquestion}

\PointSeparator

\section{钙缺乏与骨骼健康（佝偻病、骨质疏松、骨峰值）}\label{users__fpb__26_spring__pkusph_study__tmp__pdfs__prevention-medicine-past-exams__05-ux8425ux517bux5b66ux5f80ux5e74ux9898ux8003ux70b9ux6574ux7406.md__ux9499ux7f3aux4e4fux4e0eux9aa8ux9abcux5065ux5eb7ux4f5dux507bux75c5ux9aa8ux8d28ux758fux677eux9aa8ux5cf0ux503c}

\InfoLabel{考频} 4 次

\InfoLabel{对应小节} 二、钙 /（四）缺乏症；维生素 D 与健康

\InfoLabel{匹配依据}
题干涉及「Rickets（佝偻病）」名词、佝偻病的诊断/缺乏营养素及功能/改善措施，以及骨骼健康的生理与营养因素，与本节直接对应（佝偻病常伴维生素
D 缺乏）。

\ContentBand{知识点原文摘取}

佝偻病：常伴随维生素 D
缺乏；表现为生长迟缓、新骨结构异常、骨钙化不良、骨骼变形。婴幼儿：囟门闭合延迟、出牙晚、鸡胸、「肋骨串珠」、X／O
型腿等。\\
成年期缺钙：骨质软化、骨质疏松（以骨量减少及骨组织微观结构退变为特征，伴骨脆性增加）。\\
骨峰值（peak bone mass）：骨量大约在 35 岁达到峰值，之后开始降低；补钙时期应在峰值之前。营养因素：钙、维生素
D、过量饮酒、过量饮用咖啡、碳酸饮料等。

\InfoLabel{知识来源} 《营养与食品卫生学·18 柯雅蕾笔记》第 31-32 页、第 40 页

\ContentBand{对应往年题}

\begin{pastquestion}

\QuestionLabel{【往年题1｜20级完整试题回忆.docx｜名词解释】}\\
Rickets（佝偻病）

\end{pastquestion}

\begin{pastquestion}

\QuestionLabel{【往年题2｜08级考题回忆.doc｜问答】}\\
佝偻病：（1）诊断；（2）缺乏的主要营养素及其主要生理功能；（3）改善措施（膳食和其他）

\end{pastquestion}

\begin{pastquestion}

\QuestionLabel{【往年题3｜04级考题回忆.doc｜简答】}\\
你如何理解骨骼健康的生理因素和营养因素？

\end{pastquestion}

\begin{pastquestion}

\QuestionLabel{【往年题4｜09级考题回忆.docx｜简答】}\\
骨骼健康相关生理以及营养

\end{pastquestion}

\PointSeparator

\section{铁的生理功能、存在形式及吸收影响因素}\label{users__fpb__26_spring__pkusph_study__tmp__pdfs__prevention-medicine-past-exams__05-ux8425ux517bux5b66ux5f80ux5e74ux9898ux8003ux70b9ux6574ux7406.md__ux94c1ux7684ux751fux7406ux529fux80fdux5b58ux5728ux5f62ux5f0fux53caux5438ux6536ux5f71ux54cdux56e0ux7d20}

\InfoLabel{考频} 5 次

\InfoLabel{对应小节} 三、铁 /（二）生理功能；（三）吸收代谢（血红素铁与非血红素铁）

\InfoLabel{匹配依据}
题干涉及「非血红素铁」名词、铁的生理功能及在食物中的存在形式与影响吸收因素、如何提高铁利用率，与本节直接对应。

\ContentBand{知识点原文摘取}

食物中的铁有两种存在形式：血红素铁（来自动物性食品血红蛋白和肌红蛋白，吸收率
20-25\%，受膳食因素影响很小）和非血红素铁（主要存在于植物和乳制品，以 Fe(OH)₃ 络合物形式存在，吸收率不足
10\%，受多种因素影响）。\\
促进非血红素铁吸收：维生素 C（3 价→2
价）、动物肉类/肝脏（肉类因子）、有机酸等；降低吸收：植酸盐、草酸盐、钙、茶叶咖啡（鞣酸）、胃酸缺乏等。\\
铁的生理功能：构成血红素蛋白质参与 O₂/CO₂ 转运、促进红细胞形成与成熟、参与多种生理过程、调节免疫功能。

\InfoLabel{知识来源} 《营养与食品卫生学·18 柯雅蕾笔记》第 33-35 页

\ContentBand{对应往年题}

\begin{pastquestion}

\QuestionLabel{【往年题1｜04级考题回忆.doc｜名词解释】}\\
非血红素铁（全英文）

\end{pastquestion}

\begin{pastquestion}

\QuestionLabel{【往年题2｜09级考题回忆.docx｜名词解释】}\\
非血红素铁（全英文）

\end{pastquestion}

\begin{pastquestion}

\QuestionLabel{【往年题3｜04级考题回忆.doc｜简答】}\\
食物铁的存在形式和影响吸收的因素

\end{pastquestion}

\begin{pastquestion}

\QuestionLabel{【往年题4｜04级考题回忆.doc（附：其他年份回忆）｜简答】}\\
人和食物中铁的存在形式，如何提高铁的利用率

\end{pastquestion}

\begin{pastquestion}

\QuestionLabel{【往年题5｜20级完整试题回忆.docx｜简答】}\\
铁的生理功能，在食物中的存在形式及影响吸收的因素

\end{pastquestion}

\PointSeparator

\section{缺铁性贫血及贫血相关营养素}\label{users__fpb__26_spring__pkusph_study__tmp__pdfs__prevention-medicine-past-exams__05-ux8425ux517bux5b66ux5f80ux5e74ux9898ux8003ux70b9ux6574ux7406.md__ux7f3aux94c1ux6027ux8d2bux8840ux53caux8d2bux8840ux76f8ux5173ux8425ux517bux7d20}

\InfoLabel{考频} 3 次

\InfoLabel{对应小节} 三、铁 /（四）缺乏（缺铁分期、表现、治疗）

\InfoLabel{匹配依据}
题干涉及我国居民缺铁性贫血高发的原因及膳食防治措施、缺铁性贫血病例分析，以及「与贫血有关的营养素」多选，与本节直接对应。

\ContentBand{知识点原文摘取}

铁缺乏原因：①铁供给不足及生物利用率较低；②机体对铁的需要量增多（生长突增、月经丢失）；③消化道疾病、失血性疾病等；易感人群为婴幼儿、学龄儿童、孕妇、育龄女性、老年人。\\
缺铁分三期：铁减少期（ID，血清铁蛋白↓）、红细胞生成缺铁期（IDE）、缺铁性贫血期（IDA，血红蛋白下降）。\\
治疗：病因治疗（增加富含铁、维生素 C 的食物）＋铁剂治疗。与贫血有关的营养素包括铁、叶酸、维生素 B12、维生素
C、核黄素等。

\InfoLabel{知识来源} 《营养与食品卫生学·18 柯雅蕾笔记》第 34-36 页

\ContentBand{对应往年题}

\begin{pastquestion}

\QuestionLabel{【往年题1｜08级考题回忆.doc｜简答】}\\
中国居民缺铁性贫血很高的原因，以及用膳食预防和改变铁缺乏的措施

\end{pastquestion}

\begin{pastquestion}

\QuestionLabel{【往年题2｜21级营养期末考题.docx｜简答】}\\
45 岁女性头晕乏力，毛发干枯，指甲改变，有月经过多史，给了血清铁蛋白、红细胞平均体积、血红蛋白
80mg/L、血清运铁蛋白受体等，问是什么营养素缺乏、原因＋营养学纠正措施

\end{pastquestion}

\begin{pastquestion}

\QuestionLabel{【往年题3｜08级考题回忆.doc｜多选】}\\
与贫血有关的营养素：Fe、叶酸、VB12、VC、核黄素

\end{pastquestion}

\PointSeparator

\section{矿物质分类与必需微量元素（硒与克山病、碘与克汀病）}\label{users__fpb__26_spring__pkusph_study__tmp__pdfs__prevention-medicine-past-exams__05-ux8425ux517bux5b66ux5f80ux5e74ux9898ux8003ux70b9ux6574ux7406.md__ux77ffux7269ux8d28ux5206ux7c7bux4e0eux5fc5ux9700ux5faeux91cfux5143ux7d20ux7852ux4e0eux514bux5c71ux75c5ux7898ux4e0eux514bux6c40ux75c5}

\InfoLabel{考频} 3 次

\InfoLabel{对应小节} 一、概述 / 1. 矿物质分类（人体必需微量元素）

\InfoLabel{匹配依据}
题干涉及「必需微量元素」名词、「克山病缺硒」单选及「地方性克汀病」论述。笔记列出人体必需微量元素（含硒、碘），但未述及克山病、地方性克汀病的具体病因与防治，相关两题标记为未完全匹配到。

\ContentBand{知识点原文摘取}

人体必需微量元素：I、Zn、Se、Cu、Mo、Cr、Co、Fe、Mn、F（1995 年把 Mn、F 提到人体必需微量元素）。\\
微量元素含量＜0.01\%体重，需要量＜μg～mg/d。\\
（克山病-缺硒、地方性克汀病-缺碘的具体病因与防治措施未完全匹配到，参见环境健康学相关内容。）

\InfoLabel{知识来源} 《营养与食品卫生学·18 柯雅蕾笔记》第 29-30 页

\ContentBand{对应往年题}

\begin{pastquestion}

\QuestionLabel{【往年题1｜08级考题回忆.doc｜名词解释】}\\
必需微量元素（全英文）

\end{pastquestion}

\begin{pastquestion}

\QuestionLabel{【往年题2｜20级完整试题回忆.docx｜单选】}\\
克山病是由于缺（）------硒（未完全匹配到笔记原文）

\end{pastquestion}

\begin{pastquestion}

\QuestionLabel{【往年题3｜14级营养考题.pdf｜论述】}\\
地方性克汀病的病因和发病机制、发病表现和流行病学特征，我国为防治克汀病采取了什么策略和措施，效果如何（未完全匹配到笔记原文）

\end{pastquestion}

\PointSeparator

\chapter{维生素与健康}\label{users__fpb__26_spring__pkusph_study__tmp__pdfs__prevention-medicine-past-exams__05-ux8425ux517bux5b66ux5f80ux5e74ux9898ux8003ux70b9ux6574ux7406.md__ux7b2cux4e03ux7ae0-ux7ef4ux751fux7d20ux4e0eux5065ux5eb7}

\section{维生素的特性、缺乏与综合辨析}\label{users__fpb__26_spring__pkusph_study__tmp__pdfs__prevention-medicine-past-exams__05-ux8425ux517bux5b66ux5f80ux5e74ux9898ux8003ux70b9ux6574ux7406.md__ux7ef4ux751fux7d20ux7684ux7279ux6027ux7f3aux4e4fux4e0eux7efcux5408ux8fa8ux6790}

\InfoLabel{考频} 3 次

\InfoLabel{对应小节} 一、维生素的定义和分类（定义、分类、缺乏、过量）

\InfoLabel{匹配依据}
题干涉及「维生素有别于矿物质的特性及缺乏原因」论述、「维生素原发性缺乏」名词及「关于维生素说法正确的是」综合辨析，与本节直接对应。

\ContentBand{知识点原文摘取}

维生素是一类维持机体正常生理功能所必需，而体内不能合成或合成数量不能满足机体需要、必须由食物提供的有机化合物，属微量营养素。\\
分脂溶性（A、D、E、K，可贮存、过量可中毒）与水溶性（B 族、C，过量经尿排出）两类。\\
维生素缺乏原因：①摄入量不足；②机体吸收利用率降低；③需要量相对增高；④食物以外供给不足。缺乏在体内有渐进过程，先表现为亚临床缺乏（体力和脑力劳动效率下降、抗病力下降），继而出现临床缺乏症。

\InfoLabel{知识来源} 《营养与食品卫生学·18 柯雅蕾笔记》第 38-39 页

\ContentBand{对应往年题}

\begin{pastquestion}

\QuestionLabel{【往年题1｜04级考题回忆.doc（附：其他年份回忆）｜论述】}\\
维生素有别于矿物质的特性，缺乏原因

\end{pastquestion}

\begin{pastquestion}

\QuestionLabel{【往年题2｜04级考题回忆.doc（附：其他年份回忆）｜名词解释】}\\
维生素原发性缺乏（英文）

\end{pastquestion}

\begin{pastquestion}

\QuestionLabel{【往年题3｜20级完整试题回忆.docx｜单选】}\\
下面关于维生素的说法正确的是（A 维生素 B6 缺乏会导致癞皮病；B 维生素 B1 是细胞分裂增殖必需物质；C 维生素 B2
缺乏会损害神经、心血管和消化系统；D 维生素 C 可以促进类固醇激素代谢；E 维生素 E 容易缺乏）

\end{pastquestion}

\PointSeparator

\section{维生素
B1（硫胺素）：生理功能、缺乏（脚气病）与红细胞转酮醇酶活性系数}\label{users__fpb__26_spring__pkusph_study__tmp__pdfs__prevention-medicine-past-exams__05-ux8425ux517bux5b66ux5f80ux5e74ux9898ux8003ux70b9ux6574ux7406.md__ux7ef4ux751fux7d20-b1ux786bux80faux7d20ux751fux7406ux529fux80fdux7f3aux4e4fux811aux6c14ux75c5ux4e0eux7ea2ux7ec6ux80deux8f6cux916eux9187ux9176ux6d3bux6027ux7cfbux6570}

\InfoLabel{考频} 5 次

\InfoLabel{对应小节} 三、水溶性维生素 /（一）B 族维生素 / 1. 维生素 B1

\InfoLabel{匹配依据} 题干涉及维生素 B1
生理功能及缺乏表现（脚气病）、「红细胞转酮醇酶活性系数」名词及「脚气病是哪种维生素缺乏」选择，与本节直接对应（红细胞转酮醇酶活性系数为评价
B1 营养状况的指标）。

\ContentBand{知识点原文摘取}

维生素 B1 又称硫胺素（抗脚气病维生素），在能量代谢中具有关键作用，神经系统和肌肉都极为依赖硫胺素。\\
生理功能：在体内磷酸化生成焦磷酸硫胺素（TPP）作为辅酶，参与糖代谢（α-酮酸氧化脱羧和转酮基酶反应）；调节神经生理活动（抑制胆碱酯酶）；影响心脏功能和水盐代谢。\\
缺乏症：脚气病（损害神经系统、心血管系统和消化系统），分干性脚气病（多发性神经炎）、湿性脚气病（水肿）、婴儿脚气病等。

\InfoLabel{知识来源} 《营养与食品卫生学·18 柯雅蕾笔记》第 41-42 页

\ContentBand{对应往年题}

\begin{pastquestion}

\QuestionLabel{【往年题1｜04级考题回忆.doc｜简答】}\\
维生素 B1 生理功能、缺乏表现

\end{pastquestion}

\begin{pastquestion}

\QuestionLabel{【往年题2｜09级考题回忆.docx｜简答】}\\
维生素 B1 功能及缺乏症状

\end{pastquestion}

\begin{pastquestion}

\QuestionLabel{【往年题3｜04级考题回忆.doc｜名词解释】}\\
红细胞转酮醇酶活性系数（全英文）

\end{pastquestion}

\begin{pastquestion}

\QuestionLabel{【往年题4｜09级考题回忆.docx｜名词解释】}\\
红细胞转酮醇酶活力系数（全英文）

\end{pastquestion}

\begin{pastquestion}

\QuestionLabel{【往年题5｜14级营养考题.pdf｜选择】}\\
脚气病是哪种维生素缺乏

\end{pastquestion}

\PointSeparator

\section{维生素 A
缺乏的临床表现（干眼病、夜盲症）}\label{users__fpb__26_spring__pkusph_study__tmp__pdfs__prevention-medicine-past-exams__05-ux8425ux517bux5b66ux5f80ux5e74ux9898ux8003ux70b9ux6574ux7406.md__ux7ef4ux751fux7d20-a-ux7f3aux4e4fux7684ux4e34ux5e8aux8868ux73b0ux5e72ux773cux75c5ux591cux76f2ux75c7}

\InfoLabel{考频} 2 次

\InfoLabel{对应小节} 二、脂溶性维生素 /（一）维生素 A 与健康 / 4. 维生素 A 的缺乏

\InfoLabel{匹配依据} 题干为维生素 A 缺乏的临床表现及「xerophthalmia（干眼病）」名词，与本节直接对应。

\ContentBand{知识点原文摘取}

维生素 A
缺乏：（1）使眼暗适应力下降，患夜盲症和干眼病；（2）使黏膜、上皮组织结构破坏：皮肤粗糙、干燥、角质化，呼吸道/消化道/泌尿生殖道黏膜失去滋润，儿童易患支气管肺炎，泌尿系统结石；（3）生长发育受阻，特别是儿童骨骼、牙齿发育受影响；（4）味觉、嗅觉减弱，食欲下降。\\
保持角膜健康：若维生素 A 极为不足，角蛋白会聚集于角膜中，引起眼球干燥、干眼病、失明。

\InfoLabel{知识来源} 《营养与食品卫生学·18 柯雅蕾笔记》第 39-40 页

\ContentBand{对应往年题}

\begin{pastquestion}

\QuestionLabel{【往年题1｜20级完整试题回忆.docx｜简答】}\\
维生素 A 缺乏的临床表现

\end{pastquestion}

\begin{pastquestion}

\QuestionLabel{【往年题2｜14级营养考题.pdf｜名词解释】}\\
xerophthalmia（干眼病）

\end{pastquestion}

\PointSeparator

\section{维生素 D
缺乏症状及其经母乳的供给}\label{users__fpb__26_spring__pkusph_study__tmp__pdfs__prevention-medicine-past-exams__05-ux8425ux517bux5b66ux5f80ux5e74ux9898ux8003ux70b9ux6574ux7406.md__ux7ef4ux751fux7d20-d-ux7f3aux4e4fux75c7ux72b6ux53caux5176ux7ecfux6bcdux4e73ux7684ux4f9bux7ed9}

\InfoLabel{考频} 2 次

\InfoLabel{对应小节} 二、脂溶性维生素 /（二）维生素 D 与健康

\InfoLabel{匹配依据} 题干为「维生素 D
缺乏的症状」简答及「不通过母乳的维生素（VD）」选择，与本节及母乳喂养（维生素 D 较少）直接对应。

\ContentBand{知识点原文摘取}

维生素 D（抗佝偻病维生素）生理功能：在肝、肾转化为活性形式 1,25-(OH)₂VD，促进 Ca、P
在小肠的吸收，调节血钙平衡，促进牙齿和骨骼正常生长钙化。\\
维生素 D 不足：危害骨骼------儿童佝偻病；成人骨软化症。\\
母乳中除维生素 D 较少外，其他营养素可满足 4～6 个月婴儿需要，故母乳喂养需适当补充维生素 D。

\InfoLabel{知识来源} 《营养与食品卫生学·18 柯雅蕾笔记》第 40 页、第 54 页

\ContentBand{对应往年题}

\begin{pastquestion}

\QuestionLabel{【往年题1｜21级营养期末考题.docx｜简答】}\\
维生素 D 缺乏的症状

\end{pastquestion}

\begin{pastquestion}

\QuestionLabel{【往年题2｜08级考题回忆.doc｜选择】}\\
不通过母乳的维生素：VD

\end{pastquestion}

\PointSeparator

\section{维生素
C（坏血病、抗氧化）}\label{users__fpb__26_spring__pkusph_study__tmp__pdfs__prevention-medicine-past-exams__05-ux8425ux517bux5b66ux5f80ux5e74ux9898ux8003ux70b9ux6574ux7406.md__ux7ef4ux751fux7d20-cux574fux8840ux75c5ux6297ux6c27ux5316}

\InfoLabel{考频} 1 次

\InfoLabel{对应小节} 三、水溶性维生素 /（二）维生素 C 与健康

\InfoLabel{匹配依据} 题干为「坏血病是哪种维生素缺乏（VC）」选择，与本节直接对应。

\ContentBand{知识点原文摘取}

维生素
C（抗坏血酸）生理功能：作为抗氧化剂（在小肠中保护铁不被氧化、促进铁吸收）；羟化作用（参与胶原的合成和维护，脯氨酸→羟脯氨酸）；增强免疫反应、抗感染。\\
缺乏维生素 C 的最著名症状是坏血病，由胶原破坏导致。

\InfoLabel{知识来源} 《营养与食品卫生学·18 柯雅蕾笔记》第 45 页

\ContentBand{对应往年题}

\begin{pastquestion}

\QuestionLabel{【往年题1｜14级营养考题.pdf｜选择】}\\
坏血病是哪种维生素缺乏

\end{pastquestion}

\PointSeparator

\section{尼克酸（烟酸）缺乏------癞皮病}\label{users__fpb__26_spring__pkusph_study__tmp__pdfs__prevention-medicine-past-exams__05-ux8425ux517bux5b66ux5f80ux5e74ux9898ux8003ux70b9ux6574ux7406.md__ux5c3cux514bux9178ux70dfux9178ux7f3aux4e4fux765eux76aeux75c5}

\InfoLabel{考频} 1 次

\InfoLabel{对应小节} 三、水溶性维生素 /（一）B 族维生素 / 3. 尼克酸

\InfoLabel{匹配依据} 题干为「Pellagra（癞皮病）」名词解释，与本节直接对应。

\ContentBand{知识点原文摘取}

尼克酸（又称烟酸、维生素 PP、抗癞皮病因子），缺乏将导致癞皮病。\\
癞皮病常伴有其他营养素缺乏症状------4D
综合征：Diarrhea（腹泻）、Dermatitis（皮炎）、Dementia（痴呆）、Death（死亡）。\\
预防关键是尼克酸；色氨酸在体内可转化为尼克酸，玉米蛋白质中色氨酸含量极少。

\InfoLabel{知识来源} 《营养与食品卫生学·18 柯雅蕾笔记》第 43 页

\ContentBand{对应往年题}

\begin{pastquestion}

\QuestionLabel{【往年题1｜21级营养期末考题.docx｜名词解释】}\\
Pellagra（癞皮病）

\end{pastquestion}

\PointSeparator

\section{皮肤健康相关的维生素}\label{users__fpb__26_spring__pkusph_study__tmp__pdfs__prevention-medicine-past-exams__05-ux8425ux517bux5b66ux5f80ux5e74ux9898ux8003ux70b9ux6574ux7406.md__ux76aeux80a4ux5065ux5eb7ux76f8ux5173ux7684ux7ef4ux751fux7d20}

\InfoLabel{考频} 1 次

\InfoLabel{对应小节} 维生素 A、B2、尼克酸、C 相关章节

\InfoLabel{匹配依据}
题干为多选「与皮肤病有关的营养素（VA、VB2、VB3、VC\ldots\ldots）」，与上述维生素对皮肤黏膜的作用相对应。

\ContentBand{知识点原文摘取}

维生素 A 缺乏：皮肤粗糙、干燥、鳞状角质化变化。\\
维生素 B2（核黄素）：保护皮肤、毛囊黏膜及皮脂腺的功能。\\
尼克酸（VB3）缺乏：皮炎（4D 之一）。维生素 C：参与胶原合成与维护。

\InfoLabel{知识来源} 《营养与食品卫生学·18 柯雅蕾笔记》第 40、43 页

\ContentBand{对应往年题}

\begin{pastquestion}

\QuestionLabel{【往年题1｜08级考题回忆.doc｜多选】}\\
与皮肤病有关的营养素：VA、VB2、VB3、VC（另一个选项忘了）

\end{pastquestion}

\PointSeparator

\section{抗氧化营养素及其机理}\label{users__fpb__26_spring__pkusph_study__tmp__pdfs__prevention-medicine-past-exams__05-ux8425ux517bux5b66ux5f80ux5e74ux9898ux8003ux70b9ux6574ux7406.md__ux6297ux6c27ux5316ux8425ux517bux7d20ux53caux5176ux673aux7406}

\InfoLabel{考频} 1 次

\InfoLabel{对应小节} 维生素 E、维生素 C 与健康；植物化学物（抗氧化作用）

\InfoLabel{匹配依据} 题干为「抗氧化营养素有哪些？机理？」，与维生素 E、维生素 C
及植物化学物的抗氧化作用相对应。

\ContentBand{知识点原文摘取}

维生素 E 可用作抗氧化剂，通过自身氧化保护多不饱和脂肪酸和细胞膜免受破坏，在肺部具有特别重要的抗氧化作用。\\
维生素 C 作为抗氧化剂，保护铁不被氧化、保护血液敏感成分并有助于保护维生素 E。\\
植物化学物（如多酚、类胡萝卜素、植酸等）具有抗氧化作用。

\InfoLabel{知识来源} 《营养与食品卫生学·18 柯雅蕾笔记》第 41、45-46 页

\ContentBand{对应往年题}

\begin{pastquestion}

\QuestionLabel{【往年题1｜08级考题回忆.doc｜简答】}\\
抗氧化营养素有哪些？机理？

\end{pastquestion}

\PointSeparator

\chapter{食物、膳食模式与健康}\label{users__fpb__26_spring__pkusph_study__tmp__pdfs__prevention-medicine-past-exams__05-ux8425ux517bux5b66ux5f80ux5e74ux9898ux8003ux70b9ux6574ux7406.md__ux7b2cux516bux7ae0-ux98dfux7269ux81b3ux98dfux6a21ux5f0fux4e0eux5065ux5eb7}

\section{植物化学物（定义、分类、生物活性，多酚）}\label{users__fpb__26_spring__pkusph_study__tmp__pdfs__prevention-medicine-past-exams__05-ux8425ux517bux5b66ux5f80ux5e74ux9898ux8003ux70b9ux6574ux7406.md__ux690dux7269ux5316ux5b66ux7269ux5b9aux4e49ux5206ux7c7bux751fux7269ux6d3bux6027ux591aux915a}

\InfoLabel{考频} 3 次

\InfoLabel{对应小节} 一、植物化学物（定义、分类、生物活性）

\InfoLabel{匹配依据} 题干涉及多酚类物质的作用、「下面不属于植物化学物的是（辅酶 Q10）」及「bioactive food
component（生物活性食物成分）」名词，与本节直接对应（生物活性食物成分为部分匹配）。

\ContentBand{知识点原文摘取}

植物化学物（phytochemicals）：植物性食物中含有的除必需营养成分外的一些低分子量生物活性物质。\\
分类包括类胡萝卜素、多酚、单萜类、芥子油苷、皂苷、植物固醇、硫化物、植物雌激素、植酸等。\\
多酚：生物学分布最广，存在于各类植物性食物，尤其是深色水果、蔬菜和谷物；具抗肿瘤、抗微生物、抗氧化、抗血栓、免疫调节、抑制炎症、影响血压、降胆固醇、调节血糖等多种生物活性。

\InfoLabel{知识来源} 《营养与食品卫生学·18 柯雅蕾笔记》第 45-46 页

\ContentBand{对应往年题}

\begin{pastquestion}

\QuestionLabel{【往年题1｜14级营养考题.pdf｜选择】}\\
多酚类物质的作用有哪些

\end{pastquestion}

\begin{pastquestion}

\QuestionLabel{【往年题2｜20级完整试题回忆.docx｜单选】}\\
下面不属于植物化学物的是（A 辅酶 Q10）

\end{pastquestion}

\begin{pastquestion}

\QuestionLabel{【往年题3｜14级营养考题.pdf｜名词解释】}\\
bioactive food component（生物活性食物成分）

\end{pastquestion}

\PointSeparator

\section{食物营养价值评价------营养质量指数（INQ）}\label{users__fpb__26_spring__pkusph_study__tmp__pdfs__prevention-medicine-past-exams__05-ux8425ux517bux5b66ux5f80ux5e74ux9898ux8003ux70b9ux6574ux7406.md__ux98dfux7269ux8425ux517bux4ef7ux503cux8bc4ux4ef7ux8425ux517bux8d28ux91cfux6307ux6570inq}

\InfoLabel{考频} 2 次

\InfoLabel{对应小节} 二、食物的营养价值 / 3. 评价指标 /（1）营养质量指数 INQ

\InfoLabel{匹配依据} 题干为「INQ」名词及「计算 INQ 并给出评价」，与本节直接对应。

\ContentBand{知识点原文摘取}

营养质量指数（INQ）：某食物中营养素能满足人体营养需要的程度（营养素密度）与该食物能满足人体能量需要的程度（能量密度）的比值。INQ＝（某种营养素含量/该营养素参考摄入量）/（所产生的能量/能量参考摄入量）。\\
INQ＝1（0.9-1.1）营养素与能量供给能力平衡；INQ＞1 相对适合超重和肥胖者减重；INQ＜1
长期食用可能导致该营养素不足或能量过剩。

\InfoLabel{知识来源} 《营养与食品卫生学·18 柯雅蕾笔记》第 47 页

\ContentBand{对应往年题}

\begin{pastquestion}

\QuestionLabel{【往年题1｜21级营养期末考题.docx｜名词解释】}\\
INQ

\end{pastquestion}

\begin{pastquestion}

\QuestionLabel{【往年题2｜14级营养考题.pdf｜简答】}\\
计算 INQ，给出评价

\end{pastquestion}

\PointSeparator

\section{膳食模式------地中海膳食}\label{users__fpb__26_spring__pkusph_study__tmp__pdfs__prevention-medicine-past-exams__05-ux8425ux517bux5b66ux5f80ux5e74ux9898ux8003ux70b9ux6574ux7406.md__ux81b3ux98dfux6a21ux5f0fux5730ux4e2dux6d77ux81b3ux98df}

\InfoLabel{考频} 1 次

\InfoLabel{对应小节} 三、膳食模式与健康 / 3. 地中海膳食模式

\InfoLabel{匹配依据} 题干为「Mediterranean diet（地中海膳食模式）」名词解释，与本节直接对应。

\ContentBand{知识点原文摘取}

地中海膳食模式（Mediterranean
diet）膳食特点：较高摄入蔬菜、水果、全谷类、豆类和坚果；适量摄入奶制品（多为奶酪和酸奶）、红酒和鱼类等海产品；肉类及其制品摄入量较低；橄榄油为主要食用油及脂肪来源。\\
营养特点：高膳食纤维、高维生素、高单不饱和脂肪酸和低饱和脂肪。

\InfoLabel{知识来源} 《营养与食品卫生学·18 柯雅蕾笔记》第 49-50 页

\ContentBand{对应往年题}

\begin{pastquestion}

\QuestionLabel{【往年题1｜20级完整试题回忆.docx｜名词解释】}\\
Mediterranean diet

\end{pastquestion}

\PointSeparator

\chapter{生命周期营养}\label{users__fpb__26_spring__pkusph_study__tmp__pdfs__prevention-medicine-past-exams__05-ux8425ux517bux5b66ux5f80ux5e74ux9898ux8003ux70b9ux6574ux7406.md__ux7b2cux4e5dux7ae0-ux751fux547dux5468ux671fux8425ux517b}

\section{母乳喂养的优点}\label{users__fpb__26_spring__pkusph_study__tmp__pdfs__prevention-medicine-past-exams__05-ux8425ux517bux5b66ux5f80ux5e74ux9898ux8003ux70b9ux6574ux7406.md__ux6bcdux4e73ux5582ux517bux7684ux4f18ux70b9}

\InfoLabel{考频} 4 次

\InfoLabel{对应小节} 四、婴儿营养 / 3. 母乳喂养的优点

\InfoLabel{匹配依据} 题干为母乳喂养的优点（含「优点不包括」选择），与本节直接对应。

\ContentBand{知识点原文摘取}

母乳喂养优点：①母乳中营养素全面、质量好、比例合适，适应婴儿消化功能和发育需要（优质乳清蛋白、含脂酶易消化、乳糖促进钙吸收、钙磷比例合适）；②母乳富含免疫物质（初乳含淋巴细胞、抗体、IgA，非特异性免疫物质含乳铁蛋白、溶菌酶、双歧杆菌因子等）；③不容易发生过敏；④增进母婴交流、促进婴儿智力发育；⑤促进乳母体形恢复；⑥卫生、经济、方便、温度适宜、新鲜不变质。

\InfoLabel{知识来源} 《营养与食品卫生学·18 柯雅蕾笔记》第 54 页

\ContentBand{对应往年题}

\begin{pastquestion}

\QuestionLabel{【往年题1｜04级考题回忆.doc｜简答】}\\
母乳喂养优点

\end{pastquestion}

\begin{pastquestion}

\QuestionLabel{【往年题2｜04级考题回忆.doc（附：其他年份回忆）｜简答】}\\
母乳喂养的优点

\end{pastquestion}

\begin{pastquestion}

\QuestionLabel{【往年题3｜09级考题回忆.docx｜简答】}\\
母乳喂养优点

\end{pastquestion}

\begin{pastquestion}

\QuestionLabel{【往年题4｜14级营养考题.pdf｜选择】}\\
母乳喂养的优点不包括

\end{pastquestion}

\PointSeparator

\section{婴幼儿辅食添加（时间与原则）}\label{users__fpb__26_spring__pkusph_study__tmp__pdfs__prevention-medicine-past-exams__05-ux8425ux517bux5b66ux5f80ux5e74ux9898ux8003ux70b9ux6574ux7406.md__ux5a74ux5e7cux513fux8f85ux98dfux6dfbux52a0ux65f6ux95f4ux4e0eux539fux5219}

\InfoLabel{考频} 2 次

\InfoLabel{对应小节} 四、婴儿营养 / 6. 7\textasciitilde24 月龄婴儿喂养指南；7. 喂养注意事项

\InfoLabel{匹配依据} 题干为「婴儿什么时候添加辅食」及「7-24 个月婴幼儿添加辅食的原则」，与本节直接对应。

\ContentBand{知识点原文摘取}

7\textasciitilde24 月龄婴儿喂养指南：继续母乳喂养，满 6
月龄起必须添加辅食，从富含铁的泥糊状食物开始；及时引入多样化食物，重视动物性食物的添加；尽量少加糖盐，油脂适当，保持食物原味；提倡回应式喂养，鼓励但不强迫进食；注重饮食卫生和进食安全；定期监测体格指标。\\
注意事项：添加辅食不等于断奶；应在婴儿健康、消化功能正常时添加；一种辅食适应数日后再增加新品种；辅食中适量添加植物油。

\InfoLabel{知识来源} 《营养与食品卫生学·18 柯雅蕾笔记》第 55 页

\ContentBand{对应往年题}

\begin{pastquestion}

\QuestionLabel{【往年题1｜08级考题回忆.doc｜选择】}\\
婴儿什么时候添加辅食？

\end{pastquestion}

\begin{pastquestion}

\QuestionLabel{【往年题2｜21级营养期末考题.docx｜简答】}\\
7-24 个月婴幼儿添加辅食的原则（坚持母乳喂养）

\end{pastquestion}

\PointSeparator

\section{孕期营养不良对母体和胎儿的影响（胎盘生化阀）}\label{users__fpb__26_spring__pkusph_study__tmp__pdfs__prevention-medicine-past-exams__05-ux8425ux517bux5b66ux5f80ux5e74ux9898ux8003ux70b9ux6574ux7406.md__ux5b55ux671fux8425ux517bux4e0dux826fux5bf9ux6bcdux4f53ux548cux80ceux513fux7684ux5f71ux54cdux80ceux76d8ux751fux5316ux9600}

\InfoLabel{考频} 3 次

\InfoLabel{对应小节} 二、孕妇营养 / 3. 对母亲的影响；4. 对胎儿的影响

\InfoLabel{匹配依据}
题干涉及孕期营养不良对母体和胎儿的影响及「胎盘生化阀」名词。笔记述及孕期营养不良对母婴的影响，但未收录「胎盘生化阀」定义，相关两题标记为未完全匹配到。

\ContentBand{知识点原文摘取}

孕期营养不良对母亲的影响：营养性贫血（缺铁性贫血、巨幼红细胞性贫血）、骨质软化症（VD、钙）、营养不良性水肿（蛋白质、VB1）、维生素缺乏症。\\
孕期营养不良对胎儿的影响：新生儿死亡、小于胎龄儿、早产儿、低体重儿、先天畸形、巨大儿、脑发育受损。\\
（「胎盘生化阀」一词未完全匹配到笔记原文。）

\InfoLabel{知识来源} 《营养与食品卫生学·18 柯雅蕾笔记》第 51-52 页

\ContentBand{对应往年题}

\begin{pastquestion}

\QuestionLabel{【往年题1｜08级考题回忆.doc｜简答】}\\
孕期营养不良对母体和胎儿影响

\end{pastquestion}

\begin{pastquestion}

\QuestionLabel{【往年题2｜04级考题回忆.doc（附：其他年份回忆）｜名词解释】}\\
胎盘生化阀（未完全匹配到笔记原文）

\end{pastquestion}

\begin{pastquestion}

\QuestionLabel{【往年题3｜08级考题回忆.doc｜名词解释】}\\
胎盘生化阀（未完全匹配到笔记原文）

\end{pastquestion}

\PointSeparator

\section{孕期膳食原则}\label{users__fpb__26_spring__pkusph_study__tmp__pdfs__prevention-medicine-past-exams__05-ux8425ux517bux5b66ux5f80ux5e74ux9898ux8003ux70b9ux6574ux7406.md__ux5b55ux671fux81b3ux98dfux539fux5219}

\InfoLabel{考频} 1 次

\InfoLabel{对应小节} 二、孕妇营养 / 5-9. 各孕期营养与膳食指南

\InfoLabel{匹配依据} 题干为「孕期膳食原则」，与本节直接对应。

\ContentBand{知识点原文摘取}

孕早期：碳水化合物＞130g/d；叶酸
400μg/d；清淡、易消化、少食多餐。孕中期：充足能量，注意铁的补充，保证鱼禽蛋瘦肉和奶；易缺乏铁、钙、锌、碘、DHA。孕晚期：宜选体积小、营养价值高的食物，少食多餐。\\
备孕和孕期妇女膳食指南：调整孕前体重至正常范围；常吃含铁丰富的食物，选用碘盐，合理补充叶酸和维生素
D；孕中晚期适量增加奶、鱼、禽、蛋、瘦肉；经常户外活动，禁烟酒。

\InfoLabel{知识来源} 《营养与食品卫生学·18 柯雅蕾笔记》第 51-52 页

\ContentBand{对应往年题}

\begin{pastquestion}

\QuestionLabel{【往年题1｜20级完整试题回忆.docx｜简答】}\\
孕期膳食原则

\end{pastquestion}

\PointSeparator

\section{生命早期 1000
天}\label{users__fpb__26_spring__pkusph_study__tmp__pdfs__prevention-medicine-past-exams__05-ux8425ux517bux5b66ux5f80ux5e74ux9898ux8003ux70b9ux6574ux7406.md__ux751fux547dux65e9ux671f-1000-ux5929}

\InfoLabel{考频} 1 次

\InfoLabel{对应小节} 一、生命早期营养 / 3. 生命早期 1000 天

\InfoLabel{匹配依据} 题干为「生命早期 1000 天是指」，与本节直接对应。

\ContentBand{知识点原文摘取}

《中国 0～6 岁儿童营养发展报告（2012）》提出生命早期 1000 天（从怀孕到 2
岁），是决定孩子一生营养与健康状况的最关键时期。\\
DOHaD：在发育过程的早期（包括胎儿和婴幼儿时期）经历不利因素，将增加成年后罹患肥胖、糖尿病、心血管疾病等慢性代谢性疾病的几率。

\InfoLabel{知识来源} 《营养与食品卫生学·18 柯雅蕾笔记》第 50 页

\ContentBand{对应往年题}

\begin{pastquestion}

\QuestionLabel{【往年题1｜20级完整试题回忆.docx｜单选】}\\
生命早期 1000 天是指

\end{pastquestion}

\PointSeparator

\chapter{老年营养}\label{users__fpb__26_spring__pkusph_study__tmp__pdfs__prevention-medicine-past-exams__05-ux8425ux517bux5b66ux5f80ux5e74ux9898ux8003ux70b9ux6574ux7406.md__ux7b2cux5341ux7ae0-ux8001ux5e74ux8425ux517b}

\section{老年综合征}\label{users__fpb__26_spring__pkusph_study__tmp__pdfs__prevention-medicine-past-exams__05-ux8425ux517bux5b66ux5f80ux5e74ux9898ux8003ux70b9ux6574ux7406.md__ux8001ux5e74ux7efcux5408ux5f81}

\InfoLabel{考频} 1 次

\InfoLabel{对应小节} 一、老龄化现状 / 4. 老年综合征（GS）

\InfoLabel{匹配依据} 题干为「Geriatric syndrome（老年综合征）」名词解释，与本节直接对应。

\ContentBand{知识点原文摘取}

老年综合征（geriatric syndrome，GS）是指老年人由多种原因造成的同一种临床表现或问题的症候群，在 65
岁以上老年人群中发病率超过 30\%。\\
常见包括：步态异常、容易跌倒、视力/听力/睡眠障碍、痴呆、记忆力减退、抑郁症、营养不良、衰弱、多重用药等。

\InfoLabel{知识来源} 《营养与食品卫生学·18 柯雅蕾笔记》第 56 页

\ContentBand{对应往年题}

\begin{pastquestion}

\QuestionLabel{【往年题1｜20级完整试题回忆.docx｜名词解释】}\\
Geriatric syndrome

\end{pastquestion}

\PointSeparator

\section{肌肉衰减症（少肌症，Sarcopenia）}\label{users__fpb__26_spring__pkusph_study__tmp__pdfs__prevention-medicine-past-exams__05-ux8425ux517bux5b66ux5f80ux5e74ux9898ux8003ux70b9ux6574ux7406.md__ux808cux8089ux8870ux51cfux75c7ux5c11ux808cux75c7sarcopenia}

\InfoLabel{考频} 1 次

\InfoLabel{对应小节} 三、老年人营养管理 / 肌肉衰减症（少肌症）

\InfoLabel{匹配依据} 题干为「Sarcopenia（肌肉衰减症/少肌症）」名词解释，与本节直接对应。

\ContentBand{知识点原文摘取}

肌肉衰减症/少肌症（Sarcopenia）：与增龄相关的进行性骨骼肌量减少、伴有肌肉力量和（或）肌肉功能减退的综合征，会引起虚弱，增加老年人跌倒、失能、生活质量下降、死亡风险等不良结局。\\
防治策略：合理营养和抗阻运动是预防少肌症的重要手段（老年人蛋白质推荐摄入量为每日每公斤体重 1.0-1.5g）。

\InfoLabel{知识来源} 《营养与食品卫生学·18 柯雅蕾笔记》第 59 页

\ContentBand{对应往年题}

\begin{pastquestion}

\QuestionLabel{【往年题1｜21级营养期末考题.docx｜名词解释】}\\
Sarcopenia

\end{pastquestion}

\PointSeparator

\section{老年人膳食指南与老年营养要点}\label{users__fpb__26_spring__pkusph_study__tmp__pdfs__prevention-medicine-past-exams__05-ux8425ux517bux5b66ux5f80ux5e74ux9898ux8003ux70b9ux6574ux7406.md__ux8001ux5e74ux4ebaux81b3ux98dfux6307ux5357ux4e0eux8001ux5e74ux8425ux517bux8981ux70b9}

\InfoLabel{考频} 1 次

\InfoLabel{对应小节} 三、老年人营养管理 /（一）老年人膳食实践（一般/高龄老年人膳食指南）

\InfoLabel{匹配依据}
题干为「关于老年营养说法错误的是」，与本节老年人膳食指南、适宜体重、少肌症防治等要点直接对应。

\ContentBand{知识点原文摘取}

一般老年人膳食指南：食物品种丰富，动物性食物充足，常吃大豆制品；鼓励共同进餐，保持良好食欲；积极户外活动，延缓肌肉衰减，保持适宜体重；定期健康体检，测评营养状况，预防营养缺乏。\\
老年人应保障适宜体重，BMI 应维持在 21-26.9
之间。高龄老年人应选择质地细软、能量和营养素密度高的食物。合理营养和抗阻运动是预防少肌症的重要手段。

\InfoLabel{知识来源} 《营养与食品卫生学·18 柯雅蕾笔记》第 57-60 页

\ContentBand{对应往年题}

\begin{pastquestion}

\QuestionLabel{【往年题1｜20级完整试题回忆.docx｜单选】}\\
关于老年营养说法错误的是（A 食物品种丰富，动物性食物充足，常吃大豆制品；B
合理营养和抗阻运动是预防少肌症的重要手段；C 千金难买老来瘦，BMI 应维持在 18.5-23.9 之间；D 定期健康体检；E
高龄老年人应选择质地细软的食物）

\end{pastquestion}

\PointSeparator

\chapter{临床营养}\label{users__fpb__26_spring__pkusph_study__tmp__pdfs__prevention-medicine-past-exams__05-ux8425ux517bux5b66ux5f80ux5e74ux9898ux8003ux70b9ux6574ux7406.md__ux7b2cux5341ux4e00ux7ae0-ux4e34ux5e8aux8425ux517b}

\section{肠内营养（定义与禁忌症）}\label{users__fpb__26_spring__pkusph_study__tmp__pdfs__prevention-medicine-past-exams__05-ux8425ux517bux5b66ux5f80ux5e74ux9898ux8003ux70b9ux6574ux7406.md__ux80a0ux5185ux8425ux517bux5b9aux4e49ux4e0eux7981ux5fccux75c7}

\InfoLabel{考频} 3 次

\InfoLabel{对应小节} 二、临床营养基本技能 / 7. 营养治疗的途径；9. 营养支持途径选择原则

\InfoLabel{匹配依据} 题干为「Enteral nutrition（肠内营养）」名词及「肠内营养的禁忌症」，与本节直接对应。

\ContentBand{知识点原文摘取}

肠内营养（enteral
nutrition，EN）：经胃肠道用口服或管饲的途径，为营养摄入不足、不愿或吸收障碍的病人提供代谢需要的营养素。\\
营养支持途径选择原则：当胃肠道功能允许及安全前提下，首选肠内营养；肠外营养适用于急性胰腺炎、消化道活动性出血、肠梗阻等。\\
（肠内营养禁忌症对应肠外营养适应症，即胃肠道功能障碍情形。）

\InfoLabel{知识来源} 《营养与食品卫生学·18 柯雅蕾笔记》第 62-63 页

\ContentBand{对应往年题}

\begin{pastquestion}

\QuestionLabel{【往年题1｜14级营养考题.pdf｜名词解释】}\\
Enteral nutrition

\end{pastquestion}

\begin{pastquestion}

\QuestionLabel{【往年题2｜21级营养期末考题.docx｜名词解释】}\\
Enteral nutrition

\end{pastquestion}

\begin{pastquestion}

\QuestionLabel{【往年题3｜14级营养考题.pdf｜简答】}\\
肠内营养的禁忌症

\end{pastquestion}

\PointSeparator

\section{急性胰腺炎的临床营养支持方式}\label{users__fpb__26_spring__pkusph_study__tmp__pdfs__prevention-medicine-past-exams__05-ux8425ux517bux5b66ux5f80ux5e74ux9898ux8003ux70b9ux6574ux7406.md__ux6025ux6027ux80f0ux817aux708eux7684ux4e34ux5e8aux8425ux517bux652fux6301ux65b9ux5f0f}

\InfoLabel{考频} 2 次

\InfoLabel{对应小节} 二、临床营养基本技能 / 9. 营养支持途径选择原则（肠外营养）

\InfoLabel{匹配依据} 题干为「急性胰腺炎患者应选择的临床营养支持方式」，与本节肠外营养适应症直接对应。

\ContentBand{知识点原文摘取}

营养支持途径选择原则：当胃肠道功能允许及安全前提下，首选肠内营养；肠外营养（parenteral
nutrition，PN）适用于急性胰腺炎、消化道活动性出血、肠梗阻、TPN/PPN 等情形。

\InfoLabel{知识来源} 《营养与食品卫生学·18 柯雅蕾笔记》第 63 页

\ContentBand{对应往年题}

\begin{pastquestion}

\QuestionLabel{【往年题1｜14级营养考题.pdf｜选择】}\\
急性胰腺炎的临床营养方式

\end{pastquestion}

\begin{pastquestion}

\QuestionLabel{【往年题2｜20级完整试题回忆.docx｜单选】}\\
急性胰腺炎患者应选择的临床营养支持方式

\end{pastquestion}

\PointSeparator

\section{胆囊炎的临床营养（减少脂肪）}\label{users__fpb__26_spring__pkusph_study__tmp__pdfs__prevention-medicine-past-exams__05-ux8425ux517bux5b66ux5f80ux5e74ux9898ux8003ux70b9ux6574ux7406.md__ux80c6ux56caux708eux7684ux4e34ux5e8aux8425ux517bux51cfux5c11ux8102ux80aa}

\InfoLabel{考频} 1 次

\InfoLabel{对应小节} 二、临床营养基本技能（治疗膳食制定，未完全匹配到）

\InfoLabel{匹配依据}
题干为「胆囊炎的临床营养需要减少哪种营养素」。笔记给出临床营养能量代谢和营养需求的总体框架，但未单列胆囊炎的治疗膳食，故标记为未完全匹配到（按治疗原则应减少脂肪）。

\ContentBand{知识点原文摘取}

未完全匹配到完整原文。\\
可参考原文片段：治疗膳食制定依据循证指南、食物成分、患者耐受、餐食质量；脂类热量
20-30\%，1-1.5g/kg/d。（胆囊炎临床营养通常需减少脂肪摄入。）

\InfoLabel{知识来源} 《营养与食品卫生学·18 柯雅蕾笔记》第 62-63 页

\ContentBand{对应往年题}

\begin{pastquestion}

\QuestionLabel{【往年题1｜14级营养考题.pdf｜选择】}\\
胆囊炎的临床营养需要减少哪种营养素

\end{pastquestion}

\PointSeparator

\section{营养风险筛查}\label{users__fpb__26_spring__pkusph_study__tmp__pdfs__prevention-medicine-past-exams__05-ux8425ux517bux5b66ux5f80ux5e74ux9898ux8003ux70b9ux6574ux7406.md__ux8425ux517bux98ceux9669ux7b5bux67e5}

\InfoLabel{考频} 1 次

\InfoLabel{对应小节} 临床营养（营养风险筛查）；老年营养（营养诊疗程序 NRS2002）

\InfoLabel{匹配依据} 题干为「Nutrition risk screening（营养风险筛查）」名词解释，与本节直接对应。

\ContentBand{知识点原文摘取}

营养诊疗程序关键步骤：筛查、评定、干预、监测。筛查用于找到需要营养干预的患者，常用工具有营养风险筛查
2002（NRS2002）、营养不良通用筛查工具（MUST）。\\
临床营养对象由营养不良患者扩展到有营养风险者。

\InfoLabel{知识来源} 《营养与食品卫生学·18 柯雅蕾笔记》第 60、62 页

\ContentBand{对应往年题}

\begin{pastquestion}

\QuestionLabel{【往年题1｜20级完整试题回忆.docx｜名词解释】}\\
Nutrition risk screening

\end{pastquestion}

\PointSeparator

\chapter{食品安全风险分析与控制}\label{users__fpb__26_spring__pkusph_study__tmp__pdfs__prevention-medicine-past-exams__05-ux8425ux517bux5b66ux5f80ux5e74ux9898ux8003ux70b9ux6574ux7406.md__ux7b2cux5341ux4e8cux7ae0-ux98dfux54c1ux5b89ux5168ux98ceux9669ux5206ux6790ux4e0eux63a7ux5236}

\section{食品安全风险分析框架（风险评估、风险管理、风险交流）及步骤}\label{users__fpb__26_spring__pkusph_study__tmp__pdfs__prevention-medicine-past-exams__05-ux8425ux517bux5b66ux5f80ux5e74ux9898ux8003ux70b9ux6574ux7406.md__ux98dfux54c1ux5b89ux5168ux98ceux9669ux5206ux6790ux6846ux67b6ux98ceux9669ux8bc4ux4f30ux98ceux9669ux7ba1ux7406ux98ceux9669ux4ea4ux6d41ux53caux6b65ux9aa4}

\InfoLabel{考频} 4 次

\InfoLabel{对应小节} 一、基本概念和认识；二、基本步骤与方法

\InfoLabel{匹配依据} 题干涉及「风险分析」名词、「Risk management（风险管理）」「Risk
characterization（风险特征描述）」名词及「食品安全风险分析的步骤及其作用」，与本章直接对应。

\ContentBand{知识点原文摘取}

食品安全风险分析（risk analysis）由风险评估、风险管理和风险交流三部分共同构成。\\
食品安全风险评估包括危害识别、危害特征描述、暴露评估和风险特征描述 4 个步骤。\\
风险管理：依据风险评估的结果，同时考虑社会、经济等因素，对各种管理措施方案进行权衡、选择并实施，产生结果包括制定食品安全标准、准则等。\\
风险交流：风险评估者、风险管理者及社会相关团体公众之间各方面的信息交流。

\InfoLabel{知识来源} 《营养与食品卫生学·18 柯雅蕾笔记》第 4、64-65 页

\ContentBand{对应往年题}

\begin{pastquestion}

\QuestionLabel{【往年题1｜08级考题回忆.doc｜名词解释】}\\
风险分析（全英文）

\end{pastquestion}

\begin{pastquestion}

\QuestionLabel{【往年题2｜20级完整试题回忆.docx｜名词解释】}\\
Risk management（风险管理）

\end{pastquestion}

\begin{pastquestion}

\QuestionLabel{【往年题3｜21级营养期末考题.docx｜名词解释】}\\
Risk characterization（风险特征描述）

\end{pastquestion}

\begin{pastquestion}

\QuestionLabel{【往年题4｜20级完整试题回忆.docx｜简答】}\\
试述食品安全风险分析的步骤及其作用

\end{pastquestion}

\PointSeparator

\section{暴露评估}\label{users__fpb__26_spring__pkusph_study__tmp__pdfs__prevention-medicine-past-exams__05-ux8425ux517bux5b66ux5f80ux5e74ux9898ux8003ux70b9ux6574ux7406.md__ux66b4ux9732ux8bc4ux4f30}

\InfoLabel{考频} 1 次

\InfoLabel{对应小节} 二、基本步骤与方法 /（三）暴露评估

\InfoLabel{匹配依据} 题干为「暴露评估的意义及常用方法」，与本节直接对应。

\ContentBand{知识点原文摘取}

暴露评估：对通过食品或其他相关来源摄入的危害因素进行定性和/或定量评估，是风险评估的第三步。膳食暴露量＝Σ（食物消费量×化学物浓度）/体重。\\
分急性暴露评估（24 小时以内）和慢性暴露评估（每天暴露并持续终生）。\\
常用方法：筛选法（产销量法、预算法）、点评估、模型膳食、基于监测抽检数据、总膳食研究、概率评估等。

\InfoLabel{知识来源} 《营养与食品卫生学·18 柯雅蕾笔记》第 66-68 页

\ContentBand{对应往年题}

\begin{pastquestion}

\QuestionLabel{【往年题1｜21级营养期末考题.docx｜简答】}\\
暴露评估的意义及常用方法

\end{pastquestion}

\PointSeparator

\section{食品安全风险监测}\label{users__fpb__26_spring__pkusph_study__tmp__pdfs__prevention-medicine-past-exams__05-ux8425ux517bux5b66ux5f80ux5e74ux9898ux8003ux70b9ux6574ux7406.md__ux98dfux54c1ux5b89ux5168ux98ceux9669ux76d1ux6d4b}

\InfoLabel{考频} 1 次

\InfoLabel{对应小节} 一、基本概念和认识；食品安全法律法规体系（风险监测）

\InfoLabel{匹配依据}
题干为「我国食品安全风险监测制度的监测对象」（食源性疾病、食品污染、食品中的有害因素），与风险监测内容相对应（笔记未单列监测对象清单，为部分匹配）。

\ContentBand{知识点原文摘取}

食品安全技术支撑体系：实验室检验技术→风险监测→风险评估→预警、标准、交流→监管措施。\\
食品安全风险管理（食品安全风险监测、风险评估，制定并公布食品安全国家标准）主要由国家卫生健康委承担。\\
（监测对象为食源性疾病、食品污染、食品中的有害因素，属部分匹配。）

\InfoLabel{知识来源} 《营养与食品卫生学·18 柯雅蕾笔记》第 64、83 页

\ContentBand{对应往年题}

\begin{pastquestion}

\QuestionLabel{【往年题1｜20级完整试题回忆.docx｜单选】}\\
我国食品安全风险监测制度的监测对象是（D 食源性疾病、食品污染、食品中的有害因素）

\end{pastquestion}

\PointSeparator

\chapter{食品污染与食源性疾病}\label{users__fpb__26_spring__pkusph_study__tmp__pdfs__prevention-medicine-past-exams__05-ux8425ux517bux5b66ux5f80ux5e74ux9898ux8003ux70b9ux6574ux7406.md__ux7b2cux5341ux4e09ux7ae0-ux98dfux54c1ux6c61ux67d3ux4e0eux98dfux6e90ux6027ux75beux75c5}

\section{食品污染的概念、分类与生物富集作用}\label{users__fpb__26_spring__pkusph_study__tmp__pdfs__prevention-medicine-past-exams__05-ux8425ux517bux5b66ux5f80ux5e74ux9898ux8003ux70b9ux6574ux7406.md__ux98dfux54c1ux6c61ux67d3ux7684ux6982ux5ff5ux5206ux7c7bux4e0eux751fux7269ux5bccux96c6ux4f5cux7528}

\InfoLabel{考频} 2 次

\InfoLabel{对应小节} 二、食品污染（定义、分类、化学污染特点）

\InfoLabel{匹配依据} 题干为「Food contamination（食品污染）」名词及「生物富集作用」名词，与本节直接对应。

\ContentBand{知识点原文摘取}

食品污染：是指在各种条件下，导致有毒有害物质进入到食品，或食物成分本身发生化学反应而产生有毒有害物质，从而造成食品安全性、营养性和/或感官性状发生改变的过程。按污染物性质分为生物性、物理性、化学性污染物。\\
食品化学性污染的特点之一：污染物的蓄积性强，通过食物链的生物富集作用可在人体内达到很高的浓度，易对健康造成多方面的危害。

\InfoLabel{知识来源} 《营养与食品卫生学·18 柯雅蕾笔记》第 69-70、72 页

\ContentBand{对应往年题}

\begin{pastquestion}

\QuestionLabel{【往年题1｜20级完整试题回忆.docx｜名词解释】}\\
Food contamination（食品污染）

\end{pastquestion}

\begin{pastquestion}

\QuestionLabel{【往年题2｜04级考题回忆.doc（附：其他年份回忆）｜名词解释】}\\
生物富集作用

\end{pastquestion}

\PointSeparator

\section{食品的腐败变质（概念、化学过程、鉴定指标、优势微生物）}\label{users__fpb__26_spring__pkusph_study__tmp__pdfs__prevention-medicine-past-exams__05-ux8425ux517bux5b66ux5f80ux5e74ux9898ux8003ux70b9ux6574ux7406.md__ux98dfux54c1ux7684ux8150ux8d25ux53d8ux8d28ux6982ux5ff5ux5316ux5b66ux8fc7ux7a0bux9274ux5b9aux6307ux6807ux4f18ux52bfux5faeux751fux7269}

\InfoLabel{考频} 7 次

\InfoLabel{对应小节} （一）食品的微生物污染 / 7. 食品的腐败变质；食品添加剂·概述 / 2. 食品的腐败变质

\InfoLabel{匹配依据} 题干涉及「Food
spoilage（食品腐败变质）」名词、蛋白质腐败变质鉴定指标、影响蛋白质腐败因素、脂肪酸败指标、「哈喇味＝脂肪腐败」、食品腐败的主要原因及卫生学意义、「腐败变质中占优势的微生物（细菌）」，与本节直接对应。

\ContentBand{知识点原文摘取}

食品的腐败变质（food
spoilage）：指食品在一定环境因素影响下，由微生物作用而发生的食品成分与感官性状的各种变化。化学过程包括蛋白质分解（产恶臭）、脂肪酸败（产生「哈喇」味，过氧化值、碘价、酸价、羰基价等升高）、碳水化合物分解（产酸产气）。\\
鉴定指标：①感官鉴定；②物理指标；③化学鉴定（挥发性盐基总氮 TVBN、三甲胺、组胺、K
值、pH、过氧化值和酸价等）；④微生物检验（菌落总数和大肠菌群）。\\
影响因素：食品本身组成性质（含酶、营养成分、pH、水分）、微生物及环境因素（温度、湿度、氧气）。在食品腐败变质过程中通常占优势的微生物是细菌。

\InfoLabel{知识来源} 《营养与食品卫生学·18 柯雅蕾笔记》第 72、76 页

\ContentBand{对应往年题}

\begin{pastquestion}

\QuestionLabel{【往年题1｜20级完整试题回忆.docx｜名词解释】}\\
Food spoilage（食品腐败变质）

\end{pastquestion}

\begin{pastquestion}

\QuestionLabel{【往年题2｜04级考题回忆.doc（附：其他年份回忆）｜简答】}\\
蛋白质腐败变质的鉴定指标

\end{pastquestion}

\begin{pastquestion}

\QuestionLabel{【往年题3｜21级营养期末考题.docx｜简答】}\\
食品腐败的主要原因及卫生学意义

\end{pastquestion}

\begin{pastquestion}

\QuestionLabel{【往年题4｜20级完整试题回忆.docx｜单选】}\\
在食品腐败变质过程中通常占优势的微生物是（A 细菌）

\end{pastquestion}

\begin{pastquestion}

\QuestionLabel{【往年题5｜08级考题回忆.doc｜选择】}\\
有哈喇味说明：脂肪腐败

\end{pastquestion}

\begin{pastquestion}

\QuestionLabel{【往年题6｜08级考题回忆.doc｜多选】}\\
脂肪酸腐败（指标）：过氧化价、酸价、碘价、皂化价、羰基价

\end{pastquestion}

\begin{pastquestion}

\QuestionLabel{【往年题7｜08级考题回忆.doc｜多选】}\\
影响蛋白质腐败的因素：水分、PH、营养成分、环境因素、微生物

\end{pastquestion}

\PointSeparator

\section{挥发性盐基总氮（TVBN）}\label{users__fpb__26_spring__pkusph_study__tmp__pdfs__prevention-medicine-past-exams__05-ux8425ux517bux5b66ux5f80ux5e74ux9898ux8003ux70b9ux6574ux7406.md__ux6325ux53d1ux6027ux76d0ux57faux603bux6c2etvbn}

\InfoLabel{考频} 4 次

\InfoLabel{对应小节} （一）食品的微生物污染 / 食品腐败变质的化学鉴定（TVBN）

\InfoLabel{匹配依据} 题干为「挥发性盐基总氮／TVBN／total volatile basic nitrogen」名词解释，与本节直接对应。

\ContentBand{知识点原文摘取}

挥发性盐基总氮（total volatile basic
nitrogen，TVBN）：是指食品水浸液在碱性条件下能与水蒸气一起蒸馏出来的总氮量，即在此条件下能形成氨的含氮物，是食品腐败变质的化学鉴定指标之一。

\InfoLabel{知识来源} 《营养与食品卫生学·18 柯雅蕾笔记》第 72 页

\ContentBand{对应往年题}

\begin{pastquestion}

\QuestionLabel{【往年题1｜04级考题回忆.doc｜名词解释】}\\
TVBN（全英文）

\end{pastquestion}

\begin{pastquestion}

\QuestionLabel{【往年题2｜08级考题回忆.doc｜名词解释】}\\
挥发性盐基总氮（全英文）

\end{pastquestion}

\begin{pastquestion}

\QuestionLabel{【往年题3｜09级考题回忆.docx｜名词解释】}\\
挥发性盐基总氮（全英文）

\end{pastquestion}

\begin{pastquestion}

\QuestionLabel{【往年题4｜14级营养考题.pdf｜名词解释】}\\
total volatile basic nitrogen

\end{pastquestion}

\PointSeparator

\section{菌落总数与大肠菌群（最可能数
MPN、卫生学意义）}\label{users__fpb__26_spring__pkusph_study__tmp__pdfs__prevention-medicine-past-exams__05-ux8425ux517bux5b66ux5f80ux5e74ux9898ux8003ux70b9ux6574ux7406.md__ux83ccux843dux603bux6570ux4e0eux5927ux80a0ux83ccux7fa4ux6700ux53efux80fdux6570-mpnux536bux751fux5b66ux610fux4e49}

\InfoLabel{考频} 3 次

\InfoLabel{对应小节} （一）食品的微生物污染 / 5. 食品的细菌污染（菌落总数、大肠菌群）

\InfoLabel{匹配依据}
题干为「菌落总数」名词、「大肠菌群的卫生学意义」及「大肠菌群最可能数的定义和卫生学意义」，与本节直接对应。

\ContentBand{知识点原文摘取}

菌落总数：指在被检样品的单位质量（g）、容积（ml）内，在严格规定的条件下培养所形成的细菌菌落总数，以菌落形成单位（CFU）表示。卫生学意义：①作为食品被细菌污染程度即清洁状态的标志；②可用于预测食品耐保藏性。\\
大肠菌群最可能数（MPN）：当食品中大肠菌群含量较低时，采用相当于每
g（ml）食品中大肠菌群的最可能近似数标示，是基于泊松分布的间接计数法。卫生学意义：①作为食品受到人与温血动物粪便污染的指示菌；②作为肠道致病菌污染食品的指示菌。

\InfoLabel{知识来源} 《营养与食品卫生学·18 柯雅蕾笔记》第 71 页

\ContentBand{对应往年题}

\begin{pastquestion}

\QuestionLabel{【往年题1｜04级考题回忆.doc（附：其他年份回忆）｜名词解释】}\\
菌落总数

\end{pastquestion}

\begin{pastquestion}

\QuestionLabel{【往年题2｜14级营养考题.pdf｜简答】}\\
大肠菌群的卫生学意义

\end{pastquestion}

\begin{pastquestion}

\QuestionLabel{【往年题3｜20级完整试题回忆.docx｜简答】}\\
大肠菌群最可能数的定义和卫生学意义

\end{pastquestion}

\PointSeparator

\section{食源性疾病与食物中毒（定义、食品安全事故）}\label{users__fpb__26_spring__pkusph_study__tmp__pdfs__prevention-medicine-past-exams__05-ux8425ux517bux5b66ux5f80ux5e74ux9898ux8003ux70b9ux6574ux7406.md__ux98dfux6e90ux6027ux75beux75c5ux4e0eux98dfux7269ux4e2dux6bd2ux5b9aux4e49ux98dfux54c1ux5b89ux5168ux4e8bux6545}

\InfoLabel{考频} 6 次

\InfoLabel{对应小节} 三、食源性疾病（定义、范畴）

\InfoLabel{匹配依据} 题干为「食源性疾病／Foodborne
disease」「食物中毒」「食品安全事故」名词解释，与本节直接对应（食品安全事故为部分匹配）。

\ContentBand{知识点原文摘取}

食源性疾病（foodborne
disease，FBD，WHO）：通过摄食进入人体内的各种致病因子引起的，通常具有感染性质或中毒性质的一类疾病；（食品安全法）食品中致病因素进入人体引起的感染性、中毒性等疾病，包括食物中毒。\\
食物中毒（food poisoning）：由生物性、化学性病原物引起，急性或亚急性、非传染性，属食源性疾病的范畴之一。

\InfoLabel{知识来源} 《营养与食品卫生学·18 柯雅蕾笔记》第 73 页

\ContentBand{对应往年题}

\begin{pastquestion}

\QuestionLabel{【往年题1｜04级考题回忆.doc（附：其他年份回忆）｜名词解释】}\\
Food borne disease（食源性疾病）

\end{pastquestion}

\begin{pastquestion}

\QuestionLabel{【往年题2｜14级营养考题.pdf｜名词解释】}\\
Foodborne disease

\end{pastquestion}

\begin{pastquestion}

\QuestionLabel{【往年题3｜21级营养期末考题.docx｜名词解释】}\\
Foodborne disease

\end{pastquestion}

\begin{pastquestion}

\QuestionLabel{【往年题4｜04级考题回忆.doc｜名词解释】}\\
食物中毒（全英文）

\end{pastquestion}

\begin{pastquestion}

\QuestionLabel{【往年题5｜09级考题回忆.docx｜名词解释】}\\
食物中毒（全英文）

\end{pastquestion}

\begin{pastquestion}

\QuestionLabel{【往年题6｜08级考题回忆.doc｜名词解释】}\\
食品安全事故（全英文，部分匹配）

\end{pastquestion}

\PointSeparator

\section{细菌性食物中毒（沙门氏菌与葡萄球菌、预防、金黄色葡萄球菌中毒案例）}\label{users__fpb__26_spring__pkusph_study__tmp__pdfs__prevention-medicine-past-exams__05-ux8425ux517bux5b66ux5f80ux5e74ux9898ux8003ux70b9ux6574ux7406.md__ux7ec6ux83ccux6027ux98dfux7269ux4e2dux6bd2ux6c99ux95e8ux6c0fux83ccux4e0eux8461ux8404ux7403ux83ccux9884ux9632ux91d1ux9ec4ux8272ux8461ux8404ux7403ux83ccux4e2dux6bd2ux6848ux4f8b}

\InfoLabel{考频} 4 次

\InfoLabel{对应小节} 三、食源性疾病 /（一）细菌性食物中毒

\InfoLabel{匹配依据}
题干涉及沙门氏菌与葡萄球菌两种食物中毒的比较、细菌性食物中毒的预防措施，以及金黄色葡萄球菌/诺如病毒集体中毒案例分析，与本节直接对应。

\ContentBand{知识点原文摘取}

细菌性食物中毒：指摄入被致病性细菌或其毒素污染的食品而引起的食源性疾病。发病原因：致病菌污染食物、储藏方式不当、烹调加工方式不当。发病率高、死亡率一般较低，以夏秋季高发。\\
分感染型（沙门菌、变形杆菌、志贺氏菌等，伴发热）、毒素型（金黄色葡萄球菌、蜡样芽孢杆菌、肉毒梭菌等，以上消化道症状为主、发热少见）、混合型（副溶血弧菌）。\\
预防 5 Keys：保持清洁、生熟分开、彻底煮熟（＞70℃）、安全温度保存（＜5℃或＞60℃）、使用安全水和食材。

\InfoLabel{知识来源} 《营养与食品卫生学·18 柯雅蕾笔记》第 74-75 页

\ContentBand{对应往年题}

\begin{pastquestion}

\QuestionLabel{【往年题1｜04级考题回忆.doc（附：其他年份回忆）｜简答】}\\
两种食物中毒比较（沙门氏菌，葡萄球菌）

\end{pastquestion}

\begin{pastquestion}

\QuestionLabel{【往年题2｜08级考题回忆.doc｜选择】}\\
细菌性食物中毒的预防措施不包括：抗氧化？低温？加热？工作人员定期检查？

\end{pastquestion}

\begin{pastquestion}

\QuestionLabel{【往年题3｜20级完整试题回忆.docx｜论述】}\\
某市三个乡镇 5 所小学和 5 所幼儿园 1176
名学生陆续出现恶心、上腹痛、剧烈呕吐、腹泻，不发热，均饮用了某牛奶饮料（金黄色葡萄球菌中毒）。问：（1）考虑的诊断和处理；（2）依据食品安全事故流行病学调查工作流程应采取的行动；（3）如何预防类似中毒事件

\end{pastquestion}

\begin{pastquestion}

\QuestionLabel{【往年题4｜21级营养期末考题.docx｜论述】}\\
以清华大学校友聚餐诺如病毒中毒案例为原型（考生写成金黄色葡萄球菌），呕吐、腹泻、发热、乏力、肌肉酸痛。问：（1）中毒类型＋传播途径；（2）流行病学＋卫生学调查；（3）以团体餐为例校园食品安全主要风险控制点；（4）加强校园食品安全的措施

\end{pastquestion}

\PointSeparator

\section{肉毒梭菌/肉毒毒素中毒}\label{users__fpb__26_spring__pkusph_study__tmp__pdfs__prevention-medicine-past-exams__05-ux8425ux517bux5b66ux5f80ux5e74ux9898ux8003ux70b9ux6574ux7406.md__ux8089ux6bd2ux68adux83ccux8089ux6bd2ux6bd2ux7d20ux4e2dux6bd2}

\InfoLabel{考频} 3 次

\InfoLabel{对应小节} 三、食源性疾病 /（一）细菌性食物中毒（毒素型）

\InfoLabel{匹配依据}
题干为肉毒毒素中毒案例（自制豆豉/酱豆腐、神经毒性表现）的诊断、毒素、可疑食物及救治，与本节毒素型细菌性食物中毒直接对应。

\ContentBand{知识点原文摘取}

毒素型细菌性食物中毒：病原菌污染食物、大量繁殖产生肠毒素（外毒素），代表菌包括金黄色葡萄球菌、蜡样芽孢杆菌、肉毒梭菌等。\\
肉毒梭菌多见于发酵豆/谷制品（如自制豆豉、臭豆腐），产生肉毒毒素，为神经毒素，可致眼、舌、咽肌及呼吸肌麻痹。肉毒梭菌引起的食源性疾病病死率较高、病情重、恢复缓慢。地区性：肉毒梭菌------新疆、青海。

\InfoLabel{知识来源} 《营养与食品卫生学·18 柯雅蕾笔记》第 74 页

\ContentBand{对应往年题}

\begin{pastquestion}

\QuestionLabel{【往年题1｜04级考题回忆.doc｜应用题】}\\
一个案例，食物中毒表现为视力模糊、眼舌麻痹、呼吸麻痹死亡，伴消化道症状，食用自制豆豉酱。请分析中毒物质（肉毒毒素）和最可能的中毒食物，阐明原因，简述确诊和救治方法

\end{pastquestion}

\begin{pastquestion}

\QuestionLabel{【往年题2｜09级考题回忆.docx｜论述】}\\
肉毒毒素中毒案例：写出诊断依据、什么毒素、什么食物，还需做什么诊断、做什么治疗

\end{pastquestion}

\begin{pastquestion}

\QuestionLabel{【往年题3｜08级考题回忆.doc｜选择】}\\
肉毒梭菌小病例（酱豆腐，有神经毒性表现）

\end{pastquestion}

\PointSeparator

\section{真菌毒素------黄曲霉毒素}\label{users__fpb__26_spring__pkusph_study__tmp__pdfs__prevention-medicine-past-exams__05-ux8425ux517bux5b66ux5f80ux5e74ux9898ux8003ux70b9ux6574ux7406.md__ux771fux83ccux6bd2ux7d20ux9ec4ux66f2ux9709ux6bd2ux7d20}

\InfoLabel{考频} 1 次

\InfoLabel{对应小节} （一）食品的微生物污染 / 6. 真菌与真菌毒素

\InfoLabel{匹配依据}
题干为「黄曲霉毒素代谢与毒性，去毒方法」。笔记述及真菌毒素及黄曲霉毒素为重点关注对象，但未详述其代谢、毒性与去毒方法，为部分匹配。

\ContentBand{知识点原文摘取}

真菌毒素通常具有耐高温、无抗原性、主要侵害实质器官的特性。常见真菌毒素有黄曲霉毒素、镰刀菌毒素、玉米赤霉烯酮、赭曲霉毒素、展青霉素等。\\
食品安全风险评估重点关注食品如黄曲霉素-花生油。\\
（黄曲霉毒素的代谢、毒性与去毒方法未完全匹配到笔记原文。）

\InfoLabel{知识来源} 《营养与食品卫生学·18 柯雅蕾笔记》第 72、64 页

\ContentBand{对应往年题}

\begin{pastquestion}

\QuestionLabel{【往年题1｜04级考题回忆.doc（附：其他年份回忆）｜简答】}\\
黄曲霉毒素代谢与毒性，去毒方法（部分匹配）

\end{pastquestion}

\PointSeparator

\section{有毒动植物中毒（苦杏仁、鲜黄花菜、河豚毒素与小鼠单位）}\label{users__fpb__26_spring__pkusph_study__tmp__pdfs__prevention-medicine-past-exams__05-ux8425ux517bux5b66ux5f80ux5e74ux9898ux8003ux70b9ux6574ux7406.md__ux6709ux6bd2ux52a8ux690dux7269ux4e2dux6bd2ux82e6ux674fux4ec1ux9c9cux9ec4ux82b1ux83dcux6cb3ux8c5aux6bd2ux7d20ux4e0eux5c0fux9f20ux5355ux4f4d}

\InfoLabel{考频} 3 次

\InfoLabel{对应小节} 三、食源性疾病 /（四）有毒动植物中毒

\InfoLabel{匹配依据} 题干涉及苦杏仁中毒（氰苷）、鲜黄花菜毒素（类秋水仙碱）及「Mouse
unit（小鼠单位）」名词。河豚毒素以小鼠单位计效价，名词为部分匹配。

\ContentBand{知识点原文摘取}

有毒动植物中毒：常见有河鲀鱼（河豚毒素
TTX）、腐败鱼类（组胺）、毒蕈、苦杏仁/桃仁/木薯等含氰苷植物、发芽马铃薯（龙葵素）、鲜黄花菜（类秋水仙碱）、豇豆/刀豆/扁豆（皂素）等。\\
苦杏仁、木薯中的氰苷类为有毒成分。\\
河豚毒素（TTX）：神经毒素，剧毒，对人的致死量为 20
万小鼠单位；卵巢最毒，肝脏次之。中毒机制为抑制神经细胞膜对钠离子的通透性，阻断神经冲动传导。

\InfoLabel{知识来源} 《营养与食品卫生学·18 柯雅蕾笔记》第 74-75 页

\ContentBand{对应往年题}

\begin{pastquestion}

\QuestionLabel{【往年题1｜04级考题回忆.doc（附：其他年份回忆）｜论述】}\\
苦杏仁中毒的成分、机理、解救措施，不及时解救的后续症状

\end{pastquestion}

\begin{pastquestion}

\QuestionLabel{【往年题2｜08级考题回忆.doc｜选择】}\\
鲜黄花菜的毒素：类秋水仙素

\end{pastquestion}

\begin{pastquestion}

\QuestionLabel{【往年题3｜04级考题回忆.doc（附：其他年份回忆）｜名词解释】}\\
Mouse unit（小鼠单位，部分匹配）

\end{pastquestion}

\PointSeparator

\section{食源性化学性中毒------亚硝酸盐与
N-亚硝基化合物}\label{users__fpb__26_spring__pkusph_study__tmp__pdfs__prevention-medicine-past-exams__05-ux8425ux517bux5b66ux5f80ux5e74ux9898ux8003ux70b9ux6574ux7406.md__ux98dfux6e90ux6027ux5316ux5b66ux6027ux4e2dux6bd2ux4e9aux785dux9178ux76d0ux4e0e-n-ux4e9aux785dux57faux5316ux5408ux7269}

\InfoLabel{考频} 5 次

\InfoLabel{对应小节} 三、食源性疾病 /（五）食源性化学性中毒；食品的化学污染；食品添加剂·护色剂

\InfoLabel{匹配依据}
题干涉及「N-亚硝基化合物」名词、亚硝基化合物/亚硝酸盐的分类、前体、合成部位及影响因素，以及亚硝酸盐中毒病例分析。笔记述及
N-亚硝基化合物属化学污染、亚硝酸盐为亚硝胺前体，但未系统列出分类、前体、合成部位，相关题为部分匹配。

\ContentBand{知识点原文摘取}

食品化学性污染来源包括 N-亚硝基化合物（如 N-亚硝基二甲胺）。\\
护色剂（发色剂）成分为硝酸钠（钾）与亚硝酸钠（钾）；大量使用引起肠原性青紫症；NO₂⁻ 是亚硝胺的前体。\\
化学性食物中毒：由食用被有毒有害化学物质污染（如亚硝酸盐）的食品所引起。\\
（N-亚硝基化合物的分类、前体、人体合成部位及影响因素未完全匹配到笔记原文，为部分匹配。）

\InfoLabel{知识来源} 《营养与食品卫生学·18 柯雅蕾笔记》第 70、75、78 页

\ContentBand{对应往年题}

\begin{pastquestion}

\QuestionLabel{【往年题1｜08级考题回忆.doc｜名词解释】}\\
N-亚硝基化合物（全英文）

\end{pastquestion}

\begin{pastquestion}

\QuestionLabel{【往年题2｜04级考题回忆.doc｜简答】}\\
亚硝基化合物的分类、前体、在人体的合成部位及影响因素（部分匹配）

\end{pastquestion}

\begin{pastquestion}

\QuestionLabel{【往年题3｜09级考题回忆.docx｜简答】}\\
亚硝酸盐分类、前体、形成影响因素（部分匹配）

\end{pastquestion}

\begin{pastquestion}

\QuestionLabel{【往年题4｜08级考题回忆.doc｜多选】}\\
影响 N-亚硝基化合物（形成的因素）：碱性、酸性、水分、亚硝酸盐含量、硫氰根含量

\end{pastquestion}

\begin{pastquestion}

\QuestionLabel{【往年题5｜08级考题回忆.doc｜问答】}\\
亚硝酸盐中毒：（1）诊断（紫绀和消化道症状）；（2）机理及治疗措施；（3）如果你是 CDC
的工作人员，接到报告后应该做什么准备？

\end{pastquestion}

\PointSeparator

\chapter{食品添加剂及食品卫生}\label{users__fpb__26_spring__pkusph_study__tmp__pdfs__prevention-medicine-past-exams__05-ux8425ux517bux5b66ux5f80ux5e74ux9898ux8003ux70b9ux6574ux7406.md__ux7b2cux5341ux56dbux7ae0-ux98dfux54c1ux6dfbux52a0ux5242ux53caux98dfux54c1ux536bux751f}

\section{食品添加剂的定义}\label{users__fpb__26_spring__pkusph_study__tmp__pdfs__prevention-medicine-past-exams__05-ux8425ux517bux5b66ux5f80ux5e74ux9898ux8003ux70b9ux6574ux7406.md__ux98dfux54c1ux6dfbux52a0ux5242ux7684ux5b9aux4e49}

\InfoLabel{考频} 3 次

\InfoLabel{对应小节} 二、食品添加剂 / 2. 食品添加剂（food additives）

\InfoLabel{匹配依据} 题干为「Food additives（食品添加剂）」名词及「关于 WHO
对食品添加剂的定义理解，错误的是」，与本节直接对应。

\ContentBand{知识点原文摘取}

食品添加剂（food
additives，GB2760-2014）：为改善食品品质和色、香、味，以及为防腐、保鲜和加工工艺的需要而加入食品中的人工合成或者天然物质。食品用香料、胶基糖果中基础剂物质、食品工业用加工助剂也包括在内。\\
使用原则：不应对人体产生任何健康危害；不应掩盖食品的腐败变质；不应掩盖质量缺陷或以掺杂、掺假、伪造为目的；不应降低食品营养价值；在达到预期效果前提下尽可能降低使用量。

\InfoLabel{知识来源} 《营养与食品卫生学·18 柯雅蕾笔记》第 76-77 页

\ContentBand{对应往年题}

\begin{pastquestion}

\QuestionLabel{【往年题1｜20级完整试题回忆.docx｜名词解释】}\\
Food additives（食品添加剂）

\end{pastquestion}

\begin{pastquestion}

\QuestionLabel{【往年题2｜14级营养考题.pdf｜名词解释】}\\
Food additive（食品添加剂）

\end{pastquestion}

\begin{pastquestion}

\QuestionLabel{【往年题3｜20级完整试题回忆.docx｜单选】}\\
关于 WHO 对食品添加剂的定义理解，错误的是

\end{pastquestion}

\PointSeparator

\section{ADI 值与 GRAS
类}\label{users__fpb__26_spring__pkusph_study__tmp__pdfs__prevention-medicine-past-exams__05-ux8425ux517bux5b66ux5f80ux5e74ux9898ux8003ux70b9ux6574ux7406.md__adi-ux503cux4e0e-gras-ux7c7b}

\InfoLabel{考频} 4 次

\InfoLabel{对应小节} 二、食品添加剂 / 3. ADI 值；4. 联合国 FAO/WHO 对食品添加剂的管理

\InfoLabel{匹配依据} 题干为「ADI（每日容许摄入量）」名词及「关于食品添加剂 GRAS
类说法错误的是」，与本节直接对应。

\ContentBand{知识点原文摘取}

人体每日容许摄入量（acceptable daily
intake，ADI）：是指人类终生摄入该物质后而对机体不产生任何已知不良效应的剂量，以人体每公斤体重的该物质摄入量来表示（mg/kg）。\\
JECFA 将食品添加剂分四类：第一类即一般认为是安全的（GRAS）物质，可以按正常需要使用，不需建立 ADI 值；A
类（A1/A2）、B 类、C 类（原则禁止使用）。

\InfoLabel{知识来源} 《营养与食品卫生学·18 柯雅蕾笔记》第 76-77 页

\ContentBand{对应往年题}

\begin{pastquestion}

\QuestionLabel{【往年题1｜04级考题回忆.doc｜名词解释】}\\
ADI

\end{pastquestion}

\begin{pastquestion}

\QuestionLabel{【往年题2｜04级考题回忆.doc（附：其他年份回忆）｜名词解释】}\\
Acceptable daily intake（ADI）

\end{pastquestion}

\begin{pastquestion}

\QuestionLabel{【往年题3｜21级营养期末考题.docx｜名词解释】}\\
ADI

\end{pastquestion}

\begin{pastquestion}

\QuestionLabel{【往年题4｜20级完整试题回忆.docx｜单选】}\\
关于食品添加剂 GRAS 类，说法错误的是（A 可以随意使用）

\end{pastquestion}

\PointSeparator

\section{着色剂与人工色素的使用}\label{users__fpb__26_spring__pkusph_study__tmp__pdfs__prevention-medicine-past-exams__05-ux8425ux517bux5b66ux5f80ux5e74ux9898ux8003ux70b9ux6574ux7406.md__ux7740ux8272ux5242ux4e0eux4ebaux5de5ux8272ux7d20ux7684ux4f7fux7528}

\InfoLabel{考频} 1 次

\InfoLabel{对应小节} 二、食品添加剂 / 6. 常见的食品添加剂 /（3）着色剂

\InfoLabel{匹配依据} 题干为「人工色素不能用于下列哪类食物（婴儿食品）」，与本节着色剂使用相对应。

\ContentBand{知识点原文摘取}

着色剂（色素）：使食品着色后提高其感官性状的一类物质，分食用天然色素与合成色素。合成色素成本低廉、色泽鲜艳、着色力强，但有一定毒性。\\
（按规定，人工合成色素禁止用于婴幼儿食品等。）

\InfoLabel{知识来源} 《营养与食品卫生学·18 柯雅蕾笔记》第 78 页

\ContentBand{对应往年题}

\begin{pastquestion}

\QuestionLabel{【往年题1｜08级考题回忆.doc｜选择】}\\
人工色素不能用于下列哪类食物？肉类、蛋类、奶制品、婴儿食品

\end{pastquestion}

\PointSeparator

\section{防腐剂（山梨酸钾与苯甲酸钠的选择）}\label{users__fpb__26_spring__pkusph_study__tmp__pdfs__prevention-medicine-past-exams__05-ux8425ux517bux5b66ux5f80ux5e74ux9898ux8003ux70b9ux6574ux7406.md__ux9632ux8150ux5242ux5c71ux68a8ux9178ux94beux4e0eux82efux7532ux9178ux94a0ux7684ux9009ux62e9}

\InfoLabel{考频} 1 次

\InfoLabel{对应小节} 二、食品添加剂 / 6. 常见的食品添加剂 /（7）防腐剂

\InfoLabel{匹配依据}
题干为「在山梨酸钾和苯甲酸钠中选择一种作为防腐剂，哪种更好，为什么」，与本节防腐剂相对应。

\ContentBand{知识点原文摘取}

防腐剂：指为防止食品腐败、变质，延长食品保存期，抑制食品中微生物繁殖的物质。\\
酸性防腐剂：苯甲酸及其钠盐（日本进口产品受到限制）、山梨酸及其盐类；酯型防腐剂：对羟基苯甲酸酯类。

\InfoLabel{知识来源} 《营养与食品卫生学·18 柯雅蕾笔记》第 78-79 页

\ContentBand{对应往年题}

\begin{pastquestion}

\QuestionLabel{【往年题1｜14级营养考题.pdf｜简答】}\\
在山梨酸钾和苯甲酸钠中选择一种作为防腐剂时你认为哪种更好，为什么

\end{pastquestion}

\PointSeparator

\section{护色剂（发色剂）------亚硝酸盐的使用范围}\label{users__fpb__26_spring__pkusph_study__tmp__pdfs__prevention-medicine-past-exams__05-ux8425ux517bux5b66ux5f80ux5e74ux9898ux8003ux70b9ux6574ux7406.md__ux62a4ux8272ux5242ux53d1ux8272ux5242ux4e9aux785dux9178ux76d0ux7684ux4f7fux7528ux8303ux56f4}

\InfoLabel{考频} 1 次

\InfoLabel{对应小节} 二、食品添加剂 / 6. 常见的食品添加剂 /（4）护色剂

\InfoLabel{匹配依据} 题干为「亚硝酸钠可以用于以下哪些食品添加（腌制蔬菜、预制肉）」，与本节护色剂相对应。

\ContentBand{知识点原文摘取}

护色剂（发色剂）：使肉及肉制品呈现良好色泽的物质，成分为硝酸钠（钾）与亚硝酸钠（钾）。\\
作用：护色、抑制微生物增殖、对肉毒梭状芽孢杆菌有特殊抑制作用、提高腌肉风味。大量使用引起肠原性青紫症，NO₂⁻
是亚硝胺的前体。（亚硝酸盐主要用于腌腊肉制品、酱卤肉、预制肉等。）

\InfoLabel{知识来源} 《营养与食品卫生学·18 柯雅蕾笔记》第 78 页

\ContentBand{对应往年题}

\begin{pastquestion}

\QuestionLabel{【往年题1｜21级营养期末考题.docx｜选择】}\\
亚硝酸钠可以用于以下食品添加：A 面包；B 糕点；C 烘焙；D 腌制蔬菜；E 预制肉

\end{pastquestion}

\PointSeparator

\chapter{食品安全法律法规体系}\label{users__fpb__26_spring__pkusph_study__tmp__pdfs__prevention-medicine-past-exams__05-ux8425ux517bux5b66ux5f80ux5e74ux9898ux8003ux70b9ux6574ux7406.md__ux7b2cux5341ux4e94ux7ae0-ux98dfux54c1ux5b89ux5168ux6cd5ux5f8bux6cd5ux89c4ux4f53ux7cfb}

\section{食品安全监督管理机构（机构改革）}\label{users__fpb__26_spring__pkusph_study__tmp__pdfs__prevention-medicine-past-exams__05-ux8425ux517bux5b66ux5f80ux5e74ux9898ux8003ux70b9ux6574ux7406.md__ux98dfux54c1ux5b89ux5168ux76d1ux7763ux7ba1ux7406ux673aux6784ux673aux6784ux6539ux9769}

\InfoLabel{考频} 1 次

\InfoLabel{对应小节} 食品安全法律法规体系（监管机构）

\InfoLabel{匹配依据} 题干为「2018 年机构改革后管理食品卫生的机构是哪个」，与本节监管机构设置相对应。

\ContentBand{知识点原文摘取}

食品安全监管工作由国家市场监督管理总局（由原食品药品监督管理总局、国家质量监督检验检疫总局、国家工商行政管理总局组成）、农业部、公安部、出入境（海关总署）等部门承担。\\
食品安全风险管理和技术支撑工作（风险监测、风险评估，制定并公布食品安全国家标准）主要由国家卫生健康委承担。

\InfoLabel{知识来源} 《营养与食品卫生学·18 柯雅蕾笔记》第 83-84 页

\ContentBand{对应往年题}

\begin{pastquestion}

\QuestionLabel{【往年题1｜14级营养考题.pdf｜选择】}\\
2018 年机构改革后管理食品卫生的机构是哪个

\end{pastquestion}

\PointSeparator

\section{食品生产经营者依法从事生产经营}\label{users__fpb__26_spring__pkusph_study__tmp__pdfs__prevention-medicine-past-exams__05-ux8425ux517bux5b66ux5f80ux5e74ux9898ux8003ux70b9ux6574ux7406.md__ux98dfux54c1ux751fux4ea7ux7ecfux8425ux8005ux4f9dux6cd5ux4eceux4e8bux751fux4ea7ux7ecfux8425}

\InfoLabel{考频} 1 次

\InfoLabel{对应小节} 一、食品安全法律法规体系（《食品安全法》总则）

\InfoLabel{匹配依据} 题干为「食品生产经营者应当依照（）从事生产经营活动」，与本节《食品安全法》规定直接对应。

\ContentBand{知识点原文摘取}

食品生产经营者应当依照法律、法规和食品安全标准从事生产经营活动，保证食品安全，诚信自律，对社会和公众负责，接受社会监督，承担社会责任。\\
食品生产经营者对其生产经营食品的安全负责。

\InfoLabel{知识来源} 《营养与食品卫生学·18 柯雅蕾笔记》第 84 页

\ContentBand{对应往年题}

\begin{pastquestion}

\QuestionLabel{【往年题1｜20级完整试题回忆.docx｜单选】}\\
我国食品安全法规定，食品生产经营者应当依照（C 法律、法规和食品安全标准）从事生产经营活动\ldots\ldots{}

\end{pastquestion}

\PointSeparator

\section{国际食品法典（CAC/Codex）}\label{users__fpb__26_spring__pkusph_study__tmp__pdfs__prevention-medicine-past-exams__05-ux8425ux517bux5b66ux5f80ux5e74ux9898ux8003ux70b9ux6574ux7406.md__ux56fdux9645ux98dfux54c1ux6cd5ux5178caccodex}

\InfoLabel{考频} 2 次

\InfoLabel{对应小节} 二、食品安全标准 / 6. 国际食品标准（Codex 标准）

\InfoLabel{匹配依据} 题干为「CAC（国际食品法典委员会）」「国际食品法典」名词。笔记仅提及 Codex 标准，未给出
CAC 的完整定义，故为部分匹配。

\ContentBand{知识点原文摘取}

国际食品标准：Codex 标准；各国食品标准存在差异；我国已基本建立与国际接轨的食品安全国家标准体系。\\
（国际食品法典委员会（CAC）的完整定义未完全匹配到笔记原文，为部分匹配。）

\InfoLabel{知识来源} 《营养与食品卫生学·18 柯雅蕾笔记》第 85 页

\ContentBand{对应往年题}

\begin{pastquestion}

\QuestionLabel{【往年题1｜04级考题回忆.doc｜名词解释】}\\
CAC（食品法典监督局？／国际食品法典委员会，部分匹配）

\end{pastquestion}

\begin{pastquestion}

\QuestionLabel{【往年题2｜09级考题回忆.docx｜名词解释】}\\
国际食品法典（全英文，部分匹配）

\end{pastquestion}

\PointSeparator

\chapter{食品安全监督管理}\label{users__fpb__26_spring__pkusph_study__tmp__pdfs__prevention-medicine-past-exams__05-ux8425ux517bux5b66ux5f80ux5e74ux9898ux8003ux70b9ux6574ux7406.md__ux7b2cux5341ux516dux7ae0-ux98dfux54c1ux5b89ux5168ux76d1ux7763ux7ba1ux7406}

\section{食品安全监督管理的原则与内容（分段监管）}\label{users__fpb__26_spring__pkusph_study__tmp__pdfs__prevention-medicine-past-exams__05-ux8425ux517bux5b66ux5f80ux5e74ux9898ux8003ux70b9ux6574ux7406.md__ux98dfux54c1ux5b89ux5168ux76d1ux7763ux7ba1ux7406ux7684ux539fux5219ux4e0eux5185ux5bb9ux5206ux6bb5ux76d1ux7ba1}

\InfoLabel{考频} 3 次

\InfoLabel{对应小节} 三、食品安全监督管理主要内容（监管体制、监管措施）

\InfoLabel{匹配依据}
题干为食品监管的重要原则与内容、食品监督原则及内容、食品安全管理的分段及其各自负责的工作，与本节直接对应。

\ContentBand{知识点原文摘取}

基本原则：预防为主、风险管理、全程控制、社会共治。总要求：最严谨的标准，最严格的监管，最严厉的处罚，最严肃的问责。\\
分段（分部门）监管体制：①市场监督管理部门（食品、食品添加剂、食品相关产品）；②卫生行政部门（风险监测和风险评估，会同制定食品安全国家标准）；③农业行政部门（食用农产品）；④出入境检验检疫部门---海关（进出口食品）。\\
监管主要措施：行政许可、行政检查、行政处罚、监督抽检、事故处置等。

\InfoLabel{知识来源} 《营养与食品卫生学·18 柯雅蕾笔记》第 83-86 页

\ContentBand{对应往年题}

\begin{pastquestion}

\QuestionLabel{【往年题1｜04级考题回忆.doc｜简答】}\\
食品监管的重要原则和重要内容

\end{pastquestion}

\begin{pastquestion}

\QuestionLabel{【往年题2｜09级考题回忆.docx｜简答】}\\
食品监督原则以及内容

\end{pastquestion}

\begin{pastquestion}

\QuestionLabel{【往年题3｜08级考题回忆.doc｜简答】}\\
食品安全管理的分段及其各自负责的工作

\end{pastquestion}

\PointSeparator

\section{危害分析与关键控制点（HACCP）}\label{users__fpb__26_spring__pkusph_study__tmp__pdfs__prevention-medicine-past-exams__05-ux8425ux517bux5b66ux5f80ux5e74ux9898ux8003ux70b9ux6574ux7406.md__ux5371ux5bb3ux5206ux6790ux4e0eux5173ux952eux63a7ux5236ux70b9haccp}

\InfoLabel{考频} 1 次

\InfoLabel{对应小节} 三、食品安全监督管理主要内容 / 食品生产管理（HACCP）

\InfoLabel{匹配依据} 题干为「下列哪项不属于危害分析关键控制点的基本原理和组成（人员培训）」，与本节直接对应。

\ContentBand{知识点原文摘取}

危害分析与关键控制点（HACCP）：是预防性的食品安全控制体系，通过对食品生产的整个过程进行分析，找出对食品安全有影响的环节，确定关键性的控制点，并为每个关键控制点确定衡量限制和监控程序，一旦出现问题马上采取纠正和控制措施消除隐患。\\
基本原理包括危害分析、确定关键控制点、确定监控措施、建立纠偏措施等。

\InfoLabel{知识来源} 《营养与食品卫生学·18 柯雅蕾笔记》第 87 页

\ContentBand{对应往年题}

\begin{pastquestion}

\QuestionLabel{【往年题1｜20级完整试题回忆.docx｜单选】}\\
下列哪项不属于危害分析关键控制点的基本原理和组成（D 人员培训）

\end{pastquestion}

\PointSeparator

\chapter*{附：未匹配到明确考点的往年题}\label{users__fpb__26_spring__pkusph_study__tmp__pdfs__prevention-medicine-past-exams__05-ux8425ux517bux5b66ux5f80ux5e74ux9898ux8003ux70b9ux6574ux7406.md__ux9644ux672aux5339ux914dux5230ux660eux786eux8003ux70b9ux7684ux5f80ux5e74ux9898}
\addcontentsline{toc}{chapter}{附：未匹配到明确考点的往年题}

\section{环境刺激素（环境激素/内分泌干扰物）对人体健康的影响}\label{users__fpb__26_spring__pkusph_study__tmp__pdfs__prevention-medicine-past-exams__05-ux8425ux517bux5b66ux5f80ux5e74ux9898ux8003ux70b9ux6574ux7406.md__ux73afux5883ux523aux6fc0ux7d20ux73afux5883ux6fc0ux7d20ux5185ux5206ux6cccux5e72ux6270ux7269ux5bf9ux4ebaux4f53ux5065ux5eb7ux7684ux5f71ux54cd}

\InfoLabel{考频} 1 次

\InfoLabel{对应小节} 未匹配到（超出本营养学笔记范围）

\InfoLabel{匹配依据}
题干为「环境刺激素对人体健康的影响」，指环境激素（内分泌干扰物），属环境健康学内容，《营养与食品卫生学·18
柯雅蕾笔记》中无对应知识点，标记为未匹配到。

\ContentBand{知识点原文摘取}

未匹配到。建议参考《环境健康学》环境内分泌干扰物相关内容。

\InfoLabel{知识来源} 无（笔记未收录）

\ContentBand{对应往年题}

\begin{pastquestion}

\QuestionLabel{【往年题1｜08级考题回忆.doc｜简答】}\\
环境刺激素对人体健康的影响（未匹配到）

\end{pastquestion}

\PointSeparator

\chapter*{题量核对}\label{users__fpb__26_spring__pkusph_study__tmp__pdfs__prevention-medicine-past-exams__05-ux8425ux517bux5b66ux5f80ux5e74ux9898ux8003ux70b9ux6574ux7406.md__ux9898ux91cfux6838ux5bf9}
\addcontentsline{toc}{chapter}{题量核对}

按来源文件统计，共 175 题，全部纳入，无遗漏：

\begin{longtable}[]{@{}
  >{\raggedright\arraybackslash}p{(\linewidth - 2\tabcolsep) * \real{0.7000}}
  >{\raggedleft\arraybackslash}p{(\linewidth - 2\tabcolsep) * \real{0.1200}}@{}}
\toprule\noalign{}
\begin{minipage}[b]{\linewidth}\raggedright
来源文件
\end{minipage} & \begin{minipage}[b]{\linewidth}\raggedleft
题量
\end{minipage} \\
\midrule\noalign{}
\endhead
\bottomrule\noalign{}
\endlastfoot
04级考题回忆.doc（含「04 营养真题」20 题 + 「其他年份回忆」20 题） & 40 \\
05级考题回忆.docx & 4 \\
08级考题回忆.doc & 33 \\
09级考题回忆.docx & 17 \\
14级营养考题.pdf & 27 \\
20级完整试题回忆.docx & 34 \\
21级营养期末考题.docx & 20 \\
\textbf{合计} & \textbf{175} \\
\end{longtable}

\begin{quote}
说明：09 级、21
级试卷回忆中各有「还有一个简答/名词解释想不起来了」的零散记载，因无具体题干内容、无法成条，未计入上述条目；如后续补全可另行追加。
\end{quote}

\SetSubjectAccent{7A4EC9}{精神病学}

\part{精神病学}

\chapter*{说明}\label{users__fpb__26_spring__pkusph_study__tmp__pdfs__prevention-medicine-past-exams__06-ux7cbeux795eux75c5ux5b66ux5f80ux5e74ux9898ux8003ux70b9ux6574ux7406.md__ux8bf4ux660e}
\addcontentsline{toc}{chapter}{说明}

\begin{itemize}
\tightlist
\item
  本文件根据《精神病学复习整理by18级杨子铭.docx》和「精神病学/4 往年题」文件夹中全部往年题文件整理。
\item
  知识点正文均摘自《精神病学复习整理by18级杨子铭》原文，未自行扩写；原文较长时摘取与考题最相关部分。
\item
  往年题来源：06-07级史前回忆、16级考题及参考答案、20级完整试题回忆、21级考题与主观题总结。
\item
  本次共拆分 \textbf{104} 个题目条目；每题均已归入至少一个考点（含「未完全匹配到」者）。
\item
  排序依据：先按知识点总结章节顺序，再在同一章节内按考频从高到低排列；同频按收录顺序。
\item
  一道考题若考查多个知识点，可分别归入多个考点并各计 1 次。
\item
  未在总结文件中完全匹配者标记为「未完全匹配到」或归入「附」章。
\item
  已反向核对各来源文件拆分条目，确保无题目遗漏（详见文末「题量核对」）。
\end{itemize}

\PointSeparator

\chapter{绪论}\label{users__fpb__26_spring__pkusph_study__tmp__pdfs__prevention-medicine-past-exams__06-ux7cbeux795eux75c5ux5b66ux5f80ux5e74ux9898ux8003ux70b9ux6574ux7406.md__ux7b2cux4e00ux7ae0-ux7eeaux8bba}

\section{精神障碍、精神病与神经症的概念}\label{users__fpb__26_spring__pkusph_study__tmp__pdfs__prevention-medicine-past-exams__06-ux7cbeux795eux75c5ux5b66ux5f80ux5e74ux9898ux8003ux70b9ux6574ux7406.md__ux7cbeux795eux969cux788dux7cbeux795eux75c5ux4e0eux795eux7ecfux75c7ux7684ux6982ux5ff5}

\InfoLabel{考频} 4 次

\InfoLabel{对应小节} 一、绪论 / 精神障碍与精神病

\InfoLabel{匹配依据} 题干涉及「精神障碍」「精神病」「神经症」「精神病性障碍」「非精神病性」等表述，与一、绪论
/ 精神障碍与精神病对应。

\ContentBand{知识点原文摘取}

掌握：精神卫生与精神健康、精神障碍与精神病的基本概念\\
熟悉：精神病学的任务、与其他学科的关系，精神障碍病因学的有关基本概念，精神障碍分类的基本原则和目前通用的精神障碍分类系统的基本框架\\
了解：精神病学发展简史、国内外精神卫生发展；\\
拓展：精神病临床研究方法、目前趋势与一些重要问题\\
精神病学的定义和任务-熟悉/往年考试：是临床医学的一个分支，研究各种精神疾病的病因、发病机制、临床特点、疾病发展规律、治疗、预防、康复的一门科学\\
精神病学与其他学科的关系（心-身关系）-熟悉：\\
·心理因素是病因：心身疾病\\
·心理症状是结果：脑外伤导致精神病、患癌导致抑郁\\
·躯体症状是表现：躯体形式障碍、分离转换障碍\\
·临床医学必须渗透精神卫生理念：生物-心理-社会模式\\
·精神障碍-掌握：是一个临床诊断概念，一类具有诊断意义的精神方面的问题，要求症状须达到一定的严重程度和持续时间，伴有痛苦体验和/或功能损害，表现为认知、情绪、行为方面的改变（知情意）。\\
·精神病-掌握：是一类严重的精神障碍，只要出现持续的幻觉、妄想、严重精神运动兴奋或抑制表现者就是精神病。\\
例：脑或躯体疾病所致的谵妄状态；物质依赖所致的幻觉、妄想状态；精神分裂症和其他妄想性障碍；严重的情感障碍\\
·神经症-了解：包括神经衰弱、强迫症、焦虑症、恐怖症\\
精神障碍的病因-掌握：\\
1、与精神障碍致病因素有关的基本概念-掌握：\\
·素质因素：早年的、与发病时间远的有害因素，它形成易感素质。如遗传、原生家庭、幼年经历。\\
·促发因素：后来的、与发病时间近的有害因素，它与发病直接有关。如工作压力、生理因素\\
·附加因素：疾病发生之后附加于个体，使疾病加剧的事件。\\
2、常见的致病因素：年龄、性别、遗传、素质因素（原生家庭、幼年经历）、理化生因素（药物）、机体的功能状态（癌症、残疾等）\\
·双相：性别均衡、青春期发病\\
·抑郁：女性多，中青年发病\\
3、精神障碍发病的机制：生物学因素与社会心理因素综合作用

\InfoLabel{知识来源} 《精神病学复习整理by18级杨子铭》

\ContentBand{对应往年题}

\begin{pastquestion}

\QuestionLabel{【往年题1｜16级预防《精神病学》考题.pdf｜论述】}\\
请论述精神医学的概念和分类。

\end{pastquestion}

\begin{pastquestion}

\QuestionLabel{【往年题2｜06-07级 史前精神病考题.doc｜回忆/选择】}\\
神经症 症状缺乏特异性？心理因素是主要病因？

\end{pastquestion}

\begin{pastquestion}

\QuestionLabel{【往年题3｜06-07级 史前精神病考题.doc｜回忆/选择】}\\
不属于精神病范畴的是------颅脑外伤后的谵妄状态、严重的疑病症？

\end{pastquestion}

\begin{pastquestion}

\QuestionLabel{【往年题4｜06-07级 史前精神病考题.doc｜回忆/选择】}\\
在诊断精神病时，最优先考虑------器质性精神障碍？

\end{pastquestion}

\PointSeparator

\section{精神障碍的病因与发病机制}\label{users__fpb__26_spring__pkusph_study__tmp__pdfs__prevention-medicine-past-exams__06-ux7cbeux795eux75c5ux5b66ux5f80ux5e74ux9898ux8003ux70b9ux6574ux7406.md__ux7cbeux795eux969cux788dux7684ux75c5ux56e0ux4e0eux53d1ux75c5ux673aux5236}

\InfoLabel{考频} 1 次

\InfoLabel{对应小节} 一、绪论 / 精神障碍的病因

\InfoLabel{匹配依据} 题干涉及「病因」「促发因素」「附加因素」「多因素」「神经递质」等表述，与一、绪论 /
精神障碍的病因对应。

\ContentBand{知识点原文摘取}

精神障碍的病因-掌握：\\
1、与精神障碍致病因素有关的基本概念-掌握：\\
·素质因素：早年的、与发病时间远的有害因素，它形成易感素质。如遗传、原生家庭、幼年经历。\\
·促发因素：后来的、与发病时间近的有害因素，它与发病直接有关。如工作压力、生理因素\\
·附加因素：疾病发生之后附加于个体，使疾病加剧的事件。\\
2、常见的致病因素：年龄、性别、遗传、素质因素（原生家庭、幼年经历）、理化生因素（药物）、机体的功能状态（癌症、残疾等）\\
·双相：性别均衡、青春期发病\\
·抑郁：女性多，中青年发病\\
3、精神障碍发病的机制：生物学因素与社会心理因素综合作用\\
精神障碍的分类-掌握：\\
1、器质性与功能性疾病\\
2、精神病与神经症\\
3、多轴诊断：精神障碍、人格障碍、躯体疾病等多个轴\\
4、等级诊断：轻、中、重\\
5、共病：合并症\\
症状学分类基本原则-掌握：\\
·症状学诊断：根据主要症状命名的诊断\\
·属于状态性诊断，诊断名称只说明疾病当时所处的状态，如``躁狂状态''\\
·同一疾病诊断可因不同的症状群出现，分解成为几个互不相关的症状学诊断，如``双相情感障碍''可有``躁狂状态''、``抑郁状态''\\
·有利于对症治疗\\
精神病性障碍与非精神病性障碍-掌握/往年考试：\\
1、精神病性障碍：精神分裂症、情感性精神障碍（双相、抑郁）、偏执性精神障碍、器质性病变伴发精神障碍等\\
2、非精神病性障碍：神经症（神经衰弱、强迫症、焦虑症、恐怖症）、人格异常、心身疾病、病情较轻的精神发育迟滞、情绪反应等

\InfoLabel{知识来源} 《精神病学复习整理by18级杨子铭》

\ContentBand{对应往年题}

\begin{pastquestion}

\QuestionLabel{【往年题1｜06-07级 史前精神病考题.doc｜回忆/选择】}\\
关于精神病的病因说法正确的是------是一种多因素导致的疾病

\end{pastquestion}

\PointSeparator

\section{精神病学/精神医学的定义与任务}\label{users__fpb__26_spring__pkusph_study__tmp__pdfs__prevention-medicine-past-exams__06-ux7cbeux795eux75c5ux5b66ux5f80ux5e74ux9898ux8003ux70b9ux6574ux7406.md__ux7cbeux795eux75c5ux5b66ux7cbeux795eux533bux5b66ux7684ux5b9aux4e49ux4e0eux4efbux52a1}

\InfoLabel{考频} 2 次

\InfoLabel{对应小节} 一、绪论 / 精神病学的定义和任务

\InfoLabel{匹配依据}
题干涉及「精神病学」「精神医学」「psychiatry」「临床医学的分支」「4P因素」等表述，与一、绪论 /
精神病学的定义和任务对应。

\ContentBand{知识点原文摘取}

精神病学的定义和任务-熟悉/往年考试：是临床医学的一个分支，研究各种精神疾病的病因、发病机制、临床特点、疾病发展规律、治疗、预防、康复的一门科学\\
精神病学与其他学科的关系（心-身关系）-熟悉：\\
·心理因素是病因：心身疾病\\
·心理症状是结果：脑外伤导致精神病、患癌导致抑郁\\
·躯体症状是表现：躯体形式障碍、分离转换障碍\\
·临床医学必须渗透精神卫生理念：生物-心理-社会模式\\
·精神障碍-掌握：是一个临床诊断概念，一类具有诊断意义的精神方面的问题，要求症状须达到一定的严重程度和持续时间，伴有痛苦体验和/或功能损害，表现为认知、情绪、行为方面的改变（知情意）。\\
·精神病-掌握：是一类严重的精神障碍，只要出现持续的幻觉、妄想、严重精神运动兴奋或抑制表现者就是精神病。\\
例：脑或躯体疾病所致的谵妄状态；物质依赖所致的幻觉、妄想状态；精神分裂症和其他妄想性障碍；严重的情感障碍\\
·神经症-了解：包括神经衰弱、强迫症、焦虑症、恐怖症\\
精神障碍的病因-掌握：\\
1、与精神障碍致病因素有关的基本概念-掌握：\\
·素质因素：早年的、与发病时间远的有害因素，它形成易感素质。如遗传、原生家庭、幼年经历。\\
·促发因素：后来的、与发病时间近的有害因素，它与发病直接有关。如工作压力、生理因素\\
·附加因素：疾病发生之后附加于个体，使疾病加剧的事件。\\
2、常见的致病因素：年龄、性别、遗传、素质因素（原生家庭、幼年经历）、理化生因素（药物）、机体的功能状态（癌症、残疾等）\\
·双相：性别均衡、青春期发病\\
·抑郁：女性多，中青年发病\\
3、精神障碍发病的机制：生物学因素与社会心理因素综合作用\\
精神障碍的分类-掌握：\\
1、器质性与功能性疾病\\
2、精神病与神经症\\
3、多轴诊断：精神障碍、人格障碍、躯体疾病等多个轴\\
4、等级诊断：轻、中、重

\InfoLabel{知识来源} 《精神病学复习整理by18级杨子铭》

\ContentBand{对应往年题}

\begin{pastquestion}

\QuestionLabel{【往年题1｜20级完整试题回忆.docx｜论述】}\\
简述4P因素

\end{pastquestion}

\begin{pastquestion}

\QuestionLabel{【往年题2｜21级精神病学考题yys.docx / 21级本科预防二班+胡润泽+精神病+主观题总结.docx｜简答】}\\
简述并举例说明4P因素

\end{pastquestion}

\PointSeparator

\chapter{精神症状学}\label{users__fpb__26_spring__pkusph_study__tmp__pdfs__prevention-medicine-past-exams__06-ux7cbeux795eux75c5ux5b66ux5f80ux5e74ux9898ux8003ux70b9ux6574ux7406.md__ux7b2cux4e8cux7ae0-ux7cbeux795eux75c7ux72b6ux5b66}

\section{注意、记忆与智能障碍}\label{users__fpb__26_spring__pkusph_study__tmp__pdfs__prevention-medicine-past-exams__06-ux7cbeux795eux75c5ux5b66ux5f80ux5e74ux9898ux8003ux70b9ux6574ux7406.md__ux6ce8ux610fux8bb0ux5fc6ux4e0eux667aux80fdux969cux788d}

\InfoLabel{考频} 6 次

\InfoLabel{对应小节} 二、精神症状学 / 注意记忆智能

\InfoLabel{匹配依据} 题干涉及「注意」「记忆」「顺行性遗忘」「逆行性遗忘」「错构」等表述，与二、精神症状学 /
注意记忆智能对应。

\ContentBand{知识点原文摘取}

二、精神症状学\\
掌握：精神症状的共同特点和基本要素，重点掌握幻觉、妄想、抑郁、焦虑、谵妄状态等症状和综合症的含义\\
熟悉：其他临床常见精神症状的含义（定义和临床意义）、表现与实例，熟悉精神症状的类别\\
了解：精神症状学的研究方法\\
精神症状的共同特点-掌握：\\
·症状不受意识控制，如幻觉症状\\
·与客观环境不相称，如妄想症状\\
·影响社会功能\\
·多数伴有痛苦体验，如强迫症状\\
精神症状的基本要素-掌握：\\
·性质：幻觉、妄想、抑郁、焦虑、谵妄\\
·频度与强度\\
·持续时间\\
·对患者的影响\\
精神疾病的常见症状-掌握：\\
人类的正常精神活动，可以按心理学概念分为知 （认知）、情（情感）、意（意志）：\\
·知（认知）：感觉、知觉、思维、记忆、注意、智力、自知力、定向力等\\
·情（情感）：内心体验、外在表现\\
·意（意志行为）：低级意志、高级意志\\
1、认知和认知障碍\\
(1)感觉障碍：少见，包括感觉过敏、感觉减退、感觉倒错（对外界刺激做出相反的反应）、内感性不适（难以描述和定位的不适感）\\
(2)知觉障碍-核心重点：最常见\\
①错觉：有东西，你看错了\\
②幻觉-重点：没东西，你看到东西了（无中生有）。\\
·按器官分类：幻听、幻视、幻嗅、幻味、幻触、内脏性幻觉（感觉内脏扭转、穿孔）\\
·两种特殊形式：思维鸣响（想到什么就听到什么，想喝水就听到``喝水''）、功能性幻听（幻听和现实刺激同时出现和消失，打开水龙头就听见``唯物主义''）\\
·按幻觉性质分类：真性幻觉（幻象与真实事物完全相同，亲眼看到、亲耳听
到）、假性幻觉（不清晰、难定位，不用眼睛就能看到头脑里的幻象）\\
③感知综合障碍：\\
·视物变形症（感知物体大小错误）\\
·空间知觉障碍（感知距离错误）\\
·非真实感（周围事物模糊不真实）\\
·对自身躯体结构的感知综合障碍（感觉自己身体大小错误）

\InfoLabel{知识来源} 《精神病学复习整理by18级杨子铭》

\ContentBand{对应往年题}

\begin{pastquestion}

\QuestionLabel{【往年题1｜20级完整试题回忆.docx｜单项选择】}\\
多动障碍的临床表现不包括 A、品行障碍 B、智力低下 C、学习成绩差 D、注意力障碍 E、\ldots\ldots{}

\end{pastquestion}

\begin{pastquestion}

\QuestionLabel{【往年题2｜20级完整试题回忆.docx｜单项选择】}\\
阿尔兹海默病的症状 A、远期记忆首先受累 B、不出现虚构或错构 C、对记忆力下降有自知力 D、情景记忆首先受累
E、\ldots\ldots{}

\end{pastquestion}

\begin{pastquestion}

\QuestionLabel{【往年题3｜06-07级 史前精神病考题.doc｜回忆/选择】}\\
关于精神发育迟滞不正确的是------都有躯体畸形

\end{pastquestion}

\begin{pastquestion}

\QuestionLabel{【往年题4｜06-07级 史前精神病考题.doc｜回忆/选择】}\\
儿童多动症的核心------注意障碍（与活动增多区分）

\end{pastquestion}

\begin{pastquestion}

\QuestionLabel{【往年题5｜06-07级 史前精神病考题.doc｜回忆/选择】}\\
关于儿童孤独症不正确的是------痴呆面容

\end{pastquestion}

\begin{pastquestion}

\QuestionLabel{【往年题6｜06-07级 史前精神病考题.doc｜回忆/选择】}\\
最常见于脑器质性精神障碍初期的是------顺行性遗忘、近事遗忘、远事遗忘、虚构、错构？

\end{pastquestion}

\PointSeparator

\section{思维障碍（形式与内容）}\label{users__fpb__26_spring__pkusph_study__tmp__pdfs__prevention-medicine-past-exams__06-ux7cbeux795eux75c5ux5b66ux5f80ux5e74ux9898ux8003ux70b9ux6574ux7406.md__ux601dux7ef4ux969cux788dux5f62ux5f0fux4e0eux5185ux5bb9}

\InfoLabel{考频} 11 次

\InfoLabel{对应小节} 二、精神症状学 / 思维障碍

\InfoLabel{匹配依据}
题干涉及「思维奔逸」「思维迟缓」「思维贫乏」「思维破裂」「思维松弛」等表述，与二、精神症状学 / 思维障碍对应。

\ContentBand{知识点原文摘取}

(3)思维障碍\\
①思维形式障碍\\
·联想过程障碍（连贯性不好）：思维奔逸、思维迟缓（慢-抑郁）、思维贫乏（少-精分）、病理性赘述（重复不重要的）、思维不连贯、思维中断、思维松弛（无逻辑）、思维破裂、强制性思维（外来的思维，不痛苦，很快消失）\\
·逻辑障碍：病理性象征性思维（反复穿脱衣服表示``表里如一''）、语词新作（``\%''意思是离婚）\\
②思维内容障碍\\
·妄想-核心重点：\\
根据产生的条件：原发性妄想（诊断意义）、继发性妄想（继发于幻觉、心境障碍、意识障碍、痴呆状态、
酒精戒断等精神症状）\\
根据涉及的内容：\\
·关系妄想：总觉得别人说的话和自己有关系，在针对自己\\
·被害妄想：觉得有人给自己下毒\\
·夸大妄想：认为自己能力非常强\\
·嫉妒妄想：怀疑配偶出轨\\
·钟情妄想：总认为别人喜欢自己\\
·疑病妄想：认为自己有躯体疾病\\
·被洞悉感：觉得自己被看穿了\\
·影响妄想/物理影响妄想：总觉得自己被外界力量所控制（红外线）\\
·非血统妄想：觉得自己不是亲生的，而是伟人的后代\\
·超价观念：一种占主导地位的强烈情绪，坚持这种看法，与人格整合良好，有事实基础（区别于``妄想''），很超前，但逻辑上并不荒谬，可以被理解\\
·强迫观念：某一想法反复出现，是自己的观念，不是外界强加的。伴有主观的痛苦感。如强迫性回忆、强迫性怀疑、强迫性穷思竭虑。\\
(4)注意障碍\\
·主动注意（有目的 需要努力 如上课听讲）、被动注意（不需要努力就可以自然地注意，如窗外鸟叫）\\
·注意增强、注意减弱、注意缓慢、注意涣散（主动注意减弱，注意力不集中）、注意狭窄、注意转移（被动注意增加，注意力容易被吸引）\\
(5)记忆障碍\\
·顺行性遗忘：忘了最近的事，如脑震荡\\
·逆行性遗忘：忘了以前的事，如严重精神创伤。正常是逆行性。

\InfoLabel{知识来源} 《精神病学复习整理by18级杨子铭》

\ContentBand{对应往年题}

\begin{pastquestion}

\QuestionLabel{【往年题1｜20级完整试题回忆.docx｜单项选择】}\\
阿尔兹海默病容易出现的妄想类型 A、被害妄想 B、被窃妄想 C、疑病妄想

\end{pastquestion}

\begin{pastquestion}

\QuestionLabel{【往年题2｜06-07级 史前精神病考题.doc｜回忆/选择】}\\
诡异的有，思维松弛（为什么入院---母亲生病了三天，后来我也生病了三天，没有饿，就进来了）、病理性激情（书上没有，考了两次）

\end{pastquestion}

\begin{pastquestion}

\QuestionLabel{【往年题3｜06-07级 史前精神病考题.doc｜回忆/选择】}\\
2）最可能是妄想的例子（别人能够看到自己的内心、声称国家某领导人是自己的校友、感到自己的脑内有很多不属于自己的思想、在公共食堂吃饭看到别人敲饭碗，认为是在骂自己是``饭桶''）

\end{pastquestion}

\begin{pastquestion}

\QuestionLabel{【往年题4｜06-07级 史前精神病考题.doc｜回忆/选择】}\\
4）病人入院后常在窗前聆听，称有人议论她，并听到公安机关说要来抓他，患者冲窗外大喊，与之辩论，称不要诬陷好人，诊断（关系妄想、被害妄想、假性幻听、言语性幻听、命令性幻听？）------原题回忆得可能不太准确，重点区分被害妄想和幻听

\end{pastquestion}

\begin{pastquestion}

\QuestionLabel{【往年题5｜06-07级 史前精神病考题.doc｜回忆/选择】}\\
8）听见广播里一个错字，认为播音员在诱惑自己；之后听到广播认为所有播音员都在诱惑他------原发性妄想

\end{pastquestion}

\begin{pastquestion}

\QuestionLabel{【往年题6｜06-07级 史前精神病考题.doc｜回忆/选择】}\\
9）患者近2个月来常用头部撞击轮胎，解释说是为了``投胎''什么的------病理性象征思维、特殊意义妄想？？

\end{pastquestion}

\begin{pastquestion}

\QuestionLabel{【往年题7｜06-07级 史前精神病考题.doc｜回忆/选择】}\\
10）感到被人跟踪、家中安装有窃听器------被害妄想

\end{pastquestion}

\begin{pastquestion}

\QuestionLabel{【往年题8｜06-07级 史前精神病考题.doc｜回忆/选择】}\\
11）患者，20岁，称初中时偷过同学的铅笔，认为罪大恶极，应该施以惩罚------自罪妄想

\end{pastquestion}

\begin{pastquestion}

\QuestionLabel{【往年题9｜06-07级 史前精神病考题.doc｜回忆/选择】}\\
争论性幻听、思维松弛等等

\end{pastquestion}

\begin{pastquestion}

\QuestionLabel{【往年题10｜06-07级 史前精神病考题.doc｜回忆/选择】}\\
关于钟情妄想正确的是------病人坚信别人对自己钟情

\end{pastquestion}

\begin{pastquestion}

\QuestionLabel{【往年题11｜06-07级 史前精神病考题.doc｜回忆/选择】}\\
不属于思维内容障碍的是------思维贫乏

\end{pastquestion}

\PointSeparator

\section{情感与意志行为障碍}\label{users__fpb__26_spring__pkusph_study__tmp__pdfs__prevention-medicine-past-exams__06-ux7cbeux795eux75c5ux5b66ux5f80ux5e74ux9898ux8003ux70b9ux6574ux7406.md__ux60c5ux611fux4e0eux610fux5fd7ux884cux4e3aux969cux788d}

\InfoLabel{考频} 3 次

\InfoLabel{对应小节} 二、精神症状学 / 情感意志行为

\InfoLabel{匹配依据} 题干涉及「情感高涨」「情感低落」「焦虑」「恐惧」「情感淡漠」等表述，与二、精神症状学 /
情感意志行为对应。

\ContentBand{知识点原文摘取}

二、精神症状学\\
掌握：精神症状的共同特点和基本要素，重点掌握幻觉、妄想、抑郁、焦虑、谵妄状态等症状和综合症的含义\\
熟悉：其他临床常见精神症状的含义（定义和临床意义）、表现与实例，熟悉精神症状的类别\\
了解：精神症状学的研究方法\\
精神症状的共同特点-掌握：\\
·症状不受意识控制，如幻觉症状\\
·与客观环境不相称，如妄想症状\\
·影响社会功能\\
·多数伴有痛苦体验，如强迫症状\\
精神症状的基本要素-掌握：\\
·性质：幻觉、妄想、抑郁、焦虑、谵妄\\
·频度与强度\\
·持续时间\\
·对患者的影响\\
精神疾病的常见症状-掌握：\\
人类的正常精神活动，可以按心理学概念分为知 （认知）、情（情感）、意（意志）：\\
·知（认知）：感觉、知觉、思维、记忆、注意、智力、自知力、定向力等\\
·情（情感）：内心体验、外在表现\\
·意（意志行为）：低级意志、高级意志\\
1、认知和认知障碍\\
(1)感觉障碍：少见，包括感觉过敏、感觉减退、感觉倒错（对外界刺激做出相反的反应）、内感性不适（难以描述和定位的不适感）\\
(2)知觉障碍-核心重点：最常见\\
①错觉：有东西，你看错了\\
②幻觉-重点：没东西，你看到东西了（无中生有）。\\
·按器官分类：幻听、幻视、幻嗅、幻味、幻触、内脏性幻觉（感觉内脏扭转、穿孔）\\
·两种特殊形式：思维鸣响（想到什么就听到什么，想喝水就听到``喝水''）、功能性幻听（幻听和现实刺激同时出现和消失，打开水龙头就听见``唯物主义''）\\
·按幻觉性质分类：真性幻觉（幻象与真实事物完全相同，亲眼看到、亲耳听
到）、假性幻觉（不清晰、难定位，不用眼睛就能看到头脑里的幻象）\\
③感知综合障碍：\\
·视物变形症（感知物体大小错误）\\
·空间知觉障碍（感知距离错误）\\
·非真实感（周围事物模糊不真实）\\
·对自身躯体结构的感知综合障碍（感觉自己身体大小错误）

\InfoLabel{知识来源} 《精神病学复习整理by18级杨子铭》

\ContentBand{对应往年题}

\begin{pastquestion}

\QuestionLabel{【往年题1｜06-07级 史前精神病考题.doc｜回忆/选择】}\\
甲减伴精神症状比较少见的是？幻觉妄想or 精神运动性不安

\end{pastquestion}

\begin{pastquestion}

\QuestionLabel{【往年题2｜06-07级 史前精神病考题.doc｜回忆/选择】}\\
恐怖症与焦虑症的主要区别------有无特定场合

\end{pastquestion}

\begin{pastquestion}

\QuestionLabel{【往年题3｜06-07级 史前精神病考题.doc｜回忆/选择】}\\
不属于抑郁的生物学症状是------情感低落

\end{pastquestion}

\PointSeparator

\section{感知觉障碍（错觉、幻觉、感知综合障碍）}\label{users__fpb__26_spring__pkusph_study__tmp__pdfs__prevention-medicine-past-exams__06-ux7cbeux795eux75c5ux5b66ux5f80ux5e74ux9898ux8003ux70b9ux6574ux7406.md__ux611fux77e5ux89c9ux969cux788dux9519ux89c9ux5e7bux89c9ux611fux77e5ux7efcux5408ux969cux788d}

\InfoLabel{考频} 8 次

\InfoLabel{对应小节} 二、精神症状学 / 知觉障碍

\InfoLabel{匹配依据} 题干涉及「幻觉」「错觉」「幻听」「幻视」「真性幻觉」等表述，与二、精神症状学 /
知觉障碍对应。

\ContentBand{知识点原文摘取}

(2)知觉障碍-核心重点：最常见\\
①错觉：有东西，你看错了\\
②幻觉-重点：没东西，你看到东西了（无中生有）。\\
·按器官分类：幻听、幻视、幻嗅、幻味、幻触、内脏性幻觉（感觉内脏扭转、穿孔）\\
·两种特殊形式：思维鸣响（想到什么就听到什么，想喝水就听到``喝水''）、功能性幻听（幻听和现实刺激同时出现和消失，打开水龙头就听见``唯物主义''）\\
·按幻觉性质分类：真性幻觉（幻象与真实事物完全相同，亲眼看到、亲耳听
到）、假性幻觉（不清晰、难定位，不用眼睛就能看到头脑里的幻象）\\
③感知综合障碍：\\
·视物变形症（感知物体大小错误）\\
·空间知觉障碍（感知距离错误）\\
·非真实感（周围事物模糊不真实）\\
·对自身躯体结构的感知综合障碍（感觉自己身体大小错误）\\
(3)思维障碍\\
①思维形式障碍\\
·联想过程障碍（连贯性不好）：思维奔逸、思维迟缓（慢-抑郁）、思维贫乏（少-精分）、病理性赘述（重复不重要的）、思维不连贯、思维中断、思维松弛（无逻辑）、思维破裂、强制性思维（外来的思维，不痛苦，很快消失）\\
·逻辑障碍：病理性象征性思维（反复穿脱衣服表示``表里如一''）、语词新作（``\%''意思是离婚）\\
②思维内容障碍\\
·妄想-核心重点：\\
根据产生的条件：原发性妄想（诊断意义）、继发性妄想（继发于幻觉、心境障碍、意识障碍、痴呆状态、
酒精戒断等精神症状）\\
根据涉及的内容：\\
·关系妄想：总觉得别人说的话和自己有关系，在针对自己\\
·被害妄想：觉得有人给自己下毒\\
·夸大妄想：认为自己能力非常强\\
·嫉妒妄想：怀疑配偶出轨\\
·钟情妄想：总认为别人喜欢自己\\
·疑病妄想：认为自己有躯体疾病\\
·被洞悉感：觉得自己被看穿了\\
·影响妄想/物理影响妄想：总觉得自己被外界力量所控制（红外线）

\InfoLabel{知识来源} 《精神病学复习整理by18级杨子铭》

\ContentBand{对应往年题}

\begin{pastquestion}

\QuestionLabel{【往年题1｜06-07级 史前精神病考题.doc｜回忆/选择】}\\
看书时听见陌生的声音把自己正在看的内容读出来是反射性幻听还是思维鸣响？

\end{pastquestion}

\begin{pastquestion}

\QuestionLabel{【往年题2｜06-07级 史前精神病考题.doc｜回忆/选择】}\\
5）高烧患者凭空看见毒蛇、猛兽似乎要向他扑来------幻觉、错觉、虚构\ldots\ldots{}

\end{pastquestion}

\begin{pastquestion}

\QuestionLabel{【往年题3｜06-07级 史前精神病考题.doc｜回忆/选择】}\\
7）男性学生，看见房屋倾斜、电线杆变细------感知综合障碍（没有视物变形选项，坑爹）

\end{pastquestion}

\begin{pastquestion}

\QuestionLabel{【往年题4｜06-07级 史前精神病考题.doc｜回忆/选择】}\\
幻觉------虚幻的知觉体验

\end{pastquestion}

\begin{pastquestion}

\QuestionLabel{【往年题5｜06-07级 史前精神病考题.doc｜回忆/选择】}\\
酒中毒性幻觉症常见------言语性幻听、持续性幻视\ldots\ldots{}

\end{pastquestion}

\begin{pastquestion}

\QuestionLabel{【往年题6｜06-07级 史前精神病考题.doc｜回忆/选择】}\\
关于谵妄的特点正确的是------幻觉多为恐怖性幻视

\end{pastquestion}

\begin{pastquestion}

\QuestionLabel{【往年题7｜06-07级 史前精神病考题.doc｜回忆/选择】}\\
属于精神分裂症的特征性症状是------幻觉、妄想、思维鸣响\ldots\ldots{}

\end{pastquestion}

\begin{pastquestion}

\QuestionLabel{【往年题8｜06-07级 史前精神病考题.doc｜回忆/选择】}\\
焦虑症中不可能出现的是------幻听

\end{pastquestion}

\PointSeparator

\chapter{精神分裂症}\label{users__fpb__26_spring__pkusph_study__tmp__pdfs__prevention-medicine-past-exams__06-ux7cbeux795eux75c5ux5b66ux5f80ux5e74ux9898ux8003ux70b9ux6574ux7406.md__ux7b2cux4e09ux7ae0-ux7cbeux795eux5206ux88c2ux75c7}

\section{精神分裂症概述与临床表现（阳性/阴性/认知）}\label{users__fpb__26_spring__pkusph_study__tmp__pdfs__prevention-medicine-past-exams__06-ux7cbeux795eux75c5ux5b66ux5f80ux5e74ux9898ux8003ux70b9ux6574ux7406.md__ux7cbeux795eux5206ux88c2ux75c7ux6982ux8ff0ux4e0eux4e34ux5e8aux8868ux73b0ux9633ux6027ux9634ux6027ux8ba4ux77e5}

\InfoLabel{考频} 6 次

\InfoLabel{对应小节} 三、精神分裂症 / 临床表现

\InfoLabel{匹配依据}
题干涉及「精神分裂症」「阳性症状」「阴性症状」「认知症状」「瓦解症状」等表述，与三、精神分裂症 /
临床表现对应。

\ContentBand{知识点原文摘取}

精神分裂症概述：主要症状为思维、情感、行为等多方面的障碍（知情意分离），精神活动与环境不协调。通常意识清晰，多起病于青壮年，常缓慢起病，病程迁延，有慢性化倾向和衰退的可能。

精神分裂症的临床表现：\\
1、阳性症状（本质是DA亢进）：幻觉（言语性幻听）、妄想（关系、被害、非血统的原发性妄想）、被动体验（被洞悉感、物理影响妄想、强制性思维）、不协调的精神运动性兴奋（思维紊乱：思维破裂、语词新作、病理象征性思维、逻辑倒错；情感紊乱：变化莫测、情感倒错；行为紊乱：幼稚、愚蠢、装相）、紧张症行为（紧张性木僵和紧张性兴奋交替出现）。\\
2、阴性症状（本质是5-HT和谷氨酸缺乏）：思维贫乏、情感淡漠、意志缺乏、快感缺乏、社交退缩。\\
3、认知症状：注意障碍（注意涣散、注意转移困难、选择性注意障碍）、工作记忆障碍（语言和视觉记忆缺陷）、执行功能障碍（信息整合困难、应对变化能力不足）。

病因和发病机制：遗传因素、神经生化异常（多巴胺DA功能亢进、谷氨酸功能不足、5-HT功能不足）、社会心理因素（精神刺激只是诱因）、易感素质（神经发育障碍）。

\InfoLabel{知识来源} 《精神病学复习整理by18级杨子铭》

\ContentBand{对应往年题}

\begin{pastquestion}

\QuestionLabel{【往年题1｜16级预防《精神病学》考题.pdf｜简答】}\\
精神分裂症常见的阳性症状和阴性症状有哪些？

\end{pastquestion}

\begin{pastquestion}

\QuestionLabel{【往年题2｜20级完整试题回忆.docx｜单项选择】}\\
有关精神分裂症的定义哪项不正确？ A、一组病因未明的疾病 B、一般无意识障碍 C、多起病于青少年
D、常有感知、思维、情感、行为等多方面的障碍 E、精神活动的不协调

\end{pastquestion}

\begin{pastquestion}

\QuestionLabel{【往年题3｜20级完整试题回忆.docx｜单项选择】}\\
精神分裂症的五维症状不包括 A、认知损害 B、离解症状 C、焦虑抑郁 D、攻击敌意 E、\ldots\ldots{}

\end{pastquestion}

\begin{pastquestion}

\QuestionLabel{【往年题4｜20级完整试题回忆.docx｜简答】}\\
精神分裂症的临床表现

\end{pastquestion}

\begin{pastquestion}

\QuestionLabel{【往年题5｜06-07级 史前精神病考题.doc｜回忆/选择】}\\
蜡样屈曲最常见于------紧张性精神分裂症

\end{pastquestion}

\begin{pastquestion}

\QuestionLabel{【往年题6｜06-07级 史前精神病考题.doc｜回忆/选择】}\\
临床最常见的精分类型------偏执型

\end{pastquestion}

\PointSeparator

\section{抗精神病药物与治疗原则}\label{users__fpb__26_spring__pkusph_study__tmp__pdfs__prevention-medicine-past-exams__06-ux7cbeux795eux75c5ux5b66ux5f80ux5e74ux9898ux8003ux70b9ux6574ux7406.md__ux6297ux7cbeux795eux75c5ux836fux7269ux4e0eux6cbbux7597ux539fux5219}

\InfoLabel{考频} 5 次

\InfoLabel{对应小节} 三、精神分裂症 / 治疗

\InfoLabel{匹配依据} 题干涉及「抗精神病」「用药原则」「联合用药」「足量」「足疗程」等表述，与三、精神分裂症 /
治疗对应。

\ContentBand{知识点原文摘取}

精神分裂症的药物治疗原则：\\
·一旦确诊，立即开始药物治疗\\
·以单一用药为宜\\
·用药个体化\\
·小剂量起始，根据病情和治疗场所掌握药物滴定速度\\
·足量足疗程治疗\\
·定期认真评定疗效和不良反应，积极调整治疗方案\\
·依从性、知情同意、缓慢减药、防止复发

药物分类：\\
·传统/典型抗精神病药物（D2受体阻滞剂）：氟哌啶醇、氯丙嗪\\
·非典型抗精神病药（平衡阻滞5-HT和D2受体）：利培酮、奥氮平、喹硫平\\
·药物主要副作用：嗜睡、锥体外系反应（急性肌张力障碍、震颤麻痹、静坐不能、迟发性运动障碍）

抗精神病药的治疗作用：缓解阳性症状、改善阴性症状、改善认知功能、长期治疗防止复发。

\InfoLabel{知识来源} 《精神病学复习整理by18级杨子铭》

\ContentBand{对应往年题}

\begin{pastquestion}

\QuestionLabel{【往年题1｜20级完整试题回忆.docx｜单项选择】}\\
精神分裂症的用药原则错误的是 A、首选联合用药 B、足量 C、足疗程 D、个体化 E、\ldots\ldots{}

\end{pastquestion}

\begin{pastquestion}

\QuestionLabel{【往年题2｜20级完整试题回忆.docx｜单项选择】}\\
抑郁症的药物治疗原则 A、个体化 B、单一用药 C、足量 D、足疗程 E、以上都是

\end{pastquestion}

\begin{pastquestion}

\QuestionLabel{【往年题3｜06-07级 史前精神病考题.doc｜回忆/选择】}\\
氯氮平最重要的不良反应------粒细胞缺乏

\end{pastquestion}

\begin{pastquestion}

\QuestionLabel{【往年题4｜06-07级 史前精神病考题.doc｜回忆/选择】}\\
不属于锥体外系不良反应的是------恶性综合症

\end{pastquestion}

\begin{pastquestion}

\QuestionLabel{【往年题5｜06-07级 史前精神病考题.doc｜回忆/选择】}\\
精分患者急性期，明显兴奋躁动，最适合的用药------抗精神病药物、强效镇静催眠药物？

\end{pastquestion}

\PointSeparator

\section{精神分裂症诊断标准与鉴别}\label{users__fpb__26_spring__pkusph_study__tmp__pdfs__prevention-medicine-past-exams__06-ux7cbeux795eux75c5ux5b66ux5f80ux5e74ux9898ux8003ux70b9ux6574ux7406.md__ux7cbeux795eux5206ux88c2ux75c7ux8bcaux65adux6807ux51c6ux4e0eux9274ux522b}

\InfoLabel{考频} 5 次

\InfoLabel{对应小节} 三、精神分裂症 / 诊断与鉴别

\InfoLabel{匹配依据} 题干涉及「ICD-10」「病程」「一个月」「症状学标准」「鉴别」等表述，与三、精神分裂症 /
诊断与鉴别对应。

\ContentBand{知识点原文摘取}

精神分裂症诊断标准（ICD-10）：在至少1个月以上的大部分时间内肯定存在下述1～4中的任何1组（如不甚明确常需要两个或多个症状）或5～9至少2组症状群中的十分明确的症状：\\
（1）思维鸣响、思维插入、思维被撤走、思维广播\\
（2）被影响、被控制或被动妄想；妄想性知觉\\
（3）评论性、或来源于身体某一部分的幻听\\
（4）与文化不相称且根本不可能的其他类型的持续性妄想\\
（5）伴转瞬即逝或未充分形成的无明显情感内容的妄想，或伴有持久的超价观念\\
（6）思潮断裂或无关的插入语，导致言语不连贯\\
（7）紧张性行为\\
（8）阴性症状\\
（9）个人行为的某些方面发生显著而持久的总体性质的改变

精神分裂症的鉴别诊断：\\
·神经症：有自知力，有求治欲\\
·躁狂症：发作性病程，情感有协调性和感染力\\
·抑郁症：情感低落而不是淡漠\\
·人格障碍：无明确病程，无精神病性症状\\
·心因性精神障碍：症状与心因的关系\\
·躯体疾病所致精神障碍：与原发病的关系\\
·脑器质性精神障碍：急性期意识障碍，慢性期智能障碍

\InfoLabel{知识来源} 《精神病学复习整理by18级杨子铭》

\ContentBand{对应往年题}

\begin{pastquestion}

\QuestionLabel{【往年题1｜20级完整试题回忆.docx｜单项选择】}\\
ICD-10对于精神分裂症的病程诊断标准为 A、两周 B、一个月 C、三个月 D、半年 E、一年

\end{pastquestion}

\begin{pastquestion}

\QuestionLabel{【往年题2｜20级完整试题回忆.docx｜单项选择】}\\
精神分裂症的症状学标准不包括 A、思维鸣响 B、被控制妄想 C、妄想心境 D、紧张性行为 E、评论性幻听

\end{pastquestion}

\begin{pastquestion}

\QuestionLabel{【往年题3｜06-07级 史前精神病考题.doc｜回忆/选择】}\\
某病的病程特点：精神分裂不适发作性的，心境障碍不是迁延的，AD是进行性进展的

\end{pastquestion}

\begin{pastquestion}

\QuestionLabel{【往年题4｜06-07级 史前精神病考题.doc｜回忆/选择】}\\
关于心境障碍不对的说法是------病程多迁延、发作性病程、间歇期可完全恢复正常？

\end{pastquestion}

\begin{pastquestion}

\QuestionLabel{【往年题5｜06-07级 史前精神病考题.doc｜回忆/选择】}\\
躁狂的病程标准------1周

\end{pastquestion}

\PointSeparator

\chapter{双相情感障碍}\label{users__fpb__26_spring__pkusph_study__tmp__pdfs__prevention-medicine-past-exams__06-ux7cbeux795eux75c5ux5b66ux5f80ux5e74ux9898ux8003ux70b9ux6574ux7406.md__ux7b2cux56dbux7ae0-ux53ccux76f8ux60c5ux611fux969cux788d}

\section{躁狂发作与轻躁狂发作}\label{users__fpb__26_spring__pkusph_study__tmp__pdfs__prevention-medicine-past-exams__06-ux7cbeux795eux75c5ux5b66ux5f80ux5e74ux9898ux8003ux70b9ux6574ux7406.md__ux8e81ux72c2ux53d1ux4f5cux4e0eux8f7bux8e81ux72c2ux53d1ux4f5c}

\InfoLabel{考频} 3 次

\InfoLabel{对应小节} 四、双相 / 躁狂发作

\InfoLabel{匹配依据} 题干涉及「躁狂发作」「躁狂」「轻躁狂」「三高」「情绪高涨」等表述，与四、双相 /
躁狂发作对应。

\ContentBand{知识点原文摘取}

双相的临床发作类型：\\
（1）躁狂发作：持续1周或以上\\
·核心表现：情绪高涨、易激惹（情）；思维加快、夸大（知）；活动增多、精力充沛（意）\\
·伴随症状：睡眠减少、注意力不集中、鲁莽冲动、心境不稳\\
·其他：可以伴有精神病性症状，社会功能严重受损\\
（2）轻躁狂发作：肯定不伴幻觉和妄想（区别于躁狂），社会功能受损程度不严重

\InfoLabel{知识来源} 《精神病学复习整理by18级杨子铭》

\ContentBand{对应往年题}

\begin{pastquestion}

\QuestionLabel{【往年题1｜20级完整试题回忆.docx｜单项选择】}\\
\ldots\ldots 医生考虑为轻躁狂发作而非躁狂发作，本病例中重要的鉴别点在于 A、否认易激惹 B、没有明显的情感高涨
C、持续时间超过1周 D、没有明显的社交功能损害 E、没有精神病性症状

\end{pastquestion}

\begin{pastquestion}

\QuestionLabel{【往年题2｜20级完整试题回忆.docx｜单项选择】}\\
躁狂发作的核心症状不包括

\end{pastquestion}

\begin{pastquestion}

\QuestionLabel{【往年题3｜20级完整试题回忆.docx｜单项选择】}\\
躁狂发作的病程诊断标准

\end{pastquestion}

\PointSeparator

\section{双相诊断、鉴别与治疗原则}\label{users__fpb__26_spring__pkusph_study__tmp__pdfs__prevention-medicine-past-exams__06-ux7cbeux795eux75c5ux5b66ux5f80ux5e74ux9898ux8003ux70b9ux6574ux7406.md__ux53ccux76f8ux8bcaux65adux9274ux522bux4e0eux6cbbux7597ux539fux5219}

\InfoLabel{考频} 3 次

\InfoLabel{对应小节} 四、双相 / 治疗原则

\InfoLabel{匹配依据} 题干涉及「双相障碍治疗」「治疗原则」「全病程」「优先原则」「共病」等表述，与四、双相 /
治疗原则对应。

\ContentBand{知识点原文摘取}

双相的诊断原则：``症状---综合征---精神障碍''\\
·症状学标准：核心症状、伴随症状、其他\\
·严重程度标准：轻、中、重（伴有精神病性症状）\\
·病程标准：持续时间\\
·排除标准：脑器质性、躯体疾病、药物所致

双相的治疗原则：综合、全病程、全面、优先、共病、患者参与、依从、评估和监测\\
·综合治疗：药物治疗为主，心理行为治疗辅助\\
·全病程治疗：急性期、巩固期、维持期\\
·治疗药物：心境稳定剂（锂盐、丙戊酸、卡马西平）是药物治疗的核心，第二代抗精神病药SGAs辅助

\InfoLabel{知识来源} 《精神病学复习整理by18级杨子铭》

\ContentBand{对应往年题}

\begin{pastquestion}

\QuestionLabel{【往年题1｜20级完整试题回忆.docx｜单项选择】}\\
双相障碍治疗原则 A、充分评估、量化监测原则 B、提高治疗依从性原则 C、患方共同参与治疗原则 D、全病程治疗原则
E、以上都是

\end{pastquestion}

\begin{pastquestion}

\QuestionLabel{【往年题2｜20级完整试题回忆.docx｜病例分析】}\\
双相情感障碍病例：精神症状、诊断与诊断思路、治疗原则（情绪高低交替、家族史阳性、既往躁狂史）

\end{pastquestion}

\begin{pastquestion}

\QuestionLabel{【往年题3｜06-07级 史前精神病考题.doc｜回忆/选择】}\\
不属于电休克治疗的适应症------ 一种药物治疗无效

\end{pastquestion}

\PointSeparator

\section{双相障碍概述与临床特点}\label{users__fpb__26_spring__pkusph_study__tmp__pdfs__prevention-medicine-past-exams__06-ux7cbeux795eux75c5ux5b66ux5f80ux5e74ux9898ux8003ux70b9ux6574ux7406.md__ux53ccux76f8ux969cux788dux6982ux8ff0ux4e0eux4e34ux5e8aux7279ux70b9}

\InfoLabel{考频} 1 次

\InfoLabel{对应小节} 四、双相情感障碍 / 概述

\InfoLabel{匹配依据} 题干涉及「双相」「心境稳定」「发作性」「I型」「II型」等表述，与四、双相情感障碍 /
概述对应。

\ContentBand{知识点原文摘取}

双相及相关障碍：是发作性的心境障碍，由躁狂发作、轻躁狂发作、混合发作以及相关症状所定义。\\
典型特征：双相常在青春期起病，性别均衡。心境稳定剂治疗是药物干预的核心。\\
双相的临床特点：\\
（1）复杂多样：躁狂和抑郁的发作反复交替、不规则呈现\\
（2）可伴有焦虑、强迫、物质滥用，也可出现幻觉、妄想\\
（3）躁狂发作平均3个月左右；抑郁发作平均9个月\\
（4）间歇期社会功能相对恢复正常\\
双相的分型：双相I型（躁狂或混合发作为主）、双相Ⅱ型（抑郁为主，伴轻躁狂）

\InfoLabel{知识来源} 《精神病学复习整理by18级杨子铭》

\ContentBand{对应往年题}

\begin{pastquestion}

\QuestionLabel{【往年题1｜21级精神病学考题yys.docx /
21级本科预防二班+胡润泽+精神病+主观题总结.docx｜病例分析】}\\
双相I型病例：目前精神症状及举例、诊断及鉴别诊断与证据、治疗原则（22岁男性，2016抑郁-2022躁狂-2024再躁狂）

\end{pastquestion}

\PointSeparator

\section{双相病例分析}\label{users__fpb__26_spring__pkusph_study__tmp__pdfs__prevention-medicine-past-exams__06-ux7cbeux795eux75c5ux5b66ux5f80ux5e74ux9898ux8003ux70b9ux6574ux7406.md__ux53ccux76f8ux75c5ux4f8bux5206ux6790}

\InfoLabel{考频} 1 次

\InfoLabel{对应小节} 四、双相 / 病例

\InfoLabel{匹配依据}
题干涉及「病例」「双相情感障碍I型」「情绪高涨交替」「躁狂发作精神症状」「诊断与鉴别诊断」等表述，与四、双相 /
病例对应。

\ContentBand{知识点原文摘取}

双相情感障碍病例分析要点：\\
·通过病例总结归纳患者存在的临床症状（躁狂/抑郁/混合）\\
·诊断及鉴别诊断（需提供诊断依据和鉴别诊断依据）\\
·治疗原则：综合治疗、全病程治疗、心境稳定剂为核心

\InfoLabel{知识来源} 《精神病学复习整理by18级杨子铭》

\ContentBand{对应往年题}

\begin{pastquestion}

\QuestionLabel{【往年题1｜20级完整试题回忆.docx｜单项选择】}\\
一些病例分析，选择最有可能的诊断

\end{pastquestion}

\PointSeparator

\chapter{抑郁障碍}\label{users__fpb__26_spring__pkusph_study__tmp__pdfs__prevention-medicine-past-exams__06-ux7cbeux795eux75c5ux5b66ux5f80ux5e74ux9898ux8003ux70b9ux6574ux7406.md__ux7b2cux4e94ux7ae0-ux6291ux90c1ux969cux788d}

\section{抑郁障碍治疗原则与药物}\label{users__fpb__26_spring__pkusph_study__tmp__pdfs__prevention-medicine-past-exams__06-ux7cbeux795eux75c5ux5b66ux5f80ux5e74ux9898ux8003ux70b9ux6574ux7406.md__ux6291ux90c1ux969cux788dux6cbbux7597ux539fux5219ux4e0eux836fux7269}

\InfoLabel{考频} 2 次

\InfoLabel{对应小节} 五、抑郁障碍 / 治疗

\InfoLabel{匹配依据} 题干涉及「抗抑郁」「SSRI」「五朵金花」「西酞普兰」「氟西汀」等表述，与五、抑郁障碍 /
治疗对应。

\ContentBand{知识点原文摘取}

抑郁障碍的治疗原则：\\
·药物治疗是核心，心理治疗和物理治疗是辅助\\
·全病程治疗：急性期、巩固期、维持期\\
·个体化用药：安全性、有效性、经济性、适当性\\
·单一用药原则\\
·足量足疗程原则

治疗药物：5-HT再摄取抑制剂SSRI''五朵金花''（西酞普兰、氟西汀、氟伏沙明、帕罗西汀、舍曲林）、SNRI、NaSSA。\\
改良电休克治疗MECT适应症：伴强烈自伤、自杀企图及行为，有明显自责、自罪情况者为首选。

抗抑郁药的不良反应：胃肠道反应、抗胆碱能不良反应、CNS激活导致激越、镇静、性功能障碍、体重增加、静止性震颤、心血管不良反应、过敏反应、中毒。

\InfoLabel{知识来源} 《精神病学复习整理by18级杨子铭》

\ContentBand{对应往年题}

\begin{pastquestion}

\QuestionLabel{【往年题1｜16级预防《精神病学》考题.pdf｜简答】}\\
抗抑郁药物常见的不良反应及处理。

\end{pastquestion}

\begin{pastquestion}

\QuestionLabel{【往年题2｜20级完整试题回忆.docx｜单项选择】}\\
抑郁障碍的治疗分期不包括哪一期 A、急性期 B、巩固期 C、维持期 D、康复期

\end{pastquestion}

\PointSeparator

\section{抑郁障碍临床表现与诊断}\label{users__fpb__26_spring__pkusph_study__tmp__pdfs__prevention-medicine-past-exams__06-ux7cbeux795eux75c5ux5b66ux5f80ux5e74ux9898ux8003ux70b9ux6574ux7406.md__ux6291ux90c1ux969cux788dux4e34ux5e8aux8868ux73b0ux4e0eux8bcaux65ad}

\InfoLabel{考频} 2 次

\InfoLabel{对应小节} 五、抑郁障碍 / 临床表现与诊断

\InfoLabel{匹配依据} 题干涉及「抑郁障碍」「抑郁发作」「抑郁症」「三低」「心境低落」等表述，与五、抑郁障碍 /
临床表现与诊断对应。

\ContentBand{知识点原文摘取}

抑郁症的临床表现：\\
1、核心症状（情感症状）：心境低落、兴趣减退、快感缺失\\
2、附加症状：``三自''（自责、自罪、自杀）``三无''（无用、无助、无望）、焦虑、易激惹\\
3、躯体症状：疲乏、体重变化、食欲变化、睡眠紊乱（入睡困难、早醒）、植物神经紊乱、性功能障碍\\
4、神经认知症状：思考能力下降、注意力集中困难、精神运动性迟滞或激越\\
5、精神病性症状：自罪妄想、虚无妄想、幻听\\
总结：``三低''------思维迟缓、情绪低落、意志消沉

ICD-10诊断标准：核心症状3（心境低落、兴趣减退、活动减少）+ 附加症状7。轻度2+2、中度2+3、重度3+4。病程至少2周。

\InfoLabel{知识来源} 《精神病学复习整理by18级杨子铭》

\ContentBand{对应往年题}

\begin{pastquestion}

\QuestionLabel{【往年题1｜20级完整试题回忆.docx｜单项选择】}\\
相比抑郁障碍，下列哪些症状更常出现于双相障碍 A、嗜睡 B、发病年龄早 C、心境不稳 D、自伤自杀观念或行为
E、以上都是

\end{pastquestion}

\begin{pastquestion}

\QuestionLabel{【往年题2｜21级精神病学考题yys.docx / 21级本科预防二班+胡润泽+精神病+主观题总结.docx｜简答】}\\
抑郁的临床表现

\end{pastquestion}

\PointSeparator

\section{抑郁障碍鉴别与分型}\label{users__fpb__26_spring__pkusph_study__tmp__pdfs__prevention-medicine-past-exams__06-ux7cbeux795eux75c5ux5b66ux5f80ux5e74ux9898ux8003ux70b9ux6574ux7406.md__ux6291ux90c1ux969cux788dux9274ux522bux4e0eux5206ux578b}

\InfoLabel{考频} 1 次

\InfoLabel{对应小节} 五、抑郁障碍 / 鉴别与分型

\InfoLabel{匹配依据} 题干涉及「鉴别」「难治性抑郁」「恶劣心境」「心境恶劣」「2年」等表述，与五、抑郁障碍 /
鉴别与分型对应。

\ContentBand{知识点原文摘取}

抑郁障碍的鉴别诊断：\\
1、精神分裂症：精分的抑郁继发于幻听和妄想，与心境不协调\\
2、双相情感障碍：有躁狂发作，青春期起病，男女均衡\\
3、焦虑障碍：焦虑先出现、占主要地位\\
4、创伤后应激障碍：有创伤史\\
5、强迫症：抑郁表现在强迫症的基础上\\
6、躯体疾病所致的精神障碍：先有躯体疾病才抑郁

其他分型：难治性抑郁（2种不同机制药物足量足疗程后减分率＜20\%）、儿童青少年期抑郁、老年期抑郁、孕产期抑郁、心境恶劣障碍（2年内反复出现症状但达不到抑郁标准）。

\InfoLabel{知识来源} 《精神病学复习整理by18级杨子铭》

\ContentBand{对应往年题}

\begin{pastquestion}

\QuestionLabel{【往年题1｜06-07级 史前精神病考题.doc｜回忆/选择】}\\
恶劣心境病程诊断标准 2年

\end{pastquestion}

\PointSeparator

\chapter{焦虑障碍与强迫障碍}\label{users__fpb__26_spring__pkusph_study__tmp__pdfs__prevention-medicine-past-exams__06-ux7cbeux795eux75c5ux5b66ux5f80ux5e74ux9898ux8003ux70b9ux6574ux7406.md__ux7b2cux516dux7ae0-ux7126ux8651ux969cux788dux4e0eux5f3aux8febux969cux788d}

\section{强迫障碍（临床表现、诊断、治疗）}\label{users__fpb__26_spring__pkusph_study__tmp__pdfs__prevention-medicine-past-exams__06-ux7cbeux795eux75c5ux5b66ux5f80ux5e74ux9898ux8003ux70b9ux6574ux7406.md__ux5f3aux8febux969cux788dux4e34ux5e8aux8868ux73b0ux8bcaux65adux6cbbux7597}

\InfoLabel{考频} 5 次

\InfoLabel{对应小节} 六、强迫障碍

\InfoLabel{匹配依据}
题干涉及「强迫」「强迫症」「强迫观念」「强迫行为」「强迫检查」等表述，与六、强迫障碍对应。

\ContentBand{知识点原文摘取}

强迫障碍的临床特点：\\
1、强迫观念：强迫思维、强迫性穷思竭虑、强迫怀疑、强迫联想、回忆、对立思维\\
2、强迫行为：强迫检查、强迫清洗、强迫询问、强迫性仪式动作、强迫性计数、强迫性整理

ICD-10诊断：强迫思维或动作持续至少2周；必须具备4要点：①起源于内心；②令人不愉快、过分、不合理；③至少一种仍被患者徒劳加以抵制；④实施动作的想法本身也令人不愉快。影响社会功能，排除其他。

治疗：药物治疗（SSRI五朵金花）足量足疗程、个体化、缓慢减药；心理治疗（以暴露和反应预防为基础的认知行为治疗是一线心理治疗）。

\InfoLabel{知识来源} 《精神病学复习整理by18级杨子铭》

\ContentBand{对应往年题}

\begin{pastquestion}

\QuestionLabel{【往年题1｜20级完整试题回忆.docx｜单项选择】}\\
强迫障碍 A、使用SSRI治疗4周，未见明显效果也未见不良反应，医生应建议换药
B、以暴露和反应预防为基础的认知行为治疗是强迫症的一线心理治疗方法 C、患者可以选择药物治疗，心理治疗，物理治疗
D、强迫症是一种慢性精神障碍，应坚持维持期服药1-2年 E、\ldots\ldots{}

\end{pastquestion}

\begin{pastquestion}

\QuestionLabel{【往年题2｜21级精神病学考题yys.docx / 21级本科预防二班+胡润泽+精神病+主观题总结.docx｜简答】}\\
强迫症的治疗策略/治疗方法

\end{pastquestion}

\begin{pastquestion}

\QuestionLabel{【往年题3｜06-07级 史前精神病考题.doc｜回忆/选择】}\\
患者主诉不受控制的要看异性的生殖器，觉得所有人都察觉了自己的行为，感到非常痛苦，这种现象是
内心被洞悉感？强迫症？

\end{pastquestion}

\begin{pastquestion}

\QuestionLabel{【往年题4｜06-07级 史前精神病考题.doc｜回忆/选择】}\\
强迫观念是？不受自己控制的想法or 源于自己内心的思想或观念

\end{pastquestion}

\begin{pastquestion}

\QuestionLabel{【往年题5｜06-07级 史前精神病考题.doc｜回忆/选择】}\\
强迫症的强迫症状与精分的强迫症状区别主要在于前者的------``属我性''

\end{pastquestion}

\PointSeparator

\section{惊恐障碍}\label{users__fpb__26_spring__pkusph_study__tmp__pdfs__prevention-medicine-past-exams__06-ux7cbeux795eux75c5ux5b66ux5f80ux5e74ux9898ux8003ux70b9ux6574ux7406.md__ux60caux6050ux969cux788d}

\InfoLabel{考频} 2 次

\InfoLabel{对应小节} 六、焦虑障碍 / 惊恐障碍

\InfoLabel{匹配依据} 题干涉及「惊恐」「急性焦虑」「濒死感」「失控感」「5-20」等表述，与六、焦虑障碍 /
惊恐障碍对应。

\ContentBand{知识点原文摘取}

惊恐障碍（急性焦虑障碍）：突然发作的、不可预测的、反复出现的、强烈的惊恐体验，一般历时5-20min，伴濒死感/失控感、自主神经功能紊乱。\\
临床表现：突然感到紧张、害怕、恐惧，伴濒死感、失控感；肌肉紧张、全身发抖；严重的自主神经功能紊乱症状（出汗、胸闷、呼吸困难、心动过速等）；场所特异性。\\
诊断：1次发作、1个月以上的担心再次发作，影响工作和生活，排除其他。\\
治疗：抗抑郁药SSRI（五朵金花），认知重构。

\InfoLabel{知识来源} 《精神病学复习整理by18级杨子铭》

\ContentBand{对应往年题}

\begin{pastquestion}

\QuestionLabel{【往年题1｜06-07级 史前精神病考题.doc｜回忆/选择】}\\
惊恐障碍------类冠心病症状，也容易被误诊为心脏病

\end{pastquestion}

\begin{pastquestion}

\QuestionLabel{【往年题2｜06-07级 史前精神病考题.doc｜回忆/选择】}\\
1）女性，近三个月在不同场合发作，胸闷气促濒死感，每次发作不超过半小时，到急诊之前即缓解------惊恐障碍

\end{pastquestion}

\PointSeparator

\section{病理性焦虑与焦虑障碍概述}\label{users__fpb__26_spring__pkusph_study__tmp__pdfs__prevention-medicine-past-exams__06-ux7cbeux795eux75c5ux5b66ux5f80ux5e74ux9898ux8003ux70b9ux6574ux7406.md__ux75c5ux7406ux6027ux7126ux8651ux4e0eux7126ux8651ux969cux788dux6982ux8ff0}

\InfoLabel{考频} 3 次

\InfoLabel{对应小节} 六、焦虑障碍 / 概述

\InfoLabel{匹配依据}
题干涉及「病理性焦虑」「焦虑障碍」「焦虑症状」「特质焦虑」「状态焦虑」等表述，与六、焦虑障碍 / 概述对应。

\ContentBand{知识点原文摘取}

病理性焦虑：持续的紧张不安、无充分现实依据但感到即将要遇到威胁、灾难或大难临头，常伴有明显的自主神经功能紊乱或焦躁不安，伴随明显的主观痛苦或功能受损。\\
特点：持久、无现实基础、痛苦、自主神经紊乱、预感到灾难、缺乏应对能力。

焦虑与恐惧的区别：恐惧是对已知的、外在的、明确的威胁的回应；焦虑是对未知的、内在的、模糊的威胁的回应。

焦虑障碍的治疗：抗抑郁药SSRI、SNRI，5-HT受体激动剂（丁螺环酮）、苯二氮卓类、β受体阻滞剂。

\InfoLabel{知识来源} 《精神病学复习整理by18级杨子铭》

\ContentBand{对应往年题}

\begin{pastquestion}

\QuestionLabel{【往年题1｜16级预防《精神病学》考题.pdf｜简答】}\\
什么是病理性焦虑？其特点是什么？

\end{pastquestion}

\begin{pastquestion}

\QuestionLabel{【往年题2｜20级完整试题回忆.docx｜单项选择】}\\
正常焦虑与焦虑障碍的区别 A、焦虑障碍大多有具体的对象和内容 B、正常焦虑与威胁的程度成比例
C、威胁解除后正常焦虑缓解 D、焦虑障碍没有建设性 E、焦虑障碍降低效率，干扰工作社交生活

\end{pastquestion}

\begin{pastquestion}

\QuestionLabel{【往年题3｜20级完整试题回忆.docx｜单项选择】}\\
广泛性焦虑障碍的临床表现不包括 A、非现实性评估自身或他人遇到的危险 B、坐卧不宁，搓手，来回走动
C、一次发作持续数分钟

\end{pastquestion}

\PointSeparator

\section{焦虑/强迫相关鉴别}\label{users__fpb__26_spring__pkusph_study__tmp__pdfs__prevention-medicine-past-exams__06-ux7cbeux795eux75c5ux5b66ux5f80ux5e74ux9898ux8003ux70b9ux6574ux7406.md__ux7126ux8651ux5f3aux8febux76f8ux5173ux9274ux522b}

\InfoLabel{考频} 3 次

\InfoLabel{对应小节} 六、焦虑强迫 / 鉴别

\InfoLabel{匹配依据} 题干涉及「疑病症」「恐惧症」「肺炎」「癌症」「躯体变形」等表述，与六、焦虑强迫 /
鉴别对应。

\ContentBand{知识点原文摘取}

焦虑与抑郁的鉴别：\\
·焦虑的核心是过度担心\\
·抑郁的核心是负性情感增强、正性情感减少\\
·共同点是躯体症状、睡眠食欲改变

强迫障碍鉴别诊断：\\
1、抑郁障碍：主要根据哪种症状原发并占主要地位\\
2、广泛性焦虑障碍：广泛模糊的不祥预感，无自我抵抗感觉\\
3、精神分裂症：精分无自知力，强迫有自知力\\
4、强迫型人格障碍：人格特征，无明确病程界线

\InfoLabel{知识来源} 《精神病学复习整理by18级杨子铭》

\ContentBand{对应往年题}

\begin{pastquestion}

\QuestionLabel{【往年题1｜21级精神病学考题yys.docx /
21级本科预防二班+胡润泽+精神病+主观题总结.docx｜单项选择】}\\
组织让写间谍日记------问是什么症状

\end{pastquestion}

\begin{pastquestion}

\QuestionLabel{【往年题2｜21级精神病学考题yys.docx /
21级本科预防二班+胡润泽+精神病+主观题总结.docx｜单项选择】}\\
怀疑自己得肺炎天天检查------疑病症/焦虑症/恐惧症鉴别

\end{pastquestion}

\begin{pastquestion}

\QuestionLabel{【往年题3｜21级精神病学考题yys.docx /
21级本科预防二班+胡润泽+精神病+主观题总结.docx｜单项选择】}\\
阴道出血怀疑癌症伴心悸大汗------恐惧症/焦虑症/疑病症

\end{pastquestion}

\PointSeparator

\section{社交焦虑障碍}\label{users__fpb__26_spring__pkusph_study__tmp__pdfs__prevention-medicine-past-exams__06-ux7cbeux795eux75c5ux5b66ux5f80ux5e74ux9898ux8003ux70b9ux6574ux7406.md__ux793eux4ea4ux7126ux8651ux969cux788d}

\InfoLabel{考频} 1 次

\InfoLabel{对应小节} 六、焦虑障碍 / 社交焦虑

\InfoLabel{匹配依据} 题干涉及「社交焦虑」「社交恐惧」「社恐」等表述，与六、焦虑障碍 / 社交焦虑对应。

\ContentBand{知识点原文摘取}

社交焦虑障碍（社交恐惧症）：在社交场合持续紧张或恐惧、回避社交行为。\\
临床表现：显著而持续地担心在公众面前丢丑或尴尬，回避社交，脸红、手抖、不敢对视等。影响个人生活、职业功能、社会关系。\\
诊断：6个月以上社交恐惧、影响功能、排除其他。\\
鉴别：正常羞怯、躯体变形障碍、精神分裂症（被害妄想导致回避）、场所恐惧障碍（害怕独处）。\\
治疗：药物（SSRI、SNRIs、苯二氮卓类、β受体阻滞剂）、心理治疗（社交技巧训练、暴露疗法、认知治疗）。

\InfoLabel{知识来源} 《精神病学复习整理by18级杨子铭》

\ContentBand{对应往年题}

\begin{pastquestion}

\QuestionLabel{【往年题1｜21级精神病学考题yys.docx /
21级本科预防二班+胡润泽+精神病+主观题总结.docx｜单项选择】}\\
社交恐惧障碍案例分析（选项无社交恐惧）

\end{pastquestion}

\PointSeparator

\chapter{心身疾病}\label{users__fpb__26_spring__pkusph_study__tmp__pdfs__prevention-medicine-past-exams__06-ux7cbeux795eux75c5ux5b66ux5f80ux5e74ux9898ux8003ux70b9ux6574ux7406.md__ux7b2cux4e03ux7ae0-ux5fc3ux8eabux75beux75c5}

\section{躯体疾病所致精神障碍的共同特点与处理原则}\label{users__fpb__26_spring__pkusph_study__tmp__pdfs__prevention-medicine-past-exams__06-ux7cbeux795eux75c5ux5b66ux5f80ux5e74ux9898ux8003ux70b9ux6574ux7406.md__ux8eafux4f53ux75beux75c5ux6240ux81f4ux7cbeux795eux969cux788dux7684ux5171ux540cux7279ux70b9ux4e0eux5904ux7406ux539fux5219}

\InfoLabel{考频} 2 次

\InfoLabel{对应小节} 七、心身疾病 / 躯体疾病所致精神障碍

\InfoLabel{匹配依据}
题干涉及「躯体疾病所致」「共同临床特点」「处理原则」「原发病」「营养支持」等表述，与七、心身疾病 /
躯体疾病所致精神障碍对应。

\ContentBand{知识点原文摘取}

掌握：躯体疾病所致精神障碍的共同临床特点与治疗原则、综合医院常见的精神科综合征的识别与处理\\
熟悉：心身医学和会诊-联络精神病学的概念、新的医学模式及临床沟通技巧在综合医院各科临床工作中的应用\\
了解：综合医院精神科会诊的注意事项\\
心身医学-熟悉：研究心理或精神与躯体相互关系的一个医学分支，属于交叉学科，认为心理社会因素作为躯体疾病的重要病因，利用``生物-心理-社会''医学模式，强调``整体医学观''\\
会诊-联络精神病学-熟悉：与``心身医学''通用，但侧重点不同。会诊-联络精神病学侧重从技术服务层面体现精神科与临床各科的关系；而心身医学侧重从系统思想和整体医学观的层面强调心与身的相互关系、更多引发人们关注大脑、心理和躯体之间的相互作用。\\
综合医院常见的精神症状综合征-掌握：\\
·识别原则：等级诊断（重到轻：卒中、物质依赖、精分、双相、抑郁、焦虑）、构筑综合征（由症状构建综合征）\\
1、谵妄（急性脑综合征）：\\
·识别：急性、一过性、广泛性的认知障碍，尤其以意识障碍为主要特征。波动性病程、伴有意识、注意和知觉的损害，是综合医院中引起激越躁动的最常见原因，是严重躯体疾病的信号。\\
·处理：病因治疗、支持治疗、对症治疗（氟哌啶醇）、谨慎使用苯二氮卓类（特例：酒精依赖戒断反应引起的谵妄需要用苯二氮卓类）\\
2、痴呆：严重而持续的认知障碍，以缓慢出现的智能减退为主要特征（认知障碍、记忆障碍、定向障碍、思维缓慢、情绪淡漠、易激惹、社会退缩、攻击行为），伴有人格障碍\\
3、抑郁：\\
·识别：常共患躯体疾病（糖尿病、冠心病、关节炎），以躯体主诉为首诊于各个科室，对健康结局有不良影响。注意是否是疾病或药物引起的继发性抑郁\\
·处理：健康教育、转诊\\
4、焦虑：与焦虑类似的内科疾病（嗜铬细胞瘤、颞叶癫痫、药物戒断）、与内科疾病相似的焦虑（惊恐障碍）\\
5、躯体化：对健康/疾病的过度焦虑，反复就医\\
6、精神病性症状：脑炎患者可表现出类似精分症状（幻觉、妄想）\\
躯体疾病所致精神障碍的共同特点-核心重点/往年简答10分：\\
1、有明确的躯体疾病证据\\
2、时间关系：躯体疾病在先，精神障碍在后\\
3、精神障碍与躯体疾病平行发展，其病程、预后与原发躯体疾病的性质、严重程度、处理效果密切相关\\
4、意识障碍多数有昼轻夜重的特点\\
5、病因与精神障碍没有特异性：不同病因可以表现为相似的精神障碍，相同病因也可表现出不同的精神障碍

\InfoLabel{知识来源} 《精神病学复习整理by18级杨子铭》

\ContentBand{对应往年题}

\begin{pastquestion}

\QuestionLabel{【往年题1｜16级预防《精神病学》考题.pdf｜简答】}\\
躯体疾病所致的精神障碍的共同临床特点与处理原则。

\end{pastquestion}

\begin{pastquestion}

\QuestionLabel{【往年题2｜20级完整试题回忆.docx｜简答】}\\
躯体疾病所致精神障碍的共同临床特点与处理原则

\end{pastquestion}

\PointSeparator

\chapter{物质依赖}\label{users__fpb__26_spring__pkusph_study__tmp__pdfs__prevention-medicine-past-exams__06-ux7cbeux795eux75c5ux5b66ux5f80ux5e74ux9898ux8003ux70b9ux6574ux7406.md__ux7b2cux516bux7ae0-ux7269ux8d28ux4f9dux8d56}

\section{成瘾物质、依赖、耐受与戒断}\label{users__fpb__26_spring__pkusph_study__tmp__pdfs__prevention-medicine-past-exams__06-ux7cbeux795eux75c5ux5b66ux5f80ux5e74ux9898ux8003ux70b9ux6574ux7406.md__ux6210ux763eux7269ux8d28ux4f9dux8d56ux8010ux53d7ux4e0eux6212ux65ad}

\InfoLabel{考频} 3 次

\InfoLabel{对应小节} 八、物质依赖 / 基本概念

\InfoLabel{匹配依据} 题干涉及「成瘾」「物质依赖」「物质滥用」「耐受」「戒断」等表述，与八、物质依赖 /
基本概念对应。

\ContentBand{知识点原文摘取}

八、物质依赖\\
掌握：成瘾物质、物质成瘾-最重点、物质滥用、耐受性、戒断状态的概念，成瘾与疾病和道德的关系。\\
熟悉：物质成瘾的主要原因、酒精戒断反应的主要临床表现与治疗\\
了解：苯二氮卓类药物戒断反应的临床表现与治疗、成瘾的共病、成瘾的机制、成瘾的康复、兴奋剂的损害\\
精神活性物质/成瘾物质-掌握：能够影响人类情绪、行为、改变意识状态，并有致依赖作用的一类化学物质。人们使用这些物质的目的在于取得或保持某些特殊的心理、生理状态。\\
物质滥用-掌握：又称有害性使用，是指一种适应不良方式，由于反复使用成瘾物质导致了明显的不良后果。如不能完成重要的工作、学业，损害了躯体、心理健康，以及导致法律问题等。\\
·强调引起不良后果，可以没有耐受性增强或戒断（如果有 则是物质依赖）\\
精神依赖/心理依赖-熟悉：对药物的强烈渴求，用药出现欣快感和放松宁静感，满足心理需要即正性强化。停药后会产生难以忍受的痛苦和折磨，只得继续使用药物即负性强化。\\
躯体依赖/生理依赖-熟悉：反复使用成瘾物质所造成的一种适应状态，中枢神经系统发生生理、生化改变，以至于需要药物持续地存在于体内，以免发生戒断症状。\\
依赖综合征-熟悉：是一组认知、行为和生理症状群，个体明白使用成瘾物质会带来问题，但还继续使用，导致耐受性增加、戒断症状、强制性觅药行为\\
物质依赖/物质成瘾-核心重点：由于反复使用精神活性物质而出现依赖综合征。强烈渴求、耐受性增加、停用后出现戒断症状、以成瘾物质为中心、不计后果、无法控制的使用。\\
物质依赖与疾病和道德的关系-掌握/往年考论述：\\
·物质依赖是慢性、进行性、复发性的大脑疾病，
常与其它躯体及精神疾病伴发。物质可改变大脑的结构和功能，导致有害行为的发生（失去控制、强制性使用、负性情绪）\\
• 物质依赖是可以治疗的疾病，而不是道德问题、习惯问题。\\
耐受性-掌握：反复使用成瘾物质后，机体变得不敏感，使用原来剂量达不到所需效果，个体需要获得原有药理作用或心理体验，必须增加剂量，也可表现为使用方式的改变。\\
戒断综合征-掌握：停止使用药物或减少使用剂量或使用拮抗剂占据受体后所出现的特殊的、令人痛苦的心理生理症状群。其机制是由于长期用药后，突然停药引起适应性的反跳。\\
复发：物质依赖者在脱毒治疗完成，保持了一段时间的戒断状态以后，又再次使用成瘾物质。\\
物质依赖的病因-熟悉：\\
·个体因素：遗传、个体素质、人格、精神疾病共病、与某种药物或团伙的接近程度\\
·心理因素：正性强化（用药快感）、负性强化（戒断反应）、社会性强化（团伙影响）、习惯性奖赏与强化导致依赖行为定型（每天都做
习惯了）\\
·社会因素：药物可及性、同伴影响、家庭环境、社会稳定程度、文化背景、公众对物质的态度\\
酒精所致精神障碍的临床表现-了解：急性中毒、酒精依赖、戒断综合征、震颤谵妄、酒精中毒所致精神障碍、认知功能损害（维尼克脑病、柯萨科夫综合征、酒相关性痴呆）

\InfoLabel{知识来源} 《精神病学复习整理by18级杨子铭》

\ContentBand{对应往年题}

\begin{pastquestion}

\QuestionLabel{【往年题1｜20级完整试题回忆.docx｜单项选择】}\\
精神活性物质不包括 A、大麻 B、酒精 C、烟草 D、抗抑郁剂 E、海洛因

\end{pastquestion}

\begin{pastquestion}

\QuestionLabel{【往年题2｜06-07级 史前精神病考题.doc｜回忆/选择】}\\
关于戒断状态错误的是------是滥用的特征性表现之一？

\end{pastquestion}

\begin{pastquestion}

\QuestionLabel{【往年题3｜06-07级 史前精神病考题.doc｜回忆/选择】}\\
不属于精神活性物质的是------氯丙嗪

\end{pastquestion}

\PointSeparator

\section{物质依赖与疾病、道德的关系}\label{users__fpb__26_spring__pkusph_study__tmp__pdfs__prevention-medicine-past-exams__06-ux7cbeux795eux75c5ux5b66ux5f80ux5e74ux9898ux8003ux70b9ux6574ux7406.md__ux7269ux8d28ux4f9dux8d56ux4e0eux75beux75c5ux9053ux5fb7ux7684ux5173ux7cfb}

\InfoLabel{考频} 1 次

\InfoLabel{对应小节} 八、物质依赖 / 与道德和疾病

\InfoLabel{匹配依据} 题干涉及「道德」「疾病的关系」「大脑疾病」「慢性、进行性、复发性」等表述，与八、物质依赖
/ 与道德和疾病对应。

\ContentBand{知识点原文摘取}

八、物质依赖\\
掌握：成瘾物质、物质成瘾-最重点、物质滥用、耐受性、戒断状态的概念，成瘾与疾病和道德的关系。\\
熟悉：物质成瘾的主要原因、酒精戒断反应的主要临床表现与治疗\\
了解：苯二氮卓类药物戒断反应的临床表现与治疗、成瘾的共病、成瘾的机制、成瘾的康复、兴奋剂的损害\\
精神活性物质/成瘾物质-掌握：能够影响人类情绪、行为、改变意识状态，并有致依赖作用的一类化学物质。人们使用这些物质的目的在于取得或保持某些特殊的心理、生理状态。\\
物质滥用-掌握：又称有害性使用，是指一种适应不良方式，由于反复使用成瘾物质导致了明显的不良后果。如不能完成重要的工作、学业，损害了躯体、心理健康，以及导致法律问题等。\\
·强调引起不良后果，可以没有耐受性增强或戒断（如果有 则是物质依赖）\\
精神依赖/心理依赖-熟悉：对药物的强烈渴求，用药出现欣快感和放松宁静感，满足心理需要即正性强化。停药后会产生难以忍受的痛苦和折磨，只得继续使用药物即负性强化。\\
躯体依赖/生理依赖-熟悉：反复使用成瘾物质所造成的一种适应状态，中枢神经系统发生生理、生化改变，以至于需要药物持续地存在于体内，以免发生戒断症状。\\
依赖综合征-熟悉：是一组认知、行为和生理症状群，个体明白使用成瘾物质会带来问题，但还继续使用，导致耐受性增加、戒断症状、强制性觅药行为\\
物质依赖/物质成瘾-核心重点：由于反复使用精神活性物质而出现依赖综合征。强烈渴求、耐受性增加、停用后出现戒断症状、以成瘾物质为中心、不计后果、无法控制的使用。\\
物质依赖与疾病和道德的关系-掌握/往年考论述：\\
·物质依赖是慢性、进行性、复发性的大脑疾病，
常与其它躯体及精神疾病伴发。物质可改变大脑的结构和功能，导致有害行为的发生（失去控制、强制性使用、负性情绪）\\
• 物质依赖是可以治疗的疾病，而不是道德问题、习惯问题。\\
耐受性-掌握：反复使用成瘾物质后，机体变得不敏感，使用原来剂量达不到所需效果，个体需要获得原有药理作用或心理体验，必须增加剂量，也可表现为使用方式的改变。\\
戒断综合征-掌握：停止使用药物或减少使用剂量或使用拮抗剂占据受体后所出现的特殊的、令人痛苦的心理生理症状群。其机制是由于长期用药后，突然停药引起适应性的反跳。\\
复发：物质依赖者在脱毒治疗完成，保持了一段时间的戒断状态以后，又再次使用成瘾物质。\\
物质依赖的病因-熟悉：\\
·个体因素：遗传、个体素质、人格、精神疾病共病、与某种药物或团伙的接近程度\\
·心理因素：正性强化（用药快感）、负性强化（戒断反应）、社会性强化（团伙影响）、习惯性奖赏与强化导致依赖行为定型（每天都做
习惯了）\\
·社会因素：药物可及性、同伴影响、家庭环境、社会稳定程度、文化背景、公众对物质的态度\\
酒精所致精神障碍的临床表现-了解：急性中毒、酒精依赖、戒断综合征、震颤谵妄、酒精中毒所致精神障碍、认知功能损害（维尼克脑病、柯萨科夫综合征、酒相关性痴呆）

\InfoLabel{知识来源} 《精神病学复习整理by18级杨子铭》

\ContentBand{对应往年题}

\begin{pastquestion}

\QuestionLabel{【往年题1｜16级预防《精神病学》考题.pdf｜论述】}\\
请论述物质依赖与道德和疾病的关系。

\end{pastquestion}

\PointSeparator

\chapter{儿童青少年精神障碍}\label{users__fpb__26_spring__pkusph_study__tmp__pdfs__prevention-medicine-past-exams__06-ux7cbeux795eux75c5ux5b66ux5f80ux5e74ux9898ux8003ux70b9ux6574ux7406.md__ux7b2cux4e5dux7ae0-ux513fux7ae5ux9752ux5c11ux5e74ux7cbeux795eux969cux788d}

\section{孤独症谱系障碍}\label{users__fpb__26_spring__pkusph_study__tmp__pdfs__prevention-medicine-past-exams__06-ux7cbeux795eux75c5ux5b66ux5f80ux5e74ux9898ux8003ux70b9ux6574ux7406.md__ux5b64ux72ecux75c7ux8c31ux7cfbux969cux788d}

\InfoLabel{考频} 3 次

\InfoLabel{对应小节} 九、儿童青少年 / 孤独症

\InfoLabel{匹配依据} 题干涉及「孤独症」「自闭症」「ASD」「社交交流障碍」「刻板」等表述，与九、儿童青少年 /
孤独症对应。

\ContentBand{知识点原文摘取}

孤独症谱系障碍：一组自幼起病、由于神经发育异常导致的心理行为障碍。严重程度：阿斯伯格综合征＜高功能孤独症＜孤独症。\\
核心症状（二者缺一不可）：\\
①社交交流障碍：一般的交流都无法完成，无法解码\\
②局限的、重复的行为模式、兴趣或活动

诊断：2个核心症状+无明确的病理变化。\\
治疗原则：没有特效药、康复训练、早期治疗、持续治疗、个体化。

\InfoLabel{知识来源} 《精神病学复习整理by18级杨子铭》

\ContentBand{对应往年题}

\begin{pastquestion}

\QuestionLabel{【往年题1｜16级预防《精神病学》考题.pdf｜简答】}\\
孤独症的临床表现有哪些？

\end{pastquestion}

\begin{pastquestion}

\QuestionLabel{【往年题2｜20级完整试题回忆.docx｜单项选择】}\\
孤独症谱系障碍 A、孤独症诊断观察量表（ADOS）可用于孤独症的诊断
B、应用行为分析（ABA）仍是目前有充分循证医学依据的治疗方法 C、抗精神病药物可以控制核心症状
D、以交互性社交交流和社交互动持续性损害，局限及刻板重复的行为、兴趣或活动模式为临床特征
E、孤独症谱系障碍50\%伴有智力发育障碍

\end{pastquestion}

\begin{pastquestion}

\QuestionLabel{【往年题3｜20级完整试题回忆.docx｜简答】}\\
孤独症谱系障碍的核心临床特征

\end{pastquestion}

\PointSeparator

\section{注意缺陷多动障碍}\label{users__fpb__26_spring__pkusph_study__tmp__pdfs__prevention-medicine-past-exams__06-ux7cbeux795eux75c5ux5b66ux5f80ux5e74ux9898ux8003ux70b9ux6574ux7406.md__ux6ce8ux610fux7f3aux9677ux591aux52a8ux969cux788d}

\InfoLabel{考频} 1 次

\InfoLabel{对应小节} 九、儿童青少年 / 多动症

\InfoLabel{匹配依据}
题干涉及「多动」「ADHD」「注意缺陷」「中枢兴奋剂」「智力发育障碍」等表述，与九、儿童青少年 / 多动症对应。

\ContentBand{知识点原文摘取}

注意缺陷多动障碍：一组起病于儿童期的神经发育障碍，表现为与年龄不相符的注意缺陷或多动冲动症状。\\
核心症状：①注意力缺陷（上课无法专注、粗心大意）②多动冲动（无法静坐、话多、打扰别人）\\
临床分型：注意缺陷为主型、多动冲动为主型、混合型。\\
治疗：药物治疗为主（中枢兴奋剂/多巴胺）、心理治疗和社交技能培训辅助。\\
注：儿童多动症的核心症状是注意障碍（非活动过多）。

\InfoLabel{知识来源} 《精神病学复习整理by18级杨子铭》

\ContentBand{对应往年题}

\begin{pastquestion}

\QuestionLabel{【往年题1｜20级完整试题回忆.docx｜单项选择】}\\
智力发育障碍 A、严重程度分为轻度、中度、重度、极重度 B、多于发育阶段（18岁之前）起病 C、查找病因很重要
D、\ldots\ldots{} E、以上都是

\end{pastquestion}

\PointSeparator

\chapter{老年精神障碍}\label{users__fpb__26_spring__pkusph_study__tmp__pdfs__prevention-medicine-past-exams__06-ux7cbeux795eux75c5ux5b66ux5f80ux5e74ux9898ux8003ux70b9ux6574ux7406.md__ux7b2cux5341ux7ae0-ux8001ux5e74ux7cbeux795eux969cux788d}

\section{阿尔茨海默病（AD）}\label{users__fpb__26_spring__pkusph_study__tmp__pdfs__prevention-medicine-past-exams__06-ux7cbeux795eux75c5ux5b66ux5f80ux5e74ux9898ux8003ux70b9ux6574ux7406.md__ux963fux5c14ux8328ux6d77ux9ed8ux75c5ad}

\InfoLabel{考频} 2 次

\InfoLabel{对应小节} 十、老年精神障碍 / AD

\InfoLabel{匹配依据}
题干涉及「阿尔茨海默」「阿尔兹海默」「AD」「老年斑」「神经纤维缠结」等表述，与十、老年精神障碍 / AD对应。

\ContentBand{知识点原文摘取}

·常见类型：阿尔兹海默病（AD）、血管性痴呆（VaD）\\
阿尔兹海默病的核心症状（ABC）：C-认知功能障碍是核心\\
1、日常生活能力下降Activities of daily living\\
2、精神行为症状、人格改变Behavior：妄想（被窃妄想）、幻觉、抑郁、激越、抱怨（觉得家人对自己不好）、违拗、食欲改变（爱吃甜食）、昼夜颠倒。人格改变：敏感多疑、懒散、退缩、自私、暴躁\\
3、认知功能障碍Cognition：\\
·记忆障碍：忘记刚刚发生的事情\\
·视空间和定向障碍：迷路、不知年月\\
·言语障碍：找词困难、命名障碍（不知道手机怎么说）\\
·失认、失用：忘记亲人、不会用生活用品（开锁、勺子、马桶）\\
·全面智能减退：计算、理解、推理等\\
阿尔兹海默病的神经病理学改变：神经炎性斑块、神经原纤维缠结、选择性神经元及突触缺失\\
阿尔兹海默病的药物治疗：胆碱酯酶抑制剂、NMDA受体拮抗剂\\
血管性痴呆（VaD）：由脑血管病变（卒中）导致的多发梗塞性痴呆\\
血管性痴呆的临床表现：痴呆症状+血管病脑损害的局灶性症状，阶梯式发展\\
·早期症状：脑衰弱综合征（头晕、头痛、不适、记忆力减退）\\
·轻度认知障碍（局限性痴呆）\\
·局灶性神经系统症状体征（卒中引起）\\
血管性痴呆的治疗：无特效药、治疗和预防卒中危险因素（控制血压等）\\
AD和VaD的其他治疗：健康教育、心理支持、提高照护质量\\
（二）老年期抑郁\\
老年期抑郁特点：共患躯体疾病、焦虑障碍，易出现睡眠障碍\\
老年期抑郁临床表现：\\
·焦虑、抑郁和激越的混合状态

\InfoLabel{知识来源} 《精神病学复习整理by18级杨子铭》

\ContentBand{对应往年题}

\begin{pastquestion}

\QuestionLabel{【往年题1｜20级完整试题回忆.docx｜单项选择】}\\
阿尔兹海默病的ABC核心症状不包括

\end{pastquestion}

\begin{pastquestion}

\QuestionLabel{【往年题2｜20级完整试题回忆.docx｜单项选择】}\\
68岁男性，脑梗塞，治疗后好转，不久后出现一系列症状\ldots\ldots 最有可能的诊断是 A、阿尔兹海默病 B、血管性痴呆
C、帕金森病痴呆

\end{pastquestion}

\PointSeparator

\section{老年期抑郁}\label{users__fpb__26_spring__pkusph_study__tmp__pdfs__prevention-medicine-past-exams__06-ux7cbeux795eux75c5ux5b66ux5f80ux5e74ux9898ux8003ux70b9ux6574ux7406.md__ux8001ux5e74ux671fux6291ux90c1}

\InfoLabel{考频} 2 次

\InfoLabel{对应小节} 十、老年精神障碍 / 老年抑郁

\InfoLabel{匹配依据}
题干涉及「老年期抑郁」「老年抑郁」「疑病」「焦虑体验」「躯体症状」等表述，与十、老年精神障碍 / 老年抑郁对应。

\ContentBand{知识点原文摘取}

老年抑郁症的转归：1/3改善、1/3不变、1/3恶化\\
抑郁（假性痴呆）与痴呆的鉴别：\\
老年期抑郁的治疗原则-掌握：评估、时机、个体化、单一、起始剂量\\
①充分评估与监测\\
②确定药物治疗时机\\
③个体化合理用药\\
④单一用药\\
⑤起始剂量及调整原则\\
治疗药物：SSRIs-五朵金花、SNRI、NaSSA（跟成人一样），需要注意药物相互作用、不良反应、药代动力学

\InfoLabel{知识来源} 《精神病学复习整理by18级杨子铭》

\ContentBand{对应往年题}

\begin{pastquestion}

\QuestionLabel{【往年题1｜20级完整试题回忆.docx｜单项选择】}\\
老年期抑郁症状特点 A、主观情绪低落突出 B、焦虑体验非常突出 C、躯体症状突出、疑病观念强烈 D、伴有认知功能障碍
E、\ldots\ldots{}

\end{pastquestion}

\begin{pastquestion}

\QuestionLabel{【往年题2｜06-07级 史前精神病考题.doc｜回忆/选择】}\\
Cartod综合征见于 老年抑郁

\end{pastquestion}

\PointSeparator

\chapter{公共精神卫生}\label{users__fpb__26_spring__pkusph_study__tmp__pdfs__prevention-medicine-past-exams__06-ux7cbeux795eux75c5ux5b66ux5f80ux5e74ux9898ux8003ux70b9ux6574ux7406.md__ux7b2cux5341ux4e00ux7ae0-ux516cux5171ux7cbeux795eux536bux751f}

\section{公共精神卫生服务内容与三级预防}\label{users__fpb__26_spring__pkusph_study__tmp__pdfs__prevention-medicine-past-exams__06-ux7cbeux795eux75c5ux5b66ux5f80ux5e74ux9898ux8003ux70b9ux6574ux7406.md__ux516cux5171ux7cbeux795eux536bux751fux670dux52a1ux5185ux5bb9ux4e0eux4e09ux7ea7ux9884ux9632}

\InfoLabel{考频} 2 次

\InfoLabel{对应小节} 十一、公共精神卫生

\InfoLabel{匹配依据}
题干涉及「公共精神卫生」「三级预防」「一级预防」「二级预防」「精神健康促进」等表述，与十一、公共精神卫生对应。

\ContentBand{知识点原文摘取}

十一、公共精神卫生\\
公共精神卫生服务的内容有哪些?（往年考试-论述题20分）\\
1、精神疾病预防：三级预防\\
（1）一级预防（病因预防）：\\
·全人群策略-普遍性干预：健康教育、媒体正确报道、心理咨询热线、限制自杀工具的可及性\\
·高危人群策略-选择性干预：针对精神疾病的子女、有体征或症状的人、有自伤史人\\
·开展科学研究：建立全民信息库、流行病学调查、长期纵向研究\\
（2）二级预防：早发现、早诊断、早治疗\\
（3）三级预防：防止伤残、促进功能恢复、预防复发\\
2、精神健康促进：生、长、活、工作、老的健康\\
·生的健康：减少父母的不良生活方式可以降低新生儿精神疾病\\
·长的健康：原生家庭教育、学前教育、学校心理教育\\
·活的健康：减少贫困、丰富社交活动、锻炼\\
·工作健康：改善工作条件、提高自我效能、企业精神健康重视、员工心理干预\\
·老的健康：提高老年人社会接触，预防社会隔离\\
3、公共精神卫生两大关键问题：费用、消除耻感和社会歧视\\
公共精神卫生服务的需求：\\
1、总体需求：人群精神疾病患病率高、疾病负担高，不同阶段不同时期有所变化（老龄化、物质滥用、性取向的社会接纳）\\
2、不同人群需求不同\\
• 重点人群：儿童、老年、妇女、职场人群、流动人口\\
• 社会支持与患者良好的结局有关，患者在他人的建议下更容易寻求帮助\\
WHO最优精神卫生服务的``金字塔模型''：自我照料、非正式社区服务、初级卫生保健（社区）、门诊服务、住院服务

\InfoLabel{知识来源} 《精神病学复习整理by18级杨子铭》

\ContentBand{对应往年题}

\begin{pastquestion}

\QuestionLabel{【往年题1｜16级预防《精神病学》考题.pdf｜论述】}\\
公共精神卫生服务的内容有哪些？

\end{pastquestion}

\begin{pastquestion}

\QuestionLabel{【往年题2｜20级完整试题回忆.docx｜单项选择】}\\
强迫障碍的预防 A、初级预防的重点应该放在改善心理社会 B、初级预防重点关注共享环境
C、二级预防旨在早期识别、早期治疗，防止疾病发展 D、二级预防应加强对社区医生及全科医生对强迫症的识别，及时转诊
E、三级预防包括提供有效的治疗，预防复发

\end{pastquestion}

\PointSeparator

\chapter{精神康复}\label{users__fpb__26_spring__pkusph_study__tmp__pdfs__prevention-medicine-past-exams__06-ux7cbeux795eux75c5ux5b66ux5f80ux5e74ux9898ux8003ux70b9ux6574ux7406.md__ux7b2cux5341ux4e8cux7ae0-ux7cbeux795eux5eb7ux590d}

本章未在提供的往年题中匹配到明确考点。

\PointSeparator

\chapter*{附：应激相关障碍与其他未完全匹配考点}\label{users__fpb__26_spring__pkusph_study__tmp__pdfs__prevention-medicine-past-exams__06-ux7cbeux795eux75c5ux5b66ux5f80ux5e74ux9898ux8003ux70b9ux6574ux7406.md__ux9644ux5e94ux6fc0ux76f8ux5173ux969cux788dux4e0eux5176ux4ed6ux672aux5b8cux5168ux5339ux914dux8003ux70b9}
\addcontentsline{toc}{chapter}{附：应激相关障碍与其他未完全匹配考点}

\section{未完全匹配到（综合/跨章节）}\label{users__fpb__26_spring__pkusph_study__tmp__pdfs__prevention-medicine-past-exams__06-ux7cbeux795eux75c5ux5b66ux5f80ux5e74ux9898ux8003ux70b9ux6574ux7406.md__ux672aux5b8cux5168ux5339ux914dux5230ux7efcux5408ux8de8ux7ae0ux8282}

\InfoLabel{考频} 5 次

\InfoLabel{对应小节} 附：未完全匹配

\InfoLabel{匹配依据} 未能与知识点总结中的明确小节匹配，根据题干关键词暂归入本章。

\ContentBand{知识点原文摘取}

本部分收录未能精确匹配到单一知识点的综合性题目或跨章节题目，包括综合病例分析、精神障碍分类、ICD-10总体框架、多轴诊断等综合性内容。

\InfoLabel{知识来源} 《精神病学复习整理by18级杨子铭》

\ContentBand{对应往年题}

\begin{pastquestion}

\QuestionLabel{【往年题1｜06-07级 史前精神病考题.doc｜回忆/选择】}\\
控制不住地出现某种观念、冲动或动作，明知道没有必要但控制不了，为此深感痛苦

\end{pastquestion}

\begin{pastquestion}

\QuestionLabel{【往年题2｜06-07级 史前精神病考题.doc｜回忆/选择】}\\
症状学占了大半，考法有二：

\end{pastquestion}

\begin{pastquestion}

\QuestionLabel{【往年题3｜06-07级 史前精神病考题.doc｜回忆/选择】}\\
各论：最重要的是特征性症状

\end{pastquestion}

\begin{pastquestion}

\QuestionLabel{【往年题4｜06-07级 史前精神病考题.doc｜回忆/选择】}\\
某病和某病最重要的区别是

\end{pastquestion}

\begin{pastquestion}

\QuestionLabel{【往年题5｜06-07级 史前精神病考题.doc｜回忆/选择】}\\
2、定义题：

\end{pastquestion}

\PointSeparator

\section{神经症、癔症与躯体形式障碍}\label{users__fpb__26_spring__pkusph_study__tmp__pdfs__prevention-medicine-past-exams__06-ux7cbeux795eux75c5ux5b66ux5f80ux5e74ux9898ux8003ux70b9ux6574ux7406.md__ux795eux7ecfux75c7ux7654ux75c7ux4e0eux8eafux4f53ux5f62ux5f0fux969cux788d}

\InfoLabel{考频} 3 次

\InfoLabel{对应小节} 附：神经症与躯体形式障碍

\InfoLabel{匹配依据}
题干涉及「神经症」「神经衰弱」「癔症」「躯体形式」「躯体化」等表述，与附：神经症与躯体形式障碍对应。

\ContentBand{知识点原文摘取}

神经症特点：有自知力、求治欲、社会功能相对好、病程迁延、无器质性病变基础。\\
·神经症与精神病的主要鉴别在于前者有自知力\\
·癔症（歇斯底里）：由明显心理因素作用于易感个体引起，症状富有表演性、暗示性强\\
·躯体化障碍：以多样多变的躯体症状为主的神经症

癔症治疗最有效的方法是暗示治疗。

\InfoLabel{知识来源} 《精神病学复习整理by18级杨子铭》

\ContentBand{对应往年题}

\begin{pastquestion}

\QuestionLabel{【往年题1｜06-07级 史前精神病考题.doc｜回忆/选择】}\\
病历判断是什么疾病？ 某人一直感到胃肠不适，检查正常---癔症？躯体化障碍？

\end{pastquestion}

\begin{pastquestion}

\QuestionLabel{【往年题2｜06-07级 史前精神病考题.doc｜回忆/选择】}\\
3）女，25岁，近三年反复腹胀、反胃、大便次数增多、胸痛、膝关节疼痛及会阴部不适等，经检查无阳性躯体症状，不愿去精神科就诊，诊断（躯体化障碍、躯体症状自主神经紊乱、癔症性障碍\ldots\ldots？）

\end{pastquestion}

\begin{pastquestion}

\QuestionLabel{【往年题3｜06-07级 史前精神病考题.doc｜回忆/选择】}\\
6）干部，60，近半年来感到工作效率下降、脑力易疲劳、入睡困难，睡眠浅，多梦，醒后不解乏，诊断------神经衰弱、焦虑症？

\end{pastquestion}

\PointSeparator

\section{精神检查、诊断方法与量表}\label{users__fpb__26_spring__pkusph_study__tmp__pdfs__prevention-medicine-past-exams__06-ux7cbeux795eux75c5ux5b66ux5f80ux5e74ux9898ux8003ux70b9ux6574ux7406.md__ux7cbeux795eux68c0ux67e5ux8bcaux65adux65b9ux6cd5ux4e0eux91cfux8868}

\InfoLabel{考频} 1 次

\InfoLabel{对应小节} 附：检查与诊断方法

\InfoLabel{匹配依据}
题干涉及「精神检查」「临床晤谈」「心理测验」「量表」「病史采集」等表述，与附：检查与诊断方法对应。

\ContentBand{知识点原文摘取}

精神检查主要通过与患者交谈和观察来检查精神活动是否异常。\\
精神科诊断的基本步骤：病史采集、躯体检查、精神状况检查、实验室检查、对精神症状的分析。\\
临床晤谈技巧：注意观察、耐心倾听、有针对性提问、不纠正妄想、寻找兴趣话题切入。\\
量表应用：减少主观影响、反映治疗前后动态变化、便于统计分析、避免遗漏症状。但标准化评分不能纠正临床误诊。

\InfoLabel{知识来源} 《精神病学复习整理by18级杨子铭》

\ContentBand{对应往年题}

\begin{pastquestion}

\QuestionLabel{【往年题1｜06-07级 史前精神病考题.doc｜回忆/选择】}\\
精神障碍诊断最主要/常用的方法------临床晤谈（重点鉴别``心理测验''）

\end{pastquestion}

\PointSeparator

\section{创伤后应激障碍（PTSD）与急性应激}\label{users__fpb__26_spring__pkusph_study__tmp__pdfs__prevention-medicine-past-exams__06-ux7cbeux795eux75c5ux5b66ux5f80ux5e74ux9898ux8003ux70b9ux6574ux7406.md__ux521bux4f24ux540eux5e94ux6fc0ux969cux788dptsdux4e0eux6025ux6027ux5e94ux6fc0}

\InfoLabel{考频} 1 次

\InfoLabel{对应小节} 附：应激相关障碍

\InfoLabel{匹配依据}
题干涉及「PTSD」「创伤后应激」「应激反应障碍」「急性应激」「延迟性心因性」等表述，与附：应激相关障碍对应。

\ContentBand{知识点原文摘取}

PTSD（创伤后应激障碍）：应激源往往具有异常惊恐或灾难性质，常引起个体极度恐惧、害怕、无助之感。发病常在遭受创伤后数日至半年内出现。\\
临床表现：反复重现创伤性体验、持续性的警觉性增高、回避与创伤相关的刺激、性格改变（易激惹、退缩）。\\
诊断原则：有明确的创伤事件、上述症状持续存在、影响社会功能。大多数病人1年内恢复。

急性应激障碍：在强烈精神刺激之后数分钟至数小时起病，大多历时短暂，可在几天至一周内恢复。

\InfoLabel{知识来源} 《精神病学复习整理by18级杨子铭》

\ContentBand{对应往年题}

\begin{pastquestion}

\QuestionLabel{【往年题1｜21级精神病学考题yys.docx / 21级本科预防二班+胡润泽+精神病+主观题总结.docx｜论述】}\\
PTSD创伤后应激障碍的临床表现和诊断原则/诊断要点

\end{pastquestion}

\PointSeparator

\chapter*{题量核对}\label{users__fpb__26_spring__pkusph_study__tmp__pdfs__prevention-medicine-past-exams__06-ux7cbeux795eux75c5ux5b66ux5f80ux5e74ux9898ux8003ux70b9ux6574ux7406.md__ux9898ux91cfux6838ux5bf9}
\addcontentsline{toc}{chapter}{题量核对}

按来源文件统计，共 \textbf{104} 题，全部纳入：

\begin{longtable}[]{@{}
  >{\raggedright\arraybackslash}p{(\linewidth - 2\tabcolsep) * \real{0.7000}}
  >{\raggedleft\arraybackslash}p{(\linewidth - 2\tabcolsep) * \real{0.1200}}@{}}
\toprule\noalign{}
\begin{minipage}[b]{\linewidth}\raggedright
来源文件
\end{minipage} & \begin{minipage}[b]{\linewidth}\raggedleft
题量
\end{minipage} \\
\midrule\noalign{}
\endhead
\bottomrule\noalign{}
\endlastfoot
06-07级 史前精神病考题.doc & 56 \\
16级预防《精神病学》考题.pdf & 8 \\
20级完整试题回忆.docx & 31 \\
21级精神病学考题yys.docx / 21级本科预防二班+胡润泽+精神病+主观题总结.docx & 9 \\
\textbf{合计} & \textbf{104} \\
\end{longtable}

\begin{quote}
说明：20级选择题部分选项回忆不全处保留题干与可回忆选项；21级PTSD论述题为当年超考纲内容，已据回忆卷收录。21级两份文件互为补充（yys侧重选择题回忆，胡润泽侧重主观题），合并计为9题。
\end{quote}

\SetSubjectAccent{B6432F}{职业病学}

\part{职业病学}

\chapter*{说明}\label{users__fpb__26_spring__pkusph_study__tmp__pdfs__prevention-medicine-past-exams__07-ux804cux4e1aux75c5ux5b66ux5f80ux5e74ux9898ux603bux7ed3.md__ux8bf4ux660e}
\addcontentsline{toc}{chapter}{说明}

\begin{itemize}
\tightlist
\item
  本文件根据《职业病学知识点总结by17级杨贝妮.pdf》和 \texttt{职业病学/4\ 往年题/} 下所有往年题文件整理。
\item
  已读取并匹配的题源包括：04级、05级、07级、08级、14级、2015题型说明、2015-2022汇总、2020年题图、2021级考题、2022参考样题、24年完整试题回忆。
\item
  知识点正文主要摘自《职业病学知识点总结by17级杨贝妮.pdf》；题源中出现而知识点总结未覆盖或版本不一致者，标注为``未完全匹配到''或``版本提示''。
\item
  一道考题若同时考查多个知识点，分别归入多个考点；英译汉题按对应疾病/概念归入。
\item
  \texttt{2015-2022职业病往年题及答案汇总} 内含 2022参考样题、2015、2018、2020
  片段；与独立文件重复处保留题源来源，便于反向核对。
\item
  \texttt{2020年职业病考题.jpg}
  为长截图，直接按可见原图核对；截图从选择题第4题附近开始，前3题若原图未显示，则仅以 \texttt{2015-2022}
  汇总中的2020整理条目保留。
\item
  排序依据为知识点总结章节顺序；同章内优先排列高频和病例题。
\end{itemize}

\PointSeparator

\chapter{职业病总论}\label{users__fpb__26_spring__pkusph_study__tmp__pdfs__prevention-medicine-past-exams__07-ux804cux4e1aux75c5ux5b66ux5f80ux5e74ux9898ux603bux7ed3.md__ux7b2cux4e00ux7ae0-ux804cux4e1aux75c5ux603bux8bba}

\section{职业病定义、法定职业病与职业病分类目录}\label{users__fpb__26_spring__pkusph_study__tmp__pdfs__prevention-medicine-past-exams__07-ux804cux4e1aux75c5ux5b66ux5f80ux5e74ux9898ux603bux7ed3.md__ux804cux4e1aux75c5ux5b9aux4e49ux6cd5ux5b9aux804cux4e1aux75c5ux4e0eux804cux4e1aux75c5ux5206ux7c7bux76eeux5f55}

\FrequencyPill{高频}

\InfoLabel{对应小节} 第一章 / 一、职业病定义；二、职业病分类和目录；六、职业病诊断

\InfoLabel{匹配依据}
多年大题、填空、翻译和简答反复考``职业病定义''\,``法定职业病''\,``职业病目录分类''\,``十大类名称及举例''。

\ContentBand{知识点原文摘取}

职业病是指企业、事业单位和个体经济组织等用人单位的劳动者在职业活动中，因接触粉尘、放射性物质和其他有毒、有害因素而引起的疾病。

职业病分类和目录：十大类132种（2013）：职业性尘肺病及其他呼吸系统疾病，职业性皮肤病，职业性眼病，职业性耳鼻喉口腔疾病，职业性化学中毒，物理因素所致职业病，职业性放射性疾病，职业性传染病，职业性肿瘤，其他职业病。

职业病诊断法定依据：《职业病目录》，十大类132种。

\WarningLabel{版本提示}
2021级题源出现``2024年，职业病防治法中共有\_\_类职业病，\textbf{种，新增的两类分别是}''这一题，原知识点总结仍为2013版``10类132种''，需按授课年份老师口径或最新版目录另行核对。

\ContentBand{对应往年题}

\begin{pastquestion}

\QuestionLabel{【04级-职业病考题.docx｜大题】}\\
职业病的分类，每个举三个例子；职业病定义，目录，各举三例；职业病目录分类，每个举3种疾病。

\end{pastquestion}

\begin{pastquestion}

\QuestionLabel{【05级职业病考试题.doc｜大题】}\\
职业病定义及分类目录，每种列举2个疾病。

\end{pastquestion}

\begin{pastquestion}

\QuestionLabel{【08级-职业病考题.doc｜填空】}\\
法定职业病的定义：（）、（）、单位和个体经济（）的劳动者进行（）活动中\ldots\ldots；职业病共有（10）类（115）种。

\end{pastquestion}

\begin{pastquestion}

\QuestionLabel{【14级职业病考题.pdf｜英译汉】}\\
法定职业病。

\end{pastquestion}

\begin{pastquestion}

\QuestionLabel{【14级职业病考题.pdf｜填空/简答】}\\
我国职业病分类目录共\_\_类\_\_种；列举我国职业病分类的10类名称。

\end{pastquestion}

\begin{pastquestion}

\QuestionLabel{【2015年职业病考试题型.docx｜填空举例】}\\
窒息性气体分几类；尘肺的分类/分期。

\end{pastquestion}

\begin{pastquestion}

\QuestionLabel{【2015-2022职业病往年题及答案汇总.docx｜2022参考样题/问答】}\\
我国现行职业病目录共分几类？

\end{pastquestion}

\begin{pastquestion}

\QuestionLabel{【2015-2022职业病往年题及答案汇总.docx｜2018/翻译与简答】}\\
prescript occupational disease：法定职业病；列举我国职业病分类的10类名称。

\end{pastquestion}

\begin{pastquestion}

\QuestionLabel{【2020年职业病考题.jpg｜填空1】}\\
我国现行职业病分类和目录中，职业病分为\_\_大类\_\_种。

\end{pastquestion}

\begin{pastquestion}

\QuestionLabel{【2021级职业病考题yys.docx｜填空】}\\
2024年，职业病防治法中共有\_\_类职业病，\textbf{种，新增的两类分别是}。

\end{pastquestion}

\PointSeparator

\section{职业健康监护、职业健康检查与职业病诊断}\label{users__fpb__26_spring__pkusph_study__tmp__pdfs__prevention-medicine-past-exams__07-ux804cux4e1aux75c5ux5b66ux5f80ux5e74ux9898ux603bux7ed3.md__ux804cux4e1aux5065ux5eb7ux76d1ux62a4ux804cux4e1aux5065ux5eb7ux68c0ux67e5ux4e0eux804cux4e1aux75c5ux8bcaux65ad}

\FrequencyPill{中频}

\InfoLabel{对应小节} 第一章 / 四、职业健康监护；六、职业病诊断

\InfoLabel{匹配依据} 07填空、08选择/多选和24多选考健康监护内容、职业健康检查类型、职业病诊断要素。

\ContentBand{知识点原文摘取}

职业健康监护：以预防为目的，根据劳动者职业接触史，通过定期/不定期的医学健康检查和健康相关资料的收集，连续性地监测劳动者的健康状况，分析劳动者健康变化与所接触的职业病危害因素的关系，并及时将健康检查和资料分析结果报告给用人单位和劳动者本人。

职业健康监护的内容包括职业健康检查和职业健康监护档案管理。职业健康检查包括：上岗前职业健康检查、在岗期间职业健康检查、离岗时职业健康检查、离岗后健康检查、应急健康检查。

职业病诊断工作的具体内容：职业史；职业危害因素接触情况；临床情况；实验室检查。职业病集体诊断的原则：三人以上具有资格的医师作出诊断。

\ContentBand{对应往年题}

\begin{pastquestion}

\QuestionLabel{【07级-职业病考题.pdf｜填空】}\\
健康监测的内容。

\end{pastquestion}

\begin{pastquestion}

\QuestionLabel{【08级-职业病考题.doc｜单选/多选】}\\
职业病的诊断；职业病的诊断。

\end{pastquestion}

\begin{pastquestion}

\QuestionLabel{【2015-2022职业病往年题及答案汇总.docx｜2022参考样题/填空】}\\
职业病法规定职业病健康监护包括上岗前、在岗期间、离岗时的体检，离岗后的追踪观察。

\end{pastquestion}

\begin{pastquestion}

\QuestionLabel{【24年完整试题回忆.docx｜多选1】}\\
根据《职业健康检查管理办法》，职业健康检查包括：上岗前、在岗期间、离岗时、离岗后、应急。

\end{pastquestion}

\begin{pastquestion}

\QuestionLabel{【24年完整试题回忆.docx｜多选7】}\\
职业健康检查的目标指：职业禁忌症、职业病等。

\end{pastquestion}

\PointSeparator

\section{物理因素所致职业病}\label{users__fpb__26_spring__pkusph_study__tmp__pdfs__prevention-medicine-past-exams__07-ux804cux4e1aux75c5ux5b66ux5f80ux5e74ux9898ux603bux7ed3.md__ux7269ux7406ux56e0ux7d20ux6240ux81f4ux804cux4e1aux75c5}

\FrequencyPill{中频}

\InfoLabel{对应小节} 第一章 / 二、职业病分类和目录；职业性器官系统疾病

\InfoLabel{匹配依据} 07填空、2022样题填空、24简答均直接考``物理因素所致职业病包括哪些''。

\ContentBand{知识点原文摘取}

职业病分类和目录中，物理因素所致职业病7种。题源答案整理为：中暑、减压病、高原病、航空病、手臂振动病、激光所致眼损伤、冻伤。

\ContentBand{对应往年题}

\begin{pastquestion}

\QuestionLabel{【07级-职业病考题.pdf｜填空】}\\
物理因素职业病。

\end{pastquestion}

\begin{pastquestion}

\QuestionLabel{【2015-2022职业病往年题及答案汇总.docx｜2022参考样题/填空】}\\
职业病目录中的物理因素职业病包括：中暑，减压病，高原病，航空病，手臂振动病，激光所致眼损伤，冻伤，7种。

\end{pastquestion}

\begin{pastquestion}

\QuestionLabel{【2022职业病参考样题及答案整理.docx｜填空】}\\
同上。

\end{pastquestion}

\begin{pastquestion}

\QuestionLabel{【24年完整试题回忆.docx｜简答2】}\\
物理因素所致职业病包括哪些。

\end{pastquestion}

\PointSeparator

\section{职业性神经系统疾病与神经毒物}\label{users__fpb__26_spring__pkusph_study__tmp__pdfs__prevention-medicine-past-exams__07-ux804cux4e1aux75c5ux5b66ux5f80ux5e74ux9898ux603bux7ed3.md__ux804cux4e1aux6027ux795eux7ecfux7cfbux7edfux75beux75c5ux4e0eux795eux7ecfux6bd2ux7269}

\FrequencyPill{中频}

\InfoLabel{对应小节} 第一章 / 八、职业性器官系统疾病 / 职业性神经系统疾病

\InfoLabel{匹配依据} 08简答要求列举神经毒物及神经症状，24和08多选考中毒性脑病相关毒物。

\ContentBand{知识点原文摘取}

神经毒物：金属、类金属及其有机化合物（铅、汞、砷）；有机溶剂（汽油、苯、甲苯）；窒息性毒物（一氧化碳、硫化氢、氰化物）；农药（有机磷、氨基甲酸酯类、杀虫脒）；其他（氯丙烯、溴甲烷、碘甲烷）。

职业性神经系统疾病临床类型：急性中毒性脑病；慢性中毒性脑病；中毒性周围神经病；中毒性神经衰弱综合征。

\ContentBand{对应往年题}

\begin{pastquestion}

\QuestionLabel{【08级-职业病考题.doc｜简答1】}\\
列举5种神经毒物，主要的神经症状是什么。

\end{pastquestion}

\begin{pastquestion}

\QuestionLabel{【08级-职业病考题.doc｜多选】}\\
可引起中毒性脑病的毒物。

\end{pastquestion}

\begin{pastquestion}

\QuestionLabel{【24年完整试题回忆.docx｜多选2】}\\
引起中毒性脑病的毒物包括：氨、金属汞、二硫化碳、正己烷、苯。

\end{pastquestion}

\PointSeparator

\section{职业相关病概念（未完全匹配到完整原文）}\label{users__fpb__26_spring__pkusph_study__tmp__pdfs__prevention-medicine-past-exams__07-ux804cux4e1aux75c5ux5b66ux5f80ux5e74ux9898ux603bux7ed3.md__ux804cux4e1aux76f8ux5173ux75c5ux6982ux5ff5ux672aux5b8cux5168ux5339ux914dux5230ux5b8cux6574ux539fux6587}

\FrequencyPill{低频}

\InfoLabel{对应小节} 第一章 / 职业病定义与职业性疾患相关概念

\InfoLabel{匹配依据}
08级多选出现``职业相关病的概念''。《职业病学知识点总结》主要展开法定职业病定义、诊断和管理，未完整摘出``职业相关病''定义。

\ContentBand{知识点原文摘取}

未完全匹配到完整原文。可按职业卫生学通用概念理解：职业相关疾病通常指职业因素不是唯一病因，但可影响疾病发生、发展或加重的疾病；其病因多因素，不一定属于法定职业病。

\ContentBand{对应往年题}

\begin{pastquestion}

\QuestionLabel{【08级-职业病考题.doc｜多选】}\\
职业相关病的概念。

\end{pastquestion}

\PointSeparator

\chapter{肺尘埃沉着病（尘肺）}\label{users__fpb__26_spring__pkusph_study__tmp__pdfs__prevention-medicine-past-exams__07-ux804cux4e1aux75c5ux5b66ux5f80ux5e74ux9898ux603bux7ed3.md__ux7b2cux4e8cux7ae0-ux80baux5c18ux57c3ux6c89ux7740ux75c5ux5c18ux80ba}

\section{尘肺定义、分类与英译汉}\label{users__fpb__26_spring__pkusph_study__tmp__pdfs__prevention-medicine-past-exams__07-ux804cux4e1aux75c5ux5b66ux5f80ux5e74ux9898ux603bux7ed3.md__ux5c18ux80baux5b9aux4e49ux5206ux7c7bux4e0eux82f1ux8bd1ux6c49}

\FrequencyPill{高频}

\InfoLabel{对应小节} 第二章 / 一、定义；二、分类

\InfoLabel{匹配依据} 尘肺/矽肺反复出现于英译汉、分类填空、选择和病例题。

\ContentBand{知识点原文摘取}

尘肺定义：职业活动中长期吸入生产性粉尘引起的以肺组织弥漫性纤维化为主的一种全身性疾病。

分类：矽肺（游离SiO2）；硅酸盐肺（结合SiO2，如石棉、水泥、滑石、云母）；炭素尘肺（煤、石墨、活性炭、碳黑）；混合性尘肺（煤矽肺、电焊工尘肺）；其他尘肺（铝肺、磷灰石肺、铍肺）。

\ContentBand{对应往年题}

\begin{pastquestion}

\QuestionLabel{【04级-职业病考题.docx｜英译汉/大题】}\\
还有5个英文，很简单都是病名；尘肺的病理分级及X线表现；尘肺的病理类型及X线表现。

\end{pastquestion}

\begin{pastquestion}

\QuestionLabel{【05级职业病考试题.doc｜英译汉】}\\
尘肺。

\end{pastquestion}

\begin{pastquestion}

\QuestionLabel{【07级-职业病考题.pdf｜翻译】}\\
矽肺。

\end{pastquestion}

\begin{pastquestion}

\QuestionLabel{【08级-职业病考题.doc｜英译汉/多选】}\\
石棉肺；属于硅酸盐尘肺的有；属于煤工尘肺的有。

\end{pastquestion}

\begin{pastquestion}

\QuestionLabel{【14级职业病考题.pdf｜英译汉/填空】}\\
尘肺病；尘肺病的病理类型有\_\_、\textbf{、}。

\end{pastquestion}

\begin{pastquestion}

\QuestionLabel{【2015年职业病考试题型.docx｜填空举例】}\\
尘肺的分类/分期。

\end{pastquestion}

\begin{pastquestion}

\QuestionLabel{【2015-2022职业病往年题及答案汇总.docx｜2015/填空】}\\
尘肺分为：矽肺、硅酸盐肺、炭素尘肺、金属尘肺、混合型尘肺。

\end{pastquestion}

\begin{pastquestion}

\QuestionLabel{【2015-2022职业病往年题及答案汇总.docx｜2018/翻译】}\\
pneumoconiosis：尘肺病。

\end{pastquestion}

\begin{pastquestion}

\QuestionLabel{【2022职业病参考样题及答案整理.docx｜翻译】}\\
silicosis：矽肺（或硅肺）。

\end{pastquestion}

\begin{pastquestion}

\QuestionLabel{【2021级职业病考题yys.docx｜英译中】}\\
Silicosis。

\end{pastquestion}

\begin{pastquestion}

\QuestionLabel{【24年完整试题回忆.docx｜英译汉】}\\
silicosis。

\end{pastquestion}

\PointSeparator

\section{影响尘肺病发生发展的因素}\label{users__fpb__26_spring__pkusph_study__tmp__pdfs__prevention-medicine-past-exams__07-ux804cux4e1aux75c5ux5b66ux5f80ux5e74ux9898ux603bux7ed3.md__ux5f71ux54cdux5c18ux80baux75c5ux53d1ux751fux53d1ux5c55ux7684ux56e0ux7d20}

\FrequencyPill{高频}

\InfoLabel{对应小节} 第二章 / 三、影响因素

\InfoLabel{匹配依据} 04、05、07、14、2022样题、2021级均直接考``影响尘肺发病/进展/发生发展的因素''。

\ContentBand{知识点原文摘取}

空气中的粉尘浓度，特别是SiO2含量的高低；粉尘的性质（种类）；空气中粉尘的分散度，分散度越高致病性越强；粉尘颗粒是否新鲜及表面粗糙程度；接触粉尘的时间；机体一般状况，特别是上呼吸道和肺部情况；防尘措施。

\ContentBand{对应往年题}

\begin{pastquestion}

\QuestionLabel{【04级-职业病考题.docx｜大题】}\\
影响尘肺发病的因素，矽肺的并发症。

\end{pastquestion}

\begin{pastquestion}

\QuestionLabel{【05级职业病考试题.doc｜大题】}\\
影响尘肺发病及进展的因素。

\end{pastquestion}

\begin{pastquestion}

\QuestionLabel{【07级-职业病考题.pdf｜大题】}\\
尘肺发病及进展影响因素。

\end{pastquestion}

\begin{pastquestion}

\QuestionLabel{【14级职业病考题.pdf｜简答】}\\
影响尘肺病发生发展的因素。

\end{pastquestion}

\begin{pastquestion}

\QuestionLabel{【2015年职业病考试题型.docx｜简答举例】}\\
粉尘危害的决定因素有哪些。

\end{pastquestion}

\begin{pastquestion}

\QuestionLabel{【2015-2022职业病往年题及答案汇总.docx｜2022参考样题/问答】}\\
影响尘肺病发生发展的主要因素有哪些？

\end{pastquestion}

\begin{pastquestion}

\QuestionLabel{【2022职业病参考样题及答案整理.docx｜问答】}\\
影响尘肺病发生发展的主要因素有哪些？

\end{pastquestion}

\begin{pastquestion}

\QuestionLabel{【2021级职业病考题yys.docx｜论述】}\\
影响尘肺病发病的主要因素。

\end{pastquestion}

\PointSeparator

\section{尘肺X线表现、诊断分期与病例分析}\label{users__fpb__26_spring__pkusph_study__tmp__pdfs__prevention-medicine-past-exams__07-ux804cux4e1aux75c5ux5b66ux5f80ux5e74ux9898ux603bux7ed3.md__ux5c18ux80baxux7ebfux8868ux73b0ux8bcaux65adux5206ux671fux4e0eux75c5ux4f8bux5206ux6790}

\FrequencyPill{高频}

\InfoLabel{对应小节} 第二章 / X线表现；诊断

\InfoLabel{匹配依据} 04大题、08石棉肺X线、2022样题单选、2020/2021/24病例均考X线小阴影、分期、诊断依据。

\ContentBand{知识点原文摘取}

X线表现：小阴影是在X射线胸片上肺野内直径或宽度不超过10mm的阴影。圆形阴影：p型\textless1.5mm，q型1.5-3mm，r型\textgreater3mm；不规则阴影：s型\textless1.5mm，t型1.5-3mm，u型\textgreater3mm。大阴影直径和宽度大于10mm。胸膜表现包括肥厚、胸膜斑、钙化。密集度分四大级十二小级。

诊断时要写``职业性（尘肺类型）几期''。尘肺壹期：总体密集度1级的小阴影，分布范围至少达到2个肺区；或接触石棉粉尘并出现相应胸膜斑等。

\ContentBand{对应往年题}

\begin{pastquestion}

\QuestionLabel{【04级-职业病考题.docx｜大题】}\\
尘肺的病理分级及X线表现；尘肺的病理类型及X线表现。

\end{pastquestion}

\begin{pastquestion}

\QuestionLabel{【08级-职业病考题.doc｜简答/病例】}\\
石棉肺的X线表现；病例分析：可能诊断（铅中毒）之外，该年另有尘肺相关选择多选。

\end{pastquestion}

\begin{pastquestion}

\QuestionLabel{【14级职业病考题.pdf｜病例】}\\
一段很长的尘肺病病例：诊断及诊断依据；鉴别诊断；治疗原则。

\end{pastquestion}

\begin{pastquestion}

\QuestionLabel{【2015-2022职业病往年题及答案汇总.docx｜2022参考样题/单选18】}\\
尘肺病胸部X线表现中，s形小阴影是指直径小于1.5mm的小阴影。

\end{pastquestion}

\begin{pastquestion}

\QuestionLabel{【2015-2022职业病往年题及答案汇总.docx｜2022参考样题/病例】}\\
诊断：职业性煤工尘肺壹期。

\end{pastquestion}

\begin{pastquestion}

\QuestionLabel{【2020年职业病考题.jpg｜病例1】}\\
男，60岁，1968-1984年从事采掘工作，干式作业，粉尘浓度大，胸片小阴影密集度分布；问最可能诊断及依据、鉴别诊断、治疗原则。

\end{pastquestion}

\begin{pastquestion}

\QuestionLabel{【2021级职业病考题yys.docx｜病例2】}\\
男，70岁，煤矿井下采煤工30年，干式作业，粉尘很大，胸片双肺弥漫p/q型小阴影；问最可能诊断、诊断依据、鉴别诊断、治疗原则。

\end{pastquestion}

\begin{pastquestion}

\QuestionLabel{【24年完整试题回忆.docx｜病例1】}\\
男，70岁，煤矿井下采煤工27年，干式作业，胸片双肺弥漫p/q型小阴影；问诊断、诊断依据、鉴别诊断、治疗原则。

\end{pastquestion}

\PointSeparator

\section{矽肺并发症与鉴别诊断}\label{users__fpb__26_spring__pkusph_study__tmp__pdfs__prevention-medicine-past-exams__07-ux804cux4e1aux75c5ux5b66ux5f80ux5e74ux9898ux603bux7ed3.md__ux77fdux80baux5e76ux53d1ux75c7ux4e0eux9274ux522bux8bcaux65ad}

\FrequencyPill{高频}

\InfoLabel{对应小节} 第二章 / 诊断；尘肺结核早期诊断

\InfoLabel{匹配依据} 04、07、2022样题、2020图均考矽肺并发症、鉴别诊断。

\ContentBand{知识点原文摘取}

题源答案整理：矽肺常见并发症包括肺结核、支气管和肺部感染、自发性气胸、肺源性心脏病。鉴别诊断包括血行播散型肺结核、肺癌、特发性肺纤维化、结节病、含铁血黄素沉着症、肺泡微石症等。

知识总结原文提示：尘肺易并发结核；尘肺结核早期诊断注意长期低烧、咯血、血沉明显增快、痰查结核菌阳性、动态胸片肺尖和肺中上野结节影增加等。

\ContentBand{对应往年题}

\begin{pastquestion}

\QuestionLabel{【04级-职业病考题.docx｜大题】}\\
影响尘肺发病的因素，矽肺的并发症。

\end{pastquestion}

\begin{pastquestion}

\QuestionLabel{【07级-职业病考题.pdf｜大题】}\\
矽肺并发症及鉴别诊断。

\end{pastquestion}

\begin{pastquestion}

\QuestionLabel{【2015-2022职业病往年题及答案汇总.docx｜2022参考样题/问答3】}\\
矽肺有哪些并发症，应与哪些肺部疾病鉴别？

\end{pastquestion}

\begin{pastquestion}

\QuestionLabel{【2022职业病参考样题及答案整理.docx｜问答3】}\\
矽肺有哪些并发症，应与哪些肺部疾病鉴别？

\end{pastquestion}

\begin{pastquestion}

\QuestionLabel{【2020年职业病考题.jpg｜简答2】}\\
矽肺应与哪些肺部疾病进行鉴别？矽肺的常见并发症有哪些？

\end{pastquestion}

\PointSeparator

\section{石棉肺、胸膜斑和硅酸盐尘肺}\label{users__fpb__26_spring__pkusph_study__tmp__pdfs__prevention-medicine-past-exams__07-ux804cux4e1aux75c5ux5b66ux5f80ux5e74ux9898ux603bux7ed3.md__ux77f3ux68c9ux80baux80f8ux819cux6591ux548cux7845ux9178ux76d0ux5c18ux80ba}

\FrequencyPill{中高频}

\InfoLabel{对应小节} 第二章 / 分类；X线表现

\InfoLabel{匹配依据} 08、14、2022、2020、24多次考石棉肺胸膜表现、胸膜斑定义、硅酸盐尘肺。

\ContentBand{知识点原文摘取}

硅酸盐肺：结合SiO2，如石棉、水泥、滑石、云母。胸膜表现包括肥厚、胸膜斑、钙化。

题源答案整理：胸膜斑是指除肺尖和肋膈角以外的厚度大于5mm的局限性胸膜增厚或局限性钙化胸膜斑块。硅酸盐肺包括石棉肺、滑石尘肺、水泥尘肺、云母尘肺、陶工尘肺。

\ContentBand{对应往年题}

\begin{pastquestion}

\QuestionLabel{【08级-职业病考题.doc｜填空/多选/简答】}\\
胸膜斑的定义；石棉肺的主要病理改变；石棉肺的X线表现；属于硅酸盐尘肺的有。

\end{pastquestion}

\begin{pastquestion}

\QuestionLabel{【14级职业病考题.pdf｜填空/多选】}\\
胸膜斑定义；属于硅酸盐尘肺的是；属于石棉肺的病理改变的是。

\end{pastquestion}

\begin{pastquestion}

\QuestionLabel{【2015-2022职业病往年题及答案汇总.docx｜2022参考样题/多选2/填空5】}\\
石棉肺的胸膜改变包括胸膜增厚、胸膜斑、胸膜钙化；胸膜斑定义。

\end{pastquestion}

\begin{pastquestion}

\QuestionLabel{【2020年职业病考题.jpg｜选择9/填空4/填空5】}\\
石棉肺的常见病变不包括；胸膜斑定义；《职业病分类和目录》中的硅酸盐肺相关填空。

\end{pastquestion}

\begin{pastquestion}

\QuestionLabel{【24年完整试题回忆.docx｜单选18/多选10】}\\
胸膜斑见于哪种尘肺病；石棉肺的胸膜表现。

\end{pastquestion}

\PointSeparator

\section{粉尘作业职业禁忌证}\label{users__fpb__26_spring__pkusph_study__tmp__pdfs__prevention-medicine-past-exams__07-ux804cux4e1aux75c5ux5b66ux5f80ux5e74ux9898ux603bux7ed3.md__ux7c89ux5c18ux4f5cux4e1aux804cux4e1aux7981ux5fccux8bc1}

\FrequencyPill{中频}

\InfoLabel{对应小节} 第二章 / 尘肺影响因素、诊断与职业健康检查相关内容

\InfoLabel{匹配依据} 08、2022样题、24均以选择题考``粉尘作业职业禁忌证''。

\ContentBand{知识点原文摘取}

知识点总结未单列完整职业禁忌证清单；题源答案反复出现``活动性肺结核''为粉尘作业职业禁忌证。

\ContentBand{对应往年题}

\begin{pastquestion}

\QuestionLabel{【08级-职业病考题.doc｜单选】}\\
粉尘作业的职业禁忌症。

\end{pastquestion}

\begin{pastquestion}

\QuestionLabel{【2015-2022职业病往年题及答案汇总.docx｜2022参考样题/单选13】}\\
粉尘作业的职业禁忌证是：活动性肺结核。

\end{pastquestion}

\begin{pastquestion}

\QuestionLabel{【2022职业病参考样题及答案整理.docx｜单选13】}\\
粉尘作业的职业禁忌证是：活动性肺结核。

\end{pastquestion}

\begin{pastquestion}

\QuestionLabel{【24年完整试题回忆.docx｜单选19】}\\
粉尘工作的职业禁忌证。

\end{pastquestion}

\PointSeparator

\chapter{铅中毒}\label{users__fpb__26_spring__pkusph_study__tmp__pdfs__prevention-medicine-past-exams__07-ux804cux4e1aux75c5ux5b66ux5f80ux5e74ux9898ux603bux7ed3.md__ux7b2cux4e09ux7ae0-ux94c5ux4e2dux6bd2}

\section{铅接触行业、吸收排泄、诊断值和临床损害系统}\label{users__fpb__26_spring__pkusph_study__tmp__pdfs__prevention-medicine-past-exams__07-ux804cux4e1aux75c5ux5b66ux5f80ux5e74ux9898ux603bux7ed3.md__ux94c5ux63a5ux89e6ux884cux4e1aux5438ux6536ux6392ux6cc4ux8bcaux65adux503cux548cux4e34ux5e8aux635fux5bb3ux7cfbux7edf}

\FrequencyPill{高频}

\InfoLabel{对应小节} 第三章 / 铅在体内的代谢；实验室检查；诊断

\InfoLabel{匹配依据} 07、08、14、2022、2020、2021、24反复考铅接触行业、吸收途径、血铅/尿铅诊断值、损害系统。

\ContentBand{知识点原文摘取}

铅及其无机化合物：呼吸道与消化道吸收，一般不经皮肤吸收。铅的有机化合物除呼吸道与消化道吸收外，可由皮肤吸收。进入人体的铅主要由尿液排出，部分经由胃肠道排出。

靶器官：血液、神经、消化、肾脏等。

实验室检查：血铅反映近期铅接触水平，诊断值2.91μmol/L（600μg/L）；尿铅可反映近期铅吸收状况，诊断值0.58μmol/L（120μg/L）。

职业性慢性铅中毒诊断：依据明确职业接触史；有神经、消化、造血系统为主的临床表现；有关实验室检查结果；结合作业环境状况调查；排除其他原因引起的类似疾病。

\ContentBand{对应往年题}

\begin{pastquestion}

\QuestionLabel{【07级-职业病考题.pdf｜多选/单选】}\\
铅接触行业；驱铅药物；尿铅值诊断。

\end{pastquestion}

\begin{pastquestion}

\QuestionLabel{【08级-职业病考题.doc｜单选/病例】}\\
长期接触铅主要引起（）系统损害；不属于急性铅中毒的表现是（铅麻痹）；病例分析可能诊断（铅中毒）、诊断依据、为明确诊断需要做的检查、鉴别、有效治疗方法。

\end{pastquestion}

\begin{pastquestion}

\QuestionLabel{【14级职业病考题.pdf｜单选/填空/多选】}\\
尿铅的诊断值；铅中毒的早期表现；铅及其无机化合物吸收和排出；属于铅中毒影响的是。

\end{pastquestion}

\begin{pastquestion}

\QuestionLabel{【2015-2022职业病往年题及答案汇总.docx｜2022参考样题/单选10/多选7-8/填空6/病例】}\\
职业性慢性铅中毒血铅诊断值；驱铅治疗药物；接触铅的行业；铅及其无机化合物主要通过呼吸道、消化道吸收；诊断：职业性慢性轻度铅中毒。

\end{pastquestion}

\begin{pastquestion}

\QuestionLabel{【2015-2022职业病往年题及答案汇总.docx｜2018/填空】}\\
铅及其无机化合物可以通过呼吸道与消化道吸收，一般不经过皮肤吸收，大多以尿液方式排出体外，少量以粪便方式排出体外。

\end{pastquestion}

\begin{pastquestion}

\QuestionLabel{【2020年职业病考题.jpg｜选择6】}\\
职业性慢性铅中毒，尿铅的诊断值是：120μg/L。

\end{pastquestion}

\begin{pastquestion}

\QuestionLabel{【2021级职业病考题yys.docx｜单选/多选】}\\
最常见的造成血液系统损害的金属；铅损伤的系统。

\end{pastquestion}

\begin{pastquestion}

\QuestionLabel{【24年完整试题回忆.docx｜单选10/20/多选6/病例2】}\\
驱铅治疗使用；职业性铅中毒尿铅诊断值；有长期铅接触史，血铅和尿铅达到诊断值，符合哪些临床表现时诊断为中度中毒；电池厂电极板制造工20年腹绞痛病例，问诊断、依据、鉴别、辅助检查、治疗原则。

\end{pastquestion}

\PointSeparator

\section{慢性铅中毒治疗和预防}\label{users__fpb__26_spring__pkusph_study__tmp__pdfs__prevention-medicine-past-exams__07-ux804cux4e1aux75c5ux5b66ux5f80ux5e74ux9898ux603bux7ed3.md__ux6162ux6027ux94c5ux4e2dux6bd2ux6cbbux7597ux548cux9884ux9632}

\FrequencyPill{中频}

\InfoLabel{对应小节} 第三章 / 治疗；预防

\ContentBand{知识点原文摘取}

慢性中毒治疗原则：可使用CaNa2EDTA、二巯基丁二酸钠静脉滴注或二巯丁二酸（DMSA）口服驱铅；对症治疗，适当休息、治疗腹绞痛、贫血等。轻度和中度中毒驱铅治疗，重度中毒必须调离铅作业并进行驱铅和对症治疗。

预防：改善劳动条件，用无毒或低毒物质代替铅；加强个人防护；就业前体检和卫生防护教育；每年体检。

\ContentBand{对应往年题}

\begin{pastquestion}

\QuestionLabel{【04级-职业病考题.docx｜大题】}\\
急慢性铅中毒的治疗原则。

\end{pastquestion}

\begin{pastquestion}

\QuestionLabel{【05级职业病考试题.doc｜大题】}\\
慢性铅中毒如何治疗，如何预防。

\end{pastquestion}

\begin{pastquestion}

\QuestionLabel{【08级-职业病考题.doc｜病例】}\\
铅中毒病例：有效的治疗方法。

\end{pastquestion}

\begin{pastquestion}

\QuestionLabel{【24年完整试题回忆.docx｜病例2】}\\
铅中毒病例：简述治疗原则。

\end{pastquestion}

\PointSeparator

\chapter{农药中毒与有机磷农药中毒}\label{users__fpb__26_spring__pkusph_study__tmp__pdfs__prevention-medicine-past-exams__07-ux804cux4e1aux75c5ux5b66ux5f80ux5e74ux9898ux603bux7ed3.md__ux7b2cux56dbux7ae0-ux519cux836fux4e2dux6bd2ux4e0eux6709ux673aux78f7ux519cux836fux4e2dux6bd2}

\section{农药定义、分类和共同毒性}\label{users__fpb__26_spring__pkusph_study__tmp__pdfs__prevention-medicine-past-exams__07-ux804cux4e1aux75c5ux5b66ux5f80ux5e74ux9898ux603bux7ed3.md__ux519cux836fux5b9aux4e49ux5206ux7c7bux548cux5171ux540cux6bd2ux6027}

\FrequencyPill{高频}

\InfoLabel{对应小节} 第四章 / 一、农药；二、分类；三、常见农药的毒性特点

\InfoLabel{匹配依据} 07填空、2022样题填空、2021论述和24英译汉均考农药/农业化学品、共同毒性。

\ContentBand{知识点原文摘取}

农药：指用来杀灭或控制危害农作物生长和农产品贮存的各种有害生物（病毒、昆虫、杂草、鼠类、其他有害生物）的一类化合物；还包括植物生长调节剂、脱叶剂、增效剂等化学品。

按用途大致可分为：杀虫剂、杀螨剂、杀霉菌剂、除草剂、杀鼠剂、脱叶剂、不育剂、增效剂、植物生长调节剂、驱避剂。

农药共同毒性作用：皮肤粘膜刺激性；神经毒性；心脏毒性；消化系统刺激作用。

\ContentBand{对应往年题}

\begin{pastquestion}

\QuestionLabel{【07级-职业病考题.pdf｜填空】}\\
农药的一般毒性作用。

\end{pastquestion}

\begin{pastquestion}

\QuestionLabel{【2015-2022职业病往年题及答案汇总.docx｜2022参考样题/填空1】}\\
各类农药共同的毒性有：皮肤粘膜刺激性、神经毒性、心脏毒性、消化道刺激作用。

\end{pastquestion}

\begin{pastquestion}

\QuestionLabel{【2022职业病参考样题及答案整理.docx｜填空1】}\\
同上。

\end{pastquestion}

\begin{pastquestion}

\QuestionLabel{【2021级职业病考题yys.docx｜论述1】}\\
农药的定义及共同毒性。

\end{pastquestion}

\begin{pastquestion}

\QuestionLabel{【24年完整试题回忆.docx｜英译汉】}\\
agricultural chemicals。

\end{pastquestion}

\PointSeparator

\section{急性有机磷农药中毒机制、临床表现、分级和治疗}\label{users__fpb__26_spring__pkusph_study__tmp__pdfs__prevention-medicine-past-exams__07-ux804cux4e1aux75c5ux5b66ux5f80ux5e74ux9898ux603bux7ed3.md__ux6025ux6027ux6709ux673aux78f7ux519cux836fux4e2dux6bd2ux673aux5236ux4e34ux5e8aux8868ux73b0ux5206ux7ea7ux548cux6cbbux7597}

\FrequencyPill{高频}

\InfoLabel{对应小节} 第四章 / 六、有机磷农药中毒；五、急性农药中毒治疗要点

\InfoLabel{匹配依据} 04、05大题，08、14、2020、2021、24病例均考有机磷中毒机制、胆碱酯酶、临床三组症状、治疗。

\ContentBand{知识点原文摘取}

机制：有机磷中毒→胆碱酯酶失活→乙酰胆碱蓄积→胆碱能神经刺激→器官早期兴奋、晚期衰竭。

临床特征：M样/毒蕈碱样症状（腺体分泌增加、平滑肌收缩、括约肌松弛）；N样/烟碱样症状（肌束颤动、肌痉挛，晚期肌无力、肌麻痹，皮肤苍白、心率增快、血压升高）；中枢神经系统症状（烦躁焦虑、头晕头痛，晚期意识障碍、呼吸衰竭）。

分级：轻度中毒全血或红细胞胆碱酯酶活性一般在50\%-70\%；中度中毒一般在30\%-50\%；重度中毒一般在30\%以下，并可出现肺水肿、昏迷、呼吸衰竭、脑水肿。

特殊解毒剂：有机磷农药中毒采用抗胆碱药物-胆碱酯酶复活剂联合疗法。

\ContentBand{对应往年题}

\begin{pastquestion}

\QuestionLabel{【04级-职业病考题.docx｜大题/填空】}\\
有机磷中毒的并发症是什么，有机磷中毒的临床表现和治疗；有机磷中毒的定义；有机磷中毒胆碱酯抑制时的临床表现、分级、重度有机磷中毒时的并发症。

\end{pastquestion}

\begin{pastquestion}

\QuestionLabel{【05级职业病考试题.doc｜英译汉/大题】}\\
有机磷农药；有机磷中毒的发病机理及主要临床表现。

\end{pastquestion}

\begin{pastquestion}

\QuestionLabel{【07级-职业病考题.pdf｜大题/填空/多选/单选】}\\
有机磷迟发型周围性神经病发生时期及临床表现；乙酰胆碱酯酶活性下降的农药；有机磷中毒早期表现；迟发型周围神经发生时间；中间期肌无力综合征发生时间。

\end{pastquestion}

\begin{pastquestion}

\QuestionLabel{【08级-职业病考题.doc｜英译汉/多选/简答】}\\
有机磷农药中毒；急性有机磷的主要症状；重度有机磷中毒的并发症。

\end{pastquestion}

\begin{pastquestion}

\QuestionLabel{【14级职业病考题.pdf｜病例】}\\
一段很长的农药中毒病例：诊断及诊断依据；治疗原则。

\end{pastquestion}

\begin{pastquestion}

\QuestionLabel{【2015-2022职业病往年题及答案汇总.docx｜2022参考样题/单选4-5、7/多选5/填空2】}\\
有机磷引起的迟发性周围神经病仅见于恢复期；中间期肌无力综合征多发生在急性有机磷中毒后1-7天；中度有机磷中毒胆碱酯酶活性一般30\%-50\%；急性有机磷中毒主要表现为中枢神经症状、毒蕈碱样症状、烟碱样症状；能引起全血乙酰胆碱酯酶下降的两类农药是有机磷及氨基甲酸酯类。

\end{pastquestion}

\begin{pastquestion}

\QuestionLabel{【2015-2022职业病往年题及答案汇总.docx｜2020/简答与病例】}\\
急性有机磷中毒引起迟发性周围神经病的发生时间和临床表现；病例诊断：急性有机磷中毒，治疗原则：脱离接触、血液灌流、对症支持、特效解毒药（阿托品+氯磷定）。

\end{pastquestion}

\begin{pastquestion}

\QuestionLabel{【2020年职业病考题.jpg｜填空2/简答1/病例2】}\\
短时间大量接触有机磷农药后会引起体内\_\_，临床出现\_\_、\_\_和\_\_三组症状；简述急性有机磷中毒迟发性周围神经病的发生时间及临床表现；女62岁从事农药工作约3小时后气急、头晕、多汗、瞳孔针尖样、肌束颤动、胆碱酯酶活力25\%，问诊断和治疗原则。

\end{pastquestion}

\begin{pastquestion}

\QuestionLabel{【2021级职业病考题yys.docx｜病例1】}\\
妇女农药喷洒工作3h，血中乙酰胆碱酯酶活性25\%；问最可能诊断、依据、治疗措施。

\end{pastquestion}

\begin{pastquestion}

\QuestionLabel{【24年完整试题回忆.docx｜病例2】}\\
发病当日在田间从事农药工作约3小时，胆碱酯酶活力25\%；问诊断和治疗原则。

\end{pastquestion}

\PointSeparator

\section{氨基甲酸酯、有机氟、百草枯等其他农药}\label{users__fpb__26_spring__pkusph_study__tmp__pdfs__prevention-medicine-past-exams__07-ux804cux4e1aux75c5ux5b66ux5f80ux5e74ux9898ux603bux7ed3.md__ux6c28ux57faux7532ux9178ux916fux6709ux673aux6c1fux767eux8349ux67afux7b49ux5176ux4ed6ux519cux836f}

\FrequencyPill{中频}

\InfoLabel{对应小节} 第四章 / 急性农药中毒诊断步骤；治疗要点

\ContentBand{知识点原文摘取}

AchE活力下降提示有机磷、氨基甲酸酯类；活力正常提示其他农药。氨基甲酸酯类农药中毒：阿托品类抗胆碱药为特效解毒剂，无需使用胆碱酯酶复活剂。有机氟类农药中毒：乙酰胺（解氟灵）是特效解毒剂。百草枯中毒关键在于防治氧化性损伤引起的急性肺间质纤维化。

\ContentBand{对应往年题}

\begin{pastquestion}

\QuestionLabel{【08级-职业病考题.doc｜单选】}\\
百草枯的主要损害。

\end{pastquestion}

\begin{pastquestion}

\QuestionLabel{【14级职业病考题.pdf｜单选】}\\
有机氟农药中毒的治疗可以使用；引起急性肺部纤维化的农药是。

\end{pastquestion}

\begin{pastquestion}

\QuestionLabel{【2015-2022职业病往年题及答案汇总.docx｜2022参考样题/单选12】}\\
氨基甲酸酯类农药的特效解毒剂为阿托品。

\end{pastquestion}

\begin{pastquestion}

\QuestionLabel{【24年完整试题回忆.docx｜单选6/7/11】}\\
氨基甲酸酯类农药中毒的特效解毒剂；百草枯的中毒机制；什么情况下考虑用亚甲蓝治疗。

\end{pastquestion}

\PointSeparator

\chapter{汞中毒}\label{users__fpb__26_spring__pkusph_study__tmp__pdfs__prevention-medicine-past-exams__07-ux804cux4e1aux75c5ux5b66ux5f80ux5e74ux9898ux603bux7ed3.md__ux7b2cux4e94ux7ae0-ux6c5eux4e2dux6bd2}

\section{汞接触机会、吸收代谢和毒性机制}\label{users__fpb__26_spring__pkusph_study__tmp__pdfs__prevention-medicine-past-exams__07-ux804cux4e1aux75c5ux5b66ux5f80ux5e74ux9898ux603bux7ed3.md__ux6c5eux63a5ux89e6ux673aux4f1aux5438ux6536ux4ee3ux8c22ux548cux6bd2ux6027ux673aux5236}

\FrequencyPill{中高频}

\InfoLabel{对应小节} 第五章 / 接触机会；汞代谢中毒机制

\ContentBand{知识点原文摘取}

汞矿开采及冶炼、氯碱行业、乙炔法生产氯乙烯、仪表行业（温度计、血压计等）、电气行业、冶金工业、国防工业等均可接触汞。

金属汞主要以蒸气形式经呼吸道吸收入体内，吸收量达75\%-100\%；消化道对金属汞吸收量很少；皮肤不易吸收汞。体内汞大部分由尿排出，小部分经肺、皮肤、粪便、唾液、胆汁、乳汁排出。汞主要以Hg2+形式发挥毒性。

\ContentBand{对应往年题}

\begin{pastquestion}

\QuestionLabel{【04级-职业病考题.docx｜填空】}\\
接触汞的行业。

\end{pastquestion}

\begin{pastquestion}

\QuestionLabel{【08级-职业病考题.doc｜填空/单选】}\\
汞及其化合物进入人体后以（）发挥毒性；汞的蓄积部位；尿汞的诊断值。

\end{pastquestion}

\begin{pastquestion}

\QuestionLabel{【2022职业病参考样题及答案整理.docx｜多选3】}\\
慢性轻度汞中毒具有尿汞增高、口腔-牙龈炎等表现。

\end{pastquestion}

\PointSeparator

\section{急性与慢性汞中毒临床表现、分级、鉴别和治疗}\label{users__fpb__26_spring__pkusph_study__tmp__pdfs__prevention-medicine-past-exams__07-ux804cux4e1aux75c5ux5b66ux5f80ux5e74ux9898ux603bux7ed3.md__ux6025ux6027ux4e0eux6162ux6027ux6c5eux4e2dux6bd2ux4e34ux5e8aux8868ux73b0ux5206ux7ea7ux9274ux522bux548cux6cbbux7597}

\FrequencyPill{高频}

\InfoLabel{对应小节} 第五章 / 临床表现；诊断；鉴别诊断；治疗

\InfoLabel{匹配依据}
04、05大题直接考急慢性汞中毒临床表现/鉴别/治疗，07、08、14、2022、24选择填空反复考急性/慢性汞中毒表现。

\ContentBand{知识点原文摘取}

急性汞中毒：多由短期吸入高浓度汞蒸气引起，靶器官按后果严重性以肾为主，其次为肺、消化道、皮肤、脑等。可有全身症状、支气管肺炎、口腔炎、急性肾小管坏死、皮炎等。

慢性汞中毒：长期接触较低浓度汞蒸气引起，主要表现为易兴奋征、震颤、口腔-牙龈炎三大特征。慢性轻度中毒表现包括神经衰弱综合征、口腔-牙龈炎、手指震颤、近端肾小管功能障碍、尿汞增加；中度中毒可有性格情绪改变、上肢粗大震颤、明显肾脏损害；重度为慢性中毒性脑病。

治疗：急性汞中毒迅速脱离现场，使用巯基螯合剂驱汞，如二巯丙磺酸钠或二巯丁二酸钠；慢性汞中毒可用二巯丙磺酸钠、二巯丁二酸钠或口服二巯丁二酸，3日为一疗程，停药4天后开始下一疗程。

\ContentBand{对应往年题}

\begin{pastquestion}

\QuestionLabel{【04级-职业病考题.docx｜大题】}\\
急慢性汞中毒的临床表现。

\end{pastquestion}

\begin{pastquestion}

\QuestionLabel{【05级职业病考试题.doc｜英译汉/大题】}\\
慢性汞中毒；急慢性汞中毒的鉴别诊断，治疗方法。

\end{pastquestion}

\begin{pastquestion}

\QuestionLabel{【07级-职业病考题.pdf｜填空/多选/单选】}\\
汞中度中毒的标准；汞苯中毒的主要表现、分级表现（特征性的点）；慢性汞中毒肾损害的部位；驱汞的制剂。

\end{pastquestion}

\begin{pastquestion}

\QuestionLabel{【08级-职业病考题.doc｜填空/单选/多选】}\\
汞的蓄积部位；尿汞诊断值；慢性中度汞中毒表现。

\end{pastquestion}

\begin{pastquestion}

\QuestionLabel{【14级职业病考题.pdf｜单选/填空/多选】}\\
驱汞治疗首选；不属于急性轻度汞中毒的表现；慢性汞中毒三大表现；属于慢性中度汞中毒的表现。

\end{pastquestion}

\begin{pastquestion}

\QuestionLabel{【2015-2022职业病往年题及答案汇总.docx｜2022参考样题/单选11/多选3】}\\
下列哪项不是急性汞中毒的常见临床表现；慢性轻度汞中毒具有尿汞增高、口腔-牙龈炎等表现。

\end{pastquestion}

\begin{pastquestion}

\QuestionLabel{【2015-2022职业病往年题及答案汇总.docx｜2018/单选、填空】}\\
驱汞治疗首选二巯丙磺酸钠、二巯丁二酸钠；慢性中度汞中毒：间质性肺炎、明显蛋白尿；慢性汞中毒三大表现：易兴奋征、震颤、口腔牙龈炎。

\end{pastquestion}

\begin{pastquestion}

\QuestionLabel{【2015-2022职业病往年题及答案汇总.docx｜2020/填空】}\\
急性中度汞中毒是在轻度中毒基础上，具有间质性肺炎、明显蛋白尿之一者。

\end{pastquestion}

\begin{pastquestion}

\QuestionLabel{【2020年职业病考题.jpg｜填空7】}\\
急性中度汞中毒是在轻度中毒的基础上，具有\_\_、\_\_者。

\end{pastquestion}

\begin{pastquestion}

\QuestionLabel{【24年完整试题回忆.docx｜单选1】}\\
下列哪项不属于急性汞中毒的临床表现。

\end{pastquestion}

\PointSeparator

\chapter{砷中毒与砷化氢中毒}\label{users__fpb__26_spring__pkusph_study__tmp__pdfs__prevention-medicine-past-exams__07-ux804cux4e1aux75c5ux5b66ux5f80ux5e74ux9898ux603bux7ed3.md__ux7b2cux516dux7ae0-ux7837ux4e2dux6bd2ux4e0eux7837ux5316ux6c22ux4e2dux6bd2}

\section{急性砷中毒、慢性砷中毒与皮肤表现}\label{users__fpb__26_spring__pkusph_study__tmp__pdfs__prevention-medicine-past-exams__07-ux804cux4e1aux75c5ux5b66ux5f80ux5e74ux9898ux603bux7ed3.md__ux6025ux6027ux7837ux4e2dux6bd2ux6162ux6027ux7837ux4e2dux6bd2ux4e0eux76aeux80a4ux8868ux73b0}

\FrequencyPill{高频}

\InfoLabel{对应小节} 第六章 / 临床表现；诊断；治疗

\InfoLabel{匹配依据} 04填空、08单选/多选、2022样题单选/多选、24单选均考急性砷中毒、慢性砷中毒皮肤表现。

\ContentBand{知识点原文摘取}

急性砷中毒：职业性急性中毒少见，可由砷化物烟雾、粉尘经皮肤粘膜、呼吸道吸收引起；非职业性多经消化道吸收。经口服中毒早期可有急性胃肠炎，神经系统损害，中毒性心肌炎，肝肾损伤等。

慢性砷中毒皮肤损害：色素沉着（多见于非暴露部位，雨点状或花斑状）；皮肤角化过度（手掌、足底为著，典型为手掌尺侧外缘、手指根部谷粒样隆起）；疣状增生。特点为色素沉着与色素脱失可同时存在，皮肤色素沉着、角化过度及疣状增生三者可并存。

慢性砷中毒诊断原则：根据长期砷化物职业接触史，出现皮肤、神经系统及肝脏损伤为主的临床表现，结合尿砷、发砷等实验室检查结果，参考现场职业卫生调查综合分析，排除其他类似疾病后诊断。

\ContentBand{对应往年题}

\begin{pastquestion}

\QuestionLabel{【04级-职业病考题.docx｜填空】}\\
砷中毒皮肤表现（3个空）。

\end{pastquestion}

\begin{pastquestion}

\QuestionLabel{【07级-职业病考题.pdf｜翻译】}\\
砷中毒。

\end{pastquestion}

\begin{pastquestion}

\QuestionLabel{【08级-职业病考题.doc｜单选/多选】}\\
砷化合物毒性最强的是；口服导致急性砷中毒的临床表现。

\end{pastquestion}

\begin{pastquestion}

\QuestionLabel{【2015-2022职业病往年题及答案汇总.docx｜2022参考样题/单选2/多选12/翻译】}\\
急性砷中毒最突出的早期表现是急性胃肠炎；职业性慢性砷中毒表现包括皮肤过度角化、周围神经病、皮肤癌、中毒性肝病；砷中毒：arsenic
poisoning。

\end{pastquestion}

\begin{pastquestion}

\QuestionLabel{【2022职业病参考样题及答案整理.docx｜翻译/单选2/多选12】}\\
同上。

\end{pastquestion}

\begin{pastquestion}

\QuestionLabel{【24年完整试题回忆.docx｜单选3】}\\
色素沉着、角化过度、疣状增生可见于（）中毒。

\end{pastquestion}

\PointSeparator

\section{砷化氢中毒机制、临床特点和治疗}\label{users__fpb__26_spring__pkusph_study__tmp__pdfs__prevention-medicine-past-exams__07-ux804cux4e1aux75c5ux5b66ux5f80ux5e74ux9898ux603bux7ed3.md__ux7837ux5316ux6c22ux4e2dux6bd2ux673aux5236ux4e34ux5e8aux7279ux70b9ux548cux6cbbux7597}

\FrequencyPill{中频}

\InfoLabel{对应小节} 第六章 / 十二、砷化氢中毒

\ContentBand{知识点原文摘取}

砷化氢为气态毒物，主要经呼吸道侵入体内，可迅速吸收入血。绝大部分进入红细胞与氧合血红蛋白结合，迅速引起溶血。毒性机制：消耗还原型谷胱甘肽，引起急性血管内溶血；溶血产生的游离血红蛋白经肾小球滤出，在肾小管形成管型，引起急性肾小管坏死，严重者急性肾功能衰竭。

临床表现：急性血管内溶血、急性肾损伤、急性肝损伤。治疗重点包括保护肾脏功能、控制和减轻溶血反应、血液净化疗法、对症支持治疗；一般不主张使用金属络合剂。

\ContentBand{对应往年题}

\begin{pastquestion}

\QuestionLabel{【14级职业病考题.pdf｜单选】}\\
砷化氢的中毒机制。

\end{pastquestion}

\begin{pastquestion}

\QuestionLabel{【2015-2022职业病往年题及答案汇总.docx｜2018/单选】}\\
砷化氢的中毒机制：消耗还原型谷胱甘肽，引起急性血管内溶血。

\end{pastquestion}

\begin{pastquestion}

\QuestionLabel{【2015-2022职业病往年题及答案汇总.docx｜2020/选择】}\\
根据砷化氢的主要毒性，它属于：溶血性毒物。

\end{pastquestion}

\begin{pastquestion}

\QuestionLabel{【2020年职业病考题.jpg｜选择10】}\\
根据砷化氢的主要毒性，它属于：溶血性毒物（原图``我的答案''显示D，与知识点和汇总答案冲突，以知识点总结/汇总为准）。

\end{pastquestion}

\begin{pastquestion}

\QuestionLabel{【24年完整试题回忆.docx｜单选4】}\\
砷化氢按照毒性机制属于溶血性毒物。

\end{pastquestion}

\PointSeparator

\chapter{锰中毒}\label{users__fpb__26_spring__pkusph_study__tmp__pdfs__prevention-medicine-past-exams__07-ux804cux4e1aux75c5ux5b66ux5f80ux5e74ux9898ux603bux7ed3.md__ux7b2cux4e03ux7ae0-ux9530ux4e2dux6bd2}

\section{慢性锰中毒临床表现、代谢途径、诊断和治疗}\label{users__fpb__26_spring__pkusph_study__tmp__pdfs__prevention-medicine-past-exams__07-ux804cux4e1aux75c5ux5b66ux5f80ux5e74ux9898ux603bux7ed3.md__ux6162ux6027ux9530ux4e2dux6bd2ux4e34ux5e8aux8868ux73b0ux4ee3ux8c22ux9014ux5f84ux8bcaux65adux548cux6cbbux7597}

\FrequencyPill{中频}

\InfoLabel{对应小节} 第七章 / 临床表现；诊断；课堂小题

\ContentBand{知识点原文摘取}

慢性锰中毒起病缓慢，多在接触10年后或更长时间。典型表现为中枢神经系统锥体外系受损害引起的震颤麻痹综合征。早期为非特异主观症状：头痛、无力、淡漠、社交活动减少、抑郁或兴奋。诊断主要依据：职业接触史、典型临床表现。锰在慢性锰中毒病人的代谢途径主要是：经呼吸道吸收，从粪便排出。最重要治疗措施是早期发现，及时脱离接触。

\ContentBand{对应往年题}

\begin{pastquestion}

\QuestionLabel{【2015-2022职业病往年题及答案汇总.docx｜2022参考样题/单选1】}\\
慢性锰中毒的早期临床表现是头痛、无力、淡漠等非特异性主观症状。

\end{pastquestion}

\begin{pastquestion}

\QuestionLabel{【2022职业病参考样题及答案整理.docx｜单选1】}\\
同上。

\end{pastquestion}

\begin{pastquestion}

\QuestionLabel{【2020年职业病考题.jpg｜选择4】}\\
慢性锰中毒时，锰在人体的吸收代谢途径主要是：经呼吸道吸收，从粪便排出（原图学生答案显示A``从尿液排出''，与知识点总结课堂小题冲突，以知识点为准）。

\end{pastquestion}

\PointSeparator

\chapter{刺激性气体中毒}\label{users__fpb__26_spring__pkusph_study__tmp__pdfs__prevention-medicine-past-exams__07-ux804cux4e1aux75c5ux5b66ux5f80ux5e74ux9898ux603bux7ed3.md__ux7b2cux516bux7ae0-ux523aux6fc0ux6027ux6c14ux4f53ux4e2dux6bd2}

\section{刺激性气体定义、种类和英译汉}\label{users__fpb__26_spring__pkusph_study__tmp__pdfs__prevention-medicine-past-exams__07-ux804cux4e1aux75c5ux5b66ux5f80ux5e74ux9898ux603bux7ed3.md__ux523aux6fc0ux6027ux6c14ux4f53ux5b9aux4e49ux79cdux7c7bux548cux82f1ux8bd1ux6c49}

\FrequencyPill{高频}

\InfoLabel{对应小节} 第八章 / 刺激性气体概念；种类

\ContentBand{知识点原文摘取}

刺激性气体：主要是指那些由于本身理化特性，而对呼吸道及肺泡上皮有直接刺激作用的气态化合物。最主要病损是呼吸系统改变，轻者为上呼吸道刺激症状，重者喉头水肿、喉痉挛、支气管炎及急性肺损伤，严重时可发生肺炎、肺水肿，甚至ARDS。

常见刺激性气体：氨气、光气、双光气、二氧化硫、氮氧化物、甲醛、硫酸二甲酯、氯化氢、氟化氢、溴化氢、臭氧等。

\ContentBand{对应往年题}

\begin{pastquestion}

\QuestionLabel{【05级职业病考试题.doc｜英译汉】}\\
刺激性气体。

\end{pastquestion}

\begin{pastquestion}

\QuestionLabel{【07级-职业病考题.pdf｜翻译/多选/单选】}\\
刺激性气体；以下哪些是刺激性气体；哪个不是刺激性气体：硫化氢。

\end{pastquestion}

\begin{pastquestion}

\QuestionLabel{【08级-职业病考题.doc｜多选】}\\
刺激性气体的主要并发症。

\end{pastquestion}

\begin{pastquestion}

\QuestionLabel{【14级职业病考题.pdf｜英译汉/多选】}\\
刺激性气体；刺激性气体的并发症包括。

\end{pastquestion}

\begin{pastquestion}

\QuestionLabel{【2015-2022职业病往年题及答案汇总.docx｜2018/翻译】}\\
irritant gas：刺激性气体。

\end{pastquestion}

\begin{pastquestion}

\QuestionLabel{【2015-2022职业病往年题及答案汇总.docx｜2022参考样题/多选6】}\\
下面的哪些气体是刺激性气体？甲醛、氯气、光气。

\end{pastquestion}

\begin{pastquestion}

\QuestionLabel{【2020年职业病考题.jpg｜选择7】}\\
以下不是刺激性气体的是：一氧化碳。

\end{pastquestion}

\begin{pastquestion}

\QuestionLabel{【24年完整试题回忆.docx｜多选8】}\\
刺激性气体包括氨、氯气、光气等。

\end{pastquestion}

\PointSeparator

\section{化学性肺水肿分期、并发症和处理原则}\label{users__fpb__26_spring__pkusph_study__tmp__pdfs__prevention-medicine-past-exams__07-ux804cux4e1aux75c5ux5b66ux5f80ux5e74ux9898ux603bux7ed3.md__ux5316ux5b66ux6027ux80baux6c34ux80bfux5206ux671fux5e76ux53d1ux75c7ux548cux5904ux7406ux539fux5219}

\FrequencyPill{高频}

\InfoLabel{对应小节} 第八章 / 临床表现；处理原则

\ContentBand{知识点原文摘取}

化学性肺水肿分期：刺激期（肺水肿前期）、潜伏期、肺水肿期、恢复期。

并发症包括：气道粘膜坏死脱落、肺不张、纵隔气肿和自发性气胸、肺部感染、肺纤维化、阻塞性细支气管炎、中毒性心肌损害和休克、ARDS等。

处理原则：刺激性气体的主要危害是肺水肿，积极防治肺水肿是中毒抢救治疗的关键。迅速脱离现场，处理灼伤，预防肺水肿；肺水肿处理采取通气、利水、抗感染等措施。

\ContentBand{对应往年题}

\begin{pastquestion}

\QuestionLabel{【07级-职业病考题.pdf｜填空】}\\
化学性肺水肿分期（书为准）。

\end{pastquestion}

\begin{pastquestion}

\QuestionLabel{【08级-职业病考题.doc｜多选】}\\
刺激性气体的主要并发症。

\end{pastquestion}

\begin{pastquestion}

\QuestionLabel{【14级职业病考题.pdf｜填空/多选】}\\
化学性肺水肿的四期；刺激性气体的并发症包括。

\end{pastquestion}

\begin{pastquestion}

\QuestionLabel{【2015-2022职业病往年题及答案汇总.docx｜2022参考样题/多选11】}\\
刺激性气体中毒的常见并发症包括窒息、纵隔气肿和自发性气胸、肺纤维化、中毒性心肌损伤。

\end{pastquestion}

\begin{pastquestion}

\QuestionLabel{【2015-2022职业病往年题及答案汇总.docx｜2020/填空】}\\
化学性肺水肿临床分为四期：刺激期、潜伏期、肺水肿期、恢复期。

\end{pastquestion}

\begin{pastquestion}

\QuestionLabel{【2020年职业病考题.jpg｜填空3】}\\
化学性肺水肿临床分为四期，即\_\_、\textbf{、}、\_\_。

\end{pastquestion}

\begin{pastquestion}

\QuestionLabel{【24年完整试题回忆.docx｜填空7】}\\
化学性肺水肿的临床分期。

\end{pastquestion}

\begin{pastquestion}

\QuestionLabel{【2021级职业病考题yys.docx｜英译中】}\\
Acute respiratory distress syndrome。

\end{pastquestion}

\begin{pastquestion}

\QuestionLabel{【24年完整试题回忆.docx｜英译汉】}\\
latent phase（可对应``潜伏期''；在化学性肺水肿分期和急性放射病分期中均出现）。

\end{pastquestion}

\PointSeparator

\chapter{急性窒息性气体中毒与一氧化碳中毒}\label{users__fpb__26_spring__pkusph_study__tmp__pdfs__prevention-medicine-past-exams__07-ux804cux4e1aux75c5ux5b66ux5f80ux5e74ux9898ux603bux7ed3.md__ux7b2cux4e5dux7ae0-ux6025ux6027ux7a92ux606fux6027ux6c14ux4f53ux4e2dux6bd2ux4e0eux4e00ux6c27ux5316ux78b3ux4e2dux6bd2}

\section{窒息性气体定义、分类和举例}\label{users__fpb__26_spring__pkusph_study__tmp__pdfs__prevention-medicine-past-exams__07-ux804cux4e1aux75c5ux5b66ux5f80ux5e74ux9898ux603bux7ed3.md__ux7a92ux606fux6027ux6c14ux4f53ux5b9aux4e49ux5206ux7c7bux548cux4e3eux4f8b}

\FrequencyPill{高频}

\InfoLabel{对应小节} 第九章 / 窒息性气体分类

\ContentBand{知识点原文摘取}

窒息性气体分为单纯性窒息性气体和化学性窒息性气体。单纯性窒息性气体如氮气、甲烷、水蒸气等。化学性窒息性气体包括血液性窒息性气体（CO、苯胺、硝基苯等）和细胞性窒息性气体（HCN、H2S）。

\ContentBand{对应往年题}

\begin{pastquestion}

\QuestionLabel{【08级-职业病考题.doc｜英译汉/简答】}\\
窒息性气体；窒息性气体的定义、分类及举例说明。

\end{pastquestion}

\begin{pastquestion}

\QuestionLabel{【07级-职业病考题.pdf｜大题】}\\
窒息性气体分三类，各举一例。

\end{pastquestion}

\begin{pastquestion}

\QuestionLabel{【14级职业病考题.pdf｜简答/多选】}\\
窒息性气体的定义，分类，每类各举一例；属于窒息性气体的是。

\end{pastquestion}

\begin{pastquestion}

\QuestionLabel{【2015年职业病考试题型.docx｜填空举例】}\\
窒息性气体分几类。

\end{pastquestion}

\begin{pastquestion}

\QuestionLabel{【2015-2022职业病往年题及答案汇总.docx｜2015/填空】}\\
窒息性气体分为：单纯性、化学性（血液性、细胞性）窒息性气体。

\end{pastquestion}

\begin{pastquestion}

\QuestionLabel{【2015-2022职业病往年题及答案汇总.docx｜2018/简答】}\\
窒息性气体的定义，分类，每类各举一例。

\end{pastquestion}

\begin{pastquestion}

\QuestionLabel{【2015-2022职业病往年题及答案汇总.docx｜2022参考样题/单选17】}\\
以下不是窒息性气体的是氯气。

\end{pastquestion}

\begin{pastquestion}

\QuestionLabel{【2021级职业病考题yys.docx｜英译中/论述2】}\\
Asphyxiating gases；窒息性气体的定义及分类，各举一例。

\end{pastquestion}

\begin{pastquestion}

\QuestionLabel{【24年完整试题回忆.docx｜单选12/简答1】}\\
哪项不属于单纯性窒息性气体；窒息性气体是什么？按中毒机理分哪几类？各举一例。

\end{pastquestion}

\PointSeparator

\section{一氧化碳中毒诊断依据、治疗原则和相关选择题}\label{users__fpb__26_spring__pkusph_study__tmp__pdfs__prevention-medicine-past-exams__07-ux804cux4e1aux75c5ux5b66ux5f80ux5e74ux9898ux603bux7ed3.md__ux4e00ux6c27ux5316ux78b3ux4e2dux6bd2ux8bcaux65adux4f9dux636eux6cbbux7597ux539fux5219ux548cux76f8ux5173ux9009ux62e9ux9898}

\FrequencyPill{中高频}

\InfoLabel{对应小节} 第九章 / 一氧化碳；病理生理

\ContentBand{知识点原文摘取}

CO与红细胞中血红蛋白结合形成碳氧血红蛋白，使其失去携带氧气的能力。CO与Hb的亲和力约为O2的200多倍。CO中毒核心是造成机体缺氧：影响血液氧气运输，并使氧离曲线左移，影响外周组织摄氧。CO主要经肺原形排出。

\ContentBand{对应往年题}

\begin{pastquestion}

\QuestionLabel{【04级-职业病考题.docx｜大题】}\\
急性一氧化碳中毒的诊断标准和治疗原则；CO的诊断依据和治疗原则。

\end{pastquestion}

\begin{pastquestion}

\QuestionLabel{【05级职业病考试题.doc｜英译汉/大题】}\\
一氧化碳；CO中毒的诊断要点。

\end{pastquestion}

\begin{pastquestion}

\QuestionLabel{【08级-职业病考题.doc｜单选】}\\
CO与血红蛋白的亲和力比氧气与它的亲和力高（）倍。

\end{pastquestion}

\begin{pastquestion}

\QuestionLabel{【2015-2022职业病往年题及答案汇总.docx｜2022参考样题/单选14】}\\
一氧化碳进入机体血液中，可与血红蛋白中的二价铁结合生成碳氧血红蛋白。

\end{pastquestion}

\begin{pastquestion}

\QuestionLabel{【24年完整试题回忆.docx｜单选13】}\\
一氧化碳与血红蛋白结合形成。

\end{pastquestion}

\PointSeparator

\section{氰化物/氰化氢排毒途径}\label{users__fpb__26_spring__pkusph_study__tmp__pdfs__prevention-medicine-past-exams__07-ux804cux4e1aux75c5ux5b66ux5f80ux5e74ux9898ux603bux7ed3.md__ux6c30ux5316ux7269ux6c30ux5316ux6c22ux6392ux6bd2ux9014ux5f84}

\FrequencyPill{中频}

\InfoLabel{对应小节} 第九章 / 窒息性气体分类；职业病治疗中特殊解毒剂

\ContentBand{知识点原文摘取}

氰化氢属于细胞性窒息性气体。题源答案整理：氰化物在机体内最主要的解毒排除途径是在硫氰酸盐生成酶作用下，与硫生成低毒的硫氰酸盐，经尿排出。

\ContentBand{对应往年题}

\begin{pastquestion}

\QuestionLabel{【07级-职业病考题.pdf｜单选】}\\
氰化氢的主要排出途径。

\end{pastquestion}

\begin{pastquestion}

\QuestionLabel{【08级-职业病考题.doc｜单选】}\\
氰化氢的主要排毒途径。

\end{pastquestion}

\begin{pastquestion}

\QuestionLabel{【14级职业病考题.pdf｜单选】}\\
氰化氢主要排出途径。

\end{pastquestion}

\begin{pastquestion}

\QuestionLabel{【2015-2022职业病往年题及答案汇总.docx｜2022参考样题/单选15】}\\
氰化物在机体内最主要的解毒排除途径。

\end{pastquestion}

\begin{pastquestion}

\QuestionLabel{【24年完整试题回忆.docx｜单选14】}\\
氰化物的主要排除机制。

\end{pastquestion}

\PointSeparator

\chapter{有机溶剂、正己烷和三氯乙烯中毒}\label{users__fpb__26_spring__pkusph_study__tmp__pdfs__prevention-medicine-past-exams__07-ux804cux4e1aux75c5ux5b66ux5f80ux5e74ux9898ux603bux7ed3.md__ux7b2cux5341ux7ae0-ux6709ux673aux6eb6ux5242ux6b63ux5df1ux70f7ux548cux4e09ux6c2fux4e59ux70efux4e2dux6bd2}

\section{有机溶剂定义、特点、用途和共同毒性}\label{users__fpb__26_spring__pkusph_study__tmp__pdfs__prevention-medicine-past-exams__07-ux804cux4e1aux75c5ux5b66ux5f80ux5e74ux9898ux603bux7ed3.md__ux6709ux673aux6eb6ux5242ux5b9aux4e49ux7279ux70b9ux7528ux9014ux548cux5171ux540cux6bd2ux6027}

\FrequencyPill{高频}

\InfoLabel{对应小节} 第十章 / 有机溶剂概述

\ContentBand{知识点原文摘取}

有机溶剂定义：可溶解不溶于水的某些有机物质（如油脂、蜡、树脂、橡胶、染料等）的液体。有机溶剂本身也是有机化合物，分子量多不大，常温下呈液态；发生溶解过程后溶质和溶剂性质均无改变；具有较大挥发性，不少化合物易燃、易爆；用途广泛。

题源常考共同毒性：脂溶性强，具有刺激性，具有麻醉性；``有特异性解毒剂''通常为错误项。

\ContentBand{对应往年题}

\begin{pastquestion}

\QuestionLabel{【04级-职业病考题.docx｜大题】}\\
什么叫有机溶剂，它有什么用途，正己烷和三氯乙烯的毒性作用各是什么；有机溶剂定义、用途。

\end{pastquestion}

\begin{pastquestion}

\QuestionLabel{【05级职业病考试题.doc｜大题】}\\
有机溶剂的定义及分类，共同毒性包括什么。

\end{pastquestion}

\begin{pastquestion}

\QuestionLabel{【07级-职业病考题.pdf｜翻译】}\\
有机溶剂。

\end{pastquestion}

\begin{pastquestion}

\QuestionLabel{【08级-职业病考题.doc｜填空/单选】}\\
有机溶剂的定义；有机溶剂共同毒性；关于有机溶剂说法错误的是（有机溶剂中毒都有特效解毒剂）。

\end{pastquestion}

\begin{pastquestion}

\QuestionLabel{【14级职业病考题.pdf｜英译汉/单选】}\\
有机溶剂中毒；关于有机溶剂错误的是。

\end{pastquestion}

\begin{pastquestion}

\QuestionLabel{【2015-2022职业病往年题及答案汇总.docx｜2022参考样题/翻译/单选9】}\\
organic solvents：有机溶剂；关于有机溶剂以下说法不正确的是``有机溶剂中毒有特异性的解毒剂''。

\end{pastquestion}

\begin{pastquestion}

\QuestionLabel{【2020年职业病考题.jpg｜选择5】}\\
关于有机溶剂以下说法不正确的是：有特异性的解毒剂。

\end{pastquestion}

\PointSeparator

\section{正己烷慢性中毒与尿2,5-己二酮}\label{users__fpb__26_spring__pkusph_study__tmp__pdfs__prevention-medicine-past-exams__07-ux804cux4e1aux75c5ux5b66ux5f80ux5e74ux9898ux603bux7ed3.md__ux6b63ux5df1ux70f7ux6162ux6027ux4e2dux6bd2ux4e0eux5c3f25-ux5df1ux4e8cux916e}

\FrequencyPill{高频}

\InfoLabel{对应小节} 第十章 / 正己烷中毒

\ContentBand{知识点原文摘取}

正己烷慢性中毒主要表现为运动障碍为主的混合型周围神经病，可伴植物神经功能紊乱；多于接触数月后发病，脱离接触3-6个月病情仍可进展。四肢远端痛触觉减退或消失，呈手套、袜套样对称分布；可有垂腕、垂足，下肢为甚。尿2,5-己二酮测定可作为近期接触指标，脱离正己烷接触一个月左右即呈阴性。

\ContentBand{对应往年题}

\begin{pastquestion}

\QuestionLabel{【04级-职业病考题.docx｜大题】}\\
正己烷和三氯乙烯的毒性作用各是什么；列举正己烷慢性中毒的临床特点。

\end{pastquestion}

\begin{pastquestion}

\QuestionLabel{【07级-职业病考题.pdf｜单选】}\\
慢性正己烷中毒主要是周围神经病；能造成血液系统损害的有机溶剂。

\end{pastquestion}

\begin{pastquestion}

\QuestionLabel{【08级-职业病考题.doc｜单选】}\\
职业性正己烷中毒的主要途径。

\end{pastquestion}

\begin{pastquestion}

\QuestionLabel{【14级职业病考题.pdf｜单选】}\\
尿中哪个指标可以反映正己烷的近期接触；正己烷的慢性毒性主要是。

\end{pastquestion}

\begin{pastquestion}

\QuestionLabel{【2015-2022职业病往年题及答案汇总.docx｜2022参考样题/单选3】}\\
慢性正己烷中毒的主要表现是周围神经病。

\end{pastquestion}

\begin{pastquestion}

\QuestionLabel{【2015-2022职业病往年题及答案汇总.docx｜2018/单选】}\\
正己烷的慢性毒性主要是运动障碍为主的混合型周围神经病；尿中2,5-己二酮可以反映正己烷近期接触。

\end{pastquestion}

\begin{pastquestion}

\QuestionLabel{【24年完整试题回忆.docx｜单选2】}\\
哪项指标可以反映有过量正己烷接触：2,5-己二酮。

\end{pastquestion}

\PointSeparator

\section{三氯乙烯急性中毒与药疹样皮炎}\label{users__fpb__26_spring__pkusph_study__tmp__pdfs__prevention-medicine-past-exams__07-ux804cux4e1aux75c5ux5b66ux5f80ux5e74ux9898ux603bux7ed3.md__ux4e09ux6c2fux4e59ux70efux6025ux6027ux4e2dux6bd2ux4e0eux836fux75b9ux6837ux76aeux708e}

\FrequencyPill{中高频}

\InfoLabel{对应小节} 第十章 / 三氯乙烯中毒

\ContentBand{知识点原文摘取}

三氯乙烯急性中毒表现为中枢神经系统症状，肝、肾及心脏损害，颅神经症状。三氯乙烯还可引起药疹样皮炎，发病无明显剂量-效应关系，具有一定潜伏期，常以发热或皮疹为首发症状，皮肤损害可分为剥脱性皮炎、多形红斑、重症多形红斑、大疱性表皮坏死松解症。我国颁布职业性三氯乙烯药疹样皮炎诊断标准。

\ContentBand{对应往年题}

\begin{pastquestion}

\QuestionLabel{【04级-职业病考题.docx｜大题】}\\
三氯乙烯的毒性作用。

\end{pastquestion}

\begin{pastquestion}

\QuestionLabel{【08级-职业病考题.doc｜多选】}\\
三氯乙烯的职业损害。

\end{pastquestion}

\begin{pastquestion}

\QuestionLabel{【14级职业病考题.pdf｜多选】}\\
三氯乙烯的临床表现包括。

\end{pastquestion}

\begin{pastquestion}

\QuestionLabel{【2015-2022职业病往年题及答案汇总.docx｜2022参考样题/多选4】}\\
三氯乙烯的职业危害包括急性中毒、药疹样皮炎。

\end{pastquestion}

\begin{pastquestion}

\QuestionLabel{【2015-2022职业病往年题及答案汇总.docx｜2020/填空】}\\
三氯乙烯引起两种职业病：职业性急性三氯乙烯中毒、职业性三氯乙烯药疹样皮炎。

\end{pastquestion}

\begin{pastquestion}

\QuestionLabel{【2020年职业病考题.jpg｜填空6】}\\
三氯乙烯可以引起两种职业病，分别是\_\_和\_\_（知识点答案为职业性急性三氯乙烯中毒、职业性三氯乙烯药疹样皮炎；原图学生答案明显不可靠）。

\end{pastquestion}

\begin{pastquestion}

\QuestionLabel{【24年完整试题回忆.docx｜多选3】}\\
三氯乙烯的职业病诊断标准包括急性中毒、药疹样皮炎等。

\end{pastquestion}

\PointSeparator

\chapter{苯中毒}\label{users__fpb__26_spring__pkusph_study__tmp__pdfs__prevention-medicine-past-exams__07-ux804cux4e1aux75c5ux5b66ux5f80ux5e74ux9898ux603bux7ed3.md__ux7b2cux5341ux4e00ux7ae0-ux82efux4e2dux6bd2}

\section{苯中毒英译汉、急性/慢性临床表现与血液系统毒性}\label{users__fpb__26_spring__pkusph_study__tmp__pdfs__prevention-medicine-past-exams__07-ux804cux4e1aux75c5ux5b66ux5f80ux5e74ux9898ux603bux7ed3.md__ux82efux4e2dux6bd2ux82f1ux8bd1ux6c49ux6025ux6027ux6162ux6027ux4e34ux5e8aux8868ux73b0ux4e0eux8840ux6db2ux7cfbux7edfux6bd2ux6027}

\FrequencyPill{高频}

\InfoLabel{对应小节} 第十一章 / 临床表现

\ContentBand{知识点原文摘取}

急性苯中毒：苯属于有机溶剂，脂溶性强，对中枢神经系统有麻醉作用，吸入后可引起中枢神经系统先兴奋后抑制。

慢性苯中毒：主要为造血系统损伤表现，有明显量-效关系，也存在个体敏感性差异。造血系统为最具特征性表现，出现较早，为诊断慢性苯中毒的主要依据。白细胞减少最常见，也可出现血小板减少、贫血，严重者进展为再生障碍性贫血、白血病、骨髓增生异常综合征等。

\ContentBand{对应往年题}

\begin{pastquestion}

\QuestionLabel{【04级-职业病考题.docx｜大题】}\\
慢性苯中毒的临床表现，需要和哪些其他疾病鉴别。

\end{pastquestion}

\begin{pastquestion}

\QuestionLabel{【05级职业病考试题.doc｜大题】}\\
急慢性苯中毒的诊断要点及诊断标准。

\end{pastquestion}

\begin{pastquestion}

\QuestionLabel{【07级-职业病考题.pdf｜多选/单选】}\\
汞苯中毒的主要表现、分级表现（特征性点）；重症苯中毒诊断；轻度苯中毒诊断标准；能造成血液系统损害的有机溶剂。

\end{pastquestion}

\begin{pastquestion}

\QuestionLabel{【08级-职业病考题.doc｜英译汉/单选/多选】}\\
慢性苯中毒；引起血液系统损害的有机溶剂；慢性苯中毒主要损害的系统。

\end{pastquestion}

\begin{pastquestion}

\QuestionLabel{【14级职业病考题.pdf｜单选】}\\
下面毒物中没有血液系统毒性的是；主要影响血液系统的是。

\end{pastquestion}

\begin{pastquestion}

\QuestionLabel{【2015-2022职业病往年题及答案汇总.docx｜2022参考样题/翻译/单选8/多选1/多选9】}\\
benzene
poisoning：苯中毒；具有血液系统毒性的有机溶剂是苯；慢性苯中毒主要为中枢神经系统和造血系统表现；慢性重度苯中毒包括全血细胞减少症、再生障碍性贫血、骨髓增生异常综合征、白血病。

\end{pastquestion}

\begin{pastquestion}

\QuestionLabel{【2021级职业病考题yys.docx｜英译中】}\\
Occupational benzene poisoning。

\end{pastquestion}

\begin{pastquestion}

\QuestionLabel{【24年完整试题回忆.docx｜英译汉/单选8/9】}\\
occupational benzene poisoning；具有典型血液系统毒性的有机溶剂；慢性苯中毒以白细胞减少最为常见。

\end{pastquestion}

\PointSeparator

\section{职业性慢性苯中毒诊断原则、分级与鉴别诊断}\label{users__fpb__26_spring__pkusph_study__tmp__pdfs__prevention-medicine-past-exams__07-ux804cux4e1aux75c5ux5b66ux5f80ux5e74ux9898ux603bux7ed3.md__ux804cux4e1aux6027ux6162ux6027ux82efux4e2dux6bd2ux8bcaux65adux539fux5219ux5206ux7ea7ux4e0eux9274ux522bux8bcaux65ad}

\FrequencyPill{高频}

\InfoLabel{对应小节} 第十一章 / 诊断

\ContentBand{知识点原文摘取}

慢性职业性苯中毒诊断原则：具有明确的较长时期职业性苯接触史（三个月以上），现场劳动卫生学调查表明个人防护措施不良，工作人员有过量接触苯的机会；以造血系统损害为主的临床表现；参考实验室检测指标，如尿酚增加；排除其他原因引起的血象、骨髓象改变。

排除病因包括：骨髓造血障碍（电离辐射、药物、感染、骨髓恶性肿瘤、某些化学物质）；营养缺乏性造血障碍；白细胞破坏增加；分布异常性外周白细胞减少；良性家族性白细胞减少等。

慢性轻度苯中毒：白细胞大多数低于4×10\^{}9/L或中性粒细胞低于2×10\^{}9/L，或血小板大多数低于80×10\^{}9/L。慢性中度和重度进一步按血细胞减少程度及全血细胞减少、再障、MDS、白血病等分级。

\ContentBand{对应往年题}

\begin{pastquestion}

\QuestionLabel{【04级-职业病考题.docx｜大题】}\\
慢性苯中毒临床表现及其与其他主要疾病的鉴别。

\end{pastquestion}

\begin{pastquestion}

\QuestionLabel{【05级职业病考试题.doc｜大题】}\\
急慢性苯中毒的诊断要点及诊断标准。

\end{pastquestion}

\begin{pastquestion}

\QuestionLabel{【2015-2022职业病往年题及答案汇总.docx｜2020/简答】}\\
职业性慢性苯中毒诊断原则和鉴别诊断。

\end{pastquestion}

\begin{pastquestion}

\QuestionLabel{【2020年职业病考题.jpg｜简答3】}\\
简述职业性慢性苯中毒的诊断原则及鉴别的主要疾病。

\end{pastquestion}

\PointSeparator

\chapter{职业性噪声聋}\label{users__fpb__26_spring__pkusph_study__tmp__pdfs__prevention-medicine-past-exams__07-ux804cux4e1aux75c5ux5b66ux5f80ux5e74ux9898ux603bux7ed3.md__ux7b2cux5341ux4e8cux7ae0-ux804cux4e1aux6027ux566aux58f0ux804b}

\section{噪声聋概念、听力检查和耳聋分类}\label{users__fpb__26_spring__pkusph_study__tmp__pdfs__prevention-medicine-past-exams__07-ux804cux4e1aux75c5ux5b66ux5f80ux5e74ux9898ux603bux7ed3.md__ux566aux58f0ux804bux6982ux5ff5ux542cux529bux68c0ux67e5ux548cux8033ux804bux5206ux7c7b}

\FrequencyPill{中高频}

\InfoLabel{对应小节} 第十二章 / 噪声与听力；噪声聋机制；诊断原则

\ContentBand{知识点原文摘取}

职业性噪声聋：长期接触职业性噪声导致的感音神经性听力损失。耳聋分类：传导性聋（外耳、中耳），感音神经性聋（内耳、听神经、听中枢），混合性聋（中耳内耳均有病变）。纯音听阈测试可测气导和骨导听阈，是听力测试中最基本的检查。职业性噪声聋诊断步骤包括耳鼻喉科常规检查、至少3次纯音测听、年龄性别修正、声导抗除外中耳疾病；必要时进行客观检查。

\ContentBand{对应往年题}

\begin{pastquestion}

\QuestionLabel{【14级职业病考题.pdf｜英译汉/单选/多选】}\\
噪声聋；声导抗是检测哪部分病变；客观测听法包括。

\end{pastquestion}

\begin{pastquestion}

\QuestionLabel{【2015-2022职业病往年题及答案汇总.docx｜2018/翻译/单选/多选】}\\
noise-induced
deafness：噪声聋；声导抗是检测中耳病变；客观测听法包括声导抗、耳声发射、听性脑干反应、40Hz听觉相关电位。

\end{pastquestion}

\begin{pastquestion}

\QuestionLabel{【2020年职业病考题.jpg｜选择8】}\\
纯音测听中，哪种耳聋具有明显的骨导和气导差距：传导性耳聋。

\end{pastquestion}

\begin{pastquestion}

\QuestionLabel{【2021级职业病考题yys.docx｜英译中】}\\
Occupational noise-induced deafness。

\end{pastquestion}

\begin{pastquestion}

\QuestionLabel{【24年完整试题回忆.docx｜英译汉/单选5/多选4】}\\
occupational noise-induced deafness；纯音测听骨导听阈可以反映哪部分听力；客观测听法包括。

\end{pastquestion}

\PointSeparator

\section{噪声危害影响因素和职业性噪声聋诊断修正}\label{users__fpb__26_spring__pkusph_study__tmp__pdfs__prevention-medicine-past-exams__07-ux804cux4e1aux75c5ux5b66ux5f80ux5e74ux9898ux603bux7ed3.md__ux566aux58f0ux5371ux5bb3ux5f71ux54cdux56e0ux7d20ux548cux804cux4e1aux6027ux566aux58f0ux804bux8bcaux65adux4feeux6b63}

\FrequencyPill{中频}

\InfoLabel{对应小节} 第十二章 / 噪声性耳聋影响因素；诊断原则

\ContentBand{知识点原文摘取}

噪声性耳聋影响因素：噪声强度；噪声频率（2000-4000Hz最易导致耳蜗损害；窄带音或纯音较宽带声影响大；持续性噪声较断续噪声损伤性大；突然出现的噪声危害性更大；噪声伴振动较单纯噪声显著）；接触噪声时间；个体差异。年龄越大，噪声导致的内耳损伤越重。

职业性噪声聋诊断需要进行年龄、性别修正。

\ContentBand{对应往年题}

\begin{pastquestion}

\QuestionLabel{【14级职业病考题.pdf｜多选】}\\
关于影响噪声危害的因素的描述正确的是。

\end{pastquestion}

\begin{pastquestion}

\QuestionLabel{【2015-2022职业病往年题及答案汇总.docx｜2022参考样题/单选16】}\\
关于噪声导致的听力损伤，以下说法哪项不正确。

\end{pastquestion}

\begin{pastquestion}

\QuestionLabel{【24年完整试题回忆.docx｜多选5】}\\
职业性噪声聋诊断时需要根据年龄、性别等对纯音测听结果进行修正。

\end{pastquestion}

\PointSeparator

\chapter{慢性放射病}\label{users__fpb__26_spring__pkusph_study__tmp__pdfs__prevention-medicine-past-exams__07-ux804cux4e1aux75c5ux5b66ux5f80ux5e74ux9898ux603bux7ed3.md__ux7b2cux5341ux4e09ux7ae0-ux6162ux6027ux653eux5c04ux75c5}

\section{慢性放射病定义、病因和鉴别}\label{users__fpb__26_spring__pkusph_study__tmp__pdfs__prevention-medicine-past-exams__07-ux804cux4e1aux75c5ux5b66ux5f80ux5e74ux9898ux603bux7ed3.md__ux6162ux6027ux653eux5c04ux75c5ux5b9aux4e49ux75c5ux56e0ux548cux9274ux522b}

\FrequencyPill{新增中频}

\InfoLabel{对应小节} 第十三章 / 概述；诊断标准及鉴别诊断

\ContentBand{知识点原文摘取}

慢性放射病：人体在较长时间（数月-数年）内连续或间断受到超剂量限值（\textgreater20mSv/y）的电离辐射作用，达到一定累积剂量（\textgreater1.5Gy）后引起的多系统损害的全身性疾病，常以造血组织损伤为主。

病因包括放射诊断、放射治疗、工业探伤、加速器、反应堆等。慢性外照射放射病由X、γ射线及中子流体外照射引起；慢性内照射放射病由放射性核素进入体内引起。

鉴别诊断：神经衰弱症候群需与一般神经衰弱症、更年期综合征鉴别；造血组织改变需与慢性苯中毒、再生障碍性贫血等鉴别。

\ContentBand{对应往年题}

\begin{pastquestion}

\QuestionLabel{【2021级职业病考题yys.docx｜填空】}\\
内照射型放射病的原因，早期症状。

\end{pastquestion}

\begin{pastquestion}

\QuestionLabel{【24年完整试题回忆.docx｜单选15/简答3】}\\
下列哪项不为慢性放射病的病因；简述慢性放射病的定义。

\end{pastquestion}

\PointSeparator

\chapter{急性放射病}\label{users__fpb__26_spring__pkusph_study__tmp__pdfs__prevention-medicine-past-exams__07-ux804cux4e1aux75c5ux5b66ux5f80ux5e74ux9898ux603bux7ed3.md__ux7b2cux5341ux56dbux7ae0-ux6025ux6027ux653eux5c04ux75c5}

\section{电离辐射、外照射急性放射病分型和骨髓型分期}\label{users__fpb__26_spring__pkusph_study__tmp__pdfs__prevention-medicine-past-exams__07-ux804cux4e1aux75c5ux5b66ux5f80ux5e74ux9898ux603bux7ed3.md__ux7535ux79bbux8f90ux5c04ux5916ux7167ux5c04ux6025ux6027ux653eux5c04ux75c5ux5206ux578bux548cux9aa8ux9ad3ux578bux5206ux671f}

\FrequencyPill{新增中频}

\InfoLabel{对应小节} 第十四章 / 概述；临床表现

\ContentBand{知识点原文摘取}

外照射急性放射病：机体在短时间内（数秒或数日内）受到1Gy以上体外电离辐射引起的全身性疾病。电离辐射包括电磁辐射和粒子辐射。

分型：骨髓型1-10Gy（轻度、中度、重度、极重度）；肠型10-50Gy；脑型\textgreater50Gy。骨髓型急性放射病以骨髓造血组织损伤为基本病变，以白细胞数减少、感染、出血等为主要临床表现，具有典型阶段性病程：初期、假愈期、极期、恢复期。

\textbf{未完全匹配到：} 24年多选``辐射效应分类方式''在知识点总结中未展开完整分类，暂归入本章。

\ContentBand{对应往年题}

\begin{pastquestion}

\QuestionLabel{【2021级职业病考题yys.docx｜多选】}\\
电离辐射的种类：X线、β射线、γ射线、中子等。

\end{pastquestion}

\begin{pastquestion}

\QuestionLabel{【24年完整试题回忆.docx｜单选16/多选9/填空6】}\\
外照射急性放射病的分型不包括；辐射效应分类方式；骨髓型急性放射病的临床分期。

\end{pastquestion}

\begin{pastquestion}

\QuestionLabel{【24年完整试题回忆.docx｜英译汉】}\\
latent phase（假愈期/潜伏期）。

\end{pastquestion}

\PointSeparator

\chapter*{题源反向核对索引}\label{users__fpb__26_spring__pkusph_study__tmp__pdfs__prevention-medicine-past-exams__07-ux804cux4e1aux75c5ux5b66ux5f80ux5e74ux9898ux603bux7ed3.md__ux9898ux6e90ux53cdux5411ux6838ux5bf9ux7d22ux5f15}
\addcontentsline{toc}{chapter}{题源反向核对索引}

\begin{quote}
用于检查``往年题绝对不能遗漏''。以下按文件列出可见题目及其归入考点。
\end{quote}

\section{04级-职业病考题.docx}\label{users__fpb__26_spring__pkusph_study__tmp__pdfs__prevention-medicine-past-exams__07-ux804cux4e1aux75c5ux5b66ux5f80ux5e74ux9898ux603bux7ed3.md__ux7ea7-ux804cux4e1aux75c5ux8003ux9898.docx}

\begin{itemize}
\tightlist
\item
  尘肺病理分级及X线表现 → 第二章3
\item
  职业病分类/目录/定义，每类举例 → 第一章1
\item
  有机溶剂定义、用途，正己烷和三氯乙烯毒性 → 第十章1-3
\item
  慢性苯中毒临床表现及鉴别 → 第十一章1-2
\item
  有机磷中毒并发症、临床表现和治疗 → 第四章2
\item
  急慢性汞中毒临床表现 → 第五章2
\item
  急性CO中毒诊断标准和治疗原则 → 第九章2
\item
  填空：汞接触行业 → 第五章1
\item
  填空：有机磷中毒定义 → 第四章2
\item
  填空：砷中毒皮肤表现 → 第六章1
\item
  急慢性铅中毒治疗原则 → 第三章2
\item
  英文病名5个（原文未列具体词） → 按相应疾病概念保留
\end{itemize}

\section{05级职业病考试题.doc}\label{users__fpb__26_spring__pkusph_study__tmp__pdfs__prevention-medicine-past-exams__07-ux804cux4e1aux75c5ux5b66ux5f80ux5e74ux9898ux603bux7ed3.md__ux7ea7ux804cux4e1aux75c5ux8003ux8bd5ux9898.doc}

\begin{itemize}
\tightlist
\item
  英译汉：尘肺、刺激性气体、慢性汞中毒、有机磷农药、一氧化碳 → 第二章/第八章/第五章/第四章/第九章
\item
  职业病定义及分类目录，每种列举2个疾病 → 第一章1
\item
  影响尘肺发病及进展因素 → 第二章2
\item
  CO中毒诊断要点 → 第九章2
\item
  有机溶剂定义、分类、共同毒性 → 第十章1
\item
  急慢性苯中毒诊断要点及诊断标准 → 第十一章1-2
\item
  有机磷中毒发病机理及主要临床表现 → 第四章2
\item
  急慢性汞中毒鉴别诊断、治疗方法 → 第五章2
\item
  慢性铅中毒治疗、预防 → 第三章2
\end{itemize}

\section{07级-职业病考题.pdf}\label{users__fpb__26_spring__pkusph_study__tmp__pdfs__prevention-medicine-past-exams__07-ux804cux4e1aux75c5ux5b66ux5f80ux5e74ux9898ux603bux7ed3.md__ux7ea7-ux804cux4e1aux75c5ux8003ux9898.pdf}

\begin{itemize}
\tightlist
\item
  大题：职业病目录各举三例 → 第一章1
\item
  有机磷迟发型周围神经病发生时期及临床表现 → 第四章2
\item
  窒息性气体分三类各举一例 → 第九章1
\item
  尘肺发病及进展影响因素 → 第二章2
\item
  矽肺并发症及鉴别诊断 → 第二章4
\item
  填空：健康监测内容 → 第一章2
\item
  填空：物理因素职业病 → 第一章3
\item
  填空：化学性肺水肿分期 → 第八章2
\item
  填空：乙酰胆碱酯酶活性下降的农药 → 第四章2
\item
  填空：农药一般毒性作用 → 第四章1
\item
  填空：汞中度中毒标准 → 第五章2
\item
  多选：铅接触行业、驱铅药物 → 第三章1-2
\item
  多选：石棉肺X线表现、尘肺病理类型 → 第二章3/5
\item
  多选：汞/苯中毒主要表现和分级 → 第五章2/第十一章2
\item
  多选：刺激性气体、有机磷早期表现 → 第八章1/第四章2
\item
  多选：重症苯中毒诊断 → 第十一章2
\item
  单选：尿铅值诊断、慢性汞中毒肾损害部位、氰化氢排出途径、轻度苯中毒诊断标准、慢性正己烷周围神经病、血液系统损害有机溶剂、非刺激性气体、驱汞制剂、迟发周围神经/中间期肌无力发生时间
  → 对应各章
\item
  翻译：有机溶剂、矽肺、砷中毒、苯中毒、刺激性气体 → 对应各章
\end{itemize}

\section{08级-职业病考题.doc}\label{users__fpb__26_spring__pkusph_study__tmp__pdfs__prevention-medicine-past-exams__07-ux804cux4e1aux75c5ux5b66ux5f80ux5e74ux9898ux603bux7ed3.md__ux7ea7-ux804cux4e1aux75c5ux8003ux9898.doc}

\begin{itemize}
\tightlist
\item
  英译汉：石棉肺、有机磷农药中毒、铅中毒、窒息性气体、慢性苯中毒 → 对应各章
\item
  填空：胸膜斑定义、法定职业病定义、职业病类种、有机溶剂定义/共同毒性、汞毒性形式、石棉肺病理改变 →
  第一章/第二章/第十章/第五章
\item
  单选：汞蓄积部位、职业病诊断、百草枯损害、闭塞性细支气管炎刺激性气体、CO亲和力、铅损害系统、血液系统损害有机溶剂、正己烷途径、尿汞诊断值、砷化合物毒性最强、氰化氢排毒、不属于急性铅中毒、粉尘作业禁忌证、有机溶剂错误项
  → 对应各章
\item
  多选：职业病诊断、职业相关病概念、可引起中毒性脑病毒物、慢性苯中毒损害系统、刺激性气体并发症、三氯乙烯职业损害、急性有机磷主要症状、硅酸盐尘肺、煤工尘肺、口服急性砷中毒表现
  → 对应各章
\item
  简答：神经毒物、重度有机磷并发症、窒息性气体定义分类举例、石棉肺X线表现 → 第一章4/第四章2/第九章1/第二章5
\item
  病例：铅中毒诊断、依据、检查、鉴别、治疗 → 第三章1-2
\end{itemize}

\section{14级职业病考题.pdf}\label{users__fpb__26_spring__pkusph_study__tmp__pdfs__prevention-medicine-past-exams__07-ux804cux4e1aux75c5ux5b66ux5f80ux5e74ux9898ux603bux7ed3.md__ux7ea7ux804cux4e1aux75c5ux8003ux9898.pdf}

\begin{itemize}
\tightlist
\item
  英译汉：法定职业病、刺激性气体、尘肺病、噪声聋、有机溶剂中毒 → 对应各章
\item
  单选：血液系统毒性、正己烷近期接触指标、驱汞首选、急性汞表现、声导抗、中毒诊断值、铅早期表现、有机氟治疗、氰化氢排出、砷化氢机制、百草枯、正己烷慢性毒性、有机溶剂错误项
  → 对应各章
\item
  填空：职业病类种、铅吸收排出、胸膜斑、尘肺病理类型、化学性肺水肿四期、慢性汞三大表现 → 对应各章
\item
  多选：客观测听法、三氯乙烯临床表现、铅中毒影响、窒息性气体、慢性中度汞、硅酸盐尘肺、石棉肺病理改变、刺激性气体并发症、噪声危害因素
  → 对应各章
\item
  简答：职业病分类10类、尘肺影响因素、窒息性气体定义分类举例 → 第一章1/第二章2/第九章1
\item
  病例：尘肺病病例、农药中毒病例 → 第二章3/第四章2
\end{itemize}

\section{2015年职业病考试题型.docx}\label{users__fpb__26_spring__pkusph_study__tmp__pdfs__prevention-medicine-past-exams__07-ux804cux4e1aux75c5ux5b66ux5f80ux5e74ux9898ux603bux7ed3.md__ux5e74ux804cux4e1aux75c5ux8003ux8bd5ux9898ux578b.docx}

\begin{itemize}
\tightlist
\item
  题型：中英互译、填空、选择、简答、病例 → 已纳入说明
\item
  明确例题：窒息性气体分几类、尘肺分类/分期、粉尘危害决定因素、病例答诊断依据/分期/原则/治疗 →
  第九章1/第二章1-3
\end{itemize}

\section{2015-2022职业病往年题及答案汇总by
18级杨子铭.docx}\label{users__fpb__26_spring__pkusph_study__tmp__pdfs__prevention-medicine-past-exams__07-ux804cux4e1aux75c5ux5b66ux5f80ux5e74ux9898ux603bux7ed3.md__ux804cux4e1aux75c5ux5f80ux5e74ux9898ux53caux7b54ux6848ux6c47ux603bby-18ux7ea7ux6768ux5b50ux94ed.docx}

\begin{itemize}
\tightlist
\item
  2022参考样题全部题目 → 已在各章与独立2022样题共同归入
\item
  2015填空：窒息性气体分类；尘肺分类 → 第九章1/第二章1
\item
  2018翻译：法定职业病、尘肺病、噪声聋、刺激性气体、有机溶剂中毒 → 对应各章
\item
  2018单选：驱汞、砷化氢机制、正己烷慢性毒性、2,5-己二酮、声导抗 → 第五章/第六章/第十章/第十二章
\item
  2018多选：客观测听法、慢性中度汞中毒 → 第十二章/第五章
\item
  2018填空：铅吸收排出、尘肺病理类型、慢性汞三大表现 → 第三章/第二章/第五章
\item
  2018简答：职业病10类、窒息性气体定义分类举例 → 第一章/第九章
\item
  2018病例：尘肺、有机磷中毒 → 第二章/第四章
\item
  2020选择：有机溶剂错误项、尿铅诊断值、纯音测听传导性耳聋、砷化氢溶血性毒物 → 第十章/第三章/第十二章/第六章
\item
  2020填空：10类132种、化学性肺水肿四期、硅酸盐肺、三氯乙烯两种职业病、急性中度汞 →
  第一章/第八章/第二章/第十章/第五章
\item
  2020简答：有机磷迟发周围神经病、慢性苯中毒诊断原则与鉴别 → 第四章/第十一章
\item
  2020病例：煤工尘肺壹期、急性有机磷中毒 → 第二章/第四章
\end{itemize}

\section{2020年职业病考题.jpg}\label{users__fpb__26_spring__pkusph_study__tmp__pdfs__prevention-medicine-past-exams__07-ux804cux4e1aux75c5ux5b66ux5f80ux5e74ux9898ux603bux7ed3.md__ux5e74ux804cux4e1aux75c5ux8003ux9898.jpg}

\begin{itemize}
\tightlist
\item
  可见选择4：慢性锰中毒吸收代谢途径 → 第七章
\item
  选择5：有机溶剂错误项 → 第十章1
\item
  选择6：职业性慢性铅中毒尿铅诊断值 → 第三章1
\item
  选择7：不是刺激性气体 → 第八章1
\item
  选择8：纯音测听骨导气导差距明显的耳聋 → 第十二章1
\item
  选择9：石棉肺常见病变不包括 → 第二章5
\item
  选择10：砷化氢主要毒性分类 → 第六章2
\item
  填空1：职业病分类和目录类种 → 第一章1
\item
  填空2：有机磷农药中毒机制和三组症状 → 第四章2
\item
  填空3：化学性肺水肿四期 → 第八章2
\item
  填空4：胸膜斑定义 → 第二章5
\item
  填空5：硅酸盐肺相关填空 → 第二章5
\item
  填空6：三氯乙烯两种职业病 → 第十章3
\item
  填空7：急性中度汞中毒标准 → 第五章2
\item
  简答1：急性有机磷中毒迟发性周围神经病时间及临床表现 → 第四章2
\item
  简答2：矽肺鉴别诊断和常见并发症 → 第二章4
\item
  简答3：慢性苯中毒诊断原则和鉴别疾病 → 第十一章2
\item
  病例1：煤工尘肺/尘肺病例 → 第二章3
\item
  病例2：有机磷农药中毒病例 → 第四章2
\end{itemize}

\section{2022职业病参考样题及答案整理by18级杨子铭
丁滢方.docx}\label{users__fpb__26_spring__pkusph_study__tmp__pdfs__prevention-medicine-past-exams__07-ux804cux4e1aux75c5ux5b66ux5f80ux5e74ux9898ux603bux7ed3.md__ux804cux4e1aux75c5ux53c2ux8003ux6837ux9898ux53caux7b54ux6848ux6574ux7406by18ux7ea7ux6768ux5b50ux94ed-ux4e01ux6ee2ux65b9.docx}

\begin{itemize}
\tightlist
\item
  翻译：silicosis、organic solvents、benzene poisoning、噪声聋、砷中毒 → 对应各章
\item
  单选1-18：锰、砷、正己烷、有机磷迟发神经病/中间期肌无力/胆碱酯酶、有机溶剂/苯、铅诊断值、汞、氨基甲酸酯、粉尘禁忌证、CO、氰化物、噪声、窒息性气体、尘肺X线
  → 对应各章
\item
  多选1-12：苯、石棉肺、汞、三氯乙烯、有机磷、刺激性气体、驱铅、铅行业、重度苯、尘肺纤维化、刺激性气体并发症、慢性砷
  → 对应各章
\item
  填空1-6：农药共同毒性、胆碱酯酶下降农药、物理因素职业病、职业健康监护、胸膜斑、铅吸收 → 对应各章
\item
  问答1-3：职业病目录10类、尘肺影响因素、矽肺并发症和鉴别 → 第一章/第二章
\item
  病例：煤工尘肺、慢性轻度铅中毒、慢性轻度汞中毒 → 第二章/第三章/第五章
\end{itemize}

\section{2021级职业病考题yys.docx}\label{users__fpb__26_spring__pkusph_study__tmp__pdfs__prevention-medicine-past-exams__07-ux804cux4e1aux75c5ux5b66ux5f80ux5e74ux9898ux603bux7ed3.md__ux7ea7ux804cux4e1aux75c5ux8003ux9898yys.docx}

\begin{itemize}
\tightlist
\item
  英译中：Occupational benzene poisoning、Asphyxiating gases、Acute respiratory distress
  syndrome、Agricultural chemicals、Silicosis → 对应各章
\item
  单选：最常见造成血液系统损害的金属 → 第三章/第十一章
\item
  填空：新版职业病目录类种及新增类别；内照射型放射病原因、早期症状 → 第一章1/第十三章
\item
  多选：电离辐射种类；铅损伤系统 → 第十四章/第三章
\item
  论述：农药定义及共同毒性；窒息性气体定义分类举例；影响尘肺病发病因素 → 第四章/第九章/第二章
\item
  病例：有机磷中毒；煤工尘肺/尘肺病例 → 第四章/第二章
\end{itemize}

\section{24年完整试题回忆.docx}\label{users__fpb__26_spring__pkusph_study__tmp__pdfs__prevention-medicine-past-exams__07-ux804cux4e1aux75c5ux5b66ux5f80ux5e74ux9898ux603bux7ed3.md__ux5e74ux5b8cux6574ux8bd5ux9898ux56deux5fc6.docx}

\begin{itemize}
\tightlist
\item
  英译汉：occupational noise-induced deafness、silicosis、latent phase、occupational benzene
  poisoning、agricultural chemicals → 第十二章/第二章/第十四章/第十一章/第四章
\item
  单选1-20：急性汞、正己烷指标、砷皮肤表现、砷化氢、纯音测听、氨基甲酸酯、百草枯、苯、慢性苯、驱铅、亚甲蓝、单纯性窒息气体、COHb、氰化物排除、慢性放射病病因、急性放射病分型、尘肺病、胸膜斑、粉尘禁忌证、尿铅诊断值
  → 对应各章
\item
  填空1-7：铅吸收、慢性汞三大表现、有机溶剂共同毒性、2013职业病类种、耳聋分类、骨髓型急性放射病分期、化学性肺水肿分期
  → 对应各章
\item
  多选1-10：职业健康检查、中毒性脑病毒物、三氯乙烯、客观测听、噪声聋修正、铅中度中毒、职业健康检查目标、刺激性气体、辐射效应分类、石棉肺胸膜表现
  → 对应各章
\item
  简答：窒息性气体定义分类举例；物理因素所致职业病；慢性放射病定义 → 第九章/第一章/第十三章
\item
  病例：煤工尘肺/尘肺；铅中毒 → 第二章/第三章
\end{itemize}

\end{document}
